\documentclass[oneside]{amsart}
\PassOptionsToPackage{hyperfootnotes=false}{hyperref}

% Fixed PDF identifier.
\ifdefined\pdftrailerid\pdftrailerid{}\fi

\usepackage{iftex}
\ifluatex\usepackage{luatex85}\fi
\ifdefined\pdftrailerid\pdftrailerid{}\fi
% Set-up of this edition: packages, statement environments, headings and
% contents, notation, sheaf operators, page marks and cross-references.
% Written for these editions and dedicated to the public domain under CC0 1.0.
% Commutative diagrams.
\usepackage[all]{xy}
\usepackage{tikz-cd}
% OpenType text font.
\usepackage{fontspec}
\setmainfont{TeX Gyre Pagella}
% String handling and argument parsing for the reference commands below;
% etoolbox; colour.
\usepackage{xstring,xparse,etoolbox,xcolor}
% Coloured links: cross-references in deep red, citations in blue; contents
% entries linked by their page numbers.
\definecolor{egalinkred}{rgb}{0.76,0.13,0.28}
\usepackage[
  linktocpage=true,
  colorlinks=true,
  hyperindex,
  hypertexnames=false,
  citecolor=blue,
  linkcolor=egalinkred
]{hyperref}
% Margins of 1.25in at the sides and 0.75in at head and foot.
\usepackage[hmargin=1.25in,vmargin=0.75in]{geometry}
% Palatino for text and mathematics, AMS symbols and Ralph Smith's script
% alphabet; Palatino is set with its customary slightly opened leading.
\usepackage{amssymb,mathrsfs,mathpazo}\linespread{1.05}
% Headings are not left at the foot of a page; margin notes; lists, long
% tables and multiple columns; the Aurical hand (for sheaf operators);
% microtypography.
\usepackage[nobottomtitles]{titlesec}
\usepackage{marginnote}
\usepackage[shortlabels]{enumitem}
\usepackage{longtable,multicol,aurical,microtype}

% Statements. EGA numbers its statements in the text, so the environments
% are unnumbered and take the number as optional argument. For instance
% \begin{proposition}[2.1.3] starts with the bold heading "Proposition
% (2.1.3). ---" (in the language of the edition) and sets the statement in
% italics; definitions, examples, remarks and notation are set in roman;
% \begin{env}[2.1.4] starts a numbered paragraph with a bold "(2.1.4).".
\newtheoremstyle{egaitem}{}{}{\rmfamily}{}{\bfseries}{.}{ }{\thmnote{(#3)}}
\newtheoremstyle{egastatement}{}{}{\itshape}{}{\bfseries}{. --- }{ }{\thmname{#1}\thmnote{ (#3)}}
\newtheoremstyle{egadescription}{}{}{\rmfamily}{}{\bfseries}{. --- }{ }{\thmname{#1}\thmnote{ (#3)}}
\makeatletter
\newcommand*{\ega@statement}[3]{\theoremstyle{#1}\newtheorem*{#2}{#3}}
\ega@statement{egaitem}{env}{---}
\ega@statement{egastatement}{theorem}{Định lý}
\ega@statement{egastatement}{proposition}{Mệnh đề}
\ega@statement{egastatement}{lemma}{Bổ đề}
\ega@statement{egastatement}{corollary}{Hệ quả}
\ega@statement{egadescription}{definition}{Định nghĩa}
\ega@statement{egadescription}{example}{Ví dụ}
\ega@statement{egadescription}{examples}{Các ví dụ}
\ega@statement{egadescription}{remark}{Nhận xét}
\ega@statement{egadescription}{remarks}{Các nhận xét}
\ega@statement{egadescription}{scholium}{Chú giải}
\ega@statement{egadescription}{notation}{Ký hiệu}
\makeatother
\renewcommand{\proofname}{Chứng minh}

% Contents and headings. Subsection entries of the contents are indented by
% 2.5pc, with 2.2pc for the number. Section headings are amsart's centred small
% capitals, in bold at \Large size. Heading numbers end with a full stop and a
% tie, and a section number is preceded by the sign \S.
\makeatletter
\renewcommand*{\l@subsection}{\@tocline{2}{0pt}{2.5pc}{2.2pc}{}}
\renewcommand*{\section}{%
  \@startsection{section}{1}{\z@}{.7\linespacing\@plus\linespacing}%
  {.5\linespacing}{\normalfont\bfseries\Large\scshape\centering}}
\@namedef{ega@numbersign@section}{\textsection}
\newcommand*{\ega@numbersign}[1]{%
  \ifcsname ega@numbersign@#1\endcsname\csname ega@numbersign@#1\endcsname\fi}
\renewcommand*{\@seccntformat}[1]{\ega@numbersign{#1}\csname the#1\endcsname.~}
\makeatother

% Sheaves are written in the script alphabet of mathrsfs.
\renewcommand*{\mathcal}{\mathscr}
% BEGIN INLINE PREAMBLE-BASE
% Notation of the EGA editions: the abbreviations the transcription uses for alphabets, symbols, operators and
% arrows. Written for these editions and dedicated to the public domain under CC0 1.0.
\makeatletter
% a command that may or may not exist already is (re)defined in one step
\newcommand{\ega@define}[2]{\providecommand{#1}{}\renewcommand{#1}{#2}}
\newcommand{\ega@operator}[2]{\providecommand{#1}{}\renewcommand{#1}{\operatorname{#2}}}

% alphabets: sheaves, bold letters, categories
\ega@define{\sh}{\mathcal}
\ega@define{\bb}{\mathbf}
\ega@define{\cat}{\mathtt}
\ega@define{\C}{\cat{C}}
\ega@define{\Set}{\cat{Set}}
\ega@define{\CHom}{\cat{Hom}}
\ega@define{\op}{^{\cat{op}}}
\ega@define{\GL}{\bb{GL}}
\ega@define{\bbGamma}{\boldsymbol{\Gamma}}

% letters and symbols
\renewcommand{\leq}{\leqslant}
\renewcommand{\geq}{\geqslant}
\ega@define{\setmin}{-}
\ega@define{\emp}{\varnothing}
\ega@define{\vphi}{\phi}
\ega@define{\kres}{k}
\ega@define{\rad}{\mathfrak{r}}
\ega@define{\nilrad}{\mathfrak{N}}
\ega@define{\HH}{\mathrm{H}}
\ega@define{\CHH}{\check{\mathrm{H}}}
\ega@define{\RR}{\mathrm{R}}
\ega@define{\LL}{\mathrm{L}}
\ega@define{\red}{\mathrm{red}}
\providecommand{\dual}{}\renewcommand{\dual}[1]{{#1}^{\vee}}
% the tilde of the sheaf associated with a module, written after it: M\supertilde
\ega@define{\supertilde}{{\,\widetilde{\,}\,}}

% operators
\ega@operator{\rg}{rg}
\ega@operator{\gr}{gr}
\ega@operator{\Gr}{Gr}
\ega@operator{\Hom}{Hom}
\ega@operator{\Proj}{Proj}
\ega@operator{\Tor}{Tor}
\ega@operator{\Ext}{Ext}
\ega@operator{\Supp}{Supp}
\ega@operator{\Ker}{Ker}
\ega@operator{\Im}{Im}
\ega@operator{\Coker}{Coker}
\ega@operator{\Spec}{Spec}
\ega@operator{\Spm}{Spm}
\ega@operator{\Spf}{Spf}
\ega@operator{\grad}{grad}
\ega@operator{\dimc}{dimc}
\ega@operator{\codim}{codim}
\ega@operator{\id}{id}
\ega@operator{\Der}{Der}
\ega@operator{\Diff}{Diff}
\ega@operator{\Hyp}{Hyp}
\ega@operator{\Tr}{Tr}

% arrows: long in displays, short in text, scripts and subscripts
\newcommand{\ega@arrow}[2]{\mathchoice{#1}{#2}{#2}{#2}}
\let\ega@shortmapsto\mapsto
\renewcommand{\to}{\ega@arrow{\longrightarrow}{\rightarrow}}
\providecommand{\from}{}\renewcommand{\from}{\ega@arrow{\longleftarrow}{\leftarrow}}
\renewcommand{\mapsto}{\ega@arrow{\longmapsto}{\ega@shortmapsto}}
\ega@define{\isoto}{\simeq}
\ega@define{\simto}{\xrightarrow{\sim}}

% wide accents
\newcommand{\ega@wide}[2]{\renewcommand{#1}[1]{#2{##1}}}
\ega@wide{\hat}{\widehat}
\ega@wide{\tilde}{\widetilde}
\makeatother
% END INLINE PREAMBLE-BASE

% Sheaf operators (\shHom, \shExt, ...): a script capital followed by the
% rest of the name in slanted Aurical, moved 2.2pt towards the capital. The
% middle argument of \ega@sheafop is the space put after the capital.
\makeatletter
\newcommand*{\ega@sheafop}[3]{\sh{#1}#2\textup{\kern-2.2pt{\Fontauri\slshape #3}}}
\newcommand*{\shHom}{\ega@sheafop{H}{}{om}}
\newcommand*{\shProj}{\ega@sheafop{P}{}{roj}}
\newcommand*{\shExt}{\ega@sheafop{E}{}{xt}}
\newcommand*{\shTor}{\ega@sheafop{T}{}{or}}
\newcommand*{\shGr}{\ega@sheafop{G}{}{r}}
\newcommand*{\shDer}{\ega@sheafop{D}{\,}{er}}
\newcommand*{\shDiff}{\ega@sheafop{D}{\,}{if{}f}\,}
\newcommand*{\shHomcont}{\ega@sheafop{H}{}{om.\,cont}}
\newcommand*{\shAut}{\ega@sheafop{A}{}{ut}}
\makeatother

% Doubtful passages of the translation: \unsure{...} adds a red "(?)" after its
% argument, and \completelyunsure stands for a passage left open, "(???)".
\newcommand*{\unsure}[1]{#1 {\color{red}(?)}}
\newcommand*{\completelyunsure}{{\color{red}(???)}}

% Pages of the printed original, noted in the margin as "volume | page":
% \oldpage[III]{81}. \sectionbreak is a centred row of asterisks.
\newcommand{\oldpage}[2][]{{\marginnote{\normalfont\textbf{#1}~|~#2}}\ignorespaces}
\newcommand*{\sectionbreak}{\begin{center}***\end{center}}

% Give every source label its own zero-size PDF destination.  The inherited
% unnumbered theorem environments otherwise leave many labels pointing to the
% preceding subsection or display anchor rather than to their exact statement.
\let\EGAOriginalLabel\label
\renewcommand{\label}[1]{\phantomsection\EGAOriginalLabel{#1}}

% Cross-references. Labels have the form <chapter>.<number> and only the number
% is printed: \sref{0.4.1.2} gives "(4.1.2)". An optional first argument adds
% a volume in bold, \sref[I]{I.2.3.1} giving "(I, 2.3.1)"; an optional last
% argument adds a further indication, \sref{0.4.1.2}[ii] giving "(4.1.2, ii)".
% \eref{eq:0.4.1.2} refers in the same way to a numbered display. Corrections
% taken from published errata are marked \erratum[volume]{corrected text}, with
% "Err" and the volume in the margin.
\makeatletter
\newcommand{\EGAUnavailableHyperref}[2]{%
  \begingroup
  \let\@latex@warning\@gobble
  \let\@latex@warning@no@line\@gobble
  \let\G@refundefinedtrue\relax
  \hyperref[#1]{#2}%
  \endgroup
}
\newcommand*{\ega@labelnumber}[1]{\StrBehind{#1}{.}}
\newcommand*{\ega@reftext}[3]{%
  \normalfont{(\IfValueT{#1}{\textbf{#1},~}\ega@labelnumber{#2}\IfValueT{#3}{,~{#3}})}}
\NewDocumentCommand{\sref}{o m o}{\@ifundefined{r@#2}%
    {\EGAUnavailableHyperref{#2}{\ega@reftext{#1}{#2}{#3}}}%
    {\hyperref[#2]{\ega@reftext{#1}{#2}{#3}}}}
\newcommand{\eref}[1]{\hyperref[#1]{\normalfont{(\ega@labelnumber{#1})}}}
\newcommand*{\ega@errmark}[1]{\textbf{Err}\textsubscript{#1}}
\newcommand{\erratum}[2][]{{#2}\marginpar{\footnotesize\ega@errmark{#1}}}
\makeatother

\hypersetup{
  pdftitle={Các yếu tố của hình học đại số I -- Ngôn ngữ của các lược đồ},
  pdfauthor={Alexander Grothendieck và Jean Dieudonné; bản dịch tiếng Việt độc lập},
  pdfsubject={Bản tiếng Việt độc lập của EGA I}
}
\renewcommand{\contentsname}{Mục lục}
\renewcommand{\refname}{Tài liệu tham khảo}
\begin{document}
\title{Các yếu tố của hình học đại số\\EGA I: Ngôn ngữ của các lược đồ}
\author{Alexander Grothendieck và Jean Dieudonné}
\date{}
\maketitle
\phantomsection
\label{reader:ega1-vi}
\begin{center}
\large Ấn bản tiếng Việt hiệu chỉnh
\end{center}
\begin{quote}\small
Chương I, §§1–10 và các bảng tra; trang in 79–224. Phần kiến thức chuẩn bị ở Chương 0 không thuộc tài liệu này.

Tác phẩm gốc của Alexandre Grothendieck, với sự cộng tác của Jean Dieudonné: EGA I, Publications Mathématiques de l’IHÉS 4 (1960), \href{https://numdam.org/item/PMIHES\_1960\_\_4\_\_5\_0/}{NUMDAM}. Bản tham khảo tiếng Anh và cấu trúc TeX có nguồn gốc từ bản dịch do con người thực hiện tại \href{https://github.com/ryankeleti/ega/tree/3d7afe28c99d1887ccdb85acf2ca1afd93fb1971}{ryankeleti/ega}. Các cộng tác viên được ghi tên gồm Mateo Carmona, Tim Holzschuh, Tim Hosgood, Ryan Keleti, Thiago Solovera \& Nery, Wayne Yang. Phiên bản 3d7afe28c99d1887ccdb85acf2ca1afd93fb1971 xác định bản tiếng Anh đã dùng để đối chiếu; đóng góp cụ thể được nêu trong \href{https://github.com/ryankeleti/ega/blob/3d7afe28c99d1887ccdb85acf2ca1afd93fb1971/README.md}{ghi công theo từng mục}. Bản dịch tiếng Việt, các hiệu chỉnh và dàn trang do AI thực hiện, không được gán cho những người dịch tiếng Anh. Ghi công không đồng nghĩa với việc được cho phép hay được bảo chứng.

Bản dịch tiếng Việt, biên tập và dàn trang do AI thực hiện: OpenAI Codex; GPT-6 Sol và GPT-6.1 Sol đối với văn bản kế thừa, GPT-6 Astra đối với các hiệu chỉnh được đưa vào đây, mức suy luận Ultra. Không tuyên bố đã được con người rà soát.
\end{quote}
\begingroup\small
Quyền đối với các tác phẩm gốc và mọi đóng góp được bảo hộ trong các bản dịch và bản chép lại phái sinh này vẫn thuộc về các chủ thể quyền tương ứng. Dự án không đưa những tác phẩm hay đóng góp đó vào phạm vi công cộng và không cấp giấy phép đối với chúng. Các thông báo quyền của bên thứ ba và các sự cho phép đã có vẫn có hiệu lực.

Các ấn bản được cung cấp miễn phí cho mục đích phi thương mại, thuần túy giáo dục, trên giả định rằng các chủ thể quyền tương ứng sẽ không phản đối việc các bản dịch và bản chép lại được cung cấp miễn phí bên cạnh các bản gốc tiếng Pháp trên \href{https://www.numdam.org/}{NUMDAM}. Đây là một giả định, không phải tuyên bố đã được cho phép; việc truy cập miễn phí bản gốc tự nó không cho phép dịch, chuyển thể hoặc tái xuất bản.

Nếu bạn nắm giữ quyền đối với tài liệu ở đây và phản đối việc cung cấp tài liệu đó, hãy \href{https://github.com/KokunoYumeto/ega-fga-sga-vi/issues}{mở một issue trên GitHub}. Người xuất bản cam kết gỡ bỏ tài liệu bị phản đối theo yêu cầu của chủ thể quyền.

\par\endgroup\clearpage
\tableofcontents

\clearpage
\oldpage[I]{79}
\phantomsection
\label{chapter:I-vi}

\begin{center}
{\large CHƯƠNG THỨ NHẤT\par}
\vspace{1.5em}
{\Large\bfseries NGÔN NGỮ CỦA CÁC LƯỢC ĐỒ\par}
\vspace{1.8em}
{\bfseries Mục lục\par}
\end{center}

\medskip
\begin{tabular}{@{}r@{\ }l@{}}
\S~1.  & Các lược đồ affine.\\
\S~2.  & Các tiền lược đồ và cấu xạ của tiền lược đồ.\\
\S~3.  & Tích của các tiền lược đồ.\\
\S~4.  & Các tiền lược đồ con và cấu xạ nhúng.\\
\S~5.  & Các tiền lược đồ rút gọn; điều kiện tách.\\
\S~6.  & Các điều kiện hữu hạn.\\
\S~7.  & Các ánh xạ hữu tỉ.\\
\S~8.  & Các lược đồ Chevalley.\\
\S~9.  & Bổ sung về các bó quasi-coherent.\\
\S~10. & Các lược đồ hình thức.
\end{tabular}

\medskip
Các \S\S~1 đến~8 hầu như chỉ nhằm phát triển một ngôn ngữ, ngôn ngữ sẽ được
sử dụng trong toàn bộ phần tiếp theo. Tuy nhiên, cần lưu ý rằng, phù hợp với
tinh thần chung của Chuyên luận này, các \S\S~7 và~8 sẽ được sử dụng ít hơn
các phần khác và theo cách ít thiết yếu hơn; hơn nữa, chúng tôi chỉ bàn đến
các lược đồ Chevalley để tạo mối liên hệ với ngôn ngữ của
Chevalley~\cite{I-1} và Nagata~\cite{I-9}. \S~9 đưa ra các định nghĩa và kết
quả về bó quasi-coherent; một số kết quả trong đó không còn chỉ là việc chuyển
các khái niệm quen thuộc của Đại số giao hoán sang ngôn ngữ ``hình học'', mà
đã mang tính toàn cục; ngay từ các chương sau, chúng sẽ là không thể thiếu
trong việc nghiên cứu toàn cục các cấu xạ. Cuối cùng, \S~10 giới thiệu một sự
khái quát hóa của khái niệm lược đồ, sẽ đóng vai trò trung gian cho chúng ta
trong Chương~III để phát biểu và chứng minh một cách thuận tiện các kết quả cơ
bản của việc nghiên cứu đối đồng điều các cấu xạ riêng; mặt khác, cần lưu ý
rằng khái niệm lược đồ hình thức dường như không thể thiếu để diễn đạt một số
sự kiện của ``lý thuyết môđun'' (các bài toán phân loại đa tạp đại số). Các
kết quả của \S~10 sẽ không được sử dụng trước \S~3 của Chương~III, và nên bỏ
qua phần này cho đến khi đó.

\clearpage

\setcounter{section}{0}
\section{Các lược đồ affine}
\label{section:I.1}

\setcounter{subsection}{0}
\subsection{Phổ nguyên tố của một vành}
\label{subsection:I.1.1}

\begin{env}[1.1.1]
\label{I.1.1.1}
\oldpage[I]{80}
\emph{Ký hiệu}. Cho $A$ là một vành (giao hoán) và $M$ là một $A$-môđun.
Trong chương này và các chương sau, ta sẽ thường xuyên sử dụng các ký hiệu sau:
\begin{itemize}
  \item $\Spec(A)=$ \emph{tập hợp các iđêan nguyên tố} của $A$, còn được gọi là \emph{phổ nguyên tố} của $A$; với $x\in X=\Spec(A)$, thường sẽ thuận tiện khi viết $\mathfrak{j}_x$ thay cho $x$.
    Điều kiện cần và đủ để $\Spec(A)$ \emph{rỗng} là $A$ phải là vành $0$.
  \item $A_x=A_{\mathfrak{j}_x}=$ \emph{vành phân thức (địa phương) $S^{-1}A$}, trong đó $S=A\setmin\mathfrak{j}_x$.
  \item $\mathfrak{m}_x=\mathfrak{j}_x A_{\mathfrak{j}_x}=$ \emph{iđêan cực đại của $A_x$}.
  \item $\kres(x)=A_x/\mathfrak{m}_x=$ \emph{trường thặng dư của $A_x$}, đẳng cấu chính tắc với trường phân thức của miền nguyên $A/\mathfrak{j}_x$, và ta đồng nhất nó với trường ấy.
  \item $f(x)=$ \emph{lớp của $f$ môđun~$\mathfrak{j}_x$} trong $A/\mathfrak{j}_x\subset\kres(x)$, với $f\in A$ và $x\in X$.
    Ta cũng nói $f(x)$ là \emph{giá trị} của $f$ tại điểm $x\in\Spec(A)$; hai quan hệ $f(x)=0$ và $f\in\mathfrak{j}_x$ là \emph{tương đương}.
  \item $M_x=M\otimes_A A_x=$ \emph{môđun phân thức với các mẫu thuộc $A\setmin\mathfrak{j}_x$}.
  \item $\rad(E)=$ \emph{căn của iđêan trong $A$ sinh bởi một tập con $E$ của $A$}.
  \item $V(E)=$ \emph{tập hợp các $x\in X$ sao cho $E\subset\mathfrak{j}_x$} (hay tập hợp các $x\in X$ sao cho $f(x)=0$ với mọi $f\in E$), với $E\subset A$.
    Khi đó
    \[
    \label{I.1.1.1.1}
      \rad(E)=\bigcap_{x\in V(E)}\mathfrak{j}_x.
      \tag{1.1.1.1}
    \]
  \item $V(f)=V(\{f\})$ với $f\in A$.
  \item $D(f)=X\setmin V(f)=$ \emph{tập hợp các $x\in X$ tại đó $f(x)\neq 0$}.
\end{itemize}
\end{env}

\begin{proposition}[1.1.2]
\label{I.1.1.2}
Ta có các tính chất sau:
\begin{enumerate}
  \item[{\rm(i)}] $V(0)=X$, $V(1)=\emp$.
  \item[{\rm(ii)}] Quan hệ $E\subset E'$ kéo theo $V(E)\supset V(E')$.
  \item[{\rm(iii)}] Với mọi họ $(E_\lambda)$ các tập con của $A$, ta có $V(\bigcup_\lambda E_\lambda)=V(\sum_\lambda E_\lambda)=\bigcap_\lambda V(E_\lambda)$.
  \item[{\rm(iv)}] $V(EE')=V(E)\cup V(E')$.
  \item[{\rm(v)}] $V(E)=V(\rad(E))$.
\end{enumerate}
\end{proposition}

\begin{proof}
Các tính chất (i), (ii) và (iii) là hiển nhiên, còn (v) suy ra từ (ii) và công thức \sref{I.1.1.1.1}.
Hiển nhiên $V(EE')\supset V(E)\cup V(E')$;
ngược lại, nếu $x\not\in V(E)$ và $x\not\in V(E')$, thì tồn tại $f\in E$ và $f'\in E'$ sao cho $f(x)\neq 0$ và $f'(x)\neq 0$ trong $\kres(x)$, do đó $f(x)f'(x)\neq 0$, nghĩa là $x\not\in V(EE')$; điều này chứng minh (iv).
\end{proof}

Mệnh đề \sref{I.1.1.2} cho thấy, ngoài những điều khác, rằng các tập hợp có dạng $V(E)$ (khi $E$ chạy qua các tập con của $A$) là các \emph{tập đóng} của một tôpô trên $X$, mà ta sẽ gọi là \emph{tôpô phổ}
\footnote{Việc đưa tôpô này vào hình học đại số là do Zariski.
Vì vậy, tôpô này thường được gọi là ``tôpô Zariski'' của $X$.};
trừ khi có chỉ dẫn rõ ràng ngược lại, ta luôn giả sử $X=\Spec(A)$ được trang bị tôpô phổ.

\begin{env}[1.1.3]
\label{I.1.1.3}
\oldpage[I]{81}
Với mỗi tập con $Y$ của $X$, ta ký hiệu $\mathfrak{j}(Y)$ là tập hợp các $f\in A$ sao cho $f(y)=0$ với mọi $y\in Y$;
tương đương, $\mathfrak{j}(Y)$ là giao của các iđêan nguyên tố $\mathfrak{j}_y$ với $y\in Y$.
Rõ ràng quan hệ $Y\subset Y'$ kéo theo $\mathfrak{j}(Y)\supset\mathfrak{j}(Y')$, và ta có
\[
  \label{I.1.1.3.1}
  \mathfrak{j}\bigg(\bigcup_\lambda Y_\lambda\bigg)=\bigcap_\lambda\mathfrak{j}(Y_\lambda)
  \tag{1.1.3.1}
\]
với mọi họ $(Y_\lambda)$ các tập con của $X$.
Cuối cùng, ta có
\[
  \label{I.1.1.3.2}
  \mathfrak{j}(\{x\})=\mathfrak{j}_x.
  \tag{1.1.3.2}
\]
\end{env}

\begin{proposition}[1.1.4]
\label{I.1.1.4}
\medskip\noindent
\begin{enumerate}
  \item[{\rm(i)}] Với mọi tập con $E$ của $A$, ta có $\mathfrak{j}(V(E))=\rad(E)$.
  \item[{\rm(ii)}] Với mọi tập con $Y$ của $X$, $V(\mathfrak{j}(Y))=\overline{Y}$, là bao đóng của $Y$ trong $X$.
\end{enumerate}
\end{proposition}

\begin{proof}
(i) là hệ quả trực tiếp của các định nghĩa và \sref{I.1.1.1.1};
mặt khác, $V(\mathfrak{j}(Y))$ là tập đóng và chứa $Y$;
ngược lại, nếu $Y\subset V(E)$, thì $f(y)=0$ với $f\in E$ và mọi $y\in Y$, do đó $E\subset\mathfrak{j}(Y)$ và $V(E)\supset V(\mathfrak{j}(Y))$; điều này chứng minh (ii).
\end{proof}

\begin{corollary}[1.1.5]
\label{I.1.1.5}
Các tập con đóng của $X=\Spec(A)$ và các iđêan của $A$ bằng căn của chính chúng (nói cách khác, các iđêan là giao của những iđêan nguyên tố) tương ứng song ánh qua các ánh xạ đảo chiều bao hàm $Y\mapsto\mathfrak{j}(Y)$, $\mathfrak{a}\mapsto V(\mathfrak{a})$;
hợp $Y_1\cup Y_2$ của hai tập con đóng tương ứng với $\mathfrak{j}(Y_1)\cap\mathfrak{j}(Y_2)$, còn giao của một họ bất kỳ $(Y_\lambda)$ các tập con đóng tương ứng với căn của tổng các $\mathfrak{j}(Y_\lambda)$.
\end{corollary}

\begin{corollary}[1.1.6]
\label{I.1.1.6}
Nếu $A$ là một vành Noether, thì $X=\Spec(A)$ là một không gian Noether.
\end{corollary}

Lưu ý rằng mệnh đề đảo của hệ quả này là sai, như được minh họa bởi một miền nguyên không Noether có đúng một iđêan nguyên tố khác $(0)$ (chẳng hạn, một vành định giá không rời rạc hạng $1$).

Như một ví dụ về vành $A$ có phổ không phải là không gian Noether, ta có thể xét vành $\sh{C}(Y)$ các hàm thực liên tục trên một không gian compact vô hạn $Y$;
ta biết rằng, xét như một tập hợp, $Y$ được đồng nhất với tập hợp các iđêan cực đại của $A$, và dễ thấy rằng tôpô trên $Y$ cảm sinh bởi tôpô của $X=\Spec(A)$ chính là tôpô ban đầu của $Y$.
Vì $Y$ không phải là một không gian Noether, nên $X$ cũng vậy.

\begin{corollary}[1.1.7]
\label{I.1.1.7}
Với mọi $x\in X$, bao đóng của $\{x\}$ là tập hợp các $y\in X$ sao cho $\mathfrak{j}_x\subset\mathfrak{j}_y$.
Để $\{x\}$ là đóng, điều kiện cần và đủ là $\mathfrak{j}_x$ cực đại.
\end{corollary}

\begin{corollary}[1.1.8]
\label{I.1.1.8}
Không gian $X=\Spec(A)$ là một không gian Kolmogorov.
\end{corollary}

\begin{proof}
Nếu $x$ và $y$ là hai điểm phân biệt của $X$, thì hoặc $\mathfrak{j}_x\not\subset\mathfrak{j}_y$ hoặc $\mathfrak{j}_y\not\subset\mathfrak{j}_x$; do đó một trong hai điểm $x$, $y$ không thuộc bao đóng của điểm kia.
\end{proof}

\begin{env}[1.1.9]
\label{I.1.1.9}
Theo Mệnh đề \sref{I.1.1.2}[iv], với hai phần tử $f$, $g$ của $A$, ta có
\[
\label{I.1.1.9.1}
  D(fg)=D(f)\cap D(g).
  \tag{1.1.9.1}
\]
Cũng lưu ý rằng đẳng thức $D(f)=D(g)$, theo Mệnh đề~\sref{I.1.1.4}[i] và Mệnh đề~\sref{I.1.1.2}[v], có nghĩa là $\rad(f)=\rad(g)$, hoặc các iđêan nguyên tố tối tiểu chứa $(f)$ và $(g)$ là như nhau;
đặc biệt, điều này cũng đúng khi $f=ug$, trong đó $u$ khả nghịch.
\end{env}

\begin{proposition}[1.1.10]
\label{I.1.1.10}
\medskip\noindent
\oldpage[I]{82}
\begin{enumerate}
  \item[{\rm(i)}] Khi $f$ chạy qua $A$, các tập $D(f)$ tạo thành một cơ sở của tôpô trên $X$.
  \item[{\rm(ii)}] Với mọi $f\in A$, $D(f)$ là quasi-compact.
    Đặc biệt, $X=D(1)$ là quasi-compact.
\end{enumerate}
\end{proposition}

\begin{proof}
\medskip\noindent
\begin{enumerate}
  \item[(i)] Cho $U$ là một tập mở của $X$;
    theo định nghĩa, ta có $U=X\setmin V(E)$, trong đó $E$ là một tập con của $A$, và $V(E)=\bigcap_{f\in E}V(f)$, do đó $U=\bigcup_{f\in E}D(f)$.
  \item[(ii)] Theo (i), chỉ cần chứng minh rằng, nếu $(f_\lambda)_{\lambda\in L}$ là một họ phần tử của $A$ sao cho $D(f)\subset\bigcup_{\lambda\in L}D(f_\lambda)$, thì tồn tại một tập con hữu hạn $J$ của $L$ sao cho $D(f)\subset\bigcup_{\lambda\in J}D(f_\lambda)$.
    Gọi $\mathfrak{a}$ là iđêan của $A$ sinh bởi các $f_\lambda$;
    theo giả thiết, ta có $V(f)\supset V(\mathfrak{a})$, suy ra $\rad(f)\subset\rad(\mathfrak{a})$;
    vì $f\in\rad(f)$, tồn tại một số nguyên $n>0$ sao cho $f^n\in\mathfrak{a}$.
    Khi đó $f^n$ thuộc iđêan $\mathfrak{b}$ sinh bởi họ con hữu hạn $(f_\lambda)_{\lambda\in J}$, và ta có $V(f)=V(f^n)\supset V(\mathfrak{b})=\bigcap_{\lambda\in J}V(f_\lambda)$, tức là $D(f)\subset\bigcup_{\lambda\in J}D(f_\lambda)$.
\end{enumerate}
\end{proof}

\begin{proposition}[1.1.11]
\label{I.1.1.11}
Với mọi iđêan $\mathfrak{a}$ của $A$, $\Spec(A/\mathfrak{a})$ được đồng nhất chính tắc với không gian con đóng $V(\mathfrak{a})$ của $\Spec(A)$.
\end{proposition}

\begin{proof}
Ta biết rằng có một tương ứng song ánh chính tắc, bảo toàn cấu trúc thứ tự do phép bao hàm xác định, giữa các iđêan (tương ứng, các iđêan nguyên tố) của $A/\mathfrak{a}$ và các iđêan (tương ứng, các iđêan nguyên tố) của $A$ chứa $\mathfrak{a}$.
\end{proof}

Nhắc lại rằng tập $\nilrad$ gồm các phần tử lũy linh của $A$ (\emph{căn lũy linh} của $A$) là một iđêan bằng $\rad(0)$, tức giao của tất cả các iđêan nguyên tố của $A$ \sref[0]{0.1.1.1}.

\begin{corollary}[1.1.12]
\label{I.1.1.12}
Các không gian tôpô $\Spec(A)$ và $\Spec(A/\nilrad)$ đồng phôi chính tắc với nhau.
\end{corollary}

\begin{proposition}[1.1.13]
\label{I.1.1.13}
Để $X=\Spec(A)$ là bất khả quy \sref[0]{0.2.1.1}, điều kiện cần và đủ là vành $A/\nilrad$ là một miền nguyên (hay, tương đương, iđêan $\nilrad$ là nguyên tố).
\end{proposition}

\begin{proof}
Theo Hệ quả \sref{I.1.1.12}, ta có thể giới hạn ở trường hợp $\nilrad=0$.
Nếu $X$ là khả quy, thì tồn tại hai tập con đóng thực sự $Y_1$ và $Y_2$ của $X$ sao cho $X=Y_1\cup Y_2$, từ đó $\mathfrak{j}(X)=\mathfrak{j}(Y_1)\cap\mathfrak{j}(Y_2)=0$, vì các iđêan $\mathfrak{j}(Y_1)$ và $\mathfrak{j}(Y_2)$ khác $(0)$ \sref{I.1.1.5};
do đó $A$ không phải là một miền nguyên.
Ngược lại, nếu tồn tại các phần tử $f\neq 0$, $g\neq 0$ của $A$ sao cho $fg=0$, thì $V(f)\neq X$, $V(g)\neq X$ (vì giao của tất cả các iđêan nguyên tố của $A$ là $(0)$), và $X=V(fg)=V(f)\cup V(g)$.
\end{proof}

\begin{corollary}[1.1.14]
\label{I.1.1.14}
\medskip\noindent
\begin{enumerate}
  \item[{\rm(i)}] Trong tương ứng song ánh giữa các tập con đóng của $X=\Spec(A)$ và các iđêan của $A$ bằng căn của chính chúng, các tập con đóng bất khả quy của $X$ tương ứng với các iđêan nguyên tố của $A$.
    Đặc biệt, các thành phần bất khả quy của $X$ tương ứng với các iđêan nguyên tố tối tiểu của $A$.
  \item[{\rm(ii)}] Ánh xạ $x\mapsto\overline{\{x\}}$ thiết lập một tương ứng song ánh giữa $X$ và tập hợp các tập con đóng bất khả quy của $X$ (\emph{tức là,} mọi tập con đóng bất khả quy của $X$ đều có đúng một điểm tổng quát).
\end{enumerate}
\end{corollary}

\begin{proof}
(i) suy ra ngay từ \sref{I.1.1.13} và \sref{I.1.1.11};
và, để chứng minh (ii), theo \sref{I.1.1.11}, ta có thể giới hạn ở trường hợp $X$ bất khả quy;
khi đó, theo Mệnh đề~\sref{I.1.1.13}, tồn tại một iđêan nguyên tố $\nilrad$ của $A$ nằm trong mọi iđêan nguyên tố khác, tương ứng với điểm tổng quát
\oldpage[I]{83}
của $X$;
hơn nữa, $X$ có nhiều nhất một điểm tổng quát, vì nó là một không gian Kolmogorov (\sref{I.1.1.8} và \sref[0]{0.2.1.3}).
\end{proof}

\begin{proposition}[1.1.15]
\label{I.1.1.15}
Nếu $\mathfrak{J}$ là một iđêan của $A$ chứa trong căn Jacobson $\mathfrak{R}(A)$, thì lân cận duy nhất của $V(\mathfrak{J})$ trong $X=\Spec(A)$ là toàn bộ không gian $X$.
\end{proposition}

\begin{proof}
Theo định nghĩa, mọi iđêan cực đại của $A$ đều thuộc $V(\mathfrak{J})$.
Vì mọi iđêan $\mathfrak{a}\neq A$ của $A$ đều nằm trong một iđêan cực đại, ta có $V(\mathfrak{a})\cap V(\mathfrak{J})\neq\emp$, từ đó suy ra mệnh đề.
\end{proof}

\subsection{Các tính chất hàm tử của phổ nguyên tố của vành}
\label{subsection:I.1.2}

\begin{env}[1.2.1]
\label{I.1.2.1}
Cho $A$, $A'$ là hai vành và
\[
  \vphi:A'\to A
\]
là một đồng cấu vành.
Với mọi iđêan nguyên tố $x=\mathfrak{j}_x\in\Spec(A)=X$, vành $A'/\vphi^{-1}(\mathfrak{j}_x)$ đẳng cấu chính tắc với một vành con của $A/\mathfrak{j}_x$, nên là một miền nguyên; nói cách khác, $\vphi^{-1}(\mathfrak{j}_x)$ là một iđêan nguyên tố của $A'$;
ta ký hiệu nó bởi ${}^a\vphi(x)$, và như vậy đã xác định một ánh xạ
\[
  {}^a\vphi:X=\Spec(A)\to X'=\Spec(A')
\]
(cũng được ký hiệu bởi $\Spec(\vphi)$), mà ta gọi là ánh xạ \emph{liên kết} với đồng cấu $\vphi$.
Ta ký hiệu bởi $\vphi^x$ đồng cấu đơn ánh từ $A'/\vphi^{-1}(\mathfrak{j}_x)$ vào $A/\mathfrak{j}_x$ cảm sinh bởi $\vphi$ khi chuyển sang các vành thương, cũng như mở rộng chính tắc của nó thành một đơn cấu trường
\[
  \vphi^x:\kres({}^a\vphi(x))\to\kres(x);
\]
với mọi $f'\in A'$, theo định nghĩa ta có
\[
  \label{I.1.2.1.1}
  \vphi^x(f'({}^a\vphi(x)))=(\vphi(f'))(x)\qquad(x\in X).
  \tag{1.2.1.1}
\]
\end{env}

\begin{proposition}[1.2.2]
\label{I.1.2.2}
\medskip\noindent
\begin{enumerate}
  \item[{\rm(i)}] Với mọi tập con $E'$ của $A'$, ta có
    \[
    \label{I.1.2.2.1}
      {}^a\vphi^{-1}(V(E'))=V(\vphi(E')),
      \tag{1.2.2.1}
    \]
    và đặc biệt, với mọi $f'\in A'$,
    \[
    \label{I.1.2.2.2}
      {}^a\vphi^{-1}(D(f'))=D(\vphi(f')).
      \tag{1.2.2.2}
    \]
  \item[{\rm(ii)}] Với mọi iđêan $\mathfrak{a}$ của $A$, ta có
    \[
    \label{I.1.2.2.3}
      \overline{{}^a\vphi(V(\mathfrak{a}))}=V(\vphi^{-1}(\mathfrak{a})).
      \tag{1.2.2.3}
    \]
\end{enumerate}
\end{proposition}

\begin{proof}
Theo định nghĩa, quan hệ ${}^a\vphi(x)\in V(E')$ tương đương với $E'\subset\vphi^{-1}(\mathfrak{j}_x)$, do đó tương đương với $\vphi(E')\subset\mathfrak{j}_x$, và cuối cùng với $x\in V(\vphi(E'))$; điều này chứng minh (i).
Để chứng minh (ii), ta có thể giả sử $\mathfrak{a}$ bằng căn của nó, vì $V(\rad(\mathfrak{a}))=V(\mathfrak{a})$ \sref{I.1.1.2}[(v)] và $\vphi^{-1}(\rad(\mathfrak{a}))=\rad(\vphi^{-1}(\mathfrak{a}))$;
nếu đặt $Y=V(\mathfrak{a})$ và $\mathfrak{a}'=\mathfrak{j}({}^a\vphi(Y))$, thì $\overline{{}^a\vphi(Y)}=V(\mathfrak{a}')$ (\sref{I.1.1.4}[(ii)]).
Theo định nghĩa, quan hệ $f'\in\mathfrak{a}'$ tương đương với $f'(x')=0$ với mọi $x'\in{}^a\vphi(Y)$; do đó, theo công thức~\sref{I.1.2.1.1}, nó cũng tương đương với $\vphi(f')(x)=0$ với mọi $x\in Y$, hay với $\vphi(f')\in\mathfrak{j}(Y)=\mathfrak{a}$, vì $\mathfrak{a}$ bằng căn của nó;
từ đó suy ra (ii).
\end{proof}

\begin{corollary}[1.2.3]
\label{I.1.2.3}
Ánh xạ ${}^a\vphi$ là liên tục.
\end{corollary}

Lưu ý rằng, nếu $A''$ là một vành thứ ba và $\vphi'$ là một đồng cấu $A''\to A'$, thì ta có ${}^a(\vphi\circ\vphi')={}^a\vphi'\circ{}^a\vphi$;
kết quả này, cùng với Hệ quả~\sref{I.1.2.3}, có nghĩa là $\Spec(A)$ là một \emph{hàm tử phản biến} theo $A$, từ phạm trù các vành đến phạm trù các không gian tôpô.

\begin{corollary}[1.2.4]
\label{I.1.2.4}
\oldpage[I]{84}
Giả sử $\vphi$ có tính chất rằng mọi $f\in A$ đều viết được dưới dạng $f=h\vphi(f')$, trong đó $h$ khả nghịch trong $A$ (\emph{điều này đặc biệt đúng khi $\vphi$ là \emph{toàn ánh}}).
Khi đó ${}^a\vphi$ là một phép đồng phôi từ $X$ lên ${}^a\vphi(X)$.
\end{corollary}

\begin{proof}
Ta sẽ chứng minh rằng, với mọi tập con $E\subset A$, tồn tại một tập con $E'$ của $A'$ sao cho $V(E)=V(\vphi(E'))$;
theo tiên đề ($T_0$) \sref{I.1.1.8} và công thức \sref{I.1.2.2.1}, điều này trước hết suy ra ${}^a\vphi$ là đơn ánh, rồi, lại theo \sref{I.1.2.2.1}, suy ra ${}^a\vphi$ là một phép đồng phôi.
Thật vậy, với mỗi $f\in E$, chỉ cần chọn $f'\in A'$ sao cho $h\vphi(f')=f$, với $h$ khả nghịch trong $A$;
tập $E'$ gồm các phần tử $f'$ ấy chính là tập cần tìm.
\end{proof}

\begin{env}[1.2.5]
\label{I.1.2.5}
Đặc biệt, khi $\vphi$ là đồng cấu chính tắc từ $A$ lên vành thương $A/\mathfrak{a}$, ta thu lại \sref{I.1.1.11}, và ${}^a\vphi$ là \emph{phép nhúng chính tắc} của $V(\mathfrak{a})$, được đồng nhất với $\Spec(A/\mathfrak{a})$, vào $X=\Spec(A)$.
\end{env}

Một trường hợp riêng khác của \sref{I.1.2.4}:

\begin{corollary}[1.2.6]
\label{I.1.2.6}
Nếu $S$ là một tập nhân của $A$, thì phổ $\Spec(S^{-1}A)$ được đồng nhất chính tắc, cùng với tôpô của nó, với không gian con của $X=\Spec(A)$ gồm các $x$ sao cho $\mathfrak{j}_x\cap S=\emp$.
\end{corollary}

\begin{proof}
Theo \sref[0]{0.1.2.6}, ta biết rằng các iđêan nguyên tố của $S^{-1}A$ là các iđêan $S^{-1}\mathfrak{j}_x$ sao cho $\mathfrak{j}_x\cap S=\emp$, và $\mathfrak{j}_x=(i_A^S)^{-1}(S^{-1}\mathfrak{j}_x)$.
Do đó chỉ cần áp dụng Hệ quả~\sref{I.1.2.4} cho $i_A^S$.
\end{proof}

\begin{corollary}[1.2.7]
\label{I.1.2.7}
Để ${}^a\vphi(X)$ trù mật trong $X'$, điều kiện cần và đủ là mọi phần tử của hạt nhân $\Ker\vphi$ đều lũy linh.
\end{corollary}

\begin{proof}
Áp dụng công thức~\sref{I.1.2.2.3} cho iđêan $\mathfrak{a}=(0)$, ta được $\overline{{}^a\vphi(X)}=V(\Ker\vphi)$; để $V(\Ker\vphi)=X'$ đúng, điều kiện cần và đủ là $\Ker\vphi$ nằm trong mọi iđêan nguyên tố của $A'$, hay tương đương, nằm trong căn lũy linh $\nilrad$ của $A'$.
\end{proof}

\subsection{Bó liên kết với một môđun}
\label{subsection:I.1.3}
\label{I.1.3}

\begin{env}[1.3.1]
\label{I.1.3.1}
Cho $A$ là một vành giao hoán, $M$ là một $A$-môđun, $f$ là một phần tử của $A$, và $S_f$ là tập nhân gồm các $f^n$, với $n\geq 0$.
Nhắc lại rằng ta đặt $A_f=S_f^{-1}A$, $M_f=S_f^{-1}M$.
Nếu $S_f'$ là tập nhân bão hòa của $A$ gồm các $g\in A$ chia một phần tử của $S_f$, thì ta biết rằng $A_f$ và $M_f$ được đồng nhất chính tắc lần lượt với ${S_f'}^{-1}A$ và ${S_f'}^{-1}M$ \sref[0]{0.1.4.3}.
\end{env}

\begin{lemma}[1.3.2]
\label{I.1.3.2}
Các điều kiện sau tương đương:
\begin{enumerate}
  \item[{\rm(a)}] $g\in S_f'$;
  \item[{\rm(b)}] $S_g'\subset S_f'$;
  \item[{\rm(c)}] $f\in\rad(g)$;
  \item[{\rm(d)}] $\rad(f)\subset\rad(g)$;
  \item[{\rm(e)}] $V(g)\subset V(f)$;
  \item[{\rm(f)}] $D(f)\subset D(g)$.
\end{enumerate}
\end{lemma}

\begin{proof}
Điều này suy ra ngay từ các định nghĩa và \sref{I.1.1.5}.
\end{proof}

\begin{env}[1.3.3]
\label{I.1.3.3}
Nếu $D(f)=D(g)$, Bổ đề~\sref{I.1.3.2}[b] cho thấy $M_f=M_g$.
Tổng quát hơn, nếu $D(f)\supset D(g)$, thì $S_f'\subset S_g'$, và theo \sref[0]{0.1.4.1} ta biết rằng có một đồng cấu chính tắc có tính hàm tử
\[
  \rho_{g,f}:M_f\to M_g,
\]
và nếu $D(f)\supset D(g)\supset D(h)$, thì theo \sref[0]{0.1.4.4} ta có
\[
  \rho_{h,g}\circ\rho_{g,f}=\rho_{h,f}.
  \tag{1.3.3.1}
  \label{I.1.3.3.1}
\]
\end{env}

\oldpage[I]{85}
Khi $f$ chạy qua các phần tử của $A\setmin\mathfrak{j}_x$ (với một $x$ cho trước trong $X=\Spec(A)$), các tập $S_f'$ tạo thành một họ lọc tăng gồm các tập con của $A-\mathfrak{j}_x$, vì với hai phần tử $f$ và $g$ của $A\setmin\mathfrak{j}_x$, cả $S_f'$ lẫn $S_g'$ đều nằm trong $S_{fg}'$;
vì hợp của các $S_f'$ với $f\in A\setmin\mathfrak{j}_x$ là $A-\mathfrak{j}_x$, theo \sref[0]{0.1.4.5} ta suy ra rằng $A_x$-môđun $M_x$ được đồng nhất chính tắc với \emph{giới hạn quy nạp} $\varinjlim M_f$, đối với họ đồng cấu $(\rho_{g,f})$.
Ta ký hiệu bởi
\[
  \rho_x^f:M_f\to M_x
\]
đồng cấu chính tắc với $f\in A\setmin\mathfrak{j}_x$ (hay tương đương, $x\in D(f)$).

\begin{definition}[1.3.4]
\label{I.1.3.4}
Ta gọi \emph{bó cấu trúc của phổ nguyên tố} $X=\Spec(A)$ (tương ứng, \emph{bó liên kết với $A$-môđun} $M$), và ký hiệu bởi $\widetilde{A}$ hoặc $\sh{O}_X$ (tương ứng, $\widetilde{M}$), bó vành (tương ứng, $\widetilde{A}$-môđun) liên kết với tiền bó $D(f)\mapsto A_f$ (tương ứng, $D(f)\mapsto M_f$), xác định trên cơ sở $\mathfrak{B}$ của $X$ gồm các $D(f)$ với $f\in A$ (\sref{I.1.1.10}, \sref[0]{0.3.2.1} và \sref[0]{0.3.5.6}).
\end{definition}

Ta đã thấy trong \sref[0]{0.3.2.4} rằng cuống $\widetilde{A}_x$ (tương ứng, $\widetilde{M}_x$) \emph{được đồng nhất với vành $A_x$} (tương ứng, \emph{với $A_x$-môđun $M_x$});
ta ký hiệu bởi
\[
  \theta_f:A_f\to\Gamma(D(f),\widetilde{A})
\]
\[
  \text{(tương ứng, }\theta_f:M_f\to\Gamma(D(f),\widetilde{M})\text{),}
\]
ánh xạ chính tắc sao cho, với mọi $x\in D(f)$ và mọi $\xi\in M_f$, ta có
\[
  (\theta_f(\xi))_x=\rho_x^f(\xi).
  \tag{1.3.4.1}
  \label{I.1.3.4.1}
\]

\begin{proposition}[1.3.5]
\label{I.1.3.5}
$\widetilde{M}$ là một hàm tử hiệp biến khớp theo $M$, từ phạm trù các $A$-môđun đến phạm trù các $\widetilde{A}$-môđun.
\end{proposition}

\begin{proof}
Thật vậy, cho $M$, $N$ là hai $A$-môđun và $u$ là một đồng cấu $M\to N$;
với mọi $f\in A$, đồng cấu $u$ xác định chính tắc một đồng cấu $u_f$ từ $A_f$-môđun $M_f$ vào $A_f$-môđun $N_f$, và biểu đồ (với $D(g)\subset D(f)$)
\[
  \xymatrix{
    M_f\ar[r]^{u_f}\ar[d]_{\rho_{g,f}} & N_f\ar[d]^{\rho_{g,f}}\\
    M_g\ar[r]_{u_g} & N_g
  }
\]
là giao hoán \sref[0]{0.1.4.1};
do đó các đồng cấu này xác định một đồng cấu các $\widetilde{A}$-môđun $\widetilde{u}:\widetilde{M}\to\widetilde{N}$ \sref[0]{0.3.2.3}.
Hơn nữa, với mọi $x\in X$, $\widetilde{u}_x$ là giới hạn quy nạp của các $u_f$ với $x\in D(f)$ ($f\in A$), và do đó theo \sref[0]{0.1.4.5}, nếu đồng nhất chính tắc $\widetilde{M}_x$ và $\widetilde{N}_x$ lần lượt với $M_x$ và $N_x$, thì $\widetilde{u}_x$ được đồng nhất với đồng cấu $u_x$ cảm sinh chính tắc bởi $u$.
Nếu $P$ là một $A$-môđun thứ ba, $v$ là một đồng cấu $N\to P$ và $w=v\circ u$, thì hiển nhiên $w_x=v_x\circ u_x$, do đó $\widetilde{w}=\widetilde{v}\circ\widetilde{u}$.
Như vậy ta đã xác định một \emph{hàm tử hiệp biến (theo $M$)} $\widetilde{M}$, từ phạm trù các $A$-môđun đến phạm trù các $\widetilde{A}$-môđun.
\emph{Hàm tử này là khớp}, vì với mọi $x\in X$, $M_x$ là một hàm tử khớp theo $M$ \sref[0]{0.1.3.2};
hơn nữa, ta có $\Supp(M)=\Supp(\widetilde{M})$, theo định nghĩa của hai vế (\sref[0]{0.1.7.1} và \sref[0]{0.3.1.6}).
\end{proof}

\oldpage[I]{86}
\begin{proposition}[1.3.6]
\label{I.1.3.6}
Với mọi $f\in A$, tập con mở $D(f)\subset X$ được đồng nhất chính tắc với phổ nguyên tố $\Spec(A_f)$, và bó $\widetilde{M_f}$ liên kết với $A_f$-môđun $M_f$ được đồng nhất chính tắc với hạn chế $\widetilde{M}|D(f)$.
\end{proposition}

\begin{proof}
Khẳng định thứ nhất là một trường hợp riêng của \sref{I.1.2.6}.
Hơn nữa, nếu $g\in A$ sao cho $D(g)\subset D(f)$, thì $M_g$ được đồng nhất chính tắc với môđun phân thức của $M_f$ mà các mẫu số là lũy thừa của ảnh chính tắc của $g$ trong $A_f$ \sref[0]{0.1.4.6}.
Khi đó, sự đồng nhất chính tắc của $\widetilde{M_f}$ với $\widetilde{M}|D(f)$ suy ra từ các định nghĩa.
\end{proof}

\begin{theorem}[1.3.7]
\label{I.1.3.7}
Với mọi $A$-môđun $M$ và mọi $f\in A$, đồng cấu
\[
  \theta_f:M_f\to\Gamma(D(f),\widetilde{M})
\]
là song ánh \emph{(nói cách khác, tiền bó $D(f)\mapsto M_f$ là một \emph{bó})}.
Đặc biệt, $M$ được đồng nhất với $\Gamma(X,\widetilde{M})$ qua $\theta_1$.
\end{theorem}

\begin{proof}
Lưu ý rằng, nếu $M=A$, thì $\theta_f$ là một đồng cấu vành;
do đó, Định lý~\sref{I.1.3.7} sẽ suy ra rằng, nếu đồng nhất các vành $A_f$ và $\Gamma(D(f),\widetilde{A})$ qua $\theta_f$, thì đồng cấu $\theta_f:M_f\to\Gamma(D(f),\widetilde{M})$ sẽ là một đẳng cấu \emph{môđun}.

Trước hết, ta chứng minh $\theta_f$ là \emph{đơn ánh}.
Thật vậy, nếu $\xi\in M_f$ thỏa $\theta_f(\xi)=0$, thì điều này có nghĩa là với mọi iđêan nguyên tố $\mathfrak{p}$ của $A_f$, tồn tại $h\not\in\mathfrak{p}$ sao cho $h\xi=0$;
vì iđêan triệt tiêu của $\xi$ không nằm trong bất kỳ iđêan nguyên tố nào của $A_f$, nó bằng toàn bộ vành $A_f$, và do đó $\xi=0$.

Còn phải chứng minh $\theta_f$ là \emph{toàn ánh};
ta có thể quy về trường hợp $f=1$ (trường hợp tổng quát khi đó suy ra bằng cách ``địa phương hóa'', nhờ \sref{I.1.3.6}).
Cho $s$ là một nhát cắt của $\widetilde{M}$ trên $X$;
theo \sref{I.1.3.4} và \sref{I.1.1.10}[ii], tồn tại một phủ \emph{hữu hạn} $(D(f_i))_{i\in I}$ của $X$ ($f_i\in A$) sao cho, với mọi $i\in I$, hạn chế $s_i=s|D(f_i)$ có dạng $\theta_{f_i}(\xi_i)$, trong đó $\xi_i\in M_{f_i}$.
Nếu $i$, $j$ là hai chỉ số của $I$, khi biểu thị rằng các hạn chế của $s_i$ và $s_j$ lên $D(f_i)\cap D(f_j)=D(f_i f_j)$ bằng nhau, theo định nghĩa của $\widetilde{M}$ ta được
\[
  \label{I.1.3.7.1}
  \rho_{f_i f_j,f_i}(\xi_i)=\rho_{f_i f_j,f_j}(\xi_j).
  \tag{1.3.7.1}
\]
Theo định nghĩa, với mọi $i\in I$, ta có thể viết $\xi_i=z_i/f_i^{n_i}$, trong đó $z_i\in M$, và vì $I$ hữu hạn, bằng cách nhân mỗi $z_i$ với một lũy thừa của $f_i$, ta có thể giả sử mọi $n_i$ đều bằng cùng một số $n$.
Khi đó, theo định nghĩa, \sref{I.1.3.7.1} có nghĩa là tồn tại một số nguyên $m_{ij}\geq 0$ sao cho $(f_i f_j)^{m_{ij}}(f_j^n z_i-f_i^n z_j)=0$, và ta còn có thể giả sử mọi $m_{ij}$ đều bằng cùng một số nguyên $m$;
thay mỗi $z_i$ bởi $f_i^m z_i$, ta quy về trường hợp $m=0$, tức là trường hợp
\[
  \label{I.1.3.7.2}
  f_j^n z_i=f_i^n z_j
  \tag{1.3.7.2}
\]
với mọi $i$, $j$.
Ta có $D(f_i^n)=D(f_i)$, và vì các $D(f_i)$ tạo thành một phủ của $X$, iđêan sinh bởi các $f_i^n$ là $A$;
nói cách khác, tồn tại các phần tử $g_i\in A$ sao cho $\sum_i g_i f_i^n=1$.
Xét phần tử $z=\sum_i g_i z_i$ của $M$;
theo \sref{I.1.3.7.2}, ta có $f_i^n z=\sum_j g_j f_i^n z_j=(\sum_j g_j f_j^n)z_i=z_i$, do đó theo định nghĩa $\xi_i=z/1$ trong $M_{f_i}$.
Suy ra
\oldpage[I]{87}
các $s_i$ là các hạn chế trên $D(f_i)$ của $\theta_1(z)$, điều này chứng minh $s=\theta_1(z)$ và kết thúc chứng minh.
\end{proof}

\begin{corollary}[1.3.8]
\label{I.1.3.8}
Cho $M$ và $N$ là hai $A$-môđun;
đồng cấu chính tắc $u\mapsto\widetilde{u}$ từ $\Hom_A(M,N)$ vào $\Hom_{\widetilde{A}}(\widetilde{M},\widetilde{N})$ là song ánh.
Đặc biệt, các hệ thức $M=0$ và $\widetilde{M}=0$ là tương đương.
\end{corollary}

\begin{proof}
Xét đồng cấu chính tắc
\[
  v\longmapsto\Gamma(v):
  \Hom_{\widetilde{A}}(\widetilde{M},\widetilde{N})
  \longrightarrow
  \Hom_{\Gamma(\widetilde{A})}(\Gamma(\widetilde{M}),\Gamma(\widetilde{N}));
\]
môđun sau cùng được đồng nhất chính tắc với $\Hom_A(M,N)$, theo Định lý~\sref{I.1.3.7}.
Còn phải kiểm tra rằng $u\mapsto\widetilde{u}$ và $v\mapsto\Gamma(v)$ là nghịch đảo của nhau;
hiển nhiên $\Gamma(\widetilde{u})=u$ theo định nghĩa của $\widetilde{u}$;
mặt khác, nếu đặt $u=\Gamma(v)$ với $v\in\Hom_{\widetilde{A}}(\widetilde{M},\widetilde{N})$, thì ánh xạ $w:\Gamma(D(f),\widetilde{M})\to\Gamma(D(f),\widetilde{N})$ được $v$ cảm sinh chính tắc sao cho biểu đồ
\[
  \xymatrix{
    M\ar[r]^u\ar[d]_{\rho_{f,1}} & N\ar[d]^{\rho_{f,1}}\\
    M_f\ar[r]_w & N_f
  }
\]
giao hoán;
do đó, tất yếu $w=u_f$ với mọi $f\in A$ \sref[0]{0.1.2.4}, điều này chứng minh $\widetilde{\Gamma(v)}=v$.
\end{proof}

\begin{corollary}[1.3.9]
\label{I.1.3.9}
\medskip\noindent
\begin{enumerate}
  \item[{\rm(i)}] Cho $u$ là một đồng cấu từ $A$-môđun $M$ vào $A$-môđun $N$;
    khi đó các bó liên kết với $\Ker u$, $\Im u$ và $\Coker u$ lần lượt là $\Ker\widetilde{u}$, $\Im\widetilde{u}$ và $\Coker\widetilde{u}$.
    Đặc biệt, để $\widetilde{u}$ là đơn ánh (tương ứng toàn ánh, song ánh), điều kiện cần và đủ là $u$ có tính chất ấy.
  \item[{\rm(ii)}] Nếu $M$ là giới hạn quy nạp (tương ứng tổng trực tiếp) của một họ $A$-môđun $(M_\lambda)$, thì $\widetilde{M}$ là giới hạn quy nạp (tương ứng tổng trực tiếp) của họ $(\widetilde{M_\lambda})$, qua một đẳng cấu chính tắc.
\end{enumerate}
\end{corollary}

\begin{proof}
\medskip\noindent
\begin{enumerate}
  \item[(i)] Chỉ cần áp dụng sự kiện rằng $\widetilde{M}$ là một hàm tử khớp theo $M$ \sref{I.1.3.5} cho hai dãy khớp các $A$-môđun
    \[
      0\to\Ker u\to M\to\Im u\to 0,
    \]
    \[
      0\to\Im u\to N\to\Coker u\to 0.
    \]
    Khẳng định thứ hai khi đó suy ra từ Định lý~\sref{I.1.3.7}.
  \item[(ii)] Cho $(M_\lambda,g_{\mu\lambda})$ là một hệ quy nạp các $A$-môđun, có giới hạn quy nạp $M$, và cho $g_\lambda$ là đồng cấu chính tắc $M_\lambda\to M$.
    Vì $\widetilde{g_{\nu\mu}}\circ\widetilde{g_{\mu\lambda}}=\widetilde{g_{\nu\lambda}}$ và $\widetilde{g_\lambda}=\widetilde{g_\mu}\circ\widetilde{g_{\mu\lambda}}$ với $\lambda\leq\mu\leq\nu$, nên $(\widetilde{M_\lambda},\widetilde{g_{\mu\lambda}})$ là một hệ quy nạp các bó trên $X$ và, nếu ký hiệu bởi $h_\lambda$ đồng cấu chính tắc $\widetilde{M_\lambda}\to\varinjlim\widetilde{M_\lambda}$, thì tồn tại duy nhất một đồng cấu $v:\varinjlim\widetilde{M_\lambda}\to\widetilde{M}$ sao cho $v\circ h_\lambda=\widetilde{g_\lambda}$.
    Để thấy $v$ là song ánh, chỉ cần kiểm tra rằng, với mọi $x\in X$, $v_x$ là một song ánh từ $(\varinjlim\widetilde{M_\lambda})_x$ vào $\widetilde{M}_x$;
    nhưng $\widetilde{M}_x=M_x$, và
    \[
      (\varinjlim\widetilde{M_\lambda})_x=\varinjlim(\widetilde{M_\lambda})_x
      =\varinjlim(M_\lambda)_x=M_x\quad\sref[0]{0.1.3.3}.
    \]
    Mặt khác, từ các định nghĩa suy ra rằng $(\widetilde{g_\lambda})_x$ và
    $(h_\lambda)_x$ đều bằng ánh xạ chính tắc từ $(M_\lambda)_x$ vào $M_x$;
    vì $(\widetilde{g_\lambda})_x=v_x\circ(h_\lambda)_x$, nên $v_x$ là ánh xạ đồng nhất.
\oldpage[I]{88}
    Cuối cùng, nếu $M$ là tổng trực tiếp của hai $A$-môđun $N$ và $P$, thì hiển nhiên $\widetilde{M}=\widetilde{N}\oplus\widetilde{P}$;
    vì mọi tổng trực tiếp là giới hạn quy nạp của các tổng trực tiếp hữu hạn, các khẳng định trong (ii) được chứng minh.
\end{enumerate}
\end{proof}

Lưu ý rằng Hệ quả~\sref{I.1.3.8} chứng minh rằng các bó đẳng cấu với những bó liên kết với $A$-môđun tạo thành một \emph{phạm trù Abel} (T, I, 1.4).

Lưu ý thêm rằng từ Hệ quả~\sref{I.1.3.9} suy ra rằng, nếu $M$ là một $A$-môđun \emph{hữu hạn sinh} (tức là nếu tồn tại một đồng cấu toàn ánh $A^n\to M$), thì tồn tại một đồng cấu toàn ánh $\widetilde{A^n}\to\widetilde{M}$ hay, nói cách khác, $\widetilde{A}$-môđun $\widetilde{M}$ được \emph{sinh bởi một họ hữu hạn các nhát cắt trên $X$} \sref[0]{0.5.1.1}, và ngược lại.

\begin{env}[1.3.10]
\label{I.1.3.10}
Nếu $N$ là một môđun con của $A$-môđun $M$, phép nhúng chính tắc $j:N\to M$ cho, theo \sref{I.1.3.9}, một đồng cấu đơn ánh $\widetilde{N}\to\widetilde{M}$, nhờ đó có thể đồng nhất chính tắc $\widetilde{N}$ với một \emph{$\widetilde{A}$-môđun con} của $\widetilde{M}$;
ta sẽ luôn giả sử đã thực hiện phép đồng nhất này.
Nếu $N$ và $P$ là hai môđun con của $M$, thì
\[
\label{I.1.3.10.1}
  (N+P)\supertilde=\widetilde{N}+\widetilde{P},
  \tag{1.3.10.1}
\]
\[
\label{I.1.3.10.2}
  (N\cap P)\supertilde=\widetilde{N}\cap\widetilde{P},
  \tag{1.3.10.2}
\]
vì $N+P$ và $N\cap P$ lần lượt là ảnh của đồng cấu chính tắc
$N\oplus P\to M$ và hạt nhân của đồng cấu chính tắc $M\to(M/N)\oplus(M/P)$, và chỉ cần áp dụng \sref{I.1.3.9}.

Từ \sref{I.1.3.10.1} và \sref{I.1.3.10.2} suy ra rằng, nếu $\widetilde{N}=\widetilde{P}$, thì $N=P$.
\end{env}

\begin{corollary}[1.3.11]
\label{I.1.3.11}
Trong phạm trù các bó đẳng cấu với những bó liên kết với $A$-môđun, hàm tử $\Gamma$ là khớp.
\end{corollary}

\begin{proof}
Cho $\widetilde{M}\xrightarrow{\widetilde{u}}\widetilde{N}\xrightarrow{\widetilde{v}}\widetilde{P}$ là một dãy khớp tương ứng với hai đồng cấu $u:M\to N$ và $v:N\to P$ của các $A$-môđun.
Nếu $Q=\Im u$ và $R=\Ker v$, thì $\widetilde{Q}=\Im\widetilde{u}=\Ker\widetilde{v}=\widetilde{R}$ (Hệ quả~\sref{I.1.3.9}), do đó $Q=R$.
\end{proof}

\begin{corollary}[1.3.12]
\label{I.1.3.12}
Cho $M$ và $N$ là hai $A$-môđun.
\begin{enumerate}
  \item[{\rm(i)}] Bó liên kết với $M\otimes_A N$ được đồng nhất chính tắc với $\widetilde{M}\otimes_{\widetilde{A}}\widetilde{N}$.
  \item[{\rm(ii)}] Nếu thêm vào đó $M$ có một trình bày hữu hạn, thì bó liên kết với $\Hom_A(M,N)$ được đồng nhất chính tắc với $\shHom_{\widetilde{A}}(\widetilde{M},\widetilde{N})$.
\end{enumerate}
\end{corollary}

\begin{proof}
\medskip\noindent
\begin{enumerate}
  \item[(i)] Bó $\sh{F}=\widetilde{M}\otimes_{\widetilde{A}}\widetilde{N}$ liên kết với tiền bó
    \[
      U\mapsto\sh{F}(U)=\Gamma(U,\widetilde{M})\otimes_{\Gamma(U,\widetilde{A})}\Gamma(U,\widetilde{N}),
    \]
    trong đó $U$ chạy qua cơ sở \sref{I.1.1.10}[i] của $X$ gồm các $D(f)$, với $f\in A$.
    Ta biết rằng $\sh{F}(D(f))$ được đồng nhất chính tắc với $M_f\otimes_{A_f}N_f$, theo \sref{I.1.3.7} và \sref{I.1.3.6}.
    Mặt khác, ta biết $A_f$-môđun $M_f\otimes_{A_f}N_f$ đẳng cấu chính tắc với $(M\otimes_A N)_f$ \sref[0]{0.1.3.4}, và môđun sau lại đẳng cấu chính tắc với $\Gamma(D(f),(M\otimes_A N)\supertilde)$ (Định lý~\sref{I.1.3.7} và Mệnh đề~\sref{I.1.3.6}).
    Hơn nữa, kiểm tra trực tiếp cho thấy các đẳng cấu chính tắc
    \[
      \sh{F}(D(f))\isoto\Gamma(D(f),(M\otimes_A N)\supertilde)
    \]
\oldpage[I]{89}
    thu được như vậy thỏa các điều kiện tương thích đối với các toán tử hạn chế \sref[0]{0.1.4.2}, nên chúng xác định một đẳng cấu chính tắc có tính hàm tử
    \[
      \widetilde{M}\otimes_{\widetilde{A}}\widetilde{N}\isoto(M\otimes_A N)\supertilde.
    \]
  \item[(ii)] Bó $\sh{G}=\shHom_{\widetilde{A}}(\widetilde{M},\widetilde{N})$ liên kết với tiền bó
    \[
      U\mapsto\sh{G}(U)=\Hom_{\widetilde{A}|U}(\widetilde{M}|U,\widetilde{N}|U),
    \]
    trong đó $U$ chạy qua cơ sở của $X$ gồm các $D(f)$.
    Ta biết rằng $\sh{G}(D(f))$ được đồng nhất chính tắc với $\Hom_{A_f}(M_f,N_f)$ (Mệnh đề~\sref{I.1.3.6} và
    Hệ quả~\sref{I.1.3.8}), và, theo giả thiết về $M$, đối tượng này được đồng nhất chính tắc với $(\Hom_A(M,N))_f$ \sref[0]{0.1.3.5}.
    Cuối cùng, $(\Hom_A(M,N))_f$ được đồng nhất chính tắc với $\Gamma(D(f),(\Hom_A(M,N))\supertilde)$ (Mệnh đề~\sref{I.1.3.6} và Định lý~\sref{I.1.3.7}), và các đẳng cấu chính tắc $\sh{G}(D(f))\isoto\Gamma(D(f),(\Hom_A(M,N))\supertilde)$ thu được như vậy tương thích với các toán tử hạn chế \sref[0]{0.1.4.2};
    do đó chúng xác định một đẳng cấu chính tắc $\shHom_{\widetilde{A}}(\widetilde{M},\widetilde{N})\isoto(\Hom_A(M,N))\supertilde$.
\end{enumerate}
\end{proof}

\begin{env}[1.3.13]
\label{I.1.3.13}
Bây giờ cho $B$ là một $A$-đại số (giao hoán); điều này có thể được hiểu là $B$ là một $A$-môđun cùng với một phần tử $e\in B$ và một $A$-đồng cấu $\vphi:B\otimes_A B\to B$, sao cho: $1^\circ$ các biểu đồ
\[
  \xymatrix{
    B\otimes_A B\otimes_A B\ar[r]^{\vphi\otimes 1}\ar[d]_{1\otimes\vphi} &
    B\otimes_A B\ar[d]^\vphi & &
    B\otimes_A B\ar[rr]^\sigma\ar[rd]_\vphi & &
    B\otimes_A B\ar[dl]^\vphi\\
    B\otimes_A B\ar[r]_\vphi &
    B & & &
    B
  }
\]
(trong đó $\sigma$ là phép đối xứng chính tắc) giao hoán; $2^\circ$ $\vphi(e\otimes x)=\vphi(x\otimes e)=x$.
Theo Hệ quả~\sref{I.1.3.12}, đồng cấu $\widetilde{\vphi}:\widetilde{B}\otimes_{\widetilde{A}}\widetilde{B}\to\widetilde{B}$ của các $\widetilde{A}$-môđun thỏa các điều kiện tương tự, do đó xác định một cấu trúc \emph{$\widetilde{A}$-đại số} trên $\widetilde{B}$.
Tương tự, cho một $B$-môđun $N$ tương đương với cho một $A$-môđun $N$ cùng một $A$-đồng cấu $\psi:B\otimes_A N\to N$ sao cho biểu đồ
\[
  \xymatrix{
    B\otimes_A B\otimes_A N\ar[r]^{\vphi\otimes 1}\ar[d]_{1\otimes\psi} &
    B\otimes_A N\ar[d]^\psi\\
    B\otimes_A N\ar[r]_\psi &
    N
  }
\]
giao hoán và $\psi(e\otimes n)=n$;
đồng cấu $\widetilde{\psi}:\widetilde{B}\otimes_{\widetilde{A}}\widetilde{N}\to\widetilde{N}$ thỏa điều kiện tương tự, do đó xác định một cấu trúc \emph{$\widetilde{B}$-môđun} trên $\widetilde{N}$.

Cũng theo cách ấy, ta thấy rằng, nếu $u:B\to B'$ (tương ứng $v:N\to N'$) là một đồng cấu các $A$-đại số (tương ứng các $B$-môđun), thì $\widetilde{u}$ (tương ứng $\widetilde{v}$) là một đồng cấu các $\widetilde{A}$-đại số (tương ứng các $\widetilde{B}$-môđun), và $\Ker\widetilde{u}$ là một $\widetilde{B}$-iđêan (tương ứng $\Ker\widetilde{v}$, $\Coker\widetilde{v}$ và $\Im\widetilde{v}$ là các $\widetilde{B}$-môđun).
Nếu $N$ là một $B$-môđun, thì $\widetilde{N}$ là một $\widetilde{B}$-môđun hữu hạn sinh khi và chỉ khi $N$ là một $B$-môđun hữu hạn sinh \sref[0]{0.5.2.3}.

\oldpage[I]{90}
Nếu $M$ và $N$ là hai $B$-môđun, thì $\widetilde{B}$-môđun $\widetilde{M}\otimes_{\widetilde{B}}\widetilde{N}$ được đồng nhất chính tắc với $(M\otimes_B N)\supertilde$;
tương tự, $\shHom_{\widetilde{B}}(\widetilde{M},\widetilde{N})$ được đồng nhất chính tắc với $(\Hom_B(M,N))\supertilde$ khi $M$ có một trình bày hữu hạn;
các chứng minh giống như trong Hệ quả~\sref{I.1.3.12}.

Nếu $\mathfrak{J}$ là một iđêan của $B$ và $N$ là một $B$-môđun, thì $(\mathfrak{J}N)\supertilde=\widetilde{\mathfrak{J}}\cdot\widetilde{N}$.

Cuối cùng, nếu $B$ là một $A$-đại số \emph{phân bậc} bởi các $A$-môđun con $B_n$ ($n\in\bb{Z}$), thì $\widetilde{A}$-đại số $\widetilde{B}$, là tổng trực tiếp của các $\widetilde{A}$-môđun $\widetilde{B_n}$ \sref{I.1.3.9}, được phân bậc bởi các $\widetilde{A}$-môđun con ấy, vì tiên đề phân bậc biểu thị rằng ảnh của đồng cấu $B_m\otimes_A B_n\to B$ nằm trong $B_{m+n}$.
Tương tự, nếu $M$ là một $B$-môđun phân bậc bởi các môđun con $M_n$, thì $\widetilde{M}$ là một $\widetilde{B}$-môđun được phân bậc bởi các $\widetilde{M_n}$.
\end{env}

\begin{env}[1.3.14]
\label{I.1.3.14}
Nếu $B$ là một $A$-đại số và $M$ là một môđun con của $B$, thì $\widetilde{A}$-đại số con của $\widetilde{B}$ sinh bởi $\widetilde{M}$ \sref[0]{0.4.1.3} là $\widetilde{A}$-đại số con $\widetilde{C}$, trong đó $C$ ký hiệu đại số con của $B$ sinh bởi $M$.
Thật vậy, $C$ là tổng của các môđun con của $B$ vốn là ảnh của các đồng cấu $\bigotimes_A^n M\to B$ ($n\geq 0$), và chỉ cần áp dụng \sref{I.1.3.9} và \sref{I.1.3.12}.
\end{env}

\subsection{Bó quasi-coherent trên một phổ nguyên tố}
\label{subsection:I.1.4}

\begin{theorem}[1.4.1]
\label{I.1.4.1}
Cho $X$ là phổ nguyên tố của một vành $A$, $V$ là một tập con mở quasi-compact của $X$ và $\sh{F}$ là một $(\sh{O}_X|V)$-môđun.
Bốn điều kiện sau là tương đương.
\begin{enumerate}
  \item[{\rm(a)}] Tồn tại một $A$-môđun $M$ sao cho $\sh{F}$ đẳng cấu với $\widetilde{M}|V$.
  \item[{\rm(b)}] Tồn tại một phủ mở hữu hạn $(V_i)$ của $V$ bởi các tập có dạng $D(f_i)$ ($f_i\in A$) nằm trong $V$, sao cho, với mỗi $i$, $\sh{F}|V_i$ đẳng cấu với một bó có dạng $\widetilde{M_i}$, trong đó $M_i$ là một $A_{f_i}$-môđun.
  \item[{\rm(c)}] Bó $\sh{F}$ là quasi-coherent \sref[0]{0.5.1.3}.
  \item[{\rm(d)}] Hai tính chất sau được thỏa mãn:
    \begin{enumerate}
      \item[{\rm(d1)}] Với mọi $f\in A$ sao cho $D(f)\subset V$, và với mọi nhát cắt $s\in\Gamma(D(f),\sh{F})$, tồn tại một số nguyên $n\geq 0$ sao cho $f^n s$ kéo dài thành một nhát cắt của $\sh{F}$ trên $V$.
      \item[{\rm(d2)}] Với mọi $f\in A$ sao cho $D(f)\subset V$ và với mọi nhát cắt $t\in\Gamma(V,\sh{F})$ mà hạn chế của $t$ lên $D(f)$ bằng $0$, tồn tại một số nguyên $n\geq 0$ sao cho $f^n t=0$.
    \end{enumerate}
\end{enumerate}
\end{theorem}
(Trong phát biểu của các điều kiện (d1) và (d2), ta đã ngầm đồng nhất $A$ và $\Gamma(\widetilde{A})$ qua Định lý~\sref{I.1.3.7}).

\begin{proof}
Việc (a) suy ra (b) là hệ quả trực tiếp của Mệnh đề~\sref{I.1.3.6} và sự kiện các $D(f_i)$ tạo thành một cơ sở của tôpô trên $X$ \sref{I.1.1.10}.
Vì mọi $A$-môđun đẳng cấu với đối hạt nhân của một đồng cấu có dạng $A^{(I)}\to A^{(J)}$, \sref{I.1.3.9} chứng minh rằng mọi bó liên kết với một $A$-môđun đều quasi-coherent;
do đó (b) suy ra (c).
Ngược lại, nếu $\sh{F}$ là quasi-coherent, mọi $x\in V$ có một lân cận dạng $D(f)\subset V$ sao cho $\sh{F}|D(f)$ đẳng cấu với đối hạt nhân của một đồng cấu $\widetilde{A_f}^{(I)}\to\widetilde{A_f}^{(J)}$, và do đó cũng đẳng cấu với bó $\widetilde{N}$ liên kết với môđun $N$, là đối hạt nhân của đồng cấu tương ứng
$A_f^{(I)}\to A_f^{(J)}$ (các Hệ quả~\sref{I.1.3.8} và \sref{I.1.3.9});
vì $V$ là quasi-compact, hiển nhiên (c) suy ra (b).

\oldpage[I]{91}
Để chứng minh (b) suy ra (d1) và (d2), trước hết giả sử $V=D(g)$ với một $g\in A$, và $\sh{F}$ đẳng cấu với bó $\widetilde{N}$ liên kết với một $A_g$-môđun $N$;
thay $X$ bởi $V$ và $A$ bởi $A_g$ \sref{I.1.3.6}, ta có thể quy về trường hợp $g=1$.
Khi đó $\Gamma(D(f),\widetilde{N})$ và $N_f$ được đồng nhất chính tắc với nhau (Mệnh đề~\sref{I.1.3.6} và Định lý~\sref{I.1.3.7}), nên một nhát cắt $s\in\Gamma(D(f),\widetilde{N})$ được đồng nhất với một phần tử có dạng $z/f^n$, trong đó $z\in N$;
nhát cắt $f^n s$ được đồng nhất với phần tử $z/1$ của $N_f$ và do đó là hạn chế lên $D(f)$ của nhát cắt của $\widetilde{N}$ trên $X$ được đồng nhất với phần tử $z\in N$;
từ đó suy ra (d1) trong trường hợp này.
Tương tự, $t\in\Gamma(X,\widetilde{N})$ được đồng nhất với một phần tử $z'\in N$, hạn chế của $t$ lên $D(f)$ được đồng nhất với ảnh $z'/1$ của $z'$ trong $N_f$, và nói ảnh này bằng không có nghĩa là tồn tại một $n\geq 0$ sao cho $f^n z'=0$ trong $N$, hay tương đương $f^n t=0$.

Để hoàn tất việc chứng minh (b) suy ra (d1) và (d2), chỉ cần thiết lập bổ đề sau.
\begin{lemma}[1.4.1.1]
\label{I.1.4.1.1}
Giả sử $V$ là hợp hữu hạn của các tập có dạng $D(g_i)$, và mọi bó $\sh{F}|D(g_i)$ cùng $\sh{F}|(D(g_i)\cap D(g_j))=\sh{F}|D(g_i g_j)$ thỏa \emph{(d1)} và \emph{(d2)};
khi đó $\sh{F}$ có hai tính chất sau:
\begin{enumerate}
  \item[{\rm(d'1)}] Với mọi $f\in A$ và với mọi nhát cắt $s\in\Gamma(D(f)\cap V,\sh{F})$, tồn tại một số nguyên $n\geq 0$ sao cho $f^n s$ kéo dài thành một nhát cắt của $\sh{F}$ trên $V$.
  \item[{\rm(d'2)}] Với mọi $f\in A$ và với mọi nhát cắt $t\in\Gamma(V,\sh{F})$ mà hạn chế của $t$ lên $D(f)\cap V$ bằng $0$, tồn tại một số nguyên $n\geq 0$ sao cho $f^n t=0$.
\end{enumerate}
\end{lemma}

Trước hết ta chứng minh (d$'$2): vì $D(f)\cap D(g_i)=D(fg_i)$, với mỗi $i$ tồn tại một số nguyên $n_i$ sao cho hạn chế của $(fg_i)^{n_i}t$ lên $D(g_i)$ bằng không:
vì ảnh của $g_i$ trong $A_{g_i}$ là khả nghịch, hạn chế của $f^{n_i}t$ lên $D(g_i)$ cũng bằng không;
lấy $n$ là số lớn nhất trong các $n_i$, ta được (d$'$2).

Để chứng minh (d$'$1), ta áp dụng (d1) cho bó $\sh{F}|D(g_i)$: tồn tại một số nguyên $n_i\geq 0$ và một nhát cắt $s_i'$ của $\sh{F}$ trên $D(g_i)$ kéo dài hạn chế của $(fg_i)^{n_i}s$ lên $D(fg_i)$;
vì ảnh của $g_i$ trong $A_{g_i}$ là khả nghịch, tồn tại một nhát cắt $s_i$ của $\sh{F}$ trên $D(g_i)$ sao cho $s_i'=g_i^{n_i}s_i$, và $s_i$ kéo dài hạn chế của $f^{n_i}s$ lên $D(fg_i)$;
ngoài ra, ta có thể giả sử mọi $n_i$ bằng cùng một số nguyên $n$.
Theo cách xây dựng, hạn chế của $s_i-s_j$ lên $D(f)\cap D(g_i)\cap D(g_j)=D(fg_i g_j)$ bằng không;
áp dụng (d2) cho bó $\sh{F}|D(g_i g_j)$, tồn tại một số nguyên $m_{ij}\geq 0$ sao cho hạn chế lên $D(g_i g_j)$ của $(fg_i g_j)^{m_{ij}}(s_i-s_j)$ bằng không;
vì ảnh của $g_i g_j$ trong $A_{g_i g_j}$ là khả nghịch, hạn chế của $f^{m_{ij}}(s_i-s_j)$ lên $D(g_i g_j)$ bằng không.
Khi đó ta có thể giả sử mọi $m_{ij}$ bằng cùng một số nguyên $m$, và do đó tồn tại một nhát cắt $s'\in\Gamma(V,\sh{F})$ kéo dài tất cả các nhát cắt $f^m s_i$;
vì vậy nhát cắt này kéo dài $f^{n+m}s$, và (d$'$1) được chứng minh.

Còn phải chứng minh (d1) và (d2) suy ra (a).
Trước hết ta chứng minh (d1) và (d2) suy ra rằng các điều kiện này được thỏa mãn cho mọi bó $\sh{F}|D(g)$, trong đó $g\in A$ sao cho $D(g)\subset V$.
Điều này hiển nhiên đối với (d1);
mặt khác, nếu $t\in\Gamma(D(g),\sh{F})$ sao cho hạn chế của nó lên $D(f)\subset D(g)$ bằng không, thì theo (d1) tồn tại một số nguyên $m\geq 0$ sao cho $g^m t$
\oldpage[I]{92}
kéo dài thành một nhát cắt $s$ của $\sh{F}$ trên $V$;
áp dụng (d2), ta thấy tồn tại một số nguyên $n\geq 0$ sao cho $f^n g^m t=0$, và vì ảnh của $g$ trong $A_g$ là khả nghịch nên $f^n t=0$.

Do đó, vì $V$ là quasi-compact, Bổ đề~\sref{I.1.4.1.1} chứng minh rằng các điều kiện (d$'$1) và (d$'$2) được thỏa mãn.
Xét $A$-môđun $M=\Gamma(V,\sh{F})$, và định nghĩa một đồng cấu các $\widetilde{A}$-môđun $u:\widetilde{M}\to j_*(\sh{F})$, trong đó $j$ là phép nhúng chính tắc $V\to X$.
Vì các $D(f)$ tạo thành một cơ sở của tôpô trên $X$, với mỗi $f\in A$ chỉ cần định nghĩa một đồng cấu $u_f:M_f\to\Gamma(D(f),j_*(\sh{F}))=\Gamma(D(f)\cap V,\sh{F})$, cùng các điều kiện tương thích thông thường \sref[0]{0.3.2.5}.
Vì ảnh chính tắc của $f$ trong $A_f$ là khả nghịch, đồng cấu hạn chế $M=\Gamma(V,\sh{F})\to\Gamma(D(f)\cap V,\sh{F})$ phân tích thành $M\to M_f\xrightarrow{u_f}\Gamma(D(f)\cap V,\sh{F})$ \sref[0]{0.1.2.4}, và việc kiểm tra các điều kiện tương thích này đối với $D(g)\subset D(f)$ là trực tiếp.
Bây giờ ta chứng minh điều kiện (d$'$1) (tương ứng (d$'$2)) kéo theo mỗi $u_f$ là toàn ánh (tương ứng đơn ánh), điều này chứng minh $u$ là \emph{song ánh}, và do đó $\sh{F}$ là hạn chế lên $V$ của một $\widetilde{A}$-môđun đẳng cấu với $\widetilde{M}$.
Nếu $s\in\Gamma(D(f)\cap V,\sh{F})$, theo (d$'$1) tồn tại một số nguyên $n\geq 0$ sao cho $f^n s$ kéo dài thành một nhát cắt $z\in M$;
khi đó $u_f(z/f^n)=s$, nên $u_f$ là toàn ánh.
Tương tự, nếu $z\in M$ sao cho $u_f(z/1)=0$, điều này có nghĩa là trên $D(f)\cap V$, hạn chế của nhát cắt $z$ bằng không;
theo (d$'$2), tồn tại một số nguyên $n\geq 0$ sao cho $f^n z=0$, do đó $z/1=0$ trong $M_f$, và vì vậy $u_f$ là đơn ánh.
\end{proof}

\begin{corollary}[1.4.2]
\label{I.1.4.2}
Mọi bó quasi-coherent trên một tập mở quasi-compact của $X$ đều là hạn chế của một bó quasi-coherent trên $X$.
\end{corollary}

\begin{corollary}[1.4.3]
\label{I.1.4.3}
Mọi $\sh{O}_X$-đại số quasi-coherent trên $X=\Spec(A)$ đều đẳng cấu với một $\sh{O}_X$-đại số có dạng $\widetilde{B}$, trong đó $B$ là một đại số trên $A$;
mọi $\widetilde{B}$-môđun quasi-coherent đều đẳng cấu với một $\widetilde{B}$-môđun có dạng $\widetilde{N}$, trong đó $N$ là một $B$-môđun.
\end{corollary}

\begin{proof}
Thật vậy, một $\sh{O}_X$-đại số quasi-coherent là một $\sh{O}_X$-môđun quasi-coherent, do đó có dạng $\widetilde{B}$, trong đó $B$ là một $A$-môđun;
việc $B$ là một $A$-đại số suy ra từ đặc trưng của cấu trúc $\sh{O}_X$-đại số bằng đồng cấu $\widetilde{B}\otimes_{\widetilde{A}}\widetilde{B}\to\widetilde{B}$ của các $\widetilde{A}$-môđun, cùng với Hệ quả~\sref{I.1.3.12}.
Nếu $\sh{G}$ là một $\widetilde{B}$-môđun quasi-coherent, chỉ cần chứng minh rằng nó cũng là một $\widetilde{A}$-môđun quasi-coherent rồi kết luận theo cùng cách;
vì vấn đề có tính địa phương, bằng cách hạn chế lên một tập mở của $X$ có dạng $D(f)$, ta có thể giả sử $\sh{G}$ là đối hạt nhân của một đồng cấu $\widetilde{B}^{(I)}\to\widetilde{B}^{(J)}$ các $\widetilde{B}$-môđun (và \emph{a fortiori} các $\widetilde{A}$-môđun);
mệnh đề khi đó suy ra từ các Hệ quả~\sref{I.1.3.8} và \sref{I.1.3.9}.
\end{proof}

\subsection{Các bó coherent trên một phổ nguyên tố}
\label{subsection:I.1.5}

\begin{theorem}[1.5.1]
\label{I.1.5.1}
Cho $A$ là một vành \emph{Noether}, $X=\Spec(A)$ là phổ nguyên tố của nó, $V$ là một tập mở của $X$ và $\sh{F}$ là một $(\sh{O}_X|V)$-môđun.
Các điều kiện sau là tương đương.
\begin{enumerate}
  \item[{\rm(a)}] $\sh{F}$ là coherent.
  \item[{\rm(b)}] $\sh{F}$ hữu hạn sinh và quasi-coherent.
  \item[{\rm(c)}] Tồn tại một $A$-môđun hữu hạn sinh $M$ sao cho $\sh{F}$ đẳng cấu với bó $\widetilde{M}|V$.
\end{enumerate}
\end{theorem}

\begin{proof}
\oldpage[I]{93}
(a) hiển nhiên kéo theo (b).
Để thấy rằng (b) kéo theo (c), trước hết nhận xét rằng vì $V$ là quasi-compact \sref[0]{0.2.2.3}, $\sh{F}$ đẳng cấu với một bó $\widetilde{N}|V$, trong đó $N$ là một $A$-môđun \sref{I.1.4.1}.
Ta có $N=\varinjlim M_\lambda$, trong đó $M_\lambda$ chạy qua tập các $A$-môđun con hữu hạn sinh của $N$; do đó, theo \sref{I.1.3.9}, $\sh{F}=\widetilde{N}|V=\varinjlim\widetilde{M_\lambda}|V$;
nhưng vì $\sh{F}$ hữu hạn sinh và $V$ là quasi-compact, tồn tại một chỉ số $\lambda$ sao cho $\sh{F}=\widetilde{M_\lambda}|V$ \sref[0]{0.5.2.3}.

Cuối cùng, ta chứng minh rằng (c) kéo theo (a).
Rõ ràng khi đó $\sh{F}$ hữu hạn sinh (\sref{I.1.3.6} và \sref{I.1.3.9});
hơn nữa, vì vấn đề có tính địa phương, ta có thể giới hạn vào trường hợp $V=D(f)$, $f\in A$.
Vì $A_f$ là Noether, chỉ cần chứng minh rằng hạt nhân của một đồng cấu $\widetilde{A^n}\to\widetilde{M}$, trong đó $M$ là một $A$-môđun, là hữu hạn sinh.
Nhưng một đồng cấu như vậy có dạng $\widetilde{u}$, trong đó $u$ là một đồng cấu $A^n\to M$ \sref{I.1.3.8}, và nếu $P=\Ker u$ thì $\widetilde{P}=\Ker\widetilde{u}$ \sref{I.1.3.9}.
Vì $A$ là Noether, $P$ hữu hạn sinh, và như vậy chứng minh hoàn tất.
\end{proof}

\begin{corollary}[1.5.2]
\label{I.1.5.2}
Dưới các giả thiết của \sref{I.1.5.1}, bó $\sh{O}_X$ là một bó coherent các vành.
\end{corollary}

\begin{corollary}[1.5.3]
\label{I.1.5.3}
Dưới các giả thiết của \sref{I.1.5.1}, mọi bó coherent trên một tập mở của $X$ đều là hạn chế của một bó coherent trên $X$.
\end{corollary}

\begin{corollary}[1.5.4]
\label{I.1.5.4}
Dưới các giả thiết của \sref{I.1.5.1}, mọi $\sh{O}_X$-môđun quasi-coherent $\sh{F}$ là giới hạn quy nạp của các $\sh{O}_X$-môđun con coherent của $\sh{F}$.
\end{corollary}

\begin{proof}
Thật vậy, $\sh{F}=\widetilde{M}$, trong đó $M$ là một $A$-môđun, và $M$ là giới hạn quy nạp của các môđun con hữu hạn sinh của nó;
chứng minh hoàn tất bằng cách áp dụng \sref{I.1.3.9} và \sref{I.1.5.1}.
\end{proof}

\subsection{Các tính chất hàm tử của các bó quasi-coherent trên một phổ nguyên tố}
\label{subsection:I.1.6}

\begin{env}[1.6.1]
\label{I.1.6.1}
Cho $A$, $A'$ là hai vành,
\[
  \vphi:A'\to A
\]
là một đồng cấu, và
\[
  {}^a\vphi:X=\Spec(A)\to X'=\Spec(A')
\]
là ánh xạ liên tục liên kết với $\vphi$ \sref{I.1.2.1}.
Ta sẽ định nghĩa một \emph{đồng cấu chính tắc}
\[
  \widetilde{\vphi}:\sh{O}_{X'}\to{}^a\vphi_*(\sh{O}_X)
\]
của các bó vành.
Với mọi $f'\in A'$, đặt $f=\vphi(f')$;
ta có ${}^a\vphi^{-1}(D(f'))=D(f)$ \sref{I.1.2.2.2}.
Các vành $\Gamma(D(f'),\widetilde{A'})$ và $\Gamma(D(f),\widetilde{A})$ lần lượt được đồng nhất với $A_{f'}'$ và $A_f$ (\sref{I.1.3.6} và \sref{I.1.3.7}).
Đồng cấu $\vphi$ xác định một cách chính tắc một đồng cấu $\vphi_{f'}:A_{f'}'\to A_f$ \sref[0]{0.1.5.1}; nói cách khác, ta có một đồng cấu các vành
\[
  \Gamma(D(f'),\widetilde{A'})\to\Gamma({}^a\vphi^{-1}(D(f')),\widetilde{A})
  =\Gamma(D(f'),{}^a\vphi_*(\widetilde{A})).
\]
\oldpage[I]{94}
Hơn nữa, các đồng cấu này thỏa các điều kiện tương thích thông thường: đối với $D(f')\supset D(g')$, biểu đồ
\[
  \xymatrix{
    \Gamma(D(f'),\widetilde{A'})\ar[r]\ar[d] &
    \Gamma(D(f'),{}^a\vphi_*(\widetilde{A}))\ar[d]\\
    \Gamma(D(g'),\widetilde{A'})\ar[r] &
    \Gamma(D(g'),{}^a\vphi_*(\widetilde{A}))
  }
\]
là giao hoán \sref[0]{0.1.5.1};
như vậy ta đã định nghĩa một đồng cấu các $\sh{O}_{X'}$-đại số, vì các $D(f')$ tạo thành một cơ sở của tôpô trên $X'$ \sref[0]{0.3.2.3}.
Cặp $\Phi=({}^a\vphi,\widetilde{\vphi})$ do đó là một \emph{cấu xạ} của các không gian vành
\[
  \Phi:(X,\sh{O}_X)\to(X',\sh{O}_{X'}),
\]
\sref[0]{0.4.1.1}.

Ngoài ra, lưu ý rằng nếu đặt $x'={}^a\vphi(x)$, thì đồng cấu $\widetilde{\vphi}_x^\sharp$ \sref[0]{0.3.7.1} chính là đồng cấu
\[
  \vphi_x:A_{x'}'\to A_x
\]
được suy ra một cách chính tắc từ $\vphi:A'\to A$ \sref[0]{0.1.5.1}.
Thật vậy, mọi $z'\in A_{x'}'$ đều có thể viết dưới dạng $g'/f'$, trong đó $f'$, $g'$ thuộc $A'$ và $f'\not\in\mathfrak{j}_{x'}$;
$D(f')$ do đó là một lân cận của $x'$ trong $X'$, và đồng cấu $\Gamma(D(f'),\widetilde{A'})\to\Gamma({}^a\vphi^{-1}(D(f')),\widetilde{A})$ suy ra từ $\widetilde{\vphi}$ chính là $\vphi_{f'}$;
xét nhát cắt $s'\in\Gamma(D(f'),\widetilde{A'})$ tương ứng với $g'/f'\in A_{f'}'$, ta thu được $\widetilde{\vphi}_x^\sharp(z')=\vphi(g')/\vphi(f')$ trong $A_x$.
\end{env}

\begin{example}[1.6.2]
\label{I.1.6.2}
Cho $S$ là một tập nhân của $A$ và $\vphi$ là đồng cấu chính tắc $A\to S^{-1}A$;
ta đã thấy trong \sref{I.1.2.6} rằng ${}^a\vphi$ là một \emph{phép đồng phôi} từ $Y=\Spec(S^{-1}A)$ lên không gian con của $X=\Spec(A)$ gồm các $x$ sao cho $\mathfrak{j}_x\cap S=\emp$.
Hơn nữa, với mọi $x$ thuộc không gian con này, do đó có dạng ${}^a\vphi(y)$ với $y\in Y$, đồng cấu $\widetilde{\vphi}_y^\sharp:\sh{O}_x\to\sh{O}_y$ là
\emph{song ánh} \sref[0]{0.1.2.6};
nói cách khác, $\sh{O}_Y$ được đồng nhất với bó trên $Y$ thu được bằng cách hạn chế $\sh{O}_X$.
\end{example}

\begin{proposition}[1.6.3]
\label{I.1.6.3}
Với mọi $A$-môđun $M$, tồn tại một đẳng cấu chính tắc có tính hàm tử từ
$\sh{O}_{X'}$-môđun $(M_{[\vphi]})\supertilde$ lên ảnh trực tiếp
$\Phi_*(\widetilde{M})$.
\end{proposition}

\begin{proof}
Để viết gọn, đặt $M'=M_{[\vphi]}$ và, với mọi $f'\in A'$, cũng đặt
$f=\vphi(f')$.
Các môđun nhát cắt $\Gamma(D(f'),\widetilde{M'})$ và
$\Gamma(D(f),\widetilde{M})$ lần lượt được đồng nhất với các môđun $M_{f'}'$
và $M_f$ (trên $A_{f'}'$ và $A_f$, tương ứng); hơn nữa,
$A_{f'}'$-môđun $(M_f)_{[\vphi_{f'}]}$ đẳng cấu chính tắc với $M_{f'}'$
\sref[0]{0.1.5.2}.
Do đó ta có một đẳng cấu có tính hàm tử của các
$\Gamma(D(f'),\widetilde{A'})$-môđun:
$\Gamma(D(f'),\widetilde{M'})\isoto
\Gamma({}^a\vphi^{-1}(D(f')),\widetilde{M})_{[\vphi_{f'}]}$, và các đẳng cấu
này thỏa các điều kiện tương thích thông thường đối với các phép hạn chế
\sref[0]{0.1.5.6}, nên chúng xác định đẳng cấu có tính hàm tử đã nêu.
Cụ thể hơn, nếu $u:M_1\to M_2$ là một đồng cấu các $A$-môđun, thì nó có thể
được xem như một đồng cấu $(M_1)_{[\vphi]}\to(M_2)_{[\vphi]}$ của các
$A'$-môđun; nếu ký hiệu đồng cấu này bởi $u_{[\vphi]}$, thì
$\Phi_*(\widetilde{u})$ được đồng nhất với $(u_{[\vphi]})\supertilde$.
\end{proof}

Chứng minh này cũng cho thấy rằng, với mọi \emph{$A$-đại số $B$}, đẳng cấu
chính tắc có tính hàm tử
\oldpage[I]{95}
$(B_{[\vphi]})\supertilde\isoto\Phi_*(\widetilde{B})$ là một đẳng cấu các
\emph{$\sh{O}_{X'}$-đại số}; nếu $M$ là một $B$-môđun, thì đẳng cấu chính tắc
có tính hàm tử
$(M_{[\vphi]})\supertilde\isoto\Phi_*(\widetilde{M})$ là một đẳng cấu các
$\Phi_*(\widetilde{B})$-môđun.

\begin{corollary}[1.6.4]
\label{I.1.6.4}
Hàm tử ảnh trực tiếp $\Phi_*$ là khớp trên phạm trù các
$\sh{O}_X$-môđun quasi-coherent.
\end{corollary}

\begin{proof}
Thật vậy, rõ ràng $M_{[\vphi]}$ là một hàm tử khớp theo $M$ và
$\widetilde{M'}$ là một hàm tử khớp theo $M'$ \sref{I.1.3.5}.
\end{proof}

\begin{proposition}[1.6.5]
\label{I.1.6.5}
Cho $N'$ là một $A'$-môđun và $N$ là $A$-môđun
$N'\otimes_{A'}A_{[\vphi]}$; khi đó tồn tại một đẳng cấu chính tắc có tính
hàm tử từ $\sh{O}_X$-môđun $\Phi^*(\widetilde{N'})$ lên $\widetilde{N}$.
\end{proposition}

\begin{proof}
Trước hết lưu ý rằng $j:z'\mapsto z'\otimes 1$ là một $A'$-đồng cấu từ
$N'$ vào $N_{[\vphi]}$: thật vậy, theo định nghĩa, với $f'\in A'$, ta có
$(f' z')\otimes 1=z'\otimes\vphi(f')=\vphi(f')(z'\otimes 1)$.
Từ \sref{I.1.3.8} suy ra một đồng cấu
$\widetilde{j}:\widetilde{N'}\to(N_{[\vphi]})\supertilde$ của các
$\sh{O}_{X'}$-môđun và, theo \sref{I.1.6.3}, ta có thể xem $\widetilde{j}$
như ánh xạ $\widetilde{N'}$ vào $\Phi_*(\widetilde{N})$.
Đồng cấu $\widetilde{j}$ này tương ứng một cách chính tắc với một đồng cấu
$h=\widetilde{j}^\sharp$ từ $\Phi^*(\widetilde{N'})$ vào $\widetilde{N}$
\sref[0]{0.4.4.3}; ta sẽ thấy rằng, tại mỗi cuống, $h_x$ là \emph{song ánh}.
Đặt $x'={}^a\vphi(x)$ và lấy $f'\in A'$ sao cho $x'\in D(f')$; đặt
$f=\vphi(f')$.
Vành $\Gamma(D(f),\widetilde{A})$ được đồng nhất với $A_f$, còn các môđun
$\Gamma(D(f),\widetilde{N})$ và $\Gamma(D(f'),\widetilde{N'})$ lần lượt
được đồng nhất với $N_f$ và $N'_{f'}$; lấy một nhát cắt
$s'\in\Gamma(D(f'),\widetilde{N'})$, được đồng nhất với $n'/{f'}^p$
($n'\in N'$), và gọi $s$ là ảnh của nó qua $\widetilde{j}$ trong
$\Gamma(D(f),\widetilde{N})$; $s$ được đồng nhất với $(n'\otimes 1)/f^p$.
Mặt khác, lấy một nhát cắt $t\in\Gamma(D(f),\widetilde{A})$, được đồng nhất
với $g/f^q$ ($g\in A$); khi đó, theo định nghĩa,
$h_x(s'_{x'}\otimes t_x)=t_x\cdot s_x$ \sref[0]{0.4.4.3}.
Nhưng ta có thể đồng nhất một cách chính tắc $N_f$ với
$N'_{f'}\otimes_{A'_{f'}}(A_f)_{[\vphi_{f'}]}$ \sref[0]{0.1.5.4}; khi đó
$s$ tương ứng với phần tử $(n'/{f'}^p)\otimes 1$, và nhát cắt
$y\mapsto t_y\cdot s_y$ tương ứng với $(n'/{f'}^p)\otimes(g/f^q)$.
Các biểu đồ tương thích của \sref[0]{0.1.5.6} cho thấy $h_x$ chính là đẳng
cấu chính tắc
\[
  \label{I.1.6.5.1}
  N'_{x'}\otimes_{A'_{x'}}(A_x)_{[\vphi_x]}\isoto
  N_x=(N'\otimes_{A'}A_{[\vphi]})_x.
  \tag{1.6.5.1}
\]

Hơn nữa, cho $v:N_1'\to N_2'$ là một đồng cấu các $A'$-môđun; vì
$\widetilde{v}_{x'}=v_{x'}$ với mọi $x'\in X'$, từ những điều trên suy ra
ngay rằng $\Phi^*(\widetilde{v})$ được đồng nhất một cách chính tắc với
$(v\otimes 1)\supertilde$, và điều này hoàn tất chứng minh của
\sref{I.1.6.5}.
\end{proof}

Nếu $B'$ là một $A'$-đại số, thì đẳng cấu chính tắc từ
$\Phi^*(\widetilde{B'})$ lên
$(B'\otimes_{A'}A_{[\vphi]})\supertilde$ là một đẳng cấu các
$\sh{O}_X$-đại số; nếu thêm nữa $N'$ là một $B'$-môđun, thì đẳng cấu chính
tắc từ $\Phi^*(\widetilde{N'})$ lên
$(N'\otimes_{A'}A_{[\vphi]})\supertilde$ là một đẳng cấu các
$\Phi^*(\widetilde{B'})$-môđun.

\begin{corollary}[1.6.6]
\label{I.1.6.6}
Các nhát cắt của $\Phi^*(\widetilde{N'})$ thu được như ảnh chính tắc của các
nhát cắt $s'$, khi $s'$ chạy khắp $A'$-môđun $\Gamma(\widetilde{N'})$, sinh
ra $A$-môđun $\Gamma(\Phi^*(\widetilde{N'}))$.
\end{corollary}

\begin{proof}
Thật vậy, các ảnh này được đồng nhất với các phần tử $z'\otimes 1$ của $N$,
khi $N'$ và $N$ được đồng nhất lần lượt với $\Gamma(\widetilde{N'})$ và
$\Gamma(\widetilde{N})$ \sref{I.1.3.7}, và khi $z'$ chạy khắp $N'$.
\end{proof}

\begin{env}[1.6.7]
\label{I.1.6.7}
Trong chứng minh của \sref{I.1.6.5}, ta cũng đã chứng minh rằng ánh xạ chính
tắc (\textbf{0},~4.4.3.2)
$\rho:\widetilde{N'}\to\Phi_*(\Phi^*(\widetilde{N'}))$ chính là đồng cấu
$\widetilde{j}$,
\oldpage[I]{96}
trong đó $j:N'\to N'\otimes_{A'}A_{[\vphi]}$ là đồng cấu
$z'\mapsto z'\otimes 1$.
Tương tự, ánh xạ chính tắc (\textbf{0},~4.4.3.3)
$\sigma:\Phi^*(\Phi_*(\widetilde{M}))\to\widetilde{M}$ chính là
$\widetilde{p}$, trong đó
$p:M_{[\vphi]}\otimes_{A'}A_{[\vphi]}\to M$ là đồng cấu chính tắc gửi mỗi
tích tenxơ $z\otimes a$ ($z\in M$, $a\in A$) tới phần tử $a\cdot z$; điều
này suy ra ngay từ các định nghĩa (\sref[0]{0.3.7.1},
\sref[0]{0.4.4.3} và \sref{I.1.3.7}).

Suy ra (\textbf{0},~\EGAUnavailableHyperref{0.4.4.3}{4.4.3} và
\EGAUnavailableHyperref{0.3.5.4.4}{3.5.4.4}) rằng, nếu
$v:N'\to M_{[\vphi]}$ là một $A'$-đồng cấu, thì
$\widetilde{v}^\sharp=(v\otimes 1)\supertilde$.
\end{env}

\begin{env}[1.6.8]
\label{I.1.6.8}
Cho $N_1'$ và $N_2'$ là hai $A'$-môđun, và giả sử $N_1'$ có một
\emph{trình bày hữu hạn}; từ \sref{I.1.6.7} và \sref{I.1.3.12}[ii] suy ra
rằng đồng cấu chính tắc \sref[0]{0.4.4.6}
\[
  \Phi^*(\shHom_{\widetilde{A'}}(\widetilde{N_1'},\widetilde{N_2'}))
  \to\shHom_{\widetilde{A}}(\Phi^*(\widetilde{N_1'}),\Phi^*(\widetilde{N_2'}))
\]
chính là $\widetilde{\gamma}$, trong đó $\gamma$ ký hiệu đồng cấu chính tắc
của các $A$-môđun
$\Hom_{A'}(N_1',N_2')\otimes_{A'}A\to
\Hom_A(N_1'\otimes_{A'}A,N_2'\otimes_{A'}A)$.
\end{env}

\begin{env}[1.6.9]
\label{I.1.6.9}
Cho $\mathfrak{J}'$ là một iđêan của $A'$ và $M$ là một $A$-môđun; vì theo
định nghĩa, $\widetilde{\mathfrak{J}'}\widetilde{M}$ là ảnh của đồng cấu
chính tắc
$\Phi^*(\widetilde{\mathfrak{J}'})\otimes_{\widetilde{A}}\widetilde{M}
\to\widetilde{M}$, từ Mệnh đề~\sref{I.1.6.5} và
Hệ quả~\sref{I.1.3.12}[i] suy ra rằng
$\widetilde{\mathfrak{J}'}\widetilde{M}$ được đồng nhất một cách chính tắc với
$(\mathfrak{J}' M)\supertilde$; đặc biệt,
$\Phi^*(\widetilde{\mathfrak{J}'})\widetilde{A}$ được đồng nhất với
$(\mathfrak{J}' A)\supertilde$, và khi xét đến tính khớp phải của hàm tử
$\Phi^*$, $\widetilde{A}$-đại số
$\Phi^*((A'/\mathfrak{J}')\supertilde)$ được đồng nhất với
$(A/\mathfrak{J}' A)\supertilde$.
\end{env}

\begin{env}[1.6.10]
\label{I.1.6.10}
Cho $A''$ là một vành thứ ba, $\vphi'$ là một đồng cấu $A''\to A'$, và đặt
$\vphi''=\vphi\circ\vphi'$.
Từ các định nghĩa suy ra ngay rằng
${}^a\vphi''=({}^a\vphi')\circ({}^a\vphi)$ và
$\widetilde{\vphi''}=\widetilde{\vphi}\circ\widetilde{\vphi'}$
\sref[0]{0.1.5.7}.
Suy ra $\Phi''=\Phi'\circ\Phi$; nói cách khác,
$(\Spec(A),\widetilde{A})$ là một \emph{hàm tử} theo $A$ từ phạm trù các
vành đến phạm trù các không gian vành.
\end{env}

\subsection{Đặc trưng các cấu xạ giữa các lược đồ affine}
\label{subsection:I.1.7}

\begin{definition}[1.7.1]
\label{I.1.7.1}
Ta nói rằng một không gian vành $(X,\sh{O}_X)$ là một
\emph{lược đồ affine} nếu nó đẳng cấu với một không gian vành có dạng
$(\Spec(A),\widetilde{A})$, trong đó $A$ là một vành; khi đó
$\Gamma(X,\sh{O}_X)$, vốn được đồng nhất chính tắc với vành $A$
\sref{I.1.3.7}, là vành của lược đồ affine $(X,\sh{O}_X)$, và ta ký hiệu
nó bởi $A(X)$ khi không có nguy cơ nhầm lẫn.
\end{definition}

Theo lối lạm dụng ngôn ngữ, khi nói đến \emph{lược đồ affine}
$\Spec(A)$, ta luôn hiểu đó là không gian vành
$(\Spec(A),\widetilde{A})$.

\begin{env}[1.7.2]
\label{I.1.7.2}
Cho $A$, $B$ là hai vành, và cho $(X,\sh{O}_X)$, $(Y,\sh{O}_Y)$ là các
lược đồ affine tương ứng với các phổ nguyên tố $X=\Spec(A)$,
$Y=\Spec(B)$.
Ta đã thấy trong \sref{I.1.6.1} rằng ứng với mỗi đồng cấu vành
$\vphi:B\to A$ là một cấu xạ
$\Phi=({}^a\vphi,\widetilde{\vphi})=\Spec(\vphi):
(X,\sh{O}_X)\to(Y,\sh{O}_Y)$.
Cần lưu ý rằng $\vphi$ được $\Phi$ xác định hoàn toàn, vì theo định nghĩa,
$
  \vphi=\Gamma(\widetilde{\vphi}):
  \Gamma(\widetilde{B})\to
  \Gamma({}^a\vphi_*(\widetilde{A}))=\Gamma(\widetilde{A})
$.
\end{env}

\begin{theorem}[1.7.3]
\label{I.1.7.3}
\footnote{\emph{[Dịch.] Xem \hyperref[subsection:I.1.8]{(1.8)} và chú thích
chân trang tương ứng.}}
Cho $(X,\sh{O}_X)$ và $(Y,\sh{O}_Y)$ là hai lược đồ affine.
Để một cấu xạ của các không gian vành
$(\psi,\theta):(X,\sh{O}_X)\to(Y,\sh{O}_Y)$ có dạng
$({}^a\vphi,\widetilde{\vphi})$, trong đó $\vphi$ là một đồng cấu vành
$A(Y)\to A(X)$, điều kiện cần và đủ là, với mọi $x\in X$,
$\theta_x^\sharp$ là một đồng cấu địa phương:
$\sh{O}_{\psi(x)}\to\sh{O}_x$.
\end{theorem}

\begin{proof}
\oldpage[I]{97}
Đặt $A=A(X)$, $B=A(Y)$.
Điều kiện là cần, vì ta đã thấy trong \sref{I.1.6.1} rằng
$\widetilde{\vphi}_x^\sharp$ là đồng cấu từ
$B_{{}^a\vphi(x)}$ vào $A_x$ cảm sinh chính tắc bởi $\vphi$, và theo
định nghĩa, ${}^a\vphi(x)=\vphi^{-1}(\mathfrak{j}_x)$; do đó đồng cấu
này là địa phương.

Bây giờ ta chứng minh điều kiện là đủ.
Theo định nghĩa, $\theta$ là một đồng cấu
$\sh{O}_Y\to\psi_*(\sh{O}_X)$, và từ đó suy ra chính tắc một đồng cấu vành
\[
  \vphi=\Gamma(\theta):B=\Gamma(Y,\sh{O}_Y)
  \to\Gamma(Y,\psi_*(\sh{O}_X))=\Gamma(X,\sh{O}_X)=A.
\]

Giả thiết về $\theta_x^\sharp$ cho phép, bằng cách chuyển sang các vành
thương, suy ra từ đồng cấu này một đơn cấu $\theta^x$ từ trường thặng dư
$\kres(\psi(x))$ vào trường thặng dư $\kres(x)$, sao cho, với mọi nhát cắt
$f\in\Gamma(Y,\sh{O}_Y)=B$, ta có
$\theta^x(f(\psi(x)))=\vphi(f)(x)$.
Do đó quan hệ $f(\psi(x))=0$ tương đương với
$\vphi(f)(x)=0$, nghĩa là
$\mathfrak{j}_{\psi(x)}=\mathfrak{j}_{{}^a\vphi(x)}$; nói cách khác,
$\psi(x)={}^a\vphi(x)$ với mọi $x\in X$, hay
$\psi={}^a\vphi$.
Ta cũng biết rằng biểu đồ
\[
  \xymatrix{
    B=\Gamma(Y,\sh{O}_Y)\ar[r]^\vphi\ar[d] &
    \Gamma(X,\sh{O}_X)=A\ar[d]\\
    B_{\psi(x)}\ar[r]_{\theta_x^\sharp} &
    A_x
  }
\]
là giao hoán \sref[0]{0.3.7.2}, nghĩa là $\theta_x^\sharp$ bằng đồng cấu
$\vphi_x:B_{\psi(x)}\to A_x$ cảm sinh chính tắc bởi $\vphi$
\sref[0]{0.1.5.1}.
Vì các đồng cấu $\theta_x^\sharp$ xác định hoàn toàn $\theta^\sharp$, và
do đó cả $\theta$ \sref[0]{0.3.7.1}, ta kết luận rằng
$\theta=\widetilde{\vphi}$, theo định nghĩa của $\widetilde{\vphi}$
\sref{I.1.6.1}.
\end{proof}

Ta sẽ gọi một cấu xạ $(\psi,\theta)$ của các không gian vành thỏa điều kiện
của \sref{I.1.7.3} là một \emph{cấu xạ của các lược đồ affine}.

\begin{corollary}[1.7.4]
\label{I.1.7.4}
Nếu $(X,\sh{O}_X)$ và $(Y,\sh{O}_Y)$ là các lược đồ affine, thì tồn tại
một đẳng cấu chính tắc từ tập các cấu xạ của các lược đồ affine
$\Hom((X,\sh{O}_X),(Y,\sh{O}_Y))$ lên tập các đồng cấu vành từ $B$ vào
$A$, trong đó $A=\Gamma(\sh{O}_X)$ và $B=\Gamma(\sh{O}_Y)$.
\end{corollary}

Cũng có thể nói rằng các hàm tử
$A\longmapsto(\Spec(A),\widetilde{A})$ và
$(X,\sh{O}_X)\longmapsto\Gamma(X,\sh{O}_X)$ xác định một
\emph{tương đương} giữa phạm trù các vành giao hoán và phạm trù đối của
phạm trù các lược đồ affine (T, I, 1.2).

\begin{corollary}[1.7.5]
\label{I.1.7.5}
Nếu $\vphi:B\to A$ là toàn cấu, thì cấu xạ tương ứng
$({}^a\vphi,\widetilde{\vphi})$ là một đơn cấu của các không gian vành
\emph{(xem \sref{I.4.1.7})}.
\end{corollary}

\begin{proof}
Thật vậy, ta biết rằng ${}^a\vphi$ là đơn ánh \sref{I.1.2.5}; hơn nữa, vì
$\vphi$ là toàn cấu nên, với mọi $x\in X$,
$\vphi_x^\sharp:B_{{}^a\vphi(x)}\to A_x$, thu được từ $\vphi$ bằng cách
chuyển sang các vành phân thức, cũng là toàn cấu \sref[0]{0.1.5.1}; từ đó
suy ra kết luận \sref[0]{0.4.1.1}.
\end{proof}

\subsection{Các cấu xạ từ không gian vành địa phương đến lược đồ affine}
\label{subsection:I.1.8}

\oldpage[II]{217}
Theo một nhận xét của J.~Tate, các phát biểu của Định lý~\sref{I.1.7.3}
và Mệnh đề~\sref{I.2.2.4} có thể được tổng quát hóa như
sau:\footnote{\emph{[Dịch.] Mục sau đây (I.1.8) được bổ sung trong phần
đính chính của EGA~II; vì thế số trang tạm thời thay đổi và dẫn đến
EGA~II.}}

\begin{proposition}[1.8.1]
\label{I.1.8.1}
Cho $(S,\sh{O}_S)$ là một lược đồ affine và $(X,\sh{O}_X)$ là một không
gian vành địa phương.
Khi đó tồn tại một song ánh chính tắc từ tập các đồng cấu vành
\oldpage[II]{218}
$\Gamma(S,\sh{O}_S)\to\Gamma(X,\sh{O}_X)$ lên tập các cấu xạ của các
không gian vành $(\psi,\theta):(X,\sh{O}_X)\to(S,\sh{O}_S)$ sao cho, với
mọi $x\in X$, $\theta_x^\sharp$ là một đồng cấu địa phương
$\sh{O}_{\psi(x)}\to\sh{O}_x$.
\end{proposition}

\begin{proof}
Trước hết, ta lưu ý rằng nếu $(X,\sh{O}_X)$ và $(S,\sh{O}_S)$ là hai
không gian vành bất kỳ, thì một cấu xạ $(\psi,\theta)$ từ
$(X,\sh{O}_X)$ đến $(S,\sh{O}_S)$ xác định chính tắc một đồng cấu vành
$\Gamma(\theta):\Gamma(S,\sh{O}_S)\to\Gamma(X,\sh{O}_X)$, và do đó xác
định ánh xạ thứ nhất
\[
  \label{I.1.8.1.1}
  \rho:\Hom((X,\sh{O}_X),(S,\sh{O}_S))\to\Hom(\Gamma(S,\sh{O}_S),\Gamma(X,\sh{O}_X)).
  \tag{1.8.1.1}
\]
Ngược lại, dưới các giả thiết của mệnh đề, đặt
$A=\Gamma(S,\sh{O}_S)$ và xét một đồng cấu vành
$\vphi:A\to\Gamma(X,\sh{O}_X)$.
Với mọi $x\in X$, rõ ràng tập hợp các $f\in A$ sao cho
$\vphi(f)(x)=0$ là một \emph{iđêan nguyên tố} của $A$, vì
$\sh{O}_x/\mathfrak{m}_x=\kres(x)$ là một trường;
do đó nó là một phần tử của $S=\Spec(A)$, mà ta vẫn ký hiệu là
${}^a\vphi(x)$.
Hơn nữa, với mọi $f\in A$, theo định nghĩa \sref[0]{0.5.5.2}, ta có
${}^a\vphi^{-1}(D(f))=X_f$; điều này chứng tỏ rằng ${}^a\vphi$ là một
\emph{ánh xạ liên tục} $X\to S$.
Tiếp theo, ta định nghĩa một đồng cấu
\[
  \widetilde{\vphi}:\sh{O}_S\to{}^a\vphi_*(\sh{O}_X)
\]
của các $\sh{O}_S$-môđun;
với mọi $f\in A$, ta có $\Gamma(D(f),\sh{O}_S)=A_f$
\sref{I.1.3.6};
với mọi $s\in A$, ta cho phần tử $s/f\in A_f$ tương ứng với phần tử
$(\vphi(s)|X_f)(\vphi(f)|X_f)^{-1}$ của
$\Gamma(X_f,\sh{O}_X)=\Gamma(D(f),{}^a\vphi_*(\sh{O}_X))$, và kiểm tra
ngay được (bằng cách chuyển từ $D(f)$ sang $D(fg)$) rằng phép này thực
sự xác định một đồng cấu các $\sh{O}_S$-môđun, do đó xác định một cấu xạ
$({}^a\vphi,\widetilde{\vphi})$ của các không gian vành.
Hơn nữa, với cùng các ký hiệu và đặt $y={}^a\vphi(x)$ cho gọn, ta thấy
ngay theo \sref[0]{0.3.7.1} rằng
$\widetilde{\vphi}_x^\sharp(s_y/f_y)=(\vphi(s)_x)(\vphi(f)_x)^{-1}$;
vì quan hệ $s_y\in\mathfrak{m}_y$ theo định nghĩa tương đương với
$\vphi(s)_x\in\mathfrak{m}_x$, nên $\widetilde{\vphi}_x^\sharp$ là một
đồng cấu \emph{địa phương} $\sh{O}_y\to\sh{O}_x$, và như vậy ta đã xác
định ánh xạ thứ hai
$\sigma:\Hom(\Gamma(S,\sh{O}_S),\Gamma(X,\sh{O}_X))\to\mathfrak{L}$,
trong đó $\mathfrak{L}$ là tập các cấu xạ
$(\psi,\theta):(X,\sh{O}_X)\to(S,\sh{O}_S)$ sao cho
$\theta_x^\sharp$ là đồng cấu địa phương với mọi $x\in X$.
Còn phải chứng minh rằng $\sigma$ và $\rho$ (hạn chế lên
$\mathfrak{L}$) là các ánh xạ ngược nhau;
định nghĩa của $\widetilde{\vphi}$ cho thấy ngay rằng
$\Gamma(\widetilde{\vphi})=\vphi$, và do đó $\rho\circ\sigma$ là ánh xạ
đồng nhất.
Để thấy $\sigma\circ\rho$ là ánh xạ đồng nhất, bắt đầu với một cấu xạ
$(\psi,\theta)\in\mathfrak{L}$ và đặt $\vphi=\Gamma(\theta)$;
giả thiết về $\theta_x^\sharp$ cho phép ta, bằng cách chuyển sang các
thương, suy ra từ đồng cấu này một đơn cấu
$\theta^x:\kres(\psi(x))\to\kres(x)$ sao cho, với mọi nhát cắt
$f\in A=\Gamma(S,\sh{O}_S)$, ta có
$\theta^x(f(\psi(x)))=\vphi(f)(x)$;
do đó quan hệ $f(\psi(x))=0$ tương đương với $\vphi(f)(x)=0$, điều này
đã chứng tỏ rằng ${}^a\vphi=\psi$.
Mặt khác, từ các định nghĩa suy ra rằng biểu đồ
\[
  \xymatrix{
    A\ar[r]^\vphi\ar[d] &
    \Gamma(X,\sh{O}_X)\ar[d]\\
    A_{\psi(x)}\ar[r]^{\theta_x^\sharp} &
    \sh{O}_x
  }
\]
là giao hoán, và biểu đồ tương tự trong đó $\theta_x^\sharp$ được thay
bởi $\widetilde{\vphi}_x^\sharp$ cũng vậy; từ đó
$\widetilde{\vphi}_x^\sharp=\theta_x^\sharp$ \sref[0]{0.1.2.4}, và vì
thế $\widetilde{\vphi}=\theta$.
\end{proof}

\begin{env}[1.8.2]
\label{I.1.8.2}
Khi $(X,\sh{O}_X)$ và $(Y,\sh{O}_Y)$ là các không gian vành
\emph{địa phương}, ta sẽ xét các cấu xạ
$(\psi,\theta):(X,\sh{O}_X)\to(Y,\sh{O}_Y)$ sao cho, với mọi
$x\in X$, $\theta_x^\sharp$ là một đồng cấu \emph{địa phương}
$\sh{O}_{\psi(x)}\to\sh{O}_x$.
Từ nay, khi nói
\oldpage[II]{219}
về một \emph{cấu xạ của các không gian vành địa phương}, ta luôn hiểu đó
là một cấu xạ như trên;
với định nghĩa cấu xạ này, rõ ràng các không gian vành địa phương tạo
thành một \emph{phạm trù};
đối với hai đối tượng bất kỳ $X$ và $Y$ của phạm trù này, $\Hom(X,Y)$ do
đó ký hiệu tập các cấu xạ của không gian vành địa phương từ $X$ đến $Y$
(tập được ký hiệu là $\mathfrak{L}$ trong \sref{I.1.8.1});
khi xét tập các \emph{cấu xạ của không gian vành} từ $X$ đến $Y$, ta sẽ
ký hiệu nó là $\Hom_\text{rs}(X,Y)$ để tránh mọi nhầm lẫn.
Khi ấy ánh xạ \sref{I.1.8.1.1} được viết là
\[
  \label{I.1.8.2.1}
  \rho:\Hom_\text{rs}(X,Y)\to\Hom(\Gamma(Y,\sh{O}_Y),\Gamma(X,\sh{O}_X))
  \tag{1.8.2.1}
\]
và hạn chế của nó
\[
  \label{I.1.8.2.2}
  \rho':\Hom(X,Y)\to\Hom(\Gamma(Y,\sh{O}_Y),\Gamma(X,\sh{O}_X))
  \tag{1.8.2.2}
\]
là một ánh xạ có tính \emph{hàm tử} theo $X$ và $Y$ trong phạm trù các
không gian vành địa phương.
\end{env}

\begin{corollary}[1.8.3]
\label{I.1.8.3}
Cho $(Y,\sh{O}_Y)$ là một không gian vành địa phương.
Để $Y$ là một lược đồ affine, điều kiện cần và đủ là, với mọi không gian
vành địa phương $(X,\sh{O}_X)$, ánh xạ \sref{I.1.8.2.2} là song ánh.
\end{corollary}

\begin{proof}
Mệnh đề~\sref{I.1.8.1} cho thấy điều kiện là cần.
Ngược lại, giả sử điều kiện được thỏa mãn và đặt
$A=\Gamma(Y,\sh{O}_Y)$;
từ giả thiết và \sref{I.1.8.1} suy ra rằng các hàm tử
$X\mapsto\Hom(X,Y)$ và $X\mapsto\Hom(X,\Spec(A))$, từ phạm trù các không
gian vành địa phương đến phạm trù các tập hợp, là \emph{đẳng cấu}.
Ta biết rằng điều này kéo theo sự tồn tại của một đẳng cấu chính tắc
$Y\to\Spec(A)$ (xem~\textbf{0},~8).\footnote{\emph{[Dịch.] Văn bản
tiếng Pháp in ở đây $X\to\Spec(A)$; hàm tử được biểu diễn và kết luận đòi
hỏi $Y\to\Spec(A)$.}}
\end{proof}

\begin{env}[1.8.4]
\label{I.1.8.4}
Cho $S=\Spec(A)$ là một lược đồ affine;
ký hiệu bởi $(S',A')$ không gian vành có không gian nền \emph{chỉ gồm
một điểm} và có bó cấu trúc $A'$ là bó (tất nhiên là bó đơn) trên $S'$
do vành $A$ xác định.
Cho $\pi:S\to S'$ là ánh xạ duy nhất từ $S$ đến $S'$;
mặt khác, lưu ý rằng với mọi tập mở $U$ của $S$, tồn tại một ánh xạ
chính tắc $\Gamma(S',A')=\Gamma(S,\sh{O}_S)\to\Gamma(U,\sh{O}_S)$, do đó
xác định một \emph{$\pi$-cấu xạ} $\iota:A'\to\sh{O}_S$ của các bó vành.
Như vậy, ta đã xác định chính tắc một \emph{cấu xạ các không gian vành
$i=(\pi,\iota):(S,\sh{O}_S)\to(S',A')$}.
Với mọi $A$-môđun $M$, ta ký hiệu bởi $M'$ bó đơn trên $S'$ do $M$ xác
định; đó hiển nhiên là một $A'$-môđun.
Rõ ràng $i_*(\widetilde{M})=M'$ \sref{I.1.3.7}.
\end{env}

\begin{lemma}[1.8.5]
\label{I.1.8.5}
Với các ký hiệu của \sref{I.1.8.4}, với mọi $A$-môđun $M$, đồng cấu
chính tắc có tính hàm tử của các $\sh{O}_S$-môđun
\sref[0]{0.4.4.3.3}
\[
  \label{I.1.8.5.1}
  i^*(i_*(\widetilde{M}))\to\widetilde{M}
  \tag{1.8.5.1}
\]
là một đẳng cấu.
\end{lemma}

\begin{proof}
Thật vậy, hai vế của \sref{I.1.8.5.1} đều khớp phải (vì hàm tử
$M\mapsto i_*(\widetilde{M})$ hiển nhiên là khớp) và giao hoán với các
tổng trực tiếp;
xem $M$ là môđun thương theo ảnh của một đồng cấu
$A^{(I)}\to A^{(J)}$, ta quy về chứng minh bổ đề trong trường hợp $M=A$,
và trường hợp này là hiển nhiên.
\end{proof}

\begin{corollary}[1.8.6]
\label{I.1.8.6}
Cho $(X,\sh{O}_X)$ là một không gian vành và $u:X\to S$ là một cấu xạ
của các không gian vành.
\oldpage[II]{220}
Với mọi $A$-môđun $M$, ta có (với các ký hiệu của \sref{I.1.8.4}) một
đẳng cấu chính tắc có tính hàm tử của các $\sh{O}_X$-môđun
\[
  \label{I.1.8.6.1}
  u^*(\widetilde{M})\isoto u^*(i^*(M')).
  \tag{1.8.6.1}
\]
\end{corollary}

\begin{corollary}[1.8.7]
\label{I.1.8.7}
Dưới các giả thiết của \sref{I.1.8.6}, với mọi $A$-môđun $M$ và mọi
$\sh{O}_X$-môđun $\sh{F}$, ta có một đẳng cấu chính tắc, có tính hàm tử
theo $M$ và $\sh{F}$,
\[
  \label{I.1.8.7.1}
  \Hom_{\sh{O}_S}(\widetilde{M},u_*(\sh{F}))\isoto\Hom_A(M,\Gamma(X,\sh{F})).
  \tag{1.8.7.1}
\]
\end{corollary}

\begin{proof}
Theo \sref[0]{0.4.4.3} và Bổ đề~\sref{I.1.8.5}, ta có một đẳng cấu
chính tắc của các song hàm tử
\[
  \Hom_{\sh{O}_S}(\widetilde{M},u_*(\sh{F}))\isoto\Hom_{A'}(M',i_*(u_*(\sh{F})))
\]
và rõ ràng vế thứ hai chính là $\Hom_A(M,\Gamma(X,\sh{F}))$.
Lưu ý rằng đồng cấu chính tắc \sref{I.1.8.7.1} cho mỗi đồng cấu các
$\sh{O}_S$-môđun $h:\widetilde{M}\to u_*(\sh{F})$ (nói cách khác, mỗi
$u$-cấu xạ $\widetilde{M}\to\sh{F}$) tương ứng với đồng cấu các
$A$-môđun $\Gamma(h):M\to\Gamma(S,u_*(\sh{F}))=\Gamma(X,\sh{F})$.
\end{proof}

\begin{env}[1.8.8]
\label{I.1.8.8}
Với các ký hiệu của \sref{I.1.8.4}, rõ ràng theo \sref[0]{0.4.1.1} mọi
cấu xạ các không gian vành $(\psi,\theta):X\to S'$ tương đương với việc
cho một đồng cấu vành $A\to\Gamma(X,\sh{O}_X)$.
Do đó, ta có thể giải thích Mệnh đề~\sref{I.1.8.1} như việc xác định một
song ánh chính tắc $\Hom(X,S)\isoto\Hom(X,S')$ (trong đó tất nhiên vế
thứ hai là tập các cấu xạ của không gian vành, vì nói chung $A$ không
phải là một vành địa phương).
Tổng quát hơn, nếu $X$ và $Y$ là hai không gian vành địa phương và nếu
$(Y',A')$ là không gian vành có không gian nền chỉ gồm một điểm và có bó
vành $A'$ là bó đơn do vành $\Gamma(Y,\sh{O}_Y)$ xác định, thì ta có thể
giải thích \sref{I.1.8.2.1} như một ánh xạ
\[
  \label{I.1.8.8.1}
  \rho:\Hom_\text{rs}(X,Y)\to\Hom(X,Y').
  \tag{1.8.8.1}
\]
Do đó, kết quả của Hệ quả~\sref{I.1.8.3} được giải thích như sau: trong
các không gian vành địa phương, những đối tượng là lược đồ affine được
đặc trưng bởi tính chất rằng hạn chế của $\rho$ lên
$\Hom(X,Y)$:
\[
  \label{I.1.8.8.2}
  \rho':\Hom(X,Y)\to\Hom(X,Y')
  \tag{1.8.8.2}
\]
là \emph{song ánh} với \emph{mọi} không gian vành địa phương $X$.
Trong một chương sau, ta sẽ tổng quát hóa định nghĩa này; nhờ đó, với
\emph{mọi} không gian vành $Z$ (chứ không chỉ những không gian vành có
không gian nền gồm một điểm), ta có thể gắn một không gian vành địa
phương mà ta vẫn gọi là $\Spec(Z)$;
đây sẽ là điểm xuất phát của một lý thuyết ``tương đối'' về các tiền
lược đồ trên những không gian vành tùy ý, mở rộng các kết quả của
Chương~I.
\end{env}

\begin{env}[1.8.9]
\label{I.1.8.9}
Ta có thể xem các cặp $(X,\sh{F})$, gồm một không gian vành địa phương
$X$ và một $\sh{O}_X$-môđun $\sh{F}$, là các đối tượng của một phạm trù;
một \emph{cấu xạ} trong phạm trù này là một cặp $(u,h)$ gồm một cấu xạ
của các không gian vành địa phương
\oldpage[II]{221}
$u:X\to Y$ và một $u$-cấu xạ môđun $h:\sh{G}\to\sh{F}$;
các cấu xạ này (với $(X,\sh{F})$ và $(Y,\sh{G})$ cố định) tạo thành một
tập mà ta ký hiệu là $\Hom((X,\sh{F}),(Y,\sh{G}))$;
ánh xạ $(u,h)\mapsto(\rho'(u),\Gamma(h))$ là một ánh xạ chính tắc
\[
  \label{I.1.8.9.1}
  \Hom((X,\sh{F}),(Y,\sh{G}))\to\Hom((\Gamma(Y,\sh{O}_Y),\Gamma(Y,\sh{G})),(\Gamma(X,\sh{O}_X),\Gamma(X,\sh{F})))
  \tag{1.8.9.1}
\]
có tính \emph{hàm tử} theo $(X,\sh{F})$ và $(Y,\sh{G})$, trong đó vế
thứ hai là tập các song đồng cấu tương ứng với các vành và môđun đang
xét \sref[0]{0.1.0.2}.
\end{env}

\begin{corollary}[1.8.10]
\label{I.1.8.10}
Cho $Y$ là một không gian vành địa phương và $\sh{G}$ là một
$\sh{O}_Y$-môđun. Để $Y$ là một lược đồ affine và $\sh{G}$ là một
$\sh{O}_Y$-môđun quasi-coherent, điều kiện cần và đủ là, với mọi cặp
$(X,\sh{F})$ gồm một không gian vành địa phương $X$ và một
$\sh{O}_X$-môđun $\sh{F}$, ánh xạ chính tắc
\sref{I.1.8.9.1} là song ánh.
\end{corollary}

Ta dành cho độc giả việc triển khai chứng minh, theo cùng lập luận như
trong \sref{I.1.8.3} và sử dụng \sref{I.1.8.1} cùng
\sref{I.1.8.7}.

\begin{remark}[1.8.11]
\label{I.1.8.11}
Các mệnh đề \sref{I.1.7.3}, \sref{I.1.7.4} và
\sref{I.2.2.4} là các trường hợp riêng của \sref{I.1.8.1}, cũng như
định nghĩa trong \sref{I.1.6.1}; tương tự, \sref{I.2.2.5} suy ra từ
\sref{I.1.8.7}. Hệ quả~\sref{I.1.8.7} cũng kéo theo
\sref{I.1.6.3} (và do đó \sref{I.1.6.4}) như một trường hợp riêng,
vì nếu $X$ là một lược đồ affine và $\Gamma(X,\sh{F})=N$, thì các
hàm tử $M\mapsto\Hom_{\sh{O}_S}(\widetilde{M},u_*(\widetilde{N}))$ và
$M\mapsto\Hom_{\sh{O}_S}(\widetilde{M},(N_{[\vphi]})\supertilde)$
(trong đó $\vphi:A\to\Gamma(X,\sh{O}_X)$ tương ứng với $u$) đẳng cấu
với nhau theo các Hệ quả~\sref{I.1.8.7} và \sref{I.1.3.8}.
Cuối cùng, \sref{I.1.6.5} (và do đó \sref{I.1.6.6}) suy ra từ
\sref{I.1.8.6} và từ việc, với mọi $f\in A$, hai $A_f$-môđun
$N'\otimes_{A'}A_f$ và $(N'\otimes_{A'}A)_f$ (theo ký hiệu của
\sref{I.1.6.5}) đẳng cấu chính tắc với nhau.
\end{remark}

\section{Các tiền lược đồ và cấu xạ của tiền lược đồ}
\label{section:I.2}

\subsection{Định nghĩa các tiền lược đồ}
\label{subsection:I.2.1}

\begin{env}[2.1.1]
\label{I.2.1.1}
Cho một không gian vành $(X,\sh{O}_X)$, ta nói một tập con mở
$V$ của $X$ là một \emph{tập mở affine} nếu không gian vành
$(V,\sh{O}_X|V)$ là một lược đồ affine \sref{I.1.7.1}.
\end{env}

\begin{definition}[2.1.2]
\label{I.2.1.2}
Ta gọi là \emph{tiền lược đồ} một không gian vành $(X,\sh{O}_X)$ sao
cho mỗi điểm của $X$ có một lân cận mở affine.
\end{definition}

\begin{proposition}[2.1.3]
\label{I.2.1.3}
\oldpage[I]{98}
Nếu $(X,\sh{O}_X)$ là một tiền lược đồ, thì các tập mở affine của nó
tạo thành một cơ sở của tôpô trên $X$.
\end{proposition}

\begin{proof}
Nếu $V$ là một lân cận mở tùy ý của $x\in X$, thì theo giả thiết tồn
tại một lân cận mở $W$ của $x$ sao cho $(W,\sh{O}_X|W)$ là một lược
đồ affine; ký hiệu $A$ là vành của nó. Trong không gian $W$,
$V\cap W$ là một lân cận mở của $x$; do đó tồn tại $f\in A$ sao cho
$D(f)$ là một lân cận mở của $x$, nằm trong $V\cap W$
\sref{I.1.1.10}[i]. Khi đó không gian vành
$(D(f),\sh{O}_X|D(f))$ là một lược đồ affine có vành đẳng cấu với
$A_f$ \sref{I.1.3.6}, từ đó suy ra mệnh đề.
\end{proof}

\begin{proposition}[2.1.4]
\label{I.2.1.4}
Không gian nền của một tiền lược đồ là một không gian Kolmogorov.
\end{proposition}

\begin{proof}
Nếu $x$ và $y$ là hai điểm phân biệt của một tiền lược đồ $X$, thì rõ
ràng tồn tại một lân cận mở của một trong hai điểm không chứa điểm kia
khi $x$ và $y$ không thuộc cùng một tập mở affine; còn khi chúng thuộc
cùng một tập mở affine, điều này suy ra từ \sref{I.1.1.8}.
\end{proof}

\begin{proposition}[2.1.5]
\label{I.2.1.5}
Nếu $(X,\sh{O}_X)$ là một tiền lược đồ, thì mọi tập con đóng bất khả
quy của $X$ có đúng một điểm tổng quát, và ánh xạ
$x\mapsto\overline{\{x\}}$ do đó là một song ánh từ $X$ lên tập các
tập con đóng bất khả quy của nó.
\end{proposition}

\begin{proof}
Nếu $Y$ là một tập con đóng bất khả quy của $X$, $y\in Y$, và $U$ là
một lân cận mở affine của $y$ trong $X$, thì $U\cap Y$ trù mật trong
$Y$ và cũng bất khả quy (\sref[0]{0.2.1.1} và
\sref[0]{0.2.1.4}); do đó, theo Hệ quả~\sref{I.1.1.14}, $U\cap Y$ là
bao đóng trong $U$ của một điểm $x$, và vì thế
$Y\subset\overline{U}$ là bao đóng của $x$ trong $X$. Tính duy nhất
của điểm tổng quát của $Y$ suy ra từ Mệnh đề~\sref{I.2.1.4} và
\sref[0]{0.2.1.3}.
\footnote{\emph{[Dịch.] Văn bản tiếng Pháp in ở đây ``điểm tổng quát
của $X$''; mệnh đề và câu ngay trước đều nói về điểm tổng quát của
$Y$.}}
\end{proof}

\begin{env}[2.1.6]
\label{I.2.1.6}
Nếu $Y$ là một tập con đóng bất khả quy của $X$ và $y$ là điểm tổng
quát của nó, thì vành địa phương $\sh{O}_y$ (cũng được ký hiệu bởi
$\sh{O}_{X/Y}$) được gọi là \emph{vành địa phương của $X$ dọc theo
$Y$}, hoặc \emph{vành địa phương của $Y$ trong $X$}.

Nếu bản thân $X$ bất khả quy và $x$ là điểm tổng quát của nó, thì ta
còn nói $\sh{O}_x$ là \emph{vành các hàm hữu tỉ trên $X$}
(xem~\textsection\hyperref[section:I.7]{7}).
\end{env}

\begin{proposition}[2.1.7]
\label{I.2.1.7}
Nếu $(X,\sh{O}_X)$ là một tiền lược đồ, thì với mọi tập con mở $U$
của $X$, không gian vành $(U,\sh{O}_X|U)$ là một tiền lược đồ.
\end{proposition}

\begin{proof}
Điều này suy ra ngay từ Định nghĩa~\sref{I.2.1.2} và
Mệnh đề~\sref{I.2.1.3}.
\end{proof}

Ta nói $(U,\sh{O}_X|U)$ là tiền lược đồ \emph{cảm sinh} trên $U$
bởi $(X,\sh{O}_X)$, hoặc là \emph{hạn chế} của
$(X,\sh{O}_X)$ lên $U$.

\begin{env}[2.1.8]
\label{I.2.1.8}
Ta nói một tiền lược đồ $(X,\sh{O}_X)$ là \emph{bất khả quy}
(tương ứng, \emph{liên thông}) nếu không gian nền $X$ là bất khả quy
(tương ứng, liên thông). Ta nói một tiền lược đồ là \emph{nguyên} nếu
nó \emph{bất khả quy và rút gọn} (xem \sref{I.5.1.4}). Ta nói một
tiền lược đồ $(X,\sh{O}_X)$ là \emph{nguyên địa phương} nếu mọi
$x\in X$ có một lân cận mở $U$ sao cho tiền lược đồ cảm sinh trên $U$
bởi $(X,\sh{O}_X)$ là nguyên.
\end{env}

\subsection{Các cấu xạ của tiền lược đồ}
\label{subsection:I.2.2}
\label{I.2.2}

\begin{definition}[2.2.1]
\label{I.2.2.1}
Cho hai tiền lược đồ $(X,\sh{O}_X)$ và $(Y,\sh{O}_Y)$, ta gọi một
\emph{cấu xạ (của các tiền lược đồ)} từ $(X,\sh{O}_X)$ đến
$(Y,\sh{O}_Y)$ là một cấu xạ $(\psi,\theta)$ của các không gian vành sao
cho, với mọi $x\in X$, $\theta_x^\sharp$ là một đồng cấu địa phương:
$\sh{O}_{\psi(x)}\longrightarrow\sh{O}_x$.
\end{definition}

\oldpage[I]{99}

Khi chuyển sang các vành thương, ánh xạ
$\theta_x^\sharp:\sh{O}_{\psi(x)}\to\sh{O}_x$ do đó cho một đơn cấu
$\theta^x:\kres(\psi(x))\to\kres(x)$, cho phép xem $\kres(x)$ như một
\emph{mở rộng} của trường $\kres(\psi(x))$.

\begin{env}[2.2.2]
\label{I.2.2.2}
Hợp thành $(\psi'',\theta'')$ của hai cấu xạ tiền lược đồ
$(\psi,\theta)$ và $(\psi',\theta')$ vẫn là một cấu xạ tiền lược đồ,
như suy ra từ công thức
${\theta''}^\sharp=\theta^\sharp\circ\psi^*({\theta'}^\sharp)$
\sref[0]{0.3.5.5}. Từ đó suy ra rằng các tiền lược đồ tạo thành một
\emph{phạm trù}; theo ký hiệu tổng quát, ta viết $\Hom(X,Y)$ cho tập hợp
các cấu xạ từ một tiền lược đồ $X$ đến một tiền lược đồ $Y$.
\end{env}

\begin{example}[2.2.3]
\label{I.2.2.3}
Nếu $U$ là một tập con mở của $X$, thì phép nhúng chính tắc
\sref[0]{0.4.1.2} của tiền lược đồ cảm sinh
$(U,\sh{O}_X|U)$ vào $(X,\sh{O}_X)$ là một cấu xạ của các tiền lược đồ;
hơn nữa, đó là một \emph{đơn cấu} của các không gian vành (và
\emph{a fortiori} là một đơn cấu của các tiền lược đồ), như suy ra ngay
từ \sref[0]{0.4.1.1}.
\end{example}

\begin{proposition}[2.2.4]
\label{I.2.2.4}
\footnote{\emph{[Dịch.] Xem \hyperref[subsection:I.1.8]{(1.8)} và chú thích
chân trang tương ứng.}}
Cho $(X,\sh{O}_X)$ là một tiền lược đồ và $(S,\sh{O}_S)$ là một lược
đồ affine có vành $A$. Có một tương ứng song ánh chính tắc giữa các cấu
xạ từ tiền lược đồ $(X,\sh{O}_X)$ đến tiền lược đồ $(S,\sh{O}_S)$ và
các đồng cấu từ vành $A$ vào vành $\Gamma(X,\sh{O}_X)$.
\end{proposition}

\begin{proof}
Trước hết, nhận xét rằng nếu $(X,\sh{O}_X)$ và $(Y,\sh{O}_Y)$ là hai
không gian vành bất kỳ, thì một cấu xạ $(\psi,\theta)$ từ
$(X,\sh{O}_X)$ đến $(Y,\sh{O}_Y)$ xác định chính tắc một đồng cấu vành
$\Gamma(\theta):\Gamma(Y,\sh{O}_Y)\to
\Gamma(Y,\psi_*(\sh{O}_X))=\Gamma(X,\sh{O}_X)$.
Trong trường hợp đang xét, chỉ còn phải chứng minh rằng mọi đồng cấu
$\vphi:A\to\Gamma(X,\sh{O}_X)$ đều có dạng $\Gamma(\theta)$ với một
$\theta$ duy nhất.

Theo giả thiết, tồn tại một phủ $(V_\alpha)$ của $X$ bởi các tập mở
affine. Hợp thành $\vphi$ với đồng cấu hạn chế
$\Gamma(X,\sh{O}_X)\to\Gamma(V_\alpha,\sh{O}_X|V_\alpha)$, ta được một
đồng cấu
$\vphi_\alpha:A\to\Gamma(V_\alpha,\sh{O}_X|V_\alpha)$, tương ứng với một
cấu xạ duy nhất $(\psi_\alpha,\theta_\alpha)$ từ tiền lược đồ
$(V_\alpha,\sh{O}_X|V_\alpha)$ đến tiền lược đồ $(S,\sh{O}_S)$, theo
Định lý~\sref{I.1.7.3}.
Hơn nữa, với mọi cặp chỉ số $(\alpha,\beta)$, mỗi điểm của
$V_\alpha\cap V_\beta$ đều có một lân cận mở affine $W$ nằm trong
$V_\alpha\cap V_\beta$ \sref{I.2.1.3}. Rõ ràng, khi hợp thành
$\vphi_\alpha$ và $\vphi_\beta$ với các đồng cấu hạn chế xuống $W$, ta
thu được cùng một đồng cấu
$\Gamma(S,\sh{O}_S)\to\Gamma(W,\sh{O}_X|W)$; do đó, theo các hệ thức
$(\theta_\alpha^\sharp)_x=(\vphi_\alpha)_x$ với mọi $x\in V_\alpha$ và
mọi $\alpha$ \sref{I.1.6.1}, các hạn chế xuống $W$ của các cấu xạ
$(\psi_\alpha,\theta_\alpha)$ và $(\psi_\beta,\theta_\beta)$ trùng nhau.
Suy ra tồn tại duy nhất một cấu xạ của các không gian vành
$(\psi,\theta):(X,\sh{O}_X)\to(S,\sh{O}_S)$ sao cho hạn chế của nó xuống
mỗi $V_\alpha$ là $(\psi_\alpha,\theta_\alpha)$; rõ ràng đây là một cấu
xạ của các tiền lược đồ và $\Gamma(\theta)=\vphi$.

Cho $u:A\to\Gamma(X,\sh{O}_X)$ là một đồng cấu vành, và cho
$v=(\psi,\theta)$ là cấu xạ tương ứng
$(X,\sh{O}_X)\to(S,\sh{O}_S)$. Với mọi $f\in A$, ta có
\[
  \label{I.2.2.4.1}
  \psi^{-1}(D(f))=X_{u(f)}
  \tag{2.2.4.1}
\]
với ký hiệu của \sref[0]{0.5.5.2} liên quan đến bó tự do địa phương
$\sh{O}_X$. Thật vậy, chỉ cần kiểm tra công thức này khi bản thân $X$ là
affine, và khi đó nó chính là \sref{I.1.2.2.2}.
\end{proof}

\begin{proposition}[2.2.5]
\label{I.2.2.5}
Dưới các giả thiết của Mệnh đề~\sref{I.2.2.4}, cho
$\vphi:A\to\Gamma(X,\sh{O}_X)$ là một đồng cấu vành, cho
$f:(X,\sh{O}_X)\to(S,\sh{O}_S)$ là cấu xạ tiền lược đồ tương ứng,
cho $\sh{G}$ (tương ứng, $\sh{F}$) là một $\sh{O}_X$-môđun quasi-coherent
(tương ứng, một $\sh{O}_S$-môđun quasi-coherent), và đặt
$M=\Gamma(S,\sh{F})$. Khi đó có một tương ứng song ánh chính tắc
\oldpage[I]{100}
giữa các $f$-cấu xạ $\sh{F}\to\sh{G}$ \sref[0]{0.4.4.1} và các
$A$-đồng cấu $M\to(\Gamma(X,\sh{G}))_{[\vphi]}$.
\end{proposition}

\begin{proof}
Lập luận như trong Mệnh đề~\sref{I.2.2.4}, ta lập tức quy về trường hợp
$X$ là affine, và khi đó mệnh đề suy ra từ Mệnh đề~\sref{I.1.6.3}
và Hệ quả~\sref{I.1.3.8}.
\end{proof}

\begin{env}[2.2.6]
\label{I.2.2.6}
Ta nói một cấu xạ tiền lược đồ
$(\psi,\theta):(X,\sh{O}_X)\to(Y,\sh{O}_Y)$ là \emph{mở}
(tương ứng, \emph{đóng}) nếu, với mọi tập con mở $U$ của $X$
(tương ứng, mọi tập con đóng $F$ của $X$), $\psi(U)$ là một tập mở trong $Y$
(tương ứng, $\psi(F)$ là một tập đóng trong $Y$).
Ta nói $(\psi,\theta)$ là \emph{trội} nếu $\psi(X)$ trù mật trong $Y$,
và \emph{toàn ánh} nếu $\psi$ là toàn ánh.
Lưu ý rằng các điều kiện này chỉ phụ thuộc vào ánh xạ liên tục $\psi$.
\end{env}

\begin{proposition}[2.2.7]
\label{I.2.2.7}
Cho
\[
  f=(\psi,\theta):(X,\sh{O}_X)\to(Y,\sh{O}_Y),
  \qquad
  g=(\psi',\theta'):(Y,\sh{O}_Y)\to(Z,\sh{O}_Z)
\]
là hai cấu xạ tiền lược đồ.
\begin{enumerate}
  \item[\rm(i)] Nếu $f$ và $g$ đều mở (tương ứng, đóng, trội, toàn ánh),
    thì $g\circ f$ cũng vậy.
  \item[\rm(ii)] Nếu $f$ toàn ánh và $g\circ f$ đóng, thì $g$ đóng.
  \item[\rm(iii)] Nếu $g\circ f$ toàn ánh, thì $g$ toàn ánh.
\end{enumerate}
\end{proposition}

\begin{proof}
Các khẳng định (i) và (iii) là hiển nhiên.
Đặt $g\circ f=(\psi'',\theta'')$.
Nếu $F$ là một tập con đóng của $Y$, thì $\psi^{-1}(F)$ là đóng trong $X$,
do đó $\psi''(\psi^{-1}(F))$ là đóng trong $Z$; nhưng vì $\psi$ toàn ánh,
ta có $\psi(\psi^{-1}(F))=F$, do đó
$\psi''(\psi^{-1}(F))=\psi'(F)$, điều này chứng minh (ii).
\end{proof}

\begin{proposition}[2.2.8]
\label{I.2.2.8}
Cho $f=(\psi,\theta)$ là một cấu xạ
$(X,\sh{O}_X)\to(Y,\sh{O}_Y)$ và $(U_\alpha)$ là một phủ mở của $Y$.
Để $f$ là mở (tương ứng, đóng, toàn ánh, trội), điều kiện cần và đủ là
hạn chế của nó lên mỗi tiền lược đồ cảm sinh
$(\psi^{-1}(U_\alpha),\sh{O}_X|\psi^{-1}(U_\alpha))$, được xem như một cấu xạ
từ tiền lược đồ cảm sinh này vào tiền lược đồ cảm sinh
$(U_\alpha,\sh{O}_Y|U_\alpha)$, là một cấu xạ mở
(tương ứng, đóng, toàn ánh, trội).
\end{proposition}

\begin{proof}
Mệnh đề suy ra ngay từ các định nghĩa, khi lưu ý rằng một tập con $F$ của $Y$
là đóng (tương ứng, mở, trù mật) trong $Y$ nếu và chỉ nếu mỗi tập
$F\cap U_\alpha$ là đóng (tương ứng, mở, trù mật) trong $U_\alpha$.
\end{proof}

\begin{env}[2.2.9]
\label{I.2.2.9}
Cho $(X,\sh{O}_X)$ và $(Y,\sh{O}_Y)$ là hai tiền lược đồ; giả sử $X$ và
$Y$ có cùng một số \emph{hữu hạn} các thành phần bất khả quy $X_i$
(tương ứng, $Y_i$), $(1\leq i\leq n)$; cho $\xi_i$ (tương ứng, $\eta_i$)
là điểm tổng quát của $X_i$ (tương ứng, $Y_i$) \sref{I.2.1.5}.
Ta nói một cấu xạ
\[
  f=(\psi,\theta):(X,\sh{O}_X)\to(Y,\sh{O}_Y)
\]
là \emph{song hữu tỉ} nếu, với mọi $i$,
$\psi^{-1}(\eta_i)=\{\xi_i\}$ và
$\theta_{\xi_i}^{\sharp}:\sh{O}_{\eta_i}\to\sh{O}_{\xi_i}$ là một
\emph{đẳng cấu}.
Hiển nhiên một cấu xạ song hữu tỉ là \emph{trội} \sref[0]{0.2.1.8},
do đó toàn ánh nếu nó cũng đóng.
\end{env}

\begin{notation}[2.2.10]
\label{I.2.2.10}
Trong toàn bộ phần còn lại của tác phẩm này, và khi điều đó không gây nguy cơ
nhầm lẫn, ta sẽ \emph{lược bỏ} khỏi ký hiệu của một tiền lược đồ
(tương ứng, của một cấu xạ) bó cấu trúc
(tương ứng, cấu xạ giữa các bó cấu trúc).
Nếu $U$ là một tập con mở của không gian nền của một tiền lược đồ $X$, thì
khi nói đến $U$ như một tiền lược đồ, ta luôn hiểu đó là tiền lược đồ cảm sinh
trên $U$.
\end{notation}

\subsection{Phép dán các tiền lược đồ}
\label{subsection:I.2.3}

\begin{env}[2.3.1]
\label{I.2.3.1}
\oldpage[I]{101}
Từ Định nghĩa~\sref{I.2.1.2} suy ra rằng mọi không gian vành thu được bằng
\emph{phép dán} các tiền lược đồ \sref[0]{0.4.1.6} vẫn là một tiền lược đồ.
Đặc biệt, vì theo định nghĩa mọi tiền lược đồ đều có một phủ gồm các tập mở
affine, ta thấy mọi tiền lược đồ đều có thể thu được bằng
\emph{phép dán các lược đồ affine}.
\end{env}

\begin{example}[2.3.2]
\label{I.2.3.2}
Cho $K$ là một trường, $B=K[s]$ và $C=K[t]$ là hai vành đa thức một biến
trên $K$, và cho $X_1=\Spec(B)$, $X_2=\Spec(C)$ là hai lược đồ affine đẳng
cấu. Trong $X_1$ (tương ứng $X_2$), gọi $U_{12}$ (tương ứng $U_{21}$) là
tập mở affine $D(s)$ (tương ứng $D(t)$), có vành $B_s$ (tương ứng $C_t$)
gồm các phân thức hữu tỉ dạng $f(s)/s^m$ (tương ứng $g(t)/t^n$), với
$f\in B$ (tương ứng $g\in C$). Gọi $u_{12}$ là đẳng cấu tiền lược đồ
$U_{21}\to U_{12}$ tương ứng \sref{I.2.2.4} với đẳng cấu từ $B_s$ đến
$C_t$ gửi $f(s)/s^m$ tới phân thức hữu tỉ $f(1/t)/(1/t^m)$. Ta có thể dán
$X_1$ và $X_2$ dọc theo $U_{12}$ và $U_{21}$ bằng $u_{12}$, vì hiển nhiên
không có điều kiện dán nào. Sau này ta sẽ gặp lại tiền lược đồ $X$ thu được
như vậy như một trường hợp riêng của một phương pháp xây dựng tổng quát
\sref[II]{II.2.4.3}. Ở đây, ta chỉ chứng minh rằng $X$
\emph{không phải là một lược đồ affine}; điều này suy ra từ việc vành
$\Gamma(X,\sh{O}_X)$ \emph{đẳng cấu} với $K$, nên phổ của nó chỉ gồm một
điểm. Thật vậy, hạn chế của một tiết diện của $\sh{O}_X$ trên $X$ lên
$X_1$ (tương ứng $X_2$), được đồng nhất với một tập mở affine của $X$, là
một đa thức $f(s)$ (tương ứng $g(t)$); theo các định nghĩa, ta phải có
$g(t)=f(1/t)$, điều này chỉ có thể xảy ra khi $f=g\in K$.
\end{example}

\subsection{Các lược đồ địa phương}
\label{subsection:I.2.4}

\begin{env}[2.4.1]
\label{I.2.4.1}
Gọi \emph{lược đồ địa phương} là một lược đồ affine có vành $A$ là địa
phương; khi đó trong $X=\Spec(A)$ tồn tại duy nhất một \emph{điểm đóng}
$a$, và với mọi điểm khác $b\in X$, ta có
$a\in\overline{\{b\}}$ \sref{I.1.1.7}.
\end{env}

Với mọi tiền lược đồ $Y$ và mọi điểm $y\in Y$, lược đồ địa phương
$\Spec(\sh{O}_y)$ được gọi là \emph{lược đồ địa phương của $Y$ tại điểm
$y$}. Cho $V$ là một tập mở affine của $Y$ chứa $y$, và $B$ là vành của
lược đồ affine $V$; $\sh{O}_y$ được đồng nhất chính tắc với $B_y$
\sref{I.1.3.4}, và do đó đồng cấu chính tắc $B\to B_y$ tương ứng
\sref{I.1.6.1} với một cấu xạ tiền lược đồ
$\Spec(\sh{O}_y)\to V$. Hợp thành cấu xạ này với phép nhúng chính tắc
$V\to Y$, ta thu được một cấu xạ $\Spec(\sh{O}_y)\to Y$ \emph{độc lập}
với tập mở affine $V$ (chứa $y$) đã chọn: thật vậy, nếu $V'$ là một tập mở
affine thứ hai chứa $y$, thì tồn tại một tập mở affine thứ ba $W$ chứa $y$
sao cho $W\subset V\cap V'$ \sref{I.2.1.3}; do đó có thể quy về trường hợp
$V\subset V'$, và nếu $B'$ là vành của $V'$, thì chỉ cần nhận xét rằng
biểu đồ
\phantomsection\label{I.2.4.1.diagram}
\[
  \xymatrix{
    B'\ar[rr]\ar[dr] & & B\ar[dl]\\
    & \sh{O}_y
  }
\]
là giao hoán \sref[0]{0.1.5.1}. Cấu xạ
\[
  \Spec(\sh{O}_y)\to Y
\]
được định nghĩa như vậy gọi là \emph{chính tắc}.

\begin{proposition}[2.4.2]
\label{I.2.4.2}
\oldpage[I]{102}
Cho $(Y,\sh{O}_Y)$ là một tiền lược đồ; với mọi $y\in Y$, gọi
$(\psi,\theta)$ là cấu xạ chính tắc
$(\Spec(\sh{O}_y),\widetilde{\sh{O}}_y)\to(Y,\sh{O}_Y)$. Khi đó $\psi$
là một phép đồng phôi từ $\Spec(\sh{O}_y)$ lên không gian con $S_y$ của
$Y$ gồm các $z$ sao cho $y\in\overline{\{z\}}$ (nói cách khác, các
\emph{tổng quát hóa} của $y$ \sref[0]{0.2.1.2}); hơn nữa, nếu
$z=\psi(\mathfrak{p})$, thì
$\theta_z^{\sharp}:\sh{O}_z\to(\sh{O}_y)_{\mathfrak{p}}$ là một đẳng
cấu; do đó $(\psi,\theta)$ là một đơn cấu của các không gian vành.
\end{proposition}

\begin{proof}
Vì điểm đóng duy nhất $a$ của $\Spec(\sh{O}_y)$ thuộc bao đóng của mọi
điểm trong không gian này và $\psi(a)=y$, ảnh của
$\Spec(\sh{O}_y)$ qua ánh xạ liên tục $\psi$ nằm trong $S_y$. Vì $S_y$
nằm trong mọi tập mở affine chứa $y$, ta có thể quy về trường hợp $Y$ là
một lược đồ affine; nhưng trong trường hợp này, mệnh đề suy ra từ
\sref{I.1.6.2}.
\end{proof}

\emph{Do đó, ta thấy \sref{I.2.1.5} có một tương ứng song ánh giữa
$\Spec(\sh{O}_y)$ và tập hợp các tập con đóng bất khả quy của $Y$ chứa
$y$.}

\begin{corollary}[2.4.3]
\label{I.2.4.3}
Điều kiện cần và đủ để $y\in Y$ là điểm tổng quát của một thành phần bất
khả quy của $Y$ là iđêan nguyên tố duy nhất của vành địa phương
$\sh{O}_y$ chính là iđêan cực đại của nó (nói cách khác,
$\sh{O}_y$ có chiều không).
\end{corollary}

\begin{proposition}[2.4.4]
\label{I.2.4.4}
Cho $(X,\sh{O}_X)$ là một lược đồ địa phương có vành $A$, $a$ là điểm
đóng duy nhất của nó, và $(Y,\sh{O}_Y)$ là một tiền lược đồ. Mọi cấu xạ
$u=(\psi,\theta):(X,\sh{O}_X)\to(Y,\sh{O}_Y)$ phân tích duy nhất thành
$X\to\Spec(\sh{O}_{\psi(a)})\to Y$, trong đó mũi tên thứ hai chỉ cấu xạ
chính tắc, còn mũi tên thứ nhất tương ứng với một đồng cấu địa phương
$\sh{O}_{\psi(a)}\to A$. Điều này thiết lập một tương ứng song ánh chính
tắc giữa tập các cấu xạ $(X,\sh{O}_X)\to(Y,\sh{O}_Y)$ và tập các đồng
cấu địa phương $\sh{O}_y\to A$ $(y\in Y)$.
\end{proposition}

Thật vậy, với mọi $x\in X$, ta có $a\in\overline{\{x\}}$, do đó
$\psi(a)\in\overline{\{\psi(x)\}}$; điều này chứng tỏ $\psi(X)$ nằm
trong mọi tập mở affine chứa $\psi(a)$. Do đó ta có thể quy về trường
hợp $(Y,\sh{O}_Y)$ là một lược đồ affine có vành $B$; khi đó
$u=({}^a\vphi,\widetilde{\vphi})$, với $\vphi\in\Hom(B,A)$
\sref{I.1.7.3}. Hơn nữa,
$\vphi^{-1}(\mathfrak{j}_a)=\mathfrak{j}_{\psi(a)}$, vì thế ảnh qua
$\vphi$ của mọi phần tử thuộc $B\setmin\mathfrak{j}_{\psi(a)}$ đều khả
nghịch trong vành địa phương $A$; do đó phép phân tích trong mệnh đề suy
ra từ tính chất phổ dụng của các vành phân thức
\sref[0]{0.1.2.4}. Ngược lại, với mỗi đồng cấu địa phương
$\sh{O}_y\to A$ tương ứng một cấu xạ duy nhất
$(\psi,\theta):X\to\Spec(\sh{O}_y)$ sao cho $\psi(a)=y$
\sref{I.1.7.3}; hợp cấu xạ này với cấu xạ chính tắc
$\Spec(\sh{O}_y)\to Y$, ta thu được một cấu xạ $X\to Y$, và như vậy
hoàn tất chứng minh mệnh đề.

\begin{env}[2.4.5]
\label{I.2.4.5}
Các lược đồ affine có vành là một \emph{trường} $K$ đều có không gian
nền chỉ gồm một điểm. Nếu $A$ là một vành địa phương có iđêan cực đại
$\mathfrak{m}$, thì mọi đồng cấu địa phương $A\to K$ có hạt nhân bằng
$\mathfrak{m}$ và do đó phân tích thành
$A\to A/\mathfrak{m}\to K$, trong đó mũi tên thứ hai là một đơn cấu. Vì
thế, các cấu xạ $\Spec(K)\to\Spec(A)$ mà điểm duy nhất được gửi tới điểm
đóng $\mathfrak{m}$ tương ứng song ánh với các đơn cấu trường
$A/\mathfrak{m}\to K$.
\footnote{\emph{[Dịch.] Câu tiếng Pháp in ra bỏ điều kiện ``mà điểm duy
nhất được gửi tới điểm đóng $\mathfrak{m}$''; nếu không có điều kiện
này, tương ứng chỉ đúng sau khi chọn điểm ảnh và dùng trường thặng dư
của điểm ấy.}}
\end{env}

Cho $(Y,\sh{O}_Y)$ là một tiền lược đồ; với mọi $y\in Y$ và mọi iđêan
$\mathfrak{a}_y$ của $\sh{O}_y$, đồng cấu chính tắc
$\sh{O}_y\to\sh{O}_y/\mathfrak{a}_y$ xác định một cấu xạ
$\Spec(\sh{O}_y/\mathfrak{a}_y)\to\Spec(\sh{O}_y)$; hợp nó với cấu xạ
chính tắc $\Spec(\sh{O}_y)\to Y$, ta thu được một cấu xạ
$\Spec(\sh{O}_y/\mathfrak{a}_y)\to Y$, cũng gọi là \emph{chính tắc}.
Với $\mathfrak{a}_y=\mathfrak{m}_y$, nghĩa là
$\sh{O}_y/\mathfrak{a}_y=\kres(y)$, Mệnh đề~\sref{I.2.4.4} do đó cho:

\begin{corollary}[2.4.6]
\label{I.2.4.6}
\oldpage[I]{103}
Cho $(X,\sh{O}_X)$ là một lược đồ địa phương có vành là trường $K$,
$\xi$ là điểm duy nhất của $X$, và $(Y,\sh{O}_Y)$ là một tiền lược đồ.
Mọi cấu xạ $u:(X,\sh{O}_X)\to(Y,\sh{O}_Y)$ phân tích duy nhất thành
$X\to\Spec(\kres(\psi(\xi)))\to Y$, trong đó mũi tên thứ hai chỉ cấu xạ
chính tắc, còn mũi tên thứ nhất tương ứng với một đơn cấu
$\kres(\psi(\xi))\to K$. Điều này thiết lập một tương ứng song ánh chính
tắc giữa tập các cấu xạ $(X,\sh{O}_X)\to(Y,\sh{O}_Y)$ và tập các đơn
cấu $\kres(y)\to K$ $(y\in Y)$.
\end{corollary}

\begin{corollary}[2.4.7]
\label{I.2.4.7}
Với mọi $y\in Y$, mọi cấu xạ chính tắc
$\Spec(\sh{O}_y/\mathfrak{a}_y)\to Y$ là một đơn cấu của các không gian
vành.
\end{corollary}

\begin{proof}
Điều này đã được chứng minh khi $\mathfrak{a}_y=0$ \sref{I.2.4.2}, và
chỉ cần áp dụng Hệ quả~\sref{I.1.7.5}.
\end{proof}

\begin{remark}[2.4.8]
\label{I.2.4.8}
Cho $X$ là một lược đồ địa phương và $a$ là điểm đóng duy nhất của nó.
Vì mọi tập mở affine chứa $a$ nhất thiết là toàn bộ $X$, nên mọi $\sh{O}_X$-môđun \emph{khả nghịch} \sref[0]{0.5.4.1} nhất thiết \emph{đẳng cấu với $\sh{O}_X$} (hay, như người ta cũng nói, là \emph{tầm thường}).
Tính chất này nói chung không đúng đối với một lược đồ affine tùy ý $\Spec(A)$;
trong Chương~V ta sẽ thấy rằng, nếu $A$ là một vành chuẩn, thì tính chất ấy đúng khi $A$ là một \emph{miền phân tích duy nhất}.
\end{remark}

\subsection{Các tiền lược đồ trên một tiền lược đồ}
\label{subsection:I.2.5}

\begin{definition}[2.5.1]
\label{I.2.5.1}
Cho một tiền lược đồ $S$, ta nói các dữ kiện gồm một tiền lược đồ $X$ và một cấu xạ của các tiền lược đồ $\vphi:X\to S$ xác định một tiền lược đồ $X$ \emph{trên tiền lược đồ $S$}, hay một \emph{$S$-tiền lược đồ}; ta nói $S$ là \emph{tiền lược đồ cơ sở} của $S$-tiền lược đồ $X$.
Cấu xạ $\vphi$ được gọi là \emph{cấu xạ cấu trúc} của $S$-tiền lược đồ $X$.
Khi $S$ là một lược đồ affine có vành $A$, ta cũng nói rằng $X$, cùng với $\vphi$, là một tiền lược đồ trên vành $A$ (hay một $A$-tiền lược đồ).
\end{definition}

Từ \sref{I.2.2.4} suy ra rằng dữ kiện về một tiền lược đồ trên vành $A$ tương đương với dữ kiện về một tiền lược đồ $(X,\sh{O}_X)$ mà bó cấu trúc $\sh{O}_X$ là một bó các $A$-\emph{đại số}.
\emph{Do đó, mọi tiền lược đồ luôn có thể được xem một cách duy nhất như một $\bb{Z}$-tiền lược đồ.}

Nếu $\vphi:X\to S$ là cấu xạ cấu trúc của một $S$-tiền lược đồ $X$, ta nói một điểm $x\in X$ \emph{nằm trên một điểm $s\in S$} nếu $\vphi(x)=s$.
Ta nói $X$ là một \emph{$S$-tiền lược đồ trội} nếu $\vphi$ là một cấu xạ trội \sref{I.2.2.6}.

\begin{env}[2.5.2]
\label{I.2.5.2}
Cho $X$ và $Y$ là hai $S$-tiền lược đồ; ta nói một cấu xạ của các tiền lược đồ $u:X\to Y$ là một \emph{cấu xạ của các tiền lược đồ trên $S$} (hay một $S$-\emph{cấu xạ}) nếu sơ đồ
\phantomsection\label{I.2.5.2.diagram}
\[
  \xymatrix{
    X\ar[rr]^u\ar[dr] & & Y\ar[dl]\\
    & S &
  }
\]
(trong đó các mũi tên xiên là các cấu xạ cấu trúc) giao hoán; điều này kéo theo rằng, với mọi $s\in S$ và mọi $x\in X$ nằm trên $s$, $u(x)$ cũng phải nằm trên $s$.
\end{env}

Từ định nghĩa này suy ra ngay rằng hợp thành của hai $S$-cấu xạ là một $S$-cấu xạ; do đó, các $S$-tiền lược đồ tạo thành một \emph{phạm trù}.

Ta ký hiệu $\Hom_S(X,Y)$ cho tập hợp các $S$-cấu xạ từ một $S$-tiền lược đồ $X$ đến một $S$-tiền lược đồ $Y$; cấu xạ đồng nhất của một $S$-tiền lược đồ $X$ được ký hiệu bởi $1_X$.

Khi $S$ là một lược đồ affine có vành $A$, ta cũng nói $A$-\emph{cấu xạ} thay cho $S$-cấu xạ.

\begin{env}[2.5.3]
\label{I.2.5.3}
\oldpage[I]{104}
Nếu $X$ là một $S$-tiền lược đồ và $v:X'\to X$ là một cấu xạ của các tiền lược
đồ, thì cấu xạ hợp thành $X'\xrightarrow{v}X\to S$ xác định $X'$ như một
$S$-tiền lược đồ; đặc biệt, mọi tiền lược đồ cảm sinh trên một tập mở $U$ của
$X$ đều có thể được xem là một $S$-tiền lược đồ nhờ phép nhúng chính tắc.

Nếu $u:X\to Y$ là một $S$-cấu xạ của các $S$-tiền lược đồ, thì hạn chế của $u$
lên mọi tiền lược đồ cảm sinh trên một tập mở $U$ của $X$ do đó là một
$S$-cấu xạ $U\to Y$. Ngược lại, cho $(U_\alpha)$ là một phủ mở của $X$, và với
mỗi $\alpha$, cho $u_\alpha:U_\alpha\to Y$ là một $S$-cấu xạ; nếu, với mọi cặp
chỉ số $(\alpha,\beta)$, các hạn chế của $u_\alpha$ và $u_\beta$ lên
$U_\alpha\cap U_\beta$ trùng nhau, thì tồn tại duy nhất một $S$-cấu xạ
$X\to Y$ có hạn chế lên mỗi $U_\alpha$ là $u_\alpha$.

Nếu $u:X\to Y$ là một $S$-cấu xạ sao cho $u(X)\subset V$, trong đó $V$ là một
phần mở của $Y$, thì $u$, được xem như một cấu xạ từ $X$ vào $V$, vẫn là một
$S$-cấu xạ.
\end{env}

\begin{env}[2.5.4]
\label{I.2.5.4}
Cho $S'\to S$ là một cấu xạ của các tiền lược đồ; với mọi $S'$-tiền lược đồ
$X$\footnote{\emph{[Dịch.] Nguồn tiếng Pháp in «pour tout S'-préschéma», bỏ
sót đối tượng X; hợp thành và kết luận ngay sau đó buộc phải ràng buộc X là
S'-tiền lược đồ đang xét.}}, cấu xạ hợp thành $X\to S'\to S$ xác định $X$ như
một $S$-tiền lược đồ.
Ngược lại, giả sử $S'$ là tiền lược đồ cảm sinh trên một tập mở của $S$; cho
$X$ là một $S$-tiền lược đồ và giả sử cấu xạ cấu trúc $f:X\to S$ thỏa mãn
$f(X)\subset S'$; khi đó ta có thể xem $X$ như một $S'$-tiền lược đồ. Trong
trường hợp sau cùng này, nếu $Y$ là một $S$-tiền lược đồ mà cấu xạ cấu trúc của
nó cũng gửi không gian nền vào $S'$, thì mọi $S$-cấu xạ từ $X$ đến $Y$ cũng là
một $S'$-cấu xạ.
\end{env}

\begin{env}[2.5.5]
\label{I.2.5.5}
Nếu $X$ là một $S$-tiền lược đồ,\footnote{\emph{[Dịch.] Nguồn tiếng Pháp in
«Si X est un S-morphisme»; cấu xạ cấu trúc X→S và định nghĩa một nhát cắt đòi
hỏi X phải là một S-tiền lược đồ.}} và $\vphi:X\to S$ là cấu xạ cấu trúc của
nó, thì ta gọi một $S$-\emph{nhát cắt của $X$} là một $S$-cấu xạ từ $S$ đến
$X$, tức là một cấu xạ của các tiền lược đồ $\psi:S\to X$ sao cho
$\vphi\circ\psi$ là cấu xạ đồng nhất của $S$. Ta ký hiệu $\Gamma(X/S)$ cho tập
hợp các $S$-nhát cắt của $X$.
\end{env}

\section{Tích của các tiền lược đồ}
\label{section:I.3}

\subsection{Tổng của các tiền lược đồ}
\label{subsection:I.3.1}
\label{I.3.1}

\phantomsection\label{I.3.1.1}
Cho $(X_\alpha)$ là một họ tùy ý các tiền lược đồ; cho $X$ là không gian tôpô
\emph{tổng} của các không gian nền $X_\alpha$; khi đó $X$ là hợp của các không
gian con mở $X'_\alpha$ đôi một rời nhau, và với mỗi $\alpha$ có một phép đồng
phôi $\vphi_\alpha$ từ $X_\alpha$ lên $X'_\alpha$. Nếu trang bị cho mỗi
$X'_\alpha$ bó $(\vphi_\alpha)_*(\sh{O}_{X_\alpha})$, thì hiển nhiên $X$ trở
thành một tiền lược đồ, gọi là \emph{tổng} của họ tiền lược đồ $(X_\alpha)$ và
được ký hiệu bởi $\bigsqcup_\alpha X_\alpha$. Nếu $Y$ là một tiền lược đồ,
ánh xạ $f\mapsto(f\circ\vphi_\alpha)$ là một \emph{song ánh} từ tập hợp
$\Hom(X,Y)$ lên tập hợp tích $\prod_\alpha\Hom(X_\alpha,Y)$. Đặc biệt, nếu các
$X_\alpha$ là các $S$-tiền lược đồ với các cấu xạ cấu trúc $\psi_\alpha$, thì
$X$ là một $S$-tiền lược đồ đối với cấu xạ duy nhất $\psi:X\to S$ sao cho
$\psi\circ\vphi_\alpha=\psi_\alpha$ với mọi $\alpha$. Tổng của hai tiền lược
đồ $X$, $Y$ được ký hiệu bởi $X\sqcup Y$. Hiển nhiên rằng nếu $X=\Spec(A)$ và
$Y=\Spec(B)$, thì $X\sqcup Y$ được đồng nhất chính tắc với $\Spec(A\times B)$.

\subsection{Tích của các tiền lược đồ}
\label{subsection:I.3.2}

\begin{definition}[3.2.1]
\label{I.3.2.1}
Cho hai $S$-tiền lược đồ $X$, $Y$, ta nói một bộ ba $(Z,p_1,p_2)$ gồm một
$S$-tiền lược đồ $Z$ và hai $S$-cấu xạ $p_1:Z\to X$, $p_2:Z\to Y$ là một
tích của các
\oldpage[I]{105}
$S$-tiền lược đồ $X$ và $Y$ nếu, với mọi $S$-tiền lược đồ $T$, ánh xạ
$f\mapsto(p_1\circ f,p_2\circ f)$ là một song ánh từ tập hợp các $S$-cấu xạ
từ $T$ đến $Z$ lên tập hợp các cặp gồm một $S$-cấu xạ $T\to X$ và một
$S$-cấu xạ $T\to Y$ (nói cách khác, một song ánh
\[
  \Hom_S(T,Z)\isoto\Hom_S(T,X)\times\Hom_S(T,Y)).
\]
\end{definition}

Đây chính là khái niệm tổng quát về \emph{tích} của hai đối tượng trong một
phạm trù, áp dụng cho phạm trù các $S$-tiền lược đồ (T, I, 1.1); đặc biệt,
một tích của hai $S$-tiền lược đồ là \emph{duy nhất}, sai khác một
$S$-đẳng cấu duy nhất. Do tính duy nhất này, ta thường ký hiệu một tích của
hai $S$-tiền lược đồ $X$, $Y$ bởi $X\times_S Y$ (hoặc đơn giản bởi
$X\times Y$ nếu không thể xảy ra nhầm lẫn), và bỏ các cấu xạ $p_1$, $p_2$
khỏi ký hiệu (chúng được gọi là các \emph{phép chiếu chính tắc} của
$X\times_S Y$ lên $X$ và $Y$, tương ứng). Nếu $g:T\to X$, $h:T\to Y$ là hai
$S$-cấu xạ, ta ký hiệu bởi $(g,h)_S$, hoặc đơn giản bởi $(g,h)$, $S$-cấu xạ
$f:T\to X\times_S Y$ sao cho $p_1\circ f=g$, $p_2\circ f=h$. Nếu $X'$, $Y'$
là hai $S$-tiền lược đồ, $p'_1$, $p'_2$ là các phép chiếu chính tắc của
$X'\times_S Y'$ (được giả sử tồn tại), còn $u:X'\to X$, $v:Y'\to Y$ là hai
$S$-cấu xạ, ta viết $u\times_S v$ (hoặc đơn giản $u\times v$) cho
$S$-cấu xạ $(u\circ p'_1,v\circ p'_2)_S$ từ $X'\times_S Y'$ đến
$X\times_S Y$.

Khi $S$ là một lược đồ affine có vành $A$, ta thường thay $S$ bằng $A$ trong
các ký hiệu trên.

\begin{proposition}[3.2.2]
\label{I.3.2.2}
Cho $X$, $Y$, $S$ là ba lược đồ affine, và cho $B$, $C$, $A$ là các vành
tương ứng của chúng. Đặt $Z=\Spec(B\otimes_A C)$, và cho $p_1$, $p_2$ là các
$S$-cấu xạ tương ứng, theo \sref{I.2.2.4}, với các $A$-đồng cấu chính tắc
$u:b\mapsto b\otimes1$ và $v:c\mapsto1\otimes c$ từ $B$ và $C$ vào
$B\otimes_A C$; khi đó $(Z,p_1,p_2)$ là một tích của $X$ và $Y$.
\end{proposition}

\begin{proof}
Theo \sref{I.2.2.4}, chỉ cần kiểm tra rằng việc gán cho mỗi $A$-đồng cấu
$f:B\otimes_A C\to L$ (trong đó $L$ là một $A$-đại số) cặp
$(f\circ u,f\circ v)$ xác định một song ánh
\[
  \Hom_A(B\otimes_A C,L)\isoto
  \Hom_A(B,L)\times\Hom_A(C,L),
\]\footnote{Ký hiệu $\Hom_A$ ở đây chỉ tập hợp các đồng cấu của
\emph{$A$-đại số}.} điều này suy ra ngay từ các định nghĩa và hệ thức
$b\otimes c=(b\otimes1)(1\otimes c)$.
\end{proof}

\begin{corollary}[3.2.3]
\label{I.3.2.3}
Cho $T$ là một lược đồ affine có vành $D$, và cho
$\alpha=({}^a\rho,\widetilde{\rho})$ (tương ứng
$\beta=({}^a\sigma,\widetilde{\sigma})$) là một $S$-cấu xạ $T\to X$ (tương
ứng $T\to Y$), trong đó $\rho$ (tương ứng $\sigma$) là một $A$-đồng cấu từ
$B$ (tương ứng $C$) vào $D$; khi đó
$(\alpha,\beta)_S=({}^a\tau,\widetilde{\tau})$, trong đó $\tau$ là đồng cấu
$B\otimes_A C\to D$ sao cho $\tau(b\otimes c)=\rho(b)\sigma(c)$.
\end{corollary}

\begin{proposition}[3.2.4]
\label{I.3.2.4}
Cho $f:S'\to S$ là một đơn cấu của các tiền lược đồ (T, I, 1.1), và cho
$X$, $Y$ là hai $S'$-tiền lược đồ, đồng thời được xem như các $S$-tiền lược
đồ thông qua $f$. Mọi tích của các $S$-tiền lược đồ $X$, $Y$ khi đó là một
tích của các $S'$-tiền lược đồ $X$, $Y$, và ngược lại.
\end{proposition}

\begin{proof}
Cho $\vphi:X\to S'$, $\psi:Y\to S'$ là các cấu xạ cấu trúc. Nếu $T$ là một
$S$-tiền lược đồ, còn $u:T\to X$, $v:T\to Y$ là hai $S$-cấu xạ, thì theo
định nghĩa ta có $f\circ\vphi\circ u=f\circ\psi\circ v=\theta$, là cấu xạ
cấu trúc của $T$; giả thiết về $f$ kéo theo
$\vphi\circ u=\psi\circ v=\theta'$, và ta thấy có thể xem $T$ như một
$S'$-tiền lược đồ có cấu xạ cấu trúc $\theta'$, còn $u$, $v$ như các
$S'$-cấu xạ. Kết luận của mệnh đề suy ra ngay, có tính đến
\sref{I.3.2.1}.
\end{proof}

\begin{corollary}[3.2.5]
\label{I.3.2.5}
Cho $X$, $Y$ là hai $S$-tiền lược đồ, $\vphi:X\to S$, $\psi:Y\to S$ là các
cấu xạ cấu trúc của chúng, và cho $S'$ là một phần mở của $S$ sao cho
$\vphi(X)\subset S'$, $\psi(Y)\subset S'$. Mọi tích của các $S$-tiền lược đồ
$X$, $Y$ cũng là một tích của các $S'$-tiền lược đồ $X$, $Y$, và ngược lại.
\end{corollary}

\oldpage[I]{106}
Chỉ cần áp dụng \sref{I.3.2.4} cho phép nhúng chính tắc $S'\to S$.

\begin{theorem}[3.2.6]
\label{I.3.2.6}
Cho hai $S$-tiền lược đồ $X$, $Y$, tồn tại một tích $X\times_S Y$.
\end{theorem}

Ta sẽ tiến hành qua nhiều bước.

\begin{lemma}[3.2.6.1]
\label{I.3.2.6.1}
Cho $(Z,p,q)$ là một tích của $X$ và $Y$, còn $U$, $V$ lần lượt là các phần
mở của $X$, $Y$. Nếu đặt $W=p^{-1}(U)\cap q^{-1}(V)$, thì bộ ba gồm $W$ và
các hạn chế của $p$, $q$ lên $W$ \textup{(lần lượt được xem như các cấu xạ
$W\to U$, $W\to V$)} là một tích của $U$ và $V$.
\end{lemma}

Thật vậy, nếu $T$ là một $S$-tiền lược đồ, ta có thể đồng nhất các
$S$-cấu xạ $T\to W$ với các $S$-cấu xạ $T\to Z$ gửi $T$ vào $W$. Nếu
$g:T\to U$, $h:T\to V$ là hai $S$-cấu xạ bất kỳ, thì ta có thể lần lượt xem
chúng như các $S$-cấu xạ từ $T$ vào $X$ và $Y$; do giả thiết, tồn tại duy
nhất một $S$-cấu xạ $f:T\to Z$ sao cho $g=p\circ f$, $h=q\circ f$. Vì
$p(f(T))\subset U$, $q(f(T))\subset V$, ta có
\[
  f(T)\subset p^{-1}(U)\cap q^{-1}(V)=W,
\]
từ đó suy ra khẳng định.

\begin{lemma}[3.2.6.2]
\label{I.3.2.6.2}
Cho $Z$ là một $S$-tiền lược đồ, $p:Z\to X$, $q:Z\to Y$ là hai
$S$-cấu xạ, $(U_\alpha)$ là một phủ mở của $X$, $(V_\lambda)$ là một phủ mở
của $Y$. Giả sử rằng, với mọi cặp $(\alpha,\lambda)$, $S$-tiền lược đồ
$W_{\alpha\lambda}=p^{-1}(U_\alpha)\cap q^{-1}(V_\lambda)$ cùng các hạn chế
của $p$ và $q$ lên $W_{\alpha\lambda}$ tạo thành một tích của $U_\alpha$ và
$V_\lambda$. Khi đó $(Z,p,q)$ là một tích của $X$ và $Y$.
\end{lemma}

Trước hết, ta chứng minh rằng nếu $f_1$, $f_2$ là hai $S$-cấu xạ $T\to Z$,
thì các hệ thức $p\circ f_1=p\circ f_2$ và $q\circ f_1=q\circ f_2$ kéo theo
$f_1=f_2$. Thật vậy, $Z$ là hợp của các $W_{\alpha\lambda}$; do đó, các
$f_1^{-1}(W_{\alpha\lambda})$ tạo thành một phủ mở của $T$, và điều tương tự
đúng với các $f_2^{-1}(W_{\alpha\lambda})$. Hơn nữa, ta có
\[
\begin{aligned}
  f_1^{-1}(W_{\alpha\lambda})
  &=f_1^{-1}(p^{-1}(U_\alpha))\cap
    f_1^{-1}(q^{-1}(V_\lambda)) \\
  &=f_2^{-1}(p^{-1}(U_\alpha))\cap
    f_2^{-1}(q^{-1}(V_\lambda))
   =f_2^{-1}(W_{\alpha\lambda})
\end{aligned}
\]
theo giả thiết, và chỉ còn phải thấy rằng các hạn chế của $f_1$ và $f_2$ lên
$f_1^{-1}(W_{\alpha\lambda})=f_2^{-1}(W_{\alpha\lambda})$ là đồng nhất với
nhau cho mọi cặp chỉ số. Nhưng vì có thể xem các hạn chế ấy như các
$S$-cấu xạ từ $f_1^{-1}(W_{\alpha\lambda})$ vào $W_{\alpha\lambda}$, khẳng
định của ta suy ra từ các giả thiết và Định nghĩa~\sref{I.3.2.1}.

Bây giờ, giả sử cho hai $S$-cấu xạ $g:T\to X$, $h:T\to Y$. Đặt
$T_{\alpha\lambda}=g^{-1}(U_\alpha)\cap h^{-1}(V_\lambda)$; các
$T_{\alpha\lambda}$ khi đó tạo thành một phủ mở của $T$. Theo giả thiết, tồn
tại một $S$-cấu xạ $f_{\alpha\lambda}$ sao cho $p\circ f_{\alpha\lambda}$ và
$q\circ f_{\alpha\lambda}$ lần lượt là các hạn chế của $g$ và $h$ lên
$T_{\alpha\lambda}$. Hơn nữa, ta chứng minh rằng các hạn chế của
$f_{\alpha\lambda}$ và $f_{\beta\mu}$ lên
$T_{\alpha\lambda}\cap T_{\beta\mu}$ trùng nhau; điều này sẽ hoàn tất chứng
minh \sref{I.3.2.6.2}. Theo định nghĩa, ảnh của
$T_{\alpha\lambda}\cap T_{\beta\mu}$ qua $f_{\alpha\lambda}$ và
$f_{\beta\mu}$ nằm trong $W_{\alpha\lambda}\cap W_{\beta\mu}$. Vì
\[
  W_{\alpha\lambda}\cap W_{\beta\mu}
  =p^{-1}(U_\alpha\cap U_\beta)\cap
   q^{-1}(V_\lambda\cap V_\mu),
\]
từ \sref{I.3.2.6.1} suy ra rằng
$W_{\alpha\lambda}\cap W_{\beta\mu}$ cùng các hạn chế của $p$ và $q$ lên
tiền lược đồ này tạo thành một tích của $U_\alpha\cap U_\beta$ và
$V_\lambda\cap V_\mu$. Vì $p\circ f_{\alpha\lambda}$ và
$p\circ f_{\beta\mu}$ trùng nhau trên
$T_{\alpha\lambda}\cap T_{\beta\mu}$, và điều tương tự đúng với
$q\circ f_{\alpha\lambda}$ và $q\circ f_{\beta\mu}$, nên
$f_{\alpha\lambda}$ và $f_{\beta\mu}$ trùng nhau trên
$T_{\alpha\lambda}\cap T_{\beta\mu}$.\qed

\oldpage[I]{107}
\begin{lemma}[3.2.6.3]
\label{I.3.2.6.3}
Cho $(U_\alpha)$ là một phủ mở của $X$, $(V_\lambda)$ là một phủ mở của
$Y$, và giả sử rằng với mọi cặp $(\alpha,\lambda)$, tồn tại một tích của
$U_\alpha$ và $V_\lambda$; khi đó tồn tại một tích của $X$ và $Y$.
\end{lemma}

Áp dụng Bổ đề~\sref{I.3.2.6.1} cho các phần mở
$U_\alpha\cap U_\beta$ và $V_\lambda\cap V_\mu$, ta thấy tồn tại một tích
của các $S$-tiền lược đồ lần lượt do $X$ và $Y$ cảm sinh trên các phần mở
này; hơn nữa, tính duy nhất của tích cho thấy rằng, nếu đặt
$i=(\alpha,\lambda)$, $j=(\beta,\mu)$, thì có một đẳng cấu chính tắc
$h_{ij}$ (tương ứng $h_{ji}$) từ tích này lên một $S$-tiền lược đồ $W_{ij}$
(tương ứng $W_{ji}$) do $U_\alpha\times_S V_\lambda$
(tương ứng $U_\beta\times_S V_\mu$) cảm sinh trên một phần mở; do đó
$f_{ij}=h_{ij}\circ h_{ji}^{-1}$ là một đẳng cấu từ $W_{ji}$ lên $W_{ij}$.
Hơn nữa, với cặp thứ ba $k=(\gamma,\nu)$, ta có
$f_{ik}=f_{ij}\circ f_{jk}$ trên $W_{ki}\cap W_{kj}$, như suy ra khi áp
dụng \sref{I.3.2.6.1} cho các phần mở
$U_\alpha\cap U_\beta\cap U_\gamma$ và
$V_\lambda\cap V_\mu\cap V_\nu$, lần lượt trong $U_\beta$ và $V_\mu$.
Do đó tồn tại một tiền lược đồ $Z$, một phủ mở $(Z_i)$ của không gian nền
của $Z$, và với mỗi $i$ một đẳng cấu $g_i$ từ tiền lược đồ cảm sinh $Z_i$
lên tiền lược đồ $U_\alpha\times_S V_\lambda$, sao cho với mọi cặp $(i,j)$
ta có $f_{ij}=g_i\circ g_j^{-1}$ \sref{I.2.3.1}; hơn nữa,
$g_i(Z_i\cap Z_j)=W_{ij}$. Nếu $p_i$, $q_i$, $\theta_i$ là các phép chiếu và
cấu xạ cấu trúc của $S$-tiền lược đồ $U_\alpha\times_S V_\lambda$, ta thấy
ngay rằng $p_i\circ g_i=p_j\circ g_j$ trên $Z_i\cap Z_j$, và điều tương tự
đúng với hai cấu xạ còn lại. Vì vậy, ta có thể xác định các cấu xạ của các
tiền lược đồ $p:Z\to X$ (tương ứng $q:Z\to Y$, $\theta:Z\to S$) bởi điều
kiện rằng $p$ (tương ứng $q$, $\theta$) trùng với $p_i\circ g_i$
(tương ứng $q_i\circ g_i$, $\theta_i\circ g_i$) trên mỗi $Z_i$; $Z$, khi
được trang bị $\theta$, là một $S$-tiền lược đồ. Bây giờ ta chứng minh rằng
$Z'_i=p^{-1}(U_\alpha)\cap q^{-1}(V_\lambda)$ bằng $Z_i$. Với mọi chỉ số
$j=(\beta,\mu)$, ta có
$Z_j\cap Z'_i=g_j^{-1}(p_j^{-1}(U_\alpha)\cap q_j^{-1}(V_\lambda))$. Mà
\[
  p_j^{-1}(U_\alpha)\cap q_j^{-1}(V_\lambda)
  =p_j^{-1}(U_\alpha\cap U_\beta)\cap
   q_j^{-1}(V_\lambda\cap V_\mu);
\]
theo \sref{I.3.2.6.1}, các hạn chế của $p_j$ và $q_j$ lên
$p_j^{-1}(U_\alpha)\cap q_j^{-1}(V_\lambda)$ xác định trên
$S$-tiền lược đồ này một cấu trúc tích của $U_\alpha\cap U_\beta$ và
$V_\lambda\cap V_\mu$; nhưng tính duy nhất của tích khi đó kéo theo
$p_j^{-1}(U_\alpha)\cap q_j^{-1}(V_\lambda)=W_{ji}$. Vì vậy,
$Z_j\cap Z'_i=Z_j\cap Z_i$ với mọi $j$, từ đó $Z'_i=Z_i$. Theo
\sref{I.3.2.6.2}, suy ra $(Z,p,q)$ là một tích của $X$ và $Y$.

\begin{lemma}[3.2.6.4]
\label{I.3.2.6.4}
Cho $\vphi:X\to S$, $\psi:Y\to S$ là các cấu xạ cấu trúc của $X$ và
$Y$, $(S_i)$ là một phủ mở của $S$, và đặt
$X_i=\vphi^{-1}(S_i)$, $Y_i=\psi^{-1}(S_i)$. Nếu mỗi tích
$X_i\times_S Y_i$ tồn tại, thì $X\times_S Y$ tồn tại.
\end{lemma}

Theo \sref{I.3.2.6.3}, bài toán quy về chứng minh rằng các tích
$X_i\times_S Y_j$ tồn tại với mọi $i$ và $j$. Đặt
$X_{ij}=X_i\cap X_j=\vphi^{-1}(S_i\cap S_j)$,
$Y_{ij}=Y_i\cap Y_j=\psi^{-1}(S_i\cap S_j)$; theo
\sref{I.3.2.6.1}, tích $Z_{ij}=X_{ij}\times_S Y_{ij}$ tồn tại. Bây giờ
lưu ý rằng, nếu $T$ là một $S$-tiền lược đồ và nếu $g:T\to X_i$,
$h:T\to Y_j$ là các $S$-cấu xạ, thì nhất thiết
$\vphi(g(T))=\psi(h(T))\subset S_i\cap S_j$, theo định nghĩa của
$S$-cấu xạ; do đó $g(T)\subset X_{ij}$ và $h(T)\subset Y_{ij}$; khi đó
hiển nhiên $Z_{ij}$ là một tích của $X_i$ và $Y_j$.

\begin{env}[3.2.6.5]
\label{I.3.2.6.5}
Bây giờ ta có thể hoàn tất chứng minh Định lý~\sref{I.3.2.6}. Nếu $S$ là
một \emph{lược đồ affine}, thì có các phủ $(U_\alpha)$, $(V_\lambda)$ lần
lượt của $X$ và $Y$, gồm các phần mở affine; vì
$U_\alpha\times_S V_\lambda$ tồn tại theo \sref{I.3.2.2}, nên
$X\times_S Y$ cũng tồn tại theo \sref{I.3.2.6.3}. Nếu $S$ là một tiền lược
đồ bất kỳ, thì có một phủ $(S_i)$ của $S$ gồm các phần mở affine. Nếu
$\vphi:X\to S$, $\psi:Y\to S$ là các cấu xạ cấu trúc, và nếu đặt
$X_i=\vphi^{-1}(S_i)$, $Y_i=\psi^{-1}(S_i)$, thì các tích
$X_i\times_{S_i}Y_i$ tồn tại theo điều
\oldpage[I]{108}
vừa chứng minh; nhưng khi đó các tích $X_i\times_S Y_i$ cũng tồn tại
\sref{I.3.2.5}, do đó $X\times_S Y$ cũng tồn tại theo
\sref{I.3.2.6.4}.\qed
\end{env}

\begin{corollary}[3.2.7]
\label{I.3.2.7}
Cho $Z=X\times_S Y$ là tích của hai $S$-tiền lược đồ, $p$, $q$ là các phép
chiếu từ $Z$ lên $X$ và $Y$, và $\vphi$ (tương ứng $\psi$) là cấu xạ cấu
trúc của $X$ (tương ứng $Y$). Cho $S'$ là một phần mở của $S$, và $U$
(tương ứng $V$) là một phần mở của $X$ (tương ứng $Y$) nằm trong
$\vphi^{-1}(S')$ (tương ứng $\psi^{-1}(S')$). Khi đó tích
$U\times_{S'}V$ được đồng nhất chính tắc với tiền lược đồ do $Z$ cảm sinh
trên $p^{-1}(U)\cap q^{-1}(V)$ (được xem như một $S'$-tiền lược đồ). Hơn
nữa, nếu $f:T\to X$, $g:T\to Y$ là các $S$-cấu xạ sao cho
$f(T)\subset U$, $g(T)\subset V$, thì $S'$-cấu xạ $(f,g)_{S'}$ được đồng
nhất với $(f,g)_S$ khi được xem như một cấu xạ nhận giá trị trong
$p^{-1}(U)\cap q^{-1}(V)$.
\end{corollary}

\begin{proof}
Điều này suy ra từ Hệ quả~\sref{I.3.2.5} và Bổ đề~\sref{I.3.2.6.1}.
\end{proof}

\begin{env}[3.2.8]
\label{I.3.2.8}
Cho $(X_\alpha)$, $(Y_\lambda)$ là hai họ $S$-tiền lược đồ, và cho $X$
(tương ứng $Y$) là tổng của họ $(X_\alpha)$ (tương ứng $(Y_\lambda)$)
\hyperref[subsection:I.3.1]{(3.1)}. Khi đó $X\times_S Y$ được đồng nhất
với \emph{tổng} của họ $(X_\alpha\times_S Y_\lambda)$; điều này suy ra
ngay từ \sref{I.3.2.6.3}.
\end{env}

\begin{env}[3.2.9]
\label{I.3.2.9}
\footnote{\emph{[Dịch.] Mục \sref{I.3.2.9} xuất phát từ phần đính chính
của EGA~II, ở trang 221; vì thế cách đánh số trang thay đổi.}}
\oldpage[II]{221}
Từ \sref{I.1.8.1} suy ra rằng ta có thể phát biểu \sref{I.3.2.2} chính
xác hơn như sau: $Z=\Spec(B\otimes_A C)$ không những là một tích của
$X=\Spec(B)$ và $Y=\Spec(C)$ trong phạm trù các
\emph{$S$-tiền lược đồ}, mà còn là một tích trong phạm trù các
\emph{không gian vành địa phương trên $S$} (với định nghĩa các
$S$-cấu xạ được phỏng theo định nghĩa ở \sref{I.2.5.2}). Chứng minh của
\sref{I.3.2.6} còn chứng tỏ rằng, đối với hai $S$-tiền lược đồ bất kỳ
$X$ và $Y$, tiền lược đồ $X\times_S Y$ không những là tích của $X$ và
$Y$ trong phạm trù các $S$-tiền lược đồ, mà còn là tích của chúng trong
phạm trù các không gian vành địa phương trên tiền lược đồ $S$.
\end{env}

\subsection{Các tính chất hình thức của tích; thay đổi tiền lược đồ cơ sở}
\label{subsection:I.3.3}

\begin{env}[3.3.1]
\label{I.3.3.1}
Độc giả sẽ nhận thấy rằng tất cả các tính chất được phát biểu trong mục
này, ngoại trừ \sref{I.3.3.13}, \sref{I.3.3.15} và khẳng định về các
tổng trong câu cuối của \sref{I.3.3.10}, đều đúng mà không cần thay đổi
trong \emph{mọi phạm trù}, miễn là các tích xuất hiện trong các phát biểu
\emph{tồn tại} (vì rõ ràng các khái niệm đối tượng trên $S$ và
$S$-cấu xạ có thể được định nghĩa hoàn toàn như trong
\hyperref[subsection:I.2.5]{(2.5)} đối với mọi đối tượng $S$ của phạm
trù).
\end{env}

\begin{env}[3.3.2]
\label{I.3.3.2}
Trước hết, $X\times_S Y$ là một \emph{song hàm tử hiệp biến} theo $X$ và
$Y$ trong phạm trù các $S$-tiền lược đồ: thật vậy, chỉ cần nhận thấy rằng
biểu đồ
\[
  \xymatrix{
    X\times Y\ar[r]^{f\times 1}\ar[d] &
    X'\times Y\ar[r]^{f'\times 1}\ar[d] &
    X''\times Y\ar[d]\\
    X\ar[r]_f & X'\ar[r]_{f'} & X''
  }
\]
là giao hoán.
\end{env}

\begin{proposition}[3.3.3]
\label{I.3.3.3}
Với mọi $S$-tiền lược đồ $X$, phép chiếu thứ nhất (tương ứng, thứ hai) từ
$X\times_S S$ (tương ứng, $S\times_S X$) là một đẳng cấu có tính hàm tử từ
$X\times_S S$ (tương ứng, $S\times_S X$) lên $X$, mà đẳng cấu ngược của nó là
$(1_X,\vphi)_S$ (tương ứng, $(\vphi,1_X)_S$), trong đó $\vphi$ ký hiệu cấu xạ
cấu trúc $X\to S$; do đó, sai khác một đẳng cấu chính tắc, có thể viết
\[
  X\times_S S=S\times_S X=X.
\]
\end{proposition}

\begin{proof}
Chỉ cần chứng minh rằng bộ ba $(X,1_X,\vphi)$ là một tích của $X$ và $S$.
Thật vậy, nếu $T$ là một $S$-tiền lược đồ, thì $S$-cấu xạ duy nhất từ $T$
đến $S$ nhất thiết là cấu xạ cấu trúc $\psi:T\to S$. Nếu $f$ là một
$S$-cấu xạ từ $T$ đến $X$, thì nhất thiết $\psi=\vphi\circ f$, từ đó suy ra
khẳng định.
\end{proof}

\begin{corollary}[3.3.4]
\label{I.3.3.4}
Cho $X$ và $Y$ là hai $S$-tiền lược đồ, và cho $\vphi:X\to S$ và
$\psi:Y\to S$ là các cấu xạ cấu trúc của chúng. Nếu ta đồng nhất chính tắc
$X$ với $X\times_S S$ và $Y$ với $S\times_S Y$, thì các phép chiếu
$X\times_S Y\to X$ và $X\times_S Y\to Y$ lần lượt được đồng nhất với
$1_X\times\psi$ và $\vphi\times 1_Y$.
\end{corollary}

Việc kiểm tra là hiển nhiên và được dành cho độc giả.

\begin{env}[3.3.5]
\label{I.3.3.5}
Có thể định nghĩa, theo cùng cách như trong
\hyperref[subsection:I.3.2]{(3.2)}, tích của một
\oldpage[I]{109}
số hữu hạn tùy ý $n$ tiền lược đồ trên $S$; sự tồn tại của các tích này suy ra
từ \sref{I.3.2.6} bằng quy nạp theo $n$, khi nhận thấy rằng
$(X_1\times_S X_2\times_S\cdots\times_S X_{n-1})\times_S X_n$ thỏa mãn định
nghĩa của tích. Tính duy nhất của tích kéo theo, như trong mọi phạm trù, các
tính chất \emph{giao hoán} và \emph{kết hợp} của nó. Nếu, chẳng hạn, $p_1$,
$p_2$ và $p_3$ ký hiệu các phép chiếu của
$X_1\times_S X_2\times_S X_3$, và nếu đồng nhất tiền lược đồ này với
$(X_1\times_S X_2)\times_S X_3$, thì phép chiếu lên
$X_1\times_S X_2$ được đồng nhất với $(p_1,p_2)_S$.
\end{env}

\begin{env}[3.3.6]
\label{I.3.3.6}
Cho $S$ và $S'$ là hai tiền lược đồ, và cho $\vphi:S'\to S$ là một cấu xạ,
biến $S'$ thành một $S$-tiền lược đồ. Với mỗi $S$-tiền lược đồ $X$, xét tích
$X\times_S S'$, và gọi $p$ và $\pi'$ lần lượt là các phép chiếu của nó lên $X$
và $S'$. Được trang bị $\pi'$, tích này là một $S'$-tiền lược đồ; khi xét nó
như vậy, ta ký hiệu nó bởi $X_{(S')}$ hoặc $X_{(\vphi)}$, và gọi nó là tiền
lược đồ thu được bằng \emph{mở rộng tiền lược đồ cơ sở} từ $S$ đến $S'$ nhờ
cấu xạ $\vphi$, hoặc \emph{ảnh ngược} của $X$ qua $\vphi$. Lưu ý rằng, nếu
$\pi$ là cấu xạ cấu trúc của $X$, còn $\theta$ là cấu xạ cấu trúc của
$X\times_S S'$ khi được xét như một $S$-tiền lược đồ, thì biểu đồ
\[
  \xymatrix{
    X\ar[d]_\pi &
    X_{(S')}\ar[l]_p\ar[ld]_\theta\ar[d]^{\pi'}\\
    S &
    S'\ar[l]^\vphi
  }
\]
giao hoán.
\end{env}

\begin{env}[3.3.7]
\label{I.3.3.7}
Với các ký hiệu của \sref{I.3.3.6}, với mọi $S$-cấu xạ $f:X\to Y$, ta vẫn
dùng ký hiệu $f_{(S')}$ cho $S'$-cấu xạ
$f\times_S 1:X_{(S')}\to Y_{(S')}$, và gọi
$f_{(S')}$ là \emph{ảnh ngược} của cấu xạ $f$ qua $\vphi$. Do đó,
$X_{(S')}$ là một \emph{hàm tử hiệp biến} theo $X$, từ phạm trù các
$S$-tiền lược đồ đến phạm trù các $S'$-tiền lược đồ.
\end{env}

\begin{env}[3.3.8]
\label{I.3.3.8}
Tiền lược đồ $X_{(S')}$ còn có thể được xem như nghiệm của một
\emph{bài toán ánh xạ phổ dụng}: mọi $S'$-tiền lược đồ $T$ cũng là một
$S$-tiền lược đồ nhờ $\vphi$; mọi $S$-cấu xạ $g:T\to X$ khi đó được viết một
cách duy nhất dưới dạng $g=p\circ f$, trong đó $f$ là một $S'$-cấu xạ
$T\to X_{(S')}$, như suy ra từ định nghĩa của tích áp dụng cho các $S$-cấu xạ
$g$\footnote{\emph{[Dịch.] Bản tiếng Pháp in $f$ ở đây; cặp tham gia trong
định nghĩa là $(g,\psi)$.}} và $\psi:T\to S'$ (cấu xạ cấu trúc của $T$).
\end{env}

\begin{proposition}[3.3.9]
\label{I.3.3.9}
\emph{(``Tính bắc cầu của phép mở rộng các tiền lược đồ cơ sở'')}. Cho $S''$
là một tiền lược đồ và $\vphi':S''\to S'$ là một cấu xạ. Với mọi
$S$-tiền lược đồ $X$, tồn tại một đẳng cấu chính tắc có tính hàm tử từ
$S''$-tiền lược đồ $(X_{(\vphi)})_{(\vphi')}$ lên $S''$-tiền lược đồ
$X_{(\vphi\circ\vphi')}$.
\end{proposition}

\begin{proof}
Thật vậy, cho $T$ là một $S''$-tiền lược đồ, $\psi$ là cấu xạ cấu trúc của nó,
và $g$ là một $S$-cấu xạ từ $T$ đến $X$ (xét $T$ như một $S$-tiền lược đồ có
cấu xạ cấu trúc $\vphi\circ\vphi'\circ\psi$). Vì $T$ cũng là một
$S'$-tiền lược đồ có cấu xạ cấu trúc $\vphi'\circ\psi$, ta có thể viết
$g=p\circ g'$, trong đó $g'$ là một $S'$-cấu xạ
$T\to X_{(\vphi)}$, rồi $g'=p'\circ g''$, trong đó $g''$ là một
$S''$-cấu xạ $T\to(X_{(\vphi)})_{(\vphi')}$:
\[
  \xymatrix{
    X\ar[d]_\pi &
    X_{(\vphi)}\ar[l]_p\ar[d]_{\pi'} &
    (X_{(\vphi)})_{(\vphi')}\ar[l]_{p'}\ar[d]_{\pi''}\\
    S &
    S'\ar[l]^\vphi &
    S''\ar[l]^{\vphi'}
  }
\]
\oldpage[I]{110}
Do đó mệnh đề suy ra từ tính duy nhất của nghiệm của một bài toán ánh xạ phổ
dụng.
\end{proof}

Kết quả này còn được biểu thị bằng cách viết đẳng thức (sai khác một đẳng cấu
chính tắc) $(X_{(S')})_{(S'')}=X_{(S'')}$, nếu không có nguy cơ nhầm lẫn, hoặc
còn bằng
\[
\label{I.3.3.9.1}
  (X\times_S S')\times_{S'}S''=X\times_S S'';
  \tag{3.3.9.1}
\]
tính hàm tử của đẳng cấu được xác định trong \sref{I.3.3.9} cũng được biểu thị
bằng công thức về \emph{tính bắc cầu của các ảnh ngược của cấu xạ}
\[
\label{I.3.3.9.2}
  (f_{(S')})_{(S'')}=f_{(S'')}
  \tag{3.3.9.2}
\]
với mọi $S$-cấu xạ $f:X\to Y$.

\begin{corollary}[3.3.10]
\label{I.3.3.10}
Nếu $X$ và $Y$ là hai $S$-tiền lược đồ, thì tồn tại một đẳng cấu chính tắc có
tính hàm tử từ $S'$-tiền lược đồ $X_{(S')}\times_{S'}Y_{(S')}$ lên
$S'$-tiền lược đồ $(X\times_S Y)_{(S')}$. 
\end{corollary}

\begin{proof}
Ta có, sai khác các đẳng cấu chính tắc,
\[
  (X\times_S S')\times_{S'}(Y\times_S S')
  =X\times_S(Y\times_S S')=(X\times_S Y)\times_S S'
\]
theo \sref{I.3.3.9.1} và tính kết hợp của các tích của $S$-tiền lược đồ.
\end{proof}

Tính hàm tử của đẳng cấu được xác định trong Hệ quả~\sref{I.3.3.10} được biểu
thị bởi công thức
\[
\label{I.3.3.10.1}
  (u_{(S')},v_{(S')})_{S'}=((u,v)_S)_{(S')}
  \tag{3.3.10.1}
\]
với mọi cặp $S$-cấu xạ $u:T\to X$, $v:T\to Y$.

Nói cách khác, hàm tử ảnh ngược $X_{(S')}$ \emph{giao hoán với phép lấy
tích}; nó cũng giao hoán với phép lấy tổng \sref{I.3.2.8}.

\begin{corollary}[3.3.11]
\label{I.3.3.11}
Cho $Y$ là một $S$-tiền lược đồ và $f:X\to Y$ là một cấu xạ biến $X$ thành
một $Y$-tiền lược đồ (và do đó cũng thành một $S$-tiền lược đồ). Khi đó tiền
lược đồ $X_{(S')}$ được đồng nhất với tích $X\times_Y Y_{(S')}$, còn phép
chiếu $X\times_Y Y_{(S')}\to Y_{(S')}$ được đồng nhất với $f_{(S')}$. 
\end{corollary}

\begin{proof}
Gọi $\psi:Y\to S$ là cấu xạ cấu trúc của $Y$; ta có biểu đồ giao hoán
\[
  \xymatrix{
    S'\ar[d] &
    Y_{(S')}\ar[l]_{\psi_{(S')}}\ar[d] &
    X_{(S')}\ar[l]_{f_{(S')}}\ar[d]\\
    S &
    Y\ar[l]^\psi &
    X\ar[l]^f
  }
\]
Ta có $Y_{(S')}$ được đồng nhất với $S'_{(\psi)}$, và $X_{(S')}$ với
$S'_{(\psi\circ f)}$; xét đến \sref{I.3.3.9} và \sref{I.3.3.4}, ta suy ra
hệ quả.
\end{proof}

\begin{env}[3.3.12]
\label{I.3.3.12}
Cho $f:X\to X'$ và $g:Y\to Y'$ là hai $S$-cấu xạ vốn là các
\emph{đơn cấu} của các tiền lược đồ (T, I, 1.1); khi đó $f\times_S g$ là một
\emph{đơn cấu}. Thật vậy, nếu $p$ và $q$ là các phép chiếu của
$X\times_S Y$, $p'$ và $q'$ là các phép chiếu của $X'\times_S Y'$, còn $u$
và $v$ là hai $S$-cấu xạ $T\to X\times_S Y$, thì hệ thức
$(f\times_S g)\circ u=(f\times_S g)\circ v$ kéo theo
$p'\circ(f\times_S g)\circ u=p'\circ(f\times_S g)\circ v$, hay nói cách khác,
$f\circ p\circ u=f\circ p\circ v$; vì $f$ là một đơn cấu nên
$p\circ u=p\circ v$. Dùng việc $g$ là một đơn cấu, ta cũng thu được
$q\circ u=q\circ v$, do đó $u=v$.

\oldpage[I]{111}
Suy ra rằng, với mọi phép mở rộng $S'\to S$ của tiền lược đồ cơ sở,
\[
  f_{(S')}:X_{(S')}\to X'_{(S')}
\]
là một đơn cấu.\footnote{\emph{[Dịch.] Bản tiếng Pháp in $Y_{(S')}$; vì
3.3.12 giả sử $f:X\to X'$, đối miền sau phép mở rộng tiền lược đồ cơ sở là
$X'_{(S')}$.}}
\end{env}

\begin{env}[3.3.13]
\label{I.3.3.13}
Cho $S$ và $S'$ là hai lược đồ affine, và cho $A$ và $A'$ là các vành tương
ứng của chúng; do đó một cấu xạ $S'\to S$ tương ứng với một đồng cấu vành
$A\to A'$. Nếu $X$ là một $S$-tiền lược đồ, ta cũng dùng các ký hiệu
$X_{(A')}$ hoặc $X\otimes_A A'$ cho $S'$-tiền lược đồ $X_{(S')}$. Khi bản thân
$X$ là một lược
đồ affine có vành $B$, thì $X_{(A')}$ là một lược đồ affine có vành
$B_{(A')}=B\otimes_A A'$, thu được bằng cách mở rộng vô hướng của $A$-đại số
$B$ lên $A'$.
\end{env}

\begin{env}[3.3.14]
\label{I.3.3.14}
Với các ký hiệu của \sref{I.3.3.6}, với mọi \emph{$S$-cấu xạ}
$f:S'\to X$, ta có $f'=(f,1_{S'})_S$ là một $S'$-cấu xạ
$S'\to X'=X_{(S')}$ sao cho $p\circ f'=f$, $\pi'\circ f'=1_{S'}$, hay nói
cách khác, là một \emph{$S'$-nhát cắt của $X'$}; ngược lại, nếu $f'$ là một
$S'$-nhát cắt như thế, thì $f=p\circ f'$ là một $S$-cấu xạ $S'\to X$. Do đó
ta xác định một \emph{tương ứng song ánh} chính tắc
\[
  \Hom_S(S',X)\isoto\Hom_{S'}(S',X').
\]
Ta nói $f'$ là \emph{cấu xạ đồ thị} của $f$, và ký hiệu nó bởi
$\Gamma_f$.
\end{env}

\begin{env}[3.3.15]
\label{I.3.3.15}
Cho một tiền lược đồ $X$, mà ta luôn có thể xem là một
$\bb{Z}$-tiền lược đồ, từ \sref{I.3.3.14} suy ra đặc biệt rằng các
\emph{$X$-tiết diện} của $X\otimes_\bb{Z}\bb{Z}[T]$ (trong đó $T$ là một
biến bất định) tương ứng song ánh với các \emph{cấu xạ}
$X\to\bb{Z}[T]$\footnote{\emph{[Dịch.] Văn bản tiếng Pháp in
$\bb{Z}[T]\to X$; Mệnh đề~\sref{I.3.3.14} cùng các đồng cấu vành tiếp theo
$A_\alpha[T]\to A_\alpha$ đòi hỏi cấu xạ lược đồ
$X\to\bb{Z}[T]$.}}. Ta sẽ chứng minh rằng các $X$-tiết diện này cũng tương
ứng song ánh với \emph{các tiết diện của bó cấu trúc $\sh{O}_X$ trên $X$}.
Thật vậy, cho $(U_\alpha)$ là một phủ của $X$ bởi các tập mở affine; cho
$u:X\to X\otimes_\bb{Z}\bb{Z}[T]$ là một $X$-cấu xạ, và cho $u_\alpha$ là
hạn chế của nó xuống $U_\alpha$; nếu $A_\alpha$ là vành của lược đồ affine
$U_\alpha$, thì $U_\alpha\otimes_\bb{Z}\bb{Z}[T]$ là một lược đồ affine có
vành $A_\alpha[T]$ \sref{I.3.2.2}, và $u_\alpha$ tương ứng chính tắc với một
$A_\alpha$-đồng cấu $A_\alpha[T]\to A_\alpha$ \sref{I.1.7.3}. Mà một đồng
cấu như thế được xác định hoàn toàn bởi ảnh của $T$ trong $A_\alpha$, gọi là
$s_\alpha\in A_\alpha=\Gamma(U_\alpha,\sh{O}_X)$; và nếu ta viết rằng các hạn
chế của $u_\alpha$ và $u_\beta$ xuống một tập mở affine
$V\subset U_\alpha\cap U_\beta$ trùng nhau, thì thấy ngay rằng $s_\alpha$ và
$s_\beta$ trùng nhau trên $V$; do đó họ $(s_\alpha)$ gồm các hạn chế xuống
$U_\alpha$ của một tiết diện $s$ của $\sh{O}_X$ trên $X$; ngược lại, hiển
nhiên một tiết diện như thế xác định một họ $(u_\alpha)$ các cấu xạ là những
hạn chế xuống $U_\alpha$ của một $X$-cấu xạ
$X\to X\otimes_\bb{Z}\bb{Z}[T]$. Kết quả này sẽ được tổng quát hóa trong
\sref[II]{II.1.7.12}.
\end{env}

\subsection{Các điểm của một tiền lược đồ nhận giá trị trong một tiền lược đồ; các điểm hình học}
\label{subsection:I.3.4}

\begin{env}[3.4.1]
\label{I.3.4.1}
Cho $X$ là một tiền lược đồ; với mọi tiền lược đồ $T$, ta vẫn ký hiệu bởi
$X(T)$ tập $\Hom(T,X)$ các cấu xạ $T\to X$, và các phần tử của tập này cũng
được gọi là \emph{các điểm của $X$ nhận giá trị trong $T$}. Nếu với mỗi cấu
xạ $f:T\to T'$ ta gán ánh xạ $u'\mapsto u'\circ f$ từ $X(T')$ vào $X(T)$,
thì với $X$ cố định, ta thấy $X(T)$ là một \emph{hàm tử phản biến theo $T$},
từ phạm trù các tiền lược đồ đến phạm trù các tập hợp. Hơn nữa, mỗi cấu xạ
của các tiền lược đồ $g:X\to Y$ xác định một đồng cấu có tính hàm tử
$X(T)\to Y(T)$, gửi $v\in X(T)$ thành $g\circ v$.
\end{env}

\begin{env}[3.4.2]
\label{I.3.4.2}
Cho ba tập hợp $P$, $Q$ và $R$, cùng hai ánh xạ $\vphi:P\to R$ và
$\psi:Q\to R$, ta gọi \emph{tích thớ của $P$ và $Q$ trên $R$} (đối với
$\vphi$ và $\psi$) là tập con của
\oldpage[I]{112}
tập tích $P\times Q$ gồm các cặp $(p,q)$ sao cho $\vphi(p)=\psi(q)$; ta ký
hiệu nó bởi $P\times_R Q$. Định nghĩa~\sref{I.3.2.1} về tích của các
$S$-tiền lược đồ cũng có thể được diễn giải, với ký hiệu của
\sref{I.3.4.1}, bởi công thức
\[
\label{I.3.4.2.1}
  (X\times_S Y)(T)=X(T)\times_{S(T)}Y(T),
  \tag{3.4.2.1}
\]
trong đó các ánh xạ $X(T)\to S(T)$ và $Y(T)\to S(T)$ tương ứng với các cấu
xạ cấu trúc $X\to S$ và $Y\to S$.
\end{env}

\begin{env}[3.4.3]
\label{I.3.4.3}
Cho một tiền lược đồ $S$ và chỉ xét các $S$-tiền lược đồ cùng các $S$-cấu xạ,
ta vẫn ký hiệu bởi $X(T)_S$ tập hợp $\Hom_S(T,X)$ các $S$-cấu xạ $T\to X$,
và có thể bỏ chỉ số $S$ khi không có khả năng nhầm lẫn; ta vẫn gọi các phần tử
của $X(T)_S$ là các \emph{điểm} (hoặc các \emph{$S$-điểm}, khi có khả năng
nhầm lẫn) \emph{của $S$-tiền lược đồ $X$ nhận giá trị trong $S$-tiền lược đồ
$T$}. Đặc biệt, một \emph{$S$-nhát cắt} của $X$ không gì khác hơn là một
\emph{điểm của $X$ nhận giá trị trong $S$}. Khi đó công thức
\sref{I.3.4.2.1} cũng được viết thành
\[
\label{I.3.4.3.1}
  (X\times_S Y)(T)_S=X(T)_S\times Y(T)_S;
  \tag{3.4.3.1}
\]
tổng quát hơn, nếu $Z$ là một $S$-tiền lược đồ, còn $X$, $Y$ và $T$ là các
$Z$-tiền lược đồ (do đó, \emph{ipso facto}, là các $S$-tiền lược đồ), thì
\[
\label{I.3.4.3.2}
  (X\times_Z Y)(T)_S=X(T)_S\times_{Z(T)_S}Y(T)_S.
  \tag{3.4.3.2}
\]

Cần lưu ý rằng, để chứng minh một bộ ba $(W,r,s)$ gồm một $S$-tiền lược đồ
$W$ và hai $S$-cấu xạ $r:W\to X$, $s:W\to Y$ là một tích của $X$ và $Y$
(trên $Z$), theo định nghĩa chỉ cần kiểm tra rằng, với \emph{mọi $S$-tiền
lược đồ $T$}, biểu đồ
\[
  \xymatrix{
    W(T)_S\ar[r]^{r'}\ar[d]_{s'} &
    X(T)_S\ar[d]^{\vphi'}\\
    Y(T)_S\ar[r]_{\psi'} &
    Z(T)_S
  }
\]
làm cho $W(T)_S$ trở thành tích thớ của $X(T)_S$ và $Y(T)_S$ trên $Z(T)_S$,
trong đó $r'$ và $s'$ tương ứng với $r$ và $s$, còn $\vphi'$ và $\psi'$ tương
ứng với các cấu xạ cấu trúc $\vphi:X\to Z$ và $\psi:Y\to Z$.
\end{env}

\begin{env}[3.4.4]
\label{I.3.4.4}
Khi, trong phần trên, $T$ (tương ứng $S$) là một lược đồ affine có vành $B$
(tương ứng $A$), ta thay $T$ (tương ứng $S$) bằng $B$ (tương ứng $A$) trong
các ký hiệu trước đó, và gọi các phần tử của $X(B)$ là \emph{các điểm của $X$
nhận giá trị trong vành $B$}, còn các phần tử của $X(B)_A$ là \emph{các điểm
của $A$-tiền lược đồ $X$ nhận giá trị trong $A$-đại số $B$}. Ta lưu ý rằng
$X(B)$ và $X(B)_A$ lúc này là các \emph{hàm tử hiệp biến theo $B$}. Ta cũng
viết $X(T)_A$ cho tập hợp các điểm của $A$-tiền lược đồ $X$ nhận giá trị
trong $A$-tiền lược đồ $T$.
\end{env}

\begin{env}[3.4.5]
\label{I.3.4.5}
Xét riêng trường hợp $T$ có dạng $\Spec(A)$, trong đó $A$ là một vành
\emph{địa phương}; các phần tử của $X(A)$ tương ứng song ánh với các
\emph{đồng cấu địa phương} $\sh{O}_x\to A$ với $x\in X$ \sref{I.2.4.4}; ta
gọi điểm $x$ của không gian nền của $X$ là \emph{vị trí}\footnote{\emph{[Dịch.]
Ta cũng nói điểm hình học nằm trên $x$ này.}} của điểm của $X$ nhận giá trị
trong $A$ tương ứng với nó.

Cụ thể hơn, ta gọi \emph{các điểm hình học} của một tiền lược đồ $X$ là
\emph{các điểm của $X$ nhận giá trị trong một trường $K$}: việc cho một điểm
như thế tương đương với việc cho
\oldpage[I]{113}
vị trí $x$ của nó trong không gian nền của $X$ và một \emph{mở rộng} $K$ của
$\kres(x)$; $K$ được gọi là \emph{trường giá trị} của điểm hình học tương ứng,
và ta còn nói điểm hình học ấy được \emph{địa phương hóa tại $x$}. Do đó ta
xác định một ánh xạ $X(K)\to X$, gửi một điểm hình học nhận giá trị trong $K$
đến vị trí của nó.

Nếu $S'=\Spec(K)$ là một $S$-tiền lược đồ (nói cách khác, nếu $K$ được xem là
một mở rộng của trường thặng dư $\kres(s)$, với $s\in S$), và nếu $X$ là một
$S$-tiền lược đồ, thì một phần tử của $X(K)_S$, hay, như ta còn nói, một
\emph{điểm hình học của $X$ nằm trên $s$ nhận giá trị trong $K$}, gồm dữ liệu
của một $\kres(s)$-đơn cấu từ trường thặng dư $\kres(x)$ vào $K$, trong đó $x$
là một điểm của $X$ \emph{nằm trên $s$} (do đó $\kres(x)$ là một mở rộng của
$\kres(s)$).

Đặc biệt, nếu $S=\Spec(K)=\{\xi\}$, thì \emph{các điểm hình học của $X$ nhận
giá trị trong $K$ được đồng nhất với các điểm $x\in X$ sao cho $\kres(x)=K$};
ta còn gọi các điểm sau là \emph{các điểm hữu tỉ trên $K$ của $K$-tiền lược
đồ $X$}; nếu $K'$ là một mở rộng của $K$, thì các điểm hình học của $X$ nhận
giá trị trong $K'$ tương ứng song ánh với các điểm $K'$-hữu tỉ của
$X'=X_{(K')}$~\sref{I.3.3.14}.
\end{env}

\begin{lemma}[3.4.6]
\label{I.3.4.6}
Cho $X_i$ ($1\leq i\leq n$) là các $S$-tiền lược đồ, $s$ là một điểm của
$S$, và $x_i$ ($1\leq i\leq n$) là các điểm của $X_i$ nằm trên $s$. Khi đó
tồn tại một mở rộng $K$ của $\kres(s)$ và một điểm hình học của tích
$Y=X_1\times_S X_2\times_S\cdots\times_S X_n$, nhận giá trị trong $K$, sao
cho các phép chiếu của nó lên các $X_i$ được địa phương hóa tại các $x_i$.
\end{lemma}

\begin{proof}
Tồn tại các $\kres(s)$-đơn cấu $\kres(x_i)\to K$, tất cả cùng vào một mở
rộng $K$ của $\kres(s)$ (Bourbaki, \emph{Alg.}, chương~V, \textsection4,
mệnh đề~2). Các hợp thành $\kres(s)\to\kres(x_i)\to K$ đều giống nhau; do
đó các cấu xạ $\Spec(K)\to X_i$ tương ứng với các đồng cấu
$\kres(x_i)\to K$ đều là $S$-cấu xạ, và suy ra chúng xác định một cấu xạ duy
nhất $\Spec(K)\to Y$. Nếu $y$ là điểm tương ứng của $Y$, thì hiển nhiên phép
chiếu của nó lên mỗi $X_i$ là $x_i$.
\end{proof}

\begin{proposition}[3.4.7]
\label{I.3.4.7}
Cho $X_i$ ($1\leq i\leq n$) là các $S$-tiền lược đồ và, với mỗi chỉ số $i$,
cho $x_i$ là một điểm của $X_i$. Điều kiện cần và đủ để tồn tại một điểm $y$
của $Y=X_1\times_S X_2\times\ldots\times_S X_n$ có ảnh là $x_i$ qua phép
chiếu thứ $i$ với mọi $1\leq i\leq n$ là tất cả các $x_i$ đều nằm trên cùng
một điểm $s$ của $S$.
\end{proposition}

\begin{proof}
Điều kiện rõ ràng là cần; Bổ đề~\sref{I.3.4.6} chứng minh rằng nó là đủ.
\end{proof}

Nói cách khác, nếu ký hiệu $(X)$ là tập nền của $X$, ta thấy có một ánh xạ
chính tắc \emph{toàn ánh} $(X\times_S Y)\to(X)\times_{(S)}(Y)$; cần lưu ý
rằng ánh xạ này nói chung \emph{không đơn ánh}; nói cách khác, \emph{có thể
tồn tại nhiều điểm phân biệt $z$ trong $X\times_S Y$ có cùng các hình chiếu
$x\in X$ và $y\in Y$}; điều này đã thấy ngay khi $S$, $X$ và $Y$ là các phổ
nguyên tố của các trường $k$, $K$ và $K'$, tương ứng, vì tích tenxơ
$K\otimes_k K'$ nói chung có nhiều iđêan nguyên tố phân biệt
(xem~\sref{I.3.4.9}).

\begin{corollary}[3.4.8]
\label{I.3.4.8}
Cho $f:X\to Y$ là một $S$-cấu xạ, và cho
$f_{(S')}:X_{(S')}\to Y_{(S')}$ là $S'$-cấu xạ suy ra từ $f$ bằng một phép
mở rộng $S'\to S$ của tiền lược đồ cơ sở. Gọi $p$ (\emph{tương ứng, $q$}) là
phép chiếu $X_{(S')}\to X$ (\emph{tương ứng, $Y_{(S')}\to Y$}); với mọi tập
con $M$ của $X$, ta có
\[
  q^{-1}(f(M))=f_{(S')}(p^{-1}(M)).
\]
\end{corollary}

\oldpage[I]{114}
\begin{proof}
Thật vậy, theo \sref{I.3.3.11}, $X_{(S')}$ được đồng nhất với tích
$X\times_Y Y_{(S')}$ nhờ biểu đồ giao hoán
\[
  \xymatrix{
    X\ar[d]_f &
    X_{(S')}\ar[l]_p\ar[d]^{f_{(S')}}\\
    Y &
    Y_{(S')}\ar[l]^q
  }
\]
Theo \sref{I.3.4.7}, hệ thức $q(y')=f(x)$ với $x\in M$ và
$y'\in Y_{(S')}$ tương đương với sự tồn tại của một $x'\in X_{(S')}$ sao
cho $p(x')=x$ và $f_{(S')}(x')=y'$, từ đó suy ra hệ quả.
\end{proof}

Bổ đề~\sref{I.3.4.6} có thể được phát biểu chính xác hơn như sau:

\begin{proposition}[3.4.9]
\label{I.3.4.9}
Cho $X$ và $Y$ là hai $S$-tiền lược đồ, $x$ là một điểm của $X$, $y$ là một
điểm của $Y$, còn $x$ và $y$ đều nằm trên cùng một điểm $s\in S$.
Tập các điểm của $X\times_S Y$ có các hình chiếu là $x$ và $y$ tương ứng song
ánh chính tắc với tập các kiểu mở rộng hợp thành của $\kres(x)$ và $\kres(y)$,
được xét như các mở rộng của $\kres(s)$ \emph{(Bourbaki, \emph{Alg.},
chương~VIII, \textsection8, mệnh đề~2)}.
\end{proposition}

\begin{proof}
Gọi $p$ (tương ứng, $q$) là phép chiếu từ $X\times_S Y$ lên $X$ (tương ứng,
$Y$), và gọi $E$ là không gian con $p^{-1}(x)\cap q^{-1}(y)$ của không gian
nền của $X\times_S Y$.
Trước hết, lưu ý rằng các cấu xạ $\Spec(\kres(x))\to S$ và
$\Spec(\kres(y))\to S$ phân tích qua
$\Spec(\kres(x))\to\Spec(\kres(s))\to S$ và
$\Spec(\kres(y))\to\Spec(\kres(s))\to S$; vì
$\Spec(\kres(s))\to S$ là một đơn cấu \sref{I.2.4.7}, từ
\sref{I.3.2.4} suy ra rằng
\[
  P=\Spec(\kres(x))\times_S\Spec(\kres(y))=\Spec(\kres(x))\times_{\Spec(\kres(s))}\Spec(\kres(y))=\Spec(\kres(x)\otimes_{\kres(s)}\kres(y)).
\]
Ta sẽ xác định hai ánh xạ $\alpha:P_0\to E$ và $\beta:E\to P_0$ nghịch đảo
của nhau (trong đó $P_0$ ký hiệu tập nền của tiền lược đồ $P$).
Nếu $i:\Spec(\kres(x))\to X$ và $j:\Spec(\kres(y))\to Y$ là các cấu xạ chính
tắc \sref{I.2.4.5}, ta lấy $\alpha$ là ánh xạ giữa các không gian nền tương ứng
với cấu xạ $i\times_S j$.
Mặt khác, mỗi $z\in E$ xác định, theo giả thiết, hai $\kres(s)$-đơn cấu
$\kres(x)\to\kres(z)$ và $\kres(y)\to\kres(z)$, do đó một
$\kres(s)$-đồng cấu\footnote{\emph{[Dịch.] Văn bản tiếng Pháp in «đơn cấu» ở
đây. Ánh xạ cảm sinh từ tích tenxơ không nhất thiết đơn ánh; chỉ cần một đồng
cấu để xác định cấu xạ giữa các phổ tiếp theo.}}
$\kres(x)\otimes_{\kres(s)}\kres(y)\to\kres(z)$ và vì thế một cấu xạ
$\Spec(\kres(z))\to P$; $\beta(z)$ sẽ là ảnh của $z$ trong $P_0$ qua cấu xạ
này.
Việc kiểm tra rằng $\alpha\circ\beta$ và $\beta\circ\alpha$ là các ánh xạ đồng
nhất suy ra từ \sref{I.2.4.5} và định nghĩa của tích \sref{I.3.2.1}.
Cuối cùng, ta biết rằng $P_0$ tương ứng song ánh với tập các kiểu mở rộng hợp
thành của $\kres(x)$ và $\kres(y)$ (Bourbaki, \emph{Alg.}, chương~VIII,
\textsection8, mệnh đề~1).
\end{proof}

\subsection{Tính toàn ánh và tính đơn ánh}
\label{subsection:I.3.5}

\begin{env}[3.5.1]
\label{I.3.5.1}
Một cách tổng quát, xét một tính chất $\textbf{P}$ của các cấu xạ giữa các
tiền lược đồ và hai mệnh đề sau:
\begin{enumerate}
  \item[(i)] Nếu $f:X\to X'$ và $g:Y\to Y'$ là hai $S$-cấu xạ có tính chất
  $\textbf{P}$, thì $f\times_S g$ cũng có tính chất $\textbf{P}$.
  \item[(ii)] Nếu $f:X\to Y$ là một $S$-cấu xạ có tính chất $\textbf{P}$, thì
  mọi $S'$-cấu xạ $f_{(S')}:X_{(S')}\to Y_{(S')}$ thu được từ $f$ bằng một
  phép mở rộng tiền lược đồ cơ sở $S'\to S$ cũng có tính chất $\textbf{P}$.
\end{enumerate}

\oldpage[I]{115}
Vì $f_{(S')}=f\times_S 1_{S'}$, ta thấy rằng, nếu với mọi tiền lược đồ $X$,
\emph{ánh xạ đồng nhất} $1_X$ có tính chất $\textbf{P}$, thì (i) kéo theo (ii); mặt
khác, vì $f\times_S g$ là cấu xạ hợp thành
\[
  X\times_S Y\xrightarrow{f\times1_Y}X'\times_S Y\xrightarrow{1_{X'}\times g}X'\times_S Y',
\]
ta thấy rằng, nếu \emph{hợp thành} của hai cấu xạ có tính chất $\textbf{P}$
cũng có tính chất $\textbf{P}$, thì tích $f\times_S g$ cũng có tính chất ấy,
và do đó (ii) kéo theo (i).
\end{env}

Một ứng dụng đầu tiên của nhận xét này là mệnh đề sau.

\begin{proposition}[3.5.2]
\label{I.3.5.2}
\medskip\noindent
\begin{enumerate}
  \item[{\rm(i)}] Nếu $f:X\to X'$ và $g:Y\to Y'$ là các $S$-cấu xạ toàn
  ánh, thì $f\times_S g$ là toàn ánh.
  \item[{\rm(ii)}] Nếu $f:X\to Y$ là một $S$-cấu xạ toàn ánh, thì
  $f_{(S')}$ là toàn ánh đối với mọi phép mở rộng tiền lược đồ cơ sở
  $S'\to S$.
\end{enumerate}
\end{proposition}

\begin{proof}
Vì hợp thành của hai cấu xạ toàn ánh là toàn ánh, chỉ cần chứng minh (ii); nhưng
mệnh đề này suy ra ngay từ \sref{I.3.4.8} áp dụng cho $M=X$.
\end{proof}

\begin{proposition}[3.5.3]
\label{I.3.5.3}
Để một cấu xạ $f:X\to Y$ là toàn ánh, điều kiện cần và đủ là, với mọi trường
$K$ và mọi cấu xạ $\Spec(K)\to Y$, tồn tại một mở rộng $K'$ của $K$ và một
cấu xạ $\Spec(K')\to X$ sao cho sơ đồ sau giao hoán:
\[
  \xymatrix{
    X\ar[d]_f &
    \Spec(K')\ar[l]\ar[d]\\
    Y &
    \Spec(K).\ar[l]
  }
\]
\end{proposition}

\begin{proof}
Điều kiện là đủ vì, với mọi $y\in Y$, chỉ cần áp dụng nó cho một cấu xạ
$\Spec(K)\to Y$ tương ứng với một đơn cấu $\kres(y)\to K$, trong đó $K$ là
một mở rộng của $\kres(y)$ \sref{I.2.4.6}.
Ngược lại, giả sử $f$ toàn ánh, và gọi $y\in Y$ là ảnh của điểm duy nhất của
$\Spec(K)$; tồn tại một $x\in X$ sao cho $f(x)=y$; xét đơn cấu tương ứng
$\kres(y)\to\kres(x)$ \sref{I.2.2.1}; khi đó chỉ cần lấy $K'$ là một mở rộng
của $\kres(y)$ sao cho tồn tại các $\kres(y)$-đơn cấu từ $\kres(x)$ và từ
$K$ vào $K'$ (Bourbaki, \emph{Alg.}, chương~V, \textsection4, mệnh đề~2);
cấu xạ $\Spec(K')\to X$ tương ứng với $\kres(x)\to K'$ chính là cấu xạ cần
tìm.
\end{proof}

Với ngôn ngữ được đưa vào trong \sref{I.3.4.5}, ta còn có thể nói rằng
\emph{mọi điểm hình học của $Y$ nhận giá trị trong $K$ đều bắt nguồn từ một
điểm hình học của $X$ nhận giá trị trong một mở rộng của $K$}.

\begin{definition}[3.5.4]
\label{I.3.5.4}
Ta nói một cấu xạ tiền lược đồ $f:X\to Y$ là \emph{đơn ánh phổ dụng}, hay là
một \emph{cấu xạ radicial}, nếu, với mọi trường $K$, ánh xạ tương ứng
$X(K)\to Y(K)$ là đơn ánh.
\end{definition}

Từ các định nghĩa suy ra ngay rằng mọi \emph{đơn cấu tiền lược đồ} (T,~1.1)
đều là radicial.

\begin{env}[3.5.5]
\label{I.3.5.5}
Để một cấu xạ $f:X\to Y$ là radicial, chỉ cần điều kiện của Định
nghĩa~\sref{I.3.5.4} được thỏa mãn với mọi trường \emph{đóng đại số}.
Thật vậy, nếu $K$ là một trường tùy ý và $K'$ là một mở rộng của $K$ đồng
thời là trường đóng đại số, thì sơ đồ
\[
  \xymatrix{
    X(K)\ar[r]^\alpha\ar[d]_\vphi &
    Y(K)\ar[d]^{\vphi'}\\
    X(K')\ar[r]_{\alpha'} &
    Y(K')
  }
\]
\oldpage[I]{116}
giao hoán, trong đó $\vphi$ và $\vphi'$ được cảm sinh bởi cấu xạ
$\Spec(K')\to\Spec(K)$, còn $\alpha$ và $\alpha'$ tương ứng với $f$.
Mặt khác, $\vphi$ là đơn ánh, và $\alpha'$ cũng vậy theo giả thiết; do đó
$\alpha$ nhất thiết là đơn ánh.
\end{env}

\begin{proposition}[3.5.6]
\label{I.3.5.6}
Cho $f:X\to Y$ và $g:Y\to Z$ là hai cấu xạ tiền lược đồ.
\begin{enumerate}
  \item[{\rm(i)}] Nếu $f$ và $g$ là radicial, thì $g\circ f$ cũng là
  radicial.
  \item[{\rm(ii)}] Ngược lại, nếu $g\circ f$ là radicial, thì $f$ cũng là
  radicial.
\end{enumerate}
\end{proposition}

\begin{proof}
Theo Định nghĩa~\sref{I.3.5.4}, mệnh đề được quy về các khẳng định tương ứng
cho các ánh xạ $X(K)\to Y(K)\to Z(K)$, và các khẳng định ấy là hiển nhiên.
\end{proof}

\begin{proposition}[3.5.7]
\label{I.3.5.7}
\medskip\noindent
\begin{enumerate}
  \item[{\rm(i)}] Nếu các $S$-cấu xạ $f:X\to X'$ và $g:Y\to Y'$ là
  radicial, thì $f\times_S g$ cũng là radicial.
  \item[{\rm(ii)}] Nếu $S$-cấu xạ $f:X\to Y$ là radicial, thì
  $f_{(S')}:X_{(S')}\to Y_{(S')}$ cũng là radicial đối với mọi phép mở rộng
  $S'\to S$ của tiền lược đồ cơ sở.
\end{enumerate}
\end{proposition}

\begin{proof}
Theo \sref{I.3.5.1}, chỉ cần chứng minh (i).
Ta đã thấy trong \sref{I.3.4.2.1} rằng
\[
  (X\times_S Y)(K)=X(K)\times_{S(K)}Y(K),
\]
\[
  (X'\times_S Y')(K)=X'(K)\times_{S(K)}Y'(K),
\]
khi đó ánh xạ $(X\times_S Y)(K)\to(X'\times_S Y')(K)$ tương ứng với
$f\times_S g$ được đồng nhất với $(u,v)\to(f\circ u,g\circ v)$, và mệnh đề
suy ra ngay.
\end{proof}

\begin{proposition}[3.5.8]
\label{I.3.5.8}
Để một cấu xạ $f=(\psi,\theta):X\to Y$ là radicial, điều kiện cần và đủ là
$\psi$ đơn ánh và đơn cấu $\theta^x:\kres(\psi(x))\to\kres(x)$ làm cho
$\kres(x)$ thành một mở rộng radicial của $\kres(\psi(x))$ với mọi
$x\in X$.
\end{proposition}

\begin{proof}
Giả sử $f$ là radicial và trước hết chứng minh rằng đẳng thức
$\psi(x_1)=\psi(x_2)=y$ nhất thiết kéo theo $x_1=x_2$.
Thật vậy, tồn tại một trường $K$ mở rộng $\kres(y)$, cùng với các
$\kres(y)$-đơn cấu $\kres(x_1)\to K$ và $\kres(x_2)\to K$ (Bourbaki,
\emph{Alg.}, chương~V, \textsection4, mệnh đề~2); khi đó các cấu xạ tương ứng
$u_1:\Spec(K)\to X$ và $u_2:\Spec(K)\to X$ thỏa
$f\circ u_1=f\circ u_2$, nên theo giả thiết $u_1=u_2$, điều này đặc biệt
kéo theo $x_1=x_2$.
Bây giờ xét $\kres(x)$ như một mở rộng của $\kres(\psi(x))$ thông qua
$\theta^x$: nếu $\kres(x)$ không phải là một mở rộng radicial của
$\kres(\psi(x))$, thì tồn tại hai $\kres(\psi(x))$-đơn cấu phân biệt từ
$\kres(x)$ vào một trường đóng đại số $K$ mở rộng $\kres(\psi(x))$, và hai
cấu xạ tương ứng $\Spec(K)\to X$ sẽ mâu thuẫn với giả thiết.
Ngược lại, theo \sref{I.2.4.6}, hiển nhiên các điều kiện trong phát biểu là
đủ để $f$ là radicial.
\end{proof}

\begin{corollary}[3.5.9]
\label{I.3.5.9}
Nếu $A$ là một vành và $S$ là một tập nhân của $A$, thì cấu xạ chính tắc
$\Spec(S^{-1}A)\to\Spec(A)$ là radicial.
\end{corollary}

\begin{proof}
Thật vậy, cấu xạ này là một đơn cấu \sref{I.1.6.2}.
\end{proof}

\begin{corollary}[3.5.10]
\label{I.3.5.10}
Cho $f:X\to Y$ là một cấu xạ radicial, $g:Y'\to Y$ là một cấu xạ và
$X'=X_{(Y')}=X\times_Y Y'$.
Khi đó cấu xạ radicial $f_{(Y')}$ \sref{I.3.5.7}[ii] là một song ánh từ
không gian nền của $X'$ lên $g^{-1}(f(X))$; hơn nữa, với mọi trường $K$, tập
$X'(K)$ được đồng nhất với ảnh ngược của tập con $X(K)\subseteq Y(K)$ qua
ánh xạ $Y'(K)\to Y(K)$ tương ứng với $g$.
\end{corollary}

\begin{proof}
Khẳng định thứ nhất suy ra ngay từ \sref{I.3.5.8} và \sref{I.3.4.8}; khẳng
định thứ hai suy ra từ tính giao hoán của biểu đồ sau:
\oldpage[I]{117}
\[
  \xymatrix{
    X'(K)\ar[r]\ar[d] &
    Y'(K)\ar[d]\\
    X(K)\ar[r] &
    Y(K)
  }
\]
\end{proof}

\begin{remark}[3.5.11]
\label{I.3.5.11}
Ta nói một cấu xạ $f=(\psi,\theta)$ của các tiền lược đồ là
\emph{đơn ánh} nếu ánh xạ $\psi$ là đơn ánh.
Một cấu xạ $f=(\psi,\theta):X\to Y$ là radicial khi và chỉ khi, với mọi cấu
xạ $Y'\to Y$, cấu xạ $f_{(Y')}:X_{(Y')}\to Y'$ là đơn ánh (điều này giải
thích thuật ngữ cấu xạ \emph{đơn ánh phổ dụng}).
Thật vậy, điều kiện này là cần thiết theo
\sref{I.3.5.7}[ii] và \sref{I.3.5.8}.
Ngược lại, điều kiện trước hết kéo theo $\psi$ là đơn ánh; nếu, với một
$x\in X$, đơn cấu $\theta^x:\kres(\psi(x))\to\kres(x)$ không radicial, thì
sẽ tồn tại một mở rộng $K$ của $\kres(\psi(x))$ và hai cấu xạ phân biệt
$\Spec(K)\to X$ tương ứng với cùng một cấu xạ $\Spec(K)\to Y$
\sref{I.3.5.8}.
Nhưng khi đó, đặt $Y'=\Spec(K)$, sẽ có hai $Y'$-nhát cắt phân biệt của
$X_{(Y')}$ \sref{I.3.3.14}, trái với giả thiết rằng $f_{(Y')}$ là đơn ánh.
\end{remark}

\subsection{Các thớ}
\label{subsection:I.3.6}

\begin{proposition}[3.6.1]
\label{I.3.6.1}
Cho $f:X\to Y$ là một cấu xạ, $y$ là một điểm của $Y$, và $\mathfrak{a}_y$ là một iđêan định nghĩa của $\sh{O}_y$ đối với tôpô $\mathfrak{m}_y$-tiền-adic.
Khi đó phép chiếu $p:X\times_Y\Spec(\sh{O}_y/\mathfrak{a}_y)\to X$ là một phép đồng phôi từ không gian nền của tiền lược đồ $X\times_Y\Spec(\sh{O}_y/\mathfrak{a}_y)$ lên thớ $f^{-1}(y)$, được trang bị tôpô cảm sinh từ tôpô của không gian nền của $X$.
\end{proposition}

\begin{proof}
Vì $\Spec(\sh{O}_y/\mathfrak{a}_y)\to Y$ là radicial (\sref{I.3.5.4} và \sref{I.2.4.7}), vì $\Spec(\sh{O}_y/\mathfrak{a}_y)$ chỉ gồm một điểm, và vì iđêan $\mathfrak{m}_y/\mathfrak{a}_y$ là lũy linh theo giả thiết \sref{I.1.1.12}, ta đã biết (\sref{I.3.5.10} và \sref{I.3.3.4}) rằng $p$ đồng nhất, xét như các \emph{tập hợp}, không gian nền của $X\times_Y\Spec(\sh{O}_y/\mathfrak{a}_y)$ với $f^{-1}(y)$; do đó chỉ còn chứng minh rằng $p$ là một phép đồng phôi.
Theo \sref{I.3.2.7}, vấn đề có tính địa phương trên $X$ và $Y$, nên ta có thể giả sử $X=\Spec(B)$ và $Y=\Spec(A)$, với $B$ là một $A$-đại số.
Cấu xạ $p$ khi đó tương ứng với đồng cấu $1\otimes\vphi:B\to B\otimes_A A'$, trong đó $A'=A_y/\mathfrak{a}_y$ và $\vphi$ là ánh xạ chính tắc từ $A$ đến $A'$.
Mọi phần tử của $B\otimes_A A'$ khi đó đều viết được dưới dạng
\[
  \sum_i b_i\otimes\vphi(a_i)/\vphi(s)=\left(\sum_i(a_ib_i\otimes1)\right)(1\otimes\vphi(s))^{-1},
\]
trong đó $s\not\in\mathfrak{j}_y$, và áp dụng được Mệnh đề~\sref{I.1.2.4}.
\end{proof}

\begin{env}[3.6.2]
\label{I.3.6.2}
Trong toàn bộ phần còn lại của Chuyên luận này, khi ta xét một thớ $f^{-1}(y)$ của một cấu xạ như được trang bị cấu trúc $\kres(y)$-tiền lược đồ, \emph{đó luôn là tiền lược đồ thu được bằng cách chuyển cấu trúc của $X\times_Y\Spec(\kres(y))$ qua phép chiếu lên $X$}.
Ta cũng sẽ ký hiệu tích cuối cùng này bởi $X\otimes_Y\kres(y)$, hoặc $X\otimes_{\sh{O}_Y}\kres(y)$; tổng quát hơn, nếu $B$ là một $\sh{O}_y$-đại số, ta ký hiệu bởi $X\otimes_Y B$ hoặc $X\otimes_{\sh{O}_Y}B$ tích $X\times_Y\Spec(B)$.
\end{env}

Với quy ước trên, từ \sref{I.3.5.10} suy ra rằng các điểm của $X$ nhận giá trị trong một mở rộng $K$ của $\kres(y)$ được đồng nhất với các \emph{điểm của $f^{-1}(y)$ nhận giá trị trong $K$}.

\begin{env}[3.6.3]
\label{I.3.6.3}
Cho $f:X\to Y$ và $g:Y\to Z$ là hai cấu xạ, và $h=g\circ f$ là hợp thành của chúng; với mọi $z\in Z$, thớ $h^{-1}(z)$ là một tiền lược đồ đẳng cấu với
\[
  X\times_Z\Spec(\kres(z))=(X\times_Y Y)\times_Z\Spec(\kres(z))=X\times_Y g^{-1}(z)
\]
\oldpage[I]{118}
Đặc biệt, nếu $U$ là một phần mở của $X$, tiền lược đồ cảm sinh trên $U\cap f^{-1}(y)$ bởi tiền lược đồ $f^{-1}(y)$ đẳng cấu với $f^{-1}_U(y)$ (trong đó $f_U$ là hạn chế của $f$ lên $U$).
\end{env}

\begin{proposition}[3.6.4]
\label{I.3.6.4}
\emph{(``Tính bắc cầu của các thớ'')}. Cho $f:X\to Y$ và $g:Y'\to Y$ là hai cấu xạ; đặt $X'=X\times_Y Y'=X_{(Y')}$ và $f'=f_{(Y')}:X'\to Y'$.
Với mọi $y'\in Y'$, nếu đặt $y=g(y')$, tiền lược đồ $f'^{-1}(y')$ đẳng cấu với $f^{-1}(y)\otimes_{\kres(y)}\kres(y')$.
\end{proposition}

\begin{proof}
Thật vậy, chỉ cần nhận thấy rằng hai tiền lược đồ $(X\otimes_Y\kres(y))\otimes_{\kres(y)}\kres(y')$ và $(X\times_Y Y')\otimes_{Y'}\kres(y')$ đều đẳng cấu chính tắc với $X\times_Y\Spec(\kres(y'))$ theo \sref{I.3.3.9.1}.
\end{proof}

Đặc biệt, nếu $V$ là một lân cận mở của $y$ trong $Y$, và nếu ký hiệu $f_V$ là hạn chế của $f$ lên tiền lược đồ cảm sinh trên $f^{-1}(V)$, thì các tiền lược đồ $f^{-1}(y)$ và $f^{-1}_V(y)$ được đồng nhất chính tắc với nhau.

\begin{proposition}[3.6.5]
\label{I.3.6.5}
Cho $f:X\to Y$ là một cấu xạ, $y$ là một điểm của $Y$, $Z$ là tiền lược đồ địa phương $\Spec(\sh{O}_y)$, và $p=(\psi,\theta)$ là phép chiếu $X\times_Y Z\to X$; khi đó $p$ là một phép đồng phôi từ không gian nền của $X\times_Y Z$ lên không gian con $f^{-1}(Z)$ của $X$ (\emph{khi không gian nền của $Z$ được đồng nhất với một không gian con của $Y$, cf.~\sref{I.2.4.2}}), và, với mọi $t\in X\times_Y Z$, nếu đặt $z=\psi(t)$, thì $\theta_t^\sharp$ là một đẳng cấu từ $\sh{O}_z$ lên $\sh{O}_t$.
\end{proposition}

\begin{proof}
Vì $Z$ (được đồng nhất với một không gian con của $Y$) nằm trong mọi tập mở affine chứa $y$ \sref{I.2.4.2}, nên, như trong \sref{I.3.6.1}, ta có thể quy về trường hợp $X=\Spec(A)$ và $Y=\Spec(B)$ là các lược đồ affine, với $A$ là một $B$-đại số.
Khi đó $X\times_Y Z$ là phổ nguyên tố của $A\otimes_B B_y$, và vành này được đồng nhất chính tắc với $S^{-1}A$, trong đó $S$ là ảnh của $B\setmin\mathfrak{j}_y$ trong $A$ \sref[0]{0.1.5.2}; vì $p$ khi ấy tương ứng với đồng cấu chính tắc $A\to S^{-1}A$, mệnh đề suy ra từ \sref{I.1.6.2}.
\end{proof}

\subsection{Ứng dụng: rút gọn một tiền lược đồ modulo~$\mathfrak{J}$}
\label{subsection:I.3.7}
\label{I.3.7}

\emph{Tiểu mục này, sử dụng các khái niệm và kết quả ở phần sau của Chương~I và Chương~II, sẽ không được dùng trong phần còn lại của Chuyên luận này và chỉ dành cho những độc giả quen thuộc với hình học đại số cổ điển}.

\begin{env}[3.7.1]
\label{I.3.7.1}
Cho $A$ là một vành, $X$ là một $A$-tiền lược đồ và $\mathfrak{J}$ là một iđêan của $A$; khi đó $X_0=X\otimes_A(A/\mathfrak{J})$ là một $(A/\mathfrak{J})$-tiền lược đồ, và đôi khi người ta nói rằng nó thu được từ $X$ bằng phép \emph{rút gọn} modulo~$\mathfrak{J}$.
\end{env}

\begin{env}[3.7.2]
\label{I.3.7.2}
Thuật ngữ này chủ yếu được dùng khi $A$ là một \emph{vành địa phương} và $\mathfrak{J}$ là iđêan cực đại của nó, do đó $X_0$ là một tiền lược đồ trên trường thặng dư $k=A/\mathfrak{J}$ của $A$.

Khi, ngoài ra, $A$ là một miền nguyên có trường phân thức $K$, ta cũng có thể xét $K$-tiền lược đồ $X'=X\otimes_A K$.
Theo một sự lạm dụng ngôn ngữ mà ta sẽ không dùng, cho đến nay người ta thường nói rằng $X_0$ được \emph{suy ra từ $X'$} bằng phép rút gọn modulo~$\mathfrak{J}$.
Trong những trường hợp sử dụng cách nói này, $A$ là một vành địa phương chiều~$1$ (thường nhất là một vành định giá rời rạc), và người ta ngầm hiểu (một cách ít nhiều tường minh) rằng $K$-tiền lược đồ đã cho $X'$ là một tiền lược đồ con đóng của một $K$-tiền lược đồ $P'$ (thực ra là một không gian xạ ảnh kiểu $\bb{P}_K^r$, cf.~\sref[II]{II.4.1.1}), bản thân nó có dạng $P'=P\otimes_A K$, trong đó $P$ là một $A$-tiền lược đồ đã cho (thực ra là $A$-lược đồ $\bb{P}_A^r$, với các ký hiệu của \sref[II]{II.4.1.1}).
Trong ngôn ngữ của chúng ta, định nghĩa $X_0$ từ $X'$ được phát biểu như sau:

Xét lược đồ affine $Y=\Spec(A)$, gồm hai điểm, điểm đóng duy nhất
\oldpage[I]{119}
$y=\mathfrak{J}$ và điểm tổng quát $(0)$; do đó tập $U$ chỉ gồm điểm tổng quát là một tập mở $U=\Spec(K)$ trong $Y$.
Nếu $X$ là một $A$-tiền lược đồ (hay nói cách khác, một $Y$-tiền lược đồ), thì $X\otimes_A K=X'$ không gì khác hơn tiền lược đồ được cảm sinh bởi $X$ trên $\psi^{-1}(U)$, trong đó $\psi$ ký hiệu cấu xạ cấu trúc $X\to Y$.
Đặc biệt, nếu $\vphi$ là cấu xạ cấu trúc $P\to Y$, thì một tiền lược đồ con đóng $X'$ của $P'=\vphi^{-1}(U)$ là một tiền lược đồ con (đóng địa phương) của $P$.
Nếu $P$ là Noether (chẳng hạn nếu $A$ là Noether và $P$ thuộc kiểu hữu hạn trên $A$), thì tồn tại một tiền lược đồ con đóng nhỏ nhất $X=\overline{X'}$ của $P$ chứa $X'$ \sref{I.9.5.10}, và $X'$ là tiền lược đồ được cảm sinh bởi $X$ trên tập mở $\vphi^{-1}(U)\cap X$, do đó đẳng cấu với $X\otimes_A K$ \sref{I.9.5.10}.
\emph{Phép nhúng của $X'$ vào $P'=P\otimes_A K$ do đó cho phép xét một cách chính tắc $X'$ dưới dạng $X'=X\otimes_A K$, trong đó $X$ là một $A$-tiền lược đồ.}
Khi ấy ta có thể xét tiền lược đồ rút gọn modulo~$\mathfrak{J}$, $X_0=X\otimes_A k$, vốn không gì khác hơn thớ $\psi^{-1}(y)$ của điểm đóng $y$.
Cho đến nay, do thiếu một thuật ngữ thích hợp, người ta đã tránh đưa $A$-tiền lược đồ $X$ vào một cách tường minh.
Tuy nhiên, cần lưu ý rằng mọi khẳng định thường được phát biểu về tiền lược đồ “rút gọn modulo~$\mathfrak{J}$” $X_0$ phải được xem là hệ quả của những khẳng định đầy đủ hơn liên quan đến chính $X$, và chỉ có thể được phát biểu và hiểu một cách thỏa đáng khi được diễn giải như vậy.
Hơn nữa, dường như mọi giả thiết đặt lên $X_0$ luôn quy về các giả thiết đặt lên chính $X$ (không phụ thuộc vào dữ liệu trước đó của một phép nhúng $X'$ vào $\bb{P}_K^r$), điều này cho phép phát biểu những mệnh đề nội tại hơn.
\end{env}

\begin{env}[3.7.3]
\label{I.3.7.3}
Cuối cùng, ta hãy nêu một sự kiện rất đặc biệt, chắc chắn đã góp phần làm chậm
việc làm sáng tỏ về mặt khái niệm tình huống đang xét ở đây: nếu $A$ là một
vành định giá rời rạc và nếu $X$ là \emph{riêng} trên $A$ (điều này quả thực
đúng nếu $X$ là một tiền lược đồ con đóng của một $\bb{P}_A^r$ nào đó,
cf.~\sref[II]{II.5.5.4}), thì các điểm của $X$ nhận giá trị trong $A$ và các
điểm của $X'$ nhận giá trị trong $K$ tương ứng song ánh
\sref[II]{II.7.3.8}.
Vì lý do đó, người ta thường tưởng rằng đã chứng minh các kết quả liên quan
đến $X'$, trong khi thực ra lại chứng minh các mệnh đề liên quan đến $X$; các
mệnh đề ấy vẫn đúng (dưới dạng đó) khi không còn giả sử rằng vành địa phương
cơ sở có chiều~$1$.
\end{env}

\section{Các tiền lược đồ con và cấu xạ nhúng}
\label{section:I.4}

\subsection{Các tiền lược đồ con}
\label{subsection:I.4.1}

\begin{env}[4.1.1]
\label{I.4.1.1}
Vì khái niệm bó quasi-coherent \sref[0]{0.5.1.3} có tính địa phương, một $\sh{O}_X$-môđun quasi-coherent $\sh{F}$ trên một tiền lược đồ $X$ có thể được định nghĩa bởi điều kiện sau: với mọi tập mở affine $V$ của $X$, $\sh{F}|V$ đẳng cấu với bó liên kết với một $\Gamma(V,\sh{O}_X)$-môđun \sref{I.1.4.1}.
Rõ ràng, trên một tiền lược đồ $X$, bó cấu trúc $\sh{O}_X$ là quasi-coherent, và các hạt nhân, đối hạt nhân cùng ảnh của các đồng cấu giữa các $\sh{O}_X$-môđun quasi-coherent, cũng như các giới hạn quy nạp và các tổng trực tiếp của các $\sh{O}_X$-môđun quasi-coherent, đều quasi-coherent (Định lý~\sref{I.1.3.7} và Hệ quả~\sref{I.1.3.9}).
\end{env}

\begin{proposition}[4.1.2]
\label{I.4.1.2}
Cho $X$ là một tiền lược đồ và $\sh{I}$ là một bó iđêan quasi-coherent của $\sh{O}_X$.
Khi đó, giá $Y$ của bó $\sh{O}_X/\sh{I}$ là đóng, và nếu $\sh{O}_Y$ là ký hiệu cho hạn chế của $\sh{O}_X/\sh{I}$ lên $Y$, thì $(Y,\sh{O}_Y)$ là một tiền lược đồ.
\end{proposition}

\begin{proof}
\oldpage[I]{120}
Hiển nhiên chỉ cần \sref{I.2.1.3} xét trường hợp $X$ là một lược đồ affine và chứng minh rằng, trong trường hợp này, $Y$ đóng trong $X$ và là một \emph{lược đồ affine}.
Thật vậy, nếu $X=\Spec(A)$, thì $\sh{O}_X=\widetilde{A}$ và $\sh{I}=\widetilde{\mathfrak{I}}$, trong đó $\mathfrak{I}$ là một iđêan của $A$ \sref{I.1.4.1};
khi đó $Y$ bằng phần đóng $V(\mathfrak{I})$ của $X$ và được đồng nhất với phổ nguyên tố của vành $B=A/\mathfrak{I}$ \sref{I.1.1.11};
hơn nữa, nếu $\vphi$ là đồng cấu chính tắc $A\to B=A/\mathfrak{I}$, thì ảnh trực tiếp ${}^a\vphi_*(\widetilde{B})$ được đồng nhất chính tắc với bó $\widetilde{A}/\widetilde{\mathfrak{I}}=\sh{O}_X/\sh{I}$ (Mệnh đề~\sref{I.1.6.3} và Hệ quả~\sref{I.1.3.9}), điều này kết thúc chứng minh.
\end{proof}

Ta nói $(Y,\sh{O}_Y)$ là \emph{tiền lược đồ con} của $(X,\sh{O}_X)$ \emph{được xác định bởi bó iđêan $\sh{I}$};
đây là một trường hợp riêng của khái niệm tổng quát về \emph{tiền lược đồ con}:

\begin{definition}[4.1.3]
\label{I.4.1.3}
Ta nói một không gian vành $(Y,\sh{O}_Y)$ là một tiền lược đồ con của một tiền lược đồ $(X,\sh{O}_X)$ nếu:
\begin{enumerate}
  \item[1.º] $Y$ là một không gian con đóng địa phương của $X$;
  \item[2.º] nếu $U$ ký hiệu tập mở lớn nhất của $X$ chứa $Y$ và sao cho $Y$ đóng trong $U$ (\emph{tương đương}, phần bù trong $X$ của biên của $Y$ trong $\overline{Y}$), thì $(Y,\sh{O}_Y)$ là một tiền lược đồ con của $(U,\sh{O}_X|U)$ được xác định bởi một bó iđêan quasi-coherent của $\sh{O}_X|U$.
\end{enumerate}
Ta nói tiền lược đồ con $(Y,\sh{O}_Y)$ của $(X,\sh{O}_X)$ là \emph{đóng} nếu $Y$ đóng trong $X$ (khi đó $U=X$).
\end{definition}

Từ định nghĩa này và Mệnh đề~\sref{I.4.1.2} suy ra ngay rằng các tiền lược đồ con đóng của $X$ \emph{tương ứng song ánh chính tắc} với các bó iđêan quasi-coherent $\sh{J}$ của $\sh{O}_X$, vì nếu hai bó như thế $\sh{J}$ và $\sh{J}'$ có cùng giá (đóng) $Y$, và nếu các hạn chế của $\sh{O}_X/\sh{J}$ và $\sh{O}_X/\sh{J}'$ lên $Y$ trùng nhau, thì $\sh{J}'=\sh{J}$.

\begin{env}[4.1.4]
\label{I.4.1.4}
Cho $(Y,\sh{O}_Y)$ là một tiền lược đồ con của $X$, $U$ là tập mở lớn nhất của $X$ sao cho $Y$ là một không gian con đóng của $U$, còn $V$ là một tập mở của $X$ nằm trong $U$;
khi đó $V\cap Y$ đóng trong $V$.
Hơn nữa, nếu $Y$ được xác định bởi bó iđêan quasi-coherent $\sh{J}$ của $\sh{O}_X|U$, thì $\sh{J}|V$ là một bó iđêan quasi-coherent của $\sh{O}_X|V$, và hiển nhiên tiền lược đồ do $Y$ cảm sinh trên $Y\cap V$ là tiền lược đồ con đóng của $V$ được xác định bởi bó iđêan $\sh{J}|V$.
Ngược lại:
\end{env}

\begin{proposition}[4.1.5]
\label{I.4.1.5}
Cho $(Y,\sh{O}_Y)$ là một không gian vành sao cho $Y$ là một không gian con của $X$ và tồn tại một phủ $(V_\alpha)$ của $Y$ bởi các tập mở của $X$ sao cho, với mọi $\alpha$, $Y\cap V_\alpha$ đóng trong $V_\alpha$, và không gian vành $(Y\cap V_\alpha,\sh{O}_Y|(Y\cap V_\alpha))$ là một tiền lược đồ con đóng của tiền lược đồ cảm sinh trên $V_\alpha$ bởi $X$.
Khi đó $(Y,\sh{O}_Y)$ là một tiền lược đồ con của $X$.
\end{proposition}

\begin{proof}
Giả thiết kéo theo rằng $Y$ đóng địa phương trong $X$ và tập mở lớn nhất $U$ mà trong đó $Y$ đóng chứa mọi $V_\alpha$;
do đó ta có thể quy về trường hợp $U=X$ và $Y$ đóng trong $X$.
Khi đó ta định nghĩa một bó iđêan quasi-coherent $\sh{J}$ của $\sh{O}_X$ bằng cách lấy $\sh{J}|V_\alpha$ là bó iđêan của $\sh{O}_X|V_\alpha$ xác định tiền lược đồ con đóng $(Y\cap V_\alpha,\sh{O}_Y|(Y\cap V_\alpha))$, và, với mọi tập mở $W$ của $X$ không giao với $Y$, đặt $\sh{J}|W=\sh{O}_X|W$.
Từ Định nghĩa~\sref{I.4.1.3} và \sref{I.4.1.4} kiểm tra ngay được rằng tồn tại duy nhất một bó iđêan $\sh{J}$ thỏa mãn các điều kiện này và nó xác định tiền lược đồ con đóng $(Y,\sh{O}_Y)$.
\end{proof}

Đặc biệt, tiền lược đồ \emph{cảm sinh} (bởi $X$) trên một \emph{tập mở} của $X$ là một \emph{tiền lược đồ con} của $X$.

\begin{proposition}[4.1.6]
\label{I.4.1.6}
Một tiền lược đồ con (tương ứng, một tiền lược đồ con đóng) của một tiền lược đồ con
\oldpage[I]{121}
(tương ứng, một tiền lược đồ con đóng) của $X$ được đồng nhất chính tắc với một tiền lược đồ con (tương ứng, một tiền lược đồ con đóng) của $X$.
\end{proposition}

\begin{proof}
Vì một tập con đóng địa phương của một không gian con đóng địa phương của $X$ là một không gian con đóng địa phương của $X$, theo \sref{I.4.1.5} rõ ràng bài toán có tính địa phương, nên ta có thể giả sử $X$ affine;
khi đó mệnh đề suy ra từ phép đồng nhất chính tắc $A/\mathfrak{J}'$ với $(A/\mathfrak{J})/(\mathfrak{J}'/\mathfrak{J})$, trong đó $A$ là một vành, còn $\mathfrak{J}$ và $\mathfrak{J}'$ là các iđêan của $A$ sao cho $\mathfrak{J}\subset\mathfrak{J}'$.
\end{proof}

Từ đây về sau, ta luôn dùng phép đồng nhất trên.

\begin{env}[4.1.7]
\label{I.4.1.7}
Cho $Y$ là một tiền lược đồ con của một tiền lược đồ $X$, và ký hiệu bởi $\psi$ phép nhúng chính tắc $Y\to X$ của các \emph{không gian nền};
ta biết rằng ảnh ngược $\psi^*(\sh{O}_X)$ là hạn chế $\sh{O}_X|Y$ \sref[0]{0.3.7.1}.
Nếu, với mọi $y\in Y$, ta ký hiệu bởi $\omega_y$ đồng cấu chính tắc $(\sh{O}_X)_y\to(\sh{O}_Y)_y$, thì các đồng cấu này là những đồng cấu cảm sinh trên các thớ bởi một đồng cấu \emph{toàn ánh} $\omega$ của các bó vành $\sh{O}_X|Y\to\sh{O}_Y$: quả vậy, chỉ cần kiểm tra điều này một cách địa phương trên $Y$, nghĩa là ta có thể giả sử $X$ affine và tiền lược đồ con $Y$ đóng; nếu, trong trường hợp này, $\sh{I}$ là bó iđêan của $\sh{O}_X$ xác định $Y$, thì các $\omega_y$ chính là những đồng cấu cảm sinh trên các thớ bởi $\sh{O}_X|Y\to(\sh{O}_X/\sh{I})|Y$.
Như vậy ta đã xác định một \emph{đơn cấu của các không gian vành} \sref[0]{0.4.1.1} $j=(\psi,\omega^\flat)$, hiển nhiên là một cấu xạ $Y\to X$ của các tiền lược đồ \sref{I.2.2.1}, và ta gọi nó là \emph{cấu xạ nhúng chính tắc}.

Nếu $f:X\to Z$ là một cấu xạ, ta cũng gọi cấu xạ hợp thành $Y\xrightarrow{j}X\xrightarrow{f}Z$ là \emph{hạn chế} của $f$ trên tiền lược đồ con $Y$.
\end{env}

\begin{env}[4.1.8]
\label{I.4.1.8}
Theo các định nghĩa tổng quát (T, I, 1.1), ta nói một cấu xạ các tiền lược đồ $f:Z\to X$ bị \emph{chi phối}\footnote{\emph{[Dịch.] Ở đây chúng tôi dịch từ tiếng Pháp «majoré» là «bị chi phối». Đây là quan hệ tiền thứ tự được xác định bởi điều kiện $f$ phân tích qua $j$, chứ không phải một cận số.}} bởi cấu xạ nhúng $j:Y\to X$ của một tiền lược đồ con $Y$ của $X$ nếu $f$ phân tích thành $Z\xrightarrow{g}Y\xrightarrow{j}X$, trong đó $g$ là một cấu xạ các tiền lược đồ;
hơn nữa, $g$ nhất thiết \emph{duy nhất}, vì $j$ là một đơn cấu.
\end{env}

\begin{proposition}[4.1.9]
\label{I.4.1.9}
Điều kiện cần và đủ để một cấu xạ $f:Z\to X$ bị một cấu xạ nhúng $j:Y\to X$ chi phối là $f(Z)\subset Y$ và, với mọi $z\in Z$, nếu đặt $y=f(z)$, đồng cấu $(\sh{O}_X)_y\to\sh{O}_z$ tương ứng với $f$ phân tích thành $(\sh{O}_X)_y\to(\sh{O}_Y)_y\to\sh{O}_z$ (\emph{hay, tương đương,} hạt nhân của $(\sh{O}_X)_y\to\sh{O}_z$ chứa hạt nhân của $(\sh{O}_X)_y\to(\sh{O}_Y)_y$).
\end{proposition}

\begin{proof}
Các điều kiện hiển nhiên là cần.
Để thấy chúng là đủ, ta có thể quy về trường hợp $Y$ là một tiền lược đồ con \emph{đóng} của $X$, bằng cách thay $X$, nếu cần, bởi một tập mở $U$ sao cho $Y$ đóng trong $U$ \sref{I.4.1.3};
khi đó $Y$ được xác định bởi một bó iđêan quasi-coherent $\sh{I}$ của $\sh{O}_X$.
Đặt $f=(\psi,\theta)$, và cho $\sh{J}$ là bó iđêan của $\psi^*(\sh{O}_X)$, bằng hạt nhân của $\theta^\sharp:\psi^*(\sh{O}_X)\to\sh{O}_Z$;
xét đến các tính chất của hàm tử $\psi^*$ \sref[0]{0.3.7.2}, các giả thiết kéo theo rằng, với mọi $z\in Z$, ta có $(\psi^*(\sh{I}))_z\subset\sh{J}_z$ và do đó $\psi^*(\sh{I})\subset\sh{J}$.
Vì vậy, $\theta^\sharp$ phân tích thành
\[
  \psi^*(\sh{O}_X)\to\psi^*(\sh{O}_X)/\psi^*(\sh{I})=\psi^*(\sh{O}_X/\sh{I})\xrightarrow{\omega}\sh{O}_Z,
\]
trong đó mũi tên thứ nhất là đồng cấu chính tắc.
Cho $\psi'$ là ánh xạ liên tục $Z\to Y$ trùng với $\psi$;
rõ ràng ta có ${\psi'}^*(\sh{O}_Y)=\psi^*(\sh{O}_X/\sh{I})$;
mặt khác, $\omega$ hiển nhiên là một đồng cấu địa phương, nên $g=(\psi',\omega^\flat)$ là một cấu xạ
\oldpage[I]{122}
các tiền lược đồ $Z\to Y$ \sref{I.2.2.1} thỏa mãn $f=j\circ g$ theo những điều trên, điều này chứng minh mệnh đề.
\end{proof}

\begin{corollary}[4.1.10]
\label{I.4.1.10}
Điều kiện cần và đủ để một cấu xạ nhúng $Z\to X$ bị cấu xạ nhúng $Y\to X$ chi phối là $Z$ là một tiền lược đồ con của $Y$.
\end{corollary}

Khi đó ta viết $Z\leq Y$, và quan hệ này hiển nhiên là một \emph{quan hệ thứ tự} trên tập các tiền lược đồ con của $X$.

\subsection{Các cấu xạ nhúng}
\label{subsection:I.4.2}

\begin{definition}[4.2.1]
\label{I.4.2.1}
Ta nói một cấu xạ $f:Y\to X$ là một phép nhúng (tương ứng, một phép nhúng đóng, một phép nhúng mở) nếu nó phân tích thành $Y\xrightarrow{g}Z\xrightarrow{j}X$, trong đó $g$ là một đẳng cấu, $Z$ là một tiền lược đồ con của $X$ (tương ứng, một tiền lược đồ con đóng, một tiền lược đồ con cảm sinh trên một tập mở), và $j$ là cấu xạ nhúng.
\end{definition}

Khi đó tiền lược đồ con $Z$ và đẳng cấu $g$ được xác định một cách \emph{duy nhất}, vì nếu $Z'$ là một tiền lược đồ con thứ hai của $X$, $j'$ là cấu xạ nhúng $Z'\to X$, và $g'$ là một đẳng cấu $Y\to Z'$ sao cho $j\circ g=j'\circ g'$, thì suy ra $j'=j\circ g\circ{g'}^{-1}$, do đó $Z'\leq Z$ \sref{I.4.1.10}; tương tự, ta có $Z\leq Z'$, nên $Z'=Z$, và vì $j$ là một đơn cấu của các tiền lược đồ nên $g'=g$.

Ta nói $f=j\circ g$ là \emph{phân tích chính tắc} của phép nhúng $f$, và tiền lược đồ con $Z$ cùng đẳng cấu $g$ được gọi là \emph{liên kết với $f$}.

Rõ ràng một phép nhúng là một \emph{đơn cấu} của các tiền lược đồ \sref{I.4.1.7} và, \emph{a fortiori}, là một \emph{cấu xạ radicial} \sref{I.3.5.4}.

\begin{proposition}[4.2.2]
\label{I.4.2.2}
\medskip\noindent
\begin{enumerate}
  \item[{\rm(a)}] Điều kiện cần và đủ để một cấu xạ $f=(\psi,\theta):Y\to X$ là một phép nhúng mở là $\psi$ là một phép đồng phôi từ $Y$ lên một tập con mở của $X$, và, với mọi $y\in Y$, đồng cấu $\theta_y^\sharp:\sh{O}_{\psi(y)}\to\sh{O}_y$ là song ánh.
  \item[{\rm(b)}] Điều kiện cần và đủ để một cấu xạ $f=(\psi,\theta):Y\to X$ là một phép nhúng (tương ứng, một phép nhúng đóng) là $\psi$ là một phép đồng phôi từ $Y$ lên một tập con đóng địa phương (tương ứng, đóng) của $X$, và, với mọi $y\in Y$, đồng cấu $\theta_y^\sharp:\sh{O}_{\psi(y)}\to\sh{O}_y$ là toàn ánh.
\end{enumerate}
\end{proposition}

\begin{proof}
\medskip\noindent
\begin{enumerate}
  \item[(a)] Các điều kiện hiển nhiên là cần.
    Ngược lại, nếu chúng được thỏa mãn, rõ ràng $\theta^\sharp$ là một đẳng cấu từ $\psi^*(\sh{O}_X)$ lên $\sh{O}_Y$\footnote{\emph{[Dịch.] Nguồn tiếng Pháp đảo ngược nguồn và đích tại đây. Theo định nghĩa của một cấu xạ tiền lược đồ và các đồng cấu trên các thớ được nêu trong Mệnh đề~\hyperref[I.4.2.2]{4.2.2}(a), $\theta^\sharp$ có nguồn $\psi^*(\sh{O}_X)$ và đích $\sh{O}_Y$.}}, và $\psi^*(\sh{O}_X)$ là bó thu được bằng cách ``chuyển cấu trúc'' qua $\psi^{-1}$ từ $\sh{O}_X|\psi(Y)$;
    từ đó suy ra kết luận.

  \item[(b)] Các điều kiện hiển nhiên là cần; ta chứng minh chúng là đủ.
    Trước hết xét trường hợp riêng trong đó giả sử $X$ là một lược đồ affine và $Z=\psi(Y)$ \emph{đóng} trong $X$.
    Khi đó, theo \sref[0]{0.3.4.6}, $\psi_*(\sh{O}_Y)$ có giá $Z$, và hạn chế của nó lên $Z$, ký hiệu là $\sh{O}_Z'$, làm cho không gian vành $(Z,\sh{O}_Z')$ thu được từ $(Y,\sh{O}_Y)$ bằng cách chuyển cấu trúc qua phép đồng phôi $\psi$, được xem như một ánh xạ từ $Y$ lên $Z$.
    Bây giờ ta chứng minh rằng $f_*(\sh{O}_Y)=\psi_*(\sh{O}_Y)$ là một $\sh{O}_X$-môđun \emph{quasi-coherent}.
    Thật vậy, với mọi $x\not\in Z$, $\psi_*(\sh{O}_Y)$, khi hạn chế lên một lân cận thích hợp của $x$, bằng không.
    Ngược lại, nếu $x\in Z$, thì $x=\psi(y)$ với một $y\in Y$ duy nhất;
    cho $V$ là một lân cận mở affine của $y$ trong $Y$;
    khi đó $\psi(V)$ mở trong $Z$, do đó bằng giao của $Z$ với một tập mở $U$ của $X$, và hạn chế lên $U$ của $\psi_*(\sh{O}_Y)$ được đồng nhất với hạn chế lên $U$ của ảnh trực tiếp
\oldpage[I]{123}
    $(\psi_V)_*(\sh{O}_Y|V)$, trong đó $\psi_V$ là hạn chế của $\psi$ lên $V$.
    Hạn chế của cấu xạ $(\psi,\theta)$ lên $(V,\sh{O}_Y|V)$ là một cấu xạ từ tiền lược đồ này đến $(X,\sh{O}_X)$ và do đó có dạng $({}^a\vphi,\widetilde{\vphi})$, trong đó $\vphi$ là đồng cấu từ vành $A=\Gamma(X,\sh{O}_X)$ đến vành $\Gamma(V,\sh{O}_Y)$ \sref{I.1.7.3};
    suy ra $(\psi_V)_*(\sh{O}_Y|V)$ là một $\sh{O}_X$-môđun quasi-coherent \sref{I.1.6.3}, điều này chứng minh khẳng định của ta nhờ tính địa phương của các bó quasi-coherent.

    Hơn nữa, giả thiết rằng $\psi$ là một phép đồng phôi kéo theo, theo \sref[0]{0.3.4.5}, rằng với mọi $y\in Y$, $\psi_y$ là một đẳng cấu $(\psi_*(\sh{O}_Y))_{\psi(y)}\to\sh{O}_y$;
    vì biểu đồ
    \[
      \xymatrix{
        \sh{O}_{\psi(y)}\ar[r]^{\theta_{\psi(y)}}\ar[d]_{\psi_y\circ\rho_{\psi(y)}} &
        (\psi_*(\sh{O}_Y))_{\psi(y)}\ar[d]^{\psi_y}\\
        (\psi^*(\sh{O}_X))_y\ar[r]^{\theta_y^\sharp} &
        \sh{O}_y
      }
    \]
    giao hoán và các mũi tên dọc là các đẳng cấu \sref[0]{0.3.7.2}, giả thiết rằng $\theta_y^\sharp$ toàn ánh kéo theo $\theta_{\psi(y)}$ cũng toàn ánh.
    Vì giá của $\psi_*(\sh{O}_Y)$ là $Z=\psi(Y)$, $\theta$ là một đồng cấu \emph{toàn ánh} từ $\sh{O}_X=\widetilde{A}$ lên $\sh{O}_X$-môđun quasi-coherent $f_*(\sh{O}_Y)$.
    Do đó, tồn tại duy nhất một đẳng cấu $\omega$ từ một bó thương $\widetilde{A}/\widetilde{\mathfrak{J}}$ (trong đó $\mathfrak{J}$ là một iđêan của $A$) lên $f_*(\sh{O}_Y)$ sao cho, khi hợp thành với đồng cấu chính tắc $\widetilde{A}\to\widetilde{A}/\widetilde{\mathfrak{J}}$, nó cho $\theta$ \sref{I.1.3.8};
    nếu $\sh{O}_Z$ ký hiệu hạn chế của $\widetilde{A}/\widetilde{\mathfrak{J}}$ lên $Z$, thì $(Z,\sh{O}_Z)$ là một tiền lược đồ con của $(X,\sh{O}_X)$, và $f$ phân tích qua phép nhúng chính tắc của tiền lược đồ con này vào $X$ và qua đẳng cấu $(\psi_0,\omega_0)$, trong đó $\psi_0$ là $\psi$ được xem như một ánh xạ từ $Y$ lên $Z$, còn $\omega_0$ là hạn chế của $\omega$ lên $\sh{O}_Z$.

    Bây giờ ta chuyển sang trường hợp tổng quát.
    Cho $U$ là một tập mở affine của $X$ sao cho $U\cap\psi(Y)$ đóng trong $U$ và không rỗng.
    Bằng cách hạn chế $f$ lên tiền lược đồ do $Y$ cảm sinh trên tập mở $\psi^{-1}(U)$ và xem nó như một cấu xạ từ tiền lược đồ này đến tiền lược đồ do $X$ cảm sinh trên $U$, ta quy về trường hợp thứ nhất;
    do đó, hạn chế của $f$ lên $\psi^{-1}(U)$ là một phép nhúng đóng $\psi^{-1}(U)\to U$, phân tích chính tắc thành $j_U\circ g_U$, trong đó $g_U$ là một đẳng cấu từ tiền lược đồ $\psi^{-1}(U)$ lên một tiền lược đồ con $Z_U$ của $U$, và $j_U$ là phép nhúng chính tắc $Z_U\to U$.
    Cho $V$ là một tập mở affine khác của $X$ sao cho $V\subset U$;
    vì hạn chế $Z_V'$ của $Z_U$ lên $V$ là một tiền lược đồ con của tiền lược đồ $V$, nên hạn chế của $f$ lên $\psi^{-1}(V)$ phân tích thành $j_V'\circ g_V'$, trong đó $j_V'$ là phép nhúng chính tắc $Z_V'\to V$ và $g_V'$ là một đẳng cấu từ $\psi^{-1}(V)$ lên $Z_V'$.
    Theo tính duy nhất của phân tích chính tắc của một phép nhúng \sref{I.4.2.1}, nhất thiết ta có $Z_V'=Z_V$ và $g_V'=g_V$.
    Từ \sref{I.4.1.5} suy ra tồn tại một tiền lược đồ con $Z$ của $X$ có không gian nền là $\psi(Y)$ và có hạn chế lên mỗi $U\cap\psi(Y)$ bằng $Z_U$;
    khi đó các $g_U$ là các hạn chế lên $\psi^{-1}(U)$ của một đẳng cấu $g:Y\to Z$ sao cho $f=j\circ g$, trong đó $j$ là phép nhúng chính tắc $Z\to X$.
\end{enumerate}
\end{proof}

\begin{corollary}[4.2.3]
\label{I.4.2.3}
Cho $X$ là một lược đồ affine.
Để một cấu xạ $f=(\psi,\theta):Y\to X$ là một phép nhúng đóng, điều kiện cần và đủ là $Y$ là một lược đồ affine và đồng cấu $\Gamma(\theta)$\footnote{\emph{[Dịch.] Nguồn tiếng Pháp in $\Gamma(\psi)$. Vì $\psi$ là ánh xạ liên tục còn $\theta$ là đồng cấu bó, đồng cấu cảm sinh trên các nhát cắt toàn cục là $\Gamma(\theta)$.}}$:\Gamma(\sh{O}_X)\to\Gamma(\sh{O}_Y)$ là toàn ánh.
\end{corollary}

\begin{corollary}[4.2.4]
\label{I.4.2.4}
\medskip\noindent
\begin{enumerate}
  \item[{\rm(a)}] Cho $f$ là một cấu xạ $Y\to X$ và $(V_\lambda)$ là một phủ của $f(Y)$ bởi các tập mở của $X$.
    Điều kiện cần và đủ để $f$ là một phép nhúng (tương ứng, một phép nhúng mở) là
\oldpage[I]{124}
    hạn chế của nó trên mỗi tiền lược đồ cảm sinh $f^{-1}(V_\lambda)$ là một phép nhúng (tương ứng, một phép nhúng mở) vào $V_\lambda$.
  \item[{\rm(b)}] Cho $f$ là một cấu xạ $Y\to X$ và $(V_\lambda)$ là một phủ mở của $X$.
    Điều kiện cần và đủ để $f$ là một phép nhúng đóng là hạn chế của nó trên mỗi tiền lược đồ cảm sinh $f^{-1}(V_\lambda)$ là một phép nhúng đóng vào $V_\lambda$.
\end{enumerate}
\end{corollary}

\begin{proof}
Cho $f=(\psi,\theta)$;
trong trường hợp (a), $\theta_y^\sharp$ là toàn ánh (tương ứng, song ánh) với mọi $y\in Y$, và trong trường hợp (b), $\theta_y^\sharp$ là toàn ánh với mọi $y\in Y$;
do đó chỉ cần kiểm tra rằng, trong trường hợp (a), $\psi$ là một phép đồng phôi từ $Y$ lên một tập con đóng địa phương (tương ứng, mở) của $X$, và, trong trường hợp (b), $\psi$ là một phép đồng phôi từ $Y$ lên một tập con đóng của $X$.
Theo giả thiết, $\psi$ hiển nhiên là đơn ánh và biến mọi lân cận của $y$ trong $Y$ thành một lân cận của $\psi(y)$ trong $\psi(Y)$, với mọi $y\in Y$;
trong trường hợp (a), $\psi(Y)\cap V_\lambda$ đóng địa phương (tương ứng, mở) trong $V_\lambda$, nên $\psi(Y)$ đóng địa phương (tương ứng, mở) trong hợp của các $V_\lambda$, và \emph{a fortiori} trong $X$;
trong trường hợp (b), $\psi(Y)\cap V_\lambda$ đóng trong $V_\lambda$, nên $\psi(Y)$ đóng trong $X$, vì $X=\bigcup_\lambda V_\lambda$.
\end{proof}

\begin{proposition}[4.2.5]
\label{I.4.2.5}
Hợp thành của hai phép nhúng bất kỳ (tương ứng, của hai phép nhúng mở, của hai phép nhúng đóng) là một phép nhúng (tương ứng, một phép nhúng mở, một phép nhúng đóng).
\end{proposition}

\begin{proof}
Điều này suy ra hiển nhiên từ \sref{I.4.1.6}.
\end{proof}

\subsection{Tích của các phép nhúng}
\label{subsection:I.4.3}
\label{I.4.3}

\begin{proposition}[4.3.1]
\label{I.4.3.1}
Cho $\alpha:X'\to X$, $\beta:Y'\to Y$ là hai $S$-cấu xạ; nếu $\alpha$ và $\beta$ là các phép nhúng (tương ứng, các phép nhúng mở, các phép nhúng đóng), thì $\alpha\times_S\beta$ là một phép nhúng (tương ứng, một phép nhúng mở, một phép nhúng đóng).
Hơn nữa, nếu $\alpha$ (tương ứng, $\beta$) đồng nhất $X'$ (tương ứng, $Y'$) với một tiền lược đồ con $X''$ (tương ứng, $Y''$) của $X$ (tương ứng, $Y$), thì $\alpha\times_S\beta$ đồng nhất không gian nền của $X'\times_S Y'$ với không gian con $p^{-1}(X'')\cap q^{-1}(Y'')$ của không gian nền của $X\times_S Y$, trong đó $p$ và $q$ lần lượt là các phép chiếu của $X\times_S Y$ lên $X$ và $Y$.
\end{proposition}

\begin{proof}
Theo Định nghĩa~\sref{I.4.2.1}, ta có thể chỉ xét trường hợp $X'$ và $Y'$ là các tiền lược đồ con, còn $\alpha$ và $\beta$ là các cấu xạ nhúng.
Mệnh đề đã được chứng minh đối với các tiền lược đồ con cảm sinh trên các tập mở \sref{I.3.2.7};
vì mọi tiền lược đồ con đều là một tiền lược đồ con đóng của một tiền lược đồ cảm sinh trên một tập mở \sref{I.4.1.3}, ta quy về trường hợp $X'$ và $Y'$ là các tiền lược đồ con \emph{đóng}.

Trước hết, ta chứng minh rằng có thể giả sử $S$ là \emph{affine}.
Thật vậy, cho $(S_\lambda)$ là một phủ của $S$ bởi các tập mở affine;
nếu $\vphi$ và $\psi$ là các cấu xạ cấu trúc của $X$ và $Y$, đặt $X_\lambda=\vphi^{-1}(S_\lambda)$ và $Y_\lambda=\psi^{-1}(S_\lambda)$.
Hạn chế $X_\lambda'$ (tương ứng, $Y_\lambda'$) của $X'$ (tương ứng, $Y'$) lên $X_\lambda\cap X'$ (tương ứng, $Y_\lambda\cap Y'$) là một tiền lược đồ con đóng của $X_\lambda$ (tương ứng, $Y_\lambda$); các tiền lược đồ $X_\lambda$, $Y_\lambda$, $X_\lambda'$ và $Y_\lambda'$ có thể được xem là các $S_\lambda$-tiền lược đồ, và các tích $X_\lambda\times_S Y_\lambda$ và $X_\lambda\times_{S_\lambda}Y_\lambda$ (tương ứng, $X_\lambda'\times_S Y_\lambda'$ và $X_\lambda'\times_{S_\lambda}Y_\lambda'$) là đồng nhất \sref{I.3.2.5}.
Nếu mệnh đề đúng khi $S$ là affine, thì hạn chế của $\alpha\times_S\beta$ lên mỗi $X_\lambda'\times_S Y_\lambda'$ do đó là một phép nhúng \sref{I.3.2.7}.
Vì tích $X_\lambda'\times_S Y_\mu'$ (tương ứng, $X_\lambda\times_S Y_\mu$) được đồng nhất với $(X_\lambda'\cap X_\mu')\times_S(Y_\lambda'\cap Y_\mu')$ (tương ứng, $(X_\lambda\cap X_\mu)\times_S(Y_\lambda\cap Y_\mu)$) \sref{I.3.2.6.4}, nên hạn chế của $\alpha\times_S\beta$
\oldpage[I]{125}
lên mỗi $X_\lambda'\times_S Y_\mu'$ vẫn là một phép nhúng;
do đó $\alpha\times_S\beta$ cũng là một phép nhúng theo \sref{I.4.2.4}.

Thứ hai, ta chứng minh rằng cũng có thể giả sử $X$ và $Y$ là \emph{affine}.
Thật vậy, cho $(U_i)$ (tương ứng, $(V_j)$) là một phủ của $X$ (tương ứng, $Y$) bởi các tập mở affine, và gọi $X_i'$ (tương ứng, $Y_j'$) là hạn chế của $X'$ (tương ứng, $Y'$) lên $X'\cap U_i$ (tương ứng, $Y'\cap V_j$); đây là một tiền lược đồ con đóng của $U_i$ (tương ứng, $V_j$).
$U_i\times_S V_j$ được đồng nhất với hạn chế của $X\times_S Y$ lên $p^{-1}(U_i)\cap q^{-1}(V_j)$ \sref{I.3.2.7};
tương tự, nếu $p'$ và $q'$ là các phép chiếu của $X'\times_S Y'$, thì $X_i'\times_S Y_j'$ được đồng nhất với hạn chế của $X'\times_S Y'$ lên ${p'}^{-1}(X_i')\cap{q'}^{-1}(Y_j')$.
Đặt $\gamma=\alpha\times_S\beta$;
theo định nghĩa, ta có $p\circ\gamma=\alpha\circ p'$ và $q\circ\gamma=\beta\circ q'$;
vì $X_i'=\alpha^{-1}(U_i)$ và $Y_j'=\beta^{-1}(V_j)$, ta cũng có ${p'}^{-1}(X_i')=\gamma^{-1}(p^{-1}(U_i))$ và ${q'}^{-1}(Y_j')=\gamma^{-1}(q^{-1}(V_j))$, do đó
\[
  {p'}^{-1}(X_i')\cap{q'}^{-1}(Y_j')=\gamma^{-1}(p^{-1}(U_i)\cap q^{-1}(V_j))=\gamma^{-1}(U_i\times_S V_j),
\]
và khi đó ta kết luận như trong phần đầu của lập luận.

Vậy giả sử $X$, $Y$ và $S$ là affine, và cho $B$, $C$ và $A$ là các vành tương ứng của chúng.
Khi đó $B$ và $C$ là các $A$-đại số, còn $X'$ và $Y'$ là các lược đồ affine có các vành là các đại số thương $B'$ và $C'$ của $B$ và $C$, tương ứng.
Hơn nữa, ta có $\alpha=({}^a\rho,\widetilde{\rho})$ và $\beta=({}^a\sigma,\widetilde{\sigma})$, trong đó $\rho$ và $\sigma$ lần lượt là các đồng cấu chính tắc $B\to B'$ và $C\to C'$ \sref{I.1.7.3}.
Trong các điều kiện này, ta biết rằng $X\times_S Y$ (tương ứng, $X'\times_S Y'$) là một lược đồ affine có vành $B\otimes_A C$ (tương ứng, $B'\otimes_A C'$), và $\alpha\times_S\beta=({}^a\tau,\widetilde{\tau})$, trong đó $\tau$ là đồng cấu $\rho\otimes\sigma$ từ $B\otimes_A C$ vào $B'\otimes_A C'$ (Mệnh đề~\sref{I.3.2.2} và Hệ quả~\sref{I.3.2.3});
vì đồng cấu này toàn ánh, $\alpha\times_S\beta$ quả là một phép nhúng.
Hơn nữa, nếu $\mathfrak{b}$ (tương ứng, $\mathfrak{c}$) là hạt nhân của $\rho$ (tương ứng, $\sigma$), thì hạt nhân của $\tau$ là
$\operatorname{Im}(\mathfrak{b}\otimes_A C\to B\otimes_A C)+\operatorname{Im}(B\otimes_A\mathfrak{c}\to B\otimes_A C)$,
hay tương đương, tổng của các iđêan sinh bởi $u(\mathfrak{b})$ và $v(\mathfrak{c})$, trong đó $u$ (tương ứng, $v$) là đồng cấu $b\mapsto b\otimes 1$ (tương ứng, $c\mapsto 1\otimes c$).
Vì $p=({}^a u,\widetilde{u})$ và $q=({}^a v,\widetilde{v})$, hạt nhân này tương ứng, trong phổ nguyên tố của $B\otimes_A C$, với tập đóng $p^{-1}(X')\cap q^{-1}(Y')$, theo \sref{I.1.2.2.1} và Mệnh đề~\sref{I.1.1.2}[iii], qua đó kết thúc chứng minh.
\end{proof}

\begin{corollary}[4.3.2]
\label{I.4.3.2}
Nếu $f:X\to Y$ là một phép nhúng (tương ứng, một phép nhúng mở, một phép nhúng đóng) và là một $S$-cấu xạ, thì $f_{(S')}$ là một phép nhúng (tương ứng, một phép nhúng mở, một phép nhúng đóng) đối với mọi phép mở rộng $S'\to S$ của tiền lược đồ cơ sở.
\end{corollary}

\subsection{Ảnh ngược của một tiền lược đồ con}
\label{subsection:I.4.4}

\begin{proposition}[4.4.1]
\label{I.4.4.1}
Cho $f:X\to Y$ là một cấu xạ, $Y'$ là một tiền lược đồ con (tương ứng, một tiền lược đồ con đóng, một tiền lược đồ cảm sinh trên một tập mở) của $Y$, và $j:Y'\to Y$ là cấu xạ nhúng.
Khi đó phép chiếu $p:X\times_Y Y'\to X$ là một phép nhúng (tương ứng, một phép nhúng đóng, một phép nhúng mở);
không gian nền của tiền lược đồ con của $X$ liên kết với $p$ là $f^{-1}(Y')$;
hơn nữa, nếu $j'$ là cấu xạ nhúng của tiền lược đồ con này vào $X$, thì, với mọi cấu xạ $h:Z\to X$, hợp thành $f\circ h:Z\to Y$ bị $j$ chi phối khi và chỉ khi $h$ bị $j'$ chi phối.
\end{proposition}

\begin{proof}
Vì $p=1_X\times_Y j$ \sref{I.3.3.4}, nên khẳng định thứ nhất suy ra từ Mệnh đề~\sref{I.4.3.1};
khẳng định thứ hai là một trường hợp riêng của Hệ quả~\sref{I.3.5.10} (sau khi đổi vai trò của $X$ và $Y'$).
Cuối cùng, nếu $f\circ h=j\circ h'$, trong đó $h'$ là một cấu xạ $Z\to Y'$, thì từ định nghĩa của tích suy ra $h=p\circ u$, trong đó $u$ là một cấu xạ $Z\to X\times_Y Y'$; do đó suy ra khẳng định cuối cùng.
\end{proof}

Ta sẽ nói rằng tiền lược đồ con của $X$ được xác định như vậy là \emph{ảnh ngược} của tiền lược đồ con $Y'$ của $Y$ bởi cấu xạ $f$, thuật ngữ này phù hợp với thuật ngữ được đưa vào
\oldpage[I]{126}
một cách tổng quát hơn trong \sref{I.3.3.6}.
Khi nói về $f^{-1}(Y')$ như một tiền lược đồ con của $X$, ta luôn hiểu là tiền lược đồ con này.

Khi các tiền lược đồ $f^{-1}(Y')$ và $X$ là đồng nhất, $j'$ là cấu xạ đồng nhất và do đó mọi cấu xạ $h:Z\to X$ đều bị $j'$ chi phối; lấy $h=1_X$, ta thấy rằng \emph{cấu xạ $f:X\to Y$ phân tích thành $X\xrightarrow{g}Y'\xrightarrow{j}Y$}.

Khi $y$ là một điểm \emph{đóng} của $Y$ và $Y'=\Spec(\kres(y))$ là tiền lược đồ con đóng nhỏ nhất của $Y$ có $\{y\}$ làm không gian nền \sref{I.4.1.9}, tiền lược đồ con đóng $f^{-1}(Y')$ đẳng cấu chính tắc với \emph{thớ} $f^{-1}(y)$ được định nghĩa trong \sref{I.3.6.2}, và ta sẽ đồng nhất tiền lược đồ con đó với thớ này.

\begin{corollary}[4.4.2]
\label{I.4.4.2}
Cho $f:X\to Y$ và $g:Y\to Z$ là hai cấu xạ, và đặt $h=g\circ f$ là hợp thành của chúng.
Với mọi tiền lược đồ con $Z'$ của $Z$, các tiền lược đồ con $f^{-1}(g^{-1}(Z'))$ và $h^{-1}(Z')$ của $X$ là đồng nhất.
\end{corollary}

\begin{proof}
Điều này suy ra từ sự tồn tại của đẳng cấu chính tắc $X\times_Y(Y\times_Z Z')\isoto X\times_Z Z'$ \sref{I.3.3.9.1}.
\end{proof}

\begin{corollary}[4.4.3]
\label{I.4.4.3}
Cho $X'$ và $X''$ là hai tiền lược đồ con của $X$, và cho $j':X'\to X$ và $j'':X''\to X$ là các cấu xạ nhúng của chúng;
khi đó ${j'}^{-1}(X'')$ và ${j''}^{-1}(X')$ đều bằng cận dưới lớn nhất $\inf(X',X'')$ của $X'$ và $X''$ đối với quan hệ thứ tự $\leq$ giữa các tiền lược đồ con, và cận dưới lớn nhất này đẳng cấu chính tắc với $X'\times_X X''$.
\end{corollary}

\begin{proof}
Điều này suy ra ngay từ Mệnh đề~\sref{I.4.4.1} và Hệ quả~\sref{I.4.1.10}.
\end{proof}

\begin{corollary}[4.4.4]
\label{I.4.4.4}
Cho $f:X\to Y$ là một cấu xạ, và cho $Y'$ và $Y''$ là hai tiền lược đồ con của $Y$;
khi đó
\[
  f^{-1}(\inf(Y',Y''))=\inf(f^{-1}(Y'),f^{-1}(Y'')).
\]
\end{corollary}

\begin{proof}
Điều này suy ra từ sự tồn tại của đẳng cấu chính tắc giữa $(X\times_Y Y')\times_X(X\times_Y Y'')$ và $X\times_Y(Y'\times_Y Y'')$ \sref{I.3.3.9.1}.
\end{proof}

\begin{proposition}[4.4.5]
\label{I.4.4.5}
Cho $f:X\to Y$ là một cấu xạ và $Y'$ là một tiền lược đồ con đóng của $Y$ được xác định bởi một bó iđêan quasi-coherent $\sh{K}$ của $\sh{O}_Y$ \sref{I.4.1.3};
khi đó tiền lược đồ con đóng $f^{-1}(Y')$ của $X$ được xác định bởi bó iđêan quasi-coherent $f^*(\sh{K})\sh{O}_X$ của $\sh{O}_X$.
\end{proposition}

\begin{proof}
Vấn đề hiển nhiên có tính địa phương trên $X$ và $Y$;
do đó chỉ cần nhận thấy rằng, nếu \erratum[II]{$A$ là một $B$-đại số} và $\mathfrak{K}$ là một iđêan của $B$, thì ta có $A\otimes_B(B/\mathfrak{K})=A/\mathfrak{K}A$, rồi áp dụng \sref{I.1.6.9}.
\end{proof}

\begin{corollary}[4.4.6]
\label{I.4.4.6}
Cho $X'$ là một tiền lược đồ con đóng của $X$ được xác định bởi một bó iđêan quasi-coherent $\sh{J}$ của $\sh{O}_X$, và $i$ là cấu xạ nhúng $X'\to X$;
để hạn chế $f\circ i$ của $f$ lên $X'$ bị cấu xạ nhúng $j:Y'\to Y$ chi phối \emph{(nói cách khác, phân tích thành $j\circ g$, trong đó $g$ là một cấu xạ $X'\to Y'$)}, điều kiện cần và đủ là $f^*(\sh{K})\sh{O}_X\subset\sh{J}$.
\end{corollary}

\begin{proof}
Chỉ cần áp dụng Mệnh đề~\sref{I.4.4.1} cho $i$, có xét đến Mệnh đề~\sref{I.4.4.5}.
\end{proof}

\subsection{Các phép nhúng địa phương và các đẳng cấu địa phương}
\label{subsection:I.4.5}

\begin{definition}[4.5.1]
\label{I.4.5.1}
Cho $f:X\to Y$ là một cấu xạ của các tiền lược đồ.
Ta nói $f$ là một \emph{phép nhúng địa phương tại một điểm $x\in X$} nếu tồn tại một lân cận mở $U$ của $x$ trong $X$ và một lân cận mở $V$ của $f(x)$ trong $Y$ sao cho hạn chế của $f$ trên tiền lược đồ cảm sinh $U$ là một phép nhúng đóng từ $U$ vào tiền lược đồ cảm sinh $V$.
Ta nói $f$ là một phép nhúng địa phương nếu $f$ là một phép nhúng địa phương tại mọi điểm của $X$.
\end{definition}

\begin{definition}[4.5.2]
\label{I.4.5.2}
Ta nói một cấu xạ $f:X\to Y$ của các tiền lược đồ là một đẳng cấu địa phương tại
\oldpage[I]{127}
một điểm $x\in X$ nếu tồn tại một lân cận mở $U$ của $x$ trong $X$ sao cho hạn chế của $f$ trên tiền lược đồ cảm sinh $U$ là một phép nhúng mở từ $U$ vào $Y$.
Ta nói $f$ là một đẳng cấu địa phương nếu $f$ là một đẳng cấu địa phương tại mọi điểm của $X$.
\end{definition}

\begin{env}[4.5.3]
\label{I.4.5.3}
Một phép nhúng (tương ứng, một phép nhúng đóng) $f:X\to Y$ có thể được đặc trưng như một phép nhúng địa phương sao cho $f$ là một phép đồng phôi từ không gian nền của $X$ lên một tập con (tương ứng, một tập con đóng) của $Y$.
Một phép nhúng mở $f$ có thể được đặc trưng như một đẳng cấu địa phương \emph{đơn ánh}.
\end{env}

\begin{proposition}[4.5.4]
\label{I.4.5.4}
Cho $X$ là một tiền lược đồ bất khả quy và $f:X\to Y$ là một cấu xạ đơn ánh trội.
Nếu $f$ là một phép nhúng địa phương, thì $f$ là một phép nhúng và $f(X)$ mở trong $Y$.
\end{proposition}

\begin{proof}
Cho $x\in X$, và cho $U$ là một lân cận mở của $x$, $V$ là một lân cận mở của $f(x)$ trong $Y$ sao cho hạn chế của $f$ trên $U$ là một phép nhúng đóng vào $V$;
vì $U$ trù mật trong $X$, nên theo giả thiết $f(U)$ trù mật trong $Y$, do đó $f(U)=V$ và $f$ là một phép đồng phôi từ $U$ lên $V$;
giả thiết rằng $f$ là đơn ánh kéo theo $f^{-1}(V)=U$, và mệnh đề suy ra ngay.
\end{proof}

\begin{proposition}[4.5.5]
\label{I.4.5.5}
\medskip\noindent
\begin{enumerate}
  \item[{\rm(i)}] Hợp thành của hai phép nhúng địa phương (tương ứng, của hai đẳng cấu địa phương) là một phép nhúng địa phương (tương ứng, một đẳng cấu địa phương).
  \item[{\rm(ii)}] Cho $f:X\to X'$ và $g:Y\to Y'$ là hai $S$-cấu xạ.
    Nếu $f$ và $g$ là các phép nhúng địa phương (tương ứng, các đẳng cấu địa phương), thì $f\times_S g$ cũng là một phép nhúng địa phương (tương ứng, một đẳng cấu địa phương).
  \item[{\rm(iii)}] Nếu một $S$-cấu xạ $f$ là một phép nhúng địa phương (tương ứng, một đẳng cấu địa phương), thì $f_{(S')}$ cũng vậy đối với mọi mở rộng $S'\to S$ của tiền lược đồ cơ sở.
\end{enumerate}
\end{proposition}

\begin{proof}
Theo \sref{I.3.5.1}, chỉ cần chứng minh (i) và (ii).

(i) suy ra ngay từ tính bắc cầu của các phép nhúng đóng (tương ứng, mở) \sref{I.4.2.5}\footnote{\emph{[Dịch.] Nguồn tiếng Pháp dẫn (4.2.4) tại đây; tính bắc cầu là Mệnh đề~\hyperref[I.4.2.5]{4.2.5}.}} và từ sự kiện sau: nếu $f$ là một phép đồng phôi từ $X$ lên một tập con đóng của $Y$, thì với mọi tập mở $U\subset X$, $f(U)$ mở trong $f(X)$; do đó tồn tại một tập con mở $V$ của $Y$ sao cho $f(U)=V\cap f(X)$, và vì vậy $f(U)$ đóng trong $V$.

Để chứng minh (ii), cho $z\in X\times_S Y$, đặt $z'=(f\times_S g)(z)$\footnote{\emph{[Dịch.] Nguồn tiếng Pháp sử dụng $z$ và $z'$ ở dưới mà không giới thiệu; phần chèn này cung cấp điểm cần thiết và ảnh của nó.}}, gọi $p$ và $q$ là các phép chiếu của $X\times_S Y$, và $p'$ và $q'$ là các phép chiếu của $X'\times_S Y'$.
Theo giả thiết, tồn tại các lân cận mở $U$, $U'$, $V$ và $V'$ lần lượt của $x=p(z)$, $x'=p'(z')$, $y=q(z)$ và $y'=q'(z')$, sao cho các hạn chế của $f$ và $g$ trên $U$ và $V$ lần lượt là các phép nhúng đóng (tương ứng, mở) vào $U'$ và $V'$.
Vì các không gian nền của $U\times_S V$ và $U'\times_S V'$ lần lượt được đồng nhất với các lân cận mở $p^{-1}(U)\cap q^{-1}(V)$ và ${p'}^{-1}(U')\cap{q'}^{-1}(V')$ của $z$ và $z'$ \sref{I.3.2.7}, mệnh đề suy ra từ Mệnh đề~\sref{I.4.3.1}.
\end{proof}

\section{Các tiền lược đồ rút gọn; điều kiện tách}
\label{section:I.5}

\subsection{Các tiền lược đồ rút gọn}
\label{subsection:I.5.1}

\begin{proposition}[5.1.1]
\label{I.5.1.1}
Cho $(X,\sh{O}_X)$ là một tiền lược đồ và $\sh{B}$ là một $\sh{O}_X$-đại số quasi-coherent.
Khi đó tồn tại duy nhất một $\sh{O}_X$-môđun quasi-coherent $\sh{N}$ sao cho cuống $\sh{N}_x$ tại mỗi điểm $x\in X$ là căn lũy linh của vành $\sh{B}_x$.
Khi $X$ là affine và do đó $\sh{B}=\widetilde{B}$, trong đó $B$ là một đại số trên $A(X)$, ta có $\sh{N}=\widetilde{\nilrad}$, trong đó $\nilrad$ là căn lũy linh của $B$.
\end{proposition}

\begin{proof}
\oldpage[I]{128}
Vì bài toán có tính địa phương, ta quy về việc chứng minh khẳng định cuối.
Ta biết $\widetilde{\nilrad}$ là một $\sh{O}_X$-môđun quasi-coherent \sref{I.1.4.1}, và cuống của nó tại điểm $x\in X$ là iđêan $\nilrad_x$ của vành phân thức $B_x$;
do đó chỉ cần chứng minh rằng căn lũy linh của $B_x$ nằm trong $\nilrad_x$, vì bao hàm ngược lại là hiển nhiên.
Cho $z/s$ là một phần tử của căn lũy linh của $B_x$, với $z\in B$ và $s\not\in\mathfrak{j}_x$;
theo giả thiết, tồn tại một số nguyên $k$ sao cho $(z/s)^k=0$, điều này có nghĩa là tồn tại $t\not\in\mathfrak{j}_x$ sao cho $tz^k=0$.
Suy ra $(tz)^k=0$ và do đó $z/s=(tz)/(ts)\in\nilrad_x$.
\end{proof}

Ta gọi $\sh{O}_X$-môđun quasi-coherent $\sh{N}$ được xác định như trên là \emph{căn lũy linh} của $\sh{O}_X$-đại số $\sh{B}$; đặc biệt, ta ký hiệu bởi $\sh{N}_X$ căn lũy linh của $\sh{O}_X$.

\begin{corollary}[5.1.2]
\label{I.5.1.2}
Cho $X$ là một tiền lược đồ;
tiền lược đồ con đóng của $X$ được xác định bởi bó iđêan $\sh{N}_X$ là tiền lược đồ con rút gọn duy nhất \sref[0]{0.4.1.4} của $X$ có $X$ làm không gian nền;
nó cũng là tiền lược đồ con nhỏ nhất của $X$ có $X$ làm không gian nền.
\end{corollary}

\begin{proof}
Cho $Y$ là tiền lược đồ con đóng của $X$ được xác định bởi $\sh{N}_X$.
Vì bó cấu trúc của nó là $\sh{O}_X/\sh{N}_X$, hiển nhiên $Y$ rút gọn và có $X$ làm không gian nền, bởi $\sh{N}_x\neq\sh{O}_x$ với mọi $x\in X$.
Để chứng minh các khẳng định còn lại, hãy nhận xét rằng một tiền lược đồ con $Z$ của $X$ có $X$ làm không gian nền được xác định bởi một bó iđêan $\sh{I}$ \sref{I.4.1.3} sao cho $\sh{I}_x\neq\sh{O}_x$ với mọi $x\in X$.
Ta có thể giới hạn ở trường hợp $X$ affine, giả sử $X=\Spec(A)$ và $\sh{I}=\widetilde{\mathfrak{I}}$, trong đó $\mathfrak{I}$ là một iđêan của $A$;
khi đó, với mọi $x\in X$, ta có $\mathfrak{I}\subset\mathfrak{j}_x$\footnote{\emph{[Dịch.] Nguồn tiếng Pháp in ở đây $\mathfrak{I}_x\subset\mathfrak{j}_x$, một so sánh không đúng kiểu giữa một iđêan đã địa phương hóa của $A_x$ và một iđêan của $A$; lập luận đòi hỏi $\mathfrak{I}\subset\mathfrak{j}_x$.}}, do đó $\mathfrak{I}$ nằm trong mọi iđêan nguyên tố của $A$ và vì thế nằm trong giao $\nilrad$ của chúng, tức căn lũy linh của $A$.
Điều này chứng minh rằng $Y$ là tiền lược đồ con nhỏ nhất của $X$ có $X$ làm không gian nền \sref{I.4.1.9};
hơn nữa, nếu $Z$ khác $Y$, thì nhất thiết $\sh{I}_x\neq\sh{N}_x$ với ít nhất một $x\in X$, và do \sref{I.5.1.1}, $Z$ không rút gọn.
\end{proof}

\begin{definition}[5.1.3]
\label{I.5.1.3}
Ta gọi tiền lược đồ rút gọn liên kết với một tiền lược đồ $X$, và ký hiệu bởi $X_\red$, tiền lược đồ con rút gọn duy nhất của $X$ có $X$ làm không gian nền.
\end{definition}

Do đó, nói một tiền lược đồ $X$ rút gọn có nghĩa là $X=X_\red$.

\begin{proposition}[5.1.4]
\label{I.5.1.4}
Phổ nguyên tố của một vành $A$ là một tiền lược đồ rút gọn (tương ứng, nguyên) \sref{I.2.1.8}\footnote{\emph{[Dịch.] Nguồn tiếng Pháp dẫn ở đây đến (2.1.7), mục nói về hạn chế lên các tập con mở; định nghĩa thích hợp của tiền lược đồ nguyên là (2.1.8).}} khi và chỉ khi $A$ là một vành rút gọn (tương ứng, một miền nguyên).
\end{proposition}

\begin{proof}
Thật vậy, từ \sref{I.5.1.1} suy ra ngay rằng điều kiện $\sh{N}=(0)$ là cần và đủ để $X=\Spec(A)$ rút gọn;
khẳng định về các miền nguyên khi ấy là hệ quả của \sref{I.1.1.13}.
\end{proof}

Vì mọi vành phân thức khác không của một miền nguyên vẫn là miền nguyên, từ \sref{I.5.1.4} suy ra rằng, với mọi tiền lược đồ \emph{nguyên địa phương} $X$, $\sh{O}_x$ là một miền nguyên với mọi $x\in X$.
Chiều đảo đúng khi không gian nền của $X$ là \emph{Noether địa phương}:
thật vậy, khi đó $X$ rút gọn, và nếu $U$ là một tập mở affine của $X$ đồng thời là một không gian Noether, thì $U$ chỉ có hữu hạn thành phần bất khả quy, nên vành $A$ của nó chỉ có hữu hạn iđêan nguyên tố tối tiểu \sref{I.1.1.14}.
Nếu hai thành phần bất khả quy của $U$ có một điểm chung $x$, thì $\sh{O}_x$ sẽ có ít nhất hai iđêan nguyên tố tối tiểu phân biệt và do đó không phải là miền nguyên;
vì thế, các thành phần của $U$ là những tập mở đôi một rời nhau, và mỗi thành phần là nguyên.

\begin{env}[5.1.5]
\label{I.5.1.5}
Cho $f=(\psi,\theta):X\to Y$ là một cấu xạ tiền lược đồ;
\oldpage[I]{129}
đồng cấu $\theta_x^\sharp:\sh{O}_{\psi(x)}\to\sh{O}_x$ gửi mọi phần tử lũy linh của $\sh{O}_{\psi(x)}$ đến một phần tử lũy linh của $\sh{O}_x$;
khi chuyển sang các thương, $\theta^\sharp$ cảm sinh một đồng cấu
\[
  \omega:\psi^*(\sh{O}_Y/\sh{N}_Y)\to\sh{O}_X/\sh{N}_X.
\]
Rõ ràng, với mọi $x\in X$, $\omega_x:\sh{O}_{\psi(x)}/\sh{N}_{\psi(x)}\to\sh{O}_x/\sh{N}_x$ là một đồng cấu địa phương, nên $(\psi,\omega^\flat)$ là một cấu xạ tiền lược đồ $X_\red\to Y_\red$, mà ta ký hiệu bởi $f_\red$ và gọi là cấu xạ \emph{rút gọn} liên kết với $f$.
Hiển nhiên, với hai cấu xạ $f:X\to Y$ và $g:Y\to Z$, ta có $(g\circ f)_\red=g_\red\circ f_\red$, do đó đã xác định $X_\red$ như một \emph{hàm tử}, \emph{hiệp biến} theo $X$.

Định nghĩa trên cho thấy sơ đồ
\[
  \xymatrix{
    X_\red\ar[r]^{f_\red}\ar[d] &
    Y_\red\ar[d]\\
    X\ar[r]_f &
    Y
  }
\]
là giao hoán, trong đó các mũi tên thẳng đứng là các cấu xạ nhúng;
nói cách khác, $X_\red\to X$ là một cấu xạ \emph{có tính hàm tử}.
Đặc biệt, lưu ý rằng nếu $X$ rút gọn, thì mọi cấu xạ $f:X\to Y$ phân tích thành $X\xrightarrow{f_\red}Y_\red\to Y$;
nói cách khác, $f$ bị cấu xạ nhúng $Y_\red\to Y$ chi phối.
\end{env}

\begin{proposition}[5.1.6]
\label{I.5.1.6}
Cho $f:X\to Y$ là một cấu xạ.
Nếu $f$ là toàn ánh (tương ứng, radicial, một phép nhúng, một phép nhúng đóng, một phép nhúng mở, một phép nhúng địa phương, một đẳng cấu địa phương), thì $f_\red$ cũng vậy.
Ngược lại, nếu $f_\red$ là toàn ánh (tương ứng, radicial), thì $f$ cũng vậy.
\end{proposition}

\begin{proof}
Mệnh đề là hiển nhiên nếu $f$ là toàn ánh.
Nếu $f$ là radicial, điều này suy ra từ việc, với mọi $x\in X$, trường $\kres(x)$ là như nhau đối với hai tiền lược đồ $X$ và $X_\red$ \sref{I.3.5.8}.
Cuối cùng, nếu $f=(\psi,\theta)$ là một phép nhúng, một phép nhúng đóng hoặc một phép nhúng địa phương (tương ứng, một phép nhúng mở hoặc một đẳng cấu địa phương), thì mệnh đề suy ra từ việc, nếu $\theta_x^\sharp$ là toàn ánh (tương ứng, song ánh), thì đồng cấu thu được khi lấy thương theo các căn lũy linh của $\sh{O}_{\psi(x)}$ và $\sh{O}_x$ cũng vậy (\sref{I.5.1.2} và \sref{I.4.2.2}) (xem \sref{I.5.5.12}).
\end{proof}

\begin{proposition}[5.1.7]
\label{I.5.1.7}
Nếu $X$ và $Y$ là hai $S$-tiền lược đồ, thì các tiền lược đồ $X_\red\times_{S_\red}Y_\red$ và $X_\red\times_S Y_\red$ là đồng nhất và được đồng nhất chính tắc với một tiền lược đồ con của $X\times_S Y$ có cùng không gian nền với tích đó.
\end{proposition}

\begin{proof}
Sự đồng nhất chính tắc của $X_\red\times_S Y_\red$ với một tiền lược đồ con của $X\times_S Y$ có cùng không gian nền suy ra từ \sref{I.4.3.1}.
Hơn nữa, nếu $\vphi$ và $\psi$ lần lượt là các cấu xạ cấu trúc $X_\red\to S$ và $Y_\red\to S$, thì chúng phân tích qua $S_\red$ \sref{I.5.1.5}; vì $S_\red\to S$ là một đơn cấu, khẳng định thứ nhất của mệnh đề suy ra từ \sref{I.3.2.4}.
\end{proof}

\begin{corollary}[5.1.8]
\label{I.5.1.8}
Các tiền lược đồ $(X\times_S Y)_\red$ và $(X_\red\times_{S_\red}Y_\red)_\red$ được đồng nhất chính tắc với nhau.
\end{corollary}

\begin{proof}
Điều này suy ra từ \sref{I.5.1.2} và \sref{I.5.1.7}.
\end{proof}

Cần lưu ý rằng, ngay cả khi $X$ và $Y$ là các tiền lược đồ rút gọn, $X\times_S Y$ không nhất thiết rút gọn, vì tích tenxơ của hai đại số rút gọn có thể có các phần tử lũy linh.

\begin{proposition}[5.1.9]
\label{I.5.1.9}
\oldpage[I]{130}
Cho $X$ là một tiền lược đồ, và cho $\sh{I}$ là một bó iđêan quasi-coherent của $\sh{O}_X$ sao cho $\sh{I}^n=0$ với một số nguyên $n>0$ nào đó.
Cho $X_0$ là tiền lược đồ con đóng $(X,\sh{O}_X/\sh{I})$ của $X$;
điều kiện cần và đủ để $X$ là một lược đồ affine là $X_0$ là một lược đồ affine.
\end{proposition}

Điều kiện hiển nhiên là cần; ta chứng minh rằng nó là đủ.
Nếu đặt $X_k=(X,\sh{O}_X/\sh{I}^{k+1})$, chỉ cần chứng minh bằng quy nạp theo $k$ rằng các $X_k$ đều affine; do đó ta quy về trường hợp $\sh{I}^2=0$.
Đặt
\begin{align*}
  A&=\Gamma(X,\sh{O}_X)\\
  A_0&=\Gamma(X_0,\sh{O}_{X_0})=\Gamma(X,\sh{O}_X/\sh{I}).
\end{align*}
Đồng cấu chính tắc $\sh{O}_X\to\sh{O}_X/\sh{I}$ cảm sinh một đồng cấu vành $\vphi:A\to A_0$.
Sau đây ta sẽ thấy $\vphi$ là \emph{toàn ánh}, nên
\[
\label{I.5.1.9.1}
  0\to\Gamma(X,\sh{I})\to\Gamma(X,\sh{O}_X)\to\Gamma(X,\sh{O}_X/\sh{I})\to 0\tag{5.1.9.1}
\]
là một dãy \emph{khớp}.
Giả sử điểm này đã được chứng minh, và ta chứng minh rằng từ đó suy ra mệnh đề.
Lưu ý rằng $\mathfrak{K}=\Gamma(X,\sh{I})$ là một iđêan có bình phương bằng không trong $A$, nên là một môđun trên $A_0=A/\mathfrak{K}$.
Theo giả thiết, ta có $X_0=\Spec(A_0)$ và, vì các không gian nền của $X_0$ và $X$ trùng nhau, $\mathfrak{K}=\Gamma(X_0,\sh{I})$;
hơn nữa, vì $\sh{I}^2=0$, $\sh{I}$ là một $(\sh{O}_X/\sh{I})$-môđun quasi-coherent, nên $\sh{I}\cong\widetilde{\mathfrak{K}}$ và $\mathfrak{K}_x=\sh{I}_x$ với mọi $x\in X_0$ \sref{I.1.4.1}.

Trong các điều kiện này, cho $X'=\Spec(A)$, và xét cấu xạ $f=(\psi,\theta):X\to X'$ của các tiền lược đồ ứng với ánh xạ đồng nhất $A\to\Gamma(X,\sh{O}_X)$ \sref{I.2.2.4}.
Với mọi tập con mở affine $V$ của $X$, sơ đồ
\[
  \xymatrix{
    A\ar[r]\ar[d] &
    \Gamma(V,\sh{O}_X|V)\ar[d]\\
    A_0=A/\mathfrak{K}\ar[r] &
    \Gamma(V,\sh{O}_{X_0}|V)
  }
\]
giao hoán, từ đó suy ra sơ đồ
\[
  \xymatrix{
    X' &
    X\ar[l]_f\\
    X'_0\ar[u]^{j'} &
    X_0\ar[l]^{f_0}\ar[u]_j
  }
\]
cũng giao hoán (trong đó $X'_0$ là tiền lược đồ con đóng của $X'$ được xác định bởi bó iđêan quasi-coherent $\widetilde{\mathfrak{K}}$, còn $j$ và $j'$ là các cấu xạ nhúng chính tắc).
Nhưng vì $X_0$ là affine, $f_0$ là một đẳng cấu và, vì các ánh xạ liên tục nền của $j$ và $j'$ là các ánh xạ đồng nhất, ta thấy ngay rằng $\psi:X\to X'$ là một phép đồng phôi.
Hơn nữa, đẳng thức $\mathfrak{K}_x=\sh{I}_x$ cho thấy hạn chế của $\theta^\sharp:\psi^*(\sh{O}_{X'})\to\sh{O}_X$ là một \emph{đẳng cấu} từ $\psi^*(\widetilde{\mathfrak{K}})$ lên $\sh{I}$;
mặt khác, khi chuyển sang các bó thương, $\theta^\sharp$ cho một \emph{đẳng cấu} $\psi^*(\sh{O}_{X'}/\widetilde{\mathfrak{K}})\to\sh{O}_X/\sh{I}$, vì $f_0$ là một đẳng cấu;
do đó, theo bổ đề năm (M,~I,~1.1), ta suy ra ngay rằng chính $\theta^\sharp$ là một đẳng cấu, và vì vậy $f$ là một \emph{đẳng cấu}; do đó $X$ là affine.

Vì vậy, tất cả quy về việc chứng minh rằng dãy \sref{I.5.1.9.1} là khớp; điều này sẽ suy ra từ việc chứng minh $\HH^1(X,\sh{I})=0$.
\oldpage[I]{131}
Nhưng $\HH^1(X,\sh{I})=\HH^1(X_0,\sh{I})$, và ta đã thấy rằng $\sh{I}$ là một $\sh{O}_{X_0}$-môđun quasi-coherent.
Do đó, khẳng định của ta sẽ suy ra từ bổ đề sau.
\begin{lemma}[5.1.9.2]
\label{I.5.1.9.2}
Nếu $Y$ là một lược đồ affine và $\sh{F}$ là một $\sh{O}_Y$-môđun quasi-coherent, thì $\HH^1(Y,\sh{F})=0$.
\end{lemma}

\begin{proof}
Bổ đề này sẽ được chứng minh trong Chương~III, §1, như một hệ quả của định lý tổng quát hơn, theo đó $\HH^i(Y,\sh{F})=0$ với mọi $i>0$.
Để đưa ra một chứng minh độc lập, ta nhận xét rằng $\HH^1(Y,\sh{F})$ có thể được đồng nhất với môđun $\Ext_{\sh{O}_Y}^1(Y;\sh{O}_Y,\sh{F})$ gồm các lớp mở rộng của $\sh{O}_Y$-môđun $\sh{O}_Y$ bởi $\sh{O}_Y$-môđun $\sh{F}$ (T,~4.2.3);
vì vậy, tất cả quy về việc chứng minh rằng một mở rộng như thế $\sh{G}$ là tầm thường.
Thật vậy, với mọi $y\in Y$, tồn tại một lân cận $V$ của $y$ trong $Y$ sao cho $\sh{G}|V$ đẳng cấu với $\sh{F}|V\oplus\sh{O}_Y|V$\footnote{\emph{[Dịch.] Nguồn tiếng Pháp in ở đây $\sh{F}|Y$; hạn chế cần thiết trên lân cận $V$ là $\sh{F}|V$.}} \sref[0]{0.5.4.9};
từ đó suy ra $\sh{G}$ là một $\sh{O}_Y$-môđun \emph{quasi-coherent}.
Nếu $A$ là vành của $Y$, thì $\sh{F}=\widetilde{M}$ và $\sh{G}=\widetilde{N}$, trong đó $M$ và $N$ là các $A$-môđun, và theo giả thiết, $N$ là một mở rộng của $A$-môđun $A$ bởi $A$-môđun $M$ \sref{I.1.3.11}.
Vì mở rộng này nhất thiết tầm thường, bổ đề được chứng minh, và do đó \sref{I.5.1.9} cũng được chứng minh.
\end{proof}

\begin{corollary}[5.1.10]
\label{I.5.1.10}
Cho $X$ là một tiền lược đồ sao cho $\sh{N}_X$ lũy linh.
Để $X$ là một lược đồ affine, điều kiện cần và đủ là $X_\red$ là một lược đồ affine.
\end{corollary}

\subsection{Sự tồn tại của một tiền lược đồ con có không gian nền cho trước}
\label{subsection:I.5.2}

\begin{proposition}[5.2.1]
\label{I.5.2.1}
Với mọi không gian con đóng địa phương $Y$ của không gian nền của một tiền lược đồ $X$, tồn tại duy nhất một tiền lược đồ con rút gọn của $X$ có $Y$ làm không gian nền.
\end{proposition}

\begin{proof}
Tính duy nhất suy ra từ \sref{I.5.1.2}, nên chỉ còn phải chứng minh sự tồn tại của tiền lược đồ đang xét.

Nếu $X$ là affine, với vành $A$, và $Y$ đóng trong $X$, thì mệnh đề là hiển nhiên:
$\mathfrak{j}(Y)$ là iđêan lớn nhất $\mathfrak{a}\subset A$ sao cho $V(\mathfrak{a})=Y$, và nó bằng căn của chính nó \sref{I.1.1.4}[i], do đó $A/\mathfrak{j}(Y)$ là một vành rút gọn.

Trong trường hợp tổng quát, với mọi tập mở affine $U\subset X$ sao cho $U\cap Y$ đóng trong $U$, xét tiền lược đồ con đóng $Y_U$ của $U$ được xác định bởi bó iđêan liên kết với iđêan $\mathfrak{j}(U\cap Y)$ của $A(U)$.
Tiền lược đồ con này rút gọn.
Ta khẳng định rằng, nếu $V$ là một tập mở affine của $X$ chứa trong $U$, thì $Y_V$ được \emph{cảm sinh} bởi $Y_U$ trên $V\cap Y$.
Thật vậy, tiền lược đồ cảm sinh này là một tiền lược đồ con đóng rút gọn của $V$ có $V\cap Y$ làm không gian nền;
do đó, tính duy nhất của $Y_V$ chứng minh khẳng định.
\end{proof}

\begin{proposition}[5.2.2]
\label{I.5.2.2}
Cho $X$ là một tiền lược đồ rút gọn, $f:X \to Y$ là một cấu xạ, và $Z$ là một tiền lược đồ con đóng của $Y$ sao cho $f(X) \subset Z$.
Khi đó $f$ phân tích thành $X \xrightarrow{g} Z \xrightarrow{j} Y$, trong đó $j$ là cấu xạ nhúng.
\end{proposition}

\begin{proof}
Từ các giả thiết suy ra rằng tiền lược đồ con đóng $f^{-1}(Z)$ của $X$ có toàn bộ $X$ làm không gian nền \sref{I.4.4.1};
vì $X$ rút gọn, tiền lược đồ con đóng này trùng với $X$ \sref{I.5.1.2}, và do đó mệnh đề suy ra từ \sref{I.4.4.1}.
\end{proof}

\begin{corollary}[5.2.3]
\label{I.5.2.3}
Cho $X$ là một tiền lược đồ con rút gọn của một tiền lược đồ $Y$; nếu $Z$ là tiền lược đồ con đóng rút gọn của $Y$ có $\overline{X}$ làm không gian nền, thì $X$ là một tiền lược đồ con cảm sinh trên một tập mở của $Z$.
\end{corollary}

\begin{proof}
\oldpage[I]{132}
Thật vậy, tồn tại một tập mở $U$ của $Y$ sao cho $X=U\cap\overline{X}$;
vì theo \sref{I.5.2.2}, $X$ là một tiền lược đồ con rút gọn của $Z$, nên tiền lược đồ con $X$ được $Z$ cảm sinh trên không gian con mở $X$ nhờ tính duy nhất \sref{I.5.2.1}.
\end{proof}

\begin{corollary}[5.2.4]
\label{I.5.2.4}
Cho $f:X\to Y$ là một cấu xạ, và $X'$ (tương ứng, $Y'$) là một tiền lược đồ con đóng của $X$ (tương ứng, $Y$) được xác định bởi một bó iđêan quasi-coherent $\sh{J}$ (tương ứng, $\sh{K}$) của $\sh{O}_X$ (tương ứng, $\sh{O}_Y$).
Giả sử $X'$ rút gọn và $f(X')\subset Y'$.
Khi đó $f^*(\sh{K})\sh{O}_X\subset\sh{J}$.
\end{corollary}

\begin{proof}
Vì theo \sref{I.5.2.2}, hạn chế của $f$ lên $X'$ phân tích thành $X'\to Y'\to Y$, chỉ cần áp dụng \sref{I.4.4.6}.
\end{proof}

\subsection{Đường chéo; đồ thị của một cấu xạ}
\label{subsection:I.5.3}
\label{I.5.3}

\begin{env}[5.3.1]
\label{I.5.3.1}
Cho $X$ là một $S$-tiền lược đồ;
ta gọi \emph{cấu xạ đường chéo} của $X$ vào $X\times_S X$, và ký hiệu nó bởi $\Delta_{X|S}$, hoặc $\Delta_X$, hoặc thậm chí $\Delta$ nếu không thể có nhầm lẫn, là $S$-cấu xạ $(1_X,1_X)_S$, nghĩa là $S$-cấu xạ duy nhất $\Delta_X$ sao cho
\[
\label{I.5.3.1.1}
  p_1\circ\Delta_X=p_2\circ\Delta_X=1_X,
  \tag{5.3.1.1}
\]
trong đó $p_1$ và $p_2$ là các phép chiếu của $X\times_S X$ (Định nghĩa~\sref{I.3.2.1}).
Nếu $f:T\to X$ và $g:T\to Y$ là hai $S$-cấu xạ, ta kiểm tra ngay rằng
\[
\label{I.5.3.1.2}
  (f,g)_S=(f\times_S g)\circ\Delta_{T|S}.
  \tag{5.3.1.2}
\]

Người đọc lưu ý rằng định nghĩa trên và các kết quả được phát biểu từ \sref{I.5.3.1} đến \sref{I.5.3.8} đều đúng trong mọi phạm trù, \emph{miễn là các tích xuất hiện ở đó tồn tại trong phạm trù ấy}.
\end{env}

\begin{proposition}[5.3.2]
\label{I.5.3.2}
Cho $X$ và $Y$ là hai $S$-tiền lược đồ;
nếu ta đồng nhất một cách chính tắc tích $(X\times Y)\times(X\times Y)$ với $(X\times X)\times(Y\times Y)$, thì cấu xạ $\Delta_{X\times Y}$ được đồng nhất với $\Delta_X\times\Delta_Y$.
\end{proposition}

\begin{proof}
Thật vậy, nếu $p_1:X\times X\to X$ và $q_1:Y\times Y\to Y$ là các phép chiếu lên thừa số thứ nhất, thì phép chiếu lên thừa số thứ nhất $(X\times Y)\times(X\times Y)\to X\times Y$ được đồng nhất với $p_1\times q_1$, và ta có
\[
  (p_1\times q_1)\circ(\Delta_X\times\Delta_Y)=(p_1\circ\Delta_X)\times(q_1\circ\Delta_Y)=1_{X\times Y}
\]
Lập luận tương tự áp dụng cho các phép chiếu lên thừa số thứ hai.
\end{proof}

\phantomsection
\label{I.5.3.3}

\begin{corollary}[5.3.4]
\label{I.5.3.4}
Với mọi mở rộng $S'\to S$ của tiền lược đồ cơ sở, $\Delta_{X_{(S')}}$ được đồng nhất một cách chính tắc với $(\Delta_X)_{(S')}$.
\end{corollary}

\begin{proof}
Chỉ cần lưu ý rằng $(X\times_S X)_{(S')}$ được đồng nhất một cách chính tắc với $X_{(S')}\times_{S'}X_{(S')}$ \sref{I.3.3.10}.
\end{proof}

\begin{proposition}[5.3.5]
\label{I.5.3.5}
Cho $X$ và $Y$ là hai $S$-tiền lược đồ, và cho $\vphi:S\to T$ là một cấu xạ tiền lược đồ cho phép xem mọi $S$-tiền lược đồ như một $T$-tiền lược đồ.
Cho $f:X\to S$ và $g:Y\to S$\footnote{\emph{[Dịch.] Nguồn tiếng Pháp chỉ in $Y\to S$ tại đây, bỏ tên $g$; đồng nhất thức tiếp theo $\pi=f\circ p=g\circ q$ xác định cấu xạ cấu trúc thứ hai này là $g$.}} là các cấu xạ cấu trúc, $p$ và $q$ là các phép chiếu của $X\times_S Y$, và $\pi=f\circ p=g\circ q$ là cấu xạ cấu trúc $X\times_S Y\to S$.
Khi đó biểu đồ
\[
\label{I.5.3.5.1}
  \xymatrix{
    X\times_S Y\ar[r]^{(p,q)_T}\ar[d]_\pi &
    X\times_T Y\ar[d]^{f\times_T g}\\
    S\ar[r]_{\Delta_{S|T}} &
    S\times_T S
  }
  \tag{5.3.5.1}
\]
\oldpage[I]{133}
giao hoán và đồng nhất $X\times_S Y$ với tích của các $(S\times_T S)$-tiền lược đồ $S$ và $X\times_T Y$, trong đó các phép chiếu được đồng nhất với $\pi$ và $(p,q)_T$.
\end{proposition}

\begin{proof}
Theo \sref{I.3.4.3}, ta quy về việc chứng minh mệnh đề tương ứng \emph{trong phạm trù các tập hợp}, bằng cách thay $X$, $Y$ và $S$ lần lượt bởi $X(Z)_T$, $Y(Z)_T$ và $S(Z)_T$, trong đó $Z$ là một $T$-tiền lược đồ tùy ý.
Nhưng trong phạm trù các tập hợp, việc kiểm tra là hiển nhiên và được dành cho độc giả.
\end{proof}

\begin{corollary}[5.3.6]
\label{I.5.3.6}
Cấu xạ $(p,q)_T$ được đồng nhất (đặt $P=S\times_T S$) với $1_{X\times_T Y}\times_P\Delta_S$.
\end{corollary}

\begin{proof}
Điều này suy ra từ \sref{I.5.3.5} và \sref{I.3.3.4}.
\end{proof}

\begin{corollary}[5.3.7]
\label{I.5.3.7}
Nếu $f:X\to Y$ là một $S$-cấu xạ, thì sơ đồ
\[
  \xymatrix{
    X\ar[r]^{(1_X,f)_S}\ar[d]_f &
    X\times_S Y\ar[d]^{f\times_S 1_Y}\\
    Y\ar[r]_{\Delta_Y} &
    Y\times_S Y
  }
\]
giao hoán và đồng nhất $X$ với tích của các $(Y\times_S Y)$-tiền lược đồ $Y$ và $X\times_S Y$.
\end{corollary}

\begin{proof}
Chỉ cần áp dụng \sref{I.5.3.5}, thay $S$ bởi $Y$ và $T$ bởi $S$, đồng thời lưu ý rằng $X\times_Y Y=X$ \sref{I.3.3.3}.
\end{proof}

\begin{proposition}[5.3.8]
\label{I.5.3.8}
Điều kiện cần và đủ để $f:X\to Y$ là một đơn cấu của các tiền lược đồ là $\Delta_{X|Y}$ là một đẳng cấu từ $X$ lên $X\times_Y X$.
\end{proposition}

\begin{proof}
Thật vậy, $f$ là một đơn cấu khi và chỉ khi, với mọi $Y$-tiền lược đồ $Z$, ánh xạ tương ứng $f':X(Z)_Y\to Y(Z)_Y$ là đơn ánh.
Vì $Y(Z)_Y$ chỉ gồm một phần tử, điều này tương đương với việc $X(Z)_Y$ có nhiều nhất một phần tử\footnote{\emph{[Dịch.] Nguồn tiếng Pháp khẳng định rằng $X(Z)_Y$ cũng chỉ gồm đúng một phần tử; tính đơn ánh vào tập một phần tử $Y(Z)_Y$ chỉ kéo theo rằng miền có nhiều nhất một phần tử, kể cả trường hợp rỗng.}}, hay tương đương, ánh xạ đường chéo từ $X(Z)_Y$ vào $X(Z)_Y\times X(Z)_Y$ là song ánh.
Tích sau cùng này chính là tập $(X\times_Y X)(Z)_Y$ \hyperref[I.3.4.3]{(3.4.3.1)}, do đó điều này tương đương với việc $\Delta_{X|Y}$ là một đẳng cấu.
\end{proof}

\begin{proposition}[5.3.9]
\label{I.5.3.9}
Cấu xạ đường chéo $\Delta_X$ là một phép nhúng từ $X$ vào $X\times_S X$.
\end{proposition}

\begin{proof}
Chứng minh được in trong nguồn tiếng Pháp không đầy đủ, vì không chứng minh rằng $\Delta_X(X)$ đóng địa phương trong $X\times_S X$.\footnote{\emph{[Dịch.] EGA~III.2, Đính chính và bổ sung (danh sách~2), mục $\mathrm{Err}_{\mathrm{III}},10$, xác định rõ thiếu sót này và cung cấp chứng minh affine-địa phương được dịch tại đây.}}
Để có một chứng minh đúng, chỉ cần dùng \sref{I.4.2.4}[a]:
với mọi $x\in X$, hãy chọn một lân cận affine $W$ của ảnh của $x$ trong $S$, rồi chọn một lân cận affine $U$ của $x$ trong $X$ nằm trong ảnh ngược của $W$ qua cấu xạ cấu trúc;
$U\times_S U$ được đồng nhất với $U\times_W U$ và là một lân cận affine của $\Delta_X(x)$.\footnote{\emph{[Dịch.] Phát biểu được công bố trong EGA~III.2, Đính chính và bổ sung (danh sách~2), mục $\mathrm{Err}_{\mathrm{III}},10$, khẳng định điều này đối với mọi lân cận affine $U$ mà không đòi hỏi cấu xạ cấu trúc của nó phân tích qua một tập mở affine của $S$; giả thiết được nêu rõ ở đây là cần thiết cho bước này.}}
Xét đến \sref{I.5.3.16} (chứng minh của nó chỉ dùng Định nghĩa
\sref{I.5.3.1.1}), ta quy về việc chứng minh \sref{I.5.3.9} khi
$S=\Spec(B)$ và $X=\Spec(A)$ là các lược đồ affine.
Khi đó, từ \sref{I.5.3.1.1} suy ra ngay rằng $\Delta_X$ tương ứng với
đồng cấu chính tắc $A\otimes_B A\to A$ gửi $x\otimes y$ đến $xy$.
Vì đồng cấu này toàn ánh, trong trường hợp này $\Delta_X$ là một phép nhúng
đóng \sref{I.4.2.3}, và điều đó hoàn tất chứng minh.
\end{proof}

Ta gọi tiền lược đồ con của $X\times_S X$ gắn với phép nhúng $\Delta_X$ \sref{I.4.2.1} là \emph{đường chéo} của $X\times_S X$.

\begin{corollary}[5.3.10]
\label{I.5.3.10}
Dưới các giả thiết của \sref{I.5.3.5}, $(p,q)_T$ là một phép nhúng.
\end{corollary}

\begin{proof}
Điều này suy ra từ \sref{I.5.3.6} và \sref{I.4.3.1}.
\end{proof}

Ta nói (dưới các giả thiết của \sref{I.5.3.5}) rằng $(p,q)_T$ là \emph{phép nhúng chính tắc} từ $X\times_S Y$ vào $X\times_T Y$.

\begin{corollary}[5.3.11]
\label{I.5.3.11}
Cho $X$ và $Y$ là hai $S$-tiền lược đồ, và cho $f:X\to Y$ là một $S$-cấu xạ;
khi đó cấu xạ đồ thị $\Gamma_f=(1_X,f)_S$ của $f$ \sref{I.3.3.14} là một phép nhúng từ $X$ vào $X\times_S Y$.
\end{corollary}

\begin{proof}
Đây là trường hợp riêng của Hệ quả \sref{I.5.3.10}, trong đó thay $S$ bởi $Y$ và $T$ bởi $S$ (xem \sref{I.5.3.7}).
\end{proof}

\oldpage[I]{134}
Tiền lược đồ con của $X\times_S Y$ gắn với phép nhúng $\Gamma_f$ \sref{I.4.2.1} được gọi là \emph{đồ thị} của cấu xạ $f$;
các tiền lược đồ con của $X\times_S Y$ là đồ thị của các cấu xạ $X\to Y$ được đặc trưng bởi tính chất rằng, với một tiền lược đồ con như thế $G$, hạn chế lên đó của phép chiếu $p_1:X\times_S Y\to X$ là một \emph{đẳng cấu} $g$ từ $G$ lên $X$:
khi đó $G$ là đồ thị của cấu xạ $p_2\circ g^{-1}$, trong đó $p_2$ là phép chiếu $X\times_S Y\to Y$.

Đặc biệt, khi lấy $X=S$, các $S$-cấu xạ $S\to Y$ (không gì khác hơn các $S$-\emph{nhát cắt} của $Y$ \sref{I.2.5.5}) trùng với các cấu xạ đồ thị của chúng;
các tiền lược đồ con của $Y$ là đồ thị của các $S$-nhát cắt (nói cách khác, các tiền lược đồ con đẳng cấu với $S$ qua hạn chế của cấu xạ cấu trúc $Y\to S$) cũng được gọi là \emph{ảnh} của các nhát cắt này hoặc, do lạm dụng ngôn ngữ, là các $S$-\emph{nhát cắt} của $Y$.

\begin{corollary}[5.3.12]
\label{I.5.3.12}
Với các giả thiết và ký hiệu của \sref{I.5.3.11}, với mọi cấu xạ $g:S'\to S$, cho $f'$ là ảnh ngược của $f$ qua $g$ \sref{I.3.3.7};
khi đó $\Gamma_{f'}$ là ảnh ngược của $\Gamma_f$ qua $g$.
\end{corollary}

\begin{proof}
Đây là một trường hợp riêng của công thức \sref{I.3.3.10.1}.
\end{proof}

\begin{corollary}[5.3.13]
\label{I.5.3.13}
Cho $f:X\to Y$ và $g:Y\to Z$ là hai cấu xạ;
nếu $g\circ f$ là một phép nhúng (tương ứng, một phép nhúng địa phương), thì $f$ cũng vậy.
\end{corollary}

\begin{proof}
Thật vậy, $f$ phân tích thành $X\xrightarrow{\Gamma_f}X\times_Z Y\xrightarrow{p_2}Y$.
Mặt khác, $p_2$ được đồng nhất với $(g\circ f)\times_Z 1_Y$ \sref{I.3.3.4};
nếu $g\circ f$ là một phép nhúng (tương ứng, một phép nhúng địa phương), thì $p_2$ cũng vậy (\sref{I.4.3.1} và \sref{I.4.5.5}), và vì $\Gamma_f$ là một phép nhúng \sref{I.5.3.11}, ta kết luận theo \sref{I.4.2.5}\footnote{\emph{[Dịch.] Nguồn tiếng Pháp dẫn ở đây đến \sref{I.4.2.4}; EGA~III.2, Đính chính và bổ sung (danh sách~2), quy định rõ phải thay bằng \sref{I.4.2.5}.}} (tương ứng, \sref{I.4.5.5}).
\end{proof}

\begin{corollary}[5.3.14]
\label{I.5.3.14}
Cho $j:X\to Y$ và $g:X\to Z$ là hai $S$-cấu xạ.
Nếu $j$ là một phép nhúng (tương ứng, một phép nhúng địa phương), thì $(j,g)_S$ cũng vậy.
\end{corollary}

\begin{proof}
Thật vậy, nếu $p:Y\times_S Z\to Y$ là phép chiếu lên thừa số thứ nhất, thì $j=p\circ(j,g)_S$, và chỉ cần áp dụng \sref{I.5.3.13}.
\end{proof}

\begin{proposition}[5.3.15]
\label{I.5.3.15}
Nếu $f:X\to Y$ là một $S$-cấu xạ, thì sơ đồ
\[
\label{I.5.3.15.1}
  \xymatrix{
    X\ar[r]^{\Delta_X}\ar[d]_f &
    X\times_S X\ar[d]^{f\times_S f}\\
    Y\ar[r]_{\Delta_Y} &
    Y\times_S Y
  }\tag{5.3.15.1}
\]
là giao hoán \emph{(nói cách khác, $\Delta_X$ là một cấu xạ có tính hàm tử trong phạm trù các tiền lược đồ)}.
\end{proposition}

\begin{proof}
Việc kiểm tra là hiển nhiên và được dành cho độc giả.
\end{proof}

\begin{corollary}[5.3.16]
\label{I.5.3.16}
Nếu $X$ là một tiền lược đồ con của $Y$, thì đường chéo $\Delta_X(X)$ được đồng nhất với một tiền lược đồ con của $\Delta_Y(Y)$, và không gian nền của nó được đồng nhất với
\[
  \Delta_Y(Y)\cap p_1^{-1}(X)=\Delta_Y(Y)\cap p_2^{-1}(X)
\]
($p_1$ và $p_2$ là các phép chiếu của $Y\times_S Y$).
\end{corollary}

\begin{proof}
Áp dụng \sref{I.5.3.15} cho cấu xạ nhúng $f:X\to Y$, ta thấy $f\times_S f$ là một phép nhúng đồng nhất không gian nền của $X\times_S X$ với không gian con $p_1^{-1}(X)\cap p_2^{-1}(X)$ của $Y\times_S Y$ \sref{I.4.3.1};
hơn nữa, nếu $z\in\Delta_Y(Y)\cap p_1^{-1}(X)$, thì $z=\Delta_Y(y)$
\oldpage[I]{135}
và $y=p_1(z)\in X$, do đó $y=f(y)$, và $z=\Delta_Y(f(y))$ thuộc $\Delta_X(X)$ do tính giao hoán của \sref{I.5.3.15.1}.
\end{proof}

\begin{corollary}[5.3.17]
\label{I.5.3.17}
Cho $f_1:Y\to X$ và $f_2:Y\to X$ là hai $S$-cấu xạ, và cho $y$ là một điểm của $Y$ sao cho $f_1(y)=f_2(y)=x$ và các đồng cấu $\kres(x)\to\kres(y)$ tương ứng với $f_1$ và $f_2$ trùng nhau.
Khi đó, nếu $f=(f_1,f_2)_S$, điểm $f(y)$ thuộc đường chéo $\Delta_{X|S}(X)$.
\end{corollary}

\begin{proof}
Hai đồng cấu $\kres(x)\to\kres(y)$ tương ứng với $f_i$ ($i=1,2$) xác định hai $S$-cấu xạ $g_i:\Spec(\kres(y))\to\Spec(\kres(x))$ sao cho các biểu đồ
\[
  \xymatrix{
    \Spec(\kres(y))\ar[r]^{g_i}\ar[d] &
    \Spec(\kres(x))\ar[d]\\
    Y\ar[r]_{f_i} &
    X
  }
\]
là giao hoán.
Do đó, biểu đồ
\[
  \xymatrix{
    \Spec(\kres(y))\ar[rr]^{(g_1,g_2)_S}\ar[d] & &
    \Spec(\kres(x))\times_S\Spec(\kres(x))\ar[d]\\
    Y\ar[rr]_{(f_1,f_2)_S} & &
    X\times_S X
  }
\]
cũng giao hoán.
Nhưng từ đẳng thức $g_1=g_2$ suy ra rằng ảnh qua $(g_1,g_2)_S$ của điểm duy nhất của $\Spec(\kres(y))$ thuộc đường chéo của $\Spec(\kres(x))\times_S\Spec(\kres(x))$;
do đó kết luận suy ra từ \sref{I.5.3.15}.
\end{proof}

\subsection{Các cấu xạ tách và tiền lược đồ tách}
\label{subsection:I.5.4}
\label{I.5.4}

\begin{definition}[5.4.1]
\label{I.5.4.1}
Ta nói một cấu xạ của các tiền lược đồ $f:X\to Y$ là \emph{tách} nếu cấu xạ đường chéo $X\to X\times_Y X$ là một phép nhúng \emph{đóng}; khi đó ta cũng nói $X$ là một \emph{tiền lược đồ tách trên $Y$}, hoặc một \emph{$Y$-lược đồ}.
Ta nói một tiền lược đồ $X$ là tách nếu nó tách trên $\Spec(\bb{Z})$;
khi đó ta cũng gọi $X$ là một \emph{lược đồ}\footnote{\emph{[Dịch.] Ở đây chúng tôi nhắc lại lưu ý đã nêu ngay ở đầu bản dịch này: các phiên bản đầu của EGA dùng \emph{tiền lược đồ} cho đối tượng ngày nay thường được gọi là lược đồ, và dùng \emph{lược đồ} cho đối tượng ngày nay thường được gọi là lược đồ tách. Về sau chính Grothendieck nói rằng thuật ngữ hiện đại là thích hợp hơn, nhưng chúng tôi quyết định duy trì bản dịch «trung thành về mặt lịch sử» này bằng cách dùng danh pháp cũ.}} \emph{(xem \sref{I.5.5.7})}.
\end{definition}

Theo \sref{I.5.3.9}, điều kiện cần và đủ để $X$ tách trên $Y$ là $\Delta_X(X)$ là một \emph{không gian con đóng} của không gian nền của $X\times_Y X$.

\begin{proposition}[5.4.2]
\label{I.5.4.2}
Cho $S\to T$ là một cấu xạ tách.
Nếu $X$ và $Y$ là hai $S$-tiền lược đồ, thì phép nhúng chính tắc $X\times_S Y\to X\times_T Y$ \sref{I.5.3.10} là đóng.
\end{proposition}

\begin{proof}
Thật vậy, khi xét biểu đồ ở \sref{I.5.3.5.1}, ta thấy $(p,q)_T$ có thể được coi là thu được từ $\Delta_{S|T}$ bằng phép mở rộng $f\times_T g:X\times_T Y\to S\times_T S$ của tiền lược đồ cơ sở;
khi đó mệnh đề suy ra từ \sref{I.4.3.2}.
\end{proof}

\begin{corollary}[5.4.3]
\label{I.5.4.3}
Cho $Y$ là một $S$-lược đồ, và cho $f:X\to Y$ là một $S$-cấu xạ.
Khi đó cấu xạ đồ thị $\Gamma_f:X\to X\times_S Y$ \sref{I.5.3.11} là một phép nhúng đóng.
\end{corollary}

\begin{proof}
Đây là trường hợp riêng của \sref{I.5.4.2}, trong đó thay $S$ bởi $Y$ và $T$ bởi $S$.
\end{proof}

\begin{corollary}[5.4.4]
\label{I.5.4.4}
Cho $f:X\to Y$ và $g:Y\to Z$ là hai cấu xạ, với $g$ tách.
Nếu $g\circ f$ là một phép nhúng đóng, thì $f$ cũng vậy.
\end{corollary}

\begin{proof}
Chứng minh từ \sref{I.5.4.3} giống với chứng minh của \sref{I.5.3.13} từ \sref{I.5.3.11}.
\end{proof}

\oldpage[I]{136}
\begin{corollary}[5.4.5]
\label{I.5.4.5}
Cho $Z$ là một $S$-lược đồ, và cho $j:X\to Y$ và $g:X\to Z$ là hai $S$-cấu xạ.
Nếu $j$ là một phép nhúng đóng, thì $(j,g)_S:X\to Y\times_S Z$ cũng vậy.
\end{corollary}

\begin{proof}
Chứng minh từ \sref{I.5.4.4} giống với chứng minh của \sref{I.5.3.14} từ \sref{I.5.3.13}.
\end{proof}

\begin{corollary}[5.4.6]
\label{I.5.4.6}
Nếu $X$ là một $S$-lược đồ, thì mọi $S$-nhát cắt của $X$ \sref{I.2.5.5} là một phép nhúng đóng.
\end{corollary}

\begin{proof}
Nếu $\vphi:X\to S$ là cấu xạ cấu trúc, và $\psi:S\to X$ là một $S$-nhát cắt của $X$, thì chỉ cần áp dụng \sref{I.5.4.5} cho $\vphi\circ\psi=1_S$.
\end{proof}

\begin{corollary}[5.4.7]
\label{I.5.4.7}
Cho $S$ là một tiền lược đồ nguyên, $s$ là điểm tổng quát của nó, và cho $X$ là một $S$-lược đồ.
Nếu hai $S$-nhát cắt $f$ và $g$ của $X$ thỏa mãn $f(s)=g(s)$, thì $f=g$.
\end{corollary}

\begin{proof}
Thật vậy, nếu $x=f(s)=g(s)$, thì các đồng cấu $\kres(x)\to\kres(s)$ tương ứng với $f$ và $g$ nhất thiết trùng nhau.
Nếu $h=(f,g)_S$, thì từ \sref{I.5.3.17} suy ra $h(s)$ thuộc đường chéo $Z=\Delta_X(X)$;
nhưng vì $S=\overline{\{s\}}$ và $Z$ là đóng theo giả thiết, ta có $h(S)\subset Z$.
Khi đó, từ \sref{I.5.2.2} suy ra $h$ phân tích thành $S\to Z\to X\times_S X$, và do đó, theo định nghĩa của đường chéo, ta có $f=g$.
\end{proof}

\begin{remark}[5.4.8]
\label{I.5.4.8}
Nếu ngược lại giả sử kết luận của \sref{I.5.4.3} đúng khi $f=1_Y$, thì suy ra $Y$ tách trên $S$;
tương tự, nếu giả sử kết luận của \sref{I.5.4.4}\footnote{\emph{[Dịch.] Nguồn tiếng Pháp dẫn ở đây đến (5.4.5), vốn nói về một cấu xạ được xác định bởi một cặp; lập luận với $p_1\circ\Delta_Y=1_Y$ sử dụng tiêu chuẩn về hợp thành của Hệ quả~\hyperref[I.5.4.4]{5.4.4}.}} áp dụng cho hai cấu xạ $Y\xrightarrow{\Delta_Y}Y\times_Z Y\xrightarrow{p_1}Y$, thì suy ra $\Delta_Y$ là một phép nhúng đóng và do đó $Y$ tách trên $Z$;
cuối cùng, nếu giả sử kết luận của \sref{I.5.4.6} đúng đối với $Y$-nhát cắt $\Delta_Y$ của $Y$-tiền lược đồ $Y\times_S Y\to Y$, thì điều này kéo theo $Y$ tách trên $S$.
\end{remark}

\subsection{Các tiêu chuẩn về tính tách}
\label{subsection:I.5.5}

\begin{proposition}[5.5.1]
\label{I.5.5.1}
\begin{enumerate}
  \item[\rm(i)] Mọi đơn cấu của các tiền lược đồ (và đặc biệt, mọi phép
    nhúng) đều là một cấu xạ tách.
  \item[\rm(ii)] Hợp thành của hai cấu xạ tách là tách.
  \item[\rm(iii)] Nếu $f:X\to X'$ và $g:Y\to Y'$ là hai $S$-cấu xạ tách,
    thì $f\times_S g$ là tách.
  \item[\rm(iv)] Nếu $f:X\to Y$ là một $S$-cấu xạ tách, thì $S'$-cấu xạ
    $f_{(S')}$ là tách đối với mọi mở rộng $S'\to S$ của tiền lược đồ cơ sở.
  \item[\rm(v)] Nếu hợp thành $g\circ f$ của hai cấu xạ là tách, thì $f$ là
    tách.
  \item[\rm(vi)] Điều kiện cần và đủ để một cấu xạ $f$ là tách là
    $f_\red$ (\sref{I.5.1.5}) tách.
\end{enumerate}
\end{proposition}

\begin{proof}
Khẳng định~(i) suy ra ngay từ \sref{I.5.3.8}.
Nếu $f:X\to Y$ và $g:Y\to Z$ là hai cấu xạ, thì biểu đồ
\[
\label{I.5.5.1.1}
  \xymatrix{
    X\ar[rr]^{\Delta_{X|Z}}\ar[dr]_{\Delta_{X|Y}} & &
    X\times_Z X\\
    & X\times_Y X\ar[ur]_j
  }
  \tag{5.5.1.1}
\]
là giao hoán, trong đó $j$ ký hiệu phép nhúng chính tắc \sref{I.5.3.10}, như
ta kiểm tra ngay.
Nếu $f$ và $g$ tách, thì $\Delta_{X|Y}$ là một phép nhúng đóng theo định
nghĩa, còn $j$ là một phép nhúng đóng theo \sref{I.5.4.2}; do đó
$\Delta_{X|Z}$ là một phép nhúng đóng theo \sref{I.4.2.5}\footnote{\emph{[Dịch.]
Nguồn tiếng Pháp dẫn ở đây đến (4.2.4), vốn là tiêu chuẩn địa phương; tính bắc
cầu của các phép nhúng đóng được dùng trong bước này là Mệnh đề~\hyperref[I.4.2.5]{4.2.5}.}}, điều này
\oldpage[I]{137}
chứng minh~(ii).
Theo~(i) và~(ii), (iii) và~(iv) là tương đương \sref{I.3.5.1}, nên chỉ cần
chứng minh~(iv).
Thật vậy, $X_{(S')}\times_{Y_{(S')}}X_{(S')}$ được đồng nhất chính tắc với
$(X\times_Y X)\times_Y Y_{(S')}$ theo \sref{I.3.3.11} và
\hyperref[I.3.3.9.1]{(3.3.9.1)}, và ta kiểm tra ngay rằng cấu xạ đường chéo
$\Delta_{X_{(S')}}$ khi đó được đồng nhất với
$\Delta_X\times_Y 1_{Y_{(S')}}$; vì vậy mệnh đề suy ra từ \sref{I.4.3.1}.

Để chứng minh~(v), như trong \sref{I.5.3.13}, xét phân tích
$X\xrightarrow{\Gamma_f}X\times_Z Y\xrightarrow{p_2}Y$ của $f$, và lưu ý
rằng $p_2=(g\circ f)\times_Z 1_Y$; giả thiết rằng $g\circ f$ là tách kéo theo
$p_2$ là tách theo (iii) và (i), và vì $\Gamma_f$ là một phép nhúng nên
$\Gamma_f$ là tách theo (i); do đó $f$ là tách theo (ii).
Cuối cùng, để chứng minh~(vi), nhắc lại rằng các tiền lược đồ
$X_\red\times_{Y_\red}X_\red$ và $X_\red\times_Y X_\red$ được đồng nhất
chính tắc với nhau \sref{I.5.1.7}; nếu ký hiệu bởi $j$ phép nhúng
$X_\red\to X$, thì biểu đồ
\[
  \xymatrix{
    X_\red\ar[r]^{\Delta_{X_\red}}\ar[d]_j &
    X_\red\times_Y X_\red\ar[d]^{j\times_Y j}\\
    X\ar[r]_{\Delta_X} & X\times_Y X
  }
\]
là giao hoán \sref{I.5.3.15}, và mệnh đề suy ra từ việc các mũi tên thẳng
đứng là những phép đồng phôi của các không gian nền \sref{I.4.3.1}.
\end{proof}

\begin{corollary}[5.5.2]
\label{I.5.5.2}
Nếu $f:X\to Y$ là tách, thì hạn chế của $f$ trên mọi tiền lược đồ con của
$X$ là tách.
\end{corollary}

\begin{proof}
Điều này suy ra từ \sref{I.5.5.1}[i và ii].
\end{proof}

\begin{corollary}[5.5.3]
\label{I.5.5.3}
Nếu $X$ và $Y$ là hai $S$-tiền lược đồ sao cho $Y$ tách trên $S$, thì
$X\times_S Y$ tách trên $X$.
\end{corollary}

\begin{proof}
Đây là một trường hợp riêng của \sref{I.5.5.1}[iv].
\end{proof}

\begin{proposition}[5.5.4]
\label{I.5.5.4}
Cho $X$ là một tiền lược đồ, và giả sử không gian nền của nó là hợp của một họ
hữu hạn các tập đóng $X_k$ ($1\leq k\leq n$); với mỗi $k$, xét tiền lược đồ con
rút gọn của $X$ có không gian nền là $X_k$ \sref{I.5.2.1}, và vẫn ký hiệu nó
bởi $X_k$.
Cho $f:X\to Y$ là một cấu xạ, và với mỗi $k$, cho $Y_k$ là một tập đóng của
$Y$ sao cho $f(X_k)\subset Y_k$; vẫn ký hiệu bởi $Y_k$ tiền lược đồ con rút
gọn của $Y$ có không gian nền là $Y_k$, để hạn chế $X_k\to Y$ của $f$ trên
$X_k$ phân tích thành $X_k\xrightarrow{f_k}Y_k\to Y$ \sref{I.5.2.2}.
Điều kiện cần và đủ để $f$ là tách là tất cả các $f_k$ đều tách.
\end{proposition}

\begin{proof}
Tính cần thiết suy ra từ \sref{I.5.5.1}[i, ii và v].
Ngược lại, nếu điều kiện của phát biểu được thỏa mãn, thì mỗi hạn chế
$X_k\to Y$ của $f$ là tách \hyperref[I.5.5.1]{(5.5.1, (i) và (ii))}; nếu
$p_1$ và $p_2$ là các phép chiếu của $X\times_Y X$, thì không gian con
$\Delta_{X_k}(X_k)$ được đồng nhất với không gian con
$\Delta_X(X)\cap p_1^{-1}(X_k)$ của không gian nền của $X\times_Y X$
\sref{I.5.3.16}; vì các không gian con này đóng trong $X\times_Y X$, hợp của
chúng là $\Delta_X(X)$ cũng đóng.
\end{proof}

Đặc biệt, giả sử các $X_k$ là những \emph{thành phần bất khả quy} của $X$;
khi đó ta có thể giả sử các $Y_k$ là những thành phần bất khả quy của $Y$
\sref[0]{0.2.1.5}; do đó, trong trường hợp này, Mệnh đề~\sref{I.5.5.4} quy
khái niệm tách về trường hợp các tiền lược đồ \emph{nguyên}
\sref{I.2.1.8}\footnote{\emph{[Dịch.] Nguồn tiếng Pháp dẫn ở đây đến
(2.1.7), mục nói về hạn chế lên các tập con mở; định nghĩa thích hợp của tiền
lược đồ nguyên là (2.1.8).}}.

\oldpage[I]{138}

\begin{proposition}[5.5.5]
\label{I.5.5.5}
Cho $(Y_\lambda)$ là một phủ mở của một tiền lược đồ $Y$; điều kiện cần và đủ
để một cấu xạ $f:X\to Y$ là tách là mỗi hạn chế
$f^{-1}(Y_\lambda)\to Y_\lambda$ của nó đều tách.
\end{proposition}

\begin{proof}
Đặt $X_\lambda=f^{-1}(Y_\lambda)$, xét đến \sref{I.4.2.4}[b] và sự đồng nhất
của các tích $X_\lambda\times_Y X_\lambda$ và
$X_\lambda\times_{Y_\lambda}X_\lambda$ \sref{I.3.2.5}, mọi việc quy về chứng
minh rằng các $X_\lambda\times_Y X_\lambda$ tạo thành một phủ của
$X\times_Y X$.
Thật vậy, nếu đặt $Y_{\lambda\mu}=Y_\lambda\cap Y_\mu$ và
$X_{\lambda\mu}=X_\lambda\cap X_\mu=f^{-1}(Y_{\lambda\mu})$, thì
$X_\lambda\times_Y X_\mu$ được đồng nhất với tích
$X_{\lambda\mu}\times_{Y_{\lambda\mu}}X_{\lambda\mu}$
\sref{I.3.2.6.4}, do đó cũng với
$X_{\lambda\mu}\times_Y X_{\lambda\mu}$ \sref{I.3.2.5}, và cuối cùng với
một tập mở của $X_\lambda\times_Y X_\lambda$; điều này chứng minh khẳng định
của ta \sref{I.3.2.7}.
\end{proof}

Mệnh đề~\sref{I.5.5.5}\footnote{\emph{[Dịch.] Nguồn tiếng Pháp dẫn ở đây
đến (5.5.4); phép quy bằng một phủ mở của $Y$ sử dụng
Mệnh đề~\hyperref[I.5.5.5]{5.5.5}.}} cho phép, bằng cách lấy một phủ của $Y$
bởi các tập mở affine, quy việc nghiên cứu các cấu xạ tách về trường hợp các
cấu xạ tách có đích là các lược đồ affine.

\begin{proposition}[5.5.6]
\label{I.5.5.6}
Cho $Y$ là một lược đồ \emph{affine}, $X$ là một tiền lược đồ, và
$(U_\alpha)$ là một phủ của $X$ bởi các tập mở affine.
Điều kiện cần và đủ để một cấu xạ $f:X\to Y$ là tách là
$U_\alpha\cap U_\beta$ phải là một tập mở affine đối với mọi cặp chỉ số
$(\alpha,\beta)$, \emph{và} vành
$\Gamma(U_\alpha\cap U_\beta,\sh{O}_X)$ được sinh bởi hợp các ảnh chính tắc
của các vành $\Gamma(U_\alpha,\sh{O}_X)$ và
$\Gamma(U_\beta,\sh{O}_X)$.
\end{proposition}

\begin{proof}
Các $U_\alpha\times_Y U_\beta$ tạo thành một phủ mở của $X\times_Y X$
\sref{I.3.2.7}; gọi các phép chiếu của $X\times_Y X$ là $p$ và $q$, ta có
\begin{align*}
  \Delta_X^{-1}(U_\alpha\times_Y U_\beta)&=\Delta_X^{-1}(p^{-1}(U_\alpha)\cap q^{-1}(U_\beta))\\
                                         &=\Delta_X^{-1}(p^{-1}(U_\alpha))\cap\Delta_X^{-1}(q^{-1}(U_\beta))=U_\alpha\cap U_\beta;
\end{align*}
do đó mọi việc quy về chứng minh rằng hạn chế của $\Delta_X$ trên
$U_\alpha\cap U_\beta$ là một phép nhúng đóng vào
$U_\alpha\times_Y U_\beta$.
Hạn chế này không gì khác hơn $(j_\alpha,j_\beta)_Y$, trong đó $j_\alpha$
(tương ứng $j_\beta$) ký hiệu cấu xạ nhúng
$U_\alpha\cap U_\beta$ vào $U_\alpha$ (tương ứng $U_\beta$), như suy ra từ
các định nghĩa.
Vì $U_\alpha\times_Y U_\beta$ là một lược đồ affine có vành đẳng cấu chính
tắc với
$\Gamma(U_\alpha,\sh{O}_X)\otimes_{\Gamma(Y,\sh{O}_Y)}\Gamma(U_\beta,\sh{O}_X)$
\sref{I.3.2.2}, ta thấy rằng $U_\alpha\cap U_\beta$ phải là một lược đồ affine
và ánh xạ $h_\alpha\otimes h_\beta\mapsto h_\alpha h_\beta$ từ vành
$A(U_\alpha\times_Y U_\beta)$ vào
$\Gamma(U_\alpha\cap U_\beta,\sh{O}_X)$ phải toàn ánh
\sref{I.4.2.3}; điều này kết thúc chứng minh.
\end{proof}

\begin{corollary}[5.5.7]
\label{I.5.5.7}
Một lược đồ affine là tách \emph{(và do đó là một \emph{lược đồ}, điều này
biện minh cho thuật ngữ của \sref{I.5.4.1})}.
\end{corollary}

\begin{corollary}[5.5.8]
\label{I.5.5.8}
Cho $Y$ là một lược đồ affine; điều kiện cần và đủ để $f:X\to Y$ là một cấu xạ
tách là $X$ tách \emph{(nói cách khác, $X$ là một \emph{lược đồ})}.
\end{corollary}

\begin{proof}
Thật vậy, tiêu chuẩn của \sref{I.5.5.6} không phụ thuộc vào $f$.
\end{proof}

\begin{corollary}[5.5.9]
\label{I.5.5.9}
Để một cấu xạ $f:X\to Y$ là tách, điều kiện cần là tiền lược đồ cảm sinh
$f^{-1}(U)$ tách đối với mọi tập mở $U$ mà trên đó $Y$ cảm sinh một tiền lược
đồ tách, và điều kiện đủ là điều này đúng đối với mọi tập mở affine
$U\subset Y$.
\end{corollary}

\begin{proof}
Tính cần thiết của điều kiện suy ra từ \sref{I.5.5.5}\footnote{\emph{[Dịch.]
Trong hai tham chiếu đến (5.5.5) của chứng minh này, nguồn tiếng Pháp in
(5.5.4); lập luận áp dụng tính địa phương trên cơ sở được thiết lập trong
Mệnh đề~\hyperref[I.5.5.5]{5.5.5}.}} và \sref{I.5.5.1}[ii]; tính đủ suy
ra từ \sref{I.5.5.5} và \sref{I.5.5.8}, xét đến sự tồn tại của các phủ $Y$
bởi những tập mở affine.
\end{proof}

Đặc biệt, nếu $X$ và $Y$ là các lược đồ affine, thì \emph{mọi} cấu xạ
$X\to Y$ đều tách.

\begin{proposition}[5.5.10]
\label{I.5.5.10}
Cho $Y$ là một \emph{lược đồ}, và cho $f:X\to Y$ là một cấu xạ.
Với mọi tập mở affine $U$ của $X$ và mọi tập mở affine $V$ của $Y$,
$U\cap f^{-1}(V)$ là affine.
\end{proposition}

\begin{proof}
Gọi $p_1$ và $p_2$ là các phép chiếu của $X\times_\bb{Z}Y$;
không gian con $U\cap f^{-1}(V)$ là ảnh của
$\Gamma_f(X)\cap p_1^{-1}(U)\cap p_2^{-1}(V)$ qua $p_1$.
Nhưng $p_1^{-1}(U)\cap p_2^{-1}(V)$ được đồng nhất với không gian nền của
\oldpage[I]{139}
tiền lược đồ $U\times_\bb{Z}V$ \sref{I.3.2.7}, và do đó là một lược đồ affine
\sref{I.3.2.2}; vì $\Gamma_f(X)$ đóng trong $X\times_\bb{Z}Y$
\sref{I.5.4.3}, nên $\Gamma_f(X)\cap p_1^{-1}(U)\cap p_2^{-1}(V)$ đóng trong
$U\times_\bb{Z} V$, do đó tiền lược đồ do tiền lược đồ con của
$X\times_\bb{Z}Y$ gắn với $\Gamma_f$ \sref{I.4.2.1} cảm sinh trên tập mở
$\Gamma_f(X)\cap p_1^{-1}(U)\cap p_2^{-1}(V)$ của không gian nền của nó là
một tiền lược đồ con đóng của một lược đồ affine, và vì thế là một lược đồ
affine \sref{I.4.2.3}.
Mệnh đề khi đó suy ra từ việc $\Gamma_f$ là một phép nhúng.
\end{proof}

\begin{examples}[5.5.11]
\label{I.5.5.11}
Tiền lược đồ trong Ví dụ~\sref{I.2.3.2} (``đường thẳng xạ ảnh trên một trường
$K$'') là \emph{tách}, vì đối với phủ $(X_1,X_2)$ của $X$ bởi các tập mở
affine, $X_1\cap X_2=U_{12}$ là affine và
$\Gamma(U_{12},\sh{O}_X)$, vành các phân thức hữu tỉ dạng $f(s)/s^m$, với
$f\in K[s]$, được sinh bởi $K[s]$ và $1/s$, nên các điều kiện của
\sref{I.5.5.6} được thỏa mãn.

Với cùng lựa chọn $X_1$, $X_2$, $U_{12}$ và $U_{21}$ như trong
Ví dụ~\sref{I.2.3.2}, lần này ta lấy $u_{12}$ là đẳng cấu gửi $f(s)$ tới
$f(t)$;
khi đó phép dán cho ta một tiền lược đồ \emph{nguyên không tách} $X$, vì điều
kiện thứ nhất của \sref{I.5.5.6} được thỏa mãn, nhưng điều kiện thứ hai thì
không.
Ở đây hiển nhiên
$\Gamma(X,\sh{O}_X)\to\Gamma(X_1,\sh{O}_X)=K[s]$ là một đẳng cấu;
đẳng cấu ngược xác định một cấu xạ toàn ánh $f:X\to\Spec(K[s])$, và với mọi
$y\in\Spec(K[s])$ sao cho $\mathfrak{j}_y\neq(s)$, $f^{-1}(y)$ chỉ gồm một
điểm, còn khi $\mathfrak{j}_y=(s)$ thì $f^{-1}(y)$ gồm \emph{hai} điểm phân
biệt\footnote{\emph{[Dịch.] Nguồn tiếng Pháp in $(0)$ trong cả hai điều kiện
về iđêan. Đối với gốc tọa độ bị nhân đôi của $\Spec(K[s])$, iđêan nguyên tố
thích hợp là $(s)$; $(0)$ là điểm tổng quát.}} (ta nói rằng $X$ là ``đường
thẳng affine trên $K$ với điểm $0$ được nhân đôi'').

Ta cũng có thể cho một ví dụ trong đó điều kiện thứ nhất trong hai điều kiện
của \sref{I.5.5.6} không được thỏa mãn.
Trước hết, hãy nhận xét rằng trong phổ nguyên tố $Y$ của vành đa thức
$A=K[s,t]$ theo hai biến trên một trường $K$, tập mở $U$ là hợp của $D(s)$ và
$D(t)$ \emph{không phải là một tập mở affine}.
Thật vậy, nếu $z$ là một tiết diện của $\sh{O}_Y$ trên $U$, thì tồn tại hai số
nguyên $m,n\geq0$ sao cho $s^mz$ và $t^nz$ là các hạn chế của những đa thức
theo $s$ và $t$ lên $U$ \sref{I.1.4.1}; điều này hiển nhiên chỉ có thể xảy ra
nếu tiết diện $z$ kéo dài thành một tiết diện trên toàn bộ $Y$, được đồng nhất
với một đa thức theo $s$ và $t$.
Nếu $U$ là một tập mở affine, thì cấu xạ nhúng $U\to Y$ sẽ là một đẳng cấu
\sref{I.1.7.3}, điều này vô lý vì $U\neq Y$.

Theo nhận xét trên, lấy hai lược đồ affine $Y_1$ và $Y_2$, là các phổ nguyên
tố của các vành $A_1=K[s_1,t_1]$ và $A_2=K[s_2,t_2]$ tương ứng;
lấy $U_{12}=D(s_1)\cup D(t_1)$ và $U_{21}=D(s_2)\cup D(t_2)$, và lấy
$u_{12}$ là hạn chế của một đẳng cấu $Y_2\to Y_1$ lên $U_{21}$ tương ứng với
đẳng cấu vành gửi $f(s_1,t_1)$ tới $f(s_2,t_2)$;
khi đó ta thu được một ví dụ trong đó điều kiện thứ nhất của \sref{I.5.5.6}
không được thỏa mãn, còn điều kiện thứ hai thì được thỏa
mãn\footnote{\emph{[Dịch.] Nguồn tiếng Pháp khẳng định rằng không điều kiện
nào trong hai điều kiện của \sref{I.5.5.6} được thỏa mãn. Tuy nhiên, với
$U=D(s)\cup D(t)$, ta có $\Gamma(U,\sh{O}_Y)=K[s,t]$, nên điều kiện thứ hai
được thỏa mãn; chỉ có tính affine của $U$ là không thỏa mãn.}} (tiền lược đồ nguyên thu
được như vậy được gọi là ``mặt phẳng affine trên $K$ với điểm $0$ được nhân
đôi'').
\end{examples}

\begin{remark}[5.5.12]
\label{I.5.5.12}
Cho một tính chất \textbf{P} của các cấu xạ tiền lược đồ, xét các mệnh đề
sau:
\begin{enumerate}
  \item[\rm(i)] \emph{Mọi phép nhúng đóng đều có tính chất~\textbf{P}.}
  \item[\rm(ii)] \emph{Hợp thành của hai cấu xạ có tính chất~\textbf{P}
    cũng có tính chất~\textbf{P}.}
  \item[\rm(iii)] \emph{Nếu $f:X\to X'$ và $g:Y\to Y'$ là hai
    $S$-cấu xạ có tính chất~\textbf{P}, thì $f\times_S g$ có tính
    chất~\textbf{P}.}

\oldpage[I]{140}
  \item[\rm(iv)] \emph{Nếu $f:X\to Y$ là một $S$-cấu xạ có tính
    chất~\textbf{P}, thì mọi $S'$-cấu xạ $f_{(S')}$ thu được bằng một mở
    rộng $S'\to S$ của tiền lược đồ cơ sở cũng có tính chất~\textbf{P}.}
  \item[\rm(v)] \emph{Nếu hợp thành $g\circ f$ của hai cấu xạ
    $f:X\to Y$ và $g:Y\to Z$ có tính chất~\textbf{P}, và nếu $g$ tách,
    thì $f$ có tính chất~\textbf{P}.}
  \item[\rm(vi)] \emph{Nếu một cấu xạ $f:X\to Y$ có tính
    chất~\textbf{P}, thì $f_\red$ cũng có tính chất ấy
    (\sref{I.5.1.5}).}
\end{enumerate}

Trong các điều kiện đó, nếu giả sử (i) và (ii) đúng, thì (iii) và (iv) là
\emph{tương đương}, còn (v) và (vi) là các \emph{hệ quả} của (i), (ii) và
(iii).

Khẳng định thứ nhất đã được chứng minh (\sref{I.3.5.1}).
Xét phân tích (\sref{I.5.3.13})
$X\xrightarrow{\Gamma_f}X\times_Z Y\xrightarrow{p_2}Y$ của $f$; quan hệ
$p_2=(g\circ f)\times_Z 1_Y$ cho thấy rằng nếu $g\circ f$ có tính
chất~\textbf{P} thì $p_2$ cũng có tính chất ấy theo (iii); nếu $g$ tách thì
$\Gamma_f$ là một phép nhúng đóng (\sref{I.5.4.3}), nên cũng có tính
chất~\textbf{P} theo (i); cuối cùng, theo (ii), $f$ có tính
chất~\textbf{P}.

Cuối cùng, xét biểu đồ giao hoán
\[
  \xymatrix{
    X_\red\ar[r]^{f_\red}\ar[d] &
    Y_\red\ar[d]\\
    X\ar[r]_f &
    Y,
  }
\]
trong đó các mũi tên thẳng đứng là những phép nhúng đóng
(\sref{I.5.1.5}), nên có tính chất~\textbf{P} theo (i). Giả thiết rằng $f$
có tính chất~\textbf{P} kéo theo, theo (ii), rằng
$X_\red\xrightarrow{f_\red}Y_\red\to Y$ có tính chất~\textbf{P}; cuối
cùng, vì một phép nhúng đóng là tách (\sref{I.5.5.1}[i]), nên $f_\red$ có
tính chất~\textbf{P} theo (v).

Ta lưu ý rằng nếu xét các mệnh đề:
\begin{enumerate}
  \item[\rm(i')] \emph{Mọi phép nhúng đều có tính chất~\textbf{P}.}
  \item[\rm(v')] \emph{Nếu $g\circ f$ có tính chất~\textbf{P}, thì $f$
    cũng có tính chất ấy;}
\end{enumerate}
thì lập luận trên cho thấy (v') là hệ quả của (i'), (ii) và (iii).
\end{remark}

\begin{env}[5.5.13]
\label{I.5.5.13}
Ta lưu ý rằng (v) và (vi) vẫn là các hệ quả của (i), (iii) và
\begin{enumerate}
  \item[\rm(ii')] \emph{Nếu $j:X\to Y$ là một phép nhúng đóng và
    $g:Y\to Z$ là một cấu xạ có tính chất~\textbf{P}, thì $g\circ j$ có
    tính chất~\textbf{P}.}
\end{enumerate}
Tương tự, (v') là hệ quả của (i'), (iii) và
\begin{enumerate}
  \item[\rm(ii'')] \emph{Nếu $j:X\to Y$ là một phép nhúng và $g:Y\to Z$
    là một cấu xạ có tính chất~\textbf{P}, thì $g\circ j$ có tính
    chất~\textbf{P}.}
\end{enumerate}
Điều này suy ra ngay từ các lập luận trong \sref{I.5.5.12}.
\end{env}

\section{Các điều kiện hữu hạn}
\label{section:I.6}

\subsection{Tiền lược đồ Noether và Noether địa phương}
\label{subsection:I.6.1}

\begin{definition}[6.1.1]
\label{I.6.1.1}
Ta nói một tiền lược đồ $X$ là \emph{Noether} (tương ứng,
\emph{Noether địa phương}) nếu nó là hợp hữu hạn (tương ứng, hợp) của các tập
mở affine $V_\alpha$ sao cho vành của lược đồ cảm sinh trên mỗi $V_\alpha$ là
Noether.
\end{definition}

Từ \sref{I.1.5.2} suy ra ngay rằng, nếu $X$ là Noether địa phương, thì bó
cấu trúc $\sh{O}_X$ là một \emph{bó coherent các vành}, vì vấn đề có tính
địa phương. Mọi \emph{$\sh{O}_X$-môđun con quasi-coherent}
\oldpage[I]{141}
(tương ứng, mọi $\sh{O}_X$-môđun thương quasi-coherent) của một
$\sh{O}_X$-môđun \emph{coherent} $\sh{F}$ khi đó đều \emph{coherent}, vì
vấn đề lại có tính địa phương, và chỉ cần áp dụng \sref{I.1.5.1},
\sref{I.1.4.1} và \sref{I.1.3.10}, cùng với sự kiện rằng một môđun con
(tương ứng, môđun thương) của một môđun hữu hạn sinh trên một vành Noether
là hữu hạn sinh. Đặc biệt, mọi \emph{bó iđêan quasi-coherent} của
$\sh{O}_X$ đều \emph{coherent}.

Nếu một tiền lược đồ $X$ là hợp hữu hạn (tương ứng, hợp) của các tập mở
$W_\lambda$ sao cho các tiền lược đồ cảm sinh trên các $W_\lambda$ là
Noether (tương ứng, Noether địa phương), thì hiển nhiên $X$ là Noether
(tương ứng, Noether địa phương).

\begin{proposition}[6.1.2]
\label{I.6.1.2}
Để một tiền lược đồ $X$ là Noether, điều kiện cần và đủ là nó Noether địa
phương và không gian nền của nó quasi-compact. Khi đó, chính không gian nền
cũng là Noether.
\end{proposition}

\begin{proof}
Khẳng định thứ nhất suy ra ngay từ các định nghĩa và \sref{I.1.1.10}[ii].
Khẳng định thứ hai suy ra từ \sref{I.1.1.6} và sự kiện rằng mọi không gian
là hợp hữu hạn của các không gian con Noether thì cũng là Noether
\sref[0]{0.2.2.3}.
\end{proof}

\begin{proposition}[6.1.3]
\label{I.6.1.3}
Cho $X$ là một lược đồ affine có vành $A$.
Các điều kiện sau là tương đương:
\emph{(a)} $X$ là Noether;
\emph{(b)} $X$ là Noether địa phương;
\emph{(c)} $A$ là Noether.
\end{proposition}

\begin{proof}
Tính tương đương của (a) và (b) suy ra từ \sref{I.6.1.2} và sự kiện rằng
mọi lược đồ affine có không gian nền quasi-compact \sref{I.1.1.10}; ngoài
ra, hiển nhiên (c) suy ra (a). Để thấy (a) suy ra (c), hãy nhận xét rằng có
một phủ hữu hạn $(V_i)$ của $X$ bởi các tập mở affine sao cho vành $A_i$ của
tiền lược đồ cảm sinh trên $V_i$ là Noether. Bây giờ, cho
$(\mathfrak{a}_n)$ là một dãy tăng các iđêan của $A$; qua tương ứng song ánh
chính tắc \sref{I.1.3.7}, nó tương ứng với một dãy tăng
$(\widetilde{\mathfrak{a}}_n)$ các bó iđêan trong
$\widetilde{A}\equiv\sh{O}_X$. Để chứng minh dãy $(\mathfrak{a}_n)$ ổn
định, chỉ cần chứng minh dãy $(\widetilde{\mathfrak{a}}_n)$ ổn định. Nhưng
hạn chế $\widetilde{\mathfrak{a}}_n|V_i$ là một bó iđêan quasi-coherent của
$\sh{O}_X|V_i$, vì nó là ảnh ngược của
$\widetilde{\mathfrak{a}}_n$ qua phép nhúng chính tắc $V_i\to X$
\sref[0]{0.5.1.4}; do đó $\widetilde{\mathfrak{a}}_n|V_i$ có dạng
$\widetilde{\mathfrak{a}}_{ni}$, trong đó $\mathfrak{a}_{ni}$ là một iđêan
của $A_i$ \sref{I.1.3.7}. Vì $A_i$ là Noether, dãy
$(\mathfrak{a}_{ni})$ ổn định với mọi $i$, từ đó suy ra mệnh đề.
\end{proof}

Cần lưu ý rằng lập luận trên cũng chứng minh rằng \emph{nếu $X$ là một tiền
lược đồ Noether, thì mọi dãy tăng các bó coherent các iđêan của $\sh{O}_X$
đều ổn định}.

\begin{proposition}[6.1.4]
\label{I.6.1.4}
Mọi tiền lược đồ con của một tiền lược đồ Noether (tương ứng, Noether địa
phương) đều là Noether (tương ứng, Noether địa phương).
\end{proposition}

\begin{proof}
Chỉ cần chứng minh đối với một tiền lược đồ Noether $X$; hơn nữa, theo định
nghĩa \sref{I.6.1.1}, ta quy ngay về trường hợp $X$ là một lược đồ affine.
Vì mọi tiền lược đồ con của $X$ là một tiền lược đồ con đóng của một tiền
lược đồ cảm sinh trên một tập mở \sref{I.4.1.3}, ta có thể giới hạn ở trường
hợp một tiền lược đồ con $Y$ đóng hoặc cảm sinh trên một tập mở của $X$.
Trường hợp $Y$ đóng là hiển nhiên, vì nếu $A$ là vành của $X$, thì ta biết
$Y$ là một lược đồ affine có vành $A/\mathfrak{J}$, trong đó
$\mathfrak{J}$ là một iđêan của $A$ \sref{I.4.2.3}; vì $A$ là Noether
\sref{I.6.1.3}, nên $A/\mathfrak{J}$ cũng là Noether.

Bây giờ giả sử $Y$ mở trong $X$; không gian nền của $Y$ là Noether
\sref{I.6.1.2}, do đó quasi-compact và vì thế là hợp hữu hạn của các tập mở
$D(f_i)$ ($f_i\in A$);
\oldpage[I]{142}
mọi việc quy về chứng minh mệnh đề khi $Y=D(f)$, với $f\in A$. Nhưng khi đó
$Y$ là một lược đồ affine có vành đẳng cấu với $A_f$ \sref{I.1.3.6}; vì
$A$ là Noether \sref{I.6.1.3}, nên $A_f$ cũng là Noether.
\end{proof}

\begin{env}[6.1.5]
\label{I.6.1.5}
Cần lưu ý rằng \emph{tích} của hai $S$-tiền lược đồ Noether không nhất thiết
là Noether, ngay cả khi các tiền lược đồ này là affine, vì tích tenxơ của hai
đại số Noether không nhất thiết là một vành Noether (xem \sref{I.6.3.8}).
\end{env}

\begin{proposition}[6.1.6]
\label{I.6.1.6}
Nếu $X$ là một tiền lược đồ Noether, căn lũy linh $\sh{N}_X$ của
$\sh{O}_X$ là lũy linh.
\end{proposition}

\begin{proof}
Thật vậy, có thể phủ $X$ bởi một số hữu hạn tập mở affine $U_i$, và chỉ cần
chứng minh rằng tồn tại các số nguyên $n_i$ sao cho
$(\sh{N}_X|U_i)^{n_i}=0$; nếu $n$ là số lớn nhất trong các $n_i$, thì ta có
$\sh{N}_X^n=0$. Do đó, ta quy về trường hợp $X=\Spec(A)$ là affine, với $A$
là một vành Noether; theo \sref{I.5.1.1} và \sref{I.1.3.13}, chỉ cần nhận
xét rằng căn lũy linh của $A$ là lũy linh
(\cite[p.~127, cor.~4]{I-11}).
\end{proof}

\begin{corollary}[6.1.7]
\label{I.6.1.7}
Cho $X$ là một tiền lược đồ Noether;
để $X$ là một lược đồ affine, điều kiện cần và đủ là $X_\red$ là một lược đồ affine.
\end{corollary}

\begin{proof}
Điều này suy ra từ \sref{I.6.1.6} và \sref{I.5.1.10}.
\end{proof}

\begin{lemma}[6.1.8]
\label{I.6.1.8}
Cho $X$ là một không gian tôpô, $x$ một điểm của $X$, và $U$ một lân cận mở
của $x$ chỉ có hữu hạn thành phần bất khả quy. Khi đó tồn tại một lân cận $V$
của $x$ sao cho mọi lân cận mở của $x$ nằm trong $V$ đều liên thông.
\end{lemma}

\begin{proof}
Cho $U_i$ ($1\leq i\leq m$) là các thành phần bất khả quy của $U$ không chứa
$x$; phần bù (trong $U$) của hợp các $U_i$ là một lân cận mở $V$ của $x$
trong $U$, do đó cũng trong $X$; hơn nữa, nó là giao với $U$ của phần bù
(trong $X$) của hợp các thành phần bất khả quy của $X$ không chứa $x$
\sref[0]{0.2.1.6}\footnote{\emph{[Dịch.] Nguồn tiếng Pháp đồng nhất trực tiếp
$V$ ở đây với phần bù trong $X$ của hợp các thành phần bất khả quy không chứa
$x$, bỏ qua giao với $U$. Đẳng thức được in là sai, chẳng hạn nếu $X$ bất khả
quy và $U\subsetneq X$: khi ấy phần bù đó là $X$, trong khi $V=U$.}}.
Cho $W$ là một lân cận mở của $x$ nằm trong $V$. Các thành phần bất khả quy
của $W$ là các giao với $W$ của các thành phần bất khả quy của $U$ gặp $W$
\sref[0]{0.2.1.6}, do đó các thành phần ấy chứa $x$; vì chúng liên thông, $W$
cũng liên thông.
\end{proof}

\begin{corollary}[6.1.9]
\label{I.6.1.9}
Một không gian tôpô Noether địa phương là liên thông địa phương (điều này kéo
theo, trong số những hệ quả khác, rằng các thành phần liên thông của nó là mở).
\end{corollary}

\begin{proposition}[6.1.10]
\label{I.6.1.10}
Cho $X$ là một không gian tôpô Noether địa phương.
Các điều kiện sau là tương đương.
\begin{enumerate}
  \item[{\rm(a)}] Các thành phần bất khả quy của $X$ là mở.
  \item[{\rm(b)}] Các thành phần bất khả quy của $X$ chính là các thành phần
    liên thông của nó.
  \item[{\rm(c)}] Các thành phần liên thông của $X$ là bất khả quy.
  \item[{\rm(d)}] Hai thành phần bất khả quy phân biệt của $X$ có giao rỗng.
\end{enumerate}
Sau cùng, nếu $X$ là một tiền lược đồ, thì các điều kiện này cũng tương đương
với
\begin{enumerate}
  \item[{\rm(e)}] Với mọi $x\in X$, $\Spec(\sh{O}_x)$ là bất khả quy (nói cách
    khác, căn lũy linh của $\sh{O}_x$ là nguyên tố).
\end{enumerate}
\end{proposition}

\begin{proof}
Hiển nhiên (a) kéo theo (b), vì mọi không gian bất khả quy đều liên thông, và
(a) kéo theo rằng các thành phần bất khả quy của $X$ là các tập vừa mở vừa
đóng. Hiển nhiên (b) kéo theo (c); ngược lại, một tập đóng $F$ chứa
\oldpage[I]{143}
một thành phần liên thông $C$ của $X$, với $C$ khác $F$, không thể là bất khả
quy, vì do không liên thông, $F$ là hợp của hai tập không rỗng, rời nhau, vừa
mở vừa đóng trong $F$, và do đó đóng trong $X$; vì thế (c) kéo theo (b). Từ đó
suy ra ngay (c) kéo theo (d), vì hai thành phần liên thông phân biệt không có
điểm chung.

Cho đến đây, ta chưa sử dụng việc $X$ là Noether địa phương. Bây giờ giả sử
điều đó đúng và chứng minh (d) kéo theo (a): theo \sref[0]{0.2.1.6}, ta có thể
giới hạn về trường hợp không gian $X$ là Noether, và do đó chỉ có hữu hạn thành
phần bất khả quy. Vì các thành phần này đóng và đôi một rời nhau, chúng là mở.

Sau cùng, sự tương đương giữa (d) và (e) vẫn đúng ngay cả khi không giả sử
không gian nền của tiền lược đồ $X$ là Noether địa phương. Thật vậy, theo
\sref[0]{0.2.1.6}, ta có thể giới hạn về trường hợp $X=\Spec(A)$ là affine;
nói rằng $x$ nằm trong đúng một thành phần bất khả quy của $X$ tương đương với
việc $\mathfrak{j}_x$ chứa đúng một iđêan nguyên tố tối tiểu của $A$
\sref{I.1.1.14}, điều này tương đương với việc
$\mathfrak{j}_x\sh{O}_x$ chứa đúng một iđêan nguyên tố tối tiểu của
$\sh{O}_x$, từ đó suy ra kết luận.
\end{proof}

\begin{corollary}[6.1.11]
\label{I.6.1.11}
Cho $X$ là một không gian Noether địa phương. Để $X$ là bất khả quy, điều kiện
cần và đủ là $X$ liên thông và không rỗng, đồng thời hai thành phần bất khả quy
phân biệt bất kỳ của $X$ có giao rỗng. Nếu $X$ là một tiền lược đồ, thì điều
kiện sau cùng này tương đương với việc $\Spec(\sh{O}_x)$ là bất khả quy với mọi
$x\in X$.
\end{corollary}

\begin{proof}
Khẳng định thứ hai đã được chứng minh trong \sref{I.6.1.10}; do đó chỉ còn phải
chứng minh rằng các điều kiện của khẳng định thứ nhất là đủ. Nhưng theo
\sref{I.6.1.10}, các điều kiện này kéo theo rằng các thành phần bất khả quy của
$X$ chính là các thành phần liên thông của nó; vì $X$ liên thông và không rỗng,
nó là bất khả quy.
\end{proof}

\begin{corollary}[6.1.12]
\label{I.6.1.12}
Cho $X$ là một tiền lược đồ Noether địa phương. Điều kiện cần và đủ để $X$ là
nguyên là $X$ liên thông và không rỗng\footnote{\emph{[Dịch.] Nguồn tiếng
Pháp bỏ sót «và không rỗng». Nếu không có điều kiện này, trường hợp rỗng sẽ
thỏa điều kiện về mọi vành địa phương vì không có điểm nào, mặc dù không gian
rỗng không bất khả quy; hơn nữa, Hệ quả~\sref{I.6.1.11} yêu cầu rõ ràng rằng
$X$ không rỗng.}}, và $\sh{O}_x$ là một miền nguyên với mọi $x\in X$.
\end{corollary}

\begin{proposition}[6.1.13]
\label{I.6.1.13}
Cho $X$ là một tiền lược đồ Noether địa phương, và cho $x\in X$ là một điểm
sao cho căn lũy linh $\sh{N}_x$ của $\sh{O}_x$ là nguyên tố (tương ứng,
$\sh{O}_x$ là rút gọn, tương ứng, nguyên); khi đó tồn tại một lân cận mở $U$
của $x$ là bất khả quy (tương ứng, rút gọn, tương ứng, nguyên).
\end{proposition}

\begin{proof}
Chỉ cần xét hai trường hợp: $\sh{N}_x$ là nguyên tố và $\sh{N}_x=0$; giả
thiết thứ ba là hội của hai giả thiết đầu. Nếu $\sh{N}_x$ là nguyên tố, thì
$x$ chỉ thuộc một thành phần bất khả quy $Y$ duy nhất của $X$
\sref{I.6.1.10}; hợp của các thành phần bất khả quy của $X$ không chứa $x$ là
đóng (vì họ các thành phần này là hữu hạn địa phương), nên phần bù $U$ của hợp
đó là mở và nằm trong $Y$, do đó là bất khả quy \sref[0]{0.2.1.6}. Nếu
$\sh{N}_x=0$, thì cũng có $\sh{N}_y=0$ với mọi $y$ trong một lân cận của $x$,
vì $\sh{N}$ là quasi-coherent \sref{I.5.1.1}, và do đó coherent vì $X$ là
Noether địa phương; khi ấy kết luận suy ra từ \sref[0]{0.5.2.2}.
\end{proof}

\subsection{Các tiền lược đồ Artin}
\label{subsection:I.6.2}
\label{I.6.2}

\begin{definition}[6.2.1]
\label{I.6.2.1}
Ta nói một tiền lược đồ là \emph{Artin} nếu nó là affine và vành của nó là Artin.
\end{definition}

\oldpage[I]{144}
\begin{proposition}[6.2.2]
\label{I.6.2.2}
Cho một tiền lược đồ $X$, các điều kiện sau tương đương:
\begin{enumerate}
  \item[{\rm(a)}] $X$ là một lược đồ Artin;
  \item[{\rm(b)}] $X$ là Noether và không gian nền của nó là rời rạc;
  \item[{\rm(c)}] $X$ là Noether và các điểm của không gian nền của nó là đóng \emph{(điều kiện $\mathrm{T}_1$)}.
\end{enumerate}
Khi một trong các điều kiện trên được thỏa mãn, không gian nền của $X$ là hữu hạn, và vành $A$ của $X$ là tích trực tiếp của các vành địa phương (Artin) tại các điểm của $X$.
\end{proposition}

\begin{proof}
Ta biết rằng (a) kéo theo khẳng định cuối cùng (\cite[p.~205, th.~3]{I-13}); khi đó mọi iđêan nguyên tố của $A$ đều cực đại và là ảnh ngược của iđêan cực đại thuộc một trong các thành phần địa phương của $A$, do đó không gian $X$ là hữu hạn và rời rạc;
(a) do đó kéo theo (b), còn (b) hiển nhiên kéo theo (c).
Để thấy (c) kéo theo (a), trước hết ta chứng minh rằng $X$ là hữu hạn;
thật vậy, ta có thể quy về trường hợp $X$ là affine, và ta biết rằng một vành Noether mà mọi iđêan nguyên tố đều cực đại là Artin (\cite[p.~203]{I-13}), từ đó suy ra khẳng định.
Không gian nền $X$ khi đó là rời rạc, là tổng tôpô của một số hữu hạn các điểm $x_i$, và các vành địa phương $\sh{O}_{x_i}=A_i$ là Artin;
hiển nhiên $X$ đẳng cấu với lược đồ affine $\Spec(A)$, trong đó $A$ là tích trực tiếp của các $A_i$ \sref{I.1.7.3}.
\end{proof}

\subsection{Các cấu xạ kiểu hữu hạn}
\label{subsection:I.6.3}

\begin{definition}[6.3.1]
\label{I.6.3.1}
Ta nói một cấu xạ $f:X\to Y$ là \emph{kiểu hữu hạn} nếu $Y$ là hợp của một họ $(V_\alpha)$ các tập mở affine có tính chất sau:
\begin{enumerate}
  \item[(P)] $f^{-1}(V_\alpha)$ là hợp hữu hạn của các tập mở affine $U_{\alpha i}$ sao cho mỗi vành $A(U_{\alpha i})$ là một đại số kiểu hữu hạn trên $A(V_\alpha)$.
\end{enumerate}
Khi đó ta cũng nói $X$ là một tiền lược đồ kiểu hữu hạn trên $Y$, hoặc một $Y$-tiền lược đồ kiểu hữu hạn.
\end{definition}

\begin{proposition}[6.3.2]
\label{I.6.3.2}
Nếu $f:X\to Y$ là một cấu xạ kiểu hữu hạn, thì mọi tập mở affine $W$ của $Y$ đều có tính chất \emph{(P)} của \sref{I.6.3.1}.
\end{proposition}

Trước hết, ta chứng minh bổ đề sau.
\begin{lemma}[6.3.2.1]
\label{I.6.3.2.1}
Nếu $T\subset Y$ là một tập mở affine có tính chất \emph{(P)}, thì, với mọi $g\in A(T)$, $D(g)$ cũng có tính chất \emph{(P)}.
\end{lemma}

\begin{proof}
Theo giả thiết, $f^{-1}(T)$ là hợp hữu hạn của các tập mở affine $Z_j$ sao cho $A(Z_j)$ là một đại số kiểu hữu hạn trên $A(T)$;
gọi $\vphi_j:A(T)\to A(Z_j)$ là đồng cấu vành tương ứng với hạn chế của $f$ lên $Z_j$ \sref{I.2.2.4}, và đặt $g_j=\vphi_j(g)$;
khi đó $f^{-1}(D(g))\cap Z_j=D(g_j)$ \sref{I.1.2.2.2}.
Mặt khác, $A(D(g_j))=A(Z_j)_{g_j}=A(Z_j)[1/g_j]$ là kiểu hữu hạn trên $A(Z_j)$ và, \emph{a fortiori}, trên $A(T)$ theo giả thiết, do đó cũng là kiểu hữu hạn trên $A(D(g))=A(T)[1/g]$; điều này chứng minh bổ đề.
\end{proof}

\begin{proof}
Theo bổ đề trên, vì $W$ là quasi-compact \sref{I.1.1.10}, tồn tại một phủ hữu hạn của $W$ bởi các tập có dạng \erratum[II]{$D(g_i)\subset W$}, trong đó mỗi $g_i$ thuộc một vành $A(V_{\alpha(i)})$.
Mỗi $D(g_i)$, vì là quasi-compact, là hợp hữu hạn của các tập $D(h_{ik})$, trong đó $h_{ik}\in A(W)$;
nếu $\vphi_i:A(W)\to A(D(g_i))$ là ánh xạ chính tắc, thì $D(h_{ik})=D(\vphi_i(h_{ik}))$ theo \sref{I.1.2.2.2}.
Theo \sref{I.6.3.2.1}, mỗi $f^{-1}(D(h_{ik}))$ thừa nhận một phủ hữu hạn bởi các tập mở affine $U_{ijk}$ sao cho các $A(U_{ijk})$ là những đại số kiểu hữu hạn trên $A(D(h_{ik}))=A(W)[1/h_{ik}]$; từ đó suy ra mệnh đề.
\end{proof}

\oldpage[I]{145}
Do đó, ta có thể nói rằng khái niệm tiền lược đồ kiểu hữu hạn trên $Y$ là \emph{địa phương trên $Y$}.

\begin{proposition}[6.3.3]
\label{I.6.3.3}
Cho $X$ và $Y$ là hai lược đồ affine;
điều kiện cần và đủ để $X$ là kiểu hữu hạn trên $Y$ là $A(X)$ là một đại số kiểu hữu hạn trên $A(Y)$.
\end{proposition}

\begin{proof}
Vì điều kiện hiển nhiên là đủ, ta chứng minh nó là cần.
Đặt $A=A(Y)$ và $B=A(X)$;
theo \sref{I.6.3.2}, tồn tại một phủ mở affine hữu hạn $(V_i)$ của $X$ sao cho mỗi vành $A(V_i)$ là một $A$-đại số kiểu hữu hạn.
Hơn nữa, vì các $V_i$ là quasi-compact, ta có thể phủ mỗi tập bởi một số hữu hạn các tập mở có dạng $D(g_{ij})\subset V_i$, trong đó $g_{ij}\in B$;
nếu $\vphi_i$ là đồng cấu $B\to A(V_i)$ ứng với phép nhúng chính tắc $V_i\to X$, thì ta có $B_{g_{ij}}=(A(V_i))_{\vphi_i(g_{ij})}=A(V_i)[1/\vphi_i(g_{ij})]$, nên $B_{g_{ij}}$ là một $A$-đại số kiểu hữu hạn.
Do đó, ta có thể quy về trường hợp $V_i=D(g_i)$ với $g_i\in B$.
Theo giả thiết, tồn tại một tập con hữu hạn $F_i$ của $B$ và một số nguyên $n_i\geq0$ sao cho $B_{g_i}$ là đại số sinh trên $A$ bởi các phần tử $b_i/g_i^{n_i}$, trong đó $b_i$ chạy qua $F_i$.
Vì số các $g_i$ là hữu hạn, ta có thể giả sử mọi $n_i$ đều bằng cùng một số nguyên $n$.
Hơn nữa, vì các $D(g_i)$ tạo thành một phủ của $X$, iđêan của $B$ sinh bởi các $g_i$ bằng $B$; nói cách khác, tồn tại các $h_i\in B$ sao cho $\sum_i h_ig_i=1$.
Khi đó, cho $F$ là tập con hữu hạn của $B$ gồm hợp của các $F_i$, tập các $g_i$ và tập các $h_i$; ta chứng minh vành con $B'=A[F]$ của $B$ bằng $B$.
Theo giả thiết, với mọi $b\in B$ và mọi $i$, ảnh chính tắc của $b$ trong $B_{g_i}$ có dạng $b'_i/g_i^{m_i}$, trong đó $b'_i\in B'$;
nhân các $b'_i$ với các lũy thừa thích hợp của $g_i$, ta lại có thể giả sử mọi $m_i$ đều bằng cùng một số nguyên $m$.
Theo định nghĩa của các vành phân thức, do đó tồn tại một số nguyên $N$ (phụ thuộc vào $b$) sao cho $N\geq m$ và $g_i^Nb\in B'$ với mọi $i$;
nhưng, \emph{trong vành $B'$}, các $g_i^N$ sinh iđêan $B'$, vì các $g_i$ đã sinh nó (và các $h_i$ thuộc $B'$);
do đó tồn tại các $c_i\in B'$ sao cho $\sum_i c_ig_i^N=1$, suy ra $b=\sum_ic_ig_i^Nb\in B'$, Q.E.D.
\end{proof}

\begin{proposition}[6.3.4]
\label{I.6.3.4}
\medskip\noindent
\begin{enumerate}
  \item[{\rm(i)}] Mọi phép nhúng đóng đều là một cấu xạ kiểu hữu hạn.
  \item[{\rm(ii)}] Hợp thành của hai cấu xạ kiểu hữu hạn là kiểu hữu hạn.
  \item[{\rm(iii)}] Nếu $f:X\to X'$ và $g:Y\to Y'$ là hai $S$-cấu xạ kiểu hữu hạn, thì $f\times_S g$ là kiểu hữu hạn.
  \item[{\rm(iv)}] Nếu $f:X\to Y$ là một $S$-cấu xạ kiểu hữu hạn, thì $f_{(S')}$ là kiểu hữu hạn đối với mọi mở rộng $g:S'\to S$ của tiền lược đồ cơ sở.
  \item[{\rm(v)}] Nếu hợp thành $g\circ f$ của hai cấu xạ là kiểu hữu hạn, với $g$ tách, thì $f$ là kiểu hữu hạn.
  \item[{\rm(vi)}] Nếu một cấu xạ $f$ là kiểu hữu hạn, thì $f_\red$ cũng vậy.
\end{enumerate}
\end{proposition}

\begin{proof}
Theo \sref{I.5.5.12}, chỉ cần chứng minh (i), (ii) và (iv).

Để chứng minh (i), ta có thể quy về trường hợp một phép nhúng chính tắc
$X\to Y$, với $X$ là một tiền lược đồ con đóng của $Y$; hơn nữa, theo
\sref{I.6.3.2}, có thể giả sử $Y$ affine, khi đó $X$ cũng affine
\sref{I.4.2.3} và vành của nó đẳng cấu với một vành thương
$A/\mathfrak{J}$, trong đó $A$ là vành của $Y$ và $\mathfrak{J}$ là một
iđêan của $A$; vì $A/\mathfrak{J}$ là kiểu hữu hạn trên $A$, suy ra kết
luận.

Bây giờ ta chứng minh (ii). Cho $f:X\to Y$ và $g:Y\to Z$ là hai cấu xạ kiểu
hữu hạn, và cho $U$ là một tập mở affine của $Z$; $g^{-1}(U)$ có một phủ hữu
hạn bởi các tập mở affine $V_i$ sao cho mỗi $A(V_i)$ là một đại số kiểu hữu
hạn trên $A(U)$ \sref{I.6.3.2}; tương tự,
\oldpage[I]{146}
mỗi $f^{-1}(V_i)$ có một phủ hữu hạn bởi các tập mở affine $W_{ij}$ sao cho
mỗi $A(W_{ij})$ là một đại số kiểu hữu hạn trên $A(V_i)$ và do đó cũng là
một đại số kiểu hữu hạn trên $A(U)$, từ đó suy ra kết luận.

Cuối cùng, để chứng minh (iv), ta có thể quy về trường hợp $S=Y$; thật vậy,
$f_{(S')}$ cũng bằng $f_{(Y_{(S')})}$, khi xem $f$ là một $Y$-cấu xạ, và
phép mở rộng cơ sở là $Y_{(S')}\to Y$ \sref{I.3.3.9}. Khi đó, cho $p$ và
$q$ là các phép chiếu $X_{(S')}\to X$ và $X_{(S')}\to S'$. Cho $V$ là một
tập mở affine của $S$; $f^{-1}(V)$ là hợp hữu hạn của các tập mở affine
$W_i$ sao cho $A(W_i)$ là một đại số kiểu hữu hạn trên $A(V)$
\sref{I.6.3.2}. Cho $V'$ là một tập mở affine của $S'$ chứa trong
$g^{-1}(V)$; vì $f\circ p=g\circ q$, $q^{-1}(V')$ nằm trong hợp của các
$p^{-1}(W_i)$; mặt khác, giao $p^{-1}(W_i)\cap q^{-1}(V')$ được đồng nhất
với tích $W_i\times_V V'$ \sref{I.3.2.7}, là một lược đồ affine có vành
đẳng cấu với $A(W_i)\otimes_{A(V)}A(V')$ \sref{I.3.2.2}; theo giả thiết,
vành này là một đại số kiểu hữu hạn trên $A(V')$, điều này chứng minh mệnh
đề.
\end{proof}

\begin{corollary}[6.3.5]
\label{I.6.3.5}
Cho $f:X\to Y$ là một cấu xạ nhúng.
Nếu không gian nền của $Y$ (tương ứng, $X$) là Noether địa phương (tương
ứng, Noether), thì $f$ là kiểu hữu hạn.
\end{corollary}

\begin{proof}
Ta luôn có thể giả sử $Y$ affine \sref{I.6.3.2};
nếu không gian nền của $Y$ là Noether địa phương, ta còn có thể giả sử nó là
Noether, và khi đó không gian nền của $X$, vốn là một không gian con của nó,
cũng là Noether.
Nói cách khác, ta có thể giả sử $Y$ affine và không gian nền của $X$ là
Noether;
khi đó tồn tại một phủ của $X$ bởi một số hữu hạn các tập mở affine
$D(g_i)\subset Y$, trong đó $g_i\in A(Y)$, sao cho các giao
$X\cap D(g_i)$ đóng trong $D(g_i)$ (và do đó là các lược đồ affine
\sref{I.4.2.3}), vì $X$ đóng địa phương trong $Y$ \sref{I.4.1.3}.
Khi đó $A(X\cap D(g_i))$ là một đại số kiểu hữu hạn trên $A(D(g_i))$, theo
\sref{I.6.3.4}[i] và \sref{I.6.3.3}, còn
$A(D(g_i))=A(Y)_{g_i}=A(Y)[1/g_i]$ là kiểu hữu hạn trên $A(Y)$, điều này
kết thúc chứng minh.
\end{proof}

\begin{corollary}[6.3.6]
\label{I.6.3.6}
Cho $f:X\to Y$ và $g:Y\to Z$ là hai cấu xạ.
Nếu $g\circ f$ là kiểu hữu hạn và $X$ là Noether, hoặc nếu
$X\times_Z Y$ là Noether địa phương, thì $f$ là kiểu hữu hạn.
\end{corollary}

\begin{proof}
Điều này suy ra ngay từ chứng minh của \sref{I.5.5.12} và của
\sref{I.6.3.5} khi áp dụng cho cấu xạ nhúng $\Gamma_f$.
\end{proof}

\begin{proposition}[6.3.7]
\label{I.6.3.7}
Cho $f:X\to Y$ là một cấu xạ kiểu hữu hạn;
nếu $Y$ là Noether (tương ứng, Noether địa phương), thì $X$ là Noether
(tương ứng, Noether địa phương).
\end{proposition}

\begin{proof}
Ta có thể giới hạn việc chứng minh ở trường hợp $Y$ là Noether.
Khi đó $Y$ là hợp hữu hạn của các tập mở affine $V_i$ sao cho các vành
$A(V_i)$ là Noether.
Theo \sref{I.6.3.2}, mỗi $f^{-1}(V_i)$ là hợp của một số hữu hạn các tập
mở affine $W_{ij}$ sao cho các $A(W_{ij})$ là những đại số kiểu hữu hạn
trên $A(V_i)$, do đó là các vành Noether;
điều này chứng minh rằng $X$ là Noether.
\end{proof}

\begin{corollary}[6.3.8]
\label{I.6.3.8}
Cho $X$ là một tiền lược đồ kiểu hữu hạn trên $S$.
Đối với mọi mở rộng cơ sở $S'\to S$ sao cho $S'$ là Noether
(tương ứng, Noether địa phương), $X_{(S')}$ là Noether
(tương ứng, Noether địa phương).
\end{corollary}

\begin{proof}
Điều này suy ra từ \sref{I.6.3.7}, vì $X_{(S')}$ là kiểu hữu hạn trên
$S'$ theo \sref{I.6.3.4}[iv].
\end{proof}

Cũng có thể nói rằng, trong một tích $X\times_S Y$ của các $S$-tiền lược đồ,
nếu \emph{một} trong hai
\oldpage[I]{147}
thừa số $X$, $Y$ là \emph{kiểu hữu hạn} trên $S$ và
\emph{thừa số kia là Noether} (tương ứng, \emph{Noether địa phương}), thì
$X\times_S Y$ là \emph{Noether} (tương ứng, \emph{Noether địa phương}).

\begin{corollary}[6.3.9]
\label{I.6.3.9}
Cho $X$ là một tiền lược đồ kiểu hữu hạn trên một tiền lược đồ $S$ Noether
địa phương.
Khi đó mọi $S$-cấu xạ $f:X\to Y$ đều là kiểu hữu hạn.
\end{corollary}

\begin{proof}
Thật vậy, ta có thể giả sử $S$ là Noether;
nếu $\vphi:X\to S$ và $\psi:Y\to S$ là các cấu xạ cấu trúc, thì
$\vphi=\psi\circ f$, và $X$ là Noether theo \sref{I.6.3.7};
do đó, $f$ là kiểu hữu hạn theo \sref{I.6.3.6}.
\end{proof}

\begin{proposition}[6.3.10]
\label{I.6.3.10}
Cho $f:X\to Y$ là một cấu xạ kiểu hữu hạn.
Để $f$ là toàn ánh, điều kiện cần và đủ là, với mọi trường
\emph{đóng đại số} $\Omega$, ánh xạ $X(\Omega)\to Y(\Omega)$ tương ứng với
$f$ theo \sref{I.3.4.1} là toàn ánh.
\end{proposition}

\begin{proof}
Điều kiện là đủ, như thấy được bằng cách xét, với mỗi $y\in Y$, một trường
đóng đại số $\Omega$ mở rộng $\kres(y)$, và sơ đồ giao hoán
\[
  \xymatrix{
    X\ar[dd]_f\\
    & \Spec(\Omega)\ar[ul]\ar[dl]\\
    Y
  }
\]
(xem \sref{I.3.5.3}).
Ngược lại, giả sử $f$ toàn ánh và cho
$g:\{\xi\}=\Spec(\Omega)\to Y$ là một cấu xạ, trong đó $\Omega$ là một
trường đóng đại số.
Nếu xét sơ đồ
\[
  \xymatrix{
    X\ar[d]_f &
    X_{(\Omega)}\ar[l]\ar[d]^{f_{(\Omega)}}\\
    Y &
    \Spec(\Omega)\ar[l]
  }
\]
thì chỉ cần chứng minh rằng tồn tại một \emph{điểm hữu tỉ trên $\Omega$}
trong $X_{(\Omega)}$ (\sref{I.3.3.14}, \sref{I.3.4.3} và
\sref{I.3.4.4}).
Vì $f$ toàn ánh, $X_{(\Omega)}$ không rỗng \sref{I.3.5.10}; và vì $f$ là
kiểu hữu hạn, $f_{(\Omega)}$ cũng là kiểu hữu hạn \sref{I.6.3.4}[iv].
Do đó, $X_{(\Omega)}$ chứa một phần mở affine không rỗng $Z$ sao cho $A(Z)$
là một $\Omega$-đại số kiểu hữu hạn khác không.
Theo định lý không điểm của Hilbert~\cite{I-21}, tồn tại một $\Omega$-đồng cấu
$A(Z)\to\Omega$, do đó tồn tại một tiết diện của $X_{(\Omega)}$ trên
$\Spec(\Omega)$, điều này chứng minh mệnh đề.
\end{proof}

\subsection{Các tiền lược đồ đại số}
\label{subsection:I.6.4}
\label{I.6.4}

\begin{definition}[6.4.1]
\label{I.6.4.1}
Cho một trường $K$, ta gọi là \emph{$K$-tiền lược đồ đại số} một tiền lược đồ
$X$ kiểu hữu hạn trên $K$; $K$ được gọi là trường cơ sở của $X$.
Nếu hơn nữa $X$ là một lược đồ \emph{(hoặc nếu $X$ là một
\emph{$K$-lược đồ}, điều này tương đương theo \sref{I.5.5.8})}, ta cũng nói
rằng $X$ là một \emph{$K$-lược đồ đại số}.
\end{definition}

Mọi $K$-tiền lược đồ đại số đều là \emph{Noether} \sref{I.6.3.7}.

\begin{proposition}[6.4.2]
\label{I.6.4.2}
Cho $X$ là một $K$-tiền lược đồ đại số.
Để một điểm $x\in X$ là đóng, điều kiện cần và đủ là $\kres(x)$ là
một mở rộng đại số của $K$ có bậc hữu hạn.
\end{proposition}

\begin{proof}
Có thể giả sử $X$ là affine, với vành của nó là $A$, một $K$-đại số
kiểu hữu hạn.
Thật vậy, các tập mở affine $U$ của $X$ sao cho $A(U)$ là một
$K$-đại số kiểu hữu hạn tạo thành một phủ của $X$
\sref{I.6.3.1}.
Khi đó, các điểm đóng của $X$ là các điểm sao cho $\mathfrak{j}_x$ là
\oldpage[I]{148}
một iđêan cực đại của $A$, nói cách khác, sao cho $A/\mathfrak{j}_x$ là
một trường (nhất thiết bằng $\kres(x)$).
Vì $A/\mathfrak{j}_x$ là một $K$-đại số kiểu hữu hạn, ta thấy rằng, nếu
$x$ đóng, thì $\kres(x)$ là một trường đồng thời là một đại số kiểu hữu
hạn trên $K$, nên tất yếu là một $K$-đại số \emph{hạng hữu
hạn}~\cite{I-21}.
Ngược lại, nếu $\kres(x)$ có hạng hữu hạn trên $K$, thì
$A/\mathfrak{j}_x\subset\kres(x)$ cũng vậy; và vì mọi miền nguyên đồng thời là
một $K$-đại số hạng hữu hạn đều là một trường, ta có
$A/\mathfrak{j}_x=\kres(x)$, do đó $x$ đóng.
\end{proof}

\begin{corollary}[6.4.3]
\label{I.6.4.3}
Cho $K$ là một trường đóng đại số và $X$ là một $K$-tiền lược đồ đại số;
khi đó các điểm đóng của $X$ là các điểm hữu tỉ trên $K$
\sref{I.3.4.4} và được đồng nhất chính tắc với các điểm của $X$ nhận giá trị
trong $K$.
\end{corollary}

\begin{proposition}[6.4.4]
\label{I.6.4.4}
Cho $X$ là một tiền lược đồ đại số trên một trường $K$.
Các tính chất sau tương đương:
\begin{enumerate}
  \item[{\rm(a)}] $X$ là Artin.
  \item[{\rm(b)}] Không gian nền của $X$ là rời rạc.
  \item[{\rm(c)}] Không gian nền của $X$ chỉ có hữu hạn điểm đóng.
  \item[{\rm(c')}] Không gian nền của $X$ là hữu hạn.
  \item[{\rm(d)}] Các điểm của $X$ đều đóng.
  \item[{\rm(e)}] $X$ đẳng cấu với $\Spec(A)$, trong đó $A$ là một $K$-đại số có hạng hữu hạn.
\end{enumerate}
\end{proposition}

\begin{proof}
Vì $X$ là Noether, từ \sref{I.6.2.2} suy ra các điều kiện (a), (b) và (d) tương đương và kéo theo (c) và (c$'$);
mặt khác, hiển nhiên (e) kéo theo (a).
Còn phải chứng minh (c) kéo theo (d) và (e);
ta có thể quy về trường hợp $X$ là affine.
Khi đó $A(X)$ là một $K$-đại số kiểu hữu hạn \sref{I.6.3.3}, do đó là một vành Jacobson (\cite[p.~3-11 và 3-12]{I-1}), trong đó theo giả thiết chỉ có hữu hạn iđêan cực đại.
Vì giao hữu hạn của các iđêan nguyên tố chỉ có thể là một iđêan nguyên tố nếu nó bằng một trong các iđêan ấy, nên mọi iđêan nguyên tố của $A(X)$ đều cực đại, từ đó suy ra (d).
Hơn nữa, theo \sref{I.6.2.2}, $A(X)$ là một $K$-đại số Artin kiểu hữu hạn, do đó tất yếu có \emph{hạng hữu hạn}~\cite{I-21}.
\end{proof}

\begin{env}[6.4.5]
\label{I.6.4.5}
Khi các điều kiện của \sref{I.6.4.4} được thỏa mãn, ta nói $X$ là một lược đồ \emph{hữu hạn trên $K$} (xem \sref[II]{II.6.1.1}), hoặc một \emph{$K$-lược đồ hữu hạn}, có \emph{hạng} $[A:K]$, cũng được ký hiệu bởi $\rg_K(X)$;
nếu $X$ và $Y$ là các $K$-lược đồ hữu hạn, thì
\[
  \rg_K(X\sqcup Y)=\rg_K(X)+\rg_K(Y),
  \tag{6.4.5.1}
\]
\[
  \rg_K(X\times_K Y)=\rg_K(X)\rg_K(Y),
  \tag{6.4.5.2}
\]
như suy ra từ \sref{I.3.2.2}.
\end{env}

\begin{corollary}[6.4.6]
\label{I.6.4.6}
Cho $X$ là một $K$-lược đồ hữu hạn.
Với mọi mở rộng $K'$ của $K$, $X\otimes_K K'$ là một $K'$-lược đồ hữu hạn, và hạng của nó trên $K'$ bằng hạng của $X$ trên $K$.
\end{corollary}

\begin{proof}
Nếu $A=A(X)$, ta có $[A\otimes_K K':K']=[A:K]$.
\end{proof}

\begin{corollary}[6.4.7]
\label{I.6.4.7}
Cho $X$ là một lược đồ hữu hạn trên một trường $K$;
đặt $n=\sum_{x\in X}[\kres(x):K]_s$ \emph{(nhắc lại rằng, nếu $K'$ là một mở rộng của $K$, thì $[K':K]_s$ là \emph{hạng khả ly} của $K'$ trên $K$, tức hạng của mở rộng đại số khả ly lớn nhất của $K$ chứa trong $K'$)};
\oldpage[I]{149}
khi đó,
với mọi trường đóng đại số $\Omega$ mở rộng $K$, không gian nền của $X\otimes_K\Omega$ có đúng $n$ điểm, được đồng nhất với các điểm của $X$ nhận giá trị trong $\Omega$.
\end{corollary}

\begin{proof}
Hiển nhiên ta có thể quy về trường hợp vành $A=A(X)$ là \emph{địa phương} \sref{I.6.2.2};
gọi $\mathfrak{m}$ là iđêan cực đại của nó và $L=A/\mathfrak{m}$ là trường thặng dư của nó, một mở rộng đại số của $K$.
Các điểm của $X$ nhận giá trị trong $\Omega$ khi đó tương ứng song ánh với các $\Omega$-nhát cắt của $X\otimes_K\Omega$ (\sref{I.3.4.1} và \sref{I.3.3.14}), và cũng với các $K$-đồng cấu từ $L$ vào $\Omega$ (\sref{I.1.7.3}); từ đó suy ra mệnh đề (Bourbaki, \emph{Alg.}, chương~V, \textsection7, số~5, mệnh đề~8), khi tính đến \sref{I.6.4.3}.
\end{proof}

\begin{env}[6.4.8]
\label{I.6.4.8}
Số $n$ được định nghĩa trong \sref{I.6.4.7} được gọi là \emph{hạng khả ly} của $A$ (hoặc của $X$) trên $K$, hoặc còn là \emph{số điểm hình học} của $X$;
do đó, nó bằng số phần tử của $X(\Omega)_K$.
Từ định nghĩa này suy ra ngay rằng, với mọi mở rộng $K'$ của $K$, $X\otimes_K K'$ có cùng số điểm hình học như $X$.
Nếu ký hiệu số này bởi $n(X)$, rõ ràng rằng, nếu $X$ và $Y$ là hai lược đồ hữu hạn trên $K$, thì
\[
  n(X\sqcup Y)=n(X)+n(Y).
  \tag{6.4.8.1}
\]
Dưới cùng các giả thiết, ta cũng có
\[
  n(X\times_K Y)=n(X)n(Y)
  \tag{6.4.8.2}
\]
như suy ra ngay từ cách diễn giải $n(X)$ là số phần tử của $X(\Omega)_K$ và từ công thức~\hyperref[I.3.4.3.1]{(3.4.3.1)}.
\end{env}

\begin{proposition}[6.4.9]
\label{I.6.4.9}
Cho $K$ là một trường, $X$ và $Y$ là hai $K$-tiền lược đồ đại số, $f:X\to Y$ là một $K$-cấu xạ, và $\Omega$ là một trường đóng đại số mở rộng $K$, có bậc siêu việt vô hạn trên $K$.
Để $f$ là toàn ánh, điều kiện cần và đủ là ánh xạ $X(\Omega)_K\to Y(\Omega)_K$ tương ứng với $f$ theo \sref{I.3.4.1} là toàn ánh.
\end{proposition}

\begin{proof}
Tính cần thiết suy ra từ \sref{I.6.3.10}, lưu ý rằng $f$ tất yếu là kiểu hữu hạn \sref{I.6.3.9}.
Để thấy điều kiện là đủ, ta lập luận như trong \sref{I.6.3.10}, lưu ý rằng, với mọi $y\in Y$, $\kres(y)$ là một mở rộng kiểu hữu hạn của $K$ và do đó $K$-đẳng cấu với một trường con của $\Omega$.
\end{proof}

\begin{remark}[6.4.10]
\label{I.6.4.10}
Ta sẽ thấy trong Chương~IV rằng kết luận của \sref{I.6.4.9} vẫn còn đúng mà không có giả thiết nào về bậc siêu việt của $\Omega$ trên $K$.
\end{remark}

\begin{proposition}[6.4.11]
\label{I.6.4.11}
Nếu $f:X\to Y$ là một cấu xạ kiểu hữu hạn, thì, với mọi $y\in Y$, thớ $f^{-1}(y)$ là một tiền lược đồ đại số trên trường thặng dư $\kres(y)$ và, với mọi $x\in f^{-1}(y)$, $\kres(x)$ là một mở rộng kiểu hữu hạn của $\kres(y)$.
\end{proposition}

\begin{proof}
Vì $f^{-1}(y)=X\otimes_Y\kres(y)$ \sref{I.3.6.3}, mệnh đề suy ra từ \sref{I.6.3.4}[iv] và \sref{I.6.3.3}.
\end{proof}

\begin{proposition}[6.4.12]
\label{I.6.4.12}
Cho $f:X\to Y$ và $g:Y'\to Y$ là hai cấu xạ;
đặt $X'=X\times_Y Y'$ và cho $f'=f_{(Y')}:X'\to Y'$.
Cho $y'\in Y'$ và đặt $y=g(y')$;
nếu thớ $f^{-1}(y)$ là một lược đồ đại số hữu hạn trên $\kres(y)$, thì thớ $f'^{-1}(y')$ là một lược đồ đại số hữu hạn trên $\kres(y')$, có cùng hạng và cùng số điểm hình học như $f^{-1}(y)$.
\end{proposition}

\begin{proof}
Xét đến tính bắc cầu của các thớ \sref{I.3.6.5}, điều này suy ra ngay từ \sref{I.6.4.6} và \sref{I.6.4.8}.
\end{proof}

\begin{env}[6.4.13]
\label{I.6.4.13}
Mệnh đề~\sref{I.6.4.11} cho thấy các cấu xạ kiểu hữu hạn tương ứng một cách trực giác với ``các họ đại số của các đa tạp đại số'', trong đó các điểm của $Y$
\oldpage[I]{150}
đóng vai trò ``tham số'', điều này mang lại cho các cấu xạ ấy một ý nghĩa ``hình học''.
Các cấu xạ không thuộc kiểu hữu hạn về sau chủ yếu sẽ xuất hiện trong các vấn đề về ``thay đổi tiền lược đồ cơ sở'', chẳng hạn bằng địa phương hóa hoặc đầy đủ hóa.
\end{env}

\subsection{Xác định cục bộ một cấu xạ}
\label{subsection:I.6.5}
\label{I.6.5}

\begin{proposition}[6.5.1]
\label{I.6.5.1}
Cho $X$ và $Y$ là hai $S$-tiền lược đồ, trong đó $Y$ thuộc kiểu hữu hạn trên $S$;
cho $x\in X$ và $y\in Y$ nằm trên cùng một điểm $s\in S$.
\begin{enumerate}
  \item[{\rm(i)}] Nếu hai $S$-cấu xạ $f=(\psi,\theta)$ và $f'=(\psi',\theta')$ từ $X$ vào $Y$ thỏa mãn $\psi(x)=\psi'(x)=y$, và các $\sh{O}_s$-đồng cấu (địa phương) $\theta_x^\sharp$ và ${\theta'}_x^\sharp$ từ $\sh{O}_y$ vào $\sh{O}_x$ là đồng nhất, thì $f$ và $f'$ trùng nhau trên một lân cận mở của $x$.
  \item[{\rm(ii)}] Giả sử thêm rằng $S$ là Noether địa phương.
    Với mọi $\sh{O}_s$-đồng cấu địa phương $\vphi:\sh{O}_y\to\sh{O}_x$, tồn tại một lân cận mở $U$ của $x$ trong $X$ và một $S$-cấu xạ $f=(\psi,\theta)$ từ $U$ vào $Y$ sao cho $\psi(x)=y$ và $\theta_x^\sharp=\vphi$.
\end{enumerate}
\end{proposition}

\begin{proof}
\medskip\noindent
\begin{enumerate}
  \item[(i)] Vì vấn đề là cục bộ trên $S$, $X$ và $Y$, ta có thể giả sử $S$, $X$ và $Y$ là affine, với các vành tương ứng $A$, $B$ và $C$, còn $f$ và $f'$ lần lượt có dạng $({}^a\vphi,\widetilde{\vphi})$ và $({}^a\vphi',\widetilde{\vphi}')$, trong đó $\vphi$ và $\vphi'$ là hai $A$-đồng cấu từ $C$ vào $B$ thỏa mãn $\vphi^{-1}(\mathfrak{j}_x)=\vphi'^{-1}(\mathfrak{j}_x)=\mathfrak{j}_y$, và các đồng cấu $\vphi_x$ và $\vphi'_x$ từ $C_y$ vào $B_x$, được suy ra từ $\vphi$ và $\vphi'$, là đồng nhất;
    ta còn có thể giả sử $C$ là một $A$-đại số \emph{thuộc kiểu hữu hạn}.
    Cho $c_i$ ($1\leq i\leq n$) là các phần tử sinh của $A$-đại số $C$, và đặt $b_i=\vphi(c_i)$ cùng $b'_i=\vphi'(c_i)$;
    theo giả thiết, ta có $b_i/1=b'_i/1$ trong vành phân thức $B_x$ ($1\leq i\leq n$).
    Điều đó có nghĩa là tồn tại các phần tử $s_i\in B\setmin\mathfrak{j}_x$ sao cho $s_i(b_i-b'_i)=0$ với $1\leq i\leq n$, và hiển nhiên ta có thể giả sử mọi $s_i$ đều bằng cùng một phần tử $g\in B\setmin\mathfrak{j}_x$.
    Suy ra $b_i/1=b'_i/1$ với $1\leq i\leq n$ trong vành phân thức $B_g$;
    nếu $i_g$ là đồng cấu chính tắc $B\to B_g$, thì $i_g\circ\vphi=i_g\circ\vphi'$;
    do đó, các hạn chế của $f$ và $f'$ lên $D(g)$ là đồng nhất.
  \item[(ii)] Ta có thể quy về cùng tình huống như trong (i), và giả sử thêm rằng vành $A$ là noetherian.
    Cho $c_i$ ($1\leq i\leq n$) là các phần tử sinh của $A$-đại số $C$, và cho $\alpha:A[X_1,\ldots,X_n]\to C$ là đồng cấu từ đại số đa thức $A[X_1,\ldots,X_n]$ vào $C$ gửi $X_i$ thành $c_i$ với $1\leq i\leq n$.
    Ngoài ra, cho $i_y$ là đồng cấu chính tắc $C\to C_y$, và xét đồng cấu hợp thành
    \[
      \beta:A[X_1,\ldots,X_n] \xrightarrow{\alpha} C \xrightarrow{i_y} C_y \xrightarrow{\vphi} B_x.
    \]
    Ký hiệu $\mathfrak{a}$ là hạt nhân của $\beta$;
    vì $A$ là noetherian, $A[X_1,\ldots,X_n]$ cũng vậy, và do đó $\mathfrak{a}$ có một hệ sinh hữu hạn $Q_j(X_1,\ldots,X_n)$ ($1\leq j\leq m$).
    Mặt khác, mỗi phần tử $\vphi(i_y(c_i))$ có thể viết dưới dạng $b_i/s_i$, trong đó $b_i\in B$ và $s_i\not\in\mathfrak{j}_x$;
    ta còn có thể giả sử mọi $s_i$ đều bằng cùng một phần tử $g\in B\setmin\mathfrak{j}_x$.
    Trong các điều kiện ấy, theo giả thiết ta có $Q_j(b_1/g,\ldots,b_n/g)=0$ trong $B_x$; đặt
    \[
      Q_j(X_1/T,\ldots,X_n/T)=P_j(X_1,\ldots,X_n,T)/T^{k_j}
    \]
    trong đó $P_j$ là thuần nhất bậc $k_j$.
    Khi đó đặt $d_j=P_j(b_1,\ldots,b_n,g)\in B$.
    Theo giả thiết, ta có $t_jd_j=0$ với một $t_j\in B\setmin\mathfrak{j}_x$ nào đó ($1\leq j\leq m$), và hiển nhiên ta có thể giả sử mọi $t_j$
\oldpage[I]{151}
    đều bằng cùng một phần tử $h\in B\setmin\mathfrak{j}_x$;
    suy ra $P_j(hb_1,\ldots,hb_n,hg)=0$ với $1\leq j\leq m$.
    Trong các điều kiện ấy, xét đồng cấu $\rho$ từ $A[X_1,\ldots,X_n]$ vào vành phân thức $B_{hg}$ gửi $X_i$ thành $hb_i/hg$ ($1\leq i\leq n$);
    ảnh của $\mathfrak{a}$ qua đồng cấu này bằng $0$, và \emph{a fortiori}, ảnh của hạt nhân $\alpha^{-1}(0)$ qua $\rho$ cũng vậy.
    Do đó, $\rho$ phân tích thành $A[X_1,\ldots,X_n]\xrightarrow{\alpha}C\xrightarrow{\gamma}B_{hg}$, với $\gamma(c_i)=hb_i/hg$, và rõ ràng, nếu $i_x$ là đồng cấu chính tắc $B_{hg}\to B_x$, thì sơ đồ
    \[
    \label{I.6.5.1.1}
      \xymatrix{
        C\ar[r]^\gamma\ar[d]_{i_y} &
        B_{hg}\ar[d]^{i_x}\\
        C_y\ar[r]^\vphi &
        B_x
      }
      \tag{6.5.1.1}
    \]
    giao hoán; do đó $\vphi=\gamma_x$, và vì $\vphi$ là một đồng cấu địa phương, $^a\gamma(x)=y$;
    vì vậy, $f=({}^a\gamma,\widetilde{\gamma})$ là một $S$-cấu xạ từ lân cận $D(hg)$ của $x$ vào $Y$, như đã khẳng định trong mệnh đề.
\end{enumerate}
\end{proof}

\begin{corollary}[6.5.2]
\label{I.6.5.2}
Dưới các giả thiết của \sref{I.6.5.1}[ii]\footnote{\emph{[Dịch.] Nguồn tiếng Pháp dẫn ở đây tới (6.5.1, (iii)), một mục không tồn tại; giả thiết thích hợp là (6.5.1, (ii)).}}, nếu thêm nữa $X$ là kiểu hữu hạn trên $S$, ta có thể giả sử cấu xạ $f$ là kiểu hữu hạn.
\end{corollary}

\begin{proof}
Điều này suy ra từ \sref{I.6.3.6}.
\end{proof}

\begin{corollary}[6.5.3]
\label{I.6.5.3}
Giả sử các giả thiết của \sref{I.6.5.1}[ii] được thỏa mãn, và giả sử thêm rằng $Y$ là nguyên và $\vphi$ là một đồng cấu đơn ánh.
Khi đó ta có thể giả sử $f=({}^a\gamma,\widetilde{\gamma})$, trong đó $\gamma$ là đơn ánh.
\end{corollary}

\begin{proof}
Thật vậy, ta có thể giả sử $C$ là nguyên \sref{I.5.1.4}, và do đó $i_y$ là đơn ánh; khi ấy từ sơ đồ~\sref{I.6.5.1.1} suy ra $\gamma$ là đơn ánh.
\end{proof}

\begin{proposition}[6.5.4]
\label{I.6.5.4}
Cho $f=(\psi,\theta):X\to Y$ là một cấu xạ thuộc kiểu hữu hạn, $x$ là một điểm của $X$ và $y=\psi(x)$.
\begin{enumerate}
  \item[{\rm(i)}] Để $f$ là một phép nhúng địa phương tại điểm $x$ \sref{I.4.5.1}, điều kiện cần và đủ là $\theta_x^\sharp:\sh{O}_y\to\sh{O}_x$ toàn ánh.
  \item[{\rm(ii)}] Giả sử thêm rằng $Y$ là Noether địa phương.
    Để $f$ là một đẳng cấu địa phương tại điểm $x$ \sref{I.4.5.2}, điều kiện cần và đủ là $\theta_x^\sharp$ là một đẳng cấu.
\end{enumerate}
\end{proposition}

\begin{proof}
\medskip\noindent
\begin{enumerate}
  \item[(ii)] Theo \sref{I.6.5.1}, khi ấy tồn tại một lân cận mở $V$ của $y$ và một cấu xạ $g:V\to X$ sao cho $g\circ f$ (tương ứng, $f\circ g$) được xác định và trùng với đồng nhất trên một lân cận của $x$ (tương ứng, của $y$), từ đó dễ dàng suy ra rằng $f$ là một đẳng cấu địa phương.
  \item[(i)] Vì vấn đề là cục bộ trên $X$ và $Y$, ta có thể giả sử $X$ và $Y$ là affine, với các vành tương ứng $A$ và $B$;
    ta có $f=({}^a\vphi,\widetilde{\vphi})$, trong đó $\vphi$ là một đồng cấu vành $B\to A$ làm cho $A$ thành một $B$-đại số thuộc kiểu hữu hạn;
    ta có $\vphi^{-1}(\mathfrak{j}_x)=\mathfrak{j}_y$, và đồng cấu $\vphi_x:B_y\to A_x$ cảm sinh bởi $\vphi$ là \emph{toàn ánh}.
    Cho $(t_i)$ ($1\leq i\leq n$) là một hệ sinh của $B$-đại số $A$;
    giả thiết về $\vphi_x$ kéo theo rằng tồn tại các $b_i\in B$ và một $c\in B\setmin\mathfrak{j}_y$ sao cho, trong vành phân thức $A_x$, ta có $t_i/1=\vphi(b_i)/\vphi(c)$ với $1\leq i\leq n$.
    Do đó \sref{I.1.3.3}, tồn tại một $a\in A\setmin\mathfrak{j}_x$ sao cho, nếu đặt $g=a\vphi(c)$, thì ta cũng có $t_i/1=a\vphi(b_i)/g$ \emph{trong vành phân thức $A_g$}.
    Khi đó, theo giả thiết, tồn tại một đa thức $Q(X_1,\ldots,X_n)$ với các hệ số trong vành $\vphi(B)$,
\oldpage[I]{152}
    sao cho $a=Q(t_1,\ldots,t_n)$;
    đặt $Q(X_1/T,\ldots,X_n/T)=P(X_1,\ldots,X_n,T)/T^m$, trong đó $P$ thuần nhất bậc $m$.
    Trong vành $A_g$, ta có
    \[
      a/1=a^mP(\vphi(b_1),\ldots,\vphi(b_n),\vphi(c))/g^m=a^m\vphi(d)/g^m
    \]
    trong đó $d\in B$.
    Vì trong $A_g$, $g/1=(a/1)(\vphi(c)/1)$ khả nghịch theo định nghĩa, nên $a/1$ và $\vphi(c)/1$ cũng khả nghịch; do đó có thể viết $a/1=(\vphi(d)/1)(\vphi(c)/1)^{-m}$.
    Từ đó suy ra $\vphi(d)/1$ cũng khả nghịch trong $A_g$.
    Đặt $h=cd$;
    vì $\vphi(h)/1$ khả nghịch trong $A_g$, đồng cấu hợp thành $B\xrightarrow{\vphi}A\to A_g$ phân tích thành $B\to B_h\xrightarrow{\gamma}A_g$ \sref[0]{0.1.2.4}.
    Ta chứng minh $\gamma$ là \emph{toàn ánh};
    chỉ cần kiểm tra rằng ảnh của $B_h$ trong $A_g$ chứa các $t_i/1$ và $(g/1)^{-1}$.
    Thật vậy, ta có $(g/1)^{-1}=(\vphi(c)/1)^{m-1}(\vphi(d)/1)^{-1}=\gamma(c^m/h)$, và $a/1=\gamma(d^{m+1}/h^m)$, do đó $(a\vphi(b_i))/1=\gamma(b_id^{m+1}/h^m)$; vì $t_i/1=(a\vphi(b_i)/1)(g/1)^{-1}$, khẳng định của ta được chứng minh.
    Việc chọn $h$ kéo theo $\psi(D(g))\subset D(h)$, và ta cũng biết rằng hạn chế của $f$ lên $D(g)$ bằng $({}^a\gamma,\widetilde{\gamma})$;
    vì $\gamma$ toàn ánh, hạn chế này là một phép nhúng đóng từ $D(g)$ vào $D(h)$ \sref{I.4.2.3}.
\end{enumerate}
\end{proof}

\begin{corollary}[6.5.5]
\label{I.6.5.5}
Cho $f=(\psi,\theta):X\to Y$ là một cấu xạ thuộc kiểu hữu hạn.
Giả sử $X$ bất khả quy, ký hiệu $x$ là điểm tổng quát của nó và đặt $y=\psi(x)$.
\begin{enumerate}
  \item[{\rm(i)}] Để $f$ là một phép nhúng địa phương tại một điểm nào đó của $X$, điều kiện cần và đủ là $\theta_x^\sharp:\sh{O}_y\to\sh{O}_x$ toàn ánh.
  \item[{\rm(ii)}] Giả sử thêm rằng $Y$ bất khả quy và Noether địa phương.
    Để $f$ là một đẳng cấu địa phương tại một điểm nào đó của $X$, điều kiện cần và đủ là $y$ là điểm tổng quát của $Y$ \emph{(hoặc, tương đương \sref[0]{0.2.1.4}, $f$ là một cấu xạ \emph{trội})} và $\theta_x^\sharp$ là một đẳng cấu \emph{(nói cách khác, $f$ là \emph{song hữu tỉ} \sref{I.2.2.9})}.
\end{enumerate}
\end{corollary}

\begin{proof}
Hiển nhiên (i) suy ra từ \sref{I.6.5.4}[i], vì mọi tập mở khác rỗng của $X$ đều chứa $x$;
tương tự, (ii) suy ra từ \sref{I.6.5.4}[ii].
\end{proof}

\subsection{Các cấu xạ quasi-compact và các cấu xạ kiểu hữu hạn địa phương}
\label{subsection:I.6.6}

\begin{definition}[6.6.1]
\label{I.6.6.1}
Ta nói một cấu xạ $f:X\to Y$ là \emph{quasi-compact} nếu ảnh ngược của mọi tập mở quasi-compact của $Y$ qua $f$ là quasi-compact.
\end{definition}

Cho $\mathfrak{B}$ là một cơ sở của tôpô trên $Y$ gồm các tập mở quasi-compact (chẳng hạn, các tập mở affine);
để $f$ là quasi-compact, điều kiện cần và đủ là ảnh ngược của mọi phần tử của $\mathfrak{B}$ qua $f$ là quasi-compact (hoặc, tương đương, là một hợp \emph{hữu hạn} của các tập mở affine), vì mọi tập mở quasi-compact của $Y$ là hợp hữu hạn của các phần tử của $\mathfrak{B}$.
Chẳng hạn, nếu $X$ \emph{quasi-compact} và $Y$ \emph{affine}, thì \emph{mọi} cấu xạ $f: X\to Y$ đều quasi-compact:
thật vậy, $X$ là hợp hữu hạn của các tập mở affine $U_i$, và với mọi tập mở affine $V$ của $Y$, $U_i\cap f^{-1}(V)$ là affine \sref{I.5.5.10}, do đó quasi-compact.

Nếu $f: X\to Y$ là một cấu xạ quasi-compact, thì hiển nhiên, với mọi tập mở $V$ của $Y$, hạn chế của $f$ lên $f^{-1}(V)$ là một cấu xạ quasi-compact $f^{-1}(V)\to V$.
Ngược lại, nếu $(U_\alpha)$ là một phủ mở của $Y$, và $f: X\to Y$ là một cấu xạ sao cho các hạn chế $f^{-1}(U_\alpha)\to U_\alpha$ đều quasi-compact, thì $f$ là quasi-compact.

\begin{definition}[6.6.2]
\label{I.6.6.2}
Ta nói một cấu xạ $f: X\to Y$ là \emph{kiểu hữu hạn địa phương} nếu, với mọi $x\in X$, tồn tại một lân cận mở $U$ của $x$ và một lân cận mở $V\supset f(U)$ của $f(x)$\footnote{\emph{[Dịch.] Nguồn tiếng Pháp in ở đây «de $y$» mà không định nghĩa $y$; cách đọc đã hiệu chỉnh viết rõ $y=f(x)$.}} sao cho
\oldpage[I]{153}
hạn chế của $f$ trên $U$ là một cấu xạ kiểu hữu hạn từ $U$ vào $V$.
Khi đó ta cũng nói $X$ là một tiền lược đồ kiểu hữu hạn địa phương trên $Y$, hay một $Y$-tiền lược đồ kiểu hữu hạn địa phương.
\end{definition}

Từ \sref{I.6.3.2} suy ra ngay rằng, nếu $f$ là kiểu hữu hạn địa phương thì, với mọi tập mở $W$ của $Y$, hạn chế của $f$ trên $f^{-1}(W)$ là một cấu xạ $f^{-1}(W)\to W$ kiểu hữu hạn địa phương.

Nếu $Y$ là Noether địa phương và $X$ kiểu hữu hạn địa phương trên $Y$, thì $X$ là Noether địa phương theo \sref{I.6.3.7}.

\begin{proposition}[6.6.3]
\label{I.6.6.3}
Một cấu xạ $f: X\to Y$ là kiểu hữu hạn khi và chỉ khi nó quasi-compact và kiểu hữu hạn địa phương.
\end{proposition}

\begin{proof}
Tính cần thiết của các điều kiện là hiển nhiên, theo \sref{I.6.3.1} và nhận xét sau \sref{I.6.6.1}.
Ngược lại, giả sử các điều kiện này được thỏa mãn và cho $U$ là một tập mở affine của $Y$, có vành $A$;
với mọi $x\in f^{-1}(U)$, theo giả thiết tồn tại một lân cận $V(x)\subset f^{-1}(U)$ của $x$ và một lân cận $W(x)\subset U$ của $y=f(x)$, chứa $f(V(x))$, sao cho hạn chế của $f$ trên $V(x)$ là một cấu xạ $V(x)\to W(x)$ kiểu hữu hạn.
Thay $W(x)$ bằng một lân cận $W_1(x)\subset W(x)$ của $y=f(x)$\footnote{\emph{[Dịch.] Nguồn tiếng Pháp in ở đây «de $x$». Vì $W_1(x)\subset W(x)\subset U\subset Y$, nên đó phải là một lân cận của $y=f(x)$.}} có dạng $D(g)$ (với $g\in A$), và thay $V(x)$ bằng $V(x)\cap f^{-1}(W_1(x))$, ta có thể giả sử $W(x)$ có dạng $D(g)$, do đó kiểu hữu hạn trên $U$ (vì vành của nó được viết là $A[1/g]$);
do đó $V(x)$ kiểu hữu hạn trên $U$.
Hơn nữa, $f^{-1}(U)$ quasi-compact theo giả thiết, nên là hợp của một số hữu hạn các tập mở $V(x_i)$, điều này kết thúc chứng minh.
\end{proof}

\begin{proposition}[6.6.4]
\label{I.6.6.4}
\medskip\noindent
\begin{enumerate}
  \item[{\rm(i)}] Một phép nhúng $X\to Y$ là quasi-compact nếu nó là một phép nhúng đóng, hoặc nếu không gian nền của $Y$ là Noether địa phương, hoặc nếu không gian nền của $X$ là Noether.
  \item[{\rm(ii)}] Hợp thành của hai cấu xạ quasi-compact là quasi-compact.
  \item[{\rm(iii)}] Nếu $f: X\to Y$ là một $S$-cấu xạ quasi-compact, thì $f_{(S')}: X_{(S')}\to Y_{(S')}$ cũng vậy đối với mọi phép mở rộng $g: S'\to S$ của tiền lược đồ cơ sở.
  \item[{\rm(iv)}] Nếu $f:X\to X'$ và $g:Y\to Y'$ là hai $S$-cấu xạ quasi-compact, thì $f\times_S g$ là quasi-compact.
  \item[{\rm(v)}] Nếu hợp thành $g\circ f$ của hai cấu xạ $f: X\to Y$ và $g:Y\to Z$ là quasi-compact, và nếu $g$ tách hoặc không gian nền của $X$ là Noether địa phương, thì $f$ là quasi-compact.
  \item[{\rm(vi)}] Một cấu xạ $f$ là quasi-compact khi và chỉ khi $f_\red$ là quasi-compact.
\end{enumerate}
\end{proposition}

\begin{proof}
Ta lưu ý rằng (vi) là hiển nhiên, vì tính chất quasi-compact của một cấu xạ chỉ phụ thuộc vào ánh xạ liên tục tương ứng giữa các không gian nền.
Ta cũng chứng minh phần của (v) tương ứng với trường hợp không gian nền của $X$ được giả sử là Noether địa phương.
Đặt $h=g\circ f$, và cho $U$ là một tập mở quasi-compact của $Y$;
$g(U)$ là quasi-compact (nhưng không nhất thiết mở) trong $Z$, do đó được chứa trong một hợp hữu hạn các tập mở quasi-compact $V_j$ \sref{I.2.1.3}, và khi đó $f^{-1}(U)$ được chứa trong hợp của các $h^{-1}(V_j)$, vốn là các không gian con quasi-compact của $X$ và do đó là các không gian con Noether.
Vì thế, theo \sref[0]{0.2.2.3}, $f^{-1}(U)$ là một không gian Noether và, \emph{a fortiori}, quasi-compact.

Để chứng minh các khẳng định còn lại, chỉ cần chứng minh (i), (ii) và (iii) \sref{I.5.5.12}.
Thật vậy, (ii) là hiển nhiên, còn (i) suy ra từ \sref{I.6.3.5} khi không gian $Y$ là Noether địa phương hoặc không gian $X$ là Noether, và là hiển nhiên đối với một phép nhúng đóng.

Để chứng minh (iii),
\oldpage[I]{154}
ta có thể quy về trường hợp $S=Y$ \sref{I.3.3.11};
đặt $f'=f_{(S')}$, và cho $U'$ là một tập mở quasi-compact của $S'$.
Với mọi $s'\in U'$, cho $T$ là một lân cận mở affine của $g(s')$ trong $S$, và cho $W$ là một lân cận mở affine của $s'$ chứa trong $U'\cap g^{-1}(T)$;
chỉ cần chứng minh rằng $f'^{-1}(W)$ là quasi-compact;
nói cách khác, ta có thể quy về việc chứng minh rằng, khi $S$ và $S'$ là \emph{affine}, không gian nền của $X\times_S S'$ là quasi-compact.
Nhưng vì khi đó, theo giả thiết, $X$ là hợp hữu hạn của các tập mở affine $V_j$, nên $X\times_S S'$ là hợp của các không gian nền của các lược đồ affine $V_j\times_S S'$ (\sref{I.3.2.2} và \sref{I.3.2.7}), điều này chứng minh mệnh đề.
\end{proof}

Ta cũng lưu ý rằng, nếu $X=X'\sqcup X''$ là tổng của hai tiền lược đồ, thì một cấu xạ $f:X\to Y$ là quasi-compact khi và chỉ khi các hạn chế của nó lên $X'$ và $X''$ đều quasi-compact.

\begin{proposition}[6.6.5]
\label{I.6.6.5}
Cho $f:X\to Y$ là một cấu xạ quasi-compact.
Để $f$ là một cấu xạ trội, điều kiện cần và đủ là, với mọi điểm tổng quát $y$ của một thành phần bất khả quy của $Y$, $f^{-1}(y)$ chứa điểm tổng quát của một thành phần bất khả quy của $X$.
\end{proposition}

\begin{proof}
Hiển nhiên điều kiện là đủ (ngay cả khi không giả sử $f$ quasi-compact).
Để thấy điều kiện là cần, xét một lân cận mở affine $U$ của $y$;
$f^{-1}(U)$ là quasi-compact, do đó là một hợp \emph{hữu hạn} của các tập mở affine $V_i$, và giả thiết $f$ trội kéo theo $y$ thuộc bao đóng \emph{trong $U$} của một trong các $f(V_i)$.
Hiển nhiên có thể giả sử $X$ và $Y$ rút gọn;
vì bao đóng trong $X$ của một thành phần bất khả quy của $V_i$ là một thành phần bất khả quy của $X$ \sref[0]{0.2.1.6}, ta có thể thay $X$ bởi $V_i$ và $Y$ bởi tiền lược đồ con đóng rút gọn của $U$ có không gian nền là $\overline{f(V_i)}\cap U$ \sref{I.5.2.1}, và như vậy quy về chứng minh mệnh đề khi $X=\Spec(A)$ và $Y=\Spec(B)$ là affine và rút gọn.
Vì $f$ trội, khi đó $B$ là một vành con của $A$ \sref{I.1.2.7}, và mệnh đề suy ra từ việc mọi iđêan nguyên tố tối tiểu của $B$ là giao của $B$ với một iđêan nguyên tố tối tiểu của $A$ \sref[0]{0.1.5.8}.
\end{proof}

\begin{proposition}[6.6.6]
\label{I.6.6.6}
\medskip\noindent
\begin{enumerate}
  \item[{\rm(i)}] Mọi phép nhúng địa phương đều là kiểu hữu hạn địa phương.
  \item[{\rm(ii)}] Nếu hai cấu xạ $f:X\to Y$ và $g:Y\to Z$ đều là kiểu hữu hạn địa phương, thì $g\circ f$ cũng vậy.
  \item[{\rm(iii)}] Nếu $f:X\to Y$ là một $S$-cấu xạ kiểu hữu hạn địa phương, thì $f_{(S')}:X_{(S')}\to Y_{(S')}$ cũng là kiểu hữu hạn địa phương đối với mọi mở rộng $S'\to S$ của tiền lược đồ cơ sở.
  \item[{\rm(iv)}] Nếu $f:X\to X'$ và $g:Y\to Y'$ là các $S$-cấu xạ kiểu hữu hạn địa phương, thì $f\times_S g$ là kiểu hữu hạn địa phương.
  \item[{\rm(v)}] Nếu hợp thành $g\circ f$ của hai cấu xạ là kiểu hữu hạn địa phương, thì $f$ cũng là kiểu hữu hạn địa phương.
  \item[{\rm(vi)}] Nếu một cấu xạ $f$ là kiểu hữu hạn địa phương, thì $f_\red$ cũng vậy.
\end{enumerate}
\end{proposition}

\begin{proof}
Theo \sref{I.5.5.12}, chỉ cần chứng minh (i), (ii) và (iii).
Nếu $j:X\to Y$ là một phép nhúng địa phương, thì với mọi $x\in X$ tồn tại một lân cận mở $V$ của $j(x)$ trong $Y$ và một lân cận mở $U$ của $x$ trong $X$ sao cho hạn chế của $j$ trên $U$ là một phép nhúng đóng $U\to V$ \sref{I.4.5.1}; do đó hạn chế này là kiểu hữu hạn.
Để chứng minh (ii), xét một điểm $x\in X$;
theo giả thiết, tồn tại một lân cận mở $W$ của $g(f(x))$ và một lân cận mở $V$ của $f(x)$ sao cho $g(V)\subset W$ và $V$ kiểu hữu hạn trên $W$;
hơn nữa, $f^{-1}(V)$ kiểu hữu hạn địa phương trên $V$ \sref{I.6.6.2}, nên tồn tại một lân cận mở $U$ của $x$ mà
\oldpage[I]{155}
nằm trong $f^{-1}(V)$ và kiểu hữu hạn trên $V$;
do đó $g(f(U))\subset W$ và $U$ kiểu hữu hạn trên $W$ \sref{I.6.3.4}[ii].
Cuối cùng, để chứng minh (iii), ta có thể quy về trường hợp $Y=S$ \sref{I.3.3.11};
với mọi $x'\in X'=X_{(S')}$, gọi $x$ là ảnh của $x'$ trong $X$, $s$ là ảnh của $x$ trong $S$, $T$ là một lân cận mở của $s$, $T'$ là ảnh ngược của $T$ trong $S'$, và $U$ là một lân cận mở của $x$ có ảnh nằm trong $T$ và kiểu hữu hạn trên $T$;
khi đó $U\times_S T'=U\times_T T'$ là một lân cận mở của $x'$ \sref{I.3.2.7} kiểu hữu hạn trên $T'$ \sref{I.6.3.4}[iv].
\end{proof}

\begin{corollary}[6.6.7]
\label{I.6.6.7}
Cho $X$ và $Y$ là hai $S$-tiền lược đồ kiểu hữu hạn địa phương trên $S$.
Nếu $S$ là Noether địa phương, thì $X\times_S Y$ là Noether địa phương.
\end{corollary}

\begin{proof}
Thật vậy, vì $S$ là Noether địa phương và $X$ kiểu hữu hạn địa phương trên $S$, nên $X$ là Noether địa phương; hơn nữa, $X\times_S Y$ kiểu hữu hạn địa phương trên $X$, do đó cũng là Noether địa phương.
\end{proof}

\begin{remark}[6.6.8]
\label{I.6.6.8}
Mệnh đề~\sref{I.6.3.10} và chứng minh của nó lập tức mở rộng sang trường hợp chỉ giả sử cấu xạ $f$ là kiểu hữu hạn địa phương.
Tương tự, các mệnh đề~\sref{I.6.4.2} và \sref{I.6.4.9} vẫn còn đúng khi chỉ giả sử các tiền lược đồ $X$ và $Y$ xuất hiện trong phát biểu của chúng là kiểu hữu hạn địa phương trên trường $K$.
\end{remark}

\section{Ánh xạ hữu tỉ}
\label{section:I.7}

\subsection{Ánh xạ hữu tỉ và hàm hữu tỉ}
\label{subsection:I.7.1}

\begin{env}[7.1.1]
\label{I.7.1.1}
Cho $X$ và $Y$ là hai tiền lược đồ, $U$ và $V$ là hai tập mở trù mật của $X$, và $f$ (tương ứng $g$) là một cấu xạ từ $U$ (tương ứng $V$) vào $Y$; ta nói $f$ và $g$ \emph{tương đương} nếu chúng trùng nhau trên một tập mở trù mật của $U\cap V$.
Vì giao hữu hạn của các tập mở trù mật trong $X$ là một tập mở trù mật trong $X$, hiển nhiên quan hệ này là một \emph{quan hệ tương đương}.
\end{env}

\begin{definition}[7.1.2]
\label{I.7.1.2}
Cho hai tiền lược đồ $X$ và $Y$, ta gọi một ánh xạ hữu tỉ từ $X$ vào $Y$ là một lớp tương đương các cấu xạ từ các tập mở trù mật của $X$ vào $Y$, đối với quan hệ tương đương được định nghĩa trong \sref{I.7.1.1}.
Nếu $X$ và $Y$ là các $S$-tiền lược đồ, ta nói một lớp như vậy là một $S$-ánh xạ hữu tỉ nếu lớp đó có một đại diện là $S$-cấu xạ.
Ta gọi một $S$-tiết diện hữu tỉ của $X$ là mọi $S$-ánh xạ hữu tỉ từ $S$ vào $X$.
Ta gọi một hàm hữu tỉ trên tiền lược đồ $X$ là mọi $X$-tiết diện hữu tỉ của $X$-tiền lược đồ $X\otimes_{\bb{Z}}\bb{Z}[T]$ (trong đó $T$ là một biến bất định).

Theo lạm dụng ngôn ngữ, khi chỉ xét các $S$-tiền lược đồ, ta sẽ nói ``ánh xạ hữu tỉ'' thay cho ``$S$-ánh xạ hữu tỉ'' nếu không thể nảy sinh nhầm lẫn nào.

Cho $f$ là một ánh xạ hữu tỉ từ $X$ vào $Y$, và cho $U$ là một tập mở của $X$; nếu $f_1$ và $f_2$ là hai cấu xạ thuộc lớp của $f$, được xác định tương ứng trên các tập mở trù mật $V$ và $W$ của $X$, thì các hạn chế $f_1|(U\cap V)$ và $f_2|(U\cap W)$ trùng nhau trên $U\cap V\cap W$, là tập trù mật trong $U$; do đó lớp các cấu xạ $f$ xác định một ánh xạ hữu tỉ từ $U$ vào $Y$, được gọi là \emph{hạn chế của $f$ lên $U$}, và được ký hiệu bởi $f|U$.

Nếu ta cho mỗi $S$-cấu xạ $f:X\to Y$ tương ứng với $S$-ánh xạ hữu tỉ
\oldpage[I]{156}
mà $f$ thuộc vào, ta thu được một ánh xạ chính tắc từ $\Hom_S(X,Y)$ vào tập hợp các $S$-ánh xạ hữu tỉ từ $X$ vào $Y$.
Ký hiệu $\Gamma_\mathrm{rat}(X/Y)$ là tập hợp các $Y$-tiết diện hữu tỉ của $X$; do đó ta có một ánh xạ chính tắc $\Gamma(X/Y)\to\Gamma_\mathrm{rat}(X/Y)$.
Hơn nữa, hiển nhiên rằng, nếu $X$ và $Y$ là hai $S$-tiền lược đồ, thì tập hợp các $S$-ánh xạ hữu tỉ từ $X$ vào $Y$ được đồng nhất chính tắc với $\Gamma_\mathrm{rat}((X\times_S Y)/X)$ \sref{I.3.3.14}.
\end{definition}

\begin{env}[7.1.3]
\label{I.7.1.3}
Từ \sref{I.7.1.2} và \sref{I.3.3.14} cũng suy ra ngay rằng các hàm hữu tỉ trên $X$ được đồng nhất chính tắc với các \emph{lớp tương đương các tiết diện của bó cấu trúc $\sh{O}_X$} trên các tập mở trù mật của $X$, trong đó hai tiết diện như vậy tương đương nếu chúng trùng nhau trên một tập mở trù mật của $X$ nằm trong giao các tập xác định của chúng.
Nói riêng, suy ra rằng các hàm hữu tỉ trên $X$ lập thành một \emph{vành} $R(X)$.
\end{env}

\begin{env}[7.1.4]
\label{I.7.1.4}
Khi $X$ là một tiền lược đồ \emph{bất khả quy}, mọi tập mở không rỗng của $X$ đều trù mật trong $X$; ta cũng có thể nói các tập mở không rỗng của $X$ là các \emph{lân cận mở của điểm tổng quát} $x$ của $X$.
Khi đó, nói hai cấu xạ từ các tập mở không rỗng của $X$ vào $Y$ là tương đương có nghĩa là chúng có \emph{cùng một mầm} tại điểm $x$.
Nói cách khác, các ánh xạ hữu tỉ (tương ứng các $S$-ánh xạ hữu tỉ) $X\to Y$ được đồng nhất với các \emph{mầm của cấu xạ} (tương ứng \emph{$S$-cấu xạ}) từ các tập mở không rỗng của $X$ vào $Y$ tại điểm tổng quát $x$ của $X$.
Nói riêng:
\end{env}

\begin{proposition}[7.1.5]
\label{I.7.1.5}
Nếu $X$ là một tiền lược đồ bất khả quy, thì vành $R(X)$ các hàm hữu tỉ trên $X$ được đồng nhất chính tắc với vành địa phương $\sh{O}_x$ tại điểm tổng quát $x$ của $X$.
Đó là một vành địa phương có chiều~0, và do đó là một vành địa phương Artin khi $X$ là noetherian; đó là một trường khi $X$ nguyên, và khi thêm nữa $X$ là một lược đồ affine, nó được đồng nhất với trường phân thức của $A(X)$.
\end{proposition}

\begin{proof}
Theo phần trên và sự đồng nhất các hàm hữu tỉ với các tiết diện của $\sh{O}_X$ trên một tập mở trù mật của $X$, khẳng định thứ nhất chính là định nghĩa cuống của một bó tại một điểm.
Đối với các khẳng định còn lại, ta có thể quy về trường hợp $X$ affine, với vành $A$; khi đó $\mathfrak{j}_x$ là căn lũy linh của $A$, và do đó $\sh{O}_x$ có chiều~0; nếu $A$ nguyên, thì $\mathfrak{j}_x=(0)$, và do đó $\sh{O}_x$ là trường phân thức của $A$.
Cuối cùng, nếu $A$ là noetherian, ta biết (\cite[p.~127, cor.~4]{I-11}) rằng $\mathfrak{j}_x$ là lũy linh và $\sh{O}_x=A_{\mathfrak{j}_x}$ là Artin.
\end{proof}

Nếu $X$ \emph{nguyên}, thì vành $\sh{O}_z$ nguyên với mọi $z\in X$; mọi tập mở affine $U$ chứa $z$ cũng chứa $x$, và $R(U)$, bằng trường phân thức của $A(U)$, do đó được đồng nhất với $R(X)$; vì thế suy ra $R(X)$ cũng có thể được đồng nhất với \emph{trường phân thức của $\sh{O}_z$}: phép đồng nhất chính tắc $\sh{O}_z$ với một vành con của $R(X)$ gán cho mọi mầm của một tiết diện $s\in\sh{O}_z$ hàm hữu tỉ duy nhất trên $X$, là lớp của một tiết diện của $\sh{O}_X$ (tất yếu được xác định trên một tập mở trù mật của $X$) có mầm $s$ tại điểm $z$.

\begin{env}[7.1.6]
\label{I.7.1.6}
Bây giờ giả sử $X$ có một số \emph{hữu hạn} thành phần bất khả quy $X_i$ ($1\leq i\leq n$) (điều này luôn đúng khi không gian nền của $X$ là \emph{noetherian}); đặt $X'_i$ là tập mở của $X$ thu được bằng cách loại khỏi $X_i$ hợp của các $X_j\cap X_i$ với $j\neq i$; $X'_i$ bất khả quy, điểm tổng quát $x_i$ của nó là điểm tổng quát của $X_i$, và các $X'_i$ đôi một rời nhau, còn hợp của chúng trù mật trong $X$ \sref[0]{0.2.1.6}.
Với
\oldpage[I]{157}
mọi tập mở trù mật $U$ của $X$, $U_i=U\cap X'_i$ là một tập mở không rỗng trù mật trong $X'_i$, các $U_i$ đôi một rời nhau, và do đó $U'=\bigcup_i U_i$ trù mật trong $X$.
Cho một cấu xạ từ $U'$ vào $Y$ tương đương với việc cho, một cách tùy ý, một cấu xạ từ mỗi $U_i$ vào $Y$.
\end{env}

Do đó:

\begin{proposition}[7.1.7]
\label{I.7.1.7}
Cho $X$ và $Y$ là hai tiền lược đồ (tương ứng $S$-tiền lược đồ) sao cho $X$ có một số hữu hạn thành phần bất khả quy $X_i$, với các điểm tổng quát $x_i$ ($1\leq i\leq n$).
Nếu $R_i$ là tập hợp các mầm của cấu xạ (tương ứng $S$-cấu xạ) từ các tập mở của $X$ vào $Y$ tại điểm $x_i$, thì tập hợp các ánh xạ hữu tỉ (tương ứng các $S$-ánh xạ hữu tỉ) từ $X$ vào $Y$ được đồng nhất với tích các $R_i$ ($1\leq i\leq n$).
\end{proposition}

\begin{corollary}[7.1.8]
\label{I.7.1.8}
Cho $X$ là một tiền lược đồ noetherian.
Vành các hàm hữu tỉ trên $X$ là một vành Artin, có các thành phần địa phương là các vành $\sh{O}_{x_i}$ tại các điểm tổng quát $x_i$ của các thành phần bất khả quy của $X$.
\end{corollary}

\begin{corollary}[7.1.9]
\label{I.7.1.9}
Cho $A$ là một vành noetherian và $X=\Spec(A)$.
Nếu $Q$ là phần bù của hợp các iđêan nguyên tố tối tiểu của $A$, thì vành các hàm hữu tỉ trên $X$ được đồng nhất chính tắc với vành phân thức $Q^{-1}A$.
\end{corollary}

Điều này sẽ suy ra từ bổ đề sau:

\begin{lemma}[7.1.9.1]
\label{I.7.1.9.1}
Với một phần tử $f\in A$, tập $D(f)$ trù mật trong $X$ khi và chỉ khi $f\in Q$; mọi tập mở trù mật của $X$ chứa một tập mở dạng $D(f)$, trong đó $f\in Q$.
\end{lemma}

Thật vậy, giả sử bổ đề đã được chứng minh. Vì vành các tiết diện $\Gamma(D(f),\sh{O}_X)$ được đồng nhất với $A_f$ theo \sref{I.1.3.6} và \sref{I.1.3.7}, từ việc các $D(f)$ với $f\in Q$ tạo thành một họ đồng tận trong tập có thứ tự (theo quan hệ bao hàm ngược) của các tập mở trù mật của $X$, cùng với định nghĩa các hàm hữu tỉ trong \sref{I.7.1.1}, suy ra rằng vành các hàm hữu tỉ trên $X$ được đồng nhất với giới hạn quy nạp của các $A_f$ với $f\in Q$ (đối với quan hệ tiền thứ tự ``$g$ là bội của $f$''), tức là với $Q^{-1}A$ theo \sref[0]{0.1.4.5}.

\begin{proof}
Để chứng minh \hyperref[I.7.1.9.1]{(7.1.9.1)}, vẫn dùng $X_i$ ($1\leq i\leq n$) để ký hiệu các thành phần bất khả quy của $X$; nếu $D(f)$ trù mật trong $X$, thì $D(f)\cap X_i\neq\emp$ với $1\leq i\leq n$, và ngược lại; điều này có nghĩa là $f\not\in\mathfrak{p}_i$ với $1\leq i\leq n$, trong đó đặt $\mathfrak{p}_i=\mathfrak{j}(X_i)$, và vì các $\mathfrak{p}_i$ là các iđêan nguyên tố tối tiểu của $A$ theo \hyperref[I.1.1.14]{(1.1.14)}, các điều kiện $f\not\in\mathfrak{p}_i$ ($1\leq i\leq n$) tương đương với $f\in Q$, từ đó suy ra khẳng định thứ nhất của bổ đề.
Đối với khẳng định còn lại, nếu $U$ là một tập mở trù mật của $X$, thì phần bù của $U$ là một tập dạng $V(\mathfrak{a})$, trong đó $\mathfrak{a}$ là một iđêan không nằm trong bất kỳ $\mathfrak{p}_i$ nào; do đó, nó cũng không nằm trong hợp của chúng (\cite[p.~13]{I-10}), nên tồn tại một $f\in\mathfrak{a}$ thuộc $Q$; vì thế $D(f)\subset U$, và chứng minh hoàn tất.
\end{proof}

\begin{env}[7.1.10]
\label{I.7.1.10}
Giả sử lại rằng $X$ bất khả quy, với điểm tổng quát $x$.
Vì mọi tập mở không rỗng $U$ của $X$ đều chứa $x$, và do đó cũng chứa mọi $z\in X$ sao cho $x\in\overline{\{z\}}$, mọi cấu xạ $U\to Y$ đều có thể được hợp thành với cấu xạ chính tắc $\Spec(\sh{O}_x)\to X$ theo \sref{I.2.4.1}; và hai cấu xạ vào $Y$, được xác định trên hai tập mở không rỗng của $X$, trùng nhau trên một tập mở không rỗng của $X$, sẽ cho cùng một cấu xạ $\Spec(\sh{O}_x)\to Y$ sau khi hợp thành.
Nói cách khác, mỗi ánh xạ hữu tỉ từ $X$ vào $Y$ do đó xác định duy nhất một cấu xạ $\Spec(\sh{O}_x)\to Y$.
\end{env}

\oldpage[I]{158}
\begin{proposition}[7.1.11]
\label{I.7.1.11}
Cho $X$ và $Y$ là hai $S$-tiền lược đồ; giả sử $X$ bất khả quy, với điểm tổng quát $x$, và $Y$ thuộc kiểu hữu hạn trên $S$.
Hai $S$-ánh xạ hữu tỉ bất kỳ từ $X$ vào $Y$ mà tương ứng với cùng một $S$-cấu xạ $\Spec(\sh{O}_x)\to Y$ thì trùng nhau.
Nếu thêm rằng $S$ là Noether địa phương, thì với mỗi $S$-cấu xạ từ $\Spec(\sh{O}_x)$ vào $Y$ tương ứng duy nhất một $S$-ánh xạ hữu tỉ từ $X$ vào $Y$.
\end{proposition}

\begin{proof}
Vì mọi tập mở không rỗng của $X$ đều trù mật trong $X$, điều này suy ra từ \sref{I.6.5.1}.
\end{proof}

\begin{corollary}[7.1.12]
\label{I.7.1.12}
Giả sử $S$ là Noether địa phương và các giả thiết còn lại của \sref{I.7.1.11} đều được thỏa mãn.
Khi đó, các $S$-ánh xạ hữu tỉ từ $X$ vào $Y$ được đồng nhất với các điểm của $S$-tiền lược đồ $Y$ nhận giá trị trong $S$-tiền lược đồ $\Spec(\sh{O}_x)$.
\end{corollary}

\begin{proof}
Đây chính là \sref{I.7.1.11}, với thuật ngữ được đưa vào trong \sref{I.3.4.1}.
\end{proof}

\begin{corollary}[7.1.13]
\label{I.7.1.13}
Giả sử các điều kiện của \sref{I.7.1.12} được thỏa mãn.
Gọi $s$ là ảnh của $x$ trong $S$.
Việc cho một $S$-ánh xạ hữu tỉ từ $X$ vào $Y$ tương đương với việc cho một điểm $y$ của $Y$ nằm trên $s$ và một $\sh{O}_s$-đồng cấu địa phương $\sh{O}_y\to\sh{O}_x=R(X)$.
\end{corollary}

\begin{proof}
Điều này suy ra từ \sref{I.7.1.11} và \sref{I.2.4.4}.
\end{proof}

Đặc biệt:
\begin{corollary}[7.1.14]
\label{I.7.1.14}
Dưới các điều kiện của \sref{I.7.1.12}, các $S$-ánh xạ hữu tỉ từ $X$ vào $Y$ chỉ phụ thuộc (với $Y$ đã cho) vào $S$-tiền lược đồ $\Spec(\sh{O}_x)$ và, đặc biệt, vẫn không đổi khi thay $X$ bằng $\Spec(\sh{O}_z)$, với mọi $z\in X$.
\end{corollary}

\begin{proof}
Thật vậy, vì $z\in\overline{\{x\}}$, $x$ là điểm tổng quát của $Z=\Spec(\sh{O}_z)$, và $\sh{O}_{X,x}=\sh{O}_{Z,x}$.
\end{proof}

Khi $X$ nguyên, $R(X)=\sh{O}_x=\kres(x)$ là một trường \sref{I.7.1.5}; khi đó các hệ quả trên có dạng sau:
\begin{corollary}[7.1.15]
\label{I.7.1.15}
Giả sử các điều kiện của \sref{I.7.1.12} được thỏa mãn và thêm rằng $X$ nguyên.
Gọi $s$ là ảnh của $x$ trong $S$.
Khi đó các $S$-ánh xạ hữu tỉ từ $X$ vào $Y$ được đồng nhất với các điểm hình học của $Y\otimes_S\kres(s)$ nhận giá trị trong trường mở rộng $R(X)$ của $\kres(s)$; nói cách khác, mỗi ánh xạ tương đương với dữ kiện gồm một điểm $y\in Y$ nằm trên $s$ và một $\kres(s)$-đơn cấu từ $\kres(y)$ vào $\kres(x)=R(X)$.
\end{corollary}

\begin{proof}
Thật vậy, các điểm của $Y$ nằm trên $s$ được đồng nhất với các điểm của $Y\otimes_S\kres(s)$ \sref{I.3.6.3}, còn các $\sh{O}_s$-đồng cấu địa phương $\sh{O}_y\to R(X)$ được đồng nhất với các $\kres(s)$-đơn cấu $\kres(y)\to R(X)$.
\end{proof}

Cụ thể hơn:
\begin{corollary}[7.1.16]
\label{I.7.1.16}
Cho $k$ là một trường, và cho $X$, $Y$ là hai tiền lược đồ đại số trên $k$ \sref{I.6.4.1}; giả sử thêm rằng $X$ nguyên.
Khi đó các $k$-ánh xạ hữu tỉ từ $X$ vào $Y$ được đồng nhất với các điểm hình học của $Y$ nhận giá trị trong trường mở rộng $R(X)$ của $k$ \sref{I.3.4.4}.
\end{corollary}

\subsection{Miền xác định của một ánh xạ hữu tỉ}
\label{subsection:I.7.2}
\label{I.7.2}

\begin{env}[7.2.1]
\label{I.7.2.1}
Cho $X$ và $Y$ là hai tiền lược đồ, và cho $f$ là một ánh xạ hữu tỉ từ $X$ vào $Y$.
Ta nói $f$ \emph{được xác định tại một điểm $x\in X$} nếu tồn tại một tập mở trù mật $U$ của $X$ chứa $x$ và một cấu xạ $U\to Y$ thuộc lớp tương đương $f$.
Ta gọi tập hợp các điểm $x\in X$ mà tại đó $f$ được xác định là \emph{miền xác định của $f$}; hiển nhiên đây là một tập mở trù mật của $X$.
\end{env}

\oldpage[I]{159}
\begin{proposition}[7.2.2]
\label{I.7.2.2}
Cho $X$ và $Y$ là hai $S$-tiền lược đồ sao cho $X$ rút gọn và $Y$ tách trên $S$.
Cho $f$ là một $S$-ánh xạ hữu tỉ từ $X$ vào $Y$, có miền xác định $U_0$.
Khi đó tồn tại duy nhất một $S$-cấu xạ $U_0\to Y$ thuộc lớp của $f$.
\end{proposition}

Vì, đối với mọi cấu xạ $U\to Y$ thuộc lớp của $f$, ta nhất thiết có $U\subset U_0$, hiển nhiên mệnh đề sẽ là hệ quả của bổ đề sau:

\begin{lemma}[7.2.2.1]
\label{I.7.2.2.1}
Dưới các giả thiết của \sref{I.7.2.2}, cho $U_1$ và $U_2$ là hai tập mở trù mật của $X$, và cho $f_i:U_i\to Y$ ($i=1,2$) là hai $S$-cấu xạ sao cho tồn tại một tập mở $V\subset U_1\cap U_2$, trù mật trong $X$, trên đó $f_1$ và $f_2$ trùng nhau.
Khi đó $f_1$ và $f_2$ trùng nhau trên $U_1\cap U_2$.
\end{lemma}

\begin{proof}
Hiển nhiên ta có thể quy về trường hợp $X=U_1=U_2$.
Vì $X$ (và do đó $V$) rút gọn, $X$ là tiền lược đồ con đóng nhỏ nhất của $X$ chứa $V$ \sref{I.5.2.2}.
Đặt $g=(f_1,f_2)_S:X\to Y\times_S Y$; vì theo giả thiết, đường chéo $T=\Delta_Y(Y)$ là một tiền lược đồ con đóng của $Y\times_S Y$, nên $Z=g^{-1}(T)$ là một tiền lược đồ con đóng của $X$ \sref{I.4.4.1}.
Nếu $h:V\to Y$ là hạn chế chung của $f_1$ và $f_2$ lên $V$, thì hạn chế của $g$ lên $V$ là $g'=(h,h)_S$, phân tích thành $g'=\Delta_Y\circ h$; vì $\Delta_Y^{-1}(T)=Y$, ta có $g'^{-1}(T)=V$, do đó $Z$ là một tiền lược đồ con đóng của $X$ cảm sinh $V$, vì thế chứa $V$, suy ra $Z=X$.
Từ hệ thức $g^{-1}(T)=X$, theo \sref{I.4.4.1} suy ra $g$ phân tích thành $\Delta_Y\circ f$, trong đó $f$ là một cấu xạ $X\to Y$; theo định nghĩa của cấu xạ đường chéo, điều này kéo theo $f_1=f_2=f$.
\end{proof}

Hiển nhiên cấu xạ $U_0\to Y$ được xác định trong \sref{I.7.2.2} là cấu xạ duy nhất thuộc lớp của $f$ mà \emph{không thể mở rộng} thành một cấu xạ từ một tập mở của $X$ chứa thực sự $U_0$.
\emph{Dưới các giả thiết của \sref{I.7.2.2}}, do đó ta có thể \emph{đồng nhất} các ánh xạ hữu tỉ từ $X$ vào $Y$ với các cấu xạ được xác định trên các tập mở trù mật của $X$, có đích là $Y$ và \emph{không thể mở rộng} sang bất kỳ tập mở nào chứa thực sự miền xác định của chúng.
Với phép đồng nhất này, Mệnh đề~\sref{I.7.2.2} kéo theo:

\begin{corollary}[7.2.3]
\label{I.7.2.3}
Dưới các giả thiết của \sref{I.7.2.2} về $X$ và $Y$, cho $U$ là một tập mở trù mật của $X$.
Khi đó có một tương ứng song ánh chính tắc giữa các $S$-cấu xạ từ $U$ vào $Y$ và các $S$-ánh xạ hữu tỉ từ $X$ vào $Y$ được xác định tại mọi điểm của $U$.
\end{corollary}

\begin{proof}
Theo \sref{I.7.2.2}, với mọi $S$-cấu xạ $f$ từ $U$ vào $Y$, tồn tại duy nhất một $S$-ánh xạ hữu tỉ $\overline{f}$ từ $X$ vào $Y$ mở rộng $f$.
\end{proof}

\begin{corollary}[7.2.4]
\label{I.7.2.4}
Cho $S$ là một lược đồ, $X$ là một $S$-tiền lược đồ rút gọn, $Y$ là một $S$-lược đồ, và $f:U\to Y$ là một $S$-cấu xạ từ một tập mở trù mật $U$ của $X$ vào $Y$.
Nếu $\overline{f}$ là $\bb{Z}$-ánh xạ hữu tỉ từ $X$ vào $Y$ mở rộng $f$, thì $\overline{f}$ là một $S$-cấu xạ (và do đó là $S$-ánh xạ hữu tỉ từ $X$ vào $Y$ mở rộng $f$).
\end{corollary}

\begin{proof}
Thật vậy, nếu $\vphi:X\to S$ và $\psi:Y\to S$ là các cấu xạ cấu trúc, $U_0$ là miền xác định của $\overline{f}$, và $j$ là phép nhúng $U_0\to X$, thì chỉ cần chứng minh $\psi\circ\overline{f}=\vphi\circ j$; điều này suy ra từ \sref{I.7.2.2.1}, vì $f$ là một $S$-cấu xạ.
\end{proof}

\begin{corollary}[7.2.5]
\label{I.7.2.5}
Cho $X$ và $Y$ là hai $S$-tiền lược đồ; giả sử $X$ rút gọn và $X$, $Y$ tách trên $S$.
Cho $p:Y\to X$ là một $S$-cấu xạ (làm cho $Y$ trở thành một $X$-tiền lược đồ), $U$ là một tập mở trù mật của $X$, và $f$ là một $U$-nhát cắt của $Y$; khi đó ánh xạ hữu tỉ $\overline{f}$ từ $X$ vào $Y$ mở rộng $f$ là một $X$-nhát cắt hữu tỉ của $Y$.
\end{corollary}

\oldpage[I]{160}
\begin{proof}
Cần chứng minh rằng $p\circ\overline{f}$ là phép đồng nhất trên miền xác định của $\overline{f}$; vì $X$ tách trên $S$, điều này lại suy ra từ \sref{I.7.2.2.1}.
\end{proof}

\begin{corollary}[7.2.6]
\label{I.7.2.6}
Cho $X$ là một tiền lược đồ rút gọn và $U$ là một tập mở trù mật của $X$.
Khi đó có một tương ứng song ánh chính tắc giữa các tiết diện của $\sh{O}_X$ trên $U$ và các hàm hữu tỉ trên $X$ được xác định tại mọi điểm của $U$.
\end{corollary}

\begin{proof}
Theo \sref{I.7.2.3}, \sref{I.7.1.2} và \sref{I.7.1.3}, chỉ cần nhận xét rằng $X$-tiền lược đồ $X\otimes_{\bb{Z}}\bb{Z}[T]$ tách trên $X$ theo \sref{I.5.5.1}[iv].
\end{proof}

\begin{corollary}[7.2.7]
\label{I.7.2.7}
Cho $Y$ là một tiền lược đồ rút gọn, $f:X\to Y$ là một cấu xạ tách, $U$ là một tập mở trù mật của $Y$, $g:U\to f^{-1}(U)$ là một $U$-nhát cắt của $f^{-1}(U)$, và $Z$ là tiền lược đồ con rút gọn của $X$ có không gian nền là $\overline{g(U)}$ \sref{I.5.2.1}.
Để $g$ là hạn chế của một $Y$-nhát cắt của $X$ \emph{(nói cách khác, theo \sref{I.7.2.5}, để ánh xạ hữu tỉ từ $Y$ vào $X$ mở rộng $g$ được xác định khắp nơi)}, điều kiện cần và đủ là hạn chế của $f$ trên $Z$ là một đẳng cấu từ $Z$ lên $Y$.
\end{corollary}

\begin{proof}
Hạn chế của $f$ trên $f^{-1}(U)$ là một cấu xạ tách \sref{I.5.5.1}[i], do đó $g$ là một phép nhúng đóng \sref{I.5.4.6}, và vì thế $g(U)$\footnote{\emph{[Dịch.] Nguồn tiếng Pháp in $g(V)$ tại đây, nhưng $V$ không được định nghĩa; đẳng thức và dữ liệu của hệ quả đòi hỏi $g(U)$.}} $=Z\cap f^{-1}(U)$, còn tiền lược đồ con do $Z$ cảm sinh trên tập mở $g(U)$ của $Z$ trùng với tiền lược đồ con đóng của $f^{-1}(U)$ ứng với $g$ \sref{I.5.2.1}.
Khi đó hiển nhiên điều kiện trong phát biểu là đủ, vì nếu điều kiện ấy được thỏa mãn, nếu $f_Z:Z\to Y$ là hạn chế của $f$ trên $Z$, và nếu $\overline{g}:Y\to Z$ là đẳng cấu ngược, thì $\overline{g}$ mở rộng $g$.
Ngược lại, nếu $g$ là hạn chế trên $U$ của một $Y$-nhát cắt $h$ của $X$, thì $h$ là một phép nhúng đóng \sref{I.5.4.6}; do đó $h(Y)$ là đóng và, vì nó nằm trong $Z$, nó bằng $Z$; từ \sref{I.5.2.1} suy ra rằng $h$ nhất thiết là một đẳng cấu từ $Y$ lên tiền lược đồ con đóng $Z$ của $X$.
\end{proof}

\begin{env}[7.2.8]
\label{I.7.2.8}
Cho $X$ và $Y$ là hai $S$-tiền lược đồ, với $X$ rút gọn và $Y$ tách trên $S$.
Cho $f$ là một $S$-ánh xạ hữu tỉ từ $X$ vào $Y$, và cho $x$ là một điểm của $X$; ta có thể hợp thành $f$ với $S$-cấu xạ chính tắc $\Spec(\sh{O}_x)\to X$ \sref{I.2.4.1}, miễn là giao của $\Spec(\sh{O}_x)$ với miền xác định của $f$ trù mật trong $\Spec(\sh{O}_x)$ (được đồng nhất với tập hợp các $z\in X$ sao cho $x\in\overline{\{z\}}$ \sref{I.2.4.2}).
Điều này xảy ra trong các trường hợp sau:
\begin{enumerate}
  \item[1.º] $X$ là \emph{bất khả quy} (và do đó \emph{nguyên}), vì khi ấy điểm tổng quát $\xi$ của $X$ là điểm tổng quát của $\Spec(\sh{O}_x)$; vì miền xác định $U$ của $f$ chứa $\xi$, nên $U\cap\Spec(\sh{O}_x)$ chứa $\xi$ và do đó trù mật trong $\Spec(\sh{O}_x)$.
  \item[2.º] $X$ là \emph{Noether địa phương}; khi đó khẳng định của ta suy ra từ kết quả sau:
\end{enumerate}
\end{env}

\begin{lemma}[7.2.8.1]
\label{I.7.2.8.1}
Cho $X$ là một tiền lược đồ có không gian nền Noether địa phương, và cho $x$ là một điểm của $X$.
Các thành phần bất khả quy của $\Spec(\sh{O}_x)$ là các giao với $\Spec(\sh{O}_x)$ của những thành phần bất khả quy của $X$ chứa $x$.
Điều kiện cần và đủ để một tập mở $U\subset X$ thỏa $U\cap\Spec(\sh{O}_x)$ trù mật trong $\Spec(\sh{O}_x)$ là nó có giao không rỗng với mỗi thành phần bất khả quy của $X$ chứa $x$ \emph{(điều này đặc biệt xảy ra nếu $U$ \emph{trù mật} trong $X$)}.
\end{lemma}

\begin{proof}
Khẳng định thứ hai hiển nhiên suy ra từ khẳng định thứ nhất, nên chỉ cần chứng minh khẳng định thứ nhất.
Vì $\Spec(\sh{O}_x)$ nằm trong mọi tập mở affine $U$ chứa $x$, và vì các thành phần bất khả quy của $U$ chứa $x$ là các giao với $U$ của những thành phần bất khả quy của $X$ chứa $x$ \sref[0]{0.2.1.6}, ta có thể giả sử $X$ là affine, với vành $A$.
Vì các iđêan nguyên tố của $A_x$ tương ứng song ánh với các iđêan nguyên tố
\oldpage[I]{161}
của $A$ nằm trong $\mathfrak{j}_x$ \sref[0]{0.1.2.6}, nên các iđêan nguyên tố tối tiểu của $A_x$ tương ứng với các iđêan nguyên tố tối tiểu của $A$ nằm trong $\mathfrak{j}_x$; do đó suy ra bổ đề \sref{I.1.1.14}.
\end{proof}

Với những điều đó, giả sử ta ở một trong hai trường hợp nêu trong \sref{I.7.2.8}.
Nếu $U$ là miền xác định của $S$-ánh xạ hữu tỉ $f$, ký hiệu bởi $f'$ ánh xạ hữu tỉ từ $\Spec(\sh{O}_x)$ vào $Y$ mà, theo \sref{I.2.4.2}, trùng với $f$ trên $U\cap\Spec(\sh{O}_x)$; ta nói ánh xạ hữu tỉ này được \emph{cảm sinh bởi $f$}.

\begin{proposition}[7.2.9]
\label{I.7.2.9}
Cho $S$ là một tiền lược đồ Noether địa phương, $X$ là một $S$-tiền lược đồ rút gọn và $Y$ là một $S$-lược đồ kiểu hữu hạn.
Giả sử thêm rằng $X$ là bất khả quy hoặc Noether địa phương.
Cho $f$ là một $S$-ánh xạ hữu tỉ từ $X$ vào $Y$, và $x$ là một điểm của $X$.
Điều kiện cần và đủ để $f$ được xác định tại điểm $x$ là ánh xạ hữu tỉ $f'$ từ $\Spec(\sh{O}_x)$ vào $Y$, được cảm sinh bởi $f$ \sref{I.7.2.8}, là một cấu xạ.
\end{proposition}

\begin{proof}
Điều kiện này hiển nhiên là cần thiết (vì $\Spec(\sh{O}_x)$ nằm trong mọi tập mở chứa $x$); ta sẽ chứng minh rằng nó cũng là điều kiện đủ.
Theo \sref{I.6.5.1}, tồn tại một lân cận mở $U$ của $x$ trong $X$ và một $S$-cấu xạ $g$ từ $U$ vào $Y$ cảm sinh ra $f'$ trên $\Spec(\sh{O}_x)$.
Nếu $X$ bất khả quy, thì $U$ trù mật trong $X$ và, theo \sref{I.7.2.3}, ta có thể xem $g$ như một $S$-ánh xạ hữu tỉ.
Hơn nữa, điểm tổng quát của $X$ thuộc $\Spec(\sh{O}_x)$ và miền xác định của $f$, nên $f$ và $g$ trùng nhau tại điểm đó, và do đó trên một tập mở không rỗng của $X$ \sref{I.6.5.1}.
Nhưng vì $f$ và $g$ là các $S$-ánh xạ hữu tỉ, chúng trùng nhau \sref{I.7.2.3}; do đó $f$ được xác định tại $x$.

Bây giờ giả sử $X$ là Noether địa phương; ta có thể giả sử $U$ là Noether.
Khi đó chỉ có hữu hạn thành phần bất khả quy $X_i$ của $X$ chứa $x$ \sref{I.7.2.8.1}, và ta có thể giả sử chúng là những thành phần duy nhất có giao không rỗng với $U$, bằng cách thay $U$, nếu cần, bởi một tập mở nhỏ hơn (vì chỉ có hữu hạn thành phần bất khả quy của $X$ có giao không rỗng với $U$, bởi $U$ là Noether).
Như trên, khi đó ta thấy $f$ và $g$ trùng nhau trên một tập mở không rỗng của mỗi $X_i$.
Vì mỗi $X_i$ nằm trong $\overline{U}$, xét cấu xạ $f_1$, được xác định trên một tập mở trù mật của $U\cup(X\setmin\overline{U})$, bằng $g$ trên $U$ và bằng $f$ trên giao của $X\setmin\overline{U}$ với miền xác định của $f$.
Vì $U\cup(X\setmin\overline{U})$ trù mật trong $X$, nên $f_1$ và $f$ trùng nhau trên một tập mở trù mật của $X$; và vì $f$ là một ánh xạ hữu tỉ, $f$ là một mở rộng của $f_1$ \sref{I.7.2.3}, do đó được xác định tại điểm $x$.
\end{proof}

\subsection{Bó các hàm hữu tỉ}
\label{subsection:I.7.3}

\begin{env}[7.3.1]
\label{I.7.3.1}
Cho $X$ là một tiền lược đồ.
Với mỗi tập mở $U\subset X$, ký hiệu $R(U)$ là vành các hàm hữu tỉ trên $U$ \sref{I.7.1.3}; đây là một $\Gamma(U,\sh{O}_X)$-đại số.
Hơn nữa, nếu $V\subset U$ là một tập mở khác của $X$, thì mỗi tiết diện của $\sh{O}_X$ trên một tập mở trù mật trong $U$, sau khi hạn chế xuống $V$, cho một tiết diện trên một tập mở trù mật trong $V$; và nếu hai tiết diện trùng nhau trên một tập mở trù mật trong $U$, thì các hạn chế của chúng xuống $V$ trùng nhau trên một tập mở trù mật trong $V$.
Do đó ta xác định được một di-đồng cấu đại số $\rho_{V,U}:R(U)\to R(V)$, và hiển nhiên rằng, nếu $U\supset V\supset W$ là các tập mở của $X$, thì $\rho_{W,U}=\rho_{W,V}\circ\rho_{V,U}$; vì vậy các $R(U)$ xác định một \emph{tiền bó} các đại số trên $X$.
\end{env}

\oldpage[I]{162}
\begin{definition}[7.3.2]
\label{I.7.3.2}
Ta gọi bó các hàm hữu tỉ trên một tiền lược đồ $X$, ký hiệu $\sh{R}(X)$, là $\sh{O}_X$-đại số liên kết với tiền bó được tạo bởi các $R(U)$.
\end{definition}

Với mọi tiền lược đồ $X$ và mọi tập mở $U\subset X$, hiển nhiên bó hạn chế $\sh{R}(X)|U$ chính là $\sh{R}(U)$.

\begin{proposition}[7.3.3]
\label{I.7.3.3}
Cho $X$ là một tiền lược đồ sao cho họ $(X_\lambda)$ các thành phần bất khả quy của nó là hữu hạn địa phương (điều này đặc biệt đúng khi không gian nền của $X$ là noetherian địa phương).
Khi đó $\sh{O}_X$-môđun $\sh{R}(X)$ là quasi-coherent, và với mỗi tập mở $U$ của $X$ chỉ giao với một số hữu hạn các thành phần $X_\lambda$, $R(U)$ bằng $\Gamma(U,\sh{R}(X))$ và có thể được đồng nhất chính tắc với tích trực tiếp của các vành địa phương tại các điểm tổng quát của những $X_\lambda$ thỏa $U\cap X_\lambda\neq\emp$.
\end{proposition}

\begin{proof}
Rõ ràng ta có thể quy về trường hợp $X$ chỉ có hữu hạn thành phần bất khả quy $X_i$, với các điểm tổng quát $x_i$ ($1\leq i\leq n$).
Khi đó, việc $R(U)$ được đồng nhất chính tắc với tích trực tiếp của các $\sh{O}_{x_i}=R(X_i)$ thỏa $U\cap X_i\neq\emp$ suy ra từ \sref{I.7.1.7}.
Ta sẽ chứng minh tiền bó $U\to R(U)$ thỏa các tiên đề bó; điều này sẽ chứng minh $R(U)=\Gamma(U,\sh{R}(X))$.
Thật vậy, nó thỏa (F1) theo phần trình bày ở trên.
Để thấy nó thỏa (F2), xét một phủ của tập mở $U$ trong $X$ bởi các tập mở $V_\alpha\subset U$; nếu $s_\alpha\in R(V_\alpha)$ sao cho các hạn chế của $s_\alpha$ và $s_\beta$ xuống $V_\alpha\cap V_\beta$ trùng nhau với mọi cặp chỉ số, thì với mỗi chỉ số $i$ thỏa $U\cap X_i\neq\emp$, ta kết luận rằng các thành phần trong $R(X_i)$ của mọi $s_\alpha$ thỏa $V_\alpha\cap X_i\neq\emp$ đều như nhau; ký hiệu thành phần này là $t_i$, hiển nhiên phần tử của $R(U)$ có các thành phần $t_i$ sẽ hạn chế thành $s_\alpha$ trên mỗi $V_\alpha$.
Cuối cùng, để thấy $\sh{R}(X)$ là quasi-coherent, ta có thể quy về trường hợp affine $X=\Spec(A)$; cho $U$ chạy qua các tập mở affine dạng $D(f)$, với $f\in A$, từ phần trên và Định nghĩa~\sref{I.1.3.4} suy ra $\sh{R}(X)=\widetilde{M}$, trong đó $M$ là tổng trực tiếp của các $A$-môđun $A_{x_i}$.
\end{proof}

\begin{corollary}[7.3.4]
\label{I.7.3.4}
Cho $X$ là một tiền lược đồ rút gọn chỉ có hữu hạn thành phần bất khả quy, và cho $X_i$ ($1\leq i\leq n$) là các tiền lược đồ con đóng rút gọn của $X$ có các không gian nền là những thành phần bất khả quy của $X$ \sref{I.5.2.1}.
Nếu $h_i$ là phép nhúng chính tắc $X_i\to X$, thì $\sh{R}(X)$ là tích trực tiếp của các $\sh{O}_X$-đại số $(h_i)_*(\sh{R}(X_i))$.
\end{corollary}

\begin{corollary}[7.3.5]
\label{I.7.3.5}
Nếu $X$ là bất khả quy, thì mọi $\sh{R}(X)$-môđun quasi-coherent $\sh{F}$ đều là một bó đơn.
\end{corollary}

\begin{proof}
Chỉ cần chứng minh rằng mỗi $x\in X$ có một lân cận $U$ sao cho $\sh{F}|U$ là một bó đơn \sref[0]{0.3.6.2}; nói cách khác, ta quy về trường hợp $X$ là affine; hơn nữa, có thể giả sử $\sh{F}$ là cokernel của một đồng cấu $(\sh{R}(X))^{I}\to(\sh{R}(X))^{J}$ \sref[0]{0.5.1.3}, và tất cả quy về việc chứng minh $\sh{R}(X)$ là một bó đơn; nhưng điều này hiển nhiên, vì $\Gamma(U,\sh{R}(X))=R(X)$ đối với mọi tập mở khác rỗng $U$, bởi $U$ chứa điểm tổng quát của $X$.
\end{proof}

\begin{corollary}[7.3.6]
\label{I.7.3.6}
Nếu $X$ là bất khả quy, thì với mọi $\sh{O}_X$-môđun quasi-coherent $\sh{F}$, $\sh{F}\otimes_{\sh{O}_X}\sh{R}(X)$ là một bó đơn; nếu hơn nữa $X$ rút gọn (do đó nguyên), thì $\sh{F}\otimes_{\sh{O}_X}\sh{R}(X)$ đẳng cấu với một bó có dạng $(\sh{R}(X))^{(I)}$.
\end{corollary}

\begin{proof}
Khẳng định thứ hai suy ra từ việc $R(X)$ khi đó là một trường.
\end{proof}

\begin{proposition}[7.3.7]
\label{I.7.3.7}
Giả sử tiền lược đồ $X$ là nguyên địa phương hoặc noetherian
\oldpage[I]{163}
địa phương.
Khi đó $\sh{R}(X)$ là một $\sh{O}_X$-đại số quasi-coherent; nếu hơn nữa $X$ rút gọn (điều này đúng khi $X$ nguyên địa phương), thì đồng cấu chính tắc $\sh{O}_X\to\sh{R}(X)$ là đơn ánh.
\end{proposition}

\begin{proof}
Vì đây là một vấn đề địa phương, khẳng định thứ nhất suy ra từ \sref{I.7.3.3}; khẳng định thứ hai suy ra ngay từ \sref{I.7.2.3}.
\end{proof}

\begin{env}[7.3.8]
\label{I.7.3.8}
\footnote{\emph{[Dịch giả] Văn bản gốc của mục này, được xác nhận là sai, đã được thay thế toàn bộ trong phần Đính chính của EGA~II.}}
Cho $X$ và $Y$ là hai tiền lược đồ \emph{nguyên}; điều này kéo theo rằng $\sh{R}(X)$ (tương ứng $\sh{R}(Y)$) là một $\sh{O}_X$-môđun (tương ứng một $\sh{O}_Y$-môđun) quasi-coherent \sref{I.7.3.3}.
Cho $f:X\to Y$ là một cấu xạ \emph{trội}; khi đó tồn tại một đồng cấu chính tắc giữa các $\sh{O}_X$-môđun
\[
\label{I.7.3.8.1}
  \tau:f^*(\sh{R}(Y))\to\sh{R}(X).
  \tag{7.3.8.1}
\]
\end{env}

\begin{proof}
Trước hết, giả sử $X=\Spec(A)$ và $Y=\Spec(B)$ là affine, trong đó $A$ và $B$ là hai miền nguyên; do đó $f$ tương ứng với một đồng cấu đơn ánh $B\to A$, và đồng cấu này kéo dài thành một đơn cấu $L\to K$ từ trường phân thức $L$ của $B$ vào trường phân thức $K$ của $A$.
Khi đó, đồng cấu \sref{I.7.3.8.1} tương ứng với đồng cấu chính tắc $L\otimes_B A\to K$ \sref{I.1.6.5}.
Trong trường hợp tổng quát, với mỗi cặp tập mở affine khác rỗng $U\subset X$ và $V\subset Y$ sao cho $f(U)\subset V$, ta định nghĩa như trên một đồng cấu $\tau_{U,V}$; ta thấy ngay rằng nếu $U'\subset U$, $V'\subset V$, $f(U')\subset V'$, thì $\tau_{U,V}$ kéo dài $\tau_{U',V'}$, từ đó suy ra khẳng định.
Nếu $x$ và $y$ lần lượt là các điểm tổng quát của $X$ và $Y$, thì ta có $f(x)=y$ và
\[
  (f^*(\sh{R}(Y)))_x=\sh{O}_y\otimes_{\sh{O}_y}\sh{O}_x=\sh{O}_x
\]
\sref[0]{0.4.3.1}; do đó $\tau_x$ là một \emph{đẳng cấu}.
\end{proof}

\subsection{Các bó xoắn và các bó không xoắn}
\label{subsection:I.7.4}
\label{I.7.4}

\begin{env}[7.4.1]
\label{I.7.4.1}
Cho $X$ là một tiền lược đồ \emph{nguyên}.
Với mọi $\sh{O}_X$-môđun $\sh{F}$, đồng cấu chính tắc $\sh{O}_X\to\sh{R}(X)$, bằng cách lấy tích tenxơ, xác định một đồng cấu (cũng gọi là \emph{chính tắc}) $\sh{F}\to\sh{F}\otimes_{\sh{O}_X}\sh{R}(X)$, mà trên mỗi thớ chính là đồng cấu $z\to z\otimes1$ từ $\sh{F}_x$ vào $\sh{F}_x\otimes_{\sh{O}_x}R(X)$.
\emph{Hạt nhân $\sh{T}$} của đồng cấu này là một $\sh{O}_X$-môđun con của $\sh{F}$, được gọi là \emph{bó xoắn} của $\sh{F}$; nó là quasi-coherent nếu $\sh{F}$ là quasi-coherent (\sref{I.4.1.1} và \sref{I.7.3.6}).
Ta nói $\sh{F}$ là \emph{không xoắn} nếu $\sh{T}=0$, và $\sh{F}$ là một \emph{bó xoắn} nếu $\sh{T}=\sh{F}$.
Với mọi $\sh{O}_X$-môđun $\sh{F}$, $\sh{F}/\sh{T}$ không xoắn.
Từ \sref{I.7.3.5}, ta suy ra rằng:
\end{env}

\begin{proposition}[7.4.2]
\label{I.7.4.2}
Nếu $X$ là một tiền lược đồ nguyên, thì mọi $\sh{O}_X$-môđun quasi-coherent không xoắn $\sh{F}$ đều đẳng cấu với một bó con $\sh{G}$ của một bó đơn có dạng $(\sh{R}(X))^{(I)}$, được sinh (với tư cách một $\sh{R}(X)$-môđun) bởi $\sh{G}$.
\end{proposition}

Lực lượng của $I$ được gọi là \emph{hạng} của $\sh{F}$; với mọi tập mở affine khác rỗng $U$ của $X$, hạng của $\sh{F}$ bằng hạng của $\Gamma(U,\sh{F})$ với tư cách một $\Gamma(U,\sh{O}_X)$-môđun, như thấy ngay khi xét điểm tổng quát của $X$, vốn nằm trong $U$.
Đặc biệt:

\begin{corollary}[7.4.3]
\label{I.7.4.3}
Trên một tiền lược đồ nguyên $X$, mọi $\sh{O}_X$-môđun quasi-coherent không xoắn và có hạng $1$ (đặc biệt, mọi $\sh{O}_X$-môđun khả nghịch) đều đẳng cấu với một $\sh{O}_X$-môđun con của $\sh{R}(X)$, và ngược lại.
\end{corollary}

\begin{corollary}[7.4.4]
\label{I.7.4.4}
Cho $X$ là một tiền lược đồ nguyên, $\sh{L}$ và $\sh{L}'$ là hai $\sh{O}_X$-môđun không xoắn, và $f$ (tương ứng $f'$) là một nhát cắt của $\sh{L}$ (tương ứng của $\sh{L}'$) trên $X$.
Để $f\otimes f'=0$, điều kiện cần và đủ là một trong hai nhát cắt $f$ và $f'$ bằng không.
\end{corollary}

\begin{proof}
Cho $x$ là điểm tổng quát của $X$; theo giả thiết, ta có $(f\otimes f')_x=f_x\otimes f'_x=0$.
Vì $\sh{L}_x$ và $\sh{L}'_x$ có thể được đồng nhất với các $\sh{O}_x$-môđun con của trường $\sh{O}_x$, hệ thức trên kéo theo $f_x=0$ hoặc $f'_x=0$, và do đó $f=0$ hoặc $f'=0$, vì $\sh{L}$ và $\sh{L}'$ không xoắn \sref{I.7.3.5}.
\end{proof}

\begin{proposition}[7.4.5]
\label{I.7.4.5}
Cho $X$ và $Y$ là hai tiền lược đồ nguyên, và $f:X\to Y$ là một cấu xạ trội.
Với mọi $\sh{O}_X$-môđun quasi-coherent không xoắn $\sh{F}$, $f_*(\sh{F})$ là một $\sh{O}_Y$-môđun không xoắn.
\end{proposition}

\oldpage[I]{164}
\begin{proof}
Vì $f_*$ là một hàm tử khớp trái \sref[0]{0.4.2.1}, theo \sref{I.7.4.2}, chỉ cần chứng minh mệnh đề trong trường hợp $\sh{F}=(\sh{R}(X))^{(I)}$.
Nhưng mọi tập mở khác rỗng $U$ của $Y$ đều chứa điểm tổng quát của $Y$, nên $f^{-1}(U)$ chứa điểm tổng quát của $X$ \sref[0]{0.2.1.5}; do đó ta có $\Gamma(U,f_*(\sh{F}))=\Gamma(f^{-1}(U),\sh{F})=(\sh{R}(X))^{(I)}$; nói cách khác, $f_*(\sh{F})$ là bó đơn có thớ $(\sh{R}(X))^{(I)}$, được xét với tư cách một $\sh{R}(Y)$-môđun, và rõ ràng là không xoắn.
\end{proof}

\begin{proposition}[7.4.6]
\label{I.7.4.6}
Cho $X$ là một tiền lược đồ nguyên và $x$ là điểm tổng quát của nó.
Với mọi $\sh{O}_X$-môđun quasi-coherent hữu hạn sinh $\sh{F}$, các điều kiện sau đây tương đương: \emph{(a)} $\sh{F}$ là một bó xoắn; \emph{(b)} $\sh{F}_x=0$; \emph{(c)} $\Supp(\sh{F})\neq X$.
\end{proposition}

\begin{proof}
Theo \sref{I.7.3.5} và \sref{I.7.4.1}, các hệ thức $\sh{F}_x=0$ và $\sh{F}\otimes_{\sh{O}_X}\sh{R}(X)=0$ là tương đương, do đó (a) và (b) tương đương; mặt khác, $\Supp(\sh{F})$ là một tập đóng trong $X$ \sref[0]{0.5.2.2}, và vì mọi tập mở khác rỗng của $X$ đều chứa $x$, (b) và (c) tương đương.
\end{proof}

\begin{env}[7.4.7]
\label{I.7.4.7}
Ta mở rộng (do lạm dụng ngôn ngữ) các định nghĩa của \sref{I.7.4.1} sang trường hợp $X$ là một tiền lược đồ \emph{rút gọn} chỉ có \emph{hữu hạn} thành phần bất khả quy; khi đó, từ \sref{I.7.3.4} suy ra rằng tính tương đương giữa a) và c) trong \sref{I.7.4.6} vẫn đúng đối với một tiền lược đồ như vậy, nếu thay điều kiện c) bằng điều kiện $\Supp(\sh{F})$ không chứa thành phần bất khả quy nào của $X$.
\end{env}

\section{Các lược đồ Chevalley}
\label{section:I.8}

\subsection{Các vành địa phương liên hệ với nhau}
\label{subsection:I.8.1}

Với mỗi vành địa phương $A$, ta ký hiệu $\mathfrak{m}(A)$ là iđêan cực đại của $A$.
\begin{lemma}[8.1.1]
\label{I.8.1.1}
Cho $A$ và $B$ là hai vành địa phương sao cho $A\subset B$.
Khi đó các điều kiện sau đây tương đương.
\begin{enumerate}
  \item[{\rm(i)}] $\mathfrak{m}(B)\cap A=\mathfrak{m}(A)$.
  \item[{\rm(ii)}] $\mathfrak{m}(A)\subset\mathfrak{m}(B)$.
  \item[{\rm(iii)}] $1$ không thuộc iđêan của $B$ sinh bởi $\mathfrak{m}(A)$.
\end{enumerate}
\end{lemma}

\begin{proof}
Hiển nhiên (i) suy ra (ii), và (ii) suy ra (iii); cuối cùng, nếu (iii) đúng, thì $\mathfrak{m}(B)\cap A$ chứa $\mathfrak{m}(A)$ và không chứa $1$, do đó bằng $\mathfrak{m}(A)$.

Khi các điều kiện tương đương của \sref{I.8.1.1} được thỏa mãn, ta nói $B$ \emph{chi phối} $A$; điều này tương đương với việc đồng cấu đơn ánh $A\to B$ là một đồng cấu \emph{địa phương}.
Hiển nhiên, trên tập các vành con địa phương của một vành $R$, quan hệ chi phối là một quan hệ thứ tự.
\end{proof}

\begin{env}[8.1.2]
\label{I.8.1.2}
Giờ đây xét một \emph{trường} $R$.
Với mọi vành con $A$ của $R$, ta ký hiệu $L(A)$ là tập hợp các vành địa phương $A_{\mathfrak{p}}$, trong đó $\mathfrak{p}$ chạy qua phổ nguyên tố của $A$;
các vành địa phương này được đồng nhất với các vành con của $R$ chứa $A$.
Vì $\mathfrak{p}=(\mathfrak{p}A_{\mathfrak{p}})\cap A$, ánh xạ $\mathfrak{p}\mapsto A_{\mathfrak{p}}$ từ $\Spec(A)$ vào $L(A)$ là song ánh.
\end{env}

\begin{lemma}[8.1.3]
\label{I.8.1.3}
Cho $R$ là một trường và $A$ là một vành con của $R$.
Điều kiện cần và đủ để một vành con địa phương $M$ của $R$ chi phối một vành $A_{\mathfrak{p}}\in L(A)$ là $A\subset M$;
khi đó vành địa phương $A_{\mathfrak{p}}$ do $M$ chi phối là duy nhất và tương ứng với $\mathfrak{p}=\mathfrak{m}(M)\cap A$.
\end{lemma}

\begin{proof}
Thật vậy, nếu $M$ chi phối $A_{\mathfrak{p}}$, thì theo \sref{I.8.1.1} ta có $\mathfrak{m}(M)\cap A_{\mathfrak{p}}=\mathfrak{p}A_{\mathfrak{p}}$, từ đó $\mathfrak{p}$ là duy nhất;
mặt khác, nếu $A\subset M$, thì $\mathfrak{m}(M)\cap A=\mathfrak{p}$ là một iđêan nguyên tố của $A$, và vì $A-\mathfrak{p}\subset M$, ta có $A_{\mathfrak{p}}\subset M$ và $\mathfrak{p}A_{\mathfrak{p}}\subset\mathfrak{m}(M)$, do đó $M$ chi phối $A_{\mathfrak{p}}$.
\end{proof}

\oldpage[I]{165}
\begin{lemma}[8.1.4]
\label{I.8.1.4}
Cho $R$ là một trường, $M$ và $N$ là hai vành con địa phương của $R$, và $P$ là vành con của $R$ sinh bởi $M\cup N$.
Khi đó các điều kiện sau đây tương đương.
\begin{enumerate}
  \item[{\rm(i)}] Tồn tại một iđêan nguyên tố $\mathfrak{p}$ của $P$ sao cho $\mathfrak{m}(M)=\mathfrak{p}\cap M$ và $\mathfrak{m}(N)=\mathfrak{p}\cap N$.
  \item[{\rm(ii)}] Iđêan $\mathfrak{a}$ sinh trong $P$ bởi $\mathfrak{m}(M)\cup\mathfrak{m}(N)$ khác $P$.
  \item[{\rm(iii)}] Tồn tại một vành con địa phương $Q$ của $R$ chi phối đồng thời $M$ và $N$.
\end{enumerate}
\end{lemma}

\begin{proof}
Hiển nhiên (i) suy ra (ii);
ngược lại, nếu $\mathfrak{a}\neq P$, thì $\mathfrak{a}$ nằm trong một iđêan cực đại $\mathfrak{n}$ của $P$, và vì $1\not\in\mathfrak{n}$, $\mathfrak{n}\cap M$ chứa $\mathfrak{m}(M)$ và khác $M$, do đó $\mathfrak{n}\cap M=\mathfrak{m}(M)$; tương tự, $\mathfrak{n}\cap N=\mathfrak{m}(N)$.
Hiển nhiên, nếu $Q$ chi phối đồng thời $M$ và $N$, thì $P\subset Q$ và $\mathfrak{m}(M)=\mathfrak{m}(Q)\cap M=(\mathfrak{m}(Q)\cap P)\cap M$, và $\mathfrak{m}(N)=(\mathfrak{m}(Q)\cap P)\cap N$, do đó (iii) suy ra (i);
chiều ngược lại là hiển nhiên khi lấy $Q=P_{\mathfrak{p}}$.
\end{proof}

Khi các điều kiện của \sref{I.8.1.4} được thỏa mãn, theo C.~Chevalley, ta nói hai vành địa phương $M$ và $N$ \emph{liên hệ với nhau}.

\begin{proposition}[8.1.5]
\label{I.8.1.5}
Cho $A$ và $B$ là hai vành con của một trường $R$, và $C$ là vành con của $R$ sinh bởi $A\cup B$.
Khi đó các điều kiện sau đây tương đương.
\begin{enumerate}
  \item[{\rm(i)}] Với mọi vành địa phương $Q$ chứa $A$ và $B$, ta có $A_{\mathfrak{p}}=B_{\mathfrak{q}}$, trong đó $\mathfrak{p}=\mathfrak{m}(Q)\cap A$ và $\mathfrak{q}=\mathfrak{m}(Q)\cap B$.
  \item[{\rm(ii)}] Với mọi iđêan nguyên tố $\mathfrak{r}$ của $C$, ta có $A_{\mathfrak{p}}=B_{\mathfrak{q}}$, trong đó $\mathfrak{p}=\mathfrak{r}\cap A$ và $\mathfrak{q}=\mathfrak{r}\cap B$.
  \item[{\rm(iii)}] Nếu $M\in L(A)$ và $N\in L(B)$ liên hệ với nhau, thì chúng đồng nhất.
  \item[{\rm(iv)}] $L(A)\cap L(B)=L(C)$.
\end{enumerate}
\end{proposition}

\begin{proof}
Các bổ đề~\sref{I.8.1.3} và \sref{I.8.1.4} chứng minh rằng (i) và (iii) tương đương;
hiển nhiên (i) suy ra (ii) khi áp dụng nó cho $Q=C_{\mathfrak{r}}$;
ngược lại, (ii) suy ra (i), vì nếu $Q$ chứa $A\cup B$ thì nó chứa $C$, và nếu $\mathfrak{r}=\mathfrak{m}(Q)\cap C$, thì $\mathfrak{p}=\mathfrak{r}\cap A$ và $\mathfrak{q}=\mathfrak{r}\cap B$, theo \sref{I.8.1.3}.
Hiển nhiên (iv) suy ra (i), vì nếu $Q$ chứa $A\cup B$ thì theo \sref{I.8.1.3}, nó chi phối một vành địa phương $C_{\mathfrak{r}}\in L(C)$;
theo giả thiết, ta có $C_{\mathfrak{r}}\in L(A)\cap L(B)$, và \sref{I.8.1.1} cùng \sref{I.8.1.3} chứng minh rằng $C_{\mathfrak{r}}=A_{\mathfrak{p}}=B_{\mathfrak{q}}$.
Cuối cùng, ta chứng minh rằng (iii) suy ra (iv).
Cho $Q\in L(C)$;
theo \sref{I.8.1.3}, $Q$ chi phối một $M\in L(A)$ và một $N\in L(B)$, do đó $M$ và $N$, vì liên hệ với nhau, nên đồng nhất theo giả thiết.
Khi đó $C\subset M$, và theo \sref{I.8.1.3}, $M$ chi phối một $Q'\in L(C)$; do đó $Q$ chi phối $Q'$, nên theo \sref{I.8.1.3}, nhất thiết $Q=Q'=M$, và vì vậy $Q\in L(A)\cap L(B)$.
Ngược lại, nếu $Q\in L(A)\cap L(B)$, thì $C\subset Q$, nên theo \sref{I.8.1.3}, $Q$ chi phối một $Q''\in L(C)\subset L(A)\cap L(B)$;
vì $Q$ và $Q''$ liên hệ với nhau nên chúng đồng nhất, do đó $Q''=Q\in L(C)$, và chứng minh hoàn tất.
\end{proof}

\subsection{Các vành địa phương của một tiền lược đồ nguyên}
\label{subsection:I.8.2}

\begin{env}[8.2.1]
\label{I.8.2.1}
Cho $X$ là một tiền lược đồ \emph{nguyên}, và cho $R$ là trường các hàm hữu tỉ của nó, đồng nhất với vành địa phương tại điểm tổng quát $a$ của $X$;
với mọi $x\in X$, ta biết rằng $\sh{O}_x$ được đồng nhất chính tắc với một vành con của $R$ \sref{I.7.1.5}, và, với mọi hàm hữu tỉ $f\in R$, miền xác định $\delta(f)$ của $f$ là tập mở gồm các $x\in X$ sao cho
$f\in\sh{O}_x$.
Do đó, theo \sref{I.7.2.6}, với mọi tập mở $U\subset X$, ta có
\[
\label{I.8.2.1.1}
  \Gamma(U,\sh{O}_X)=\bigcap_{x\in U}\sh{O}_x.
  \tag{8.2.1.1}
\]
\end{env}

\begin{proposition}[8.2.2]
\label{I.8.2.2}
Cho $X$ là một tiền lược đồ nguyên và $R$ là trường các hàm hữu tỉ của nó. Điều kiện cần và đủ để $X$ là một lược đồ là: quan hệ «$\sh{O}_x$ và $\sh{O}_y$ liên hệ với nhau» \sref{I.8.1.4}, với các điểm $x$ và $y$ của $X$, kéo theo $x=y$.
\end{proposition}

\begin{proof}
Giả sử điều kiện này được thỏa mãn; ta sẽ chứng minh rằng $X$ là một tiền lược đồ tách.
Cho $U$ và $V$ là hai tập mở affine phân biệt của $X$, có các vành tương ứng là $A$ và $B$, được đồng nhất với các vành con của $R$;
$U$ (tương ứng $V$) vì thế được đồng nhất \sref{I.8.1.2} với $L(A)$ (tương ứng $L(B)$), và giả thiết kéo theo \sref{I.8.1.5} rằng, nếu $C$ là vành con của $R$ sinh bởi $A\cup B$, thì $W=U\cap V$ được đồng nhất với $L(A)\cap L(B)=L(C)$.
Hơn nữa, ta biết (\cite{I-1}, p.~5-03, prop.~4~\emph{bis}) rằng mọi vành con $E$ của $R$ bằng giao của các vành địa phương thuộc $L(E)$;
do đó $C$ được đồng nhất với giao của các vành $\sh{O}_z$ với $z\in W$, nói cách khác \sref{I.8.2.1.1}, với $\Gamma(W,\sh{O}_X)$.
Bây giờ xét tiền lược đồ con của $X$ cảm sinh trên $W$;
đồng cấu đồng nhất $\vphi: C\to\Gamma(W,\sh{O}_X)$ tương ứng \sref{I.2.2.4} với một cấu xạ $\Phi=(\psi,\theta):W\to\Spec(C)$;
ta sẽ thấy rằng $\Phi$ là một \emph{đẳng cấu} giữa các tiền lược đồ, từ đó suy ra $W$ là một tập mở \emph{affine}.
Việc đồng nhất $W$ với $L(C)=\Spec(C)$ cho thấy $\psi$ là \emph{song ánh}.
Mặt khác, với mọi $x\in W$, $\theta_x^\sharp$ là đồng cấu đơn ánh $C_{\mathfrak{r}}\to\sh{O}_x$, trong đó $\mathfrak{r}=\mathfrak{m}_x\cap C$, và theo định nghĩa, $C_{\mathfrak{r}}$ được đồng nhất với $\sh{O}_x$, nên $\theta_x^\sharp$ là song ánh.
Do đó chỉ còn phải chứng minh rằng $\psi$ là một \emph{phép đồng phôi}, nói cách khác, với mọi tập con đóng $F\subset W$, $\psi(F)$ đóng trong
$\Spec(C)$.
Mà $F$ là giao của $W$ với một tập con đóng của $U$ có dạng $V(\mathfrak{a})$, trong đó $\mathfrak{a}$ là một iđêan của $A$;
ta sẽ chứng minh rằng $\psi(F)=V(\mathfrak{a}C)$, điều này sẽ chứng minh khẳng định.
Thật vậy, các iđêan nguyên tố của $C$ chứa $\mathfrak{a}C$ chính là các iđêan nguyên tố của $C$ chứa $\mathfrak{a}$, nên chúng có dạng $\psi(x)=\mathfrak{m}_x\cap C$, trong đó $\mathfrak{a}\subset\mathfrak{m}_x$ và $x\in W$;
vì, với $x\in W$, $\mathfrak{a}\subset\mathfrak{m}_x$ tương đương với $x\in V(\mathfrak{a})$, và do đó với $x\in F=W\cap V(\mathfrak{a})$, ta quả thực có $\psi(F)=V(\mathfrak{a}C)$\footnote{\emph{[Dịch giả] Văn bản tiếng Pháp in $x\in V(\mathfrak a)=W\cap F$ đối với $x\in U$. Vì $F$ vừa được định nghĩa là giao của $V(\mathfrak a)$ với $W$, đẳng thức cần dùng là $F=W\cap V(\mathfrak a)$, và bước này xét $x\in W$.}}.

Suy ra $X$ là một tiền lược đồ tách, vì $U\cap V$ là affine và vành $C$ của nó được sinh bởi hợp $A\cup B$ của các vành của $U$ và $V$ \sref{I.5.5.6}.

Ngược lại, giả sử $X$ là một tiền lược đồ tách, và cho $x$ và $y$ là hai điểm của $X$ sao cho $\sh{O}_x$ và $\sh{O}_y$ liên hệ với nhau.
Cho $U$ (tương ứng $V$) là một tập mở affine chứa $x$ (tương ứng $y$), có vành $A$ (tương ứng $B$);
khi đó ta biết rằng $U\cap V$ là affine và vành $C$ của nó được sinh bởi $A\cup B$ \sref{I.5.5.6}.
Nếu $\mathfrak{p}=\mathfrak{m}_x\cap A$ và $\mathfrak{q}=\mathfrak{m}_y\cap B$, thì $A_{\mathfrak{p}}=\sh{O}_x$ và $B_{\mathfrak{q}}=\sh{O}_y$; vì $A_{\mathfrak{p}}$ và $B_{\mathfrak{q}}$ liên hệ với nhau, tồn tại một iđêan nguyên tố $\mathfrak{r}$ của $C$ sao cho $\mathfrak{p}=\mathfrak{r}\cap A$ và $\mathfrak{q}=\mathfrak{r}\cap B$ \sref{I.8.1.4}.
Nhưng khi đó tồn tại một điểm $z\in U\cap V$ sao cho $\mathfrak{r}=\mathfrak{m}_z\cap C$, vì $U\cap V$ là affine; do đó hiển nhiên $x=z$ và $y=z$, suy ra $x=y$.
\end{proof}

\begin{corollary}[8.2.3]
\label{I.8.2.3}
Cho $X$ là một lược đồ \emph{nguyên}, và cho $x$ và $y$ là hai điểm của $X$.
Điều kiện cần và đủ để $x\in\overline{\{y\}}$ là
$\sh{O}_x\subset\sh{O}_y$, hay tương đương, mọi hàm hữu tỉ được xác định
tại $x$ cũng được xác định tại $y$.
\end{corollary}

\begin{proof}
Điều kiện này hiển nhiên là cần thiết, vì miền xác định
$\delta(f)$ của một hàm hữu tỉ $f\in R$ là tập mở; bây giờ ta chứng minh
điều kiện đó là đủ.
Nếu $\sh{O}_x\subset\sh{O}_y$, tồn tại một iđêan nguyên tố $\mathfrak{p}$
của $\sh{O}_x$ sao cho $\sh{O}_y$ chi phối
$(\sh{O}_x)_{\mathfrak{p}}$ \sref{I.8.1.3}; nhưng theo \sref{I.2.4.2}, tồn tại
một điểm $z\in X$ sao cho $x\in\overline{\{z\}}$ và
$\sh{O}_z=(\sh{O}_x)_{\mathfrak{p}}$; vì $\sh{O}_z$ và $\sh{O}_y$ liên hệ
với nhau, theo \sref{I.8.2.2} ta có $z=y$, từ đó suy ra hệ quả.
\end{proof}

\begin{corollary}[8.2.4]
\label{I.8.2.4}
Cho $X$ là một lược đồ \emph{nguyên}; khi đó ánh xạ $x\mapsto\sh{O}_x$ là đơn ánh; nói cách khác, nếu $x$ và $y$ là hai điểm phân biệt của $X$, tồn tại một hàm hữu tỉ được xác định tại một trong hai điểm ấy mà không được xác định tại điểm kia.
\end{corollary}

\begin{proof}
Điều này suy ra từ \sref{I.8.2.3} và tiên đề $(T_0)$ \sref{I.2.1.4}.
\end{proof}

\begin{corollary}[8.2.5]
\label{I.8.2.5}
Cho $X$ là một lược đồ \emph{nguyên} có không gian nền là noetherian;
khi $f$ chạy qua trường $R$ các hàm hữu tỉ trên $X$, các tập $\delta(f)$ sinh ra tôpô của $X$.
\end{corollary}

Thật vậy, mọi tập con đóng của $X$ khi đó là hợp hữu hạn của các tập con đóng bất khả quy, nghĩa là các tập con có dạng $\overline{\{y\}}$ \sref{I.2.1.5}.
Mà nếu $x\not\in\overline{\{y\}}$, tồn tại một hàm hữu tỉ $f$ được xác định tại $x$ nhưng không được xác định tại $y$ \sref{I.8.2.3}; nói cách khác, ta có $x\in\delta(f)$ và $\delta(f)$ không giao với $\overline{\{y\}}$.
Do đó, phần bù của $\overline{\{y\}}$ là hợp của các tập có dạng $\delta(f)$, và theo nhận xét đầu tiên, mọi tập mở của $X$ là hợp của các giao hữu hạn của những tập mở có dạng $\delta(f)$.

\begin{env}[8.2.6]
\label{I.8.2.6}
Hệ quả~\sref{I.8.2.5} cho thấy tôpô của $X$ được đặc trưng hoàn toàn bởi họ các vành địa phương $(\sh{O}_x)_{x\in X}$ có $R$ làm trường phân thức.
Điều này cũng tương đương với việc nói rằng các tập con đóng của $X$ được xác định như sau: cho một tập con hữu hạn $\{x_1,\ldots,x_n\}$ của $X$, xét tập hợp các $y\in X$ sao cho $\sh{O}_y\subset\sh{O}_{x_i}$ với ít nhất một chỉ số $i$; các tập này (ứng với mọi cách chọn $\{x_1,\ldots,x_n\}$) chính là các tập con đóng của $X$.
Hơn nữa, một khi tôpô của $X$ đã biết, bó cấu trúc $\sh{O}_X$ cũng được xác định bởi họ các $\sh{O}_x$, vì $\Gamma(U,\sh{O}_X)=\bigcap_{x\in U}\sh{O}_x$, theo \sref{I.8.2.1.1}.
Do đó, họ $(\sh{O}_x)_{x\in X}$ xác định hoàn toàn tiền lược đồ $X$ khi $X$ là một lược đồ nguyên có không gian nền là noetherian.
\end{env}

\begin{proposition}[8.2.7]
\label{I.8.2.7}
Cho $X$ và $Y$ là hai lược đồ nguyên, $f:X\to Y$ là một cấu xạ trội \sref{I.2.2.6}, và $K$ (tương ứng $L$) là trường các hàm hữu tỉ trên $X$ (tương ứng trên $Y$).
Khi đó $L$ có thể được đồng nhất với một trường con của $K$, và với mọi $x\in X$, $\sh{O}_{f(x)}$ là vành địa phương duy nhất của $Y$ do $\sh{O}_x$ chi phối.
\end{proposition}

\begin{proof}
Thật vậy, nếu $f=(\psi,\theta)$ và $a$ là điểm tổng quát của $X$, thì $\psi(a)$ là điểm tổng quát của $Y$ \sref[0]{0.2.1.5};
khi đó $\theta_a^\sharp$ là một đơn cấu từ trường $L=\sh{O}_{\psi(a)}$ vào trường $K=\sh{O}_a$.
Vì mọi tập mở affine khác rỗng $U$ của $Y$ đều chứa $\psi(a)$, theo \sref{I.2.2.4}, đồng cấu $\Gamma(U,\sh{O}_Y)\to\Gamma(\psi^{-1}(U),\sh{O}_X)$ tương ứng với $f$ là hạn chế của $\theta_a^\sharp$ trên $\Gamma(U,\sh{O}_Y)$.
Do đó, với mọi $x\in X$, $\theta_x^\sharp$ là hạn chế trên $\sh{O}_{\psi(x)}$ của $\theta_a^\sharp$, nên là một đơn cấu.
Ta còn biết rằng $\theta_x^\sharp$ là một đồng cấu địa phương, nên, nếu đồng nhất $L$ với một trường con của $K$ qua $\theta_a^\sharp$, thì $\sh{O}_{\psi(x)}$ do $\sh{O}_x$ chi phối \sref{I.8.1.1};
hơn nữa, đó là vành địa phương duy nhất của $Y$ do $\sh{O}_x$ chi phối, vì hai vành địa phương của $Y$ liên hệ với nhau thì đồng nhất \sref{I.8.2.2}.
\end{proof}

\begin{proposition}[8.2.8]
\label{I.8.2.8}
Cho $X$ là một tiền lược đồ \emph{bất khả quy}, $f:X\to Y$ là một phép nhúng địa phương (tương ứng, một đẳng cấu địa phương), và giả sử thêm rằng $f$ là tách. Khi đó $f$ là một phép nhúng (tương ứng, một phép nhúng mở).
\end{proposition}

\begin{proof}
Đặt $f=(\psi,\theta)$;
trong cả hai trường hợp, chỉ cần chứng minh rằng $\psi$ là một \emph{phép đồng phôi} từ $X$ lên $\psi(X)$ \sref{I.4.5.3}.
Thay $f$ bằng $f_\text{red}$ (\sref{I.5.1.6} và \sref{I.5.5.1}[vi]), ta có thể giả sử rằng $X$ và $Y$ đều \emph{rút gọn}.
Nếu $Y'$ là tiền lược đồ con đóng rút gọn của $Y$ có không gian nền là $\overline{\psi(X)}$, thì $f$ phân tích thành $X\xrightarrow{f'}Y'\xrightarrow{j}Y$, trong đó $j$ là phép nhúng chính tắc \sref{I.5.2.2}.
Từ \sref{I.5.5.1}[v] suy ra rằng $f'$ vẫn là một cấu xạ tách; hơn nữa, $f'$ vẫn là
\oldpage[I]{168}
một phép nhúng địa phương (tương ứng, một đẳng cấu địa phương), vì vấn đề có tính địa phương trên $X$ và $Y$, nên để chứng minh điều này, ta có thể giới hạn ở trường hợp $f$ là một phép nhúng đóng (tương ứng, một phép nhúng mở), và khẳng định của ta khi đó suy ra ngay từ \sref{I.4.2.2}.

Do đó, ta có thể giả sử rằng $f$ là một cấu xạ \emph{trội}, kéo theo $Y$ cũng bất khả quy \sref[0]{0.2.1.5}, và vì vậy cả $X$ lẫn $Y$ đều \emph{nguyên}.
Hơn nữa, vì vấn đề có tính địa phương trên $Y$, ta có thể giả sử rằng $Y$ là một lược đồ affine;
do $f$ là tách, $X$ là một lược đồ \sref{I.5.5.1}[ii], và cuối cùng ta ở trong các giả thiết của Mệnh đề~\sref{I.8.2.7}.
Khi đó, với mọi $x\in X$, $\theta_x^\sharp$ là đơn ánh;
nhưng giả thiết rằng $f$ là một phép nhúng địa phương kéo theo $\theta_x^\sharp$ là toàn ánh \sref{I.4.2.2}, nên $\theta_x^\sharp$ là song ánh, tức là, với phép đồng nhất trong Mệnh đề~\sref{I.8.2.7}, ta có $\sh{O}_{\psi(x)}=\sh{O}_x$.
Theo Hệ quả~\sref{I.8.2.4}, điều này kéo theo $\psi$ là một ánh xạ \emph{đơn ánh}, và như vậy đã chứng minh mệnh đề khi $f$ là một đẳng cấu địa phương \sref{I.4.5.3}.
Nếu chỉ giả sử rằng $f$ là một phép nhúng địa phương, thì với mọi $x\in X$, tồn tại một lân cận mở $U$ của $x$ trong $X$ và một lân cận mở $V$ của $\psi(x)$ trong $Y$ sao cho hạn chế của $\psi$ trên $U$ là một phép đồng phôi từ $U$ lên một tập con \emph{đóng} của $V$.
Nhưng $U$ trù mật trong $X$, nên $\psi(U)$ trù mật trong $Y$ và \emph{do đó cũng trù mật} trong $V$, chứng tỏ rằng $\psi(U)=V$;
vì $\psi$ là đơn ánh, $\psi^{-1}(V)=U$, và điều này hoàn tất việc chứng minh rằng $\psi$ là một phép đồng phôi từ $X$ lên $\psi(X)$.
\end{proof}

\subsection{Các lược đồ Chevalley}
\label{subsection:I.8.3}

\begin{env}[8.3.1]
\label{I.8.3.1}
Cho $X$ là một lược đồ nguyên \emph{noetherian}, và $R$ là trường các hàm hữu tỉ của nó;
ký hiệu $X'$ là tập các vành con địa phương $\sh{O}_x\subset R$, trong đó $x$ chạy qua mọi điểm của $X$.
Tập $X'$ thỏa mãn ba điều kiện sau.
\begin{enumerate}
  \item[(Sch.~1)] \emph{Với mọi $M\in X'$, $R$ là trường phân thức của $M$.}
  \item[(Sch.~2)] \emph{Tồn tại một tập hữu hạn các vành con noetherian $A_i$ của $R$ sao cho $X'=\bigcup_i L(A_i)$ và, với mọi cặp chỉ số $i$, $j$, vành con $A_{ij}$ của $R$ sinh bởi $A_i\cup A_j$ là một đại số kiểu hữu hạn trên $A_i$.}
  \item[(Sch.~3)] \emph{Hai phần tử bất kỳ $M$ và $N$ của $X'$ liên hệ với nhau thì trùng nhau.}
\end{enumerate}
\end{env}

Ta đã thấy trong \sref{I.8.2.1} rằng (Sch.~1) được thỏa mãn, còn (Sch.~3) suy ra từ \sref{I.8.2.2}.
Để chứng minh (Sch.~2), chỉ cần phủ $X$ bởi hữu hạn tập mở affine $U_i$ có các vành noetherian, và đặt $A_i=\Gamma(U_i,\sh{O}_X)$;
giả thiết rằng $X$ là một lược đồ kéo theo $U_i\cap U_j$ là affine, đồng thời $\Gamma(U_i\cap U_j,\sh{O}_X)=A_{ij}$ \sref{I.5.5.6};
hơn nữa, vì không gian $U_i$ là noetherian, phép nhúng $U_i\cap U_j\to U_i$ là kiểu hữu hạn \sref{I.6.3.5}, nên $A_{ij}$ là một đại số kiểu hữu hạn trên $A_i$ \sref{I.6.3.3}.

\begin{env}[8.3.2]
\label{I.8.3.2}
Các cấu trúc có các tiên đề (Sch.~1), (Sch.~2) và (Sch.~3) tổng quát hóa các ``lược đồ'' theo nghĩa của C.~Chevalley, người còn giả sử rằng $R$ là một mở rộng kiểu hữu hạn của một trường $K$, và các $A_i$ là những đại số kiểu hữu hạn trên $K$ (điều này làm cho một phần của (Sch.~2) trở nên thừa) \cite{I-1}.
Ngược lại, nếu có một cấu trúc như vậy trên một tập $X'$, thì ta có thể gắn với nó một lược đồ nguyên $X$ bằng cách sử dụng các nhận xét trong \sref{I.8.2.6}: không gian nền của $X$ bằng $X'$ được trang bị tôpô xác định trong \sref{I.8.2.6}, cùng với bó $\sh{O}_X$
\oldpage[I]{169}
sao cho $\Gamma(U,\sh{O}_X)=\bigcap_{x\in U}\sh{O}_x$ với mọi tập mở $U\subset X$, với định nghĩa hiển nhiên của các đồng cấu hạn chế.
Ta dành cho người đọc việc kiểm tra rằng bằng cách này ta thu được một lược đồ nguyên có các vành địa phương là những phần tử của $X'$;
ta sẽ không sử dụng kết quả này trong phần tiếp theo.
\end{env}

\section{Bổ sung về các bó quasi-coherent}
\label{section:I.9}

\subsection{Tích tenxơ của các bó quasi-coherent}
\label{subsection:I.9.1}

\begin{proposition}[9.1.1]
\label{I.9.1.1}
Cho $X$ là một tiền lược đồ (tương ứng, một tiền lược đồ noetherian địa phương). Cho $\sh{F}$ và
$\sh{G}$ là hai $\sh{O}_X$-môđun quasi-coherent (tương ứng, coherent); khi đó
$\sh{F}\otimes_{\sh{O}_X}\sh{G}$ là quasi-coherent (tương ứng, coherent); hơn nữa, nó là
kiểu hữu hạn nếu $\sh{F}$ và $\sh{G}$ đều là kiểu hữu hạn. Nếu
$\sh{F}$ có một trình bày hữu hạn và $\sh{G}$ là quasi-coherent
(tương ứng, coherent), thì $\shHom_{\sh{O}_X}(\sh{F},\sh{G})$ là quasi-coherent
(tương ứng, coherent).
\end{proposition}

\begin{proof}
Vì vấn đề có tính địa phương, ta có thể giả sử rằng $X$ là affine (tương ứng, affine
noetherian); hơn nữa, nếu $\sh{F}$ là coherent, ta có thể giả sử rằng nó là
đối hạt nhân của một đồng cấu $\sh{O}_X^m\to\sh{O}_X^n$. Các khẳng định về
bó quasi-coherent suy ra từ \sref{I.1.3.12} và \sref{I.1.3.9}; các khẳng định
về bó coherent suy ra từ \sref{I.1.5.1} và từ việc
nếu $M$ và $N$ là các môđun kiểu hữu hạn trên một vành noetherian $A$,
thì $M\otimes_A N$ và $\Hom_A(M,N)$ đều là các $A$-môđun kiểu hữu hạn.
\end{proof}

\begin{definition}[9.1.2]
\label{I.9.1.2}
Cho $X$ và $Y$ là hai $S$-tiền lược đồ, $p$ và
$q$ là các phép chiếu của $X\times_S Y$, và $\sh{F}$ (tương ứng, $\sh{G}$) là một
$\sh{O}_X$-môđun quasi-coherent (tương ứng, một $\sh{O}_Y$-môđun quasi-coherent). Ta gọi
tích tenxơ của $\sh{F}$ và $\sh{G}$ trên $\sh{O}_S$ (\emph{hoặc} trên $S$),
ký hiệu $\sh{F}\otimes_{\sh{O}_S}\sh{G}$
(\emph{hoặc} $\sh{F}\otimes_S\sh{G}$), là tích tenxơ
$p^*(\sh{F})\otimes_{\sh{O}_{X\times_S Y}}q^*(\sh{G})$ trên tiền lược đồ $X\times_S Y$.
\end{definition}

Nếu các $X_i$ ($1\leq i\leq n$) là các $S$-tiền lược đồ, và các $\sh{F}_i$ ($1\leq i\leq n$) là
các $\sh{O}_{X_i}$-môđun quasi-coherent, thì ta định nghĩa tương tự tích tenxơ
$\sh{F}_1\otimes_S\sh{F}_2\otimes_S\cdots\otimes_S\sh{F}_n$ trên
tiền lược đồ $Z=X_1\times_S X_2\times_S\cdots\times_S X_n$; đây là một
$\sh{O}_Z$-môđun \emph{quasi-coherent} theo \sref{I.9.1.1} và
\sref[0]{0.5.1.4}; nó là \emph{coherent} nếu các $\sh{F}_i$ đều coherent \emph{và}
$Z$ là \emph{noetherian địa phương}, theo \sref{I.9.1.1},
\sref[0]{0.5.3.11} và \sref{I.6.1.1}.

Chú ý rằng, nếu lấy $X=Y=S$, thì Định nghĩa~\sref{I.9.1.2} cho lại tích
tenxơ của các $\sh{O}_S$-môđun. Hơn nữa, vì $q^*(\sh{O}_Y)=\sh{O}_{X\times_S Y}$
\sref[0]{0.4.3.4}, tích $\sh{F}\otimes_S\sh{O}_Y$ được đồng nhất
chính tắc với $p^*(\sh{F})$, và tương tự,
$\sh{O}_X\otimes_S\sh{G}$ được đồng nhất chính tắc với $q^*(\sh{G})$. Đặc
biệt, nếu lấy $Y=S$ và ký hiệu $f$ là cấu xạ cấu trúc $X\to Y$,
thì ta có $\sh{O}_X\otimes_Y\sh{G}=f^*(\sh{G})$: tích tenxơ thông thường
và ảnh ngược do đó xuất hiện như những trường hợp riêng của tích
tenxơ tổng quát.

Định nghĩa~\sref{I.9.1.2} suy ra ngay rằng, với $X$ và
$Y$ cố định, $\sh{F}\otimes_S\sh{G}$ là một \emph{song hàm tử hiệp biến, cộng tính và khớp phải} theo $\sh{F}$ và $\sh{G}$.

\begin{proposition}[9.1.3]
\label{I.9.1.3}
Cho $S$, $X$ và $Y$ là ba lược đồ affine có các vành
tương ứng $A$, $B$ và $C$, trong đó $B$ và $C$ là các đại số trên $A$. Cho $M$
(tương ứng, $N$) là một $B$-môđun (tương ứng, một $C$-môđun), và
$\sh{F}=\widetilde{M}$ (tương ứng, $\sh{G}=\widetilde{N}$) là
bó quasi-coherent liên kết; khi đó $\sh{F}\otimes_S\sh{G}$
đẳng cấu chính tắc với bó liên kết với $(B\otimes_A C)$-môđun
$M\otimes_A N$.
\end{proposition}

\begin{proof}
\oldpage[I]{170}
Thật vậy, theo \sref{I.1.6.5}, $\sh{F}\otimes_S\sh{G}$
đẳng cấu chính tắc với bó liên kết với $(B\otimes_A C)$-môđun
\[
  \big(M\otimes_B(B\otimes_A C)\big)\otimes_{B\otimes_A C}\big((B\otimes_A C)\otimes_C N\big)
\]
và, nhờ các đẳng cấu chính tắc giữa các tích tenxơ, môđun sau cùng này đẳng cấu với
\[
  M\otimes_B(B\otimes_A C)\otimes_C N=(M\otimes_B B)\otimes_A(C\otimes_C N)=M\otimes_A N.
\]
\end{proof}

\begin{proposition}[9.1.4]
\label{I.9.1.4}
Cho $f:T\to X$ và $g:T\to Y$ là hai $S$-cấu xạ, và $\sh{F}$ (tương ứng, $\sh{G}$) là một $\sh{O}_X$-môđun quasi-coherent (tương ứng, một $\sh{O}_Y$-môđun quasi-coherent).
Khi đó
\[
  (f,g)^*_S(\sh{F}\otimes_S\sh{G})=f^*(\sh{F})\otimes_{\sh{O}_T}g^*(\sh{G}).
\]
\end{proposition}

\begin{proof}
Nếu $p$, $q$ là các phép chiếu của $X\times_S Y$, thì công thức suy ra
từ các đẳng thức $(f,g)^*_S\circ p^*=f^*$ và
$(f,g)^*_S\circ q^*=g^*$ \sref[0]{0.3.5.5}, và từ việc ảnh
ngược của một tích tenxơ các bó đại số là tích tenxơ các ảnh
ngược của chúng \sref[0]{0.4.3.3}.
\end{proof}

\begin{corollary}[9.1.5]
\label{I.9.1.5}
Cho $f:X\to X'$ và $g:Y\to Y'$
là hai $S$-cấu xạ, và $\sh{F}'$ (tương ứng, $\sh{G}'$) là một
$\sh{O}_{X'}$-môđun quasi-coherent (tương ứng, một $\sh{O}_{Y'}$-môđun quasi-coherent). Khi đó
\[
  (f\times_S g)^*(\sh{F}'\otimes_S\sh{G}')=f^*(\sh{F}')\otimes_S g^*(\sh{G}').
\]
\end{corollary}

\begin{proof}
Điều này suy ra từ \sref{I.9.1.4} và từ việc $f\times_S g=(f\circ p, g\circ q)_S$,
trong đó $p$ và $q$ là các phép chiếu của $X\times_S Y$.
\end{proof}

\begin{corollary}[9.1.6]
\label{I.9.1.6}
Cho $X$, $Y$ và $Z$ là ba $S$-tiền lược đồ, và $\sh{F}$ (tương ứng, $\sh{G}$, $\sh{H}$) là một
$\sh{O}_X$-môđun quasi-coherent (tương ứng, một $\sh{O}_Y$-môđun quasi-coherent, một
$\sh{O}_Z$-môđun quasi-coherent); khi đó bó $\sh{F}\otimes_S\sh{G}\otimes_S\sh{H}$ là ảnh ngược
của $(\sh{F}\otimes_S\sh{G})\otimes_S\sh{H}$ qua đẳng cấu chính tắc từ
$X\times_S Y\times_S Z$ đến $(X\times_S Y)\times_S Z$.
\end{corollary}

\begin{proof}
Đẳng cấu này được viết là $(p_1,p_2)_S\times_S p_3$, trong đó $p_1$, $p_2$ và $p_3$
là các phép chiếu của $X\times_S Y\times_S Z$.

Tương tự, ảnh ngược của $\sh{G}\otimes_S\sh{F}$ qua đẳng cấu chính tắc từ
$X\times_S Y$ đến $Y\times_S X$ là $\sh{F}\otimes_S\sh{G}$.
\end{proof}

\begin{corollary}[9.1.7]
\label{I.9.1.7}
Nếu $X$ là một $S$-tiền lược đồ, thì mọi $\sh{O}_X$-môđun quasi-coherent $\sh{F}$ đều là ảnh
ngược của $\sh{F}\otimes_S\sh{O}_S$ qua đẳng cấu chính tắc từ $X$ đến $X\times_S S$
\sref{I.3.3.3}.
\end{corollary}

\begin{proof}
Đẳng cấu này là $(1_X,\vphi)_S$, trong đó $\vphi$ là cấu xạ cấu trúc
$X\to S$, và hệ quả suy ra từ \sref{I.9.1.4} và từ việc
$\vphi^*(\sh{O}_S)=\sh{O}_X$.
\end{proof}

\begin{env}[9.1.8]
\label{I.9.1.8}
Cho $X$ là một $S$-tiền lược đồ, $\sh{F}$ là một $\sh{O}_X$-môđun
quasi-coherent và $\vphi:S'\to S$ là một cấu xạ; ta ký hiệu
$\sh{F}_{(\vphi)}$ hoặc $\sh{F}_{(S')}$ là bó quasi-coherent
$\sh{F}\otimes_S\sh{O}_{S'}$ trên $X\times_S S'=X_{(\vphi)}=X_{(S')}$; do đó
$\sh{F}_{(S')}=p^*(\sh{F})$, trong đó $p$ là phép chiếu $X_{(S')}\to X$.
\end{env}

\begin{proposition}[9.1.9]
\label{I.9.1.9}
Cho $\vphi':S''\to S'$ là một cấu xạ.
Với mọi $\sh{O}_X$-môđun quasi-coherent $\sh{F}$ trên $S$-tiền lược đồ $X$,
$(\sh{F}_{(\vphi)})_{(\vphi')}$ là ảnh ngược của
$\sh{F}_{(\vphi\circ\vphi')}$ qua đẳng cấu chính tắc
$(X_{(\vphi)})_{(\vphi')}\isoto X_{(\vphi\circ\vphi')}$
\sref{I.3.3.9}.
\end{proposition}

\begin{proof}
Điều này suy ra ngay từ các định nghĩa và từ \sref{I.3.3.9}, và còn được
viết dưới dạng
\[
\label{I.9.1.9.1}
  (\sh{F}\otimes_S\sh{O}_{S'})\otimes_{S'}\sh{O}_{S''}=\sh{F}\otimes_S\sh{O}_{S''}.
  \tag{9.1.9.1}
\]
\end{proof}

\begin{proposition}[9.1.10]
\label{I.9.1.10}
Cho $Y$ là một $S$-tiền lược đồ, và cho $f:X\to Y$ là một $S$-cấu xạ.
Với mọi $\sh{O}_Y$-môđun quasi-coherent $\sh{G}$ và mọi cấu xạ
$S'\to S$, ta có
$(f_{(S')})^*(\sh{G}_{(S')})=(f^*(\sh{G}))_{(S')}$.
\end{proposition}

\begin{proof}
Điều này suy ra ngay từ tính giao hoán của biểu đồ
\oldpage[I]{171}
\[
  \begin{tikzcd}
    X_{(S')}\arrow[r,"f_{(S')}"]\arrow[d] &
    Y_{(S')}\arrow[d]\\
    X\arrow[r,"f"'] &
    Y.
  \end{tikzcd}
\]
\end{proof}

\begin{corollary}[9.1.11]
\label{I.9.1.11}
Cho $X$ và $Y$ là hai $S$-tiền lược đồ, và cho
$\sh{F}$ (tương ứng, $\sh{G}$) là một $\sh{O}_X$-môđun quasi-coherent
(tương ứng, một $\sh{O}_Y$-môđun quasi-coherent). Ảnh ngược của bó
$(\sh{F}_{(S')})\otimes_{S'}(\sh{G}_{(S')})$ qua đẳng cấu chính tắc
$(X\times_S Y)_{(S')}\isoto(X_{(S')})\times_{S'}(Y_{(S')})$
\sref{I.3.3.10} bằng $(\sh{F}\otimes_S\sh{G})_{(S')}$.
\end{corollary}

\begin{proof}
Nếu $p$ và $q$ là các phép chiếu của $X\times_S Y$, thì đẳng cấu đang xét
chính là $(p_{(S')}, q_{(S')})_{S'}$; khi đó hệ quả suy ra từ các
Mệnh đề~\sref{I.9.1.4} và \sref{I.9.1.10}.
\end{proof}

\begin{proposition}[9.1.12]
\label{I.9.1.12}
Với các ký hiệu trong Định nghĩa~\sref{I.9.1.2}, cho $z$
là một điểm của $X\times_S Y$, và đặt $x=p(z)$, $y=q(z)$; thớ
$(\sh{F}\otimes_S\sh{G})_z$ đẳng cấu với
$(\sh{F}_x\otimes_{\sh{O}_x}\sh{O}_z)\otimes_{\sh{O}_z}(\sh{G}_y\otimes_{\sh{O}_y}\sh{O}_z)
  =\sh{F}_x\otimes_{\sh{O}_x}\sh{O}_z\otimes_{\sh{O}_y}\sh{G}_y$.
\end{proposition}

\begin{proof}
Vì ta có thể quy về trường hợp affine, mệnh đề suy ra từ
công thức~\sref{I.1.6.5.1}.
\end{proof}

\begin{corollary}[9.1.13]
\label{I.9.1.13}
Nếu $\sh{F}$ và $\sh{G}$ có kiểu hữu hạn, thì ta có
\[
  \Supp(\sh{F}\otimes_S\sh{G})=p^{-1}(\Supp(\sh{F}))\cap q^{-1}(\Supp(\sh{G})).
\]
\end{corollary}

\begin{proof}
Vì $p^*(\sh{F})$ và $q^*(\sh{G})$ đều có kiểu hữu hạn trên
$\sh{O}_{X\times_S Y}$, theo Mệnh đề~\sref{I.9.1.12} và
\sref[0]{0.1.7.5}, ta quy về trường hợp riêng $\sh{G}=\sh{O}_Y$, nghĩa
là cần chứng minh công thức sau:\footnote{\emph{[ND] Bản
gốc tiếng Pháp in $\Supp(p^{-1}(\sh{F}))$ trong (9.1.13.1). Phép quy về
vừa thực hiện đồng nhất $\sh{F}\otimes_S\sh{O}_Y$ với $p^*(\sh{F})$, và phép
kiểm tra tiếp theo $\sh{O}_z/\mathfrak{m}_x\sh{O}_z\neq0$ là phép tính
trên các thớ của ảnh ngược này; vì vậy ta viết $p^*(\sh{F})$.
Công thức in trong bản gốc là đúng đối với ảnh ngược của các bó nhóm abel; ở đây ta phát biểu khẳng định về ảnh ngược của môđun.}}
\[
\label{I.9.1.13.1}
  \Supp(p^*(\sh{F}))=p^{-1}(\Supp(\sh{F})).
  \tag{9.1.13.1}
\]

Lập luận tương tự như trong \sref[0]{0.1.7.5} quy việc chứng minh về kiểm tra rằng, với mọi
$z\in X\times_S Y$, ta có $\sh{O}_z/\mathfrak{m}_x\sh{O}_z\neq0$ (với $x=p(z)$),
điều này suy ra từ việc đồng cấu $\sh{O}_x\to\sh{O}_z$ là \emph{địa phương},
theo giả thiết.
\end{proof}

Ta dành cho người đọc việc mở rộng sang một tích với số thừa số tùy ý
các kết quả đã chứng minh trong mục này cho hai thừa số.

\subsection{Ảnh trực tiếp của một bó quasi-coherent}
\label{subsection:I.9.2}

\begin{proposition}[9.2.1]
\label{I.9.2.1}
Cho $f:X\to Y$ là một cấu xạ các tiền lược đồ. Giả sử tồn tại một phủ
$(Y_\alpha)$ của $Y$ bởi các tập mở affine có tính chất sau: mỗi
$f^{-1}(Y_\alpha)$ thừa nhận một phủ \emph{hữu hạn} $(X_{\alpha i})$
bởi các tập mở affine nằm trong $f^{-1}(Y_\alpha)$, sao cho mỗi giao
$X_{\alpha i}\cap X_{\alpha j}$ lại là một hợp \emph{hữu hạn} của các
tập mở affine. Khi đó, với mọi $\sh{O}_X$-môđun quasi-coherent
$\sh{F}$, $f_*(\sh{F})$ là một $\sh{O}_Y$-môđun quasi-coherent.
\end{proposition}

\begin{proof}
Vì vấn đề có tính địa phương trên $Y$, ta có thể giả sử $Y$ bằng một
trong các $Y_\alpha$ và do đó bỏ các chỉ số $\alpha$.
\begin{enumerate}
  \item[(a)] Trước hết, giả sử các $X_i\cap X_j$ tự chúng là những tập
    mở \emph{affine}. Đặt $\sh{F}_i=\sh{F}|X_i$ và
    $\sh{F}_{ij}=\sh{F}|(X_i\cap X_j)$, và gọi $\sh{F}'_i$ và
    $\sh{F}'_{ij}$ là các ảnh trực tiếp của $\sh{F}_i$ và
    $\sh{F}_{ij}$ qua các hạn chế của $f$ trên $X_i$ và trên
    $X_i\cap X_j$, tương ứng; ta biết $\sh{F}'_i$ và $\sh{F}'_{ij}$
    là quasi-coherent \sref{I.1.6.3}. Đặt
    $\sh{G}=\bigoplus_i\sh{F}'_i$ và
    $\sh{H}=\bigoplus_{i,j}\sh{F}'_{ij}$; $\sh{G}$ và $\sh{H}$ là các
    $\sh{O}_Y$-môđun quasi-coherent. Ta sẽ định nghĩa một đồng cấu
    $u:\sh{G}\to\sh{H}$ sao cho $f_*(\sh{F})$ là \emph{hạt nhân} của
    $u$; từ đó suy ra $f_*(\sh{F})$ là quasi-coherent
    \sref{I.1.3.9}. Chỉ cần định nghĩa $u$
\oldpage[I]{172}
    như một đồng cấu tiền bó; xét đến các định nghĩa của $\sh{G}$ và
    $\sh{H}$, với mỗi tập mở $W\subset Y$, chỉ cần định nghĩa một đồng
    cấu
    \[
      u_W:\bigoplus_i\Gamma(f^{-1}(W)\cap X_i,\sh{F})\to
      \bigoplus_{i,j}\Gamma(f^{-1}(W)\cap X_i\cap X_j,\sh{F})
    \]
    thỏa các điều kiện tương thích thông thường khi $W$ thay đổi. Nếu,
    với mỗi tiết diện $s_i\in\Gamma(f^{-1}(W)\cap X_i,\sh{F})$, ký hiệu
    $s_{i|j}$ là hạn chế của nó trên
    $f^{-1}(W)\cap X_i\cap X_j$, ta đặt
    \[
      u_W\big((s_i)\big)=(s_{i|j}-s_{j|i})
    \]
    và các điều kiện tương thích hiển nhiên được thỏa mãn. Để chứng
    minh hạt nhân $\sh{R}$ của $u$ là $f_*(\sh{F})$, ta định nghĩa một
    đồng cấu từ $f_*(\sh{F})$ vào $\sh{R}$ bằng cách gán cho mỗi tiết
    diện $s\in\Gamma(f^{-1}(W),\sh{F})$ họ $(s_i)$, trong đó $s_i$ là
    hạn chế của $s$ trên $f^{-1}(W)\cap X_i$; các tiên đề (F1) và (F2)
    của bó (G,~II,~1.1) kéo theo đồng cấu này là \emph{song ánh}, và
    như vậy kết thúc chứng minh trong trường hợp này.
  \item[(b)] Trong trường hợp tổng quát, ta áp dụng cùng lập luận sau
    khi chứng minh được các $\sh{F}'_{ij}$ là quasi-coherent. Thật vậy,
    theo giả thiết, $X_i\cap X_j$ là một hợp hữu hạn của các tập mở
    affine $X_{ijk}$; và vì các $X_{ijk}$ là các tập mở affine
    \emph{trong một lược đồ}, giao của hai tập bất kỳ trong số đó vẫn
    là một tập mở affine \sref{I.5.5.6}. Do đó ta quy về trường hợp đầu
    tiên, và Mệnh đề~\sref{I.9.2.1} được chứng minh.
\end{enumerate}
\end{proof}

\begin{corollary}[9.2.2]
\label{I.9.2.2}
Kết luận của Mệnh đề~\sref{I.9.2.1} vẫn đúng trong mỗi trường hợp sau:
\begin{enumerate}
  \item[{\rm(a)}] $f$ là tách và quasi-compact;
  \item[{\rm(b)}] $f$ là tách và kiểu hữu hạn;
  \item[{\rm(c)}] $f$ là quasi-compact, và không gian nền của $X$ là noetherian địa phương.
\end{enumerate}
\end{corollary}

\begin{proof}
Trong trường hợp~(a), các $X_{\alpha i}\cap X_{\alpha j}$ là affine \sref{I.5.5.6}.
Trường hợp~(b) là một trường hợp riêng của trường hợp~(a) \sref{I.6.6.3}. Cuối cùng, trong trường hợp~(c), ta có thể quy về trường hợp $Y$ là affine và không gian nền của $X$ là noetherian; khi đó $X$ thừa nhận một phủ hữu hạn $(X_i)$ bởi các tập mở affine, và vì các $X_i\cap X_j$ là quasi-compact, chúng là các hợp hữu hạn của các tập mở affine \sref{I.2.1.3}.
\end{proof}

\subsection{Mở rộng các tiết diện của bó quasi-coherent}
\label{subsection:I.9.3}

\begin{theorem}[9.3.1]
\label{I.9.3.1}
Cho $X$ là một tiền lược đồ có không gian nền noetherian, hoặc một lược đồ có không gian
nền quasi-compact. Cho $\sh{L}$ là một $\sh{O}_X$-môđun khả nghịch \sref[0]{0.5.4.1}, $f$ là một
tiết diện của $\sh{L}$ trên $X$, $X_f$ là tập mở gồm các $x\in X$ sao cho $f(x)\neq0$
\sref[0]{0.5.5.1}, và $\sh{F}$ là một $\sh{O}_X$-môđun quasi-coherent.
\begin{enumerate}
  \item[{\rm(i)}] Nếu $s\in\Gamma(X,\sh{F})$ thỏa $s|X_f=0$, thì tồn tại một số nguyên $n>0$ sao cho $s\otimes f^{\otimes n}=0$.
  \item[{\rm(ii)}] Với mọi tiết diện $s\in\Gamma(X_f,\sh{F})$, tồn tại một số nguyên $n>0$ sao cho $s\otimes f^{\otimes n}$ mở rộng thành một tiết diện của $\sh{F}\otimes\sh{L}^{\otimes n}$ trên $X$.
\end{enumerate}
\end{theorem}

\begin{proof}
\medskip\noindent
\begin{enumerate}
  \item[(i)] Vì không gian nền của $X$ là quasi-compact, và do đó là hợp của
    hữu hạn tập mở affine $U_i$ sao cho $\sh{L}|U_i$ đẳng cấu với
    $\sh{O}_X|U_i$, ta có thể quy về trường hợp $X$ là affine và $\sh{L}=\sh{O}_X$.
    Trong trường hợp này, $f$ được đồng nhất với một phần tử của $A(X)$ và ta có
    $X_f=D(f)$; $s$ được đồng nhất với một phần tử của một $A(X)$-môđun $M$, còn
    $s|X_f$ với phần tử tương ứng của $M_f$; khi đó kết quả là hiển nhiên,
    theo định nghĩa của phép địa phương hóa môđun.
\oldpage[I]{173}
  \item[(ii)] Một lần nữa, $X$ là hợp hữu hạn của các tập mở affine $U_i$ ($1\leq i\leq r$)
    sao cho $\sh{L}|U_i\cong\sh{O}_X|U_i$, và, với mọi $i$,
    $(s\otimes f^{\otimes n})|(U_i\cap X_f)$ được đồng nhất (qua đẳng cấu nói trên)
    với $(f|(U_i\cap X_f))^n(s|(U_i\cap X_f))$. Khi đó ta biết
    \sref{I.1.4.1} rằng tồn tại một số nguyên $n>0$ sao cho, với mọi
    $i$, $(s\otimes f^{\otimes n})|(U_i\cap X_f)$ mở rộng thành một tiết diện $s_i$ của
    $\sh{F}\otimes\sh{L}^{\otimes n}$ trên $U_i$. Gọi $s_{i|j}$ là hạn chế
    của $s_i$ trên $U_i\cap U_j$; theo định nghĩa, ta có $s_{i|j}-s_{j|i}=0$ trên
    $X_f\cap U_i\cap U_j$. Nhưng nếu không gian nền của $X$ là noetherian, thì $U_i\cap U_j$ là
    quasi-compact; nếu $X$ là một lược đồ, thì $U_i\cap U_j$ là một tập mở affine
    \sref{I.5.5.6}, và do đó cũng quasi-compact. Theo (i), vì thế
    tồn tại một số nguyên $m$ (độc lập với $i$ và $j$) sao cho
    $(s_{i|j}-s_{j|i})\otimes f^{\otimes m}=0$. Suy ra ngay rằng
    tồn tại một tiết diện $s'$ của $\sh{F}\otimes\sh{L}^{\otimes(n+m)}$ trên $X$ có
    hạn chế bằng $s_i\otimes f^{\otimes m}$ trên mỗi $U_i$, và hạn chế bằng
    $s\otimes f^{\otimes(n+m)}$ trên $X_f$.
\end{enumerate}
\end{proof}

Các hệ quả sau đây cho một cách diễn giải Định lý~\sref{I.9.3.1} bằng ngôn ngữ
đại số hơn:
\begin{corollary}[9.3.2]
\label{I.9.3.2}
Với các giả thiết của \sref{I.9.3.1}, xét vành phân bậc $A_*=\Gamma_*(\sh{L})$
và $A_*$-môđun phân bậc $M_*=\Gamma_*(\sh{L},\sh{F})$ \sref[0]{0.5.4.6}. Nếu $f\in A_n$,
trong đó $n\in\bb{Z}$, thì ta có một đẳng cấu chính tắc
$\Gamma(X_f,\sh{F})\isoto((M_*)_f)_0$ (\emph{nhóm con của môđun đã địa phương hóa
$(M_*)_f$ gồm các phần tử bậc $0$}).
\end{corollary}

\begin{corollary}[9.3.3]
\label{I.9.3.3}
Giả sử các giả thiết của \sref{I.9.3.1} được thỏa mãn, và giả sử thêm rằng
$\sh{L}=\sh{O}_X$. Khi đó, nếu đặt $A=\Gamma(X,\sh{O}_X)$ và $M=\Gamma(X,\sh{F})$, thì $A_f$-môđun
$\Gamma(X_f,\sh{F})$ đẳng cấu chính tắc với $M_f$.
\end{corollary}

\begin{proposition}[9.3.4]
\label{I.9.3.4}
Cho $X$ là một tiền lược đồ noetherian, $\sh{F}$ là một $\sh{O}_X$-môđun coherent, và
$\sh{J}$ là một bó coherent các iđêan của $\sh{O}_X$, sao cho giá của $\sh{F}$ được chứa
trong giá của $\sh{O}_X/\sh{J}$. Khi đó tồn tại một số nguyên $n>0$ sao cho $\sh{J}^n\sh{F}=0$.
\end{proposition}

\begin{proof}
Vì $X$ là hợp hữu hạn của các tập mở affine có các vành noetherian, ta có thể giả sử
$X$ là affine với vành noetherian $A$; khi đó $\sh{F}=\widetilde{M}$, trong đó
$M=\Gamma(X,\sh{F})$ là một $A$-môđun hữu hạn sinh, và $\sh{J}=\widetilde{\mathfrak{J}}$,
trong đó $\mathfrak{J}=\Gamma(X,\sh{J})$ là một iđêan của $A$ (\sref{I.1.4.1} và
\sref{I.1.5.1}). Vì $A$ là noetherian, $\mathfrak{J}$ có một hệ sinh hữu hạn
$f_i$ ($1\leq i\leq m$). Theo giả thiết, mọi tiết diện của $\sh{F}$ trên
$X$ đều bằng không trên mỗi $D(f_i)$; nếu $s_j$ ($1\leq j\leq q$) là các tiết diện của
$\sh{F}$ sinh $M$, thì tồn tại một số nguyên $h$, độc lập với $i$ và $j$,
sao cho $f_i^h s_j=0$ \sref{I.1.4.1}, do đó $f_i^h s=0$ với mọi $s\in M$. Suy ra rằng,
nếu $n=mh$, thì $\mathfrak{J}^n M=0$, và do đó $\sh{O}_X$-môđun tương ứng
$\sh{J}^n\sh{F}=\widetilde{\mathfrak{J}^n M}$ \sref{I.1.3.13} bằng không.
\end{proof}

\begin{corollary}[9.3.5]
\label{I.9.3.5}
Theo các giả thiết của \sref{I.9.3.4}, tồn tại một tiền lược đồ con đóng $Y$ của $X$
có không gian nền là giá của $\sh{O}_X/\sh{J}$, sao cho nếu $j:Y\to X$ là phép nhúng
chính tắc thì $\sh{F}=j_*(j^*(\sh{F}))$.
\end{corollary}

\begin{proof}
Trước hết, lưu ý rằng các giá của $\sh{O}_X/\sh{J}$ và $\sh{O}_X/\sh{J}^n$ là như nhau,
vì nếu $\sh{J}_x=\sh{O}_x$ thì cũng có $\sh{J}_x^n=\sh{O}_x$, và mặt khác
$\sh{J}_x^n\subset\sh{J}_x$ với mọi $x\in X$. Do đó, nhờ \sref{I.9.3.4}, ta có thể giả
sử $\sh{J}\sh{F}=0$; khi ấy ta lấy $Y$ là tiền lược đồ con đóng của $X$ được xác định
bởi $\sh{J}$, và vì $\sh{F}$ khi ấy là một $(\sh{O}_X/\sh{J})$-môđun, kết luận là hiển nhiên.
\end{proof}

\oldpage[I]{174}
\subsection{Mở rộng các bó quasi-coherent}
\label{subsection:I.9.4}

\begin{env}[9.4.1]
\label{I.9.4.1}
Cho $X$ là một không gian tôpô, $\sh{F}$ là một bó các tập hợp (tương ứng, các nhóm, các
vành) trên $X$, $U$ là một tập mở của $X$, $\psi:U\to X$ là phép nhúng chính tắc, và
$\sh{G}$ là một bó con của $\sh{F}|U=\psi^*(\sh{F})$. Vì $\psi_*$ là một hàm tử khớp
trái, $\psi_*(\sh{G})$ là một bó con của $\psi_*(\psi^*(\sh{F}))$; ta ký hiệu $\rho$ là
đồng cấu chính tắc $\sh{F}\to\psi_*(\psi^*(\sh{F}))$ \sref[0]{0.3.5.3}, và ký hiệu
$\overline{\sh{G}}$ là bó con $\rho^{-1}(\psi_*(\sh{G}))$ của $\sh{F}$. Từ các định nghĩa
suy ra ngay rằng, với mọi tập mở $V$ của $X$, $\Gamma(V,\overline{\sh{G}})$ gồm các tiết
diện $s\in\Gamma(V,\sh{F})$ mà hạn chế xuống $V\cap U$ là một tiết diện của $\sh{G}$ trên
$V\cap U$. Do đó, $\overline{\sh{G}}|U=\psi^*(\overline{\sh{G}})=\sh{G}$, và
$\overline{\sh{G}}$ là bó con \emph{lớn nhất} của $\sh{F}$ có hạn chế bằng $\sh{G}$ trên
$U$; ta gọi $\overline{\sh{G}}$ là \emph{mở rộng chính tắc} của bó con $\sh{G}$ của
$\sh{F}|U$ thành một bó con của $\sh{F}$.
\end{env}

\begin{proposition}[9.4.2]
\label{I.9.4.2}
Cho $X$ là một tiền lược đồ và $U$ là một tập mở của $X$ sao cho phép nhúng chính tắc
$j:U\to X$ là một cấu xạ quasi-compact \emph{(điều này đúng với \emph{mọi} $U$ nếu
không gian nền của $X$ là \emph{noetherian địa phương}
\sref{I.6.6.4}[i])}. Khi đó:
\begin{enumerate}
  \item[{\rm(i)}] với mọi $(\sh{O}_X|U)$-môđun quasi-coherent $\sh{G}$, $j_*(\sh{G})$
    là một $\sh{O}_X$-môđun quasi-coherent, và
    $j_*(\sh{G})|U=j^*(j_*(\sh{G}))=\sh{G}$;
  \item[{\rm(ii)}] với mọi $\sh{O}_X$-môđun quasi-coherent $\sh{F}$ và mọi
    $(\sh{O}_X|U)$-môđun con quasi-coherent $\sh{G}$, mở rộng chính tắc
    $\overline{\sh{G}}$ của $\sh{G}$ \sref{I.9.4.1} là một
    $\sh{O}_X$-môđun con quasi-coherent của $\sh{F}$.
\end{enumerate}
\end{proposition}

\begin{proof}
Nếu $j=(\psi,\theta)$ (với $\psi$ là phép nhúng $U\to X$ của các không gian nền), thì
theo định nghĩa, $j_*(\sh{G})=\psi_*(\sh{G})$ với mọi
$(\sh{O}_X|U)$-môđun $\sh{G}$; hơn nữa,
$j^*(\sh{H})=\psi^*(\sh{H})=\sh{H}|U$ với mọi $\sh{O}_X$-môđun $\sh{H}$, theo định
nghĩa của tiền lược đồ cảm sinh trên một tập mở. Do đó, (i) là một trường hợp riêng
của \sref{I.9.2.2}[a]; vì cùng lý do, $j_*(j^*(\sh{F}))$ là quasi-coherent, và vì
$\overline{\sh{G}}$ là ảnh ngược của $j_*(\sh{G})$ qua đồng cấu
$\rho:\sh{F}\to j_*(j^*(\sh{F}))$, (ii) suy ra từ \sref{I.4.1.1}.
\end{proof}

Lưu ý rằng giả thiết cấu xạ $j:U\to X$ là quasi-compact cũng được thỏa mãn khi tập mở
$U$ là \emph{quasi-compact} và $X$ là một \emph{lược đồ}: thật vậy, khi đó $U$ là hợp
hữu hạn của các tập mở affine $U_i$, và với mọi tập mở affine $V$ của $X$,
$V\cap U_i$ là một tập mở affine \sref{I.5.5.6}, do đó quasi-compact.

\begin{corollary}[9.4.3]
\label{I.9.4.3}
Cho $X$ là một tiền lược đồ và $U$ là một tập mở quasi-compact của $X$ sao cho cấu xạ
nhúng $j:U\to X$ là quasi-compact. Giả sử thêm rằng mọi $\sh{O}_X$-môđun
quasi-coherent đều là giới hạn quy nạp của các $\sh{O}_X$-môđun con quasi-coherent
hữu hạn sinh của nó \emph{(điều này đúng khi $X$ là một \emph{lược đồ affine})}.
Cho $\sh{F}$ là một $\sh{O}_X$-môđun quasi-coherent và $\sh{G}$ là một
$(\sh{O}_X|U)$-môđun con quasi-coherent hữu hạn sinh của $\sh{F}|U$. Khi đó tồn tại
một $\sh{O}_X$-môđun con quasi-coherent hữu hạn sinh $\sh{G}'$ của $\sh{F}$ sao cho
$\sh{G}'|U=\sh{G}$.
\end{corollary}

\begin{proof}
Thật vậy, $\sh{G}=\overline{\sh{G}}|U$, và $\overline{\sh{G}}$ là quasi-coherent theo
\sref{I.9.4.2}, nên là giới hạn quy nạp của các $\sh{O}_X$-môđun con quasi-coherent
hữu hạn sinh $\sh{H}_\lambda$ của nó. Do đó, $\sh{G}$ là giới hạn quy nạp của các
$\sh{H}_\lambda|U$, và vì thế bằng một trong các $\sh{H}_\lambda|U$, do nó hữu hạn
sinh \sref[0]{0.5.2.3}.
\end{proof}

\begin{remark}[9.4.4]
\label{I.9.4.4}
Giả sử rằng với \emph{mọi} tập mở affine $U\subset X$, cấu xạ nhúng $U\to X$ là
quasi-compact. Khi đó, nếu kết luận của \sref{I.9.4.3} đúng với mọi tập mở affine $U$
và mọi $(\sh{O}_X|U)$-môđun con quasi-coherent hữu hạn sinh $\sh{G}$ của $\sh{F}|U$,
\oldpage[I]{175}
thì suy ra $\sh{F}$ là giới hạn quy nạp của các $\sh{O}_X$-môđun con quasi-coherent
hữu hạn sinh của nó. Thật vậy, với mọi tập mở affine $U\subset X$, ta có
$\sh{F}|U=\widetilde{M}$, trong đó $M$ là một $A(U)$-môđun, và vì môđun này là giới hạn
quy nạp của các môđun con hữu hạn sinh của nó, nên $\sh{F}|U$ là giới hạn quy nạp của
các $(\sh{O}_X|U)$-môđun con quasi-coherent hữu hạn sinh của nó \sref{I.1.3.9}.
Nhưng theo giả thiết, mỗi môđun con ấy được cảm sinh trên $U$ bởi một
$\sh{O}_X$-môđun con quasi-coherent hữu hạn sinh $\sh{G}_{\lambda,U}$ của $\sh{F}$.
Các tổng hữu hạn của những $\sh{G}_{\lambda,U}$ vẫn là các $\sh{O}_X$-môđun
quasi-coherent hữu hạn sinh, bởi vì vấn đề có tính địa phương và trường hợp $X$ là
affine đã được xét trong \sref{I.1.3.10}; khi đó rõ ràng $\sh{F}$ là giới hạn quy nạp
của các tổng hữu hạn này, suy ra khẳng định.
\end{remark}

\begin{corollary}[9.4.5]
\label{I.9.4.5}
Dưới các giả thiết của \sref{I.9.4.3}, với mọi $(\sh{O}_X|U)$-môđun quasi-coherent
hữu hạn sinh $\sh{G}$, tồn tại một $\sh{O}_X$-môđun quasi-coherent hữu hạn sinh
$\sh{G}'$ sao cho $\sh{G}'|U=\sh{G}$.
\end{corollary}

\begin{proof}
Vì $\sh{F}=j_*(\sh{G})$ là quasi-coherent \sref{I.9.4.2} và
$\sh{F}|U=\sh{G}$, chỉ cần áp dụng \sref{I.9.4.3} cho $\sh{F}$.
\end{proof}

\begin{lemma}[9.4.6]
\label{I.9.4.6}
Cho $X$ là một tiền lược đồ, $L$ là một tập được sắp thứ tự tốt,
$(V_\lambda)_{\lambda\in L}$ là một phủ của $X$ bởi các tập mở affine, và $U$ là một
tập mở của $X$; với mọi $\lambda\in L$, đặt
$W_\lambda=\bigcup_{\mu<\lambda}V_\mu$. Giả sử rằng:
\begin{enumerate}
  \item[1\textsuperscript{o}] Với mọi $\lambda\in L$,
    $V_\lambda\cap W_\lambda$ là quasi-compact;
  \item[2\textsuperscript{o}] Cấu xạ nhúng $U\to X$ là quasi-compact.
\end{enumerate}
Khi đó, với mọi $\sh{O}_X$-môđun quasi-coherent $\sh{F}$ và mọi
$(\sh{O}_X|U)$-môđun con quasi-coherent hữu hạn sinh $\sh{G}$ của $\sh{F}|U$, tồn tại
một $\sh{O}_X$-môđun con quasi-coherent hữu hạn sinh $\sh{G}'$ của $\sh{F}$ sao cho
$\sh{G}'|U=\sh{G}$.
\end{lemma}

\begin{proof}
Cho $\kappa$ là kiểu thứ tự của $L$ và đánh chỉ số phủ đã cho bởi
$\lambda<\kappa$.\footnote{Chú thích biên tập: phép quy nạp trong bản in bỏ sót bước đầu và, nếu $L$ có phần tử lớn nhất, bước mở rộng cuối cùng. Các bước $0\leq\lambda\leq\kappa$ hoàn chỉnh chứng minh mà không thay đổi các giả thiết.}
Với $0\leq\lambda\leq\kappa$, đặt
$W_\lambda=\bigcup_{\mu<\lambda}V_\mu$ và $U_\lambda=U\cup W_\lambda$.
Bằng quy nạp siêu hạn, ta xây dựng một
$(\sh{O}_X|U_\lambda)$-môđun con quasi-coherent hữu hạn sinh $\sh{G}'_\lambda$ của
$\sh{F}|U_\lambda$, sao cho $\sh{G}'_\lambda|U_\mu=\sh{G}'_\mu$ với $\mu<\lambda$ và
$\sh{G}'_\lambda|U=\sh{G}$.
Ta bắt đầu với $\sh{G}'_0=\sh{G}$ trên $U_0=U$; vì $U_\kappa=X$, môđun con
$\sh{G}'=\sh{G}'_\kappa$ sẽ cho phép mở rộng cần tìm.
Bây giờ, giả sử $0<\lambda\leq\kappa$ và các $\sh{G}'_\mu$ đã được xây dựng
với các tính chất này cho $\mu<\lambda$.
Nếu $\lambda$ là một số thứ tự giới hạn, ta lấy $\sh{G}'_\lambda$ là
$(\sh{O}_X|U_\lambda)$-môđun con duy nhất của $\sh{F}|U_\lambda$ sao cho
$\sh{G}'_\lambda|U_\mu=\sh{G}'_\mu$ với mọi $\mu<\lambda$.
Khi đó các $U_\mu$ với $\mu<\lambda$ phủ $U_\lambda$, nên ta có thể dán
các môđun con tương thích \sref[0]{0.3.3.1}; tính quasi-coherent và tính hữu hạn sinh
là các tính chất địa phương.
Ngược lại, nếu $\lambda=\mu+1$, thì
$U_\lambda=U_\mu\cup V_\mu$, và chỉ cần định nghĩa một
$(\sh{O}_X|V_\mu)$-môđun con quasi-coherent hữu hạn sinh $\sh{G}''_\mu$ của
$\sh{F}|V_\mu$ sao cho
\[
  \sh{G}''_\mu|(U_\mu\cap V_\mu)=\sh{G}'_\mu|(U_\mu\cap V_\mu);
\]
sau đó lấy $\sh{G}'_\lambda$ là $(\sh{O}_X|U_\lambda)$-môđun con của
$\sh{F}|U_\lambda$ sao cho $\sh{G}'_\lambda|U_\mu=\sh{G}'_\mu$ và
$\sh{G}'_\lambda|V_\mu=\sh{G}''_\mu$ \sref[0]{0.3.3.1}. Nhưng vì $V_\mu$ là affine,
sự tồn tại của $\sh{G}''_\mu$ được bảo đảm bởi \sref{I.9.4.3} ngay khi ta chứng minh
được rằng $U_\mu\cap V_\mu$ là quasi-compact; mà $U_\mu\cap V_\mu$ là hợp của
$U\cap V_\mu$ và $W_\mu\cap V_\mu$, cả hai đều quasi-compact theo các giả thiết.
\end{proof}

\begin{theorem}[9.4.7]
\label{I.9.4.7}
Cho $X$ là một tiền lược đồ và $U$ là một tập mở của $X$. Giả sử một trong các điều
kiện sau được thỏa mãn:
\begin{enumerate}
  \item[{\rm(a)}] Không gian nền của $X$ là noetherian địa phương.
  \item[{\rm(b)}] $X$ là một lược đồ quasi-compact và $U$ là một tập mở quasi-compact.
\end{enumerate}
Khi đó, với mọi $\sh{O}_X$-môđun quasi-coherent $\sh{F}$ và mọi
$(\sh{O}_X|U)$-môđun con quasi-coherent hữu hạn sinh $\sh{G}$ của $\sh{F}|U$, tồn tại
một $\sh{O}_X$-môđun con quasi-coherent hữu hạn sinh $\sh{G}'$ của $\sh{F}$ sao cho
$\sh{G}'|U=\sh{G}$.
\end{theorem}

\oldpage[I]{176}
\begin{proof}
Cho $(V_\lambda)_{\lambda\in L}$ là một phủ của $X$ bởi các tập mở affine, trong đó
$L$ được giả sử hữu hạn ở trường hợp (b); trang bị cho $L$ một cấu trúc thứ tự tốt,
chỉ cần kiểm tra rằng các điều kiện của \sref{I.9.4.6} được thỏa mãn. Điều này hiển
nhiên dưới giả thiết (a), vì các không gian $V_\lambda$ là noetherian. Dưới giả thiết
(b), các $V_\lambda\cap V_\mu$ là affine \sref{I.5.5.6}, nên quasi-compact; vì $L$
hữu hạn, $V_\lambda\cap W_\lambda$ là quasi-compact. Do đó định lý được chứng minh.
\end{proof}

\begin{corollary}[9.4.8]
\label{I.9.4.8}
Dưới các giả thiết của \sref{I.9.4.7}, với mọi $(\sh{O}_X|U)$-môđun quasi-coherent
hữu hạn sinh $\sh{G}$, tồn tại một $\sh{O}_X$-môđun quasi-coherent hữu hạn sinh
$\sh{G}'$ sao cho $\sh{G}'|U=\sh{G}$.
\end{corollary}

\begin{proof}
Chỉ cần áp dụng \sref{I.9.4.7} cho $\sh{F}=j_*(\sh{G})$, là môđun quasi-coherent
\sref{I.9.4.2} và thỏa mãn $\sh{F}|U=\sh{G}$.
\end{proof}

\begin{corollary}[9.4.9]
\label{I.9.4.9}
Cho $X$ là một tiền lược đồ có không gian nền noetherian địa phương, hoặc là một lược
đồ quasi-compact. Khi đó mọi $\sh{O}_X$-môđun quasi-coherent đều là giới hạn quy nạp
của các $\sh{O}_X$-môđun con quasi-coherent hữu hạn sinh của nó.
\end{corollary}

\begin{proof}
Điều này suy ra từ \sref{I.9.4.7} và nhận xét \sref{I.9.4.4}.
\end{proof}

\begin{corollary}[9.4.10]
\label{I.9.4.10}
Dưới các giả thiết của \sref{I.9.4.9}, nếu một $\sh{O}_X$-môđun quasi-coherent
$\sh{F}$ có tính chất là mọi $\sh{O}_X$-môđun con quasi-coherent hữu hạn sinh của $\sh{F}$
đều được sinh bởi các tiết diện của nó trên $X$, thì $\sh{F}$ được sinh bởi các
tiết diện của nó trên $X$.
\end{corollary}

\begin{proof}
Cho $U$ là một lân cận mở affine của một điểm $x\in X$, và cho $s$ là một tiết diện của
$\sh{F}$ trên $U$; $\sh{O}_X$-môđun con $\sh{G}$ của $\sh{F}|U$ được sinh bởi $s$ là
quasi-coherent và hữu hạn sinh, do đó tồn tại một $\sh{O}_X$-môđun con quasi-coherent
hữu hạn sinh $\sh{G}'$ của $\sh{F}$ sao cho $\sh{G}'|U=\sh{G}$ \sref{I.9.4.7}. Theo
giả thiết, khi đó tồn tại một số hữu hạn tiết diện $t_i$ của $\sh{G}'$ trên $X$ và
các tiết diện $a_i$ của $\sh{O}_X$ trên một lân cận $V\subset U$ của $x$ sao cho
$s|V=\sum_i a_i(t_i|V)$, điều này chứng minh hệ quả.
\end{proof}

\subsection{Ảnh đóng của một tiền lược đồ; bao đóng của một tiền lược đồ con}
\label{subsection:I.9.5}

\begin{proposition}[9.5.1]
\label{I.9.5.1}
Cho $f:X\to Y$ là một cấu xạ giữa các tiền lược đồ sao cho $f_*(\sh{O}_X)$ là một
$\sh{O}_Y$-môđun quasi-coherent (điều này xảy ra nếu $f$ quasi-compact và, ngoài ra,
hoặc $f$ tách được hoặc $X$ là noetherian địa phương \sref{I.9.2.2}). Khi đó tồn tại
một tiền lược đồ con đóng nhỏ nhất $Y'$ của $Y$ sao cho $f$ phân tích qua phép nhúng chính tắc
$j:Y'\to Y$ (\emph{hoặc, tương đương \sref{I.4.4.1}, sao cho tiền lược đồ con
$f^{-1}(Y')$ của $X$ là \emph{đồng nhất} với $X$}).
\end{proposition}

Cụ thể hơn:
\begin{corollary}[9.5.2]
\label{I.9.5.2}
Dưới các điều kiện của \sref{I.9.5.1}, cho $f=(\psi,\theta)$, và cho $\sh{J}$ là
hạt nhân (quasi-coherent) của đồng cấu $\theta:\sh{O}_Y\to f_*(\sh{O}_X)$. Khi đó tiền
lược đồ con đóng $Y'$ của $Y$ xác định bởi $\sh{J}$ thỏa mãn các điều kiện của
\sref{I.9.5.1}.
\end{corollary}

\begin{proof}
Vì $\psi^*$ là một hàm tử khớp, phân tích chính tắc
$\theta:\sh{O}_Y\to\sh{O}_Y/\sh{J}\xrightarrow{\theta'}\psi_*(\sh{O}_X)$ cho (\textbf{0},~3.5.4.3)
một phân tích
$\theta^\sharp:\psi^*(\sh{O}_Y)\to\psi^*(\sh{O}_Y)/\psi^*(\sh{J})
  \xrightarrow{{\theta'}^\sharp}\sh{O}_X$; vì $\theta_x^\sharp$ là một đồng cấu địa phương với
mọi $x\in X$, điều tương tự đúng với ${\theta_x'}^\sharp$; nếu ký hiệu $\psi_0$ là
ánh xạ liên tục $\psi$ được xem như ánh xạ từ $X$ vào \erratum[II]{$Y'$}, và $\theta_0$ là hạn chế
$\theta'|Y':(\sh{O}_Y/\sh{J})|Y'\to\psi_*(\sh{O}_X)|Y'=(\psi_0)_*(\sh{O}_X)$, ta thấy rằng
$f_0=(\psi_0,\theta_0)$ là một cấu xạ của các tiền lược đồ $X\to Y'$ \sref{I.2.2.1} sao cho
$f=j\circ f_0$. Bây giờ, nếu $Y''$
\oldpage[I]{177}
là một tiền lược đồ con đóng thứ hai của $Y$, được xác định bởi một bó iđêan quasi-coherent $\sh{J}'$ của
$\sh{O}_Y$, và sao cho $f$ phân tích qua phép nhúng $j':Y''\to Y$, trước hết phải có
$\psi(X)\subset Y''$ và do đó $Y'\subset Y''$, vì $Y''$ là đóng.
Hơn nữa, với mọi $y\in Y''$, $\theta$ phải phân tích thành
$\sh{O}_y\to\sh{O}_y/\sh{J}'_y\to(\psi_*(\sh{O}_X))_y$, điều này theo định nghĩa kéo theo
$\sh{J}'_y\subset\sh{J}_y$, và do đó $Y'$ là một tiền lược đồ con đóng của $Y''$
\sref{I.4.1.10}.
\end{proof}

\begin{definition}[9.5.3]
\label{I.9.5.3}
Khi tồn tại một tiền lược đồ con đóng nhỏ nhất $Y'$ của $Y$ sao cho $f$ phân tích qua
phép nhúng chính tắc $j:Y'\to Y$, ta gọi $Y'$ là tiền lược đồ
\emph{ảnh đóng} của $X$ dưới cấu xạ $f$.
\end{definition}

\begin{proposition}[9.5.4]
\label{I.9.5.4}
Nếu $f_*(\sh{O}_X)$ là một $\sh{O}_Y$-môđun quasi-coherent, thì không gian nền của
ảnh đóng của $X$ dưới cấu xạ $f$ là bao đóng $\overline{f(X)}$ trong $Y$.
\end{proposition}

\begin{proof}
Vì giá của $f_*(\sh{O}_X)$ nằm trong $\overline{f(X)}$, ta có (với ký hiệu của
\sref{I.9.5.2}) $\sh{J}_y=\sh{O}_y$ với $y\not\in\overline{f(X)}$, do đó giá của
$\sh{O}_Y/\sh{J}$ nằm trong $\overline{f(X)}$. Ngoài ra, giá này đóng và chứa
$f(X)$: thật vậy, nếu $y\in f(X)$, phần tử đơn vị của vành
$(\psi_*(\sh{O}_X))_y$ khác không, vì là mầm tại $y$ của tiết diện
\[
  1\in\Gamma(X,\sh{O}_X)=\Gamma(Y,\psi_*(\sh{O}_X));
\]
vì nó là ảnh qua $\theta$ của phần tử đơn vị của $\sh{O}_y$, phần tử sau không
thuộc $\sh{J}_y$, do đó $\sh{O}_y/\sh{J}_y\neq 0$; chứng minh kết thúc.
\end{proof}

\begin{proposition}[9.5.5]
\label{I.9.5.5}
\emph{(Tính bắc cầu của các ảnh đóng).} Cho $f:X\to Y$ và $g:Y\to Z$ là hai cấu xạ
giữa các tiền lược đồ; giả sử ảnh đóng $Y'$ của $X$ dưới $f$ tồn tại, và nếu $g'$ là
hạn chế của $g$ lên $Y'$, thì ảnh đóng $Z'$ của $Y'$ dưới $g'$ tồn tại. Khi đó ảnh đóng
của $X$ dưới $g\circ f$ tồn tại và bằng $Z'$.
\end{proposition}

\begin{proof}
Chỉ cần dùng \sref{I.9.5.1} để chứng minh rằng $Z'$ là tiền lược đồ con đóng nhỏ nhất
$Z_1$ của $Z$ sao cho tiền lược đồ con đóng $(g\circ f)^{-1}(Z_1)$ của $X$
(bằng $f^{-1}(g^{-1}(Z_1))$ theo Hệ quả~\sref{I.4.4.2}) bằng $X$; điều này tương đương
với việc nói rằng $Z'$ là tiền lược đồ con đóng nhỏ nhất của $Z$ sao cho $f$ phân tích qua
phép nhúng $g^{-1}(Z_1)\to Y$ \sref{I.4.4.1}. Do sự tồn tại của ảnh đóng $Y'$, mọi $Z_1$
có tính chất này đều sao cho $g^{-1}(Z_1)$ phân tích qua $Y'$, điều này tương đương với
việc $j^{-1}(g^{-1}(Z_1))=g'^{-1}(Z_1)=Y'$, trong đó $j$ là phép nhúng $Y'\to Y$.
Theo định nghĩa của $Z'$, ta kết luận rằng $Z'$ là tiền lược đồ con đóng nhỏ nhất của $Z$
thỏa mãn điều kiện trước đó.
\end{proof}

\begin{corollary}[9.5.6]
\label{I.9.5.6}
Cho $f:X\to Y$ là một cấu xạ trên $S$ sao cho $Y$ là ảnh đóng của $X$ dưới $f$.
Cho $Z$ là một lược đồ trên $S$; nếu hai cấu xạ trên $S$ $g_1$, $g_2$ từ $Y$ vào $Z$
thỏa mãn $g_1\circ f=g_2\circ f$, thì $g_1=g_2$.
\end{corollary}

\begin{proof}
Cho $h=(g_1,g_2)_S:Y\to Z\times_S Z$; vì đường chéo $T=\Delta_Z(Z)$ là một
tiền lược đồ con đóng của $Z\times_S Z$, $Y'=h^{-1}(T)$ là một tiền lược đồ con đóng của $Y$
\sref{I.4.4.1}. Đặt $u=g_1\circ f=g_2\circ f$; khi đó, theo định nghĩa của tích,
$h'=h\circ f=(u,u)_S$, nên $h\circ f=\Delta_Z\circ u$; vì
$\Delta_Z^{-1}(T)=Z$, ta có $h'^{-1}(T)=u^{-1}(Z)=X$, do đó $f^{-1}(Y')=X$.
Từ đó, theo \sref{I.4.4.1}, ta kết luận rằng $f$ phân tích qua phép nhúng chính tắc $Y'\to Y$,
do đó $Y'=Y$ theo giả thiết; tiếp theo, theo \sref{I.4.4.1}, $h$ phân tích thành
$\Delta_Z\circ v$, trong đó $v$ là một cấu xạ $Y\to Z$, suy ra $g_1=g_2=v$.
\end{proof}

\begin{remark}[9.5.7]
\label{I.9.5.7}
Nếu $X$ và $Y$ là các lược đồ trên $S$, mệnh đề~\sref{I.9.5.6} cho thấy rằng,
\oldpage[I]{178}
khi $Y$ là ảnh đóng của $X$ dưới $f$, $f$ là một \emph{toàn cấu} trong phạm trù các \emph{lược đồ trên $S$} (T,~1.1). Chúng ta sẽ chứng minh trong Chương~V rằng, ngược lại, nếu ảnh đóng $Y'$ của $X$ dưới $f$ tồn tại và $f$ là một toàn cấu của các lược đồ trên $S$, thì tất yếu $Y'=Y$.
\end{remark}

\begin{proposition}[9.5.8]
\label{I.9.5.8}
Giả sử các giả thiết của \sref{I.9.5.1} được thỏa mãn, và cho $Y'$ là ảnh đóng của $X$ qua $f$. Với mọi tập mở $V$ của $Y$, đặt
$f_V:f^{-1}(V)\to V$ là hạn chế của $f$; khi đó ảnh đóng của
$f^{-1}(V)$ qua $f_V$ trong $V$ tồn tại và bằng tiền lược đồ cảm sinh bởi
$Y'$ trên tập mở $V\cap Y'$ của $Y'$ \emph{(nói cách khác, bằng tiền lược đồ con
$\inf(V,Y')$ của $Y$ \sref{I.4.4.3})}.
\end{proposition}

\begin{proof}
Đặt $X'=f^{-1}(V)$; vì ảnh trực tiếp của $\sh{O}_{X'}$ qua
$f_V$ chính là hạn chế của $f_*(\sh{O}_X)$ trên $V$, nên rõ ràng
hạt nhân $\sh{J}'$ của đồng cấu
$\sh{O}_V\to(f_V)_*(\sh{O}_{X'})$ là hạn chế của
$\sh{J}$ trên $V$, từ đó mệnh đề suy ra ngay.
\end{proof}

Ta sẽ thấy rằng kết quả này có thể được hiểu là nói rằng việc lấy ảnh đóng giao hoán với một mở rộng $Y_1\to Y$ của tiền lược đồ cơ sở, là một \emph{phép nhúng mở}.
Trong Chương~IV, ta sẽ thấy điều tương tự xảy ra đối với một mở rộng $Y_1\to Y$ là một cấu xạ \emph{phẳng}, với điều kiện $f$ tách được và quasi-compact.

\begin{proposition}[9.5.9]
\label{I.9.5.9}
Cho $f:X\to Y$ là một cấu xạ sao cho ảnh đóng $Y'$ của $X$ qua $f$ tồn tại.
\begin{enumerate}
  \item[{\rm(i)}] Nếu $X$ là giảm, thì $Y'$ cũng vậy.
  \item[{\rm(ii)}] Nếu các giả thiết của Mệnh đề~\sref{I.9.5.1} được thỏa mãn và $X$ là bất khả quy (tương ứng, nguyên), thì $Y'$ cũng vậy.
\end{enumerate}
\end{proposition}

\begin{proof}
Theo giả thiết, cấu xạ $f$ phân tích qua $X\xrightarrow{g}Y'\xrightarrow{j}Y$, trong đó $j$ là phép nhúng chính tắc.
Vì $X$ là giảm, $g$ phân tích qua $X\xrightarrow{h}Y'_{\mathrm{red}}\xrightarrow{j'}Y'$, trong đó $j'$ là phép nhúng chính tắc \sref{I.5.2.2}; từ định nghĩa của $Y'$ suy ra $Y'_{\mathrm{red}}=Y'$.
Hơn nữa, nếu các điều kiện của Mệnh đề~\sref{I.9.5.1} được thỏa mãn, thì từ \sref{I.9.5.4} suy ra $f(X)$ trù mật trong $Y'$; nếu $X$ bất khả quy, thì $Y'$ cũng vậy \sref[0]{0.2.1.5}.
Khẳng định về các tiền lược đồ nguyên suy ra từ việc kết hợp hai khẳng định kia.
\end{proof}

\begin{proposition}[9.5.10]
\label{I.9.5.10}
Cho $Y$ là một tiền lược đồ con của một tiền lược đồ $X$, sao cho phép nhúng chính tắc $i:Y\to X$ là một cấu xạ quasi-compact.
Khi đó tồn tại một tiền lược đồ con đóng nhỏ nhất $\overline{Y}$ của $X$ chứa $Y$; không gian nền của nó là bao đóng của không gian nền của $Y$; không gian sau mở trong bao đóng của nó, và tiền lược đồ $Y$ được cảm sinh trên tập mở này bởi $\overline{Y}$.
\end{proposition}

\begin{proof}
Chỉ cần áp dụng Mệnh đề~\sref{I.9.5.1} cho phép nhúng $i$\footnote{\emph{[Dịch.] Nguồn tiếng Pháp in $j$ ở đây, nhưng mệnh đề định nghĩa phép nhúng chính tắc là $i:Y\to X$.}}, vốn là cấu xạ tách được \sref{I.5.5.1} và quasi-compact theo giả thiết; do đó \sref{I.9.5.1} chứng minh sự tồn tại của $\overline{Y}$, và \sref{I.9.5.4} cho thấy không gian nền của nó là bao đóng của $Y$ trong $X$; vì $Y$ đóng địa phương trong $X$, nó mở trong $\overline{Y}$, và khẳng định cuối cùng là kết quả của việc áp dụng \sref{I.9.5.8} cho một tập mở $V$ của $X$ sao cho $Y$ đóng trong $V$.
\end{proof}

Với các ký hiệu trên, nếu phép nhúng $V\to X$ là quasi-compact, và nếu $\sh{J}$ là bó iđêan quasi-coherent của $\sh{O}_X|V$ xác định tiền lược đồ con đóng $Y$ của $V$, thì từ Mệnh đề~\sref{I.9.5.1} suy ra rằng bó iđêan quasi-coherent của $\sh{O}_X$ xác định $\overline{Y}$ là mở rộng chính tắc \sref{I.9.4.1} $\overline{\sh{J}}$ của $\sh{J}$, vì rõ ràng đó là bó con quasi-coherent iđêan lớn nhất của $\sh{O}_X$ cảm sinh $\sh{J}$ trên $V$.

\oldpage[I]{179}
\begin{corollary}[9.5.11]
\label{I.9.5.11}
Dưới các giả thiết của Mệnh đề~\sref{I.9.5.10}, mọi tiết diện của
$\sh{O}_{\overline{Y}}$ trên một tập mở $V$ của $\overline{Y}$ mà bằng không trên
$V\cap Y$ thì bằng không.
\end{corollary}

\begin{proof}
Theo Mệnh đề~\sref{I.9.5.8}, ta có thể quy về trường hợp $V=\overline{Y}$.
Nếu lưu ý rằng các tiết diện của $\sh{O}_{\overline{Y}}$ trên $\overline{Y}$
tương ứng một cách chính tắc với các tiết diện trên $\overline{Y}$ của
$\overline{Y}\otimes_{\mathbb Z}\mathbb Z[T]$ \sref{I.3.3.15} và rằng đối tượng
sau tách được trên $\overline{Y}$, thì hệ quả này là một trường hợp riêng của
\sref{I.9.5.6}.
\end{proof}

Khi tồn tại một tiền lược đồ con đóng nhỏ nhất $Y'$ của $X$ chứa một tiền lược đồ
con $Y$ của $X$, ta gọi $Y'$ là \emph{bao đóng} của $Y$ trong $X$, khi không gây
nhầm lẫn.

\subsection{Các bó đại số quasi-coherent; thay đổi bó cấu trúc}
\label{subsection:I.9.6}
\label{I.9.6}

\begin{proposition}[9.6.1]
\label{I.9.6.1}
Cho $X$ là một tiền lược đồ và $\sh{B}$ là một $\sh{O}_X$-đại số quasi-coherent \sref[0]{0.5.1.3}.
Để một $\sh{B}$-môđun $\sh{F}$ là quasi-coherent (trên không gian vành $(X,\sh{B})$), điều kiện cần và đủ là $\sh{F}$ là một $\sh{O}_X$-môđun quasi-coherent.
\end{proposition}

\begin{proof}
Vì vấn đề là địa phương, ta có thể giả sử $X$ là affine, được cho bởi một vành $A$, và do đó $\sh{B}=\widetilde{B}$, trong đó $B$ là một $A$-đại số \sref{I.1.4.3}.
Nếu $\sh{F}$ là quasi-coherent trên không gian vành $(X,\sh{B})$, ta cũng có thể giả sử $\sh{F}$ là đối hạt nhân của một $\sh{B}$-đồng cấu $\sh{B}^{(I)}\to\sh{B}^{(J)}$; vì đồng cấu này cũng là một $\sh{O}_X$-đồng cấu của các $\sh{O}_X$-môđun, và $\sh{B}^{(I)}$, $\sh{B}^{(J)}$ là các $\sh{O}_X$-môđun quasi-coherent \sref{I.1.3.9}[ii], nên $\sh{F}$ cũng là một $\sh{O}_X$-môđun quasi-coherent \sref{I.1.3.9}[i].

Ngược lại, nếu $\sh{F}$ là một $\sh{O}_X$-môđun quasi-coherent, thì $\sh{F}=\widetilde{M}$, trong đó $M$ là một $B$-môđun \sref{I.1.4.3}; $M$ đẳng cấu với đối hạt nhân của một $B$-đồng cấu $B^{(I)}\to B^{(J)}$, nên $\sh{F}$ là một $\sh{B}$-môđun đẳng cấu với đối hạt nhân của đồng cấu tương ứng $\sh{B}^{(I)}\to\sh{B}^{(J)}$ \sref{I.1.3.13}, điều này kết thúc chứng minh.
\end{proof}

Đặc biệt, nếu $\sh{F}$ và $\sh{G}$ là hai $\sh{B}$-môđun quasi-coherent, thì $\sh{F}\otimes_{\sh{B}}\sh{G}$ là một $\sh{B}$-môđun quasi-coherent; tương tự đối với $\shHom_{\sh{B}}(\sh{F},\sh{G})$ nếu thêm vào đó ta giả sử $\sh{F}$ có một trình bày hữu hạn \sref{I.1.3.13}.

\begin{env}[9.6.2]
\label{I.9.6.2}
Cho một tiền lược đồ $X$, ta nói một $\sh{O}_X$-đại số quasi-coherent $\sh{B}$ là thuộc \emph{loại hữu hạn} nếu với mọi $x\in X$, tồn tại một lân cận mở \emph{affine} $U$ của $x$ sao cho $\Gamma(U,\sh{B})=B$ là một đại số thuộc loại hữu hạn trên $\Gamma(U,\sh{O}_X)=A$.
Khi đó $\sh{B}|U=\widetilde{B}$ và với mọi $f\in A$, $(\sh{O}_X|D(f))$-đại số cảm sinh $\sh{B}|D(f)$ thuộc loại hữu hạn, vì nó đẳng cấu với $(B_f)^{\sim}$, và $B_f=B\otimes_A A_f$ hiển nhiên là một đại số thuộc loại hữu hạn trên $A_f$.
Vì các $D(f)$ tạo thành một cơ sở của tôpô của $U$, ta kết luận rằng nếu $\sh{B}$ là một $\sh{O}_X$-đại số quasi-coherent thuộc loại hữu hạn thì với mọi tập mở $V$ của $X$, $\sh{B}|V$ là một $(\sh{O}_X|V)$-đại số quasi-coherent thuộc loại hữu hạn.
\end{env}

\begin{proposition}[9.6.3]
\label{I.9.6.3}
Cho $X$ là một tiền lược đồ địa phương noetherian.
Khi đó mọi $\sh{O}_X$-đại số quasi-coherent $\sh{B}$ thuộc loại hữu hạn là một bó coherent các vành \sref[0]{0.5.3.7}.
\end{proposition}

\begin{proof}
Có thể lại hạn chế về trường hợp $X$ là một lược đồ affine cho bởi một vành noetherian $A$, và $\sh{B}=\widetilde{B}$, trong đó $B$ là một $A$-đại số thuộc loại hữu hạn; khi đó $B$ là một vành noetherian.
Với điều này, cần chứng minh rằng hạt nhân $\sh{N}$ của một $\sh{B}$-đồng cấu $\sh{B}^m\to\sh{B}$ là một $\sh{B}$-
\oldpage[I]{180}
môđun thuộc loại hữu hạn; nhưng nó đẳng cấu (với tư cách $\sh{B}$-môđun) với $\widetilde{N}$, trong đó $N$ là hạt nhân của đồng cấu tương ứng giữa các $B$-môđun $B^m\to B$ \sref{I.1.3.13}.
Vì $B$ là noetherian, môđun con $N$ của $B^m$ là một $B$-môđun thuộc loại hữu hạn, nên tồn tại một đồng cấu $B^p\to B^m$ có ảnh là $N$; vì dãy $B^p\to B^m\to B$ là khớp, dãy tương ứng $\sh{B}^p\to\sh{B}^m\to\sh{B}$ cũng khớp \sref{I.1.3.5}, và vì $\sh{N}$ là ảnh của $\sh{B}^p\to\sh{B}^m$ \sref{I.1.3.9}[i], điều này chứng minh mệnh đề.
\end{proof}

\begin{corollary}[9.6.4]
\label{I.9.6.4}
Dưới các giả thiết của \sref{I.9.6.3}, để một $\sh{B}$-môđun $\sh{F}$ là coherent, điều kiện cần và đủ là nó là một $\sh{O}_X$-môđun quasi-coherent và một $\sh{B}$-môđun thuộc loại hữu hạn.
Nếu vậy, và nếu $\sh{G}$ là một $\sh{B}$-môđun con hoặc một $\sh{B}$-môđun thương của $\sh{F}$, thì để $\sh{G}$ là một $\sh{B}$-môđun coherent, điều kiện cần và đủ là nó là một $\sh{O}_X$-môđun quasi-coherent.
\end{corollary}

\begin{proof}
Xét \sref{I.9.6.1}, các điều kiện trên $\sh{F}$ rõ ràng là cần thiết; ta sẽ chứng minh chúng là đủ.
Ta có thể hạn chế về trường hợp $X$ là affine, cho bởi một vành noetherian $A$, trong đó $\sh{B}=\widetilde{B}$, với $B$ là một $A$-đại số thuộc loại hữu hạn, $\sh{F}=\widetilde{M}$, với $M$ là một $B$-môđun, và tồn tại một $\sh{B}$-đồng cấu toàn ánh $\sh{B}^m\to\sh{F}\to0$.
Khi đó ta có dãy khớp tương ứng $B^m\to M\to0$, nên $M$ là một $B$-môđun thuộc loại hữu hạn; hơn nữa, hạt nhân $P$ của đồng cấu $B^m\to M$ khi đó là một $B$-môđun thuộc loại hữu hạn, vì $B$ là noetherian.
Từ \sref{I.1.3.13} suy ra $\sh{F}$ là đối hạt nhân của một $\sh{B}$-đồng cấu $\sh{B}^n\to\sh{B}^m$, và do đó là coherent, vì $\sh{B}$ là một bó coherent các vành \sref[0]{0.5.3.4}.
Cùng lập luận đó cho thấy một $\sh{B}$-môđun con quasi-coherent (resp. một $\sh{B}$-môđun thương) của $\sh{F}$ thuộc loại hữu hạn, từ đó suy ra phần thứ hai của hệ quả.
\end{proof}

\begin{proposition}[9.6.5]
\label{I.9.6.5}
Cho $X$ là một lược đồ quasi-compact, hoặc một tiền lược đồ có không gian nền noetherian.
Với mọi $\sh{O}_X$-đại số quasi-coherent $\sh{B}$ thuộc loại hữu hạn, tồn tại một $\sh{O}_X$-môđun con quasi-coherent $\sh{E}$ của $\sh{B}$ thuộc loại hữu hạn sao cho $\sh{E}$ sinh ra \sref[0]{0.4.1.4} $\sh{O}_X$-đại số $\sh{B}$.
\end{proposition}

\begin{proof}
Thật vậy, theo giả thiết, tồn tại một phủ hữu hạn $(U_i)$ của $X$ gồm các tập mở affine sao cho $\Gamma(U_i,\sh{B})=B_i$ là một đại số thuộc loại hữu hạn trên $\Gamma(U_i,\sh{O}_X)=A_i$; lấy $E_i$ là một $A_i$-môđun con thuộc loại hữu hạn của $B_i$ sinh ra $A_i$-đại số $B_i$; nhờ \sref{I.9.4.7}, tồn tại một $\sh{O}_X$-môđun con $\sh{E}_i$ của $\sh{B}$, quasi-coherent và thuộc loại hữu hạn, sao cho $\sh{E}_i|U_i=\widetilde{E_i}$.
Rõ ràng tổng $\sh{E}$ của các $\sh{E}_i$ là đối tượng cần tìm.
\end{proof}

\begin{proposition}[9.6.6]
\label{I.9.6.6}
Cho $X$ là một tiền lược đồ có không gian nền địa phương noetherian, hoặc một lược đồ quasi-compact.
Khi đó mọi $\sh{O}_X$-đại số quasi-coherent $\sh{B}$ là giới hạn quy nạp của các $\sh{O}_X$-đại số con quasi-coherent thuộc loại hữu hạn của nó.
\end{proposition}

\begin{proof}
Thật vậy, từ \sref{I.9.4.9} suy ra rằng $\sh{B}$ là giới hạn quy nạp (với tư cách $\sh{O}_X$-môđun) của các $\sh{O}_X$-môđun con quasi-coherent thuộc loại hữu hạn của nó; các môđun này sinh ra các \emph{$\sh{O}_X$-đại số con} quasi-coherent thuộc loại hữu hạn của $\sh{B}$ \sref{I.1.3.14}, và do đó $\sh{B}$ \emph{a fortiori} là giới hạn quy nạp của chúng.
\end{proof}

\section{Lược đồ hình thức}
\label{section:I.10}

\subsection{Lược đồ affine hình thức}
\label{subsection:I.10.1}

\begin{env}[10.1.1]
\label{I.10.1.1}
Cho $A$ là một vành tôpô \emph{admissible} \sref[0]{0.7.1.2}; với mỗi ideal định nghĩa $\mathfrak{J}$ của $A$, $\Spec(A/\mathfrak{J})$ có thể được đồng nhất với không gian con đóng $V(\mathfrak{J})$ của $\Spec(A)$ \sref{I.1.1.11}, là tập hợp các iđêan nguyên tố \emph{mở} của $A$; không gian tôpô này không phụ thuộc
\oldpage[I]{181}
vào ideal định nghĩa $\mathfrak{J}$ được xét; ký hiệu không gian tôpô này là $\mathfrak{X}$. Cho $(\mathfrak{J}_\lambda)$ là một hệ cơ bản các lân cận của $0$ trong $A$, gồm các ideal định nghĩa, và với mỗi $\lambda$, cho $\sh{O}_\lambda$ là bó cấu trúc của $\Spec(A/\mathfrak{J}_\lambda)$; bó này được cảm sinh trên $\mathfrak{X}$ bởi $\widetilde{A}/\widetilde{\mathfrak{J}_\lambda}$ (bằng không ngoài $\mathfrak{X}$).
Với $\mathfrak{J}_\mu\subset\mathfrak{J}_\lambda$, đồng cấu chính tắc $A/\mathfrak{J}_\mu\to A/\mathfrak{J}_\lambda$ do đó xác định một đồng cấu $u_{\lambda\mu}:\sh{O}_\mu\to\sh{O}_\lambda$ của các bó vành \sref{I.1.6.1}, và $(\sh{O}_\lambda)$ là một \emph{hệ xạ ảnh các bó vành} đối với các đồng cấu này.
Vì tôpô của $\mathfrak{X}$ có một cơ sở gồm các tập mở quasi-compact, ta có thể gắn với mỗi $\sh{O}_\lambda$ một \emph{bó giả rời rạc các vành tôpô} \sref[0]{0.3.8.1}, có bó vành nền (không có tôpô) là $\sh{O}_\lambda$, và ta vẫn ký hiệu bó đó là $\sh{O}_\lambda$; các $\sh{O}_\lambda$ lại tạo thành một \emph{hệ xạ ảnh các bó vành tôpô} \sref[0]{0.3.8.2}.
Ta ký hiệu $\sh{O}_\mathfrak{X}$ là \emph{bó các vành tôpô} trên $\mathfrak{X}$, giới hạn xạ ảnh của hệ $(\sh{O}_\lambda)$; với mọi tập mở \emph{quasi-compact} $U$ của $\mathfrak{X}$, $\Gamma(U,\sh{O}_\mathfrak{X})$ do đó là vành tôpô, giới hạn xạ ảnh của hệ các vành \emph{rời rạc} $\Gamma(U,\sh{O}_\lambda)$ \sref[0]{0.3.2.6}.
\end{env}

\begin{definition}[10.1.2]
\label{I.10.1.2}
Cho $A$ là một vành tôpô \emph{admissible}, ta gọi phổ hình thức của $A$, ký hiệu là $\Spf(A)$, là không gian con đóng $\mathfrak{X}$ của $\Spec(A)$ gồm các iđêan nguyên tố \emph{mở} của $A$.
Ta nói rằng một không gian vành tôpô là một \emph{lược đồ affine hình thức} nếu nó đẳng cấu với một phổ hình thức $\Spf(A)=\mathfrak{X}$ được trang bị bó các vành tôpô $\sh{O}_\mathfrak{X}$, là giới hạn xạ ảnh của các bó các vành tôpô giả rời rạc $(\widetilde{A}/\widetilde{\mathfrak{J}_\lambda})|\mathfrak{X}$, trong đó $\mathfrak{J}_\lambda$ chạy qua tập lọc các ideal định nghĩa của $A$.
\end{definition}

Khi nói về một \emph{phổ hình thức $\mathfrak{X}=\Spf(A)$} như một lược đồ affine hình thức, thì luôn hiểu đó là không gian vành tôpô $(\mathfrak{X},\sh{O}_\mathfrak{X})$, trong đó $\sh{O}_\mathfrak{X}$ được định nghĩa như trên.

Ta nhận thấy mọi \emph{lược đồ affine} $X=\Spec(A)$ đều có thể được xem như một lược đồ affine hình thức theo đúng một cách, bằng cách xem $A$ là một vành tôpô rời rạc: các vành tôpô $\Gamma(U,\sh{O}_X)$ khi đó là rời rạc bất cứ khi nào $U$ là quasi-compact (nhưng nói chung không phải khi $U$ là một tập mở tùy ý của $X$).

\begin{proposition}[10.1.3]
\label{I.10.1.3}
Nếu $\mathfrak{X}=\Spf(A)$, trong đó $A$ là một vành admissible, thì $\Gamma(\mathfrak{X},\sh{O}_\mathfrak{X})$ đẳng cấu tôpô với $A$.
\end{proposition}

\begin{proof}
Thật vậy, vì $\mathfrak{X}$ đóng trong $\Spec(A)$, nó quasi-compact, và do đó $\Gamma(\mathfrak{X},\sh{O}_\mathfrak{X})$ đẳng cấu tôpô với giới hạn xạ ảnh của các vành rời rạc $\Gamma(\mathfrak{X},\sh{O}_\lambda)$; nhưng $\Gamma(\mathfrak{X},\sh{O}_\lambda)$ đẳng cấu với $A/\mathfrak{J}_\lambda$ \sref{I.1.3.7}; vì $A$ tách và đầy đủ, nó đẳng cấu tôpô với $\varprojlim A/\mathfrak{J}_\lambda$ \sref[0]{0.7.2.1}, từ đó suy ra mệnh đề.
\end{proof}

\begin{proposition}[10.1.4]
\label{I.10.1.4}
Cho $A$ là một vành admissible, $\mathfrak{X}=\Spf(A)$, và với mọi $f\in A$, đặt $\mathfrak{D}(f)=D(f)\cap\mathfrak{X}$; khi đó không gian vành tôpô $(\mathfrak{D}(f),\sh{O}_\mathfrak{X}|\mathfrak{D}(f))$ đẳng cấu với phổ affine hình thức $\Spf(A_{\{f\}})$ \sref[0]{0.7.6.15}.
\end{proposition}

\begin{proof}
Với mọi ideal định nghĩa $\mathfrak{J}$ của $A$, vành rời rạc $S_f^{-1}A/ S_f^{-1}\mathfrak{J}$ được đồng nhất chính tắc với $A_{\{f\}}/\mathfrak{J}_{\{f\}}$ \sref[0]{0.7.6.9}, do đó, theo \sref{I.1.2.5} và \sref{I.1.2.6}, không gian tôpô $\Spf(A_{\{f\}})$ được đồng nhất chính tắc với $\mathfrak{D}(f)$.
Hơn nữa, với mọi tập mở quasi-compact $U$ của $\mathfrak{X}$ nằm trong $\mathfrak{D}(f)$, $\Gamma(U,\sh{O}_\lambda)$ có thể được đồng nhất với môđun các tiết diện của bó cấu trúc của $\Spec(S_f^{-1}A/ S_f^{-1}\mathfrak{J}_\lambda)$ trên $U$ \sref{I.1.3.6}; do đó, đặt $\mathfrak{Y}=\Spf(A_{\{f\}})$, $\Gamma(U,\sh{O}_\mathfrak{X})$ có thể được đồng nhất với môđun các tiết diện $\Gamma(U,\sh{O}_\mathfrak{Y})$, điều này chứng minh mệnh đề.
\end{proof}

\oldpage[I]{182}
\begin{env}[10.1.5]
\label{I.10.1.5}
Với tư cách là bó các vành \emph{không có tôpô}, bó cấu trúc $\sh{O}_\mathfrak{X}$ của $\Spf(A)$, với mỗi $x\in\mathfrak{X}$, có một thớ; theo \sref{I.10.1.4}, thớ đó có thể được đồng nhất với giới hạn quy nạp $\varinjlim_{f\notin\mathfrak{j}_x}A_{\{f\}}$. Do đó, theo \sref[0]{0.7.6.17} và \sref[0]{0.7.6.18}:
\end{env}

\begin{proposition}[10.1.6]
\label{I.10.1.6}
Với mọi $x\in\mathfrak{X}=\Spf(A)$, thớ $\sh{O}_x$ là một vành địa phương có trường dư đẳng cấu với $\kres(x)=A_x/\mathfrak{j}_xA_x$.
Nếu thêm nữa, $A$ là adic và noetherian, thì $\sh{O}_x$ là một vành noetherian.
\end{proposition}

Vì $\kres(x)$ không phải là vành không, từ kết quả này suy ra rằng \emph{giá} của bó các vành $\sh{O}_\mathfrak{X}$ \emph{bằng $\mathfrak{X}$}.

\subsection{Các cấu xạ của lược đồ affine hình thức}
\label{subsection:I.10.2}

\begin{env}[10.2.1]
\label{I.10.2.1}
Cho $A$, $B$ là hai vành admissible, và cho $\vphi:B\to A$ là một đồng cấu \emph{liên tục}.
Ánh xạ liên tục ${}^a\vphi:\Spec(A)\to\Spec(B)$ \sref{I.1.2.1} khi đó ánh xạ $\mathfrak{X}=\Spf(A)$ vào $\mathfrak{Y}=\Spf(B)$, vì ảnh ngược qua $\vphi$ của một ideal nguyên tố mở của $A$ là một ideal nguyên tố mở của $B$.
Mặt khác, với mọi $g\in B$, $\vphi$ xác định một đồng cấu liên tục
\[
  \Gamma(\mathfrak{D}(g),\sh{O}_\mathfrak{Y})
  \longrightarrow
  \Gamma(\mathfrak{D}(\vphi(g)),\sh{O}_\mathfrak{X})
\]
theo \sref{I.10.1.4}, \sref{I.10.1.3} và \sref[0]{0.7.6.7}; vì các đồng cấu này thỏa mãn các điều kiện tương thích đối với những hạn chế tương ứng với việc chuyển từ $g$ sang một bội của $g$, và vì $\mathfrak{D}(\vphi(g))={}^a\vphi^{-1}(\mathfrak{D}(g))$, chúng xác định một đồng cấu \emph{liên tục} của các bó các vành tôpô $\sh{O}_\mathfrak{Y}\to{}^a\vphi_*(\sh{O}_\mathfrak{X})$ \sref[0]{0.3.2.5}, mà ta vẫn ký hiệu là $\widetilde{\vphi}$; do đó ta đã xác định một cấu xạ $\Phi=({}^a\vphi,\widetilde{\vphi})$ của các không gian vành tôpô $\mathfrak{X}\to\mathfrak{Y}$.
Ta lưu ý rằng, với tư cách là một đồng cấu của các bó vành không có tôpô, $\widetilde{\vphi}$ xác định một đồng cấu $\widetilde{\vphi}_x^\sharp:\sh{O}_{{}^a\vphi(x)}\to\sh{O}_x$ trên các thớ, với mọi $x\in\mathfrak{X}$.
\end{env}

\begin{proposition}[10.2.2]
\label{I.10.2.2}
Cho $A$, $B$ là hai vành tôpô admissible, và cho $\mathfrak{X}=\Spf(A)$, $\mathfrak{Y}=\Spf(B)$. Để một cấu xạ $u=(\psi,\theta):\mathfrak{X}\to\mathfrak{Y}$ của các không gian vành tôpô có dạng $({}^a\vphi,\widetilde{\vphi})$, trong đó $\vphi$ là một đồng cấu vành liên tục $B\to A$, điều kiện cần và đủ là, với mọi $x\in\mathfrak{X}$, $\theta_x^\sharp$ là một đồng cấu địa phương $\sh{O}_{\psi(x)}\to\sh{O}_x$.
\end{proposition}

\begin{proof}
Điều kiện là cần: lấy $\mathfrak{p}=\mathfrak{j}_x\in\Spf(A)$ và $\mathfrak{q}=\vphi^{-1}(\mathfrak{j}_x)$; nếu $g\notin\mathfrak{q}$, thì $\vphi(g)\notin\mathfrak{p}$, và hiển nhiên đồng cấu $B_{\{g\}}\to A_{\{\vphi(g)\}}$ do $\vphi$ cảm sinh \sref[0]{0.7.6.7} gửi $\mathfrak{q}_{\{g\}}$ vào một tập con của $\mathfrak{p}_{\{\vphi(g)\}}$; chuyển qua giới hạn quy nạp, có (tính đến \sref{I.10.1.5} và \sref[0]{0.7.6.17}) rằng $\widetilde{\vphi}_x^\sharp$ là một đồng cấu địa phương.

Ngược lại, cho $(\psi,\theta)$ là một cấu xạ thỏa mãn điều kiện của mệnh đề; theo \sref{I.10.1.3}, $\theta$ xác định một đồng cấu vành liên tục
\[
  \vphi=\Gamma(\theta):B=\Gamma(\mathfrak{Y},\sh{O}_\mathfrak{Y})\to\Gamma(\mathfrak{X},\sh{O}_\mathfrak{X})=A.
\]
Theo giả thiết về $\theta$, để tiết diện $\vphi(g)$ của $\sh{O}_\mathfrak{X}$ trên $\mathfrak{X}$ có một mầm khả nghịch tại điểm $x$, điều kiện cần và đủ là $g$ có một mầm khả nghịch tại điểm $\psi(x)$. Nhưng theo \sref[0]{0.7.6.17}, các tiết diện của $\sh{O}_\mathfrak{X}$ (tương ứng $\sh{O}_\mathfrak{Y}$) trên $\mathfrak{X}$ (tương ứng $\mathfrak{Y}$) có mầm không khả nghịch tại điểm $x$ (tương ứng $\psi(x)$) chính là các phần tử của $\mathfrak{j}_x$
\oldpage[I]{183}
(tương ứng $\mathfrak{j}_{\psi(x)}$); nhận xét trên do đó cho thấy ${}^a\vphi=\psi$. Cuối cùng, với mọi $g\in B$, biểu đồ
\[
  \xymatrix{
    B=\Gamma(\mathfrak{Y},\sh{O}_\mathfrak{Y})
      \ar[r]^\vphi\ar[d] &
    \Gamma(\mathfrak{X},\sh{O}_\mathfrak{X})=A\ar[d]\\
    B_{\{g\}}=\Gamma(\mathfrak{D}(g),\sh{O}_\mathfrak{Y})
      \ar[r]_{\Gamma(\theta_{\mathfrak{D}(g)})} &
    \Gamma(\mathfrak{D}(\vphi(g)),\sh{O}_\mathfrak{X})
      =A_{\{\vphi(g)\}}
  }
\]
giao hoán; theo tính chất phổ dụng của các vành phân thức hoàn chỉnh \sref[0]{0.7.6.6}, $\theta_{\mathfrak{D}(g)}$ bằng $\widetilde{\vphi}_{\mathfrak{D}(g)}$ với mọi $g\in B$, và theo \sref[0]{0.3.2.5} ta có $\theta=\widetilde{\vphi}$.
\end{proof}

Ta nói rằng một cấu xạ $(\psi,\theta)$ của các không gian vành tôpô thỏa mãn điều kiện của Mệnh đề~\sref{I.10.2.2} là một \emph{cấu xạ của lược đồ affine hình thức}. Có thể nói rằng các hàm tử $A\mapsto\Spf(A)$ và $\mathfrak{X}\mapsto\Gamma(\mathfrak{X},\sh{O}_\mathfrak{X})$ xác định một \emph{tương đương} giữa phạm trù các vành admissible và phạm trù đối của các lược đồ affine hình thức (T, I, 1.2).

\begin{env}[10.2.3]
\label{I.10.2.3}
Như một trường hợp riêng của \sref{I.10.2.2}, lưu ý rằng, với $f\in A$, phép nhúng chính tắc của lược đồ affine hình thức do $\mathfrak{X}$ cảm sinh trên $\mathfrak{D}(f)$ tương ứng với đồng cấu chính tắc liên tục $A\to A_{\{f\}}$.
Dưới các giả thiết của Mệnh đề~\sref{I.10.2.2}, cho $h$ là một phần tử của $B$ và $g$ là một phần tử của $A$, bội của $\vphi(h)$;
khi đó $\psi(\mathfrak{D}(g))\subset\mathfrak{D}(h)$; hạn chế của $u$ lên $\mathfrak{D}(g)$, được xem như một cấu xạ từ $\mathfrak{D}(g)$ vào $\mathfrak{D}(h)$, là cấu xạ duy nhất $v$ làm cho sơ đồ sau giao hoán
\[
  \xymatrix{
    \mathfrak{D}(g)\ar[r]^v\ar[d] &
    \mathfrak{D}(h)\ar[d]\\
    \mathfrak{X}\ar[r]_u &
    \mathfrak{Y}
  }
\]
Cấu xạ này tương ứng với đồng cấu liên tục duy nhất $\vphi':B_{\{h\}}\to A_{\{g\}}$ \sref[0]{0.7.6.7} làm cho sơ đồ sau giao hoán
\[
  \xymatrix{
    A\ar[d] &
    B\ar[l]_\vphi\ar[d]\\
    A_{\{g\}} &
    B_{\{h\}}\ar[l]_{\vphi'}
  }
\]
\end{env}

\subsection{Các iđêan định nghĩa của một lược đồ affine hình thức}
\label{subsection:I.10.3}

\begin{env}[10.3.1]
\label{I.10.3.1}
Cho $A$ là một vành khả chấp, $\mathfrak{J}$ là một iđêan mở của $A$, và $\mathfrak{X}$ là lược đồ affine hình thức $\Spf(A)$.
Cho $(\mathfrak{J}_\lambda)$ là tập các iđêan định nghĩa của $A$ nằm trong $\mathfrak{J}$; khi đó $\widetilde{\mathfrak{J}}/\widetilde{\mathfrak{J}}_\lambda$ là một bó iđêan của $\widetilde{A}/\widetilde{\mathfrak{J}}_\lambda$.
Ký hiệu bởi $\mathfrak{J}^\Delta$ giới hạn xạ ảnh của các bó được cảm sinh trên $\mathfrak{X}$ bởi $\widetilde{\mathfrak{J}}/\widetilde{\mathfrak{J}}_\lambda$, giới hạn này được đồng nhất với một \emph{bó iđêan} của $\sh{O}_\mathfrak{X}$ \sref[0]{0.3.2.6}.
Với mọi $f\in A$, $\Gamma(\mathfrak{D}(f),\mathfrak{J}^\Delta)$ là giới hạn xạ ảnh của các $S_f^{-1}\mathfrak{J}/S_f^{-1}\mathfrak{J}_\lambda$, nói cách khác, được đồng nhất với iđêan mở $\mathfrak{J}_{\{f\}}$ của vành $A_{\{f\}}$ \sref[0]{0.7.6.9}, và đặc biệt $\Gamma(\mathfrak{X},\mathfrak{J}^\Delta)=\mathfrak{J}$; ta suy ra (vì các $\mathfrak{D}(f)$ lập thành một cơ sở của tôpô trên $\mathfrak{X}$) rằng
\[
\label{I.10.3.1.1}
  \mathfrak{J}^\Delta|\mathfrak{D}(f)=(\mathfrak{J}_{\{f\}})^\Delta.
  \tag{10.3.1.1}
\]
\end{env}

\oldpage[I]{184}
\begin{env}[10.3.2]
\label{I.10.3.2}
Với các ký hiệu của \sref{I.10.3.1}, với mọi $f\in A$, ánh xạ chính tắc từ $A_{\{f\}}=\Gamma(\mathfrak{D}(f),\sh{O}_\mathfrak{X})$ đến $\Gamma(\mathfrak{D}(f),(\widetilde{A}/\widetilde{\mathfrak{J}})|\mathfrak{X})=S_f^{-1}A/S_f^{-1}\mathfrak{J}$ là \emph{toàn ánh} và có hạt nhân $\Gamma(\mathfrak{D}(f),\mathfrak{J}^\Delta)=\mathfrak{J}_{\{f\}}$ \sref[0]{0.7.6.9}; do đó các ánh xạ này xác định một đồng cấu liên tục \emph{toàn ánh}, gọi là \emph{chính tắc}, từ bó các vành tôpô $\sh{O}_\mathfrak{X}$ lên bó các vành rời rạc $(\widetilde{A}/\widetilde{\mathfrak{J}})|\mathfrak{X}$, có hạt nhân là $\mathfrak{J}^\Delta$; đồng cấu này chính là $\widetilde{\vphi}$ \sref{I.10.2.1}, trong đó $\vphi$ là đồng cấu liên tục $A\to A/\mathfrak{J}$; cấu xạ $({}^a\vphi,\widetilde{\vphi}):\Spec(A/\mathfrak{J})\to\mathfrak{X}$ của các lược đồ affine hình thức (trong đó ${}^a\vphi$ là phép đồng phôi đồng nhất từ $\mathfrak{X}$ lên chính nó) cũng được gọi là \emph{chính tắc}.
Do đó, theo những điều trên, ta có một \emph{đẳng cấu chính tắc}
\[
\label{I.10.3.2.1}
  \sh{O}_\mathfrak{X}/\mathfrak{J}^\Delta\isoto(\widetilde{A}/\widetilde{\mathfrak{J}})|\mathfrak{X}.
  \tag{10.3.2.1}
\]

Hiển nhiên (do $\Gamma(\mathfrak{X},\mathfrak{J}^\Delta)=\mathfrak{J}$) ánh xạ $\mathfrak{J}\to\mathfrak{J}^\Delta$ là \emph{tăng nghiêm ngặt}; theo những điều trên, với $\mathfrak{J}\subset\mathfrak{J}'$, bó ${\mathfrak{J}'}^\Delta/\mathfrak{J}^\Delta$ đẳng cấu chính tắc với $\widetilde{\mathfrak{J}'}/\widetilde{\mathfrak{J}}=(\mathfrak{J}'/\mathfrak{J})\supertilde$.
\end{env}

\begin{env}[10.3.3]
\label{I.10.3.3}
Vẫn với các giả thiết và ký hiệu của \sref{I.10.3.1}, ta nói rằng một bó iđêan $\sh{J}$ của $\sh{O}_\mathfrak{X}$ là một \emph{bó iđêan định nghĩa của $\mathfrak{X}$} (hay một \emph{iđêan định nghĩa của $\mathfrak{X}$}) nếu, với mọi $x\in\mathfrak{X}$, tồn tại một lân cận mở của $x$ có dạng $\mathfrak{D}(f)$, trong đó $f\in A$, sao cho $\sh{J}|\mathfrak{D}(f)$ có dạng $\mathfrak{H}^\Delta$, với $\mathfrak{H}$ là một iđêan định nghĩa của $A_{\{f\}}$.
\end{env}

\begin{proposition}[10.3.4]
\label{I.10.3.4}
Với mọi $f\in A$, mỗi iđêan định nghĩa của $\mathfrak{X}$ cảm sinh một iđêan định nghĩa của $\mathfrak{D}(f)$.
\end{proposition}

\begin{proof}
Điều này suy ra từ \sref{I.10.3.1.1}.
\end{proof}

\begin{proposition}[10.3.5]
\label{I.10.3.5}
Nếu $A$ là một vành khả chấp, thì mọi iđêan định nghĩa của $\mathfrak{X}=\Spf(A)$ đều có dạng $\mathfrak{J}^\Delta$, trong đó $\mathfrak{J}$ là một iđêan định nghĩa của $A$, được xác định duy nhất.
\end{proposition}

\begin{proof}
Thật vậy, cho $\sh{J}$ là một iđêan định nghĩa của $\mathfrak{X}$; theo giả thiết, và vì $\mathfrak{X}$ là quasi-compact, tồn tại hữu hạn phần tử $f_i\in A$ sao cho các $\mathfrak{D}(f_i)$ phủ $\mathfrak{X}$ và $\sh{J}|\mathfrak{D}(f_i)=\mathfrak{H}_i^\Delta$, trong đó $\mathfrak{H}_i$ là một iđêan định nghĩa của $A_{\{f_i\}}$.
Với mọi $i$, do đó tồn tại một iđêan mở $\mathfrak{K}_i$ của $A$ sao cho $(\mathfrak{K}_i)_{\{f_i\}}=\mathfrak{H}_i$ \sref[0]{0.7.6.9}; chọn $\mathfrak{K}$ là một iđêan định nghĩa của $A$ nằm trong mọi $\mathfrak{K}_i$.
Ảnh chính tắc của $\sh{J}/\mathfrak{K}^\Delta$ trong bó cấu trúc $(A/\mathfrak{K})\supertilde$ của $\Spec(A/\mathfrak{K})$ \sref{I.10.3.2} do đó có hạn chế trên $\mathfrak{D}(f_i)$ bằng hạn chế của $(\mathfrak{K}_i/\mathfrak{K})\supertilde$;
ta suy ra ảnh chính tắc này là một bó \emph{quasi-coherent} trên $\Spec(A/\mathfrak{K})$, nên có dạng $(\mathfrak{J}/\mathfrak{K})\supertilde$, trong đó $\mathfrak{J}$ là một iđêan của $A$ chứa $\mathfrak{K}$ \sref{I.1.4.1}, từ đó $\sh{J}=\mathfrak{J}^\Delta$ \sref{I.10.3.2};
hơn nữa, vì với mỗi $i$ tồn tại một số nguyên $n_i$ sao cho $\mathfrak{H}_i^{n_i}\subset\mathfrak{K}_{\{f_i\}}$, đặt $n$ là số lớn nhất trong các $n_i$, ta có $(\sh{J}/\mathfrak{K}^\Delta)^n=0$, và do đó theo \sref{I.10.3.2}, $((\mathfrak{J}/\mathfrak{K})\supertilde)^n=0$, cuối cùng suy ra $(\mathfrak{J}/\mathfrak{K})^n=0$ \sref{I.1.3.13}; điều này chứng minh $\mathfrak{J}$ là một iđêan định nghĩa của $A$ \sref[0]{0.7.1.4}.
\end{proof}

\begin{proposition}[10.3.6]
\label{I.10.3.6}
Cho $A$ là một vành adic và $\mathfrak{J}$ là một iđêan định nghĩa của $A$ sao cho $\mathfrak{J}/\mathfrak{J}^2$ là một $(A/\mathfrak{J})$-môđun hữu hạn sinh. Khi đó, với mọi số nguyên $n>0$, ta có $(\mathfrak{J}^\Delta)^n=(\mathfrak{J}^n)^\Delta$.
\end{proposition}

\begin{proof}
Thật vậy, với mọi $f\in A$, ta có (vì $\mathfrak{J}^n$ là một iđêan mở)
\[
  (\Gamma(\mathfrak{D}(f),\mathfrak{J}^\Delta))^n=(\mathfrak{J}_{\{f\}})^n=(\mathfrak{J}^n)_{\{f\}}=\Gamma(\mathfrak{D}(f^n),(\mathfrak{J}^n)^\Delta)
\]
\oldpage[I]{185}
theo \sref{I.10.3.1.1} và \sref[0]{0.7.6.12}.
Vì $(\mathfrak{J}^\Delta)^n$ được liên kết với tiền bó $U\mapsto(\Gamma(U,\mathfrak{J}^\Delta))^n$ \sref[0]{0.4.1.6}, mệnh đề suy ra\footnote{\emph{[Dịch.] Nguồn tiếng Pháp in «le corollaire» tại đây, mặc dù 10.3.6 được phát biểu là một mệnh đề.}} từ việc các $\mathfrak{D}(f)$ lập thành một cơ sở của tôpô trên $\mathfrak{X}$.
\end{proof}

\begin{env}[10.3.7]
\label{I.10.3.7}
Ta nói một họ $(\sh{J}_\lambda)$ các iđêan định nghĩa của $\mathfrak{X}$ là một \emph{hệ cơ bản các iđêan định nghĩa} nếu mỗi iđêan định nghĩa của $\mathfrak{X}$ chứa một trong các $\sh{J}_\lambda$; vì $\sh{J}_\lambda=\mathfrak{J}_\lambda^\Delta$, điều này tương đương với việc các $\mathfrak{J}_\lambda$ lập thành một \emph{hệ cơ bản các lân cận của $0$} trong $A$.
Cho $(f_\alpha)$ là một họ các phần tử của $A$ sao cho các $\mathfrak{D}(f_\alpha)$ phủ $\mathfrak{X}$.
Nếu $(\sh{J}_\lambda)$ là một họ lọc giảm dần các bó iđêan của $\sh{O}_\mathfrak{X}$ sao cho, với mọi $\alpha$, họ $(\sh{J}_\lambda|\mathfrak{D}(f_\alpha))$ là một hệ cơ bản các iđêan định nghĩa của $\mathfrak{D}(f_\alpha)$, thì $(\sh{J}_\lambda)$ là một hệ cơ bản các iđêan định nghĩa của $\mathfrak{X}$.
Thật vậy, với mỗi iđêan định nghĩa $\sh{J}$ của $\mathfrak{X}$, tồn tại một phủ hữu hạn của $\mathfrak{X}$ bởi các $\mathfrak{D}(f_i)$ sao cho, với mọi $i$, $\sh{J}_{\lambda_i}|\mathfrak{D}(f_i)$ là một iđêan định nghĩa của $\mathfrak{D}(f_i)$ nằm trong $\sh{J}|\mathfrak{D}(f_i)$.
Nếu $\mu$ là một chỉ số sao cho $\sh{J}_\mu\subset\sh{J}_{\lambda_i}$ với mọi $i$, thì từ \sref{I.10.3.3} suy ra $\sh{J}_\mu$ là một iđêan định nghĩa của $\mathfrak{X}$, hiển nhiên nằm trong $\sh{J}$; vậy khẳng định được chứng minh.
\end{env}

\subsection{Tiền lược đồ hình thức và các cấu xạ của tiền lược đồ hình thức}
\label{subsection:I.10.4}

\begin{env}[10.4.1]
\label{I.10.4.1}
Cho một không gian vành tôpô $\mathfrak{X}$, ta nói một tập mở $U\subset\mathfrak{X}$ là một \emph{tập mở affine hình thức} (tương ứng, một \emph{tập mở affine hình thức adic}, tương ứng, một \emph{tập mở affine hình thức noetherian}) nếu không gian vành tôpô do $\mathfrak{X}$ cảm sinh trên $U$ là một lược đồ affine hình thức (tương ứng, một lược đồ như thế có vành adic, tương ứng, adic và noetherian).
\end{env}

\begin{definition}[10.4.2]
\label{I.10.4.2}
Ta gọi một không gian vành tôpô $\mathfrak{X}$ là một \emph{tiền lược đồ hình thức} nếu mỗi điểm của nó có một lân cận là tập mở affine hình thức.
Ta nói tiền lược đồ hình thức $\mathfrak{X}$ là adic (tương ứng, noetherian địa phương) nếu mỗi điểm của $\mathfrak{X}$ có một lân cận là tập mở affine hình thức adic (tương ứng, noetherian).
Ta nói $\mathfrak{X}$ là noetherian nếu nó là noetherian địa phương và không gian nền của nó là quasi-compact (và do đó noetherian).
\end{definition}

\begin{proposition}[10.4.3]
\label{I.10.4.3}
Nếu $\mathfrak{X}$ là một tiền lược đồ hình thức (tương ứng, một tiền lược đồ hình thức noetherian địa phương), thì các tập mở affine hình thức (tương ứng, affine noetherian) tạo thành một cơ sở của tôpô trên $\mathfrak{X}$.
\end{proposition}

\begin{proof}
Điều này suy ra từ \sref{I.10.4.2} và \sref{I.10.1.4}, khi lưu ý rằng nếu $A$ là một vành adic noetherian thì $A_{\{f\}}$ cũng vậy với mọi $f\in A$ \sref[0]{0.7.6.11}.
\end{proof}

\begin{corollary}[10.4.4]
\label{I.10.4.4}
Nếu $\mathfrak{X}$ là một tiền lược đồ hình thức (tương ứng, một tiền lược đồ hình thức noetherian địa phương, tương ứng, noetherian), thì không gian vành tôpô cảm sinh trên mỗi tập mở của $\mathfrak{X}$ vẫn là một tiền lược đồ hình thức (tương ứng, một tiền lược đồ hình thức noetherian địa phương, tương ứng, noetherian).
\end{corollary}

\begin{definition}[10.4.5]
\label{I.10.4.5}
Cho hai tiền lược đồ hình thức $\mathfrak{X}$ và $\mathfrak{Y}$, ta gọi một cấu xạ (của các tiền lược đồ hình thức) từ $\mathfrak{X}$ đến $\mathfrak{Y}$ là một cấu xạ $(\psi,\theta)$ của các không gian vành tôpô sao cho, với mọi $x\in\mathfrak{X}$, $\theta_x^\sharp$ là một đồng cấu địa phương $\sh{O}_{\psi(x)}\to\sh{O}_x$.
\end{definition}

Hiển nhiên hợp thành của hai cấu xạ giữa các tiền lược đồ hình thức vẫn là một cấu xạ như thế; do đó các tiền lược đồ hình thức tạo thành một \emph{phạm trù}, và ta ký hiệu $\Hom(\mathfrak{X},\mathfrak{Y})$ cho tập hợp các cấu xạ từ một tiền lược đồ hình thức $\mathfrak{X}$ đến một tiền lược đồ hình thức $\mathfrak{Y}$.

\oldpage[I]{186}
Nếu $U$ là một phần mở của $\mathfrak{X}$, thì phép nhúng chính tắc vào $\mathfrak{X}$ của tiền lược đồ hình thức do $\mathfrak{X}$ cảm sinh trên $U$ là một cấu xạ của các tiền lược đồ hình thức (và thậm chí là một \emph{đơn cấu} của các không gian vành tôpô \sref[0]{0.4.1.1}).

\begin{proposition}[10.4.6]
\label{I.10.4.6}
Cho $\mathfrak{X}$ là một tiền lược đồ hình thức và $\mathfrak{S}=\Spf(A)$ là một lược đồ affine hình thức.
Có một tương ứng song ánh chính tắc giữa các cấu xạ từ tiền lược đồ hình thức $\mathfrak{X}$ đến tiền lược đồ hình thức $\mathfrak{S}$ và các đồng cấu liên tục từ vành $A$ vào vành tôpô $\Gamma(\mathfrak{X},\sh{O}_\mathfrak{X})$.
\end{proposition}

\begin{proof}
Chứng minh giống như chứng minh của \sref{I.2.2.4}, bằng cách thay ``đồng cấu'' bởi ``đồng cấu liên tục'', ``tập mở affine'' bởi ``tập mở affine hình thức'', và dùng \sref{I.10.2.2} thay cho \sref{I.1.7.3}; ta để các chi tiết cho người đọc.
\end{proof}

\begin{env}[10.4.7]
\label{I.10.4.7}
Cho một tiền lược đồ hình thức $\mathfrak{S}$, một tiền lược đồ hình thức $\mathfrak{X}$ và một cấu xạ $\vphi:\mathfrak{X}\to\mathfrak{S}$, ta nói các dữ kiện này xác định một tiền lược đồ hình thức \emph{$\mathfrak{X}$ trên $\mathfrak{S}$}, hay một \emph{$\mathfrak{S}$-tiền lược đồ hình thức}; $\vphi$ được gọi là \emph{cấu xạ cấu trúc} của $\mathfrak{S}$-tiền lược đồ $\mathfrak{X}$.
Nếu $\mathfrak{S}=\Spf(A)$, trong đó $A$ là một vành khả chấp, thì ta cũng nói $\mathfrak{S}$-tiền lược đồ hình thức $\mathfrak{X}$ là một \emph{$A$-tiền lược đồ hình thức}, hay một tiền lược đồ hình thức \emph{trên $A$}.
Mọi tiền lược đồ hình thức đều có thể được xem là một tiền lược đồ hình thức trên $\bb{Z}$ (được trang bị tôpô rời rạc).

Nếu $\mathfrak{X}$ và $\mathfrak{Y}$ là hai $\mathfrak{S}$-tiền lược đồ hình thức, ta nói một cấu xạ $u:\mathfrak{X}\to\mathfrak{Y}$ là một \emph{$\mathfrak{S}$-cấu xạ} nếu sơ đồ
\[
  \xymatrix{
    \mathfrak{X}\ar[rr]^u\ar[rd] & &
    \mathfrak{Y}\ar[ld]\\
    & \mathfrak{S}
  }
\]
(trong đó các mũi tên xiên là các cấu xạ cấu trúc) giao hoán.
Với định nghĩa này, các $\mathfrak{S}$-tiền lược đồ hình thức (với $\mathfrak{S}$ cố định) tạo thành một \emph{phạm trù}.
Ta ký hiệu $\Hom_\mathfrak{S}(\mathfrak{X},\mathfrak{Y})$ cho tập hợp các $\mathfrak{S}$-cấu xạ từ $\mathfrak{S}$-tiền lược đồ hình thức $\mathfrak{X}$ đến $\mathfrak{S}$-tiền lược đồ hình thức $\mathfrak{Y}$.
Khi $\mathfrak{S}=\Spf(A)$, ta cũng nói \emph{$A$-cấu xạ} thay cho \emph{$\mathfrak{S}$-cấu xạ}.
\end{env}

\begin{env}[10.4.8]
\label{I.10.4.8}
Vì mọi lược đồ affine đều có thể được xem là một lược đồ affine hình thức \sref{I.10.1.2}, nên mọi tiền lược đồ (thông thường) đều có thể được xem là một tiền lược đồ hình thức.
Hơn nữa, từ \sref{I.10.4.5} suy ra rằng, đối với các tiền lược đồ \emph{thông thường}, các cấu xạ (tương ứng, $S$-cấu xạ) của tiền lược đồ \emph{hình thức} trùng với các cấu xạ (tương ứng, $S$-cấu xạ) được định nghĩa trong \textsection\hyperref[section:I.2]{2}.
\end{env}

\subsection{Các iđêan định nghĩa của tiền lược đồ hình thức}
\label{subsection:I.10.5}

\begin{env}[10.5.1]
\label{I.10.5.1}
Cho $\mathfrak{X}$ là một tiền lược đồ hình thức; ta nói một $\sh{O}_\mathfrak{X}$-iđêan $\sh{J}$ là một \emph{bó iđêan định nghĩa} (hay một \emph{iđêan định nghĩa}) của $\mathfrak{X}$ nếu mỗi $x\in\mathfrak{X}$ có một lân cận mở affine hình thức $U$ sao cho $\sh{J}|U$ là một iđêan định nghĩa của lược đồ affine hình thức do $\mathfrak{X}$ cảm sinh trên $U$ \sref{I.10.3.3}; theo \sref{I.10.3.1.1} và \sref{I.10.4.3}, với mọi tập mở $V\subset\mathfrak{X}$, $\sh{J}|V$ khi đó là một iđêan định nghĩa của tiền lược đồ hình thức do $\mathfrak{X}$ cảm sinh trên $V$.

Ta nói một họ $(\sh{J}_\lambda)$ các iđêan định nghĩa của $\mathfrak{X}$ là một \emph{hệ cơ bản}
\oldpage[I]{187}
\emph{các iđêan định nghĩa} nếu tồn tại một phủ $(U_\alpha)$ của $\mathfrak{X}$ bởi các tập mở affine hình thức sao cho, với mọi $\alpha$, họ các $\sh{J}_\lambda|U_\alpha$ là một hệ cơ bản các iđêan định nghĩa \sref{I.10.3.7}\footnote{\emph{[Dịch.] Nguồn tiếng Pháp in dẫn ở đây đến 10.3.6; định nghĩa được viện dẫn là 10.3.7.}} của lược đồ affine hình thức do $\mathfrak{X}$ cảm sinh trên $U_\alpha$.
Từ nhận xét cuối của \sref{I.10.3.7} suy ra rằng, khi $\mathfrak{X}$ là một lược đồ affine hình thức, định nghĩa này trùng với định nghĩa được nêu trong \sref{I.10.3.7}.
Với mọi tập mở $V$ của $\mathfrak{X}$, các hạn chế $\sh{J}_\lambda|V$ khi đó tạo thành một hệ cơ bản các iđêan định nghĩa của tiền lược đồ hình thức cảm sinh trên $V$, theo \sref{I.10.3.1.1}.
Nếu $\mathfrak{X}$ là một tiền lược đồ hình thức \emph{noetherian địa phương} và $\sh{J}$ là một iđêan định nghĩa của $\mathfrak{X}$, thì từ \sref{I.10.3.6} suy ra các lũy thừa $\sh{J}^n$ tạo thành một hệ cơ bản các iđêan định nghĩa của $\mathfrak{X}$.
\end{env}

\begin{env}[10.5.2]
\label{I.10.5.2}
Cho $\mathfrak{X}$ là một tiền lược đồ hình thức và $\sh{J}$ là một iđêan định nghĩa của $\mathfrak{X}$.
Khi đó không gian vành $(\mathfrak{X},\sh{O}_\mathfrak{X}/\sh{J})$ là một \emph{tiền lược đồ} (thông thường), affine (tương ứng, noetherian địa phương, tương ứng, noetherian) khi $\mathfrak{X}$ là một lược đồ affine hình thức (tương ứng, một tiền lược đồ hình thức noetherian địa phương, tương ứng, noetherian); thật vậy, ta quy ngay về trường hợp affine, và khi ấy mệnh đề đã được chứng minh trong \sref{I.10.3.2}.
Hơn nữa, nếu $\theta:\sh{O}_\mathfrak{X}\to\sh{O}_\mathfrak{X}/\sh{J}$ là đồng cấu chính tắc, thì $u=(1_\mathfrak{X},\theta)$ là một \emph{cấu xạ} (gọi là \emph{chính tắc}) của các tiền lược đồ hình thức $(\mathfrak{X},\sh{O}_\mathfrak{X}/\sh{J})\to(\mathfrak{X},\sh{O}_\mathfrak{X})$, vì, cũng như trên, điều này đã được thấy trong trường hợp affine \sref{I.10.3.2}, mà ta quy ngay về.
\end{env}

\begin{proposition}[10.5.3]
\label{I.10.5.3}
Cho $\mathfrak{X}$ là một tiền lược đồ hình thức và $(\sh{J}_\lambda)$ là một hệ cơ bản các iđêan định nghĩa của $\mathfrak{X}$.
Khi đó bó các vành tôpô $\sh{O}_\mathfrak{X}$ là giới hạn xạ ảnh của các bó vành giả rời rạc \sref[0]{0.3.8.1} $\sh{O}_\mathfrak{X}/\sh{J}_\lambda$.
\end{proposition}

\begin{proof}
Vì tôpô của $\mathfrak{X}$ có một cơ sở gồm các tập mở affine hình thức quasi-compact \sref{I.10.4.3}, ta quy về trường hợp affine, trong đó mệnh đề là hệ quả của Mệnh đề \sref{I.10.3.5}, của \sref{I.10.3.2} và của Định nghĩa \sref{I.10.1.1}.
\end{proof}

Không chắc rằng mọi tiền lược đồ hình thức đều có iđêan định nghĩa.
Tuy nhiên:
\begin{proposition}[10.5.4]
\label{I.10.5.4}
Cho $\mathfrak{X}$ là một tiền lược đồ hình thức noetherian địa phương.
Tồn tại một iđêan định nghĩa lớn nhất $\sh{T}$ của $\mathfrak{X}$; đó là iđêan định nghĩa duy nhất $\sh{J}$ sao cho tiền lược đồ $(\mathfrak{X},\sh{O}_\mathfrak{X}/\sh{J})$ là rút gọn.
Nếu $\sh{J}$ là một iđêan định nghĩa của $\mathfrak{X}$, thì $\sh{T}$ là nghịch ảnh, qua $\sh{O}_\mathfrak{X}\to\sh{O}_\mathfrak{X}/\sh{J}$, của căn không của $\sh{O}_\mathfrak{X}/\sh{J}$.
\end{proposition}

\begin{proof}
Trước hết giả sử $\mathfrak{X}=\Spf(A)$, trong đó $A$ là một vành adic noetherian.
Sự tồn tại và các tính chất của $\sh{T}$ suy ra ngay từ \sref{I.10.3.5}\footnote{\emph{[Dịch.] Nguồn tiếng Pháp in dẫn đến 10.3.4; bản đính chính đã công bố trong EGA~II dẫn đến 10.3.5.}} và \sref{I.5.1.1}, khi tính đến sự tồn tại và các tính chất của iđêan định nghĩa lớn nhất của $A$ (\sref[0]{0.7.1.6} và \sref[0]{0.7.1.7}).

Để chứng minh sự tồn tại và các tính chất của $\sh{T}$ trong trường hợp tổng quát, chỉ cần chỉ ra rằng nếu $U\supset V$ là hai tập mở affine hình thức noetherian của $\mathfrak{X}$, thì iđêan định nghĩa lớn nhất $\sh{T}_U$ của $U$ cảm sinh iđêan định nghĩa lớn nhất $\sh{T}_V$ của $V$; nhưng vì $(V,(\sh{O}_\mathfrak{X}|V)/(\sh{T}_U|V))$ là rút gọn, điều này suy ra từ phần trên.
\end{proof}

Ta ký hiệu $\mathfrak{X}_\text{red}$ là tiền lược đồ (thông thường) rút gọn $(\mathfrak{X},\sh{O}_\mathfrak{X}/\sh{T}$).

\begin{corollary}[10.5.5]
\label{I.10.5.5}
Cho $\mathfrak{X}$ là một tiền lược đồ hình thức noetherian địa phương và $\sh{T}$ là iđêan định nghĩa lớn nhất của $\mathfrak{X}$; với mọi tập mở $V$ của $\mathfrak{X}$, $\sh{T}|V$ là iđêan định nghĩa lớn nhất của tiền lược đồ hình thức do $\mathfrak{X}$ cảm sinh trên $V$.
\end{corollary}

\oldpage[I]{188}

\begin{proposition}[10.5.6]
\label{I.10.5.6}
Cho $\mathfrak{X}$ và $\mathfrak{Y}$ là hai tiền lược đồ hình thức, $\sh{J}$ (tương ứng, $\sh{K}$) là một iđêan định nghĩa của $\mathfrak{X}$ (tương ứng, $\mathfrak{Y}$), và $f:\mathfrak{X}\to\mathfrak{Y}$ là một cấu xạ của các tiền lược đồ hình thức.
\begin{enumerate}
  \item[\rm{(i)}] Nếu $f^*(\sh{K})\sh{O}_\mathfrak{X}\subset\sh{J}$, thì tồn tại duy nhất một cấu xạ $f':(\mathfrak{X},\sh{O}_\mathfrak{X}/\sh{J})\to(\mathfrak{Y},\sh{O}_\mathfrak{Y}/\sh{K})$ của các tiền lược đồ thông thường làm cho sơ đồ
    \[
    \label{I.10.5.6.1}
      \xymatrix{
        (\mathfrak{X},\sh{O}_\mathfrak{X})\ar[r]^f &
        (\mathfrak{Y},\sh{O}_\mathfrak{Y})\\
        (\mathfrak{X},\sh{O}_\mathfrak{X}/\sh{J})\ar[r]_{f'}\ar[u] &
        (\mathfrak{Y},\sh{O}_\mathfrak{Y}/\sh{K})\ar[u]
      }
      \tag{10.5.6.1}
    \]
    giao hoán, trong đó các mũi tên dọc là các cấu xạ chính tắc.
  \item[\rm{(ii)}] Giả sử $\mathfrak{X}=\Spf(A)$ và $\mathfrak{Y}=\Spf(B)$ là các lược đồ affine hình thức, $\sh{J}=\mathfrak{J}^\Delta$ và $\sh{K}=\mathfrak{K}^\Delta$, trong đó $\mathfrak{J}$ (tương ứng, $\mathfrak{K}$) là một iđêan định nghĩa của $A$ (tương ứng, $B$), và $f=({}^a\vphi,\widetilde{\vphi})$, trong đó $\vphi:B\to A$ là một đồng cấu liên tục;
    để có $f^*(\sh{K})\sh{O}_\mathfrak{X}\subset\sh{J}$, điều kiện cần và đủ là $\vphi(\mathfrak{K})\subset\mathfrak{J}$, và khi đó $f'$ là cấu xạ $({}^a\vphi',\widetilde{\vphi'})$, trong đó $\vphi':B/\mathfrak{K}\to A/\mathfrak{J}$ là đồng cấu suy ra từ $\vphi$ bằng cách chuyển sang các thương.
\end{enumerate}
\end{proposition}

\begin{proof}
\medskip\noindent
\begin{enumerate}
  \item[\rm{(i)}] Nếu $f=(\psi,\theta)$, thì giả thiết kéo theo rằng ảnh qua $\theta^\sharp:\psi^*(\sh{O}_\mathfrak{Y})\to\sh{O}_\mathfrak{X}$ của bó iđêan $\psi^*(\sh{K})$ trong $\psi^*(\sh{O}_\mathfrak{Y})$ nằm trong $\sh{J}$ \sref[0]{0.4.3.5}.
    Do đó, bằng cách chuyển sang các thương, ta suy ra từ $\theta^\sharp$ một đồng cấu các bó vành
    \[
      \omega:\psi^*(\sh{O}_\mathfrak{Y}/\sh{K})=\psi^*(\sh{O}_\mathfrak{Y})/\psi^*(\sh{K})\to\sh{O}_\mathfrak{X}/\sh{J};
    \]
    hơn nữa, vì với mọi $x\in\mathfrak{X}$, $\theta_x^\sharp$ là một đồng cấu \emph{địa phương}, nên $\omega_x$ cũng vậy.
    Do đó cấu xạ các không gian vành $(\psi,\omega^\flat)$ là \sref{I.2.2.1} cấu xạ duy nhất $f'$ của các không gian vành đáp ứng yêu cầu.
  \item[\rm{(ii)}] Tương ứng chính tắc có tính hàm tử giữa các cấu xạ của lược đồ affine hình thức và các đồng cấu liên tục của vành \sref{I.10.2.2} cho thấy rằng, trong trường hợp đang xét, quan hệ $f^*(\sh{K})\sh{O}_\mathfrak{X}\subset\sh{J}$ kéo theo $f'=({}^a\vphi',\widetilde{\vphi'})$, trong đó $\vphi':B/\mathfrak{K}\to A/\mathfrak{J}$ là đồng cấu duy nhất làm cho sơ đồ
    \[
    \label{I.10.5.6.2}
      \xymatrix{
        B\ar[r]^\vphi\ar[d] &
        A\ar[d]\\
        B/\mathfrak{K}\ar[r]_{\vphi'} &
        A/\mathfrak{J}
      }
      \tag{10.5.6.2}
    \]
    giao hoán.
    Do đó sự tồn tại của $\vphi'$ kéo theo $\vphi(\mathfrak{K})\subset\mathfrak{J}$.
    Ngược lại, nếu điều kiện này được thỏa mãn, ký hiệu $\vphi'$ là đồng cấu duy nhất làm cho sơ đồ \sref{I.10.5.6.2} giao hoán và đặt $f'=({}^a\vphi',\widetilde{\vphi'})$, thì hiển nhiên sơ đồ \sref{I.10.5.6.1} giao hoán; khi đó việc xét các đồng cấu ${}^a\vphi^*(\sh{O}_\mathfrak{Y})\to\sh{O}_\mathfrak{X}$ và ${}^a{\vphi'}^*(\sh{O}_\mathfrak{Y}/\sh{K})\to\sh{O}_\mathfrak{X}/\sh{J}$ tương ứng với $f$ và $f'$ cho thấy điều đó kéo theo quan hệ $f^*(\sh{K})\sh{O}_\mathfrak{X}\subset\sh{J}$.
\end{enumerate}
\end{proof}

Hiển nhiên tương ứng $f\mapsto f'$ được xác định ở trên có tính \emph{hàm tử}.

\subsection{Tiền lược đồ hình thức như giới hạn quy nạp của các tiền lược đồ}
\label{subsection:I.10.6}

\begin{env}[10.6.1]
\label{I.10.6.1}
Cho $\mathfrak{X}$ là một tiền lược đồ hình thức và $(\sh{J}_\lambda)$ là một hệ cơ bản các iđêan định nghĩa của $\mathfrak{X}$; với mỗi $\lambda$, gọi $f_\lambda$ là cấu xạ chính tắc $(\mathfrak{X},\sh{O}_\mathfrak{X}/\sh{J}_\lambda)\to\mathfrak{X}$ \sref{I.10.5.2}; khi $\sh{J}_\mu\subset\sh{J}_\lambda$, đồng cấu chính tắc $\sh{O}_\mathfrak{X}/\sh{J}_\mu\to\sh{O}_\mathfrak{X}/\sh{J}_\lambda$ xác định một cấu xạ chính tắc
\oldpage[I]{189}
$f_{\mu\lambda}:(\mathfrak{X},\sh{O}_\mathfrak{X}/\sh{J}_\lambda)\to(\mathfrak{X},\sh{O}_\mathfrak{X}/\sh{J}_\mu)$ của các tiền lược đồ (thông thường) sao cho $f_\lambda=f_\mu\circ f_{\mu\lambda}$.
Do đó các tiền lược đồ $X_\lambda=(\mathfrak{X},\sh{O}_\mathfrak{X}/\sh{J}_\lambda)$ và các cấu xạ $f_{\mu\lambda}$ tạo thành (theo \sref{I.10.4.8}) một \emph{hệ quy nạp} trong phạm trù các tiền lược đồ hình thức.
\end{env}

\begin{proposition}[10.6.2]
\label{I.10.6.2}
Với các ký hiệu của \sref{I.10.6.1}, tiền lược đồ hình thức $\mathfrak{X}$ và các cấu xạ $f_\lambda$ tạo thành một giới hạn quy nạp (T,~I,~1.8) của hệ $(X_\lambda,f_{\mu\lambda})$ trong phạm trù các tiền lược đồ hình thức.
\end{proposition}

\begin{proof}
Cho $\mathfrak{Y}$ là một tiền lược đồ hình thức và, với mỗi chỉ số $\lambda$, cho
\[
  g_\lambda=(\psi_\lambda,\theta_\lambda):X_\lambda\to\mathfrak{Y}
\]
là một cấu xạ sao cho $g_\lambda=g_\mu\circ f_{\mu\lambda}$ khi $\sh{J}_\mu\subset\sh{J}_\lambda$.
Điều kiện sau cùng này và định nghĩa các $X_\lambda$ trước hết kéo theo rằng tất cả các $\psi_\lambda$ đều đồng nhất với một ánh xạ liên tục duy nhất $\psi:\mathfrak{X}\to\mathfrak{Y}$ của các không gian nền; hơn nữa, các đồng cấu $\theta_\lambda^\sharp:\psi^*(\sh{O}_\mathfrak{Y})\to\sh{O}_{X_\lambda}=\sh{O}_\mathfrak{X}/\sh{J}_\lambda$ tạo thành một \emph{hệ xạ ảnh} các đồng cấu bó vành.
Chuyển sang giới hạn xạ ảnh, ta suy ra một đồng cấu $\omega:\psi^*(\sh{O}_\mathfrak{Y})\to\varprojlim\sh{O}_\mathfrak{X}/\sh{J}_\lambda=\sh{O}_\mathfrak{X}$, và hiển nhiên cấu xạ $g=(\psi,\omega^\flat)$ của các \emph{không gian vành} là cấu xạ \emph{duy nhất} làm cho các sơ đồ
\[
\label{I.10.6.2.1}
  \xymatrix{
    X_\lambda\ar[rr]^{g_\lambda}\ar[rd]_{f_\lambda} & &
    \mathfrak{Y}\\
    & \mathfrak{X}\ar[ru]_g
  }
  \tag{10.6.2.1}
\]
giao hoán.
Còn phải chứng minh rằng $g$ là một cấu xạ của các \emph{tiền lược đồ hình thức}; vì vấn đề là địa phương trên $\mathfrak{X}$ và $\mathfrak{Y}$, ta có thể giả sử $\mathfrak{X}=\Spf(A)$ và $\mathfrak{Y}=\Spf(B)$, trong đó $A$ và $B$ là các vành admissible, và $\sh{J}_\lambda=\mathfrak{J}_\lambda^\Delta$, trong đó $(\mathfrak{J}_\lambda)$ là một hệ cơ bản các iđêan định nghĩa của $A$ \sref{I.10.3.5}; vì $A=\varprojlim A/\mathfrak{J}_\lambda$, sự tồn tại của một cấu xạ $g$ của các lược đồ affine hình thức làm cho các sơ đồ \sref{I.10.6.2.1} giao hoán khi đó suy ra từ tương ứng song ánh \sref{I.10.2.2} giữa các cấu xạ của lược đồ affine hình thức và các đồng cấu vành liên tục, cùng với định nghĩa của giới hạn xạ ảnh.
Nhưng tính duy nhất của $g$ như một cấu xạ của các không gian vành cho thấy nó trùng với cấu xạ được ký hiệu cùng như vậy ở đầu chứng minh.
\end{proof}

Mệnh đề sau đây, dưới một số điều kiện bổ sung, thiết lập sự tồn tại của giới hạn quy nạp của một hệ quy nạp tiền lược đồ (thông thường) đã cho trong phạm trù các tiền lược đồ hình thức:
\begin{proposition}[10.6.3]
\label{I.10.6.3}
Cho $\mathfrak{X}$ là một không gian tôpô và $(\sh{O}_i,u_{ji})$ là một hệ xạ ảnh các bó vành trên $\mathfrak{X}$, với tập chỉ số là $\bb{N}$.
Đặt $\sh{J}_i$ là hạt nhân của $u_{0i}:\sh{O}_i\to\sh{O}_0$.
Giả sử rằng:
\begin{enumerate}
  \item[\rm{(a)}] không gian vành $(\mathfrak{X},\sh{O}_i)$ là một tiền lược đồ $X_i$;
  \item[\rm{(b)}] với mọi $x\in\mathfrak{X}$\footnote{\emph{[Dịch.] Nguồn tiếng Pháp in ghi $x\in X$; bản đính chính đã công bố trong EGA~II thay $X$ bằng $\mathfrak{X}$.}} và mọi $i$, tồn tại một lân cận mở $U_i$ của $x$ trong $\mathfrak{X}$ sao cho hạn chế $\sh{J}_i|U_i$ là lũy linh; và
  \item[\rm{(c)}] các đồng cấu $u_{ji}$ là toàn cấu.
\end{enumerate}

\oldpage[I]{190}
Cho $\sh{O}_\mathfrak{X}$ là bó các vành tôpô được cho bởi giới hạn xạ ảnh của các bó vành giả rời rạc $\sh{O}_i$, và cho $u_i:\sh{O}_\mathfrak{X}\to\sh{O}_i$ là đồng cấu chính tắc.
Khi đó không gian vành tôpô $(\mathfrak{X},\sh{O}_\mathfrak{X})$ là một tiền lược đồ hình thức; các đồng cấu $u_i$ là toàn cấu; các hạt nhân $\sh{J}^{(i)}$ của chúng tạo thành một hệ cơ bản các bó iđêan định nghĩa của $\mathfrak{X}$, và $\sh{J}^{(0)}$ là giới hạn xạ ảnh của các bó iđêan $\sh{J}_i$.
\end{proposition}

\begin{proof}
Trước hết lưu ý rằng, trên mỗi thớ, $u_{ji}$ là một đồng cấu toàn cấu và \emph{a fortiori} là một đồng cấu địa phương; do đó $v_{ij}=(1_\mathfrak{X},u_{ji})$ là một cấu xạ tiền lược đồ $X_j\to X_i$ ($i\geq j$) \sref{I.2.2.1}.
Trước hết giả sử mỗi $X_i$ là một lược đồ affine của vành $A_i$.
Tồn tại một đồng cấu \emph{vành} $\vphi_{ji}:A_i\to A_j$ sao cho $u_{ji}=\widetilde{\vphi_{ji}}$ \sref{I.1.7.3}; do đó \sref{I.1.6.3}, bó $\sh{O}_j$ là một $\sh{O}_i$-môđun quasi-coherent trên $X_i$ (đối với phép nhân vô hướng được xác định bởi $u_{ji}$), liên kết với $A_j$ được xem là một $A_i$-môđun nhờ $\vphi_{ji}$.
Với mọi $f\in A_i$, đặt $f'=\vphi_{ji}(f)$; theo giả thiết, các tập mở $D(f)$ và $D(f')$ trùng nhau trong $\mathfrak{X}$, và đồng cấu từ $\Gamma(D(f),\sh{O}_i)=(A_i)_f$ vào $\Gamma(D(f),\sh{O}_j)=(A_j)_{f'}$ tương ứng với $u_{ji}$ chính là $(\vphi_{ji})_f$ \sref{I.1.6.1}.
Nhưng khi xem $A_j$ là một $A_i$-môđun, $(A_j)_{f'}$ là $(A_i)_f$-môđun $(A_j)_f$, nên ta cũng có $u_{ji}=\widetilde{\vphi_{ji}}$, trong đó lần này $\vphi_{ji}$ được xem là đồng cấu các \emph{$A_i$-môđun}.
Khi đó, vì $u_{ji}$ là toàn cấu, suy ra $\vphi_{ji}$ cũng là toàn cấu \sref{I.1.3.9}; và nếu $\mathfrak{J}_{ji}$ là hạt nhân của $\vphi_{ji}$, thì hạt nhân của $u_{ji}$ là một $\sh{O}_i$-môđun quasi-coherent bằng $\widetilde{\mathfrak{J}_{ji}}$.
Đặc biệt, ta có $\sh{J}_i=\widetilde{\mathfrak{J}_i}$, trong đó $\mathfrak{J}_i$ là hạt nhân của $\vphi_{0i}:A_i\to A_0$.
Giả thiết~(b) kéo theo $\sh{J}_i$ là \emph{lũy linh}: thật vậy, vì $\mathfrak{X}$ là quasi-compact, ta có thể phủ $\mathfrak{X}$ bởi hữu hạn tập mở $U_k$ sao cho $(\sh{J}_i|U_k)^{n_k}=0$ và, lấy $n$ là số lớn nhất trong các $n_k$, ta có $\sh{J}_i^n=0$.
Do đó suy ra $\mathfrak{J}_i$ là lũy linh \sref{I.1.3.13}.
Khi đó vành $A=\varprojlim A_i$ là khả chấp \sref[0]{0.7.2.2}, đồng cấu chính tắc $\vphi_i:A\to A_i$ là toàn cấu và hạt nhân $\mathfrak{J}^{(i)}$ của nó bằng giới hạn xạ ảnh của các $\mathfrak{J}_{ik}$ với $k\geq i$; các $\mathfrak{J}^{(i)}$ tạo thành một hệ cơ bản các lân cận của $0$ trong $A$.
Các khẳng định của Mệnh đề~\sref{I.10.6.3} trong trường hợp này suy ra từ \sref{I.10.1.1} và \sref{I.10.3.2}, vì $(\mathfrak{X},\sh{O}_\mathfrak{X})$ chính là $\Spf(A)$.

Vẫn trong trường hợp riêng này, lưu ý rằng nếu $f=(f_i)$ là một phần tử của giới hạn xạ ảnh $A=\varprojlim A_i$, thì tất cả các tập mở $D(f_i)$ (là các tập mở affine trong $X_i$) đều đồng nhất được với tập mở $\mathfrak{D}(f)$ của $\mathfrak{X}$, và do đó tiền lược đồ do $X_i$ cảm sinh trên $\mathfrak{D}(f)$ đồng nhất được với lược đồ affine $\Spec((A_i)_{f_i})$.

Trong trường hợp tổng quát, trước hết nhận xét rằng, với mọi tập mở quasi-compact $U$ của $\mathfrak{X}$, mỗi $\sh{J}_i|U$ đều lũy linh, như lập luận trên đã chỉ ra.
Ta sẽ chỉ ra rằng, với mọi $x\in\mathfrak{X}$, tồn tại một lân cận mở $U$ của $x$ trong $\mathfrak{X}$ là một \emph{tập mở affine} đối với \emph{mọi} $X_i$.
Thật vậy, lấy $U$ là một tập mở affine đối với $X_0$ và nhận xét rằng $\sh{O}_{X_0}=\sh{O}_{X_i}/\sh{J}_i$.
Vì, theo phần trên, $\sh{J}_i|U$ là lũy linh, nên $U$ cũng là tập mở affine đối với mỗi $X_i$, theo Mệnh đề~\sref{I.5.1.9}.
Do đó, với mọi $U$ thỏa mãn các điều kiện trên, việc xét trường hợp affine ở trên cho thấy $(U,\sh{O}_\mathfrak{X}|U)$ là một tiền lược đồ hình thức mà các $\sh{J}^{(i)}|U$ tạo thành một hệ cơ bản các bó iđêan định nghĩa, và $\sh{J}^{(0)}|U$ là giới hạn xạ ảnh của các $\sh{J}_i|U$; từ đó có kết luận.
\end{proof}

\begin{corollary}[10.6.4]
\label{I.10.6.4}
Giả sử rằng, với $i\geq j$, hạt nhân của $u_{ji}$ là $\sh{J}_i^{j+1}$ và $\sh{J}_1/\sh{J}_1^2$
\oldpage[I]{191}
là một $\sh{O}_0=\sh{O}_1/\sh{J}_1$-môđun hữu hạn sinh.
Khi đó $\mathfrak{X}$ là một tiền lược đồ hình thức adic và, nếu $\sh{J}^{(n)}$ là hạt nhân của $\sh{O}_\mathfrak{X}\to\sh{O}_n$, thì $\sh{J}^{(n)}=\sh{J}^{n+1}$ và $\sh{J}/\sh{J}^2$ đẳng cấu với $\sh{J}_1$.
Nếu thêm nữa $X_0$ là noetherian địa phương (tương ứng, noetherian), thì $\mathfrak{X}$ là noetherian địa phương (tương ứng, noetherian).
\end{corollary}

\begin{proof}
Vì các không gian nền của $\mathfrak{X}$ và $X_0$ là như nhau, vấn đề có tính địa phương và ta có thể giả sử mọi $X_i$ đều affine; xét đến các hệ thức $\sh{J}_{ij}=\widetilde{\mathfrak{J}_{ji}}$ (với ký hiệu của Mệnh đề~\sref{I.10.6.3}), bài toán được quy ngay về các khẳng định tương ứng của Mệnh đề~\sref[0]{0.7.2.7} và Hệ quả~\sref[0]{0.7.2.8}, khi lưu ý rằng $\mathfrak{J}_1/\mathfrak{J}_1^2$ khi đó là một $A_0$-môđun hữu hạn sinh \sref{I.1.3.9}.
\end{proof}

Đặc biệt, \emph{mọi tiền lược đồ hình thức noetherian địa phương $\mathfrak{X}$} là giới hạn quy nạp của một dãy $(X_n)$ các tiền lược đồ (thông thường) noetherian địa phương thỏa mãn các điều kiện của Mệnh đề~\sref{I.10.6.3} và Hệ quả~\sref{I.10.6.4}: chỉ cần xét một bó iđêan định nghĩa $\sh{J}$ của $\mathfrak{X}$ \sref{I.10.5.4} và lấy $X_n=(\mathfrak{X},\sh{O}_\mathfrak{X}/\sh{J}^{n+1})$ (\sref{I.10.5.1} và Mệnh đề~\sref{I.10.6.2}).

\begin{corollary}[10.6.5]
\label{I.10.6.5}
Cho $A$ là một vành khả chấp.
Điều kiện cần và đủ để lược đồ affine hình thức $\mathfrak{X}=\Spf(A)$ là noetherian là $A$ adic và noetherian.
\end{corollary}

\begin{proof}
Điều kiện hiển nhiên là đủ.
Ngược lại, giả sử $\mathfrak{X}$ là noetherian, và cho $\mathfrak{J}$ là một iđêan định nghĩa của $A$, còn $\sh{J}=\mathfrak{J}^\Delta$ là bó iđêan định nghĩa tương ứng của $\mathfrak{X}$.
Các tiền lược đồ (thông thường) $X_n=(\mathfrak{X},\sh{O}_\mathfrak{X}/\sh{J}^{n+1})$ khi đó đều affine và noetherian, nên các vành $A_n=A/\mathfrak{J}^{n+1}$ là noetherian \sref{I.6.1.3}; từ đó suy ra $\mathfrak{J}/\mathfrak{J}^2$ là một $A/\mathfrak{J}$-môđun hữu hạn sinh.
Vì các $\sh{J}^n$ tạo thành một hệ cơ bản các bó iđêan định nghĩa của $\mathfrak{X}$ \sref{I.10.5.1}, ta có $\sh{O}_\mathfrak{X}=\varprojlim\sh{O}_\mathfrak{X}/\sh{J}^n$ \sref{I.10.5.3}; do đó suy ra \sref{I.10.1.3} rằng $A$ đẳng cấu tôpô với $\varprojlim A/\mathfrak{J}^n$ và vì thế là adic và noetherian \sref[0]{0.7.2.8}.
\end{proof}

\begin{remark}[10.6.6]
\label{I.10.6.6}
Với ký hiệu của Mệnh đề~\sref{I.10.6.3}, cho $\sh{F}_i$ là một $\sh{O}_i$-môđun, và giả sử với $i\geq j$ đã cho một $v_{ij}$-cấu xạ $\theta_{ji}:\sh{F}_i\to\sh{F}_j$, sao cho $\theta_{kj}\circ\theta_{ji}=\theta_{ki}$ với $k\leq j\leq i$.
Vì ánh xạ liên tục nền của $v_{ij}$ là ánh xạ đồng nhất, $\theta_{ji}$ là một đồng cấu các bó nhóm abel trên không gian $\mathfrak{X}$; hơn nữa, nếu $\sh{F}$ là giới hạn xạ ảnh của hệ xạ ảnh $(\sh{F}_i)$ các bó nhóm abel, thì việc các $\theta_{ji}$ là $v_{ij}$-cấu xạ cho phép ta xác định trên $\sh{F}$ một cấu trúc $\sh{O}_\mathfrak{X}$-môđun bằng cách chuyển sang giới hạn xạ ảnh; khi được trang bị cấu trúc này, ta nói rằng $\sh{F}$ là \emph{giới hạn xạ ảnh} (đối với các $\theta_{ji}$) của hệ các $\sh{O}_i$-môđun $(\sh{F}_i)$.
Trong trường hợp riêng khi $v_{ij}^*(\sh{F}_i)=\sh{F}_j$ và $\theta_{ji}$ là \emph{đồng nhất}, để viết gọn ta nói rằng $\sh{F}$ là giới hạn xạ ảnh của một hệ $(\sh{F}_i)$ sao cho $v_{ij}^*(\sh{F}_i)=\sh{F}_j$ với $j\leq i$ (không nhắc đến các $\theta_{ji}$).
\end{remark}

\begin{env}[10.6.7]
\label{I.10.6.7}
Cho $\mathfrak{X}$ và $\mathfrak{Y}$ là hai tiền lược đồ hình thức, $\sh{J}$ (tương ứng, $\sh{K}$) là một iđêan định nghĩa của $\mathfrak{X}$ (tương ứng, $\mathfrak{Y}$), và $f:\mathfrak{X}\to\mathfrak{Y}$ là một cấu xạ sao cho $f^*(\sh{K})\sh{O}_\mathfrak{X}\subset\sh{J}$.
Khi đó, với mọi số nguyên $n>0$, ta có $f^*(\sh{K}^n)\sh{O}_\mathfrak{X}=(f^*(\sh{K})\sh{O}_\mathfrak{X})^n\subset\sh{J}^n$; do đó $f$ cảm sinh \sref{I.10.5.6} một cấu xạ các tiền lược đồ (thông thường) $f_n:X_n\to Y_n$, bằng cách đặt $X_n=(\mathfrak{X},\sh{O}_\mathfrak{X}/\sh{J}^{n+1})$ và $Y_n=(\mathfrak{Y},\sh{O}_\mathfrak{Y}/\sh{K}^{n+1})$, và từ các định nghĩa suy ra ngay rằng các sơ đồ
\[
\label{I.10.6.7.1}
  \xymatrix{
    X_m\ar[r]^{f_m}\ar[d] &
    Y_m\ar[d]\\
    X_n\ar[r]_{f_n} &
    Y_n
  }
  \tag{10.6.7.1}
\]
\oldpage[I]{192}
giao hoán với $m\leq n$; nói cách khác, $(f_n)$ là một \emph{hệ quy nạp} các cấu xạ.
\end{env}

\begin{env}[10.6.8]
\label{I.10.6.8}
Ngược lại, cho $(X_n)$ (tương ứng, $(Y_n)$) là một hệ quy nạp các tiền lược đồ (thông thường) thỏa mãn các điều kiện~(b) và (c) của Mệnh đề~\sref{I.10.6.3}, và cho $\mathfrak{X}$ (tương ứng, $\mathfrak{Y}$) là giới hạn quy nạp của nó.
Theo định nghĩa của giới hạn quy nạp, mọi dãy $(f_n)$ các cấu xạ $X_n\to Y_n$ tạo thành một hệ quy nạp đều có một \emph{giới hạn quy nạp $f:\mathfrak{X}\to\mathfrak{Y}$}, là cấu xạ duy nhất của các tiền lược đồ hình thức làm cho các sơ đồ
\[
  \xymatrix{
    X_n\ar[r]^{f_n}\ar[d] &
    Y_n\ar[d]\\
    \mathfrak{X}\ar[r]_f &
    \mathfrak{Y}
  }
\]
giao hoán.
\end{env}

\begin{proposition}[10.6.9]
\label{I.10.6.9}
Cho $\mathfrak{X}$ và $\mathfrak{Y}$ là hai tiền lược đồ hình thức noetherian địa phương, và $\sh{J}$ (tương ứng, $\sh{K}$) là một iđêan định nghĩa của $\mathfrak{X}$ (tương ứng, $\mathfrak{Y}$); ánh xạ $f\mapsto(f_n)$ được định nghĩa trong \sref{I.10.6.7} là một song ánh từ tập các cấu xạ $f:\mathfrak{X}\to\mathfrak{Y}$ thỏa mãn $f^*(\sh{K})\sh{O}_\mathfrak{X}\subset\sh{J}$ lên tập các dãy $(f_n)$ cấu xạ làm cho các sơ đồ \sref{I.10.6.7.1} giao hoán.
\end{proposition}

\begin{proof}
Nếu $f$ là giới hạn quy nạp của một dãy như vậy, còn phải chứng minh rằng $f^*(\sh{K})\sh{O}_\mathfrak{X}\subset\sh{J}$.
Vì vấn đề có tính địa phương trên $\mathfrak{X}$ và $\mathfrak{Y}$, ta có thể quy về trường hợp $\mathfrak{X}=\Spf(A)$ và $\mathfrak{Y}=\Spf(B)$ là affine, với $A$ và $B$ là các vành adic noetherian, và $\sh{J}=\mathfrak{J}^\Delta$, $\sh{K}=\mathfrak{K}^\Delta$, trong đó $\mathfrak{J}$ (tương ứng, $\mathfrak{K}$) là một iđêan định nghĩa của $A$ (tương ứng, $B$).
Khi đó $X_n=\Spec(A_n)$, $Y_n=\Spec(B_n)$, với $A_n=A/\mathfrak{J}^{n+1}$ và $B_n=B/\mathfrak{K}^{n+1}$, theo \sref{I.10.3.6} và \sref{I.10.3.2}; hơn nữa, $f_n=({}^a\vphi_n,\widetilde{\vphi_n})$, trong đó các đồng cấu $\vphi_n:B_n\to A_n$ tạo thành một hệ xạ ảnh, và do đó $f=({}^a\vphi,\widetilde{\vphi})$, với $\vphi=\varprojlim\vphi_n$.
Tính giao hoán của sơ đồ \sref{I.10.6.7.1} khi $m=0$ khi đó cho điều kiện $\vphi_n(\mathfrak{K}/\mathfrak{K}^{n+1})\subset\mathfrak{J}/\mathfrak{J}^{n+1}$ với mọi $n$; chuyển sang giới hạn xạ ảnh, ta thu được $\vphi(\mathfrak{K})\subset\mathfrak{J}$, kéo theo $f^*(\sh{K})\sh{O}_\mathfrak{X}\subset\sh{J}$ \sref{I.10.5.6}[ii].
\end{proof}

\begin{corollary}[10.6.10]
\label{I.10.6.10}
Cho $\mathfrak{X}$ và $\mathfrak{Y}$ là hai tiền lược đồ hình thức noetherian địa phương, và cho $\sh{T}$ là iđêan định nghĩa lớn nhất của $\mathfrak{X}$ \sref{I.10.5.4}.
\begin{enumerate}
  \item[\rm{(i)}] Với mọi iđêan định nghĩa $\sh{K}$ của $\mathfrak{Y}$ và mọi cấu xạ $f:\mathfrak{X}\to\mathfrak{Y}$, ta có $f^*(\sh{K})\sh{O}_\mathfrak{X}\subset\sh{T}$.
  \item[\rm{(ii)}] Có một tương ứng song ánh chính tắc giữa $\Hom(\mathfrak{X},\mathfrak{Y})$ và tập các dãy $(f_n)$ cấu xạ làm cho các sơ đồ \sref{I.10.6.7.1} giao hoán, trong đó $X_n=(\mathfrak{X},\sh{O}_\mathfrak{X}/\sh{T}^{n+1})$ và $Y_n=(\mathfrak{Y},\sh{O}_\mathfrak{Y}/\sh{K}^{n+1})$.
\end{enumerate}
\end{corollary}

\begin{proof}
(ii) suy ra ngay từ (i) và Mệnh đề~\sref{I.10.6.9}.
Để chứng minh (i), ta có thể quy về trường hợp $\mathfrak{X}=\Spf(A)$ và $\mathfrak{Y}=\Spf(B)$, với $A$ và $B$ noetherian, và $\sh{T}=\mathfrak{T}^\Delta$, $\sh{K}=\mathfrak{K}^\Delta$, trong đó $\mathfrak{T}$ là iđêan định nghĩa lớn nhất của $A$ và $\mathfrak{K}$ là một iđêan định nghĩa của $B$.
Cho $f=({}^a\vphi,\widetilde{\vphi})$, trong đó $\vphi:B\to A$ là một đồng cấu liên tục; vì các phần tử của $\mathfrak{K}$ là lũy linh tôpô \sref[0]{0.7.1.4}[ii], các phần tử của $\vphi(\mathfrak{K})$ cũng vậy, nên $\vphi(\mathfrak{K})\subset\mathfrak{T}$ vì $\mathfrak{T}$ là tập các phần tử lũy linh tôpô của $A$ \sref[0]{0.7.1.6}; từ đó có kết luận theo \sref{I.10.5.6}[ii].
\end{proof}

\begin{corollary}[10.6.11]
\label{I.10.6.11}
Cho $\mathfrak{S}$, $\mathfrak{X}$ và $\mathfrak{Y}$ là ba tiền lược đồ hình thức noetherian địa phương, và cho $f:\mathfrak{X}\to\mathfrak{S}$ và $g:\mathfrak{Y}\to\mathfrak{S}$ là các cấu xạ làm cho $\mathfrak{X}$ và $\mathfrak{Y}$ thành các tiền lược đồ hình thức trên $\mathfrak{S}$.
Cho $\sh{J}$ (tương ứng, $\sh{K}$, $\sh{L}$) là một iđêan định nghĩa của $\mathfrak{S}$ (tương ứng, $\mathfrak{X}$, $\mathfrak{Y}$), và giả sử $f^*(\sh{J})\sh{O}_\mathfrak{X}\subset\sh{K}$ và $g^*(\sh{J})\sh{O}_\mathfrak{Y}=\sh{L}$; đặt $S_n=(\mathfrak{S},\sh{O}_\mathfrak{S}/\sh{J}^{n+1})$, $X_n=(\mathfrak{X},\sh{O}_\mathfrak{X}/\sh{K}^{n+1})$ và $Y_n=(\mathfrak{Y},\sh{O}_\mathfrak{Y}/\sh{L}^{n+1})$.
Khi đó có một tương ứng song ánh chính tắc
\oldpage[I]{193}
giữa $\Hom_\mathfrak{S}(\mathfrak{X},\mathfrak{Y})$ và tập các dãy $(u_n)$ gồm các $S_n$-cấu xạ $u_n:X_n\to Y_n$ làm cho các sơ đồ \sref{I.10.6.7.1} giao hoán.
\end{corollary}

\begin{proof}
Với mọi $\mathfrak{S}$-cấu xạ $u:\mathfrak{X}\to\mathfrak{Y}$, theo định nghĩa ta có $f=g\circ u$, do đó
\[
  u^*(\sh{L})\sh{O}_\mathfrak{X}=u^*(g^*(\sh{J})\sh{O}_\mathfrak{Y})\sh{O}_\mathfrak{X}=f^*(\sh{J})\sh{O}_\mathfrak{X}\subset\sh{K}
\]
và hệ quả suy ra từ Mệnh đề~\sref{I.10.6.9}.
\end{proof}

Lưu ý rằng, với $m\leq n$, dữ kiện về một cấu xạ $f_n:X_n\to Y_n$ xác định duy nhất một cấu xạ $f_m:X_m\to Y_m$ làm cho sơ đồ \sref{I.10.6.7.1} giao hoán, như thấy ngay bằng cách quy về trường hợp affine; như vậy ta đã xác định một ánh xạ $\vphi_{mn}:\Hom_{S_n}(X_n,Y_n)\to\Hom_{S_m}(X_m,Y_m)$, và các $\Hom_{S_n}(X_n,Y_n)$ cùng với các $\vphi_{mn}$ tạo thành một \emph{hệ xạ ảnh các tập hợp}; Hệ quả~\sref{I.10.6.11} cũng có thể phát biểu rằng tồn tại một song ánh chính tắc
\[
  \Hom_\mathfrak{S}(\mathfrak{X},\mathfrak{Y})\isoto\varprojlim_n\Hom_{S_n}(X_n,Y_n).
\]

\subsection{Tích của các tiền lược đồ hình thức}
\label{subsection:I.10.7}

\begin{env}[10.7.1]
\label{I.10.7.1}
Cho $\mathfrak{S}$ là một tiền lược đồ hình thức; vì các tiền lược đồ hình thức trên $\mathfrak{S}$ tạo thành một phạm trù, ta có thể định nghĩa khái niệm \emph{tích} của các tiền lược đồ hình thức trên $\mathfrak{S}$.
\end{env}

\begin{proposition}[10.7.2]
\label{I.10.7.2}
Cho $\mathfrak{X}=\Spf(B)$ và $\mathfrak{Y}=\Spf(C)$ là hai lược đồ affine hình thức trên một lược đồ affine hình thức $\mathfrak{S}=\Spf(A)$.
Cho $\mathfrak{Z}=\Spf(B\widehat{\otimes}_A C)$, và cho $p_1$, $p_2$ là các $\mathfrak{S}$-cấu xạ tương ứng, theo \sref{I.10.2.2}, với các $A$-đồng cấu chính tắc (liên tục) $\rho:B\to B\widehat{\otimes}_A C$ và $\sigma:C\to B\widehat{\otimes}_A C$; khi đó $(\mathfrak{Z},p_1,p_2)$ là một tích của hai lược đồ affine hình thức trên $\mathfrak{S}$ là $\mathfrak{X}$ và $\mathfrak{Y}$.
\end{proposition}

\begin{proof}
Theo \sref{I.10.4.6}, chỉ cần kiểm tra rằng ánh xạ gán cho mỗi $A$-đồng cấu liên tục $\vphi:B\widehat{\otimes}_A C\to D$ (trong đó $D$ là một vành khả chấp đồng thời là một $A$-đại số tôpô) cặp $(\vphi\circ\rho,\vphi\circ\sigma)$ là một song ánh
\[
  \Hom_A(B\widehat{\otimes}_A C,D)\isoto\Hom_A(B,D)\times\Hom_A(C,D),
\]
đây chính là tính chất phổ dụng của tích tenxơ đầy đủ \sref[0]{0.7.7.6}.
\end{proof}

\begin{proposition}[10.7.3]
\label{I.10.7.3}
Cho hai $\mathfrak{S}$-tiền lược đồ hình thức $\mathfrak{X}$ và $\mathfrak{Y}$, tích $\mathfrak{X}\times_\mathfrak{S}\mathfrak{Y}$ của chúng tồn tại.
\end{proposition}

\begin{proof}
Chứng minh giống hệt chứng minh của \sref{I.3.2.6}, trong đó thay các lược đồ affine (tương ứng, các tập mở affine) bằng các lược đồ affine hình thức (tương ứng, các tập mở affine hình thức), và thay Mệnh đề \sref{I.3.2.2} bằng Mệnh đề \sref{I.10.7.2}.
\end{proof}

Mọi tính chất hình thức của tích các tiền lược đồ (\sref{I.3.2.7} và \sref{I.3.2.8}, từ \sref{I.3.3.1} đến \sref{I.3.3.12}) đều đúng mà không cần sửa đổi đối với tích các tiền lược đồ hình thức.

\begin{env}[10.7.4]
\label{I.10.7.4}
Cho $\mathfrak{S}$, $\mathfrak{X}$ và $\mathfrak{Y}$ là ba tiền lược đồ hình thức, và cho $f:\mathfrak{X}\to\mathfrak{S}$, $g:\mathfrak{Y}\to\mathfrak{S}$ là hai cấu xạ.
Giả sử trong $\mathfrak{S}$, $\mathfrak{X}$ và $\mathfrak{Y}$ lần lượt tồn tại ba hệ cơ bản các bó iđêan định nghĩa $(\sh{J}_\lambda)$, $(\sh{K}_\lambda)$ và $(\sh{L}_\lambda)$, cùng có tập chỉ số $I$, sao cho $f^*(\sh{J}_\lambda)\sh{O}_\mathfrak{X}\subset\sh{K}_\lambda$ và $g^*(\sh{J}_\lambda)\sh{O}_\mathfrak{Y}\subset\sh{L}_\lambda$ với mọi $\lambda$.
Đặt $S_\lambda=(\mathfrak{S},\sh{O}_\mathfrak{S}/\sh{J}_\lambda)$, $X_\lambda=(\mathfrak{X},\sh{O}_\mathfrak{X}/\sh{K}_\lambda)$ và $Y_\lambda=(\mathfrak{Y},\sh{O}_\mathfrak{Y}/\sh{L}_\lambda)$; nếu $\sh{J}_\mu\subset\sh{J}_\lambda$, $\sh{K}_\mu\subset\sh{K}_\lambda$ và $\sh{L}_\mu\subset\sh{L}_\lambda$, thì $S_\lambda$ (tương ứng, $X_\lambda$, $Y_\lambda$) là một tiền lược đồ con đóng của $S_\mu$ (tương ứng, $X_\mu$, $Y_\mu$) có \emph{cùng}
\oldpage[I]{194}
không gian nền \sref{I.10.6.1}.
Vì $S_\lambda\to S_\mu$ là một đơn cấu tiền lược đồ, trước hết ta thấy các tích $X_\lambda\times_{S_\lambda}Y_\lambda$ và $X_\lambda\times_{S_\mu}Y_\lambda$ là đồng nhất \sref{I.3.2.4}, rồi $X_\lambda\times_{S_\mu}Y_\lambda$ được đồng nhất với một tiền lược đồ con đóng của $X_\mu\times_{S_\mu}Y_\mu$ có \emph{cùng} không gian nền \sref{I.4.3.1}.
Do đó, tích $\mathfrak{X}\times_\mathfrak{S}\mathfrak{Y}$ là \emph{giới hạn quy nạp} của các tiền lược đồ thông thường $X_\lambda\times_{S_\lambda}Y_\lambda$: thật vậy, như trong \sref{I.10.6.2}, ta có thể quy về trường hợp $\mathfrak{S}$, $\mathfrak{X}$ và $\mathfrak{Y}$ là các lược đồ affine hình thức.
Tính đến \sref{I.10.5.6}[ii] và các giả thiết về những hệ cơ bản các bó iđêan định nghĩa của $\mathfrak{S}$, $\mathfrak{X}$ và $\mathfrak{Y}$, ta thấy ngay rằng khẳng định trên suy ra từ định nghĩa tích tenxơ đầy đủ của hai đại số \sref[0]{0.7.7.1}.

Hơn nữa, cho $\mathfrak{Z}$ là một $\mathfrak{S}$-tiền lược đồ hình thức, cho $(\sh{M}_\lambda)$ là một hệ cơ bản các iđêan định nghĩa của $\mathfrak{Z}$ có tập chỉ số $I$, và cho $u:\mathfrak{Z}\to\mathfrak{X}$, $v:\mathfrak{Z}\to\mathfrak{Y}$ là hai $\mathfrak{S}$-cấu xạ sao cho $u^*(\sh{K}_\lambda)\sh{O}_\mathfrak{Z}\subset\sh{M}_\lambda$ và $v^*(\sh{L}_\lambda)\sh{O}_\mathfrak{Z}\subset\sh{M}_\lambda$ với mọi $\lambda$.
Nếu đặt $Z_\lambda=(\mathfrak{Z},\sh{O}_\mathfrak{Z}/\sh{M}_\lambda)$, và nếu $u_\lambda:Z_\lambda\to X_\lambda$, $v_\lambda:Z_\lambda\to Y_\lambda$ là các $S_\lambda$-cấu xạ tương ứng với $u$ và $v$ \sref{I.10.5.6}, thì ta có ngay $(u,v)_\mathfrak{S}$ là giới hạn quy nạp của các $S_\lambda$-cấu xạ $(u_\lambda,v_\lambda)_{S_\lambda}$.

Các lập luận của mục này đặc biệt áp dụng khi $\mathfrak{S}$, $\mathfrak{X}$ và $\mathfrak{Y}$ là địa phương noetherian, bằng cách lấy các hệ gồm những lũy thừa của một bó iđêan định nghĩa \sref{I.10.5.1} làm các hệ cơ bản các bó iđêan định nghĩa.
Tuy nhiên, cần lưu ý rằng $\mathfrak{X}\times_\mathfrak{S}\mathfrak{Y}$ không nhất thiết địa phương noetherian (dù vậy, xem \sref{I.10.13.5}).
\end{env}

\subsection{Đầy đủ hóa hình thức của một tiền lược đồ dọc theo một tập con đóng}
\label{subsection:I.10.8}

\begin{env}[10.8.1]
\label{I.10.8.1}
Cho $X$ là một tiền lược đồ (thông thường) \emph{địa phương noetherian}, và cho $X'$ là một tập con đóng của không gian nền của $X$; ký hiệu $\Phi$ là tập các bó \emph{coherent} các iđêan $\sh{J}$ của $\sh{O}_X$ sao cho giá của $\sh{O}_X/\sh{J}$ là $X'$.
Tập $\Phi$ không rỗng (\sref{I.5.2.1}, \sref{I.4.1.4}, \sref{I.6.1.1}); ta sắp thứ tự nó bởi quan hệ $\supset$.
\end{env}

\begin{lemma}[10.8.2]
\label{I.10.8.2}
Tập có thứ tự $\Phi$ là một tập lọc; nếu $X$ là noetherian thì, với mọi $\sh{J}_0\in\Phi$, tập các lũy thừa $\sh{J}_0^n$ ($n>0$) là đồng tận trong $\Phi$.
\end{lemma}

\begin{proof}
Thật vậy, nếu $\sh{J}_1$ và $\sh{J}_2$ thuộc $\Phi$, và nếu đặt $\sh{J}=\sh{J}_1\cap\sh{J}_2$, thì $\sh{J}$ coherent vì $\sh{O}_X$ coherent (\sref{I.6.1.1} và \sref[0]{0.5.3.4}), và ta có $\sh{J}_x=(\sh{J}_1)_x\cap(\sh{J}_2)_x$ với mọi $x\in X$; do đó $\sh{J}_x=\sh{O}_x$ khi $x\not\in X'$, còn $\sh{J}_x\neq\sh{O}_x$ khi $x\in X'$, điều này chứng minh $\sh{J}\in\Phi$.
Mặt khác, nếu $X$ là noetherian và $\sh{J}_0$, $\sh{J}$ thuộc $\Phi$, thì tồn tại một số nguyên $n>0$ sao cho $\sh{J}_0^n(\sh{O}_X/\sh{J})=0$ \sref{I.9.3.4}, kéo theo $\sh{J}_0^n\subset\sh{J}$.
\end{proof}

\begin{env}[10.8.3]
\label{I.10.8.3}
Cho $\sh{F}$ là một $\sh{O}_X$-môđun \emph{coherent}; với mọi $\sh{J}\in\Phi$, $\sh{F}\otimes_{\sh{O}_X}(\sh{O}_X/\sh{J})$ là một $\sh{O}_X$-môđun coherent \sref{I.9.1.1}, có giá nằm trong $X'$, và ta thường đồng nhất nó với hạn chế của nó lên $X'$.
Khi $\sh{J}$ chạy qua $\Phi$, các bó này tạo thành một \emph{hệ xạ ảnh} các bó nhóm abel.
\end{env}

\begin{definition}[10.8.4]
\label{I.10.8.4}
Cho một tập con đóng $X'$ của một tiền lược đồ địa phương noetherian $X$ và một $\sh{O}_X$-môđun coherent $\sh{F}$, ta gọi \emph{đầy đủ hóa của $\sh{F}$ dọc theo $X'$}, ký hiệu bởi $\sh{F}_{/X'}$ (hoặc bởi $\widehat{\sh{F}}$ khi không có nguy cơ nhầm lẫn), là hạn chế lên $X'$ của bó
\oldpage[I]{195}
$\varprojlim_\Phi(\sh{F}\otimes_{\sh{O}_X}(\sh{O}_X/\sh{J}))$; ta gọi các tiết diện của nó trên $X'$ là các \emph{tiết diện hình thức của $\sh{F}$ dọc theo $X'$}.
\end{definition}

Hiển nhiên rằng, với mọi tập mở $U\subset X$, ta có $(\sh{F}|U)_{/(U\cap X')}=(\sh{F}_{/X'})|(U\cap X')$.

Khi lấy giới hạn xạ ảnh, rõ ràng $(\sh{O}_X)_{/X'}$ là một bó vành, và $\sh{F}_{/X'}$ có thể được xem như một $(\sh{O}_X)_{/X'}$-môđun.
Hơn nữa, vì tôpô của $X'$ có một cơ sở gồm các tập mở quasi-compact, ta có thể xem $(\sh{O}_X)_{/X'}$ (tương ứng, $\sh{F}_{/X'}$) như một \emph{bó các vành tôpô} (tương ứng, một bó \emph{các nhóm tôpô}), là giới hạn xạ ảnh của các bó vành (tương ứng, nhóm) \emph{giả rời rạc} $\sh{O}_X/\sh{J}$ (tương ứng, $\sh{F}\otimes_{\sh{O}_X}(\sh{O}_X/\sh{J})=\sh{F}/\sh{J}\sh{F}$), và, khi lấy giới hạn xạ ảnh, $\sh{F}_{/X'}$ khi đó trở thành một \emph{$(\sh{O}_X)_{/X'}$-môđun tôpô} (\sref[0]{0.3.8.1} và \sref[0]{0.3.8.2}); nhắc lại rằng, với mọi tập mở \emph{quasi-compact} $U\subset X$, $\Gamma(U\cap X',(\sh{O}_X)_{/X'})$ (tương ứng, $\Gamma(U\cap X',\sh{F}_{/X'})$) khi đó là giới hạn xạ ảnh của các vành (tương ứng, nhóm) rời rạc $\Gamma(U,\sh{O}_X/\sh{J})$ (tương ứng, $\Gamma(U,\sh{F}/\sh{J}\sh{F})$).

Nếu $u:\sh{F}\to\sh{G}$ là một đồng cấu các $\sh{O}_X$-môđun, thì ta suy ra chính tắc các đồng cấu $u_\sh{J}:\sh{F}\otimes_{\sh{O}_X}(\sh{O}_X/\sh{J})\to\sh{G}\otimes_{\sh{O}_X}(\sh{O}_X/\sh{J})$ với mọi $\sh{J}\in\Phi$, và các đồng cấu này tạo thành một hệ xạ ảnh.
Khi lấy giới hạn xạ ảnh và hạn chế lên $X'$, chúng cho một $(\sh{O}_X)_{/X'}$-đồng cấu liên tục $\sh{F}_{/X'}\to\sh{G}_{/X'}$, ký hiệu $u_{/X'}$ hoặc $\widehat{u}$, gọi là \emph{đầy đủ hóa của đồng cấu $u$ dọc theo $X'$}.
Rõ ràng, nếu $v:\sh{G}\to\sh{H}$ là một đồng cấu thứ hai của các $\sh{O}_X$-môđun, thì $(v\circ u)_{/X'}=(v_{/X'})\circ(u_{/X'})$; do đó $\sh{F}_{/X'}$ là một \emph{hàm tử hiệp biến cộng tính} theo $\sh{F}$, từ phạm trù các $\sh{O}_X$-môđun coherent đến phạm trù các $(\sh{O}_X)_{/X'}$-môđun tôpô.

\begin{proposition}[10.8.5]
\label{I.10.8.5}
Giá của $(\sh{O}_X)_{/X'}$ là $X'$; không gian vành tôpô $(X',(\sh{O}_X)_{/X'})$ là một tiền lược đồ hình thức địa phương noetherian, và nếu $\sh{J}\in\Phi$ thì $\sh{J}_{/X'}$ là một bó iđêan định nghĩa của tiền lược đồ hình thức này.
Nếu $X=\Spec(A)$ là một lược đồ affine có vành noetherian $A$, $\sh{J}=\widetilde{\mathfrak{J}}$, trong đó $\mathfrak{J}$ là một iđêan của $A$, và $X'=V(\mathfrak{J})$, thì $(X',(\sh{O}_X)_{/X'})$ được đồng nhất chính tắc với $\Spf(\widehat{A})$, trong đó $\widehat{A}$ là đầy đủ hóa tách của $A$ đối với tôpô $\mathfrak{J}$-tiền adic.
\end{proposition}

\begin{proof}
Hiển nhiên chỉ cần chứng minh khẳng định sau cùng.
Theo \sref[0]{0.7.3.3}, đầy đủ hóa tách $\widehat{\mathfrak{J}}$ của $\mathfrak{J}$ đối với tôpô $\mathfrak{J}$-tiền adic được đồng nhất với iđêan $\mathfrak{J}\widehat{A}$ của $\widehat{A}$, và $\widehat{A}$ là một vành noetherian $\widehat{\mathfrak{J}}$-adic sao cho $\widehat{A}/\widehat{\mathfrak{J}}^n=A/\mathfrak{J}^n$ \sref[0]{0.7.2.6}.
Đẳng thức sau cùng này cho thấy các iđêan nguyên tố mở của $\widehat{A}$ là các iđêan $\widehat{\mathfrak{p}}=\mathfrak{p}\widehat{A}$, trong đó $\mathfrak{p}$ là một iđêan nguyên tố của $A$ chứa $\mathfrak{J}$, và $\widehat{\mathfrak{p}}\cap A=\mathfrak{p}$; do đó $\Spf(\widehat{A})=X'$.
Vì $\sh{O}_X/\sh{J}^n=(A/\mathfrak{J}^n)\supertilde$, mệnh đề suy ra ngay từ các định nghĩa.
\end{proof}

Ta gọi tiền lược đồ hình thức vừa định nghĩa là \emph{đầy đủ hóa của $X$ dọc theo $X'$}, và ký hiệu nó bởi $X_{/X'}$ hoặc $\widehat{X}$ khi không có nguy cơ nhầm lẫn.
Khi lấy $X'=X$, ta có thể lấy $\sh{J}=0$, và do đó $X_{/X}=X$.

Rõ ràng, nếu $U$ là một tiền lược đồ con cảm sinh trên một tập mở của $X$, thì $U_{/(U\cap X')}$ được đồng nhất chính tắc với tiền lược đồ con hình thức do $X_{/X'}$ cảm sinh trên tập mở $U\cap X'$ của $X'$.

\begin{corollary}[10.8.6]
\label{I.10.8.6}
Tiền lược đồ (thông thường) $\widehat{X}_\red$ là tiền lược đồ con rút gọn duy nhất của $X$ có không gian nền là $X'$ \sref{I.5.2.1}.
Để $\widehat{X}$ là noetherian, điều kiện cần và đủ là $\widehat{X}_\red$ là noetherian, và chỉ cần $X$ là noetherian.
\end{corollary}

\oldpage[I]{196}

\begin{proof}
Vì $\widehat{X}_\red$ được xác định cục bộ \sref{I.10.5.4}, ta có thể giả sử $X$ là một lược đồ affine của một vành noetherian $A$; với các ký hiệu của \sref{I.10.8.5}, iđêan $\mathfrak{T}$ gồm các phần tử lũy linh tôpô của $\widehat{A}$ là ảnh ngược, qua ánh xạ chính tắc $\widehat{A}\to\widehat{A}/\widehat{\mathfrak{J}}=A/\mathfrak{J}$, của căn lũy linh của $A/\mathfrak{J}$ \sref[0]{0.7.1.3}; do đó $\widehat{A}/\mathfrak{T}$ đẳng cấu với thương của $A/\mathfrak{J}$ theo căn lũy linh của nó.
Khẳng định thứ nhất suy ra từ \sref{I.10.5.4} và \sref{I.5.1.1}.
Nếu $\widehat{X}_\red$ là noetherian thì không gian nền $X'$ của nó cũng vậy; do đó các $X_n'=\Spec(\sh{O}_X/\sh{J}^n)$ là noetherian \sref{I.6.1.2}, và $\widehat{X}$ cũng vậy \sref{I.10.6.4}; chiều ngược lại là hiển nhiên theo \sref{I.6.1.2}.
\end{proof}

\begin{env}[10.8.7]
\label{I.10.8.7}
Các đồng cấu chính tắc $\sh{O}_X\to\sh{O}_X/\sh{J}$ (với $\sh{J}\in\Phi$) tạo thành một hệ xạ ảnh và do đó, khi chuyển sang giới hạn xạ ảnh, cho một đồng cấu các bó vành $\theta:\sh{O}_X\to\psi_*((\sh{O}_X)_{/X'})=\varprojlim_\Phi(\sh{O}_X/\sh{J})$, trong đó $\psi$ ký hiệu phép nhúng chính tắc $X'\to X$ của các không gian nền.
Ta ký hiệu bởi $i$ (hoặc $i_X$) cấu xạ (gọi là \emph{chính tắc})
\[
  (\psi,\theta):X_{/X'}\to X
\]
của các không gian vành.

Bằng cách lấy tích tenxơ, với mọi $\sh{O}_X$-môđun coherent $\sh{F}$, các đồng cấu chính tắc $\sh{O}_X\to\sh{O}_X/\sh{J}$ cho các đồng cấu $\sh{F}\to\sh{F}\otimes_{\sh{O}_X}(\sh{O}_X/\sh{J})$ của các $\sh{O}_X$-môđun; các đồng cấu này lại tạo thành một hệ xạ ảnh và do đó, khi chuyển sang giới hạn xạ ảnh, cho một đồng cấu chính tắc có tính hàm tử $\gamma:\sh{F}\to\psi_*(\sh{F}_{/X'})$ của các $\sh{O}_X$-môđun.
\end{env}

\begin{proposition}[10.8.8]
\label{I.10.8.8}
\medskip\noindent
\begin{enumerate}
  \item[\rm{(i)}] Hàm tử $\sh{F}_{/X'}$ (theo $\sh{F}$) là khớp.
  \item[\rm{(ii)}] Đồng cấu có tính hàm tử $\gamma^\sharp:i^*(\sh{F})\to\sh{F}_{/X'}$ của các $(\sh{O}_X)_{/X'}$-môđun là một đẳng cấu.
\end{enumerate}
\end{proposition}

\begin{proof}
\medskip\noindent
\begin{enumerate}
  \item[\rm{(i)}] Chỉ cần chứng minh rằng, nếu $0\to\sh{F}'\to\sh{F}\to\sh{F}''\to 0$ là một dãy khớp các $\sh{O}_X$-môđun coherent và nếu $U$ là một tập mở affine của $X$ có vành noetherian $A$, thì dãy
\[
  0\to\Gamma(U\cap X',\sh{F}'_{/X'})\to\Gamma(U\cap X',\sh{F}_{/X'})\to\Gamma(U\cap X',\sh{F}''_{/X'})\to 0
\]
là khớp.
Ta có $\sh{F}|U=\widetilde{M}$, $\sh{F}'|U=\widetilde{M'}$ và $\sh{F}''|U=\widetilde{M''}$, trong đó $M$, $M'$ và $M''$ là ba $A$-môđun thuộc loại hữu hạn sao cho dãy $0\to M'\to M\to M''\to 0$ là khớp (\sref{I.1.5.1} và \sref{I.1.3.11}); lấy $\sh{J}\in\Phi$ và một iđêan $\mathfrak{J}$ của $A$ sao cho $\sh{J}|U=\widetilde{\mathfrak{J}}$.
Khi ấy
\[
  \Gamma(U\cap X',\sh{F}\otimes_{\sh{O}_X}\sh{O}_X/\sh{J}^n)=M\otimes_A(A/\mathfrak{J}^n)
\]
\sref{I.1.3.12}; do đó, theo định nghĩa của giới hạn xạ ảnh, ta có
\[
  \Gamma(U\cap X',\sh{F}_{/X'})=\varprojlim_n(M\otimes_A(A/\mathfrak{J}^n))=\widehat{M},
\]
là đầy đủ hóa tách của $M$ đối với tôpô $\mathfrak{J}$-tiền adic, và tương tự
\[
  \Gamma(U\cap X',\sh{F}'_{/X'})=\widehat{M'},\ \Gamma(U\cap X',\sh{F}''_{/X'})=\widehat{M''};
\]
khẳng định của ta khi đó suy ra từ việc, do $A$ noetherian, hàm tử $\widehat{M}$ theo $M$ là khớp trên phạm trù các $A$-môđun thuộc loại hữu hạn \sref[0]{0.7.3.3}.
\oldpage[I]{197}
  \item[\rm{(ii)}] Bài toán có tính cục bộ, nên ta có thể giả sử có một dãy khớp $\sh{O}_X^m\to\sh{O}_X^n\to\sh{F}\to 0$ \sref[0]{0.5.3.2}; vì $\gamma^\sharp$ có tính hàm tử và các hàm tử $i^*(\sh{F})$ và $\sh{F}_{/X'}$ là khớp phải (theo (i) và \sref[0]{0.4.3.1}), ta có sơ đồ giao hoán
\[
\label{I.10.8.8.1}
  \xymatrix{
    i^*(\sh{O}_X^m)\ar[r]\ar[d]_{\gamma^\sharp} &
    i^*(\sh{O}_X^n)\ar[r]\ar[d]_{\gamma^\sharp} &
    i^*(\sh{F})\ar[r]\ar[d]_{\gamma^\sharp} &
    0\\
    (\sh{O}_X^m)_{/X'}\ar[r] &
    (\sh{O}_X^n)_{/X'}\ar[r] &
    \sh{F}_{/X'}\ar[r] &
    0
  }
  \tag{10.8.8.1}
\]
có các hàng khớp.
Hơn nữa, các hàm tử $i^*(\sh{F})$ và $\sh{F}_{/X'}$ giao hoán với các tổng trực tiếp hữu hạn (\sref[0]{0.3.2.6} và \sref[0]{0.4.3.2}), nên ta quy về việc chứng minh khẳng định cho $\sh{F}=\sh{O}_X$.
Ta có $i^*(\sh{O}_X)=(\sh{O}_X)_{/X'}=\sh{O}_{\widehat{X}}$ \sref[0]{0.4.3.4}, và $\gamma^\sharp$ là một đồng cấu của các \emph{$\sh{O}_{\widehat{X}}$-môđun}; do đó chỉ cần kiểm tra rằng $\gamma^\sharp$ đưa tiết diện đơn vị của $\sh{O}_{\widehat{X}}$ trên một tập mở của $X'$ vào chính nó, điều này là hiển nhiên và cho thấy trong trường hợp này $\gamma^\sharp$ là đồng nhất.
\end{enumerate}
\end{proof}

\begin{corollary}[10.8.9]
\label{I.10.8.9}
Cấu xạ $i:X_{/X'}\to X$ của các không gian vành là phẳng.
\end{corollary}

\begin{proof}
Điều này suy ra từ \sref[0]{0.6.7.3} và \sref{I.10.8.8}[i].
\end{proof}

\begin{corollary}[10.8.10]
\label{I.10.8.10}
Nếu $\sh{F}$ và $\sh{G}$ là các $\sh{O}_X$-môđun coherent, thì tồn tại các đẳng cấu chính tắc có tính hàm tử (theo $\sh{F}$ và $\sh{G}$)
\[
\label{I.10.8.10.1}
  (\sh{F}_{/X'})\otimes_{(\sh{O}_X)_{/X'}}(\sh{G}_{/X'})\isoto(\sh{F}\otimes_{\sh{O}_X}\sh{G})_{/X'},
  \tag{10.8.10.1}
\]
\[
\label{I.10.8.10.2}
  (\shHom_{\sh{O}_X}(\sh{F},\sh{G}))_{/X'}\isoto\shHom_{(\sh{O}_X)_{/X'}}(\sh{F}_{/X'},\sh{G}_{/X'}).
  \tag{10.8.10.2}
\]
\end{corollary}

\begin{proof}
Điều này suy ra từ sự đồng nhất chính tắc của $i^*(\sh{F})$ với $\sh{F}_{/X'}$; sự tồn tại của đẳng cấu thứ nhất khi đó là một kết quả đúng với mọi cấu xạ của các không gian vành \sref[0]{0.4.3.3.1}, còn sự tồn tại của đẳng cấu thứ hai là một kết quả đúng với mọi cấu xạ phẳng \sref[0]{0.6.7.6}, theo \sref{I.10.8.9}.
\end{proof}

\begin{proposition}[10.8.11]
\label{I.10.8.11}
Với mọi $\sh{O}_X$-môđun coherent $\sh{F}$, hạt nhân của đồng cấu chính tắc $\Gamma(X,\sh{F})\to\Gamma(X',\sh{F}_{/X'})$ cảm sinh bởi $\sh{F}\to\sh{F}_{/X'}$ gồm các tiết diện bằng không trên một lân cận nào đó của $X'$.
\end{proposition}

\begin{proof}
Từ định nghĩa của $\sh{F}_{/X'}$ suy ra ảnh chính tắc của một tiết diện như vậy bằng không.
Ngược lại, nếu $s\in\Gamma(X,\sh{F})$ có ảnh bằng không trong $\Gamma(X',\sh{F}_{/X'})$, thì chỉ cần thấy rằng mọi $x\in X'$ đều có một lân cận trong $X$ mà trên đó $s$ bằng không; do đó ta có thể quy về trường hợp $X=\Spec(A)$ là affine, $A$ noetherian, $X'=V(\mathfrak{J})$ với một iđêan $\mathfrak{J}$ nào đó của $A$, và $\sh{F}=\widetilde{M}$ với một $A$-môđun $M$ thuộc loại hữu hạn nào đó.
Khi ấy $\Gamma(X',\sh{F}_{/X'})$ là đầy đủ hóa tách $\widehat{M}$ của $M$ đối với tôpô $\mathfrak{J}$-tiền adic, và đồng cấu $\Gamma(X,\sh{F})\to\Gamma(X',\sh{F}_{/X'})$ là đồng cấu chính tắc $M\to\widehat{M}$.
Theo \sref[0]{0.7.3.7}, hạt nhân của đồng cấu này là tập các $z\in M$ bị triệt tiêu bởi một phần tử của $1+\mathfrak{J}$.
Do đó $(1+f)s=0$ với một $f\in\mathfrak{J}$ nào đó; với mọi $x\in X'$ suy ra $(1_x+f_x)s_x=0$, và vì $1_x+f_x$ khả nghịch trong $\sh{O}_x$ (do $\mathfrak{J}_x\sh{O}_x$ nằm trong iđêan cực đại của $\sh{O}_x$), ta có $s_x=0$, điều này chứng minh mệnh đề.
\end{proof}

\begin{corollary}[10.8.12]
\label{I.10.8.12}
Giá của $\sh{F}_{/X'}$ bằng $\Supp(\sh{F})\cap X'$.
\end{corollary}

\begin{proof}
Hiển nhiên $\sh{F}_{/X'}$ là một $(\sh{O}_X)_{/X'}$-môđun thuộc loại hữu hạn (\sref{I.10.8.8}[ii] và \sref[0]{0.5.2.4}),
\oldpage[I]{198}
do đó giá của nó là đóng \sref[0]{0.5.2.2} và hiển nhiên nằm trong $\Supp(\sh{F})\cap X'$.
Để chứng minh nó bằng tập sau cùng này, ta quy ngay về việc chứng minh rằng đẳng thức $\Gamma(X',\sh{F}_{/X'})=0$ kéo theo $\Supp(\sh{F})\cap X'=\emp$; điều này suy ra từ \sref{I.10.8.11} và \sref{I.1.4.1}.
\end{proof}

\begin{corollary}[10.8.13]
\label{I.10.8.13}
Cho $u:\sh{F}\to\sh{G}$ là một đồng cấu của các $\sh{O}_X$-môđun coherent.
Để $u_{/X'}:\sh{F}_{/X'}\to\sh{G}_{/X'}$ bằng không, điều kiện cần và đủ là $u$ bằng không trên một lân cận của $X'$.
\end{corollary}

\begin{proof}
Thật vậy, theo \sref{I.10.8.8}[ii], $u_{/X'}$ được đồng nhất với $i^*(u)$; do đó, nếu xem $u$ như một tiết diện trên $X$ của bó $\sh{H}=\shHom_{\sh{O}_X}(\sh{F},\sh{G})$, thì $u_{/X'}$ là tiết diện trên $X'$ của $i^*(\sh{H})=\sh{H}_{/X'}$ tương ứng chính tắc với nó (\sref{I.10.8.10.2} và \sref[0]{0.4.4.6}).
Vì vậy chỉ cần áp dụng \sref{I.10.8.11} cho $\sh{O}_X$-môđun coherent $\sh{H}$.
\end{proof}

\begin{corollary}[10.8.14]
\label{I.10.8.14}
Cho $u:\sh{F}\to\sh{G}$ là một đồng cấu của các $\sh{O}_X$-môđun coherent.
Để $u_{/X'}$ là một đơn cấu (tương ứng, một toàn cấu), điều kiện cần và đủ là $u$ là một đơn cấu (tương ứng, một toàn cấu) trên một lân cận của $X'$.
\end{corollary}

\begin{proof}
Gọi $\sh{P}$ và $\sh{N}$ lần lượt là đối hạt nhân và hạt nhân của $u$, sao cho ta có dãy khớp $0\to\sh{N}\xrightarrow{v}\sh{F}\xrightarrow{u}\sh{G}\xrightarrow{w}\sh{P}\to 0$, và do đó, theo \sref{I.10.8.8}[i], có dãy khớp
\[
  0\to\sh{N}_{/X'}\xrightarrow{v_{/X'}}\sh{F}_{/X'}\xrightarrow{u_{/X'}}\sh{G}_{/X'}\xrightarrow{w_{/X'}}\sh{P}_{/X'}\to 0.
\]
Nếu $u_{/X'}$ là một đơn cấu (tương ứng, một toàn cấu), thì $v_{/X'}=0$ (tương ứng, $w_{/X'}=0$); do đó tồn tại một lân cận của $X'$ mà trên đó $v=0$ (tương ứng, $w=0$), theo \sref{I.10.8.13}.
\end{proof}

\subsection{Mở rộng một cấu xạ đến các đầy đủ hóa}
\label{subsection:I.10.9}

\begin{env}[10.9.1]
\label{I.10.9.1}
Cho $X$ và $Y$ là hai tiền lược đồ (thông thường) Noether địa phương, $f:X\to Y$ là một cấu xạ, và $X'$ (tương ứng, $Y'$) là một tập con đóng của không gian nền $X$ (tương ứng, $Y$), sao cho $f(X')\subset Y'$.
Cho $\sh{J}$ (tương ứng, $\sh{K}$) là một bó iđêan của $\sh{O}_X$ (tương ứng, $\sh{O}_Y$) sao cho giá của $\sh{O}_X/\sh{J}$ (tương ứng, $\sh{O}_Y/\sh{K}$) là $X'$ (tương ứng, $Y'$) và $f^*(\sh{K})\sh{O}_X\subset\sh{J}$; lưu ý rằng luôn tồn tại các bó iđêan như thế, vì chẳng hạn có thể lấy $\sh{J}$ là bó iđêan lớn nhất của $\sh{O}_X$ xác định một tiền lược đồ con của $X$ có không gian nền là $X'$ \sref{I.5.2.1}, và giả thiết $f(X')\subset Y'$ khi đó kéo theo $f^*(\sh{K})\sh{O}_X\subset\sh{J}$ \sref{I.5.2.4}.
Với mọi số nguyên $n>0$, ta có $f^*(\sh{K}^n)\sh{O}_X\subset\sh{J}^n$ \sref[0]{0.4.3.5}; do đó, theo \sref{I.4.4.6}, nếu đặt $X_n'=(X',\sh{O}_X/\sh{J}^{n+1})$ và $Y_n'=(Y',\sh{O}_Y/\sh{K}^{n+1})$, thì $f$ cảm sinh một cấu xạ $f_n:X_n'\to Y_n'$, và hiển nhiên các $f_n$ tạo thành một hệ quy nạp.
Ta ký hiệu giới hạn quy nạp của nó \sref{I.10.6.8} bởi $\widehat{f}:X_{/X'}\to Y_{/Y'}$, và nói (bằng cách lạm dụng ngôn ngữ) rằng $\widehat{f}$ là \emph{mở rộng của $f$ đến các đầy đủ hóa của $X$ và $Y$ dọc theo $X'$ và $Y'$}.
Có thể kiểm tra ngay rằng cấu xạ này không phụ thuộc vào việc chọn các bó iđêan $\sh{J}$ và $\sh{K}$ thỏa mãn các điều kiện trên.
Thật vậy, chỉ cần xét trường hợp $X$ và $Y$ là các lược đồ affine Noether có các vành lần lượt là $A$ và $B$; khi đó $\sh{J}=\widetilde{\mathfrak{J}}$ và $\sh{K}=\widetilde{\mathfrak{K}}$, trong đó $\mathfrak{J}$ (tương ứng, $\mathfrak{K}$) là một iđêan của $A$ (tương ứng, $B$), và $f$ tương ứng với một đồng cấu vành $\vphi:B\to A$ sao cho $\vphi(\mathfrak{K})\subset\mathfrak{J}$ (\sref{I.4.4.6} và \sref{I.1.7.4}); \emph{khi đó $\widehat{f}$ là cấu xạ tương ứng \sref{I.10.2.2} với đồng cấu liên tục $\widehat{\vphi}:\widehat{B}\to\widehat{A}$}, trong đó $\widehat{A}$ (tương ứng, $\widehat{B}$) là đầy đủ hóa tách của $A$ (tương ứng, $B$) đối với tôpô $\mathfrak{J}$-tiền adic (tương ứng, $\mathfrak{K}$-tiền adic) \sref{I.10.6.8}; ta biết rằng nếu thay $\sh{J}$ bằng một
\oldpage[I]{199}
bó iđêan khác $\sh{J}'=\widetilde{\mathfrak{J}'}$ sao cho giá của $\sh{O}_X/\sh{J}'$ vẫn là $X'$, thì các tôpô $\mathfrak{J}$-tiền adic và $\mathfrak{J}'$-tiền adic trên $A$ là như nhau \sref{I.10.8.2}.

Lưu ý rằng, theo định nghĩa, ánh xạ liên tục $X'\to Y'$ của các không gian nền của $X_{/X'}$ và $Y_{/Y'}$ tương ứng với $\widehat{f}$ chính là hạn chế của $f$ lên $X'$.
\end{env}

\begin{env}[10.9.2]
\label{I.10.9.2}
Từ định nghĩa trên suy ra ngay rằng sơ đồ các cấu xạ của không gian vành
\[
  \label{I.10.9.2.1}
  \xymatrix{
    \widehat{X}\ar[r]^{\widehat{f}}\ar[d]_{i_X} &
    \widehat{Y}\ar[d]^{i_Y}\\
    X\ar[r]_f &
    Y
  }
\]
là giao hoán, trong đó các mũi tên thẳng đứng là các cấu xạ chính tắc \sref{I.10.8.7}.
\end{env}

\begin{env}[10.9.3]
\label{I.10.9.3}
Cho $Z$ là một tiền lược đồ thứ ba, $g:Y\to Z$ là một cấu xạ và $Z'$ là một tập con đóng của $Z$ sao cho $g(Y')\subset Z'$.
Nếu $\widehat{g}$ chỉ đầy đủ hóa của cấu xạ $g$ dọc theo $Y'$ và $Z'$, thì từ \sref{I.10.9.1} suy ra ngay rằng
$\widehat{(g\circ f)}=\widehat{g}\circ\widehat{f}$.
\end{env}

\begin{proposition}[10.9.4]
\label{I.10.9.4}
Cho $X$ và $Y$ là hai $S$-tiền lược đồ Noether địa phương, trong đó $Y$ hữu hạn kiểu trên $S$.
Cho $f$ và $g$ là hai $S$-cấu xạ từ $X$ đến $Y$ sao cho $f(X')\subset Y'$ và $g(X')\subset Y'$.
Ta có $\widehat{f}=\widehat{g}$ khi và chỉ khi $f$ và $g$ trùng nhau trên một lân cận của $X'$.
\end{proposition}

\begin{proof}
Điều kiện hiển nhiên là đủ (không cần giả thiết hữu hạn đối với $Y$).
Để thấy điều kiện này là cần, trước hết lưu ý rằng giả thiết $\widehat{f}=\widehat{g}$ kéo theo $f(x)=g(x)$ với mọi $x\in X'$.
Mặt khác, vì bài toán có tính địa phương, ta có thể giả sử $X$ và $Y$ lần lượt là các lân cận mở affine của $x$ và của $y=f(x)=g(x)$, với các vành Noether, rằng $S$ là affine và $\Gamma(Y,\sh{O}_Y)$ là một $\Gamma(S,\sh{O}_S)$-đại số hữu hạn kiểu \sref{I.6.3.3}.
Khi đó $f$ và $g$ tương ứng với hai $\Gamma(S,\sh{O}_S)$-đồng cấu $\rho$ và $\sigma$ từ $\Gamma(Y,\sh{O}_Y)$ vào $\Gamma(X,\sh{O}_X)$ \sref{I.1.7.3}, và theo giả thiết, các mở rộng liên tục của những đồng cấu này đến đầy đủ hóa tách của $\Gamma(Y,\sh{O}_Y)$ là như nhau.
Từ \sref{I.10.8.11} suy ra rằng, với mọi tiết diện $s\in\Gamma(Y,\sh{O}_Y)$, các tiết diện $\rho(s)$ và $\sigma(s)$ trùng nhau trên một lân cận của $X'$ (phụ thuộc vào $s$); vì $\Gamma(Y,\sh{O}_Y)$ là một đại số hữu hạn kiểu trên $\Gamma(S,\sh{O}_S)$, suy ra ngay rằng tồn tại một lân cận $V$ của $X'$ sao cho $\rho(s)$ và $\sigma(s)$ trùng nhau trên $V$ với \emph{mọi} tiết diện $s\in\Gamma(Y,\sh{O}_Y)$.
Nếu $h\in\Gamma(X,\sh{O}_X)$\footnote{\emph{[Dịch.] Nguồn tiếng Pháp in $\Gamma(Y,\sh{O}_X)$ tại đây; miền phải là $X$ để $D(h)\subset X$ có nghĩa.}} sao cho $D(h)$ là một lân cận của $x$ nằm trong $V$, thì từ điều trên và \sref{I.1.4.1}[d] suy ra rằng $f$ và $g$ trùng nhau trên $D(h)$.
\end{proof}

\begin{proposition}[10.9.5]
\label{I.10.9.5}
Dưới các giả thiết của \sref{I.10.9.1}, với mọi $\sh{O}_Y$-môđun coherent $\sh{G}$, tồn tại một đẳng cấu chính tắc có tính hàm tử của các $(\sh{O}_X)_{/X'}$-môđun
\[
  (f^*(\sh{G}))_{/X'}\isoto\widehat{f}^*(\sh{G}_{/Y'}).
\]
\end{proposition}

\begin{proof}
Nếu ta đồng nhất chính tắc $(f^*(\sh{G}))_{/X'}$ với $i_X^*(f^*(\sh{G}))$, và $\widehat{f}^*(\sh{G}_{/Y'})$ với $\widehat{f}^*(i_Y^*(\sh{G}))$ \sref{I.10.8.8}, thì mệnh đề suy ra ngay từ tính giao hoán của sơ đồ trong \sref{I.10.9.2}.
\end{proof}

\begin{env}[10.9.6]
\label{I.10.9.6}
Bây giờ cho $\sh{F}$ là một $\sh{O}_X$-môđun coherent và $\sh{G}$ là một $\sh{O}_Y$-môđun coherent.
Nếu $u:\sh{G}\to\sh{F}$ là một $f$-cấu xạ từ $\sh{G}$ đến $\sh{F}$, thì nó tương ứng với một $\sh{O}_X$-đồng cấu $u^\sharp:f^*(\sh{G})\to\sh{F}$, và do đó, bằng cách đầy đủ hóa, với một $(\sh{O}_X)_{/X'}$-đồng cấu liên tục $(u^\sharp)_{/X'}:(f^*(\sh{G}))_{/X'}\to\sh{F}_{/X'}$; theo \sref{I.10.9.5}, tồn tại một $\widehat{f}$-cấu xạ $v:\sh{G}_{/Y'}\to\sh{F}_{/X'}$,
\oldpage[I]{200}
và chỉ một, sao cho $v^\sharp=(u^\sharp)_{/X'}$.
Nếu coi các bộ ba $(\sh{F},X,X')$ (trong đó $\sh{F}$ là một $\sh{O}_X$-môđun coherent và $X'$ là một tập con đóng của $X$) như một \emph{phạm trù}, mà các cấu xạ $(\sh{F},X,X')\to(\sh{G},Y,Y')$ gồm một cấu xạ tiền lược đồ $f:X\to Y$ sao cho $f(X')\subset Y'$ và một $f$-cấu xạ $u:\sh{G}\to\sh{F}$, thì có thể nói rằng $(X_{/X'},\sh{F}_{/X'})$ là một \emph{hàm tử} theo $(\sh{F},X,X')$, nhận giá trị trong phạm trù các cặp $(\mathfrak{Z},\sh{H})$ gồm một tiền lược đồ hình thức Noether địa phương $\mathfrak{Z}$ và một $\sh{O}_\mathfrak{Z}$-môđun $\sh{H}$, trong đó các cấu xạ của phạm trù sau là các cặp gồm một cấu xạ $g$ của các tiền lược đồ hình thức và một $g$-cấu xạ.
\end{env}

\begin{proposition}[10.9.7]
\label{I.10.9.7}
Cho $S$, $X$ và $Y$ là ba tiền lược đồ Noether địa phương, $g:X\to S$ và $h:Y\to S$ là hai cấu xạ, $S'$ là một tập con đóng của $S$, và $X'$ (tương ứng, $Y'$) là một tập con đóng của $X$ (tương ứng, $Y$) sao cho $g(X')\subset S'$ (tương ứng, $h(Y')\subset S'$); đặt $Z=X\times_S Y$; giả sử $Z$ Noether địa phương, và đặt $Z'=p^{-1}(X')\cap q^{-1}(Y')$, trong đó $p$ và $q$ là các phép chiếu của $X\times_S Y$.
Trong các điều kiện này, đầy đủ hóa $Z_{/Z'}$ được đồng nhất với tích $(X_{/X'})\times_{S_{/S'}}(Y_{/Y'})$ của các $S_{/S'}$-tiền lược đồ hình thức, trong đó các cấu xạ cấu trúc được đồng nhất với $\widehat{g}$ và $\widehat{h}$, còn các phép chiếu được đồng nhất với $\widehat{p}$ và $\widehat{q}$.
\end{proposition}

\begin{proof}
Hiển nhiên bài toán có tính địa phương đối với $S$, $X$ và $Y$, nên ta quy về trường hợp $S=\Spec(A)$, $X=\Spec(B)$, $Y=\Spec(C)$, $S'=V(\mathfrak{J})$, $X'=V(\mathfrak{K})$ và $Y'=V(\mathfrak{L})$, trong đó $\mathfrak{J}$, $\mathfrak{K}$ và $\mathfrak{L}$ là các iđêan sao cho $\vphi(\mathfrak{J})\subset\mathfrak{K}$ và $\psi(\mathfrak{J})\subset\mathfrak{L}$, với $\vphi$ và $\psi$ lần lượt ký hiệu các đồng cấu $A\to B$ và $A\to C$ tương ứng với $g$ và $h$.
Khi đó ta biết rằng $Z=\Spec(B\otimes_A C)$ và $Z'=V(\mathfrak{M})$, trong đó $\mathfrak{M}$ là iđêan $\Im(\mathfrak{K}\otimes_A C)+\Im(B\otimes_A\mathfrak{L})$.
Kết luận suy ra từ \sref{I.10.7.2} và từ việc tích tenxơ đầy đủ $(\widehat{B}\otimes_{\widehat{A}}\widehat{C})^\wedge$ (trong đó $\widehat{A}$, $\widehat{B}$ và $\widehat{C}$ lần lượt là các đầy đủ hóa tách của $A$, $B$ và $C$ đối với các tôpô $\mathfrak{J}$-, $\mathfrak{K}$- và $\mathfrak{L}$-tiền adic) là đầy đủ hóa tách của tích tenxơ $B\otimes_A C$ đối với tôpô $\mathfrak{M}$-tiền adic \sref[0]{0.7.7.2}.
\end{proof}

Ngoài ra, lưu ý rằng nếu $T$ là một $S$-tiền lược đồ Noether địa phương, $u:T\to X$ và $v:T\to Y$ là hai $S$-cấu xạ, và $T'$ là một tập con đóng của $T$ sao cho $u(T')\subset X'$ và $v(T')\subset Y'$, thì mở rộng đến các đầy đủ hóa $((u,v)_S)^\wedge$ được đồng nhất với $(\widehat{u},\widehat{v})_{S_{/S'}}$.

\begin{corollary}[10.9.8]
\label{I.10.9.8}
Cho $X$ và $Y$ là hai $S$-tiền lược đồ Noether địa phương sao cho $X\times_S Y$ Noether địa phương; cho $S'$ là một tập con đóng của $S$, và $X'$ (tương ứng, $Y'$) là một tập con đóng của $X$ (tương ứng, $Y$) có ảnh trong $S$ nằm trong $S'$.
Với mọi $S$-cấu xạ $f:X\to Y$ sao cho $f(X')\subset Y'$, cấu xạ đồ thị $\Gamma_{\widehat f}$ được đồng nhất với mở rộng $(\Gamma_f)^\wedge$ của cấu xạ đồ thị của $f$.
\end{corollary}

\begin{corollary}[10.9.9]
\label{I.10.9.9}
Cho $X$ và $Y$ là hai tiền lược đồ Noether địa phương, $f:X\to Y$ là một cấu xạ, $Y'$ là một tập con đóng của $Y$, và $X'=f^{-1}(Y')$.
Khi đó tiền lược đồ hình thức $X_{/X'}$ được đồng nhất, qua sơ đồ giao hoán
\[
  \label{I.10.9.9.1}
  \xymatrix{
    X\ar[d]_f &
    X_{/X'}\ar[l]\ar[d]^{\widehat f}\\
    Y &
    Y_{/Y'}\ar[l]
  }
\]
với tích $X\times_Y(Y_{/Y'})$ của các tiền lược đồ hình thức.
\end{corollary}

\begin{proof}
Chỉ cần áp dụng \sref{I.10.9.7}, thay $S$ và $S'$ bằng $Y$, và thay $X$ và $X'$ bằng $X$.
\end{proof}

\begin{remark}[10.9.10]
\label{I.10.9.10}
\oldpage[I]{201}
Nếu $X$ là hợp rời $X_1\amalg X_2$ \hyperref[subsection:I.3.1]{(3.1)}, và $X'$ là hợp $X'_1\cup X'_2$, trong đó $X'_i$ là một tập con đóng của $X_i$ ($i=1,2$), thì ta thấy ngay rằng
\[
  X_{/X'}={X_1}_{/X'_1}\amalg {X_2}_{/X'_2}.
\]
\end{remark}

\subsection{Ứng dụng cho các bó coherent trên các lược đồ affine hình thức}
\label{subsection:I.10.10}

\begin{env}[10.10.1]
\label{I.10.10.1}
Trong toàn bộ mục này, $A$ sẽ ký hiệu một \emph{vành Noether adic}, còn $\mathfrak{J}$ là một iđêan định nghĩa của $A$.
Đặt $X=\Spec(A)$ và $\mathfrak{X}=\Spf(A)$; khi đó đối tượng thứ hai được đồng nhất với tập con đóng $V(\mathfrak{J})$ của $X$ \sref{I.10.1.2}.
Ngoài ra, các Định nghĩa \sref{I.10.1.2} và \sref{I.10.8.4} cho thấy \emph{lược đồ affine hình thức $\mathfrak{X}$} đồng nhất với \emph{đầy đủ hóa $X_{/\mathfrak{X}}$} của lược đồ affine $X$ dọc theo tập con đóng $\mathfrak{X}$ của không gian nền của nó.
Do đó, với mỗi $\sh{O}_X$-môđun coherent $\sh{F}$ tương ứng một $\sh{O}_\mathfrak{X}$-môđun hữu hạn sinh $\sh{F}_{/\mathfrak{X}}$; đây còn là một bó các môđun tôpô trên bó các vành tôpô $\sh{O}_\mathfrak{X}$.
Nhưng mọi $\sh{O}_X$-môđun coherent $\sh{F}$ đều có dạng $\widetilde{M}$, trong đó $M$ là một $A$-môđun hữu hạn sinh \sref{I.1.5.1}; ta đặt $(\widetilde{M})_{/\mathfrak{X}}=M^\Delta$.
Ngoài ra, nếu $u:M\to N$ là một $A$-đồng cấu giữa các $A$-môđun hữu hạn sinh, thì nó tương ứng với một đồng cấu $\widetilde{u}:\widetilde{M}\to\widetilde{N}$, và do đó cũng với một đồng cấu liên tục $\widetilde{u}_{/\mathfrak{X}}:(\widetilde{M})_{/\mathfrak{X}}\to(\widetilde{N})_{/\mathfrak{X}}$, mà ta ký hiệu bởi $u^\Delta$.
Hiển nhiên $(v\circ u)^\Delta=v^\Delta\circ u^\Delta$; như vậy ta đã xác định một \emph{hàm tử cộng tính hiệp biến $M^\Delta$} từ phạm trù các $A$-môđun hữu hạn sinh đến phạm trù các $\sh{O}_\mathfrak{X}$-môđun hữu hạn sinh.
Khi $A$ là một vành \emph{rời rạc}, ta có $M^\Delta=\widetilde{M}$.
\end{env}

\begin{proposition}[10.10.2]
\label{I.10.10.2}
\medskip\noindent
\begin{enumerate}
  \item[\rm{(i)}] $M^\Delta$ là một hàm tử khớp theo $M$, và tồn tại một đẳng cấu chính tắc có tính hàm tử của các $A$-môđun $\Gamma(\mathfrak{X},M^\Delta)\isoto M$.
  \item[\rm{(ii)}] Nếu $M$ và $N$ là hai $A$-môđun hữu hạn sinh, thì tồn tại các đẳng cấu chính tắc có tính hàm tử
    \[
    \label{I.10.10.2.1}
      (M\otimes_A N)^\Delta\isoto M^\Delta\otimes_{\sh{O}_\mathfrak{X}}N^\Delta,
      \tag{10.10.2.1}
    \]
    \[
    \label{I.10.10.2.2}
      (\Hom_A(M,N))^\Delta\isoto\shHom_{\sh{O}_\mathfrak{X}}(M^\Delta,N^\Delta).
      \tag{10.10.2.2}
    \]
  \item[\rm{(iii)}] Ánh xạ $u\mapsto u^\Delta$ là một đẳng cấu có tính hàm tử
    \[
    \label{I.10.10.2.3}
      \Hom_A(M,N)\isoto\Hom_{\sh{O}_\mathfrak{X}}(M^\Delta,N^\Delta).
      \tag{10.10.2.3}
    \]
\end{enumerate}
\end{proposition}

\begin{proof}
Tính khớp của $M^\Delta$ suy ra từ tính khớp của các hàm tử $\widetilde{M}$ \sref{I.1.3.5} và $\sh{F}_{/\mathfrak{X}}$ \sref{I.10.8.8}.
Theo định nghĩa, $\Gamma(\mathfrak{X},M^\Delta)$ là đầy đủ hóa tách của $A$-môđun $\Gamma(X,\widetilde{M})=M$ đối với tôpô $\mathfrak{J}$-tiền adic; nhưng, vì $A$ đầy đủ và $M$ hữu hạn sinh, ta biết \sref[0]{0.7.3.6} rằng $M$ tách và đầy đủ, điều này hoàn tất chứng minh (i).
Đẳng cấu \sref{I.10.10.2.1} (tương ứng, \sref{I.10.10.2.2}) thu được bằng cách hợp thành các đẳng cấu \sref{I.1.3.12}[i] và \sref{I.10.8.10.1} (tương ứng, \sref{I.1.3.12}[ii] và \sref{I.10.8.10.2}).
Cuối cùng, vì $\Hom_A(M,N)$ là một $A$-môđun hữu hạn sinh, ta có thể áp dụng (i) cho nó, qua đó đồng nhất $\Gamma(\mathfrak{X},(\Hom_A(M,N))^\Delta)$ với $\Hom_A(M,N)$, và sử dụng \sref{I.10.10.2.2}, điều này chứng minh đồng cấu \sref{I.10.10.2.3} là một đẳng cấu.
\end{proof}

Từ Mệnh đề~\sref{I.10.10.2}, ta suy ra cả một loạt hệ quả tương tự như các hệ quả suy ra từ Định lý~\sref{I.1.3.7} và Hệ quả~\sref{I.1.3.12}, mà ta dành cho độc giả việc phát biểu.

\oldpage[I]{202}
Lưu ý rằng tính khớp của $M^\Delta$, khi áp dụng cho dãy khớp $0\to\mathfrak{J}\to A\to A/\mathfrak{J}\to 0$, cho thấy bó các iđêan của $\sh{O}_\mathfrak{X}$ được ký hiệu ở đây bởi $\mathfrak{J}^\Delta$ trùng với bó trước đó đã được ký hiệu theo cùng cách trong \sref{I.10.3.1}, nhờ \sref{I.10.3.2}.

\begin{proposition}[10.10.3]
\label{I.10.10.3}
Dưới các giả thiết của \sref{I.10.10.1}, $\sh{O}_\mathfrak{X}$ là một bó vành coherent.
\end{proposition}

\begin{proof}
Nếu $f\in A$, ta biết $A_{\{f\}}$ là một vành adic Noether \sref[0]{0.7.6.11}; vì vấn đề mang tính địa phương, ta quy về \sref{I.10.1.4} việc chứng minh rằng hạt nhân của một đồng cấu $v:\sh{O}_\mathfrak{X}^n\to\sh{O}_\mathfrak{X}$ là một $\sh{O}_\mathfrak{X}$-môđun hữu hạn sinh.
Khi đó $v=u^\Delta$, trong đó $u$ là một $A$-đồng cấu $A^n\to A$ \sref{I.10.10.2}; vì $A$ là Noether, hạt nhân của $u$ hữu hạn sinh, hay tương đương, tồn tại một đồng cấu $A^m\xrightarrow{w}A^n$ sao cho dãy $A^m\xrightarrow{w}A^n\xrightarrow{u}A$ khớp.
Suy ra từ \sref{I.10.10.2} rằng dãy $\sh{O}_\mathfrak{X}^m\xrightarrow{w^\Delta}\sh{O}_\mathfrak{X}^n\xrightarrow{v}\sh{O}_\mathfrak{X}$ khớp, điều này chứng minh hạt nhân của $v$ hữu hạn sinh.
\end{proof}

\begin{env}[10.10.4]
\label{I.10.10.4}
Với các ký hiệu trên, đặt $A_n=A/\mathfrak{J}^{n+1}$, và cho $X_n$ là lược đồ affine $\Spec(A_n)=(\mathfrak{X},\sh{O}_\mathfrak{X}/\sh{J}^{n+1})$, trong đó $\sh{J}=\mathfrak{J}^\Delta$ là bó các iđêan định nghĩa của $\sh{O}_\mathfrak{X}$ tương ứng với iđêan $\mathfrak{J}$.
Cho $u_{mn}$ là cấu xạ tiền lược đồ $X_m\to X_n$ tương ứng với đồng cấu chính tắc $A_n\to A_m$ khi $m\leq n$; lược đồ hình thức $\mathfrak{X}$ là giới hạn quy nạp của các $X_n$ theo các $u_{mn}$ \sref{I.10.6.3}.
\end{env}

\begin{proposition}[10.10.5]
\label{I.10.10.5}
Dưới các giả thiết của \sref{I.10.10.1}, cho $\sh{F}$ là một $\sh{O}_\mathfrak{X}$-môđun.
Các điều kiện sau là tương đương:
\begin{enumerate}
  \item[\rm{(a)}] $\sh{F}$ là một $\sh{O}_\mathfrak{X}$-môđun coherent;
  \item[\rm{(b)}] $\sh{F}$ đẳng cấu với giới hạn xạ ảnh \sref{I.10.6.6} của một dãy $(\sh{F}_n)$ các $\sh{O}_{X_n}$-môđun coherent sao cho $u_{mn}^*(\sh{F}_n)=\sh{F}_m$; và
  \item[\rm{(c)}] tồn tại một $A$-môđun hữu hạn sinh $M$ (xác định sai khác một đẳng cấu chính tắc bởi Mệnh đề~\sref{I.10.10.2}[i]) sao cho $\sh{F}$ đẳng cấu với $M^\Delta$.
\end{enumerate}
\end{proposition}

\begin{proof}
Trước hết, ta chứng minh (b) kéo theo (c).
Ta có $\sh{F}_n=\widetilde{M_n}$, trong đó $M_n$ là một $A_n$-môđun hữu hạn sinh, và giả thiết kéo theo $M_m=M_n\otimes_{A_n}A_m$ khi $m\leq n$ \sref{I.1.6.5}; do đó các $M_n$ tạo thành một hệ xạ ảnh đối với các đi-đồng cấu chính tắc $M_n\to M_m$ ($m\leq n$), và từ định nghĩa của các $A_n$ suy ra ngay rằng hệ xạ ảnh này thỏa mãn các điều kiện của \sref[0]{0.7.2.9}; vì thế, giới hạn xạ ảnh $M$ của nó là một $A$-môđun hữu hạn sinh sao cho $M_n=M\otimes_A A_n$ với mọi $n$.
Suy ra $\sh{F}_n$ được cảm sinh trên $X_n$ bởi $\widetilde{M}\otimes_{\sh{O}_X}(\sh{O}_X/\widetilde{\mathfrak{J}}^{n+1})$, và do đó $\sh{F}=M^\Delta$ theo định nghĩa \sref{I.10.8.4}.
\end{proof}

Ngược lại, (c) kéo theo (b); thật vậy, nếu $u_n$ là cấu xạ nhúng $X_n\to X$, thì $u_n^*(\widetilde{M})=(M\otimes_A A_n)\supertilde$ được cảm sinh trên $X_n$ bởi $\widetilde{M}\otimes_{\sh{O}_X}(\sh{O}_X/\widetilde{\mathfrak{J}}^{n+1})$, và $M^\Delta=\varprojlim u_n^*(\widetilde{M})$ theo định nghĩa \sref{I.10.8.4}; vì $u_m=u_n\circ u_{mn}$ khi $m\leq n$, các $\sh{F}_n=u_n^*(\widetilde{M})$ thỏa mãn các điều kiện của (b), từ đó suy ra khẳng định.

Bây giờ ta chứng minh (c) kéo theo (a): thật vậy, theo định nghĩa ta có $\sh{O}_\mathfrak{X}=A^\Delta$; vì $M$ là cokernel của một đồng cấu $A^m\to A^n$, từ Mệnh đề~\sref{I.10.10.2} suy ra $M^\Delta$ là cokernel của một đồng cấu $\sh{O}_\mathfrak{X}^m\to\sh{O}_\mathfrak{X}^n$; và vì bó vành $\sh{O}_\mathfrak{X}$ coherent \sref{I.10.10.3}, $M^\Delta$ cũng coherent \sref[0]{0.5.3.4}.
\oldpage[I]{203}

Cuối cùng, (a) kéo theo (b).
Xét như một $\sh{O}_\mathfrak{X}$-môđun, ta có $\sh{O}_{X_n}=\sh{O}_\mathfrak{X}/\sh{J}^{n+1}=A_n^\Delta$; khi ấy $\sh{F}_n=\sh{F}\otimes_{\sh{O}_\mathfrak{X}}\sh{O}_{X_n}$ là một $\sh{O}_\mathfrak{X}$-môđun coherent \sref[0]{0.5.3.5}, và vì nó cũng là một $\sh{O}_{X_n}$-môđun còn $\sh{J}^{n+1}$ coherent, ta suy ra $\sh{F}_n$ là một $\sh{O}_{X_n}$-môđun coherent \sref[0]{0.5.3.10}; hơn nữa, hiển nhiên $u_{mn}^*(\sh{F}_n)=\sh{F}_m$ khi $m\leq n$ (nhớ rằng ánh xạ liên tục $X_m\to X_n$ của các không gian nền là ánh xạ đồng nhất trên $\mathfrak{X}$).
Do đó, bó $\sh{G}=\varprojlim\sh{F}_n$ là một $\sh{O}_\mathfrak{X}$-môđun coherent, vì ta đã thấy (b) kéo theo (a). Các đồng cấu chính tắc $\sh{F}\to\sh{F}_n$ tạo thành một hệ xạ ảnh; khi chuyển sang giới hạn, hệ này cho một đồng cấu chính tắc $w:\sh{F}\to\sh{G}$, và chỉ còn phải chứng minh $w$ là song ánh.
Vì vấn đề bây giờ mang tính \emph{địa phương}, ta có thể giới hạn trong trường hợp $\sh{F}$ là cokernel của một đồng cấu $\sh{O}_\mathfrak{X}^p\to\sh{O}_\mathfrak{X}^q$; vì đồng cấu này có dạng $v^\Delta$, trong đó $v$ là một đồng cấu $A^p\to A^q$\footnote{\emph{[Dịch.] Nguồn tiếng Pháp in $A^m\to A^n$ ở đây; các hạng của đồng cấu ngay trước đó đòi hỏi $A^p\to A^q$.}}, $\sh{F}$ đẳng cấu với $M^\Delta$, trong đó $M=\Coker{v}$ \sref{I.10.10.2}.
Khi đó, theo Mệnh đề~\sref{I.10.10.2}, ta có $\sh{F}_n=M^\Delta\otimes_{\sh{O}_\mathfrak{X}}A_n^\Delta=(M\otimes_A A_n)^\Delta$; và vì tôpô $\mathfrak{J}$-adic trên $M\otimes_A A_n$ là rời rạc, ta có $(M\otimes_A A_n)^\Delta=(M\otimes_A A_n)\supertilde$ (như một $\sh{O}_{X_n}$-môđun); ở trên ta đã thấy $M^\Delta=\varprojlim\sh{F}_n$, nên trong trường hợp này $w$ chính là ánh xạ đồng nhất.

\begin{corollary}[10.10.6]
\label{I.10.10.6}
Nếu $\sh{F}$ thỏa mãn điều kiện \emph{(b)} của Mệnh đề~\sref{I.10.10.5}, thì hệ xạ ảnh $(\sh{F}_n)$ đẳng cấu với hệ gồm các $\sh{F}\otimes_{\sh{O}_\mathfrak{X}}\sh{O}_{X_n}$.
\end{corollary}

\begin{env}[10.10.7]
\label{I.10.10.7}
Bây giờ cho $A$ và $B$ là hai vành Noether adic, và $\vphi:B\to A$ là một đồng cấu liên tục; ta ký hiệu bởi $\mathfrak{J}$ (tương ứng, $\mathfrak{K}$) một iđêan định nghĩa của $A$ (tương ứng, của $B$) sao cho $\vphi(\mathfrak{K})\subset\mathfrak{J}$, và đặt $X=\Spec(A)$, $Y=\Spec(B)$, $\mathfrak{X}=\Spf(A)$, $\mathfrak{Y}=\Spf(B)$.
Cho $f:X\to Y$ là cấu xạ tiền lược đồ tương ứng với $\vphi$ \sref{I.1.6.1}, và $\widehat{f}:\mathfrak{X}\to\mathfrak{Y}$ là mở rộng của nó đến các đầy đủ hóa \sref{I.10.9.1}, đồng thời cũng là cấu xạ tiền lược đồ hình thức tương ứng với $\vphi$ \sref{I.10.2.2}.
\end{env}

\begin{proposition}[10.10.8]
\label{I.10.10.8}
Với mọi $B$-môđun $N$ hữu hạn sinh, tồn tại một đẳng cấu $\sh{O}_\mathfrak{X}$-môđun chính tắc và có tính hàm tử
\[
  \widehat{f}^*(N^\Delta)\isoto(N\otimes_B A)^\Delta.
\]
\end{proposition}

\begin{proof}
Ký hiệu bởi $i_X:\mathfrak{X}\to X$ và $i_Y:\mathfrak{Y}\to Y$ các cấu xạ chính tắc, theo \sref{I.10.8.8} ta có, xét đến các đẳng cấu chính tắc có tính hàm tử, $N^\Delta=i_Y^*(\widetilde{N})$ và
\[
  (N\otimes_B A)^\Delta=i_X^*((N\otimes_B A)\supertilde)=i_X^*(f^*(\widetilde{N}))
\]
\sref{I.1.6.5}; do đó mệnh đề suy ra từ tính giao hoán của biểu đồ trong \sref{I.10.9.2}.
\end{proof}

\begin{corollary}[10.10.9]
\label{I.10.10.9}
Với mọi iđêan $\mathfrak{b}$ của $B$, ta có $\widehat{f}^*(\mathfrak{b}^\Delta)\sh{O}_\mathfrak{X}=(\mathfrak{b}A)^\Delta$.
\end{corollary}

\begin{proof}
Cho $j$ là phép nhúng chính tắc $\mathfrak{b}\to B$, tương ứng với nó là phép nhúng chính tắc $j^\Delta:\mathfrak{b}^\Delta\to\sh{O}_\mathfrak{Y}$ của các bó $\sh{O}_\mathfrak{Y}$-môđun; theo định nghĩa, $\widehat{f}^*(\mathfrak{b}^\Delta)\sh{O}_\mathfrak{X}$ là ảnh của đồng cấu $\widehat{f}^*(j^\Delta):\widehat{f}^*(\mathfrak{b}^\Delta)\to\sh{O}_\mathfrak{X}=\widehat{f}^*(\sh{O}_\mathfrak{Y})$; nhưng đồng cấu này đồng nhất với $(j\otimes 1)^\Delta:(\mathfrak{b}\otimes_B A)^\Delta\to\sh{O}_\mathfrak{X}=(B\otimes_B A)^\Delta$ theo Mệnh đề~\sref{I.10.10.8}.
Vì ảnh của $j\otimes 1$ là iđêan $\mathfrak{b}A$ của $A$, ảnh của $(j\otimes 1)^\Delta$ do đó là $(\mathfrak{b}A)^\Delta$ theo Mệnh đề~\sref{I.10.10.2}, suy ra kết luận.
\end{proof}

\subsection{Các bó coherent trên các tiền lược đồ hình thức}
\label{subsection:I.10.11}

\begin{proposition}[10.11.1]
\label{I.10.11.1}
\oldpage[I]{204}
Nếu $\mathfrak{X}$ là một tiền lược đồ hình thức Noether địa phương, thì bó các vành $\sh{O}_\mathfrak{X}$ là coherent, và mọi bó các iđêan định nghĩa của $\mathfrak{X}$ đều coherent.
\end{proposition}

\begin{proof}
Bài toán có tính địa phương, nên ta quy về trường hợp một lược đồ hình thức affine Noether, và khi đó mệnh đề suy ra từ các Mệnh đề~\sref{I.10.10.3} và~\sref{I.10.10.5}.
\end{proof}

\begin{env}[10.11.2]
\label{I.10.11.2}
Cho $\mathfrak{X}$ là một tiền lược đồ hình thức Noether địa phương, $\sh{J}$ là một bó các iđêan định nghĩa của $\mathfrak{X}$, và $X_n$ là tiền lược đồ (thông thường) Noether địa phương $(\mathfrak{X},\sh{O}_\mathfrak{X}/\sh{J}^{n+1})$, sao cho $\mathfrak{X}$ là \emph{giới hạn quy nạp} của dãy $(X_n)$ theo các cấu xạ chính tắc $u_{mn}:X_m\to X_n$ \sref{I.10.6.3}.
Với các ký hiệu này:
\end{env}

\begin{theorem}[10.11.3]
\label{I.10.11.3}
Để một $\sh{O}_\mathfrak{X}$-môđun $\sh{F}$ là coherent, điều kiện cần và đủ là nó đẳng cấu với giới hạn xạ ảnh của một dãy $(\sh{F}_n)$, trong đó các $\sh{F}_n$ là những $\sh{O}_{X_n}$-môđun coherent sao cho $u_{mn}^*(\sh{F}_n)=\sh{F}_m$ với $m\leq n$ \sref{I.10.6.6}.
Khi đó hệ xạ ảnh $(\sh{F}_n)$ đẳng cấu với hệ gồm các $u_n^*(\sh{F})=\sh{F}\otimes_{\sh{O}_\mathfrak{X}}\sh{O}_{X_n}$, trong đó $u_n$ là cấu xạ chính tắc $X_n\to\mathfrak{X}$.
\end{theorem}

\begin{proof}
Bài toán có tính địa phương, nên ta quy về trường hợp $\mathfrak{X}$ là một lược đồ hình thức affine Noether, và khi đó định lý là hệ quả của Mệnh đề~\sref{I.10.10.5} và Hệ quả~\sref{I.10.10.6}.
\end{proof}

Do đó có thể nói rằng \emph{việc cho một $\sh{O}_\mathfrak{X}$-môđun coherent tương đương với việc cho một hệ xạ ảnh $(\sh{F}_n)$ các $\sh{O}_{X_n}$-môđun coherent sao cho $u_{mn}^*(\sh{F}_n)=\sh{F}_m$ với $m\leq n$}.

\begin{corollary}[10.11.4]
\label{I.10.11.4}
Nếu $\sh{F}$ và $\sh{G}$ là hai $\sh{O}_\mathfrak{X}$-môđun coherent, thì (với các ký hiệu của Định lý~\sref{I.10.11.3}) ta có thể xác định một đẳng cấu chính tắc có tính hàm tử
\[
\label{I.10.11.4.1}
  \Hom_{\sh{O}_\mathfrak{X}}(\sh{F},\sh{G})\isoto\varprojlim_n\Hom_{\sh{O}_{X_n}}(\sh{F}_n,\sh{G}_n).
  \tag{10.11.4.1}
\]
\end{corollary}

\begin{proof}
Giới hạn xạ ảnh ở vế phải được hiểu theo các ánh xạ $\theta_n\mapsto u_{mn}^*(\theta_n)$ ($m\leq n$) từ $\Hom_{\sh{O}_{X_n}}(\sh{F}_n,\sh{G}_n)$ vào $\Hom_{\sh{O}_{X_m}}(\sh{F}_m,\sh{G}_m)$.
Đồng cấu \sref{I.10.11.4.1} gửi một phần tử $\theta\in\Hom_{\sh{O}_\mathfrak{X}}(\sh{F},\sh{G})$ tới dãy $(u_n^*(\theta))$; ta xác định đồng cấu nghịch đảo của nó bằng cách gửi hệ xạ ảnh $(\theta_n)\in\varprojlim_n\Hom_{\sh{O}_{X_n}}(\sh{F}_n,\sh{G}_n)$ tới giới hạn xạ ảnh của hệ này trong $\Hom_{\sh{O}_\mathfrak{X}}(\sh{F},\sh{G})$, có xét đến Định lý~\sref{I.10.11.3}.
\end{proof}

\begin{corollary}[10.11.5]
\label{I.10.11.5}
Để một đồng cấu $\theta:\sh{F}\to\sh{G}$ là toàn ánh, điều kiện cần và đủ là đồng cấu tương ứng $\theta_0=u_0^*(\theta):\sh{F}_0\to\sh{G}_0$ toàn ánh.
\end{corollary}

\begin{proof}
Bài toán có tính địa phương, nên ta quy về trường hợp $\mathfrak{X}=\Spf(A)$, với $A$ là một vành Noether adic, $\sh{F}=M^\Delta$, $\sh{G}=N^\Delta$ và $\theta=u^\Delta$, trong đó $M$ và $N$ là các $A$-môđun hữu hạn sinh và $u$ là một đồng cấu $M\to N$; hơn nữa ta có $\theta_0=\widetilde{u_0}$, trong đó $u_0$ là đồng cấu $u\otimes 1:M\otimes_A A/\mathfrak{J}\to N\otimes_A A/\mathfrak{J}$; kết luận suy ra từ việc $\theta$ và $u$ (tương ứng, $\theta_0$ và $u_0$) đồng thời toàn ánh (\sref{I.1.3.9} và \sref{I.10.10.2}) và từ việc $u$ và $u_0$ đồng thời toàn ánh \sref[0]{0.7.1.14}.
\end{proof}

\begin{env}[10.11.6]
\label{I.10.11.6}
Định lý~\sref{I.10.11.3} cho thấy ta có thể xem mọi $\sh{O}_\mathfrak{X}$-môđun coherent $\sh{F}$ như một \emph{$\sh{O}_\mathfrak{X}$-môđun tôpô}, bằng cách xem nó là giới hạn xạ ảnh của các bó nhóm \emph{giả rời rạc} $\sh{F}_n$ \sref[0]{0.3.8.1}.
Khi đó, từ Hệ quả~\sref{I.10.11.4} suy ra mọi đồng cấu $u:\sh{F}\to\sh{G}$ của các $\sh{O}_\mathfrak{X}$-môđun coherent đều tự động \emph{liên tục}
\oldpage[I]{205}
\sref[0]{0.3.8.2}.
Hơn nữa, nếu $\sh{H}$ là một $\sh{O}_\mathfrak{X}$-môđun con coherent của một $\sh{O}_\mathfrak{X}$-môđun coherent $\sh{F}$, thì với mọi tập mở $U\subset\mathfrak{X}$, $\Gamma(U,\sh{H})$ là một nhóm con \emph{đóng} của nhóm tôpô $\Gamma(U,\sh{F})$, bởi vì hàm tử $\Gamma$ là khớp trái và $\Gamma(U,\sh{H})$ là hạt nhân của đồng cấu $\Gamma(U,\sh{F})\to\Gamma(U,\sh{F}/\sh{H})$, đồng cấu này \emph{liên tục} theo điều vừa nói vì $\sh{F}/\sh{H}$ là coherent \sref[0]{0.5.3.4}; khẳng định suy ra từ việc $\Gamma(U,\sh{F}/\sh{H})$ là một nhóm tôpô tách.
\end{env}

\begin{proposition}[10.11.7]
\label{I.10.11.7}
Cho $\sh{F}$ và $\sh{G}$ là hai $\sh{O}_\mathfrak{X}$-môđun coherent.
Ta có thể xác định (với các ký hiệu của Định lý~\sref{I.10.11.3}) các đẳng cấu chính tắc có tính hàm tử của các $\sh{O}_\mathfrak{X}$-môđun tôpô \sref{I.10.11.6}
\[
\label{I.10.11.7.1}
  \sh{F}\otimes_{\sh{O}_\mathfrak{X}}\sh{G}\isoto\varprojlim_n(\sh{F}_n\otimes_{\sh{O}_{X_n}}\sh{G}_n),
  \tag{10.11.7.1}
\]
\[
\label{I.10.11.7.2}
  \shHom_{\sh{O}_\mathfrak{X}}(\sh{F},\sh{G})\isoto\varprojlim_n\shHom_{\sh{O}_{X_n}}(\sh{F}_n,\sh{G}_n).
  \tag{10.11.7.2}
\]
\end{proposition}

\begin{proof}
Sự tồn tại của đẳng cấu \sref{I.10.11.7.1} suy ra từ công thức
\[
  \sh{F}_n\otimes_{\sh{O}_{X_n}}\sh{G}_n=(\sh{F}\otimes_{\sh{O}_\mathfrak{X}}\sh{O}_{X_n})\otimes_{\sh{O}_{X_n}}(\sh{G}\otimes_{\sh{O}_\mathfrak{X}}\sh{O}_{X_n})=(\sh{F}\otimes_{\sh{O}_\mathfrak{X}}\sh{G})\otimes_{\sh{O}_\mathfrak{X}}\sh{O}_{X_n}
\]
và từ Định lý~\sref{I.10.11.3}.
Đẳng cấu \sref{I.10.11.7.2}, trong đó hai vế được xem như các bó môđun không có tôpô, suy ra từ định nghĩa các tiết diện của $\shHom_{\sh{O}_\mathfrak{X}}(\sh{F},\sh{G})$ và $\shHom_{\sh{O}_{X_n}}(\sh{F}_n,\sh{G}_n)$, cùng với sự tồn tại của đẳng cấu \sref{I.10.11.4.1}, áp dụng cho tiền lược đồ cảm sinh trên một tập mở affine hình thức Noether tùy ý của $\mathfrak{X}$.
Còn phải chứng minh rằng đẳng cấu \sref{I.10.11.7.2} liên tục theo cả hai chiều trên một tập quasi-compact; do đó ta quy về trường hợp $\mathfrak{X}=\Spf(A)$ với $A$ là một vành Noether adic, và rồi, theo \sref{I.10.10.5}, về trường hợp $\sh{F}=M^\Delta$ và $\sh{G}=N^\Delta$, trong đó $M$ và $N$ là hai $A$-môđun hữu hạn sinh; xét đến \sref{I.10.10.2.1}, \sref{I.10.10.2.3} và Hệ quả~\sref{I.1.3.12}[ii], ta quy về việc chứng minh rằng đẳng cấu chính tắc $\Hom_A(M,N)\isoto\varprojlim_n\Hom_{A_n}(M_n,N_n)$ (với $M_n=M\otimes_A A_n$ và $N_n=N\otimes_A A_n$) là liên tục, điều này đã được chứng minh trong \sref[0]{0.7.8.2}.
\end{proof}

\begin{env}[10.11.8]
\label{I.10.11.8}
Vì $\Hom_{\sh{O}_\mathfrak{X}}(\sh{F},\sh{G})$ là nhóm các tiết diện của bó nhóm tôpô $\shHom_{\sh{O}_\mathfrak{X}}(\sh{F},\sh{G})$, nó được trang bị một tôpô nhóm.
Nếu $\mathfrak{X}$ là \emph{Noether}, thì từ \sref{I.10.11.7.2} suy ra rằng các nhóm con $\Hom_{\sh{O}_\mathfrak{X}}(\sh{F},\sh{J}^n\sh{G})$ (với $n$ tùy ý) tạo thành một hệ cơ bản các lân cận của $0$ trong nhóm này.
\end{env}

\begin{proposition}[10.11.9]
\label{I.10.11.9}
Cho $\mathfrak{X}$ là một tiền lược đồ hình thức Noether, và $\sh{F}$, $\sh{G}$ là hai $\sh{O}_\mathfrak{X}$-môđun coherent.
Trong nhóm tôpô $\Hom_{\sh{O}_\mathfrak{X}}(\sh{F},\sh{G})$, các đồng cấu toàn ánh (tương ứng, đơn ánh, song ánh) tạo thành một tập mở.
\end{proposition}

\begin{proof}
Theo Hệ quả~\sref{I.10.11.5}, tập các đồng cấu toàn ánh trong $\Hom_{\sh{O}_\mathfrak{X}}(\sh{F},\sh{G})$ là ảnh ngược, qua ánh xạ liên tục $\Hom_{\sh{O}_\mathfrak{X}}(\sh{F},\sh{G})\to\Hom_{\sh{O}_{X_0}}(\sh{F}_0,\sh{G}_0)$, của một tập con của nhóm rời rạc $\Hom_{\sh{O}_{X_0}}(\sh{F}_0,\sh{G}_0)$, suy ra khẳng định thứ nhất.
Để chứng minh khẳng định thứ hai, ta phủ $\mathfrak{X}$ bởi một số hữu hạn các tập mở affine hình thức Noether $U_i$.
Để $\theta\in\Hom_{\sh{O}_\mathfrak{X}}(\sh{F},\sh{G})$ là đơn ánh, điều kiện cần và đủ là mọi ảnh của nó qua các ánh xạ hạn chế (liên tục) $\Hom_{\sh{O}_\mathfrak{X}}(\sh{F},\sh{G})\to\Hom_{\sh{O}_\mathfrak{X}|U_i}(\sh{F}|U_i,\sh{G}|U_i)$ đều đơn ánh; do đó ta quy về trường hợp affine, và điều này đã được chứng minh trong \sref[0]{0.7.8.3}.
\end{proof}

\subsection{Các cấu xạ adic của các tiền lược đồ hình thức}
\label{subsection:I.10.12}

\begin{env}[10.12.1]
\label{I.10.12.1}
\oldpage[I]{206}
Cho $\mathfrak{X}$ và $\mathfrak{S}$ là hai tiền lược đồ hình thức \emph{Noether địa phương};
ta nói một cấu xạ $f:\mathfrak{X}\to\mathfrak{S}$ là \emph{adic} nếu tồn tại một iđêan định nghĩa $\sh{J}$ của $\mathfrak{S}$ sao cho $\sh{K}=f^*(\sh{J})\sh{O}_\mathfrak{X}$ là một iđêan định nghĩa của $\mathfrak{X}$;
khi đó ta cũng nói $\mathfrak{X}$ là một \emph{$\mathfrak{S}$-tiền lược đồ adic} (đối với $f$).
Khi điều này xảy ra, với \emph{mọi} iđêan định nghĩa $\sh{J}_1$ của $\mathfrak{S}$, $\sh{K}_1=f^*(\sh{J}_1)\sh{O}_\mathfrak{X}$ là một iđêan định nghĩa của $\mathfrak{X}$.
Thật vậy, vì vấn đề là địa phương, ta có thể giả sử $\mathfrak{X}$ và $\mathfrak{S}$ là affine và noetherian;
khi đó tồn tại một số nguyên $n$ sao cho $\sh{J}^n\subset\sh{J}_1$ và $\sh{J}_1^n\subset\sh{J}$ (\sref{I.10.3.6} và \sref[0]{0.7.1.4}), do đó $\sh{K}^n\subset\sh{K}_1$ và $\sh{K}_1^n\subset\sh{K}$.
Quan hệ thứ nhất chứng tỏ rằng $\sh{K}_1=\mathfrak{K}_1^\Delta$, trong đó $\mathfrak{K}_1$ là một iđêan mở của $A=\Gamma(\mathfrak{X},\sh{O}_\mathfrak{X})$, còn quan hệ thứ hai chứng tỏ rằng $\mathfrak{K}_1$ là một iđêan định nghĩa của $A$ \sref[0]{0.7.1.4}; từ đó suy ra khẳng định của ta.
\end{env}

Từ những điều trên suy ra ngay rằng, nếu $\mathfrak{X}$ và $\mathfrak{Y}$ là hai $\mathfrak{S}$-tiền lược đồ adic, thì \emph{mọi $\mathfrak{S}$-cấu xạ $u:\mathfrak{X}\to\mathfrak{Y}$ đều là adic}:
thật vậy, nếu $f:\mathfrak{X}\to\mathfrak{S}$ và $g:\mathfrak{Y}\to\mathfrak{S}$ là các cấu xạ cấu trúc, và $\sh{J}$ là một iđêan định nghĩa của $\mathfrak{S}$, thì $f=g\circ u$, nên $u^*(g^*(\sh{J})\sh{O}_\mathfrak{Y})\sh{O}_\mathfrak{X}=f^*(\sh{J})\sh{O}_\mathfrak{X}$ là một iđêan định nghĩa của $\mathfrak{X}$ và, theo giả thiết, $g^*(\sh{J})\sh{O}_\mathfrak{Y}$ là một iđêan định nghĩa của $\mathfrak{Y}$.

\begin{env}[10.12.2]
\label{I.10.12.2}
Trong phần tiếp theo, ta cố định một tiền lược đồ hình thức Noether địa phương $\mathfrak{S}$ và một iđêan định nghĩa $\sh{J}$ của $\mathfrak{S}$;
ta đặt $S_n=(\mathfrak{S},\sh{O}_\mathfrak{S}/\sh{J}^{n+1})$.
Các $\mathfrak{S}$-tiền lược đồ adic (Noether địa phương) hiển nhiên tạo thành một \emph{phạm trù}.
Ta gọi một hệ quy nạp $(X_n)$ gồm các $S_n$-tiền lược đồ (thông thường) Noether địa phương là một \emph{$(S_n)$-hệ quy nạp adic} nếu các cấu xạ cấu trúc $f_n:X_n\to S_n$ thỏa mãn rằng, với $m\leq n$, các sơ đồ
\[
\label{I.10.12.2.1}
  \xymatrix{
    X_n \ar[d]_{f_n}
    & X_m \ar[l] \ar[d]^{f_m}\\
    S_n
    & S_m \ar[l]
  }
  \tag{10.12.2.1}
\]
giao hoán và \emph{đồng nhất hóa $X_m$ với tích $X_n\times_{S_n}S_m=(X_n)_{(S_m)}$}.
Các hệ quy nạp adic tạo thành một \emph{phạm trù}:
thật vậy, chỉ cần định nghĩa một cấu xạ $(X_n)\to(Y_n)$ giữa các hệ như vậy là một \emph{hệ quy nạp các $S_n$-cấu xạ $u_n:X_n\to Y_n$} sao cho $u_m$ được đồng nhất với $(u_n)_{(S_m)}$ khi $m\leq n$.
Với quy ước này:
\end{env}

\begin{theorem}[10.12.3]
\label{I.10.12.3}
Có một tương đương chính tắc giữa phạm trù các $\mathfrak{S}$-tiền lược đồ adic và phạm trù các $(S_n)$-hệ quy nạp adic.
\end{theorem}

Tương đương đang xét được thiết lập như sau:
nếu $\mathfrak{X}$ là một $\mathfrak{S}$-tiền lược đồ adic và $f:\mathfrak{X}\to\mathfrak{S}$ là cấu xạ cấu trúc, thì $\sh{K}=f^*(\sh{J})\sh{O}_\mathfrak{X}$ là một iđêan định nghĩa của $\mathfrak{X}$, và ta cho $\mathfrak{X}$ tương ứng với hệ quy nạp gồm các $X_n=(\mathfrak{X},\sh{O}_\mathfrak{X}/\sh{K}^{n+1})$, trong đó cấu xạ cấu trúc $f_n:X_n\to S_n$ tương ứng với $f$ \sref{I.10.5.6}.
Trước hết, ta chứng minh rằng $(X_n)$ là một \emph{hệ quy nạp adic}:
nếu $f=(\psi,\theta)$, thì $\psi^*(\sh{J})\sh{O}_\mathfrak{X}=\sh{K}$, do đó $\psi^*(\sh{J}^n)\sh{O}_\mathfrak{X}=\sh{K}^n$ với mọi $n$, và (do tính khớp của hàm tử $\psi^*$) $\sh{K}^{m+1}/\sh{K}^{n+1}=\psi^*(\sh{J}^{m+1}/\sh{J}^{n+1})(\sh{O}_\mathfrak{X}/\sh{K}^{n+1})$ khi $m\leq n$;
vì thế kết luận của ta suy ra từ \sref{I.4.4.5}.
Hơn nữa, ta kiểm tra ngay rằng một $\mathfrak{S}$-cấu xạ $u:\mathfrak{X}\to\mathfrak{Y}$ giữa các $\mathfrak{S}$-tiền lược đồ adic tương ứng (với các ký hiệu hiển nhiên) với
\oldpage[I]{207}
một hệ quy nạp các $S_n$-cấu xạ $u_n:X_n\to Y_n$ sao cho $u_m$ được đồng nhất với $(u_n)_{(S_m)}$ khi $m\leq n$.

Việc ta thực sự đã xác định một tương đương theo cách này sẽ suy ra từ mệnh đề chính xác hơn sau đây:

\begin{proposition}[10.12.3.1]
\label{I.10.12.3.1}
Cho $(X_n)$ là một hệ quy nạp các $S_n$-tiền lược đồ;
giả sử các cấu xạ cấu trúc $f_n:X_n\to S_n$ có tính chất là các sơ đồ trong \sref{I.10.12.2.1} giao hoán và đồng nhất hóa $X_m$ với $X_n\times_{S_n}S_m$ khi $m\leq n$.
Khi đó hệ quy nạp $(X_n)$ thỏa mãn các điều kiện \emph{(b)} và \emph{(c)} của \sref{I.10.6.3};
gọi $\mathfrak{X}$ là giới hạn quy nạp của nó và $f:\mathfrak{X}\to\mathfrak{S}$ là cấu xạ giới hạn quy nạp của hệ quy nạp $(f_n)$.
Khi đó, nếu $X_0$ là Noether địa phương, thì $\mathfrak{X}$ là Noether địa phương và $f$ là một cấu xạ adic.
\end{proposition}

\begin{proof}
Vì bó iđêan của $\sh{O}_{S_n}$ xác định tiền lược đồ con $S_m$ của $S_n$ là lũy linh, nên theo \sref{I.4.4.5}, bó iđêan của $\sh{O}_{X_n}$ xác định tiền lược đồ con $X_m$ của $X_n$ cũng lũy linh; do đó các điều kiện của \sref{I.10.6.3} được thỏa mãn.
Vấn đề vì thế là địa phương trên $\mathfrak{X}$ và $\mathfrak{S}$, và ta có thể giả sử $\mathfrak{S}=\Spf(A)$, $\sh{J}=\mathfrak{J}^\Delta$ (với $A$ là một vành Noether $\mathfrak{J}$-adic) và $X_n=\Spec(B_n)$;
nếu $A_n=A/\mathfrak{J}^{n+1}$, thì giả thiết kéo theo $B_0$ là Noether và, nếu đặt $\mathfrak{J}_n=\mathfrak{J}/\mathfrak{J}^{n+1}$, thì $B_m=B_n/\mathfrak{J}_n^{m+1}B_n$.
Do đó hạt nhân của $B_n\to B_0$ là $\mathfrak{K}_n=\mathfrak{J}_nB_n$, và hạt nhân của $B_n\to B_m$ là $\mathfrak{K}_n^{m+1}$ khi $m\leq n$;
hơn nữa, vì $A_1$ là Noether nên $\mathfrak{J}_1$ có kiểu hữu hạn trên $A_1$, do đó $\mathfrak{K}_1/\mathfrak{K}_1^2$ có kiểu hữu hạn trên $B_1$ và, \emph{a fortiori}, có kiểu hữu hạn trên $B_0=B_1/\mathfrak{K}_1$;
khi đó việc $\mathfrak{X}$ là Noether suy ra từ \sref{I.10.6.4};
nếu $B=\varprojlim B_n$, thì $\mathfrak{X}=\Spf(B)$ và, nếu $\mathfrak{K}$ là hạt nhân của $B\to B_0$, thì $B_n=B/\mathfrak{K}^{n+1}$.
Nếu $\rho_n:A/\mathfrak{J}^{n+1}\to B/\mathfrak{K}^{n+1}$ là đồng cấu tương ứng với $f_n$, thì
\[
    \mathfrak{K}/\mathfrak{K}^{n+1} = (B/\mathfrak{K}^{n+1})\rho_n(\mathfrak{J}/\mathfrak{J}^{n+1})
\]
Vì đồng cấu $\rho:A\to B$ tương ứng với $f$ bằng $\varprojlim\rho_n$, nên iđêan $\mathfrak{J}B$ của $B$ trù mật trong $\mathfrak{K}$; vì mọi iđêan của $B$ đều đóng \sref[0]{0.7.3.5}, ta cũng có $\mathfrak{K}=\mathfrak{J}B$.
Nếu $\sh{K}=\mathfrak{K}^\Delta$, thì đẳng thức $f^*(\sh{J})\sh{O}_\mathfrak{X}=\sh{K}$ suy ra từ \sref{I.10.10.9} và kết thúc chứng minh.
\end{proof}

\begin{env}[10.12.3.2]
\label{I.10.12.3.2}
Tương đương trên cung cấp, đối với hai $\mathfrak{S}$-tiền lược đồ adic $\mathfrak{X}$ và $\mathfrak{Y}$, một \emph{song ánh chính tắc}
\[
  \Hom_\mathfrak{S}(\mathfrak{X},\mathfrak{Y})\isoto\varprojlim_n\Hom_{S_n}(X_n,Y_n)
\]
trong đó giới hạn xạ ảnh được lấy theo các ánh xạ $u_n\to(u_n)_{(S_m)}$ khi $m\leq n$.
\end{env}

\subsection{Các cấu xạ kiểu hữu hạn}
\label{subsection:I.10.13}

\begin{proposition}[10.13.1]
\label{I.10.13.1}
Cho $\mathfrak{Y}$ là một tiền lược đồ hình thức Noether địa phương, $\sh{K}$ là một iđêan định nghĩa của $\mathfrak{Y}$, và $f:\mathfrak{X}\to\mathfrak{Y}$ là một cấu xạ các tiền lược đồ hình thức.
Các điều kiện sau là tương đương.
\begin{enumerate}
  \item[\rm{(a)}] $\mathfrak{X}$ là Noether địa phương, $f$ là một cấu xạ adic \sref{I.10.12.1} và, nếu đặt $\sh{J}=f^*(\sh{K})\sh{O}_\mathfrak{X}$, thì cấu xạ $f_0:(\mathfrak{X},\sh{O}_\mathfrak{X}/\sh{J})\to(\mathfrak{Y},\sh{O}_\mathfrak{Y}/\sh{K})$ cảm sinh bởi $f$ là kiểu hữu hạn.
  \item[\rm{(b)}] $\mathfrak{X}$ là Noether địa phương và là giới hạn quy nạp của một $(Y_n)$-hệ quy nạp adic $(X_n)$ sao cho cấu xạ $X_0\to Y_0$ là kiểu hữu hạn.
\oldpage[I]{208}
  \item[\rm{(c)}] Mọi điểm của $\mathfrak{Y}$ đều có một lân cận mở affine hình thức Noether $V$ thỏa mãn tính chất sau:
\begin{enumerate}
  \item[(\textbf{Q})] $f^{-1}(V)$ là hợp của một họ hữu hạn các tập mở affine hình thức Noether $U_i$ sao cho vành adic Noether $\Gamma(U_i,\sh{O}_\mathfrak{X})$ đẳng cấu tôpô với thương của một đại số chuỗi hình thức hạn chế \sref[0]{0.7.5.1} trên $\Gamma(V,\sh{O}_\mathfrak{Y})$ bởi một iđêan (tất yếu đóng).
\end{enumerate}
\end{enumerate}
\end{proposition}

\begin{proof}
Hiển nhiên (a) kéo theo (b), theo \sref{I.10.12.3}.
Để chứng minh (b) kéo theo (c), vì vấn đề là địa phương trên $\mathfrak{Y}$, ta có thể giả sử $\mathfrak{Y}=\Spf(B)$, trong đó $B$ là Noether và adic;
đặt $\sh{K}=\mathfrak{K}^\Delta$, với $\mathfrak{K}$ là một iđêan định nghĩa của $B$.
Theo giả thiết, vì $X_0$ là kiểu hữu hạn trên $Y_0$, nên $X_0$ là hợp của một số hữu hạn các tập mở affine $U_i$ sao cho vành $A_{i0}$ của lược đồ affine do $X_0$ cảm sinh trên $U_i$ là một đại số kiểu hữu hạn trên vành $B/\mathfrak{K}$ của $Y_0$ \sref{I.6.3.2}.
Theo \sref{I.5.1.9}, $U_i$ cũng là một tập mở affine trong mỗi tiền lược đồ Noether $X_n$ và, nếu $A_{in}$ là vành của lược đồ affine do $X_n$ cảm sinh trên $U_i$, thì giả thiết (b) kéo theo, với $m\leq n$, rằng $A_{im}$ đẳng cấu với $A_{in}/\mathfrak{K}^{m+1}A_{in}$.
Do đó, tiền lược đồ hình thức do $\mathfrak{X}$ cảm sinh trên $U_i$ đẳng cấu với $\Spf(A_i)$, trong đó $A_i=\varprojlim_n A_{in}$ \sref{I.10.6.4};
$A_i$ là một vành $\mathfrak{K}A_i$-adic, và $A_i/\mathfrak{K}A_i$, vốn đẳng cấu với $A_{i0}$, là một đại số kiểu hữu hạn trên $B/\mathfrak{K}$.
Do đó, từ \sref[0]{0.7.5.5} suy ra $A_i$ đẳng cấu tôpô với một thương của một đại số chuỗi hình thức hạn chế trên $B$ (bởi một iđêan tất yếu đóng, vì một đại số như vậy là Noether \sref[0]{0.7.5.4}).

Để chứng minh (c) kéo theo (a), ta có thể giới hạn ở trường hợp $\mathfrak{X}=\Spf(A)$ cũng là affine, với $A$ là một vành adic Noether đẳng cấu với một thương của một đại số chuỗi hình thức hạn chế trên $B$ bởi một iđêan đóng.
Khi đó, theo \sref[0]{0.7.5.5}, $A/\mathfrak{K}A$ là một đại số kiểu hữu hạn trên $B/\mathfrak{K}$ và $\mathfrak{K}A=\mathfrak{J}$ là một iđêan định nghĩa của $A$; do đó, theo \sref{I.10.10.9}, các điều kiện của (a) được thỏa mãn.
\end{proof}

Ta lưu ý rằng, nếu các điều kiện của Mệnh đề~\sref{I.10.13.1} được thỏa mãn, thì tính chất~(a) đúng với \emph{mọi} iđêan định nghĩa $\sh{K}$ của $\mathfrak{Y}$ (theo \emph{(c)}) và do đó, trong tính chất~(b), \emph{tất cả} các $f_n$ đều là các cấu xạ kiểu hữu hạn.

\begin{corollary}[10.13.2]
\label{I.10.13.2}
Nếu các điều kiện của \sref{I.10.13.1} được thỏa mãn, thì mọi tập mở affine hình thức Noether $V$ của $\mathfrak{Y}$ đều có tính chất~(\textbf{Q}) và, nếu $\mathfrak{Y}$ là Noether, thì $\mathfrak{X}$ cũng vậy.
\end{corollary}

\begin{proof}
Điều này suy ra ngay từ \sref{I.10.13.1} và \sref{I.6.3.2}.
\end{proof}

\begin{definition}[10.13.3]
\label{I.10.13.3}
Khi các tính chất tương đương (a), (b) và (c) của \sref{I.10.13.1} được thỏa mãn, ta nói cấu xạ $f$ là \emph{kiểu hữu hạn}, hoặc $\mathfrak{X}$ là một \emph{$\mathfrak{Y}$-tiền lược đồ hình thức kiểu hữu hạn}, hoặc một \emph{tiền lược đồ hình thức kiểu hữu hạn trên $\mathfrak{Y}$}.
\end{definition}

\begin{corollary}[10.13.4]
\label{I.10.13.4}
Cho $\mathfrak{X}=\Spf(A)$ và $\mathfrak{Y}=\Spf(B)$ là hai lược đồ affine hình thức Noether;
để $\mathfrak{X}$ là kiểu hữu hạn trên $\mathfrak{Y}$, điều kiện cần và đủ là vành adic Noether $A$ đẳng cấu với thương của một đại số chuỗi hình thức hạn chế trên $B$ bởi một iđêan đóng.
\end{corollary}

\begin{proof}
Với các ký hiệu của \sref{I.10.13.1}, nếu $\mathfrak{X}$ là kiểu hữu hạn trên $\mathfrak{Y}$, thì $A/\mathfrak{K}A$ là một $(B/\mathfrak{K})$-đại số kiểu hữu hạn theo \sref{I.6.3.3}, và $\mathfrak{K}A$ là một iđêan định nghĩa của $A$ \sref{I.10.10.9}.
Khi đó kết luận suy ra từ \sref[0]{0.7.5.5}.
\end{proof}

\begin{proposition}[10.13.5]
\label{I.10.13.5}
\medskip\noindent
\begin{enumerate}
  \item[\rm{(i)}] Hợp thành của hai cấu xạ kiểu hữu hạn giữa các tiền lược đồ hình thức là kiểu hữu hạn.
\oldpage[I]{209}
  \item[\rm{(ii)}] Cho $\mathfrak{X}$, $\mathfrak{S}$ và $\mathfrak{S}'$ là ba tiền lược đồ hình thức Noether địa phương (tương ứng, Noether), và cho $f:\mathfrak{X}\to\mathfrak{S}$ và $g:\mathfrak{S}'\to\mathfrak{S}$ là hai cấu xạ.
    Nếu $f$ là kiểu hữu hạn, thì $\mathfrak{X}\times_\mathfrak{S}\mathfrak{S}'$ là Noether địa phương (tương ứng, Noether) và là kiểu hữu hạn trên $\mathfrak{S}'$.
  \item[\rm{(iii)}] Cho $\mathfrak{S}$ là một tiền lược đồ hình thức Noether địa phương, và cho $\mathfrak{X}'$, $\mathfrak{Y}'$ là hai $\mathfrak{S}$-tiền lược đồ hình thức Noether địa phương sao cho $\mathfrak{X}'\times_\mathfrak{S}\mathfrak{Y}'$ là Noether địa phương.
    Nếu $\mathfrak{X}$ và $\mathfrak{Y}$ là các $\mathfrak{S}$-tiền lược đồ hình thức Noether địa phương, và $f:\mathfrak{X}\to\mathfrak{X}'$, $g:\mathfrak{Y}\to\mathfrak{Y}'$ là hai $\mathfrak{S}$-cấu xạ kiểu hữu hạn, thì $\mathfrak{X}\times_\mathfrak{S}\mathfrak{Y}$ là Noether địa phương và $f\times_\mathfrak{S}g$ là một $\mathfrak{S}$-cấu xạ kiểu hữu hạn.
\end{enumerate}
\end{proposition}

\begin{proof}
Theo lập luận hình thức của \sref{I.3.5.1}, (iii) suy ra từ (i) và (ii); do đó chỉ cần chứng minh (i) và (ii).

Cho $\mathfrak{X}$, $\mathfrak{Y}$ và $\mathfrak{Z}$ là ba tiền lược đồ hình thức Noether địa phương, và cho $f:\mathfrak{X}\to\mathfrak{Y}$ và $g:\mathfrak{Y}\to\mathfrak{Z}$ là hai cấu xạ kiểu hữu hạn.
Nếu $\sh{L}$ là một iđêan định nghĩa của $\mathfrak{Z}$, thì $\sh{K}=g^*(\sh{L})\sh{O}_\mathfrak{Y}$ là một iđêan định nghĩa của $\mathfrak{Y}$ và $\sh{J}=f^*(g^*(\sh{L}))\sh{O}_\mathfrak{X}$ là một iđêan định nghĩa của $\mathfrak{X}$.
Đặt $X_0=(\mathfrak{X},\sh{O}_\mathfrak{X}/\sh{J})$, $Y_0=(\mathfrak{Y},\sh{O}_\mathfrak{Y}/\sh{K})$ và $Z_0=(\mathfrak{Z},\sh{O}_\mathfrak{Z}/\sh{L})$, và cho $f_0:X_0\to Y_0$, $g_0:Y_0\to Z_0$ là các cấu xạ tương ứng với $f$ và $g$.
Theo giả thiết, vì $f_0$ và $g_0$ là kiểu hữu hạn nên $g_0\circ f_0$ cũng là kiểu hữu hạn \sref{I.6.3.4}, và nó tương ứng với $g\circ f$;
do đó $g\circ f$ là kiểu hữu hạn theo \sref{I.10.13.1}.

Trong các điều kiện của (ii), $\mathfrak{S}$ (tương ứng, $\mathfrak{X}$, $\mathfrak{S}'$) là giới hạn quy nạp của một dãy $(S_n)$ (tương ứng, $(X_n)$, $(S'_n)$) các tiền lược đồ Noether địa phương, và ta có thể giả sử theo \sref{I.10.13.1} rằng $X_m=X_n\times_{S_n}S_m$ khi $m\leq n$.
Khi đó tiền lược đồ hình thức $\mathfrak{X}\times_\mathfrak{S}\mathfrak{S}'$ là giới hạn quy nạp của các tiền lược đồ $X_n\times_{S_n}S'_n$ \sref{I.10.7.4}, và ta có
\[
    X_m\times_{S_m}S'_m = (X_n\times_{S_n}S_m)\times_{S_m}S'_m = (X_n\times_{S_n}S'_n)\times_{S'_n}S'_m.
\]
Hơn nữa, $X_0\times_{S_0}S'_0$ là Noether địa phương, vì $X_0$ là kiểu hữu hạn trên $S_0$ \sref{I.6.3.8}.
Trước hết, từ \sref{I.10.12.3.1} suy ra $\mathfrak{X}\times_\mathfrak{S}\mathfrak{S}'$ là Noether địa phương;
hơn nữa, vì $X_0\times_{S_0}S'_0$ là kiểu hữu hạn trên $S'_0$ \sref{I.6.3.8}, nên từ \sref{I.10.12.3.1} và \sref{I.10.13.1} suy ra $\mathfrak{X}\times_\mathfrak{S}\mathfrak{S}'$ là kiểu hữu hạn trên $\mathfrak{S}'$, kết thúc chứng minh (ii) (khẳng định về các tiền lược đồ Noether là hệ quả trực tiếp của \sref{I.6.3.8}).
\end{proof}

\begin{corollary}[10.13.6]
\label{I.10.13.6}
Dưới các giả thiết của \sref{I.10.9.9}, nếu $f$ là một cấu xạ kiểu hữu hạn, thì phần mở rộng $\widehat{f}$ của nó đến các đầy đủ hóa cũng là kiểu hữu hạn.
\end{corollary}

\subsection{Các tiền lược đồ con đóng của tiền lược đồ hình thức}
\label{subsection:I.10.14}

\begin{proposition}[10.14.1]
\label{I.10.14.1}
Cho $\mathfrak{X}$ là một tiền lược đồ hình thức Noether địa phương và $\sh{A}$ là một bó coherent các iđêan của $\sh{O}_\mathfrak{X}$.
Nếu $\mathfrak{Y}$ là giá (đóng) của $\sh{O}_\mathfrak{X}/\sh{A}$, thì không gian vành tôpô $(\mathfrak{Y},(\sh{O}_\mathfrak{X}/\sh{A})|\mathfrak{Y})$ là một tiền lược đồ hình thức Noether địa phương, và là Noether nếu $\mathfrak{X}$ là Noether.
\end{proposition}

\begin{proof}
Lưu ý rằng $\sh{O}_\mathfrak{X}/\sh{A}$ là coherent theo \sref{I.10.10.3} và \sref[0]{0.5.3.4}; do đó giá $\mathfrak{Y}$ của nó là đóng \sref[0]{0.5.2.2}.
Cho $\sh{J}$ là một iđêan định nghĩa của $\mathfrak{X}$ và đặt $X_n=(\mathfrak{X},\sh{O}_\mathfrak{X}/\sh{J}^{n+1})$;
bó vành $\sh{O}_\mathfrak{X}/\sh{A}$ là giới hạn xạ ảnh của các bó $\sh{O}_\mathfrak{X}/(\sh{A}+\sh{J}^{n+1})=(\sh{O}_\mathfrak{X}/\sh{A})\otimes_{\sh{O}_\mathfrak{X}}(\sh{O}_\mathfrak{X}/\sh{J}^{n+1})$ \sref{I.10.11.3}, tất cả đều có giá là $\mathfrak{Y}$.
Bó $(\sh{A}+\sh{J}^{n+1})/\sh{J}^{n+1}$ là một $\sh{O}_\mathfrak{X}$-môđun coherent, vì $\sh{J}^{n+1}$ là coherent; do đó $(\sh{A}+\sh{J}^{n+1})/\sh{J}^{n+1}$ cũng là một $(\sh{O}_\mathfrak{X}/\sh{J}^{n+1})$-môđun coherent \sref[0]{0.5.3.10};
nếu $Y_n$ là tiền lược đồ con đóng của $X_n$ được xác định bởi bó iđêan này, thì hiển nhiên $(\mathfrak{Y},(\sh{O}_\mathfrak{X}/\sh{A})|\mathfrak{Y})$
\oldpage[I]{210}
là tiền lược đồ hình thức giới hạn quy nạp của các $Y_n$ và, vì các điều kiện của \sref{I.10.6.4} được thỏa mãn, điều này chứng minh rằng tiền lược đồ hình thức ấy là Noether địa phương, và là Noether nếu $\mathfrak{X}$ là Noether (vì khi đó $Y_0$ là Noether theo \sref{I.6.1.4}).
\end{proof}

\begin{definition}[10.14.2]
\label{I.10.14.2}
Ta gọi một \emph{tiền lược đồ con đóng} của một tiền lược đồ hình thức $\mathfrak{X}$ là bất kỳ tiền lược đồ hình thức nào có dạng $(\mathfrak{Y},(\sh{O}_\mathfrak{X}/\sh{A})|\mathfrak{Y})$, trong đó $\sh{A}$ là một \erratum[II]{iđêan coherent của $\sh{O}_\mathfrak{X}$};
ta nói tiền lược đồ này là tiền lược đồ con được xác định bởi $\sh{A}$.
\end{definition}

Rõ ràng phép tương ứng vừa được xác định giữa các \erratum[II]{iđêan coherent của $\sh{O}_\mathfrak{X}$} và các tiền lược đồ con đóng của $\mathfrak{X}$ là song ánh.

Cấu xạ của các không gian vành tôpô $j=(\psi,\theta):\mathfrak{Y}\to\mathfrak{X}$, trong đó $\psi$ là phép nhúng $\mathfrak{Y}\to\mathfrak{X}$ và $\theta^\sharp$ là đồng cấu chính tắc $\sh{O}_\mathfrak{X}\to\sh{O}_\mathfrak{X}/\sh{A}$, hiển nhiên là một cấu xạ các tiền lược đồ hình thức theo \sref{I.10.4.5}, và được gọi là \emph{phép nhúng chính tắc} từ $\mathfrak{Y}$ vào $\mathfrak{X}$.
Lưu ý rằng, nếu $\mathfrak{X}=\Spf(A)$, trong đó $A$ là Noether và adic, thì $\sh{A}=\mathfrak{a}^\Delta$, với $\mathfrak{a}$ là một iđêan của $A$ \sref{I.10.10.5}; từ những điều trên suy ra ngay rằng $\mathfrak{Y}=\Spf(A/\mathfrak{a})$, sai khác một đẳng cấu, và $j$ tương ứng theo \sref{I.10.2.2} với đồng cấu chính tắc $A\to A/\mathfrak{a}$.

Ta nói một cấu xạ $f:\mathfrak{Z}\to\mathfrak{X}$ của các tiền lược đồ hình thức Noether địa phương là một \emph{phép nhúng đóng} nếu nó phân tích thành $\mathfrak{Z}\xrightarrow{g}\mathfrak{Y}\xrightarrow{j}\mathfrak{X}$, trong đó $g$ là một đẳng cấu từ $\mathfrak{Z}$ lên một tiền lược đồ con đóng $\mathfrak{Y}$ của $\mathfrak{X}$ và $j$ là phép nhúng chính tắc.
Vì $j$ là một đơn cấu của các không gian vành, $g$ và $\mathfrak{Y}$ nhất thiết là \emph{duy nhất}.

\begin{proposition}[10.14.3]
\label{I.10.14.3}
Một phép nhúng đóng là một cấu xạ kiểu hữu hạn.
\end{proposition}

\begin{proof}
Ta quy ngay về trường hợp $\mathfrak{X}$ là một lược đồ affine hình thức $\Spf(A)$ và $\mathfrak{Y}=\Spf(A/\mathfrak{a})$;
khi đó mệnh đề suy ra từ Mệnh đề~\sref{I.10.13.1}[\emph{c}].
\end{proof}

\begin{lemma}[10.14.4]
\label{I.10.14.4}
Cho $f:\mathfrak{Y}\to\mathfrak{X}$ là một cấu xạ của các tiền lược đồ hình thức Noether địa phương, và cho $(U_\alpha)$ là một phủ của $f(\mathfrak{Y})$ bởi các tập mở affine hình thức Noether của $\mathfrak{X}$ sao cho các $f^{-1}(U_\alpha)$ là những tập mở affine hình thức Noether của $\mathfrak{Y}$.
Để $f$ là một phép nhúng đóng, điều kiện cần và đủ là $f(\mathfrak{Y})$ là một tập con đóng của $\mathfrak{X}$ và, với mọi $\alpha$, hạn chế của $f$ trên $f^{-1}(U_\alpha)$ tương ứng theo \sref{I.10.4.6} với một đồng cấu toàn ánh $\Gamma(U_\alpha,\sh{O}_\mathfrak{X})\to\Gamma(f^{-1}(U_\alpha),\sh{O}_\mathfrak{Y})$.
\end{lemma}

\begin{proof}
Các điều kiện hiển nhiên là cần thiết.
Ngược lại, nếu chúng được thỏa mãn, và nếu ký hiệu $\mathfrak{a}_\alpha$ là hạt nhân của $\Gamma(U_\alpha,\sh{O}_\mathfrak{X})\to\Gamma(f^{-1}(U_\alpha),\sh{O}_\mathfrak{Y})$, ta định nghĩa một bó coherent các iđêan $\sh{A}$ của $\sh{O}_\mathfrak{X}$ bằng cách đặt $\sh{A}|U_\alpha=\mathfrak{a}_\alpha^\Delta$ và lấy $\sh{A}$ bằng $\sh{O}_\mathfrak{X}$ trên phần bù của hợp các $U_\alpha$.
\footnote{\emph{[Dịch.] Nguồn tiếng Pháp in rằng $\sh{A}$ bằng không trên phần bù này và, trong câu kế tiếp, rằng chính $\mathfrak{a}_\alpha^\Delta$ có giá $U_\alpha\cap f(\mathfrak{Y})$. Iđêan xác định phải là iđêan đơn vị bên ngoài $f(\mathfrak{Y})$, còn bó thương của nó có giá ấy; cả hai hiệu chỉnh đều cần thiết để các bó dán thành tiền lược đồ con đóng có không gian nền là $f(\mathfrak{Y})$.}}
Thật vậy, vì $f(\mathfrak{Y})$ đóng và giá của bó thương $\bigl(\sh{O}_\mathfrak{X}|U_\alpha\bigr)/\mathfrak{a}_\alpha^\Delta$ là $U_\alpha\cap f(\mathfrak{Y})$, mọi việc quy về chứng minh rằng $\mathfrak{a}_\alpha^\Delta$ và $\mathfrak{a}_\beta^\Delta$ cảm sinh cùng một bó trên mọi tập mở affine hình thức Noether $V\subset U_\alpha\cap U_\beta$.
Nhưng hạn chế trên $f^{-1}(U_\alpha)$ của $f$ là một phép nhúng đóng của tiền lược đồ hình thức này vào $U_\alpha$, $f^{-1}(V)$ là một tập mở affine hình thức Noether của $f^{-1}(U_\alpha)$, và hạn chế của $f$ trên $f^{-1}(V)$ là một phép nhúng đóng;
nếu $\mathfrak{b}$ là hạt nhân của đồng cấu toàn ánh $\Gamma(V,\sh{O}_\mathfrak{X})\to\Gamma(f^{-1}(V),\sh{O}_\mathfrak{Y})$ tương ứng với hạn chế này, thì theo \sref{I.10.10.2}, hiển nhiên $\mathfrak{a}_\alpha^\Delta$ cảm sinh $\mathfrak{b}^\Delta$ trên $V$.
Sau khi bó iđêan $\sh{A}$ được định nghĩa như vậy, rõ ràng $f=j\circ g$\footnote{\emph{[Dịch.] Nguồn tiếng Pháp in $f=g\circ j$; với $g:\mathfrak{Y}\to\mathfrak{Z}$ và $j:\mathfrak{Z}\to\mathfrak{X}$, hợp thành đúng kiểu là $j\circ g$.}}, trong đó $j:\mathfrak{Z}\to\mathfrak{X}$ là phép nhúng chính tắc của tiền lược đồ con đóng $\mathfrak{Z}$ của $\mathfrak{X}$ được xác định bởi $\sh{A}$, và $g$ là một đẳng cấu từ $\mathfrak{Y}$ lên $\mathfrak{Z}$.
\end{proof}

\begin{proposition}[10.14.5]
\label{I.10.14.5}
\medskip\noindent
\begin{enumerate}
  \item[\rm{(i)}] Nếu $f:\mathfrak{Z}\to\mathfrak{Y}$ và $g:\mathfrak{Y}\to\mathfrak{X}$ là các phép nhúng đóng của các tiền lược đồ hình thức Noether địa phương, thì $g\circ f$ là một phép nhúng đóng.
\oldpage[I]{211}
  \item[\rm{(ii)}] Cho $\mathfrak{X}$, $\mathfrak{Y}$ và $\mathfrak{S}$ là ba tiền lược đồ hình thức Noether địa phương, cho $f:\mathfrak{X}\to\mathfrak{S}$ là một phép nhúng đóng và $g:\mathfrak{Y}\to\mathfrak{S}$ là một cấu xạ.
    Khi đó cấu xạ $\mathfrak{X}\times_\mathfrak{S}\mathfrak{Y}\to\mathfrak{Y}$ là một phép nhúng đóng.
  \item[\rm{(iii)}] Cho $\mathfrak{S}$ là một tiền lược đồ hình thức Noether địa phương, và cho $\mathfrak{X}'$, $\mathfrak{Y}'$ là hai $\mathfrak{S}$-tiền lược đồ hình thức Noether địa phương sao cho $\mathfrak{X}'\times_\mathfrak{S}\mathfrak{Y}'$ là Noether địa phương.
    Nếu $\mathfrak{X}$ và $\mathfrak{Y}$ là các $\mathfrak{S}$-tiền lược đồ hình thức Noether địa phương, và $f:\mathfrak{X}\to\mathfrak{X}'$, $g:\mathfrak{Y}\to\mathfrak{Y}'$ là các $\mathfrak{S}$-cấu xạ đồng thời là phép nhúng đóng, thì $f\times_\mathfrak{S}g$ là một phép nhúng đóng.
\end{enumerate}
\end{proposition}

\begin{proof}
Theo \sref{I.3.5.1}, một lần nữa chỉ cần chứng minh (i) và (ii).

Để chứng minh (i), có thể giả sử $\mathfrak{Y}$ (tương ứng, $\mathfrak{Z}$) là một tiền lược đồ con đóng của $\mathfrak{X}$ (tương ứng, $\mathfrak{Y}$), được xác định bởi một bó coherent $\sh{J}$ (tương ứng, $\sh{K}$) các iđêan của $\sh{O}_\mathfrak{X}$ (tương ứng, $\sh{O}_\mathfrak{Y}$);
nếu $\psi$ là phép nhúng $\mathfrak{Y}\to\mathfrak{X}$ của các không gian nền, thì $\psi_*(\sh{K})$ là một bó coherent các iđêan của $\psi_*(\sh{O}_\mathfrak{Y})=\sh{O}_\mathfrak{X}/\sh{J}$ \sref[0]{0.5.3.12}, và do đó cũng là một $\sh{O}_\mathfrak{X}$-môđun coherent \sref[0]{0.5.3.10};
vì vậy hạt nhân $\sh{K}_1$ của $\sh{O}_\mathfrak{X}\to(\sh{O}_\mathfrak{X}/\sh{J})/\psi_*(\sh{K})$ là một bó coherent các iđêan của $\sh{O}_\mathfrak{X}$ \sref[0]{0.5.3.4}, và $\sh{O}_\mathfrak{X}/\sh{K}_1$ đẳng cấu với $\psi_*(\sh{O}_\mathfrak{Y}/\sh{K})$, điều này chứng minh rằng $\mathfrak{Z}$ đẳng cấu với một tiền lược đồ con đóng của $\mathfrak{X}$.

Để chứng minh (ii), có thể quy ngay về trường hợp $\mathfrak{S}=\Spf(A)$, $\mathfrak{X}=\Spf(B)$ và $\mathfrak{Y}=\Spf(C)$, trong đó $A$ là một vành Noether $\mathfrak{J}$-adic, $B=A/\mathfrak{a}$ (với $\mathfrak{a}$ là một iđêan của $A$) và $C$ là một $A$-đại số tôpô adic Noether.
Mọi việc khi đó quy về chứng minh rằng đồng cấu $C\to C\widehat{\otimes}_A(A/\mathfrak{a})$ là \emph{toàn ánh}:
nhưng $A/\mathfrak{a}$ là một $A$-môđun kiểu hữu hạn, và tôpô của nó là tôpô $\mathfrak{J}$-adic;
do đó, từ \sref[0]{0.7.7.8} suy ra $C\widehat{\otimes}_A(A/\mathfrak{a})$ được đồng nhất với $C\otimes_A(A/\mathfrak{a})=C/\mathfrak{a}C$, từ đó có khẳng định.
\end{proof}

\begin{corollary}[10.14.6]
\label{I.10.14.6}
Dưới các giả thiết của \sref{I.10.14.5}[ii], cho $p:\mathfrak{X}\times_\mathfrak{S}\mathfrak{Y}\to\mathfrak{X}$ và $q:\mathfrak{X}\times_\mathfrak{S}\mathfrak{Y}\to\mathfrak{Y}$ là các phép chiếu, sao cho biểu đồ
\[
  \xymatrix{
    \mathfrak{X}\ar[d]_f &
    \mathfrak{X}\times_\mathfrak{S}\mathfrak{Y}\ar[l]_p\ar[d]^q\\
    \mathfrak{S} &
    \mathfrak{Y}\ar[l]^g
  }
\]
giao hoán.
Với mọi $\sh{O}_\mathfrak{X}$-môđun coherent $\sh{F}$, khi đó ta có một đẳng cấu chính tắc của các $\sh{O}_\mathfrak{Y}$-môđun
\[
\label{I.10.14.6.1}
  u:g^*f_*(\sh{F})\isoto q_*p^*(\sh{F}).
  \tag{10.14.6.1}
\]
\end{corollary}

\begin{proof}
Ta biết rằng việc định nghĩa một đồng cấu $g^*f_*(\sh{F}) \to q_*p^*(\sh{F})$ tương đương với việc định nghĩa một đồng cấu $f_*(\sh{F}) \to g_*q_*p^*(\sh{F}) = f_*p_*p^*(\sh{F})$ \sref[0]{0.4.4.3}:
ta lấy $u$ là đồng cấu tương ứng qua phép liên hợp này với $f_*(\rho)$,\footnote{\emph{[Dịch.] Nguồn tiếng Pháp viết $u=f_*(\rho)$; các ánh xạ này sống trên những không gian khác nhau và có các bó nguồn và đích khác nhau, trong đó $f_*(\rho)$ là ánh xạ liên hợp của $u$.}} với $\rho$ là đồng cấu chính tắc $\sh{F}\to p_*p^*(\sh{F})$ \sref[0]{0.4.4.3}.
Để thấy $u$ là một đẳng cấu, có thể quy ngay về trường hợp $\mathfrak{S}$, $\mathfrak{X}$ và $\mathfrak{Y}$ lần lượt là các phổ hình thức của các vành adic Noether $A$, $B$ và $C$, thỏa mãn các điều kiện trong \sref{I.10.14.5}[ii];
khi đó $\sh{F}=M^\Delta$, trong đó $M$ là một $(A/\mathfrak{a})$-môđun kiểu hữu hạn \sref{I.10.10.5}, và hai vế của \sref{I.10.14.6.1} lần lượt được đồng nhất, theo \sref{I.10.10.8}, với $(C\otimes_A M)^\Delta$ và $((C/\mathfrak{a}C)\otimes_{A/\mathfrak{a}}M)^\Delta$; từ đó suy ra hệ quả, vì $(C/\mathfrak{a}C)\otimes_{A/\mathfrak{a}}M = (C\otimes_A(A/\mathfrak{a}))\otimes_{A/\mathfrak{a}}M$ được đồng nhất chính tắc với $C\otimes_A M$.
\end{proof}

\oldpage[I]{212}
\begin{corollary}[10.14.7]
\label{I.10.14.7}
Cho $X$ là một tiền lược đồ thông thường Noether địa phương, $Y$ là một tiền lược đồ con đóng của $X$, $j$ là phép nhúng chính tắc $Y\to X$, $X'$ là một tập con đóng của $X$ và $Y'=Y\cap X'$;
khi đó $\widehat{j}:Y_{/Y'}\to X_{/X'}$ là một phép nhúng đóng và, với mọi $\sh{O}_Y$-môđun coherent $\sh{F}$, ta có
\[
    \widehat{j}_*(\sh{F}_{/Y'}) = (j_*(\sh{F}))_{/X'}.
\]
\end{corollary}

\begin{proof}
Vì $Y'=j^{-1}(X')$, chỉ cần sử dụng \sref{I.10.9.9} và áp dụng \sref{I.10.14.5} cùng \sref{I.10.14.6}.
\end{proof}

\subsection{Các tiền lược đồ hình thức tách}
\label{subsection:I.10.15}

\begin{definition}[10.15.1]
\label{I.10.15.1}
Cho $\mathfrak{S}$ là một tiền lược đồ hình thức, $\mathfrak{X}$ là một $\mathfrak{S}$-tiền lược đồ hình thức và $f:\mathfrak{X}\to\mathfrak{S}$ là cấu xạ cấu trúc.
Ta gọi cấu xạ đường chéo $\Delta_{\mathfrak{X}|\mathfrak{S}}:\mathfrak{X}\to\mathfrak{X}\times_\mathfrak{S}\mathfrak{X}$ (cũng được ký hiệu $\Delta_\mathfrak{X}$) là cấu xạ $(1_\mathfrak{X},1_\mathfrak{X})_\mathfrak{S}$.
Ta nói $\mathfrak{X}$ tách trên $\mathfrak{S}$, hoặc là một $\mathfrak{S}$-lược đồ hình thức, hoặc $f$ là một cấu xạ tách, nếu ảnh của không gian nền của $\mathfrak{X}$ qua $\Delta_\mathfrak{X}$ là một tập con đóng của không gian nền của $\mathfrak{X}\times_\mathfrak{S}\mathfrak{X}$.
Ta nói một tiền lược đồ hình thức $\mathfrak{X}$ là tách, hoặc là một lược đồ hình thức, nếu nó tách trên $\bb{Z}$.
\end{definition}

\begin{proposition}[10.15.2]
\label{I.10.15.2}
Giả sử các tiền lược đồ hình thức $\mathfrak{S}$ và $\mathfrak{X}$ lần lượt là các giới hạn quy nạp của những dãy $(S_n)$ và $(X_n)$ các tiền lược đồ thông thường, và cấu xạ $f:\mathfrak{X}\to\mathfrak{S}$ là giới hạn quy nạp của một dãy cấu xạ $f_n:X_n\to S_n$.
Để $f$ là tách, điều kiện cần và đủ là cấu xạ $f_0:X_0\to S_0$ tách.
\end{proposition}

\begin{proof}
Thật vậy, khi đó $\Delta_{\mathfrak{X}|\mathfrak{S}}$ là giới hạn quy nạp của dãy các cấu xạ $\Delta_{X_n|S_n}$ \sref{I.10.7.4}; các không gian nền của $\mathfrak{X}$ và $\mathfrak{X}\times_\mathfrak{S}\mathfrak{X}$ lần lượt đồng nhất với các không gian nền của $X_0$ và $X_0\times_{S_0}X_0$, và, dưới những đồng nhất này, ảnh của $\mathfrak{X}$ qua $\Delta_{\mathfrak{X}|\mathfrak{S}}$ đồng nhất với ảnh của $X_0$ qua $\Delta_{X_0|S_0}$;
từ đó có kết luận.
\end{proof}

\begin{proposition}[10.15.3]
\label{I.10.15.3}
Trong phần tiếp theo, giả sử mọi tiền lược đồ hình thức (tương ứng, cấu xạ của các tiền lược đồ hình thức) được xét đều là giới hạn quy nạp của các dãy tiền lược đồ thông thường (tương ứng, của các cấu xạ tiền lược đồ thông thường).
\begin{enumerate}
    \item[\rm{(i)}] Hợp thành của hai cấu xạ tách là tách.
    \item[\rm{(ii)}] Nếu $f:\mathfrak{X}\to\mathfrak{X}'$ và $g:\mathfrak{Y}\to\mathfrak{Y}'$ là hai $\mathfrak{S}$-cấu xạ tách, thì $f\times_\mathfrak{S}g$ là tách.
    \item[\rm{(iii)}] Nếu $f:\mathfrak{X}\to\mathfrak{Y}$ là một $\mathfrak{S}$-cấu xạ tách, thì $\mathfrak{S}'$-cấu xạ $f_{(\mathfrak{S}')}$ là tách đối với mọi mở rộng $\mathfrak{S}'\to\mathfrak{S}$ của tiền lược đồ hình thức cơ sở.
    \item[\rm{(iv)}] Nếu hợp thành $g\circ f$ của hai cấu xạ là tách, thì $f$ là tách.
\end{enumerate}
\emph{(Trong phát biểu này, ta ngầm hiểu rằng nếu cùng một tiền lược đồ hình thức $\mathfrak{Z}$ xuất hiện nhiều lần trong cùng một mệnh đề, thì ở mọi nơi nó xuất hiện ta đều xem nó là giới hạn quy nạp của \emph{cùng một} dãy $(Z_n)$ các tiền lược đồ thông thường, và các cấu xạ từ $\mathfrak{Z}$ đến một tiền lược đồ hình thức khác (tương ứng, từ một tiền lược đồ hình thức đến $\mathfrak{Z}$) là các giới hạn quy nạp của các cấu xạ từ $Z_n$ đến những tiền lược đồ thông thường (tương ứng, từ những tiền lược đồ thông thường đến $Z_n$).)}
\end{proposition}

\begin{proof}
Với các ký hiệu của \sref{I.10.15.2}, thật vậy ta có $(g\circ f)_0=g_0\circ f_0$ và $(f\times_\mathfrak{S}g)_0=f_0\times_{S_0}g_0$;
các khẳng định của \sref{I.10.15.3} khi đó là hệ quả trực tiếp của \sref{I.10.15.2} và các khẳng định tương ứng trong \sref{I.5.5.1} đối với các tiền lược đồ thông thường.
\end{proof}

Ta để độc giả phát biểu, cho cùng loại tiền lược đồ hình thức và cấu xạ như trong \sref{I.10.15.3}, các mệnh đề tương ứng với \sref{I.5.5.5},
\oldpage[I]{213}
\sref{I.5.5.9} và \sref{I.5.5.10} (thay «tập mở affine» bằng «tập mở affine hình thức thỏa mãn điều kiện (b) của \sref{I.10.6.3}»).

Một lập luận tương tự cũng cho thấy mọi lược đồ affine hình thức \emph{Noether} đều tách, điều này biện minh cho thuật ngữ.

\begin{proposition}[10.15.4]
\label{I.10.15.4}
Cho $\mathfrak{S}$ là một tiền lược đồ hình thức Noether địa phương, và $\mathfrak{X}$, $\mathfrak{Y}$ là hai $\mathfrak{S}$-tiền lược đồ hình thức Noether địa phương sao cho $\mathfrak{X}$ hoặc $\mathfrak{Y}$ là kiểu hữu hạn trên $\mathfrak{S}$ (để $\mathfrak{X}\times_\mathfrak{S}\mathfrak{Y}$ là Noether địa phương) và $\mathfrak{Y}$ tách trên $\mathfrak{S}$.
Cho $f:\mathfrak{X}\to\mathfrak{Y}$ là một $\mathfrak{S}$-cấu xạ; khi đó cấu xạ đồ thị $\Gamma_f=(1_\mathfrak{X},f)_\mathfrak{S}:\mathfrak{X}\to\mathfrak{X}\times_\mathfrak{S}\mathfrak{Y}$ là một phép nhúng đóng.
\end{proposition}

\begin{proof}
Có thể giả sử $\mathfrak{S}$ là giới hạn quy nạp của một dãy $(S_n)$ các tiền lược đồ Noether địa phương, $\mathfrak{X}$ (tương ứng, $\mathfrak{Y}$) là giới hạn quy nạp của một dãy $(X_n)$ (tương ứng, $(Y_n)$) các $S_n$-tiền lược đồ, và $f$ là giới hạn quy nạp của một dãy $(f_n:X_n\to Y_n)$ các $S_n$-cấu xạ;
khi đó $\mathfrak{X}\times_\mathfrak{S}\mathfrak{Y}$ là giới hạn quy nạp của dãy $(X_n\times_{S_n}Y_n)$, và $\Gamma_f$ là giới hạn quy nạp của dãy $(\Gamma_{f_n})$ \sref{I.10.7.4};
theo giả thiết, $Y_0$ tách trên $S_0$ \sref{I.10.15.2}, nên không gian $\Gamma_{f_0}(X_0)$ là một không gian con đóng của $X_0\times_{S_0}Y_0$;
vì các không gian nền của $\mathfrak{X}\times_\mathfrak{S}\mathfrak{Y}$ (tương ứng, $\Gamma_f(\mathfrak{X})$) và $X_0\times_{S_0}Y_0$ (tương ứng, $\Gamma_{f_0}(X_0)$) là đồng nhất, ta đã thấy $\Gamma_f(\mathfrak{X})$ là một không gian con \emph{đóng} của $\mathfrak{X}\times_\mathfrak{S}\mathfrak{Y}$.
Bây giờ lưu ý rằng, khi $(U,V)$ chạy qua tập các cặp gồm một tập mở affine hình thức Noether $U$ (tương ứng, $V$) của $\mathfrak{X}$ (tương ứng, $\mathfrak{Y}$) sao cho $f(U)\subset V$, các tập mở $U\times_\mathfrak{S} V$ lập thành một phủ của $\Gamma_f(\mathfrak{X})$ trong $\mathfrak{X}\times_\mathfrak{S}\mathfrak{Y}$ và, nếu $f_U:U\to V$ là hạn chế của $f$ trên $U$, thì $\Gamma_{f_U}:U\to U\times_\mathfrak{S}V$ là hạn chế của $\Gamma_f$ trên $U$.
Nếu chứng minh được $\Gamma_{f_U}$ là một phép nhúng đóng, thì $\Gamma_f$ sẽ là một phép nhúng đóng \sref{I.10.14.4}; nói cách khác, ta được quy về trường hợp $\mathfrak{S}=\Spf(A)$, $\mathfrak{X}=\Spf(B)$ và $\mathfrak{Y}=\Spf(C)$ là affine (với $A$, $B$, $C$ là các vành adic Noether), và $f$ tương ứng với một $A$-đồng cấu liên tục $\vphi:C\to B$;
khi đó $\Gamma_f$ tương ứng với đồng cấu liên tục duy nhất $\omega: B\widehat{\otimes}_A C\to B$ mà khi hợp thành với các đồng cấu chính tắc $B\to B\widehat{\otimes}_A C$ và $C\to B\widehat{\otimes}_A C$ thì lần lượt cho đồng nhất và $\vphi$.
Nhưng rõ ràng $\omega$ là \emph{toàn ánh}, từ đó có khẳng định.
\end{proof}

\begin{corollary}[10.15.5]
\label{I.10.15.5}
Cho $\mathfrak{S}$ là một tiền lược đồ hình thức Noether địa phương và $\mathfrak{X}$ là một $\mathfrak{S}$-tiền lược đồ kiểu hữu hạn;
để $\mathfrak{X}$ tách trên $\mathfrak{S}$, điều kiện cần và đủ là cấu xạ đường chéo $\mathfrak{X}\to\mathfrak{X}\times_\mathfrak{S}\mathfrak{X}$ là một phép nhúng đóng.
\end{corollary}

\begin{proposition}[10.15.6]
\label{I.10.15.6}
Một phép nhúng đóng $j:\mathfrak{Y}\to\mathfrak{X}$ của các tiền lược đồ hình thức Noether địa phương là một cấu xạ tách.
\end{proposition}

\begin{proof}
Với các ký hiệu của \sref{I.10.14.2}, $j_0:Y_0\to X_0$ là một phép nhúng đóng, do đó là một cấu xạ tách, và chỉ cần áp dụng \sref{I.10.15.2}.
\end{proof}

\begin{proposition}[10.15.7]
\label{I.10.15.7}
Cho $X$ là một tiền lược đồ (thông thường) Noether địa phương, $X'$ là một tập con đóng của $X$ và $\widehat{X}=X_{/X'}$.
Để $\widehat{X}$ là tách, điều kiện cần và đủ là $\widehat{X}_\red$ tách, và chỉ cần $X$ tách.
\end{proposition}

\begin{proof}
Với các ký hiệu của \sref{I.10.8.5}, để $\widehat{X}$ là tách, điều kiện cần và đủ là $X'_0$ tách \sref{I.10.15.2}; vì $\widehat{X}_\red=(X'_0)_\red$, điều này tương đương với việc $\widehat{X}_\red$ tách \sref{I.5.5.1}[vi].
\end{proof}

% Bản chiếu TeX lấy nguồn làm chuẩn của bảng ký hiệu EGA I.
% Dòng bản thảo: Markdown đối chiếu của Jacob C. Reinhold; bản quét tiếng Pháp NUMDAM có thẩm quyền.
\clearpage
\section*{Bảng ký hiệu}
\phantomsection
\label{section:index-notation}
\begingroup
\small
\setlength{\columnsep}{2em}
\begin{multicols}{2}
\raggedcolumns
\newcommand{\EGAIndexLine}[2]{\par\noindent\hangindent=1em\hangafter=1 #1\ \dotfill\ \textup{#2}\par}

\phantomsection\label{section:index-notation-chapter0}
\medskip\noindent\textbf{Chương 0}\par\nopagebreak
\EGAIndexLine{$M_{[\phi]}$}{(0.1.0.2)}
\EGAIndexLine{$B\mathfrak{J}$, $\mathfrak{J}B$}{(0.1.0.3)}
\EGAIndexLine{$r(\mathfrak{a})$}{(0.1.1.1)}
\EGAIndexLine{$\mathfrak{R}(A)$}{(0.1.1.2)}
\EGAIndexLine{$S_{f}$}{(0.1.2.1)}
\EGAIndexLine{$S^{-1}A$, $S^{-1}M$, $m/s$, $i_A^S$, $i_M^S$, $i^S$}{(0.1.2.2)}
\EGAIndexLine{$A_{f}$, $M_{f}$, $A_{\mathfrak{p}}$, $M_{\mathfrak{p}}$}{(0.1.2.3)}
\EGAIndexLine{$S^{-1}u$}{(0.1.3.1)}
\EGAIndexLine{$\rho^{T,S}_{A}$, $\rho^{T,S}_{M}$, $\rho^{T,S}$}{(0.1.4.1)}
\EGAIndexLine{$\operatorname{Supp}(M)$}{(0.1.7.1)}
\EGAIndexLine{$\mathcal{F}|U$, $u|U$}{(0.3.1.5)}
\EGAIndexLine{$\mathcal{F}_{x}$, $s_{x}$, $\Gamma(U, \mathcal{F})$, $u(s)$, $\operatorname{Supp}(\mathcal{F})$}{(0.3.1.6)}
\EGAIndexLine{$\psi_{*}(\mathcal{F})$}{(0.3.4.1)}
\EGAIndexLine{$\psi_{*}(u)$}{(0.3.4.2)}
\EGAIndexLine{$\psi_{x}$}{(0.3.4.4)}
\EGAIndexLine{$\mathcal{G} \to \mathcal{F}$}{(0.3.5.1)}
\EGAIndexLine{$u^{\sharp}$, $v^{\flat}$, $\rho_{\mathcal{G}}$}{(0.3.5.3)}
\EGAIndexLine{$\psi^*(\mathcal{G})$, $\psi^*(v)$, $\sigma_{\mathcal{F}}$}{(0.3.5.5)}
\EGAIndexLine{$\mathcal{O}_{X}$, $\mathcal{O}_{X,x}$, $\mathcal{O}_{x}$, $1$, $e$}{(0.4.1.1)}
\EGAIndexLine{$\check{\mathcal{F}}$, $\bigwedge^{p} \mathcal{F}$}{(0.4.1.4)}
\EGAIndexLine{$\mathcal{J}\mathcal{F}$}{(0.4.1.5)}
\EGAIndexLine{$\Psi_{*}(\mathcal{F})$, $\Psi_{*}(u)$}{(0.4.2.1)}
\EGAIndexLine{$\Psi_{*}(\mathcal{C})$}{(0.4.2.4)}
\EGAIndexLine{$\Psi^*(\mathcal{G})$, $\Psi^*(v)$}{(0.4.3.1)}
\EGAIndexLine{$\Psi^*(\mathcal{C})$}{(0.4.3.4)}
\EGAIndexLine{$\Psi^*(\mathcal{J})\mathcal{O}_{X}$, $\mathcal{J}^{c}_{X}$}{(0.4.3.5)}
\EGAIndexLine{$\mathcal{G} \to \mathcal{F}$}{(0.4.4.1)}
\EGAIndexLine{$u_{\theta}^{\sharp}$, $u^{\sharp}$, $v_{\theta}^{\flat}$, $v^{\flat}$, $\rho_{\mathcal{G}}$, $\sigma_{\mathcal{F}}$}{(0.4.4.3)}
\EGAIndexLine{$u_{1} \otimes u_{2}$}{(0.4.4.4)}
\EGAIndexLine{$\mathcal{L}^{-1}$}{(0.5.4.3)}
\EGAIndexLine{$\mathcal{L}^{\otimes n}$}{(0.5.4.4)}
\EGAIndexLine{$\Gamma_{*}(\mathcal{L})$, $\Gamma_{*}(\mathcal{L}, \mathcal{F})$}{(0.5.4.6)}
\EGAIndexLine{$\mathcal{O}_{X}^{*}$}{(0.5.4.7)}
\EGAIndexLine{$\mathfrak{m}_{x}$, $\kappa(x)$, $f(x)$}{(0.5.5.1)}
\EGAIndexLine{$X_{f}$}{(0.5.5.2)}
\EGAIndexLine{$\operatorname{Im}(M' \otimes_{A} N')$}{(0.6.0)}
\EGAIndexLine{$\widehat A$, $\widehat M$}{(0.7.2.3) và (0.7.3.1)}
\EGAIndexLine{$A\{T_{1}, \cdots, T_{r}\}$}{(0.7.5.1)}
\EGAIndexLine{$A\{S^{-1}\}$}{(0.7.6.1)}
\EGAIndexLine{$\mathfrak{a}\{S^{-1}\}$}{(0.7.6.9)}
\EGAIndexLine{$A_{\{f\}}$, $\mathfrak{a}_{\{f\}}$}{(0.7.6.14)}
\EGAIndexLine{$A_{\{S\}}$}{(0.7.6.14)}
\EGAIndexLine{$(M \otimes_{A} N)^{\wedge}$, $M \widehat{\otimes}_{A} N$}{(0.7.7.1)}
\EGAIndexLine{$u \widehat{\otimes} v$}{(0.7.7.3)}

\phantomsection\label{section:index-notation-chapterI}
\medskip\noindent\textbf{Chương I}\par\nopagebreak
\EGAIndexLine{$\operatorname{Spec}(A)$, $\mathfrak{j}_{x}$, $\mathfrak{m}_{x}$, $\kappa(x)$, $f(x)$, $M_{x}$, $r(E)$, $V(E)$, $V(f)$, $D(f)$}{(I, \hyperref[I.1.1.1]{1.1.1})}
\EGAIndexLine{$\mathfrak{j}(Y)$}{(I, \hyperref[I.1.1.3]{1.1.3})}
\EGAIndexLine{${}^{a}\phi$}{(I, \hyperref[I.1.2.1]{1.2.1})}
\EGAIndexLine{$S_{f}'$}{(I, \hyperref[I.1.3.1]{1.3.1})}
\EGAIndexLine{$\rho_{g,f}$}{(I, \hyperref[I.1.3.3]{1.3.3})}
\EGAIndexLine{$\widetilde A$, $\widetilde{M}$, $\theta_{f}$}{(I, \hyperref[I.1.3.4]{1.3.4})}
\EGAIndexLine{$\widetilde u$}{(I, \hyperref[I.1.3.5]{1.3.5})}
\EGAIndexLine{$\widetilde{\phi}$}{(I, \hyperref[I.1.6.1]{1.6.1})}
\EGAIndexLine{$A(X)$}{(I, \hyperref[I.1.7.1]{1.7.1})}
\EGAIndexLine{$\mathcal{O}_{X/Y}$}{(I, \hyperref[I.2.1.6]{2.1.6})}
\EGAIndexLine{$\operatorname{Hom}(X, Y)$}{(I, \hyperref[I.2.2.1]{2.2.1})}
\EGAIndexLine{$\operatorname{Hom}_{S}(X, Y)$, $1_X$}{(I, \hyperref[I.2.5.2]{2.5.2})}
\EGAIndexLine{$\Gamma(X/S)$}{(I, \hyperref[I.2.5.5]{2.5.5})}
\EGAIndexLine{$X \amalg Y$}{(I, \hyperref[I.3.1.1]{3.1.1})}
\EGAIndexLine{$X \times_{S} Y$, $X \times Y$, $(g, h)_{S}$, $u \times_{S} v$, $u \times v$}{(I, \hyperref[I.3.2.1]{3.2.1})}
\EGAIndexLine{$X \times_{A} Y$, $X \otimes_{A} B$, $(g, h)_{A}$, $u \times_{A} v$}{(I, \hyperref[I.3.2.1]{3.2.1})}
\EGAIndexLine{$X_{(S')}$}{(I, \hyperref[I.3.3.6]{3.3.6})}
\EGAIndexLine{$f_{(S')}$}{(I, \hyperref[I.3.3.7]{3.3.7})}
\EGAIndexLine{$\Gamma_{f}$}{(I, \hyperref[I.3.3.14]{3.3.14})}
\EGAIndexLine{$X(T)$}{(I, \hyperref[I.3.4.1]{3.4.1})}
\EGAIndexLine{$P \times_{R} Q$}{(I, \hyperref[I.3.4.2]{3.4.2})}
\EGAIndexLine{$X(T)_{S}$}{(I, \hyperref[I.3.4.3]{3.4.3})}
\EGAIndexLine{$X(B)$, $X(B)_{A}$}{(I, \hyperref[I.3.4.4]{3.4.4})}
\EGAIndexLine{$X \otimes_{y} B$, $X \otimes_{\mathcal{O}_y} B$}{(I, \hyperref[I.3.6.3]{3.6.3})}
\EGAIndexLine{$Z \leq Y$}{(I, \hyperref[I.4.1.10]{4.1.10})}
\EGAIndexLine{$f^{-1}(Y)$}{(I, \hyperref[I.4.4.1]{4.4.1})}
\EGAIndexLine{$\mathcal{N}_{X}$}{(I, \hyperref[I.5.1.1]{5.1.1})}
\EGAIndexLine{$X_{\mathrm{red}}$}{(I, \hyperref[I.5.1.3]{5.1.3})}
\EGAIndexLine{$f_{\mathrm{red}}$}{(I, \hyperref[I.5.1.5]{5.1.5})}
\EGAIndexLine{$\Delta_{X/S}$, $\Delta_{X}$, $\Delta$}{(I, \hyperref[I.5.3.1]{5.3.1})}
\EGAIndexLine{$\operatorname{rg}_{K}(X)$}{(I, \hyperref[I.6.4.5]{6.4.5})}
\EGAIndexLine{$n(X)$}{(I, \hyperref[I.6.4.8]{6.4.8})}
\EGAIndexLine{$\Gamma_{\mathrm{rat}}(X/Y)$}{(I, \hyperref[I.7.1.2]{7.1.2})}
\EGAIndexLine{$R(X)$}{(I, \hyperref[I.7.1.3]{7.1.3})}
\EGAIndexLine{$\mathcal{R}(X)$}{(I, \hyperref[I.7.3.1]{7.3.1})}
\EGAIndexLine{$L(A)$}{(I, \hyperref[I.8.1.2]{8.1.2})}
\EGAIndexLine{$\delta(f)$}{(I, \hyperref[I.8.2.1]{8.2.1})}
\EGAIndexLine{$\mathcal{F} \otimes_{\mathcal{O}_{S}} \mathcal{G}$, $\mathcal{F} \otimes_{S} \mathcal{G}$}{(I, \hyperref[I.9.1.2]{9.1.2})}
\EGAIndexLine{$\overline{\mathcal{G}}$}{(I, \hyperref[I.9.4.1]{9.4.1})}
\EGAIndexLine{$\overline Y$}{(I, \hyperref[I.9.5.11]{9.5.11})}
\EGAIndexLine{$\operatorname{Spf}(A)$, $\mathcal{O}_{\mathfrak{X}}$ ($\mathfrak{X}=\operatorname{Spf}(A)$)}{(I, \hyperref[I.10.1.2]{10.1.2})}
\EGAIndexLine{$\mathfrak{D}(f)$}{(I, \hyperref[I.10.1.4]{10.1.4})}
\EGAIndexLine{${}^{a}\phi$, $\widetilde{\phi}$}{(I, \hyperref[I.10.2.1]{10.2.1})}
\EGAIndexLine{$\mathfrak{J}^{\Delta}$}{(I, \hyperref[I.10.3.1]{10.3.1})}
\EGAIndexLine{$\mathfrak{X} \times_{\mathfrak{S}} \mathfrak{Y}$}{(I, \hyperref[I.10.7.3]{10.7.3})}
\EGAIndexLine{$\mathcal{F}_{/X'}$, $\widehat{\mathcal{F}}$, $u_{/X'}$, $\widehat{u}$}{(I, \hyperref[I.10.8.4]{10.8.4})}
\EGAIndexLine{$X_{/X'}$, $\widehat{X}$}{(I, \hyperref[I.10.8.5]{10.8.5})}
\EGAIndexLine{$\widehat{f}$}{(I, \hyperref[I.10.9.1]{10.9.1})}
\EGAIndexLine{$M^{\Delta}$, $u^{\Delta}$}{(I, \hyperref[I.10.10.1]{10.10.1})}
\EGAIndexLine{$\Delta_{\mathfrak{S}|\mathfrak{X}}$, $\Delta_{\mathfrak{X}}$}{(I, \hyperref[I.10.15.1]{10.15.1})}
\end{multicols}
\endgroup

% Bản chiếu tiếng Việt lấy nguồn làm chuẩn của chỉ mục thuật ngữ EGA I.
% Bản quét tiếng Pháp NUMDAM có thẩm quyền. Các mục giữ nguyên thứ tự của chỉ mục nguồn
% để từng dòng có thể được đối chiếu trực tiếp với các trang in 219--224.
\clearpage
\section*{Chỉ mục thuật ngữ}
\phantomsection
\label{section:index-terminology}
\begingroup
\footnotesize
\setlength{\columnsep}{2em}
\begin{multicols}{2}
\raggedcolumns
\newcommand{\EGATermIndexLine}[2]{\par\noindent\hangindent=1em\hangafter=1 #1\ \dotfill\ \textup{#2}\par}

\noindent\emph{Thứ tự các mục theo chỉ mục của nguồn tiếng Pháp;
vì thế các bản dịch tiếng Việt không phải lúc nào cũng theo thứ tự chữ cái tiếng Việt.}\par\medskip

\EGATermIndexLine{Bao đóng của một tiền lược đồ con}{(I, \hyperref[I.9.5.11]{9.5.11})}
\EGATermIndexLine{$\mathcal O_X$-đại số}{(0.4.1.3)}
\EGATermIndexLine{Vành adic; vành $\mathfrak J$-adic}{(0.7.1.9)}
\EGATermIndexLine{Vành khả chấp}{(0.7.1.2)}
\EGATermIndexLine{Vành phân thức đầy đủ}{(0.7.6.5)}
\EGATermIndexLine{Vành phân thức}{(0.1.2.2)}
\EGATermIndexLine{Vành các hàm hữu tỉ}{(I, \hyperref[I.2.1.6]{2.1.6} và \hyperref[I.7.1.3]{7.1.3})}
\EGATermIndexLine{Vành của một lược đồ affine}{(I, \hyperref[I.1.7.1]{1.7.1})}
\EGATermIndexLine{Vành nguyên}{(0.1.0.6)}
\EGATermIndexLine{Vành được tôpô hóa tuyến tính}{(0.7.1.1)}
\EGATermIndexLine{Vành địa phương}{(0.1.0.7)}
\EGATermIndexLine{Vành địa phương chi phối}{(I, \hyperref[I.8.1.1]{8.1.1})}
\EGATermIndexLine{Vành địa phương của $X$ dọc theo $Y$; vành địa phương của $Y$ trong $X$}{(I, \hyperref[I.2.1.6]{2.1.6})}
\EGATermIndexLine{Vành tiền-adic; vành $\mathfrak J$-tiền-adic}{(0.7.1.9)}
\EGATermIndexLine{Vành tiền-khả chấp}{(0.7.1.2)}
\EGATermIndexLine{Vành rút gọn}{(0.1.1.1)}
\EGATermIndexLine{Vành chính quy}{(0.4.1.3)}
\EGATermIndexLine{Các vành địa phương liên hệ với nhau}{(I, \hyperref[I.8.1.4]{8.1.4})}
\EGATermIndexLine{Linh hóa tử của một $\mathcal O_X$-môđun}{(0.5.3.7)}
\EGATermIndexLine{Ánh xạ phổ liên kết với một đồng cấu vành}{(I, \hyperref[I.1.2.1]{1.2.1})}
\EGATermIndexLine{Ánh xạ hữu tỉ; $S$-ánh xạ hữu tỉ}{(I, \hyperref[I.7.1.2]{7.1.2})}
\EGATermIndexLine{Ánh xạ hữu tỉ xác định tại một điểm}{(I, \hyperref[I.7.2.1]{7.2.1})}
\EGATermIndexLine{Ánh xạ hữu tỉ cảm sinh trên một tập con mở}{(I, \hyperref[I.7.1.2]{7.1.2})}
\EGATermIndexLine{Ánh xạ hữu tỉ cảm sinh trên $\operatorname{Spec}(\mathcal O_x)$}{(I, \hyperref[I.7.2.8]{7.2.8})}

\EGATermIndexLine{$\mathcal O_X$-môđun coherent}{(0.5.3.1)}
\EGATermIndexLine{$\mathcal O_X$-đại số coherent}{(0.5.3.6)}
\EGATermIndexLine{Đầy đủ hóa một $\mathcal O_X$-môđun và một đồng cấu các $\mathcal O_X$-môđun dọc theo một tập con đóng}{(I, \hyperref[I.10.8.4]{10.8.4})}
\EGATermIndexLine{Đầy đủ hóa một tiền lược đồ dọc theo một tập con đóng}{(I, \hyperref[I.10.8.5]{10.8.5})}
\EGATermIndexLine{Thành phần bất khả quy}{(0.2.1.5)}
\EGATermIndexLine{Hợp thành của một $\psi$-cấu xạ và một $\psi'$-cấu xạ}{(0.3.5.2)}
\EGATermIndexLine{Hợp thành của một $\Psi$-cấu xạ và một $\Psi'$-cấu xạ}{(0.4.4.2)}
\EGATermIndexLine{Điều kiện dán}{(0.3.3.1 và 0.4.1.6)}
\EGATermIndexLine{Trường cơ sở của một tiền lược đồ đại số}{(I, \hyperref[I.6.4.1]{6.4.1})}
\EGATermIndexLine{Trường thặng dư của một điểm hình học}{(I, \hyperref[I.3.4.5]{3.4.5})}

\EGATermIndexLine{Đường chéo của $X\times_S X$}{(I, \hyperref[I.5.3.9]{5.3.9})}
\EGATermIndexLine{Di-đồng cấu}{(0.1.0.2)}
\EGATermIndexLine{Miền xác định của một ánh xạ hữu tỉ}{(I, \hyperref[I.7.2.1]{7.2.1})}
\EGATermIndexLine{Đối ngẫu của một $\mathcal O_X$-môđun}{(0.4.1.4)}

\EGATermIndexLine{Phần tử lũy linh tôpô}{(0.7.1.1)}
\EGATermIndexLine{$\mathcal O_X$-môđun sinh bởi một họ nhát cắt}{(0.5.1.2)}
\EGATermIndexLine{Tập không của một nhát cắt}{(0.5.5.1)}
\EGATermIndexLine{Đại số nguyên; đại số nguyên hữu hạn}{(0.1.0.5)}
\EGATermIndexLine{Không gian vành}{(0.4.1.1)}
\EGATermIndexLine{Không gian vành địa phương}{(0.5.5.1)}
\EGATermIndexLine{Không gian vành cảm sinh trên một tập con mở}{(0.4.1.2)}
\EGATermIndexLine{Không gian vành chuẩn, rút gọn hoặc chính quy}{(0.4.1.3)}
\EGATermIndexLine{Không gian vành thu được bằng phép dán}{(0.4.1.6)}
\EGATermIndexLine{Không gian Kolmogorov}{(0.2.1.2)}
\EGATermIndexLine{Không gian bất khả quy}{(0.2.1.1)}
\EGATermIndexLine{Không gian Noether}{(0.2.2.1)}
\EGATermIndexLine{Không gian quasi-compact}{(0.2.2.4)}
\EGATermIndexLine{Không gian nền của một không gian vành}{(0.4.1.1)}
\EGATermIndexLine{Không gian vành tôpô}{(0.4.1.1)}

\EGATermIndexLine{Bó đại số trên một không gian vành}{(0.4.1.3)}
\EGATermIndexLine{Bó liên kết với một tiền bó}{(0.3.5.7)}
\EGATermIndexLine{Bó liên kết với một $A$-môđun trên $\operatorname{Spec}(A)$}{(I, \hyperref[I.1.3.4]{1.3.4})}
\EGATermIndexLine{Bó nhận giá trị trong một phạm trù}{(0.3.1.1)}
\EGATermIndexLine{Bó coherent các vành}{(0.5.3.5)}
\EGATermIndexLine{Bó vành chuẩn tại một điểm, chuẩn, rút gọn tại một điểm, rút gọn, chính quy tại một điểm hoặc chính quy}{(0.4.1.3)}
\EGATermIndexLine{Bó phân bậc các vành}{(0.4.1.3)}
\EGATermIndexLine{Bó các hàm hữu tỉ}{(I, \hyperref[I.7.3.1]{7.3.1})}
\EGATermIndexLine{Bó xoắn}{(I, \hyperref[I.7.4.1]{7.4.1})}
\EGATermIndexLine{Bó các iđêan}{(0.4.1.3)}
\EGATermIndexLine{Bó các iđêan định nghĩa}{(I, \hyperref[I.10.3.3]{10.3.3} và \hyperref[I.10.5.1]{10.5.1})}
\EGATermIndexLine{Bó cảm sinh}{(0.3.7.1)}
\EGATermIndexLine{Bó đơn địa phương}{(0.3.6.1)}
\EGATermIndexLine{Bó thu được bằng phép dán}{(0.3.3.1)}
\EGATermIndexLine{Bó giả rời rạc}{(0.3.8.1)}
\EGATermIndexLine{Bó đơn}{(0.3.6.1)}
\EGATermIndexLine{Bó cấu trúc của một không gian vành}{(0.4.1.1)}
\EGATermIndexLine{Bó cấu trúc của một lược đồ affine}{(I, \hyperref[I.1.3.4]{1.3.4})}
\EGATermIndexLine{Bó trên một cơ sở các tập con mở}{(0.3.2.2)}
\EGATermIndexLine{Cuống của một bó}{(0.3.1.6)}
\EGATermIndexLine{Thớ của một cấu xạ tiền lược đồ}{(I, \hyperref[I.3.6.3]{3.6.3})}
\EGATermIndexLine{Đại số hữu hạn}{(0.1.0.5)}
\EGATermIndexLine{Hàm hữu tỉ}{(I, \hyperref[I.7.1.2]{7.1.2})}

\EGATermIndexLine{Tổng quát hóa của một điểm}{(0.2.1.2)}
\EGATermIndexLine{$\mathcal O_X$-môđun phân bậc}{(0.4.1.3)}
\EGATermIndexLine{Đồ thị của một cấu xạ}{(I, \hyperref[I.5.3.11]{5.3.11})}

\EGATermIndexLine{Đồng cấu liên tục của các bó vành tôpô}{(0.3.1.4)}
\EGATermIndexLine{Đồng cấu xác định bởi một nhát cắt}{(0.5.1.1)}
\EGATermIndexLine{Đồng cấu địa phương của các vành địa phương}{(0.1.0.7)}
\EGATermIndexLine{$\phi$-đồng cấu các môđun}{(0.1.0.2)}

\EGATermIndexLine{Iđêan định nghĩa của một vành khả chấp}{(0.7.1.2)}
\EGATermIndexLine{$\mathcal O_X$-iđêan}{(0.4.1.3)}
\EGATermIndexLine{$\mathcal O_{\mathfrak X}$-iđêan định nghĩa của một tiền lược đồ hình thức $\mathfrak X$}{(I, \hyperref[I.10.3.3]{10.3.3} và \hyperref[I.10.5.1]{10.5.1})}
\EGATermIndexLine{Iđêan nguyên tố}{(0.1.0.6)}
\EGATermIndexLine{Ảnh trực tiếp của một $\mathcal O_X$-môđun}{(0.4.2.1)}
\EGATermIndexLine{Ảnh trực tiếp của một tiền bó}{(0.3.4.1)}
\EGATermIndexLine{Ảnh của một $S$-nhát cắt}{(I, \hyperref[I.5.3.11]{5.3.11})}
\EGATermIndexLine{Ảnh đóng của một tiền lược đồ qua một cấu xạ}{(I, \hyperref[I.9.5.3]{9.5.3})}
\EGATermIndexLine{Ảnh ngược của một $\mathcal O_Y$-môđun}{(0.4.3.1)}
\EGATermIndexLine{Ảnh ngược của một $S$-cấu xạ}{(I, \hyperref[I.3.3.7]{3.3.7})}
\EGATermIndexLine{Ảnh ngược của một tiền bó}{(0.3.5.3)}
\EGATermIndexLine{Ảnh ngược của một $S$-tiền lược đồ}{(I, \hyperref[I.3.3.6]{3.3.6})}
\EGATermIndexLine{Ảnh ngược của một tiền lược đồ con}{(I, \hyperref[I.4.4.1]{4.4.1})}
\EGATermIndexLine{Phép nhúng, phép nhúng đóng hoặc phép nhúng mở của các tiền lược đồ}{(I, \hyperref[I.4.2.1]{4.2.1})}
\EGATermIndexLine{Phép nhúng đóng của các tiền lược đồ hình thức}{(I, \hyperref[I.10.14.2]{10.14.2})}
\EGATermIndexLine{Phép nhúng địa phương của các tiền lược đồ}{(I, \hyperref[I.4.5.1]{4.5.1})}
\EGATermIndexLine{Phép nhúng chính tắc của một không gian vành cảm sinh trên một tập con mở}{(0.4.1.2)}
\EGATermIndexLine{Phép nhúng chính tắc của một tiền lược đồ con}{(I, \hyperref[I.4.1.7]{4.1.7})}
\EGATermIndexLine{Phép nhúng chính tắc của một tiền lược đồ con hình thức}{(I, \hyperref[I.10.14.2]{10.14.2})}
\EGATermIndexLine{Phép nhúng hình học}{(I, \hyperref[I.3.5.4]{3.5.4})}
\EGATermIndexLine{Nghịch đảo của một $\mathcal O_X$-môđun khả nghịch}{(0.5.4.3)}
\EGATermIndexLine{$\mathcal O_X$-môđun khả nghịch}{(0.5.4.1)}
\EGATermIndexLine{Đẳng cấu liên kết với một phép nhúng}{(I, \hyperref[I.4.2.1]{4.2.1})}
\EGATermIndexLine{Đẳng cấu địa phương của các tiền lược đồ}{(I, \hyperref[I.4.5.1]{4.5.1})}

\EGATermIndexLine{Giới hạn xạ ảnh của các $\mathcal O_{X_n}$-môđun}{(I, \hyperref[I.10.6.6]{10.6.6})}
\EGATermIndexLine{$\mathcal O_X$-môđun tự do địa phương}{(0.5.4.1)}
\EGATermIndexLine{Tính địa phương của một điểm nhận giá trị trong một vành địa phương}{(I, \hyperref[I.3.4.5]{3.4.5})}

\EGATermIndexLine{Môđun phân thức}{(0.1.2.2)}
\EGATermIndexLine{Môđun phẳng trung thành}{(0.6.4.1)}
\EGATermIndexLine{Môđun phẳng}{(0.6.1.1)}
\EGATermIndexLine{Môđun $A$-phẳng}{(0.6.2)}
\EGATermIndexLine{Môđun quasi-finite}{(0.7.4.1)}
\EGATermIndexLine{$\mathcal O_X$-môđun}{(0.4.1.3)}
\EGATermIndexLine{Cấu xạ của các không gian vành hoặc các không gian vành tôpô}{(0.4.1.1)}
\EGATermIndexLine{Cấu xạ phẳng trung thành}{(0.6.7.8)}
\EGATermIndexLine{Cấu xạ phẳng}{(0.6.7.1)}
\EGATermIndexLine{Cấu xạ của các tiền bó xác định trên một cơ sở các tập con mở}{(0.3.2.3)}
\EGATermIndexLine{Cấu xạ của các tiền lược đồ}{(I, \hyperref[I.2.2.1]{2.2.1})}
\EGATermIndexLine{Cấu xạ của các $S$-tiền lược đồ hoặc các $A$-tiền lược đồ}{(I, \hyperref[I.2.5.2]{2.5.2})}
\EGATermIndexLine{Cấu xạ song hữu tỉ}{(I, 6.5.6)}
\EGATermIndexLine{Cấu xạ kiểu hữu hạn}{(I, \hyperref[I.6.3.1]{6.3.1})}
\EGATermIndexLine{Cấu xạ đường chéo}{(I, \hyperref[I.5.3.1]{5.3.1})}
\EGATermIndexLine{Cấu xạ trội}{(I, \hyperref[I.2.2.6]{2.2.6})}
\EGATermIndexLine{Cấu xạ đóng}{(I, \hyperref[I.2.2.6]{2.2.6})}
\EGATermIndexLine{Cấu xạ đồ thị của một cấu xạ}{(I, \hyperref[I.3.3.14]{3.3.14})}
\EGATermIndexLine{Cấu xạ kiểu hữu hạn địa phương}{(I, \hyperref[I.6.6.2]{6.6.2})}
\EGATermIndexLine{Cấu xạ bị một cấu xạ khác chi phối}{(I, \hyperref[I.4.1.8]{4.1.8})}
\EGATermIndexLine{Cấu xạ mở}{(I, \hyperref[I.2.2.6]{2.2.6})}
\EGATermIndexLine{Cấu xạ quasi-compact}{(I, \hyperref[I.6.6.1]{6.6.1})}
\EGATermIndexLine{Cấu xạ radicial}{(I, \hyperref[I.3.5.4]{3.5.4})}
\EGATermIndexLine{Cấu xạ rút gọn}{(I, \hyperref[I.5.1.5]{5.1.5})}
\EGATermIndexLine{Cấu xạ tách}{(I, \hyperref[I.5.4.1]{5.4.1})}
\EGATermIndexLine{Cấu xạ cấu trúc của một $S$-tiền lược đồ}{(I, \hyperref[I.2.5.1]{2.5.1})}
\EGATermIndexLine{Cấu xạ toàn ánh}{(I, \hyperref[I.2.2.6]{2.2.6})}
\EGATermIndexLine{Cấu xạ đơn ánh phổ dụng}{(I, \hyperref[I.3.5.4]{3.5.4})}
\EGATermIndexLine{Các cấu xạ tương đương}{(I, \hyperref[I.7.1.1]{7.1.1})}
\EGATermIndexLine{$A$-cấu xạ hoặc $S$-cấu xạ của các tiền lược đồ}{(I, \hyperref[I.2.5.2]{2.5.2})}
\EGATermIndexLine{Cấu xạ của các tiền lược đồ hình thức}{(I, \hyperref[I.10.4.5]{10.4.5})}
\EGATermIndexLine{Cấu xạ adic của các tiền lược đồ hình thức}{(I, \hyperref[I.10.12.1]{10.12.1})}
\EGATermIndexLine{Cấu xạ đường chéo của các tiền lược đồ hình thức}{(I, \hyperref[I.10.15.1]{10.15.1})}
\EGATermIndexLine{Cấu xạ kiểu hữu hạn của các tiền lược đồ hình thức}{(I, \hyperref[I.10.13.3]{10.13.3})}
\EGATermIndexLine{Cấu xạ tách của các tiền lược đồ hình thức}{(I, \hyperref[I.10.15.1]{10.15.1})}
\EGATermIndexLine{Cấu xạ cấu trúc của một $\mathfrak S$-tiền lược đồ hình thức}{(I, \hyperref[I.10.4.7]{10.4.7})}
\EGATermIndexLine{$A$-cấu xạ hoặc $\mathfrak S$-cấu xạ của các tiền lược đồ hình thức}{(I, \hyperref[I.10.4.7]{10.4.7})}
\EGATermIndexLine{$\psi$-cấu xạ của các tiền bó}{(0.3.5.1)}
\EGATermIndexLine{$\Psi$-cấu xạ từ một $\mathcal O_Y$-môđun đến một $\mathcal O_X$-môđun}{(0.4.4.1)}

\EGATermIndexLine{Căn lũy linh của một vành}{(0.1.1.1)}
\EGATermIndexLine{Căn lũy linh của một $\mathcal O_X$-đại số}{(I, \hyperref[I.5.1.1]{5.1.1})}
\EGATermIndexLine{Số điểm hình học của một tiền lược đồ}{(I, \hyperref[I.6.4.8]{6.4.8})}

\EGATermIndexLine{Tập con mở affine}{(I, \hyperref[I.2.1.1]{2.1.1})}
\EGATermIndexLine{Tập con mở affine hình thức; tập con mở affine hình thức adic; tập con mở affine hình thức Noether}{(I, \hyperref[I.10.4.1]{10.4.1})}

\EGATermIndexLine{Tập con nhân của một vành}{(0.1.2.1)}
\EGATermIndexLine{Tập con nhân bão hòa}{(0.1.4.3)}
\EGATermIndexLine{Môđun $f$-phẳng trên $\mathcal O_X$}{(0.6.7.1)}
\EGATermIndexLine{Điểm của một tiền lược đồ nhận giá trị trong một vành}{(I, \hyperref[I.3.4.4]{3.4.4})}
\EGATermIndexLine{Điểm của một tiền lược đồ nhận giá trị trong một tiền lược đồ}{(I, \hyperref[I.3.4.1]{3.4.1})}
\EGATermIndexLine{Điểm của một $A$-tiền lược đồ nhận giá trị trong một $A$-đại số}{(I, \hyperref[I.3.4.4]{3.4.4})}
\EGATermIndexLine{Điểm của một $S$-tiền lược đồ nằm trên một điểm $s\in S$}{(I, \hyperref[I.2.5.1]{2.5.1})}
\EGATermIndexLine{Điểm của một $S$-tiền lược đồ nhận giá trị trong một $S$-tiền lược đồ}{(I, \hyperref[I.3.4.3]{3.4.3})}
\EGATermIndexLine{Điểm đóng}{(0.2.2.6)}
\EGATermIndexLine{Điểm tổng quát}{(0.2.1.2)}
\EGATermIndexLine{Điểm hình học của một tiền lược đồ}{(I, \hyperref[I.3.4.5]{3.4.5})}
\EGATermIndexLine{Điểm hình học nằm trên $s$}{(I, \hyperref[I.3.4.5]{3.4.5})}
\EGATermIndexLine{Điểm hình học địa phương hóa tại $x$}{(I, \hyperref[I.3.4.5]{3.4.5})}
\EGATermIndexLine{Điểm $K$-hữu tỉ}{(I, \hyperref[I.3.4.5]{3.4.5})}
\EGATermIndexLine{Tiền bó hằng}{(0.3.6.1)}
\EGATermIndexLine{Tiền bó trên một cơ sở các tập con mở}{(0.3.2.1)}
\EGATermIndexLine{Tiền lược đồ}{(I, \hyperref[I.2.1.2]{2.1.2})}
\EGATermIndexLine{Tiền lược đồ Artin}{(I, \hyperref[I.6.2.1]{6.2.1})}
\EGATermIndexLine{Tiền lược đồ liên thông}{(I, \hyperref[I.2.1.7]{2.1.7})}
\EGATermIndexLine{Tiền lược đồ cơ sở}{(I, \hyperref[I.2.5.1]{2.5.1})}
\EGATermIndexLine{Tiền lược đồ thu được bằng phép rút gọn modulo $\mathfrak J$}{(I, \hyperref[I.3.7.1]{3.7.1})}
\EGATermIndexLine{Tiền lược đồ cảm sinh trên một tập con mở}{(I, \hyperref[I.2.1.8]{2.1.8})}
\EGATermIndexLine{Tiền lược đồ nguyên}{(I, \hyperref[I.2.1.7]{2.1.7})}
\EGATermIndexLine{Tiền lược đồ bất khả quy}{(I, \hyperref[I.2.1.7]{2.1.7})}
\EGATermIndexLine{Tiền lược đồ nguyên địa phương}{(I, \hyperref[I.2.1.7]{2.1.7})}
\EGATermIndexLine{Tiền lược đồ Noether địa phương}{(I, \hyperref[I.6.1.1]{6.1.1})}
\EGATermIndexLine{Tiền lược đồ rút gọn liên kết với một tiền lược đồ}{(I, \hyperref[I.5.1.3]{5.1.3})}
\EGATermIndexLine{$A$-tiền lược đồ; tiền lược đồ trên $A$}{(I, \hyperref[I.2.5.1]{2.5.1})}
\EGATermIndexLine{$K$-tiền lược đồ đại số}{(I, \hyperref[I.6.4.5]{6.4.5})}
\EGATermIndexLine{$S$-tiền lược đồ; tiền lược đồ trên $S$}{(I, \hyperref[I.2.5.1]{2.5.1})}
\EGATermIndexLine{$S$-tiền lược đồ trội}{(I, \hyperref[I.2.5.1]{2.5.1})}
\EGATermIndexLine{Tiền lược đồ kiểu hữu hạn trên $S$; $S$-tiền lược đồ kiểu hữu hạn}{(I, \hyperref[I.6.3.1]{6.3.1})}
\EGATermIndexLine{Tiền lược đồ thu được bằng mở rộng tiền lược đồ cơ sở}{(I, \hyperref[I.3.3.6]{3.3.6})}
\EGATermIndexLine{Tiền lược đồ tách trên $S$}{(I, \hyperref[I.5.4.1]{5.4.1})}
\EGATermIndexLine{Tiền lược đồ hình thức}{(I, \hyperref[I.10.4.2]{10.4.2})}
\EGATermIndexLine{Tiền lược đồ hình thức adic}{(I, \hyperref[I.10.4.2]{10.4.2})}
\EGATermIndexLine{Tiền lược đồ hình thức Noether địa phương}{(I, \hyperref[I.10.4.2]{10.4.2})}
\EGATermIndexLine{Tiền lược đồ hình thức Noether}{(I, \hyperref[I.10.4.2]{10.4.2})}
\EGATermIndexLine{Tiền lược đồ hình thức trên $A$; $A$-tiền lược đồ hình thức}{(I, \hyperref[I.10.4.7]{10.4.7})}
\EGATermIndexLine{Tiền lược đồ hình thức trên $\mathfrak S$; $\mathfrak S$-tiền lược đồ hình thức}{(I, \hyperref[I.10.4.7]{10.4.7})}
\EGATermIndexLine{Tiền lược đồ hình thức kiểu hữu hạn trên $\mathfrak S$; $\mathfrak S$-tiền lược đồ hình thức kiểu hữu hạn}{(I, \hyperref[I.10.13.3]{10.13.3})}
\EGATermIndexLine{Tiền lược đồ hình thức tách trên $S$}{(I, \hyperref[I.10.15.1]{10.15.1})}
\EGATermIndexLine{Trình bày hữu hạn (môđun có một trình bày như vậy)}{(0.1.0.5)}
\EGATermIndexLine{Trình bày hữu hạn ($\mathcal O_X$-môđun có một trình bày như vậy)}{(0.5.2.5)}
\EGATermIndexLine{Nguyên lý quy nạp Noether}{(0.2.2.2)}
\EGATermIndexLine{Tích của các $S$-tiền lược đồ}{(I, \hyperref[I.3.2.1]{3.2.1})}
\EGATermIndexLine{Tích của các $\mathfrak S$-tiền lược đồ hình thức}{(I, \hyperref[I.10.7.1]{10.7.1})}
\EGATermIndexLine{Tích thớ của các tập hợp}{(I, \hyperref[I.3.4.2]{3.4.2})}
\EGATermIndexLine{Tích tenxơ của các bó trên các tiền lược đồ khác nhau}{(I, \hyperref[I.9.1.2]{9.1.2})}
\EGATermIndexLine{Tích tenxơ đầy đủ hóa của các đại số}{(0.7.7.5)}
\EGATermIndexLine{Tích tenxơ đầy đủ hóa của các môđun}{(0.7.7.2)}
\EGATermIndexLine{Các phép chiếu chính tắc của một tích}{(I, \hyperref[I.3.2.1]{3.2.1})}
\EGATermIndexLine{Mở rộng chính tắc của một môđun con}{(I, \hyperref[I.9.4.1]{9.4.1})}
\EGATermIndexLine{Mở rộng một cấu xạ tới các đầy đủ hóa}{(I, \hyperref[I.10.9.1]{10.9.1})}
\EGATermIndexLine{Lũy thừa ngoài thứ $p$ của một $\mathcal O_X$-môđun}{(0.4.1.4)}

\EGATermIndexLine{$\mathcal O_X$-môđun quasi-coherent}{(0.5.1.3)}
\EGATermIndexLine{$\mathcal O_X$-đại số quasi-coherent}{(0.5.1.3)}

\EGATermIndexLine{Căn của một iđêan}{(0.1.1.1)}
\EGATermIndexLine{Căn của một vành}{(0.1.1.2)}
\EGATermIndexLine{Hạng của một $\mathcal O_X$-môđun tự do địa phương}{(0.5.4.1)}
\EGATermIndexLine{Hạng của một $\mathcal O_X$-môđun không xoắn}{(I, \hyperref[I.7.4.2]{7.4.2})}
\EGATermIndexLine{Hạng của một $K$-lược đồ hữu hạn}{(I, \hyperref[I.6.4.5]{6.4.5})}
\EGATermIndexLine{Hạng khả ly của một $K$-lược đồ hữu hạn}{(I, \hyperref[I.6.4.8]{6.4.8})}
\EGATermIndexLine{Hạn chế của một không gian vành lên một tập con mở}{(0.4.1.2)}
\EGATermIndexLine{Hạn chế của một cấu xạ lên một tập con mở}{(0.4.1.2)}
\EGATermIndexLine{Hạn chế của một cấu xạ tiền lược đồ lên một tiền lược đồ con}{(I, \hyperref[I.4.1.7]{4.1.7})}
\EGATermIndexLine{Hạn chế của một tiền lược đồ lên một tập con mở}{(I, \hyperref[I.2.1.8]{2.1.8})}
\EGATermIndexLine{Hạn chế của một ánh xạ hữu tỉ lên một tập con mở}{(I, \hyperref[I.7.1.2]{7.1.2})}

\EGATermIndexLine{Lược đồ}{(I, \hyperref[I.5.4.1]{5.4.1})}
\EGATermIndexLine{Lược đồ affine}{(I, \hyperref[I.1.7.1]{1.7.1})}
\EGATermIndexLine{Lược đồ địa phương; lược đồ địa phương tại một điểm của một tiền lược đồ}{(I, \hyperref[I.2.4.1]{2.4.1})}
\EGATermIndexLine{$K$-lược đồ đại số}{(I, \hyperref[I.6.4.1]{6.4.1})}
\EGATermIndexLine{Lược đồ hữu hạn trên $K$; $K$-lược đồ hữu hạn}{(I, \hyperref[I.6.4.5]{6.4.5})}
\EGATermIndexLine{$S$-lược đồ}{(I, \hyperref[I.5.4.1]{5.4.1})}
\EGATermIndexLine{Lược đồ affine hình thức}{(I, \hyperref[I.10.1.2]{10.1.2})}
\EGATermIndexLine{$\mathfrak S$-lược đồ hình thức}{(I, \hyperref[I.10.15.1]{10.15.1})}
\EGATermIndexLine{Các nhát cắt của một bó}{(0.3.1.6)}
\EGATermIndexLine{$S$-nhát cắt của một $S$-tiền lược đồ}{(I, \hyperref[I.2.5.5]{2.5.5} và \hyperref[I.5.3.11]{5.3.11})}
\EGATermIndexLine{$S$-nhát cắt hữu tỉ của một $S$-tiền lược đồ}{(I, \hyperref[I.7.1.2]{7.1.2})}
\EGATermIndexLine{Nhát cắt đơn vị của $\mathcal O_X$}{(0.4.1.1)}
\EGATermIndexLine{Chuỗi lũy thừa hình thức hạn chế}{(0.7.5.1)}
\EGATermIndexLine{Tổng của các tiền lược đồ}{(I, \hyperref[I.3.1.1]{3.1.1}; chỉ mục nguồn in 3.6.1)}
\EGATermIndexLine{$\mathcal O_X$-đại số con sinh bởi một môđun con}{(0.4.1.3)}
\EGATermIndexLine{Tiền lược đồ con}{(I, \hyperref[I.4.1.3]{4.1.3})}
\EGATermIndexLine{Tiền lược đồ con liên kết với một phép nhúng}{(I, \hyperref[I.4.2.1]{4.2.1})}
\EGATermIndexLine{Tiền lược đồ con đóng}{(I, \hyperref[I.4.1.3]{4.1.3})}
\EGATermIndexLine{Tiền lược đồ con đóng được xác định bởi một bó iđêan}{(I, \hyperref[I.4.1.2]{4.1.2})}
\EGATermIndexLine{Tiền lược đồ con hình thức đóng}{(I, \hyperref[I.10.14.2]{10.14.2})}
\EGATermIndexLine{Chuyên biệt hóa của một điểm}{(0.2.1.2)}
\EGATermIndexLine{Phổ của một vành}{(I, \hyperref[I.1.1.1]{1.1.1})}
\EGATermIndexLine{Phổ hình thức của một vành khả chấp}{(I, \hyperref[I.10.1.2]{10.1.2})}
\EGATermIndexLine{Giá của một môđun}{(0.1.7.1)}
\EGATermIndexLine{Giá của một bó nhóm}{(0.3.1.6)}
\EGATermIndexLine{Hệ cơ bản các iđêan định nghĩa}{(I, \hyperref[I.10.3.6]{10.3.6} và \hyperref[I.10.5.1]{10.5.1})}
\EGATermIndexLine{Hệ quy nạp adic các tiền lược đồ}{(I, \hyperref[I.10.12.2]{10.12.2})}

\EGATermIndexLine{Tôpô $\mathfrak J$-adic; tôpô $\mathfrak J$-tiền-adic}{(0.7.1.9 và 0.7.2.3)}
\EGATermIndexLine{Tôpô phổ}{(I, \hyperref[I.1.1.2]{1.1.2})}
\EGATermIndexLine{$\mathcal O_X$-môđun khả nghịch tầm thường}{(I, \hyperref[I.2.4.8]{2.4.8})}
\EGATermIndexLine{$\mathcal O_X$-môđun kiểu hữu hạn}{(0.5.2.1)}

\EGATermIndexLine{Giá trị của một nhát cắt tại một điểm}{(0.5.5.1)}

\medskip\noindent\textit{Ghi chú thuật ngữ.} Nguồn chỉ dùng “adic”
cho các vành hoặc môđun tách và đầy đủ; cách dùng này hơi khác thông lệ.
\end{multicols}
\endgroup

\clearpage
\nocite{*}
\bibliography{the}
\bibliographystyle{amsalpha}
\end{document}
